\chapter*{How to read this book}
The LeanFun RevOLVER is a book about the mathematics behind Chebfun \cite{chebfun} and
ApproxFun \cite{approxfun-jl,approxfun}: Chebyshev and Jacobi expansions, their approximation theory, and the
banded operators of sparse spectral methods such as the ultraspherical spectral method \cite{olver-townsend}. Every result in the main chapters is
formalised in Lean on top of Mathlib.

\textbf{Reading paths.} Part~\ref{part:chebfun} covers Chebyshev approximation,
calculus, quadrature and rootfinding, starting from samples and coefficients.
Part~\ref{part:op} develops orthogonal-polynomial theory for general measures
and the classical families. Part~\ref{part:operators} uses those families in
sparse differential operators and studies convergence of finite solves.
Part~\ref{part:beyond} treats measure modification and separated bases in
several variables, including an annular Dirichlet solver. For Chebfun,
Part~\ref{part:chebfun} can be read independently. For ApproxFun, continue
through the classical bases to Part~\ref{part:operators}.

\textbf{Scope.} The book formalises formulas and convergence theorems used by
these methods. It does not verify either software package as a whole: the FFT
and convergence of the Remez algorithm, for example, are not formalised here.
Chapter~\ref{chap:solvers} states the assumptions needed for solver convergence.
The annular construction identifies its energy completion with the standard
$H^1_0$ space, giving weak Poisson solvability and convergence of the finite
spectral solves.

\textbf{Finding results.} Advanced search includes family formulas and
individual results, even those in collapsed boxes. Each chapter opens with
links to its principal results. Part~\ref{part:foundations} collects technical
definitions for reference; it can be consulted as needed while reading
Part~\ref{part:chebfun}. Integral equations are in
Part~\ref{part:integral-operators}.

\textbf{Layout.} Each mathematical statement has a collapsed \emph{Lean
Formalisation} box below it. The signature is extracted from Lean and includes
all hypotheses, including section variables and instances. Parentheses denote
explicit arguments, braces implicit arguments, and square brackets instances
such as a field structure. A hypothesis may therefore appear in the displayed
signature even if it was introduced before the declaration's source line.
Definitions include their Lean source. The book omits written proofs; the
checked Lean source is available in the \emph{Lean proof} box. Auxiliary
results appear in collapsed \emph{Supporting results} boxes. The dependency
graph records which results depend on which, with colours indicating
formalisation status.

\textbf{What is checked.} Lean checks every proof. Whether a Lean statement says
what the prose claims is a human question. Statements checked by hand carry a
person icon, in the book and in the dependency graph, where their nodes also have
bold borders. The review is invalidated automatically if the statement changes.
Statements with a robot icon are formalised and checked by Lean but not yet
human-reviewed. Results supplied by Mathlib carry no icon: they have a purple
border, as in the dependency graph, and a \emph{Mathlib} badge linking to the
Mathlib source file at the pinned revision.

\textbf{References.} Sources are cited with a locator, as in
\cite[Table 18.3.1]{dlmf}; each citation links to its entry in the References at
the end of the book.

\textbf{Conventions.}
\begin{itemize}
\item New families use natural-number indices; Mathlib's Chebyshev polynomials
  keep their integer indices.
\item Algebraic identities hold over any commutative $\Q$-algebra with
  denominators cleared: division is replaced by multiplication by the relevant
  factors. Such identities remain meaningful at degenerate parameters. Analytic
  results state their parameter ranges and nondegeneracy hypotheses explicitly.
\item Lean's Bochner integral is a total operation: for a nonintegrable function
  it takes the default value zero. An integral expression alone does not assert
  integrability; the hypotheses of each analytic theorem matter.
\item An $L^2$ element is an almost-everywhere equivalence class, rather than a
  chosen pointwise function. Lean's \texttt{HasSum} asserts convergence of a
  series to the specified value. Neither notation alone supplies a uniform or
  pointwise convergence claim.
\item \texttt{Polynomial.Sequence} describes a degree-graded polynomial family;
  orthogonality is an additional property. Results and definitions supplied by
  Mathlib carry the Mathlib badge.
\item Chebyshev expansions keep the full constant term, $f=a_0+\sum_{n\ge1}a_nT_n$.
\item Operators are linear maps; matrices are their coordinates in chosen bases.
\item Orthogonality allows free scaling; uniqueness needs a nondegenerate measure.
\end{itemize}

\textbf{Notation across chapters.} Some letters have different meanings in
different chapters:
\begin{itemize}
\item $S_N$ is the half-endpoint sampling sum in
  Chapter~\ref{chap:interpolation}; $S_Nf$ denotes a polynomial partial sum in
  Chapter~\ref{chap:jacobi-approx}. The shifted Legendre family $S_n$ occurs in
  Chapter~\ref{chap:classical}, while $S$ denotes basis conversion in the
  operator chapters.
\item $J(z)$ is the Joukowski map in Chapter~\ref{chap:analytic}; $J$ is a
  multiplication matrix in the operator and quadrature chapters. The map
  $J_\mu$ sends a polynomial to its $L^2(\mu)$ class.
\item $P_N$ can denote interpolation or a finite projection; $P_n$ denotes a
  polynomial family, with Jacobi parameters written as superscripts when
  needed. $\mathcal Q_N$ is a quadrature rule, whereas $Q_n$ in coefficient
  calculus is a primitive and $Q_{n,k}$ on the triangle is a Proriol polynomial.
\item $a_n(f)$ is a Chebyshev coefficient; the $a_k(x)$ in a differential
  equation are its polynomial coefficients. Recurrence coefficients $C_n$
  are distinct from ultraspherical polynomials $C_n^{(\lambda)}$.
\end{itemize}

\part{Reference foundations}\label{part:foundations}
\rvsummary{Technical conventions and special functions, available independently of the reading paths.}
\chapter{Integration and classical domains}\label{chap:foundations}

\rvsummary{The precise domains behind integrals, derivatives and analytic statements.}

This chapter fixes the domains and notation used in the analytic results. It is a reference: the approximation path begins in Chapter~\ref{chap:weights}.

\section{Integrable functions and weights}\label{sec:textbook-foundations-integrals}

The notation $\mathcal I(\mu)$ records when an integral exists. A weight on an open interval defines a measure on $\R$ by assigning no mass outside that interval.

\begin{definition}[Integrable real functions]
\label{def:integrable-input}
\lean{MeasureTheory.IntegrableInput}\leanok
For a measure $\mu$ on a measurable space $X$, let $\mathcal I(\mu)$ be the set of functions $f:X\to\R$ that are $\mu$-integrable.
\end{definition}

\begin{definition}[Measures from open-interval densities]
\label{def:open-interval-weight-measure}
\lean{MeasureTheory.openIntervalWeightMeasure}\leanok
For $a,b\in\R$ and a continuous $w:(a,b)\to\R_{\ge0}$, form $w\,dx$ on the open interval $(a,b)$, then push it forward to $\R$ by inclusion. Denote this measure by $\mu_w$.
\end{definition}

\begin{theorem}[Open-interval density characterisation]
\label{thm:open-interval-weight-measure}
\uses{def:open-interval-weight-measure,def:integrable-input}
\lean{MeasureTheory.integrable_openIntervalWeightMeasure_iff, MeasureTheory.integral_openIntervalWeightMeasure}\leanok
A function $f:\R\to\R$ is $\mu_w$-integrable if and only if $x\mapsto w(x)f(x)$ is Lebesgue-integrable on the open interval $(a,b)$. For $f\in\mathcal I(\mu_w)$,
\[\int_{\R} f\,d\mu_w=\int_{x\in(a,b)}w(x)f(x)\,dx.\]
Both integrals exist as finite real integrals.
\end{theorem}
\begin{proof}
\leanok
\end{proof}

\begin{supporting}[Integrability of functions in the domain]

\begin{theorem}[Integrability of functions in $\mathcal I(\mu)$]
\label{thm:integrable-input}
\uses{def:integrable-input}
\lean{MeasureTheory.IntegrableInput.integrable}\leanok
Every $f\in\mathcal I(\mu)$ is $\mu$-integrable.
\end{theorem}
\begin{proof}
\leanok
\end{proof}

\end{supporting}

\section{Derivatives and their domains}\label{sec:textbook-foundations-derivatives}

Derivatives on open sets are local. For absolutely continuous real functions, the derivative representative need only agree with the derivative almost everywhere.

\begin{definition}[Derivatives on open domains]
\label{def:open-domain-derivative}
\lean{HasDerivAtOnOpen}\leanok
Let $\mathbb K$ be a nontrivially normed field (for example $\mathbb R$ or $\mathbb C$). Let $U\subseteq\mathbb K$ be open, $f:U\to\mathbb K$, $x\in U$ and $d\in\mathbb K$.
Say that $f$ has derivative $d$ at $x$ if there exists $F:\mathbb K\to\mathbb K$
agreeing with $f$ on $U$ such that $F'(x)$ exists and equals $d$.
\end{definition}

\begin{theorem}[Open-domain derivative characterisation]
\label{thm:open-domain-derivative}\uses{def:open-domain-derivative}
\lean{hasDerivAtOnOpen_iff, HasDerivAtOnOpen.unique, HasDerivAtOnOpen.mul, HasDerivAtOnOpen.const_mul}\leanok
Let $\mathbb K$ be a nontrivially normed field (for example $\mathbb R$ or $\mathbb C$).
For every extension $F$ agreeing with $f$ on $U$, the open-domain derivative
condition is equivalent to $F'(x)=d$ with the derivative existing.
The derivative value is unique. If $f'(x)=d$ and $g'(x)=e$, then
$(fg)'(x)=d\,g(x)+f(x)e$ and $(cf)'(x)=cd$ for every $c\in\mathbb K$.
\end{theorem}
\begin{proof}
\leanok
\end{proof}

\begin{definition}[Integrable derivative representatives]
\label{def:integrable-derivative}\uses{def:integrable-input}
\lean{IntegrableDerivative}\leanok
For $F:\mathbb R\to\mathbb R$ and a measure $\mu$ on $\mathbb R$, let
$\mathcal D_\mu(F)$ consist of the integrable functions $h\in\mathcal I(\mu)$
for which $F$ has derivative $h(x)$ at $\mu$-almost every $x$.
This is an almost-everywhere derivative representation; no derivative is
asserted at the exceptional points.
\end{definition}

\begin{theorem}[Existence and uniqueness of derivative representatives]
\label{thm:integrable-derivative}\uses{def:integrable-derivative}
\lean{IntegrableDerivative.integrable,IntegrableDerivative.unique,IntegrableDerivative.nonempty_of_absolutelyContinuous}\leanok
Every member of $\mathcal D_\mu(F)$ is integrable, and any two members agree
$\mu$-almost everywhere. If $F$ is absolutely continuous on the closed interval
with endpoints $a,b$, then $\mathcal D_\mu(F)$ is nonempty for Lebesgue measure
restricted to that interval.
\end{theorem}
\begin{proof}
\leanok
\end{proof}

\begin{definition}[Iterated derivatives on open domains]
\label{def:open-domain-iterated-derivative}\uses{def:open-domain-derivative}
\lean{HasIteratedDerivOnOpen}\leanok
Let $\mathbb K$ be a nontrivially normed field (for example $\mathbb R$ or $\mathbb C$).
For open $U\subseteq\mathbb K$, write $D_U^n f=g$ when there is a chain
$f_0=f,f_1,\ldots,f_n=g$ of functions $U\to\mathbb K$ such that $f_{k+1}(x)$
is the derivative of $f_k$ at every $x\in U$, for every $k<n$.
\end{definition}

\begin{definition}[Classical iterated derivative values]
\label{def:open-domain-derivative-value}\uses{def:open-domain-iterated-derivative}
\lean{IteratedDerivInput, iteratedDerivOnOpen}\leanok
Let $\mathbb K$ be a nontrivially normed field (for example $\mathbb R$ or $\mathbb C$).
For an open set $U\subseteq\mathbb K$ and $n\ge0$, let $\mathcal D^n(U)$
consist of functions $f:U\to\mathbb K$ possessing a derivative chain of
order $n$. For $f\in\mathcal D^n(U)$ define $D_U^n f$ to be the unique
final function of that chain. No value is assigned without such a chain.
\end{definition}

\begin{definition}[Analyticity on an open domain]
\label{def:open-domain-analytic}
\lean{AnalyticOnOpen}\leanok
Let $\mathbb K$ be a nontrivially normed field (for example $\mathbb R$ or $\mathbb C$), let $U\subseteq\mathbb K$ be open,
and let $f:U\to\mathbb K$. We call $f$ analytic on $U$ if, for every
$x\in U$, there exist a function $F:\mathbb K\to\mathbb K$, coefficients
$a_n\in\mathbb K$ for $n\in\mathbb N$, and $r\in(0,\infty]$ such that:
\begin{enumerate}
\item $F(y)=f(y)$ for every $y\in U$;
\item $r\le R$, where the supremum is taken in $[0,\infty]$ and
\[
 R=\sup\left\{s\ge0:\ \exists C\in\mathbb R,\quad
       \|a_n\|s^n\le C\text{ for every }n\in\mathbb N\right\};
\]
\item for every $h\in\mathbb K$ with $\|h\|<r$, the partial sums
\[
 \sum_{n=0}^{N-1}a_nh^n
\]
converge to $F(x+h)$ as $N\to\infty$.
\end{enumerate}
\end{definition}

\begin{theorem}[Analytic functions possess derivative chains]
\label{thm:open-domain-analytic-derivatives}\uses{thm:open-domain-analytic,def:open-domain-derivative-value,thm:open-domain-iterated-derivative}
\lean{AnalyticOnOpen.exists_iteratedDeriv, IteratedDerivInput.ofAnalytic}\leanok
Let $\mathbb K$ be a complete nontrivially normed field (for example $\mathbb R$ or $\mathbb C$), and let $U\subseteq\mathbb K$ be open.
Every analytic function $f:U\to\mathbb K$ belongs to $\mathcal D^n(U)$
for every $n\ge0$.
\end{theorem}
\begin{proof}
\leanok
\end{proof}

\begin{supporting}[Derivative values, restriction and local jets]

\begin{theorem}[Iterated-derivative characterisation]
\label{thm:open-domain-iterated-derivative}\uses{def:open-domain-iterated-derivative,thm:open-domain-derivative}
\lean{hasIteratedDerivOnOpen_zero, hasIteratedDerivOnOpen_succ, HasIteratedDerivOnOpen.unique}\leanok
Let $\mathbb K$ be a nontrivially normed field (for example $\mathbb R$ or $\mathbb C$).
The relation $D_U^0f=g$ is equivalent to $f=g$.
The relation $D_U^{n+1}f=g$ is equivalent to the existence of $h$ with
$D_U^nf=h$ and $h'(x)=g(x)$ at every point of $U$.
Whenever such a derivative chain exists, its final function is unique.
\end{theorem}
\begin{proof}
\leanok
\end{proof}

\begin{theorem}[Classical derivative value characterisation]
\label{thm:open-domain-derivative-value}\uses{def:open-domain-derivative-value,thm:open-domain-iterated-derivative}
\lean{iteratedDerivOnOpen_spec, iteratedDerivOnOpen_eq_iff, iteratedDerivOnOpen_zero}\leanok
Let $\mathbb K$ be a nontrivially normed field (for example $\mathbb R$ or $\mathbb C$).
For $f\in\mathcal D^n(U)$, the function $D_U^nf$ belongs to an actual
order-$n$ derivative chain starting at $f$. A function $g:U\to\mathbb K$
equals $D_U^nf$ if and only if it is the final function of such a chain.
In particular, $D_U^0f=f$.
\end{theorem}
\begin{proof}
\leanok
\end{proof}

\begin{theorem}[Open-domain analyticity characterisation]
\label{thm:open-domain-analytic}\uses{def:open-domain-analytic}
\lean{analyticOnOpen_iff}\leanok
Let $\mathbb K$ be a nontrivially normed field (for example $\mathbb R$ or $\mathbb C$).
For every extension $F$ agreeing with $f$ on $U$, analyticity of $f$ on
its open domain is equivalent to analyticity of $F$ at every point of $U$.
Thus the condition is independent of the extension away from $U$.
\end{theorem}
\begin{proof}
\leanok
\end{proof}

\begin{theorem}[Restriction and truncation of derivative chains]
\label{thm:open-domain-derivative-restriction}\uses{thm:open-domain-derivative-value}
\lean{HasDerivAtOnOpen.restrict, HasIteratedDerivOnOpen.restrict, HasIteratedDerivOnOpen.exists_lower, IteratedDerivInput.truncate, IteratedDerivInput.truncate_val}\leanok
Let $\mathbb K$ be a nontrivially normed field (for example $\mathbb R$ or $\mathbb C$).
Derivative chains restrict to smaller open domains. An order-$n$ chain restricts to an order-$k$ chain for every $k\le n$, without changing its initial function.
\end{theorem}
\begin{proof}
\leanok
\end{proof}

\begin{definition}[Iterated derivatives at an individual point]
\label{def:open-domain-pointwise-derivative}\uses{def:open-domain-iterated-derivative}
\lean{HasIteratedDerivAtOnOpen}\leanok
Let $\mathbb K$ be a nontrivially normed field (for example $\mathbb R$ or $\mathbb C$).
Let $f:U\to\mathbb K$, with $U$ open, and $x\in U$. Its order-zero derivative
at $x$ is $f(x)$. For $n=k+1$, say that its order-$n$ derivative at $x$ is $d$
if some open neighbourhood $V\subseteq U$ of $x$ carries an order-$k$
derivative chain of $f|_V$, whose final function has derivative $d$ at $x$.
The highest derivative is required only at $x$.
\end{definition}

\begin{theorem}[Uniqueness of pointwise iterated derivatives]
\label{thm:open-domain-pointwise-derivative}\uses{def:open-domain-pointwise-derivative,thm:open-domain-derivative-restriction}
\lean{HasIteratedDerivAtOnOpen.unique, HasIteratedDerivOnOpen.hasAt}\leanok
Let $\mathbb K$ be a nontrivially normed field (for example $\mathbb R$ or $\mathbb C$).
The pointwise iterated derivative, when it exists, is unique. An order-$n$
derivative chain on $U$ gives that pointwise derivative at every $x\in U$.
\end{theorem}
\begin{proof}
\leanok
\end{proof}

\begin{definition}[Finite derivative jets]
\label{def:open-domain-derivative-jet}\uses{def:open-domain-pointwise-derivative}
\lean{DerivativeJetPoint, derivativeJet}\leanok
Let $\mathbb K$ be a nontrivially normed field (for example $\mathbb R$ or $\mathbb C$).
For $f:U\to\mathbb K$ and $m\ge0$, a jet point is $x\in U$ at which every
derivative of order $0,\ldots,m$ exists in the preceding pointwise sense.
At such a point define $J_kf(x)$ to be its unique order-$k$ derivative value.
\end{definition}

\begin{theorem}[Derivative-jet characterisation]
\label{thm:open-domain-derivative-jet}\uses{def:open-domain-derivative-jet,thm:open-domain-pointwise-derivative}
\lean{derivativeJet_spec, derivativeJet_eq_iff, DerivativeJetPoint.ofIterated, derivativeJet_ofIterated}\leanok
Let $\mathbb K$ be a nontrivially normed field (for example $\mathbb R$ or $\mathbb C$).
For each jet point and $0\le k\le m$, $J_kf(x)=d$ if and only if the
order-$k$ derivative exists at $x$ with value $d$. If $f$ admits an order-$m$ derivative chain on $U$, every $x\in U$ is a jet point of order $m$, and its truncated chain values agree
with these jet values.
\end{theorem}
\begin{proof}
\leanok
\end{proof}

\end{supporting}

\section{Circle integrals}\label{sec:textbook-foundations-circles}

The contour formulas in Chapter~\ref{chap:analytic} use the following parametrized circle integral.

\begin{definition}[Integrable functions on a circle]
\label{def:circle-integral-input}
\lean{Complex.CircleIntegralInput, Complex.circleIntegralOn, Complex.CircleIntegralInput.ofContinuous, Complex.CircleIntegralInput.add}\leanok
For $c\in\mathbb C$ and $r\in\mathbb R$, let
$S=\{z:|z-c|=|r|\}$ and $\gamma(\theta)=c+r e^{i\theta}$.
Let $\mathcal C_{c,r}$ consist of functions $f:S\to\mathbb C$ for which
$\theta\mapsto i r e^{i\theta}f(\gamma(\theta))$ is integrable on $[0,2\pi]$.
For $f\in\mathcal C_{c,r}$, define
\[
\mathcal I_{c,r}(f)=\int_0^{2\pi}i r e^{i\theta}f(\gamma(\theta))\,d\theta.
\]
Addition in $\mathcal C_{c,r}$ is pointwise addition of functions.
Reference: \cite[proof of Theorem 8.1]{atap}.
\end{definition}

\begin{theorem}[Circle-integral norm bound]
\label{thm:circle-integral-bound}\uses{def:circle-integral-input}
\lean{Complex.norm_circleIntegralOn_le}\leanok
If $f\in\mathcal C_{c,r}$ and $|f(z)|\le M$ for every $z\in S$, then
\[
|\mathcal I_{c,r}(f)|\le 2\pi |r| M.
\]
Reference: \cite[proof of Theorem 8.1]{atap}.
\end{theorem}
\begin{proof}
\leanok
\end{proof}

\begin{supporting}[Parametrization and linearity]

\begin{theorem}[Circle-integral characterisation]
\label{thm:circle-integral-input}\uses{def:circle-integral-input}
\lean{Complex.circleIntegralOn_spec, Complex.CircleIntegralInput.ofContinuous_val, Complex.circleIntegralOn_radius_zero, Complex.circleIntegralOn_add}\leanok
For $f\in\mathcal C_{c,r}$, the function
$\theta\mapsto i r e^{i\theta}f(c+r e^{i\theta})$ is integrable on $[0,2\pi]$,
and its integral equals $\mathcal I_{c,r}(f)$.
Every continuous $f:S\to\mathbb C$ belongs to $\mathcal C_{c,r}$, with its
values on $S$ unchanged. For $r=0$ the curve is constant and
$\mathcal I_{c,0}(f)=0$. For $f,g\in\mathcal C_{c,r}$,
\[
\mathcal I_{c,r}(f+g)=\mathcal I_{c,r}(f)+\mathcal I_{c,r}(g).
\]
\end{theorem}
\begin{proof}
\leanok
\end{proof}

\end{supporting}

\chapter{Gamma and hypergeometric functions}\label{chap:special-functions}

\rvsummary{Classical Gamma and convergent hypergeometric functions used by the polynomial families.}

Gamma functions give the normalization constants of the classical families. Terminating hypergeometric series give their polynomial representations.

\section{Gamma and reciprocal Gamma}\label{sec:textbook-special-functions-gamma}

Write $(a)_0=1$ and $(a)_n=a(a+1)\cdots(a+n-1)$ for the rising factorial. The Gamma recurrence relates this finite product to ratios of Gamma values.

\begin{definition}[Gamma on its classical domain]
\label{def:gamma-function}
\lean{Complex.GammaPoint, Complex.GammaOn, Real.GammaPoint, Real.GammaOn}\leanok
Let $D=\mathbb C\setminus\{0,-1,-2,\ldots\}$. For $z\in D$, let
$N=\lfloor\max(1-\Re z,0)\rfloor$ and define
\[
 \Gamma(z)=\frac{1}{(z)_N}\int_0^\infty e^{-t}t^{z+N-1}\,dt,
\]
where $t$ ranges over the strictly positive reals and the power uses the real
logarithm of $t$. The real Gamma function is its restriction to
$\mathbb R\setminus\{0,-1,-2,\ldots\}$. Gamma has no assigned value at a pole.

Reference: \cite[5.2.1, 5.5.1]{dlmf}.
\end{definition}

\begin{theorem}[Gamma characterisation on the pole-free domain]
\label{thm:gamma-function}\uses{def:gamma-function}
\lean{Complex.GammaPoint.ofRePos, Complex.GammaPoint.addNat, Complex.GammaPoint.addNat_val, Complex.GammaPoint.exists_positive_shift, Complex.GammaOn_one, Complex.GammaOn_ne_zero, Complex.GammaOn_euler, Complex.GammaOn_add_one, Complex.GammaOn_add_nat_div, Complex.GammaOn_euler_shift, Real.GammaPoint.ofPos, Real.GammaPoint.toComplex, Real.GammaOn_ne_zero, Real.GammaOn_pos, Real.GammaOn_euler, Real.GammaOn_toComplex}\leanok
Every $z\in D$ admits a natural shift into the right half-plane. For every
such shift, the displayed Euler integral is integrable, $(z)_N\ne0$, and
the formula gives the same value $\Gamma(z)$. Gamma never vanishes on $D$,
$\Gamma(1)=1$, and
\[
 \Gamma(z+1)=z\Gamma(z),\qquad
 \frac{\Gamma(z+N)}{\Gamma(z)}=(z)_N.
\]
The right half-plane and all natural shifts of $D$ lie in $D$.
For real $x>0$, the real Euler integral is integrable and equals
$\Gamma(x)>0$; the real and complex values agree on real non-poles.
\end{theorem}
\begin{proof}
\leanok
\end{proof}

\begin{definition}[Reciprocal Gamma]
\label{def:reciprocal-gamma}\uses{def:gamma-function}
\lean{Complex.reciprocalGamma}\leanok
Define $G:\mathbb C\to\mathbb C$ by $G(z)=1/\Gamma(z)$ for $z\in D$ and
$G(-n)=0$ for $n\ge0$.

Reference: \cite[5.2(i)]{dlmf}.
\end{definition}

\begin{theorem}[Reciprocal Gamma characterisation]
\label{thm:reciprocal-gamma}\uses{def:reciprocal-gamma,thm:gamma-function,def:polynomial-linear-maps}
\lean{Complex.reciprocalGamma_mul_GammaOn, Complex.reciprocalGamma_eq_zero_iff, Complex.differentiable_reciprocalGamma}\leanok
The function $G$ is complex differentiable everywhere. Its zeros are exactly
$0,-1,-2,\ldots$, and $G(z)\Gamma(z)=1$ for every $z\in D$.
\end{theorem}
\begin{proof}
\leanok
\end{proof}

\section{The Gauss series}\label{sec:textbook-special-functions-gauss}

The coefficient definition includes terminating series. Its parameter condition ensures that every nonzero term has a nonzero denominator.

\begin{definition}[Admissible Gauss coefficients]
\label{def:gauss-coefficients}
\lean{GaussParameters,ordinaryHypergeometricCoefficientOn}\leanok
Over a field of characteristic zero, call $(a,b,c)$ admissible when
$(c)_n=0$ implies $(a)_n(b)_n=0$ for every $n\ge0$. Its classical Gauss
coefficient is zero after numerator termination; otherwise it is
\[
 u_n=\frac{(a)_n(b)_n}{(c)_n n!}.
\]
 In a terminating case,
termination is taken before any later denominator pole; this includes coincident
terminating parameters.

Reference: \cite[15.2.1, 15.2.4--15.2.6]{dlmf}.
\end{definition}

\begin{definition}[Ordinary hypergeometric function]
\label{def:ordinary-hypergeometric}\uses{def:gauss-coefficients}
\lean{GaussSeriesPoint,ordinaryHypergeometricOn}\leanok
For admissible real or complex parameters, the classical Gauss series is
\[{}_2F_1(a,b;c;z)=\sum_{n=0}^{\infty}u_n z^n.\]
Its argument domain consists exactly of the points where the partial sums in
increasing natural index converge. Its value is their unique limit.  The domain contains only arguments where these partial sums converge.

Reference: \cite[15.2.1, 15.2.4--15.2.5]{dlmf}.
\end{definition}

\begin{theorem}[Characterisation of classical Gauss values]
\label{thm:ordinary-hypergeometric}\uses{def:ordinary-hypergeometric,thm:gauss-coefficients}
\lean{tendsto_ordinaryHypergeometricOn,ordinaryHypergeometricOn_eq_of_tendsto,GaussSeriesPoint.ofHasSum,ordinaryHypergeometricOn_ofHasSum,ordinaryHypergeometricCoefficientOn_eq_zero_of_termination,GaussSeriesPoint.ofTerminating,ordinaryHypergeometricOn_terminating,GaussParameters.toComplex,GaussParameters.coefficient_toComplex,GaussSeriesPoint.toComplex,Complex.ofReal_ordinaryHypergeometric}\leanok
The natural partial sums converge to the stated value, and every other limit of
those sums agrees with it. An unconditional sum supplies this value. If either
numerator is $-k$, all coefficients after degree $k$ vanish: every argument is in
the domain and the value equals $\sum_{n=0}^{k}u_nz^n$, allowing earlier termination.
Real admissible parameters and convergent real arguments embed into the complex
domains, preserving coefficients and values, including conditionally convergent
and terminating series.
\end{theorem}
\begin{proof}
\leanok
\end{proof}

\begin{theorem}[Gauss series calculus and differential equation]
\label{thm:hypergeometric-calculus}\uses{thm:ordinary-hypergeometric,def:open-domain-analytic,def:open-domain-iterated-derivative,thm:open-domain-derivative}
\lean{GaussSeriesPoint.ofUnitDisk,hasSum_ordinaryHypergeometricOn,one_le_ordinaryHypergeometricOn_radius,hasFPowerSeriesOnBall_ordinaryHypergeometricOn,analyticOnOpen_ordinaryHypergeometricOn,hasDerivAtOnOpen_ordinaryHypergeometricOn,hasDerivAtOnOpen_ordinaryHypergeometricOn_constant,hasIteratedDerivOnOpen_ordinaryHypergeometricOn_two,hasIteratedDerivOnOpen_ordinaryHypergeometricOn_two_constant,hasIteratedDerivOnOpen_ordinaryHypergeometricOn_two_linear,ordinaryHypergeometricCoefficientOn_recurrence,eq_ordinaryHypergeometricCoefficientOn_of_recurrence,ordinaryHypergeometricOn_ode}\leanok
For admissible real or complex parameters, the classical Gauss series has radius
at least one and converges on $|z|<1$, where it is analytic with its defining
power-series representation. Write $F(z)={}_2F_1(a,b;c;z)$.
If $ab=0$, its first and second derivatives vanish. Otherwise $c\ne0$, the
shifted parameters are admissible, and
\[
 F'(z)=\frac{ab}{c}{}_2F_1(a+1,b+1;c+1;z).
\]
In this nonconstant case, if $(a+1)(b+1)=0$, the second derivative vanishes.
Otherwise the twice-shifted
parameters are admissible, $c+1\ne0$, and
\[
 F''(z)=\frac{ab}{c}\frac{(a+1)(b+1)}{c+1}
 {}_2F_1(a+2,b+2;c+2;z).
\]
All derivatives here exist on the open disk. When $c+j\ne0$ for every natural
$j$, the normalized coefficient recurrence
\[
 (n+1)(c+n)u_{n+1}=(a+n)(b+n)u_n,\qquad u_0=1
\]
uniquely determines the coefficients, and the genuine derivatives satisfy
\[
 z(1-z)F''+(c-(a+b+1)z)F'-abF=0.
\]
Terminating cases are included.
Reference: \cite[15.2.1, 15.5.1, 15.10.1]{dlmf}.
\end{theorem}
\begin{proof}
\leanok
\end{proof}

\begin{supporting}[Admissible parameters and coefficient recurrences]

\begin{theorem}[Characterisation of admissible Gauss coefficients]
\label{thm:gauss-coefficients}\uses{def:gauss-coefficients}
\lean{GaussParameters.denominator_ne_zero,ordinaryHypergeometricCoefficientOn_eq_zero,ordinaryHypergeometricCoefficientOn_eq_div,GaussParameters.ofNoPole,GaussParameters.ofTerminatingRight,GaussParameters.swap,GaussParameters.ofTerminatingLeft,GaussParameters.c_ne_zero,GaussParameters.shift,ordinaryHypergeometricCoefficientOn_succ_of_constant,ordinaryHypergeometricCoefficientOn_shift,GaussParameters.pochhammer_ne_zero_of_le,GaussParameters.ofTerminating,GaussParameters.quadratic_terminates,GaussParameters.quadratic,GaussParameters.quadratic_has_termination,GaussParameters.pochhammer_neg_nat_ne_zero}\leanok
A denominator parameter outside $\{0,-1,-2,\ldots\}$ always gives admissible
parameters. If a numerator is $-k$, nonvanishing denominators through degree $k$
suffice; exchanging numerator parameters preserves admissibility. Nonzero
numerator coefficients have nonzero denominators and are the displayed quotients.
If $ab=0$, all positive-degree coefficients vanish. If $ab\ne0$, then $c\ne0$,
$(a+1,b+1,c+1)$ is admissible, and
\[
 (n+1)u_{n+1}(a,b;c)=\frac{ab}{c}u_n(a+1,b+1;c+1).
\]

Nonvanishing of $(c)_k$ implies nonvanishing of $(c)_j$ for $j\le k$, so a
single finite Pochhammer condition certifies a numerator terminating at $k$.
The pair $(-n/2,(1-n)/2)$ has one numerator equal to $-\lfloor n/2\rfloor$;
hence $(c)_{\lfloor n/2\rfloor}\ne0$ gives admissible quadratic data.
For natural $m,k$ with $k\le m$, $(-m)_k\ne0$.
\end{theorem}
\begin{proof}
\leanok
\end{proof}

\end{supporting}

\section{Regularized series}\label{sec:textbook-special-functions-regularized}

Reciprocal Gamma gives a coefficient formula that is defined even at denominator parameters where the ordinary Gauss series is unavailable.

\begin{definition}[Regularized hypergeometric coefficients]
\label{def:regularized-hypergeometric-coefficients}\uses{def:reciprocal-gamma}
\lean{Complex.regularizedHGFunCoeffOn}\leanok
For finite multisets $A,B$ of complex parameters, define
\[
 v_n(A;B)=\frac{\prod_{a\in A}(a)_n\prod_{b\in B}G(b+n)}{n!},
\]
where $G$ is the reciprocal Gamma function of
Definition~\ref{def:reciprocal-gamma}.

Reference: \cite[15.2.2, 16.2]{dlmf}.
\end{definition}

\begin{definition}[Regularized generalized hypergeometric function]
\label{def:regularized-hypergeometric}\uses{def:regularized-hypergeometric-coefficients}
\lean{Complex.RegularizedHGSeriesPoint,Complex.regularizedHGFunOn,Complex.regularizedGaussHGFunOn}\leanok
The regularized generalized hypergeometric value is the limit of
$\sum_{n=0}^{N-1}v_n(A;B)z^n$ as $N\to\infty$, on exactly the domain where this
limit exists. Its Gauss specialization $\mathbf F(a,b;c;z)$ takes
$A=\{a,b\}$ and $B=\{c\}$. The domain includes conditional convergence and excludes divergent series.
Reference: \cite[15.2.2, 16.2]{dlmf}.
\end{definition}

\begin{theorem}[Regularized and ordinary Gauss functions]
\label{thm:regularized-gauss}\uses{thm:regularized-hypergeometric,thm:ordinary-hypergeometric,thm:gamma-function,thm:regularized-hypergeometric-coefficients}
\lean{Complex.regularizedHGFunCoeffOn_eq_gauss_div_Gamma,Complex.RegularizedHGSeriesPoint.ofGauss,Complex.ordinaryHypergeometricOn_div_Gamma_eq}\leanok
If $c$ is not a nonpositive integer and the classical ordinary Gauss series
converges at $z$, then the regularized series also converges there and
$\mathbf F(a,b;c;z)={}_2F_1(a,b;c;z)/\Gamma(c)$.
The same normalization identity holds coefficient by coefficient.
Reference: \cite[15.2.2]{dlmf}.
\end{theorem}
\begin{proof}
\leanok
\end{proof}

\begin{supporting}[Regularized coefficients and convergence]

\begin{theorem}[Regularized Gauss coefficients]
\label{thm:regularized-hypergeometric-coefficients}\uses{def:regularized-hypergeometric-coefficients}
\lean{Complex.regularizedHGFunCoeffOn_gauss}\leanok
For $A=\{a,b\}$ and $B=\{c\}$,
$ v_n=(a)_n(b)_nG(c+n)/n!$.
\end{theorem}
\begin{proof}
\leanok
\end{proof}

\begin{theorem}[Characterisation of regularized hypergeometric values]
\label{thm:regularized-hypergeometric}\uses{def:regularized-hypergeometric}
\lean{Complex.tendsto_regularizedHGFunOn,Complex.regularizedHGFunOn_eq_of_tendsto,Complex.RegularizedHGSeriesPoint.ofHasSum,Complex.regularizedHGFunOn_ofHasSum}\leanok
The natural partial sums converge to the regularized value, uniquely. Every
unconditional sum is also this natural-order sum and supplies a point of the domain.
\end{theorem}
\begin{proof}
\leanok
\end{proof}

\end{supporting}

\part{Chebfun on an interval}\label{part:chebfun}
\rvsummary{Chebyshev expansions on $[-1,1]$, from orthogonality to the convergence theory, interpolation, calculus, quadrature and rootfinding behind Chebfun.}

\chapter{Chebyshev weights and orthogonality}\label{chap:weights}

\rvsummary{Chebyshev polynomials, their weights and orthogonality relations.}

Chebyshev expansions represent functions on $[-1,1]$ using polynomials tied to cosine series. The two Chebyshev weights give the orthogonality relations from which coefficients are computed.

\section{The polynomials and their trigonometric form}\label{sec:textbook-weights-polynomials}

The three-term recurrence defines polynomials; the cosine formula explains their use on $[-1,1]$.

\begin{definition}[Chebyshev polynomials]
\label{def:chebyshev-mathlib}
\lean{Polynomial.Chebyshev.T, Polynomial.Chebyshev.U}\mathlibok
For a commutative ring $R$ and every integer $n$, define $T_n,U_n\in R[X]$ by $T_0=U_0=1$, $T_1=X$, $U_1=2X$, and the common recurrence
$p_{n+2}=2Xp_{n+1}-p_n$, run forwards for nonnegative and backwards for
negative indices.

Reference: \cite[Table 18.9.1]{dlmf}.
\end{definition}

\begin{theorem}[Trigonometric form]
\label{thm:chebyshev-trig}\uses{def:chebyshev-mathlib}
\lean{Polynomial.Chebyshev.T_real_cos, Polynomial.Chebyshev.U_real_cos}\mathlibok
For real $\theta$ and integer $n$, $T_n(\cos\theta)=\cos n\theta$ and
$U_n(\cos\theta)\sin\theta=\sin(n+1)\theta$.

Reference: \cite[18.5.1--18.5.2]{dlmf}.
\end{theorem}
\begin{proof}
\mathlibok
\end{proof}

\section{The first-kind weight}\label{sec:textbook-weights-first-kind}

The substitution $x=\cos\theta$ converts integration against $w_T$ into ordinary integration over $(0,\pi)$. The change-of-variable theorem below states this precisely.

\begin{definition}[Classical Chebyshev densities]
\label{def:classical-chebyshev-weights}
\lean{Polynomial.Chebyshev.weightU, Polynomial.Chebyshev.weightT}\leanok
On the open interval $I=(-1,1)$, define the continuous nonnegative weights $w_U(x)=\sqrt{1-x^2}$ and $w_T(x)=1/\sqrt{1-x^2}$.
\end{definition}

\begin{definition}[First-kind weight measure]
\label{def:measureT}\uses{def:open-interval-weight-measure,def:classical-chebyshev-weights}
\lean{Polynomial.Chebyshev.classicalMeasureT}\leanok
Let $\mu_T=\mu_{w_T}$ be the measure of Definition~\ref{def:open-interval-weight-measure}
for the weight $w_T$ of Definition~\ref{def:classical-chebyshev-weights}.

Reference: \cite[Table 18.3.1]{dlmf}.
\end{definition}

\begin{theorem}[First-kind integrability characterisation]
\label{thm:measureT-integrability}
\uses{def:measureT,thm:open-interval-weight-measure}
\lean{Polynomial.Chebyshev.integrable_classicalMeasureT_iff, Polynomial.Chebyshev.integral_classicalMeasureT}\leanok
A real function $f$ is integrable for $\mu_T$ exactly when $x\mapsto w_T(x)f(x)$ is Lebesgue-integrable on $(-1,1)$. For $f\in\mathcal I(\mu_T)$,
\[\int f\,d\mu_T=\int_{x\in(-1,1)}w_T(x)f(x)\,dx.\]
\end{theorem}
\begin{proof}
\leanok
\end{proof}

\begin{theorem}[Classical cosine transport and coefficient integrability]
\label{thm:T-integrable-transport}\uses{thm:measureT-integrability,def:chebyshev-mathlib,thm:classical-chebyshev-weights}
\lean{Polynomial.Chebyshev.integrable_classicalMeasureT, Polynomial.Chebyshev.intervalIntegrable_classicalMeasureT_cos, Polynomial.Chebyshev.integral_classicalMeasureT_eq_integral_cos, Polynomial.Chebyshev.integrable_classicalMeasureT_mul_T}\leanok
Every real function continuous on $[-1,1]$ is $\mu_T$-integrable.
For every $f\in\mathcal I(\mu_T)$, the function $\theta\mapsto f(\cos\theta)$
is integrable on $(0,\pi)$ and
\[\int f\,d\mu_T=\int_0^\pi f(\cos\theta)\,d\theta.\]
For every integer $n$, $fT_n$ is also $\mu_T$-integrable. No continuity of $f$
is required; integrable endpoint singularities are included.

Reference: \cite[Chapter 3]{atap}.
\end{theorem}
\begin{proof}
\leanok
\end{proof}

\begin{theorem}[First-kind orthogonality and norms]
\label{thm:T-orthogonal}
\uses{thm:T-integrable-transport}
\lean{Polynomial.Chebyshev.classicalMeasureT_orthogonality}
\leanok
For natural $n,m$, $T_nT_m$ is integrable. For distinct $n,m$, $\int T_nT_m\,d\mu_T=0$.
The squared norm of $T_0$ is $\pi$, and that of $T_n$ for $n>0$ is $\pi/2$.

Reference: \cite[Table 18.3.1]{dlmf}.
\end{theorem}
\begin{proof}
\leanok
\end{proof}

\begin{supporting}[Density formulas and first-kind moments]

\begin{theorem}[Chebyshev density characterisation]
\label{thm:classical-chebyshev-weights}
\uses{def:classical-chebyshev-weights}
\lean{Polynomial.Chebyshev.weightUT_apply}\leanok
For every $x\in I$, $w_U(x)=\sqrt{1-x^2}$ and $w_T(x)=1/\sqrt{1-x^2}$.
\end{theorem}
\begin{proof}
\leanok
\end{proof}

\begin{theorem}[First-kind moments]
\label{thm:T-moments}
\uses{thm:T-integrable-transport,thm:chebyshev-trig}
\lean{Polynomial.Chebyshev.classicalMeasureT_moments}
\leanok
For integer $k$, $T_k$ is integrable and $\int T_k\,d\mu_T=\pi$ if $k=0$, and is zero otherwise.

Reference: \cite[Table 18.3.1]{dlmf}.
\end{theorem}
\begin{proof}
\leanok
\end{proof}

\end{supporting}

\section{The second-kind weight}\label{sec:textbook-weights-second-kind}

The second-kind family uses $w_U$ instead. Its squared norms are the same at every nonnegative degree.

\begin{definition}[Second-kind weight measure]
\label{def:measureU}
\uses{def:open-interval-weight-measure,def:classical-chebyshev-weights}
\lean{Polynomial.Chebyshev.measureU}\leanok
Let $\mu_U=\mu_{w_U}$ be the measure of Definition~\ref{def:open-interval-weight-measure}
for the weight $w_U$ of Definition~\ref{def:classical-chebyshev-weights}.

Reference: \cite[Table 18.3.1]{dlmf}.
\end{definition}

\begin{lemma}[Weighted integral formula]
\label{lem:U-integral}
\uses{def:measureU,thm:open-interval-weight-measure}
\lean{Polynomial.Chebyshev.integral_measureU}\leanok
For $f\in\mathcal I(\mu_U)$,
\[\int f\,d\mu_U=\int_{x\in(-1,1)} f(x)\sqrt{1-x^2}\,dx.\]
The function $x\mapsto f(x)\sqrt{1-x^2}$ is Lebesgue-integrable on $(-1,1)$.
Reference: \cite[Table 18.3.1]{dlmf}.
\end{lemma}
\begin{proof}
\leanok
\end{proof}

\begin{theorem}[Second-kind orthogonality]
\label{thm:U-orthogonal}
\lean{Polynomial.Chebyshev.integral_eval_U_real_mul_eval_U_real_measureU_of_ne}
\leanok
\uses{lem:U-product}
For distinct $n,m\in\N$,
\[\int U_n(x)U_m(x)\,d\mu_U(x)=0.\]

Reference: \cite[Table 18.3.1]{dlmf}.
\end{theorem}
\begin{proof}
\leanok
\end{proof}

\begin{theorem}[Second-kind squared norms]
\label{thm:U-norm}
\lean{Polynomial.Chebyshev.integral_eval_U_real_mul_self_measureU}
\leanok
\uses{lem:U-product}
For every $n\in\N$,
\[\int U_n(x)^2\,d\mu_U(x)=\frac\pi2.\]

Reference: \cite[Table 18.3.1]{dlmf}.
\end{theorem}
\begin{proof}
\leanok
\end{proof}

\begin{supporting}[Integrability, comparison of weights and product reduction]

\begin{theorem}[Second-kind integrability characterisation]
\label{thm:measureU-integrability}
\uses{def:measureU,thm:open-interval-weight-measure}
\lean{Polynomial.Chebyshev.integrable_measureU_iff}\leanok
A real function $f$ is $\mu_U$-integrable exactly when $x\mapsto\sqrt{1-x^2}f(x)$ is Lebesgue-integrable on the open interval $(-1,1)$.
\end{theorem}
\begin{proof}
\leanok
\end{proof}

\begin{lemma}[Integrability]
\label{lem:U-integrable}
\lean{Polynomial.Chebyshev.integrable_measureU}
\leanok
\uses{thm:measureU-integrability}
Every real-valued function continuous on $[-1,1]$ is integrable for $\mu_U$.

Reference: the second-kind weight in \cite[Table 18.3.1]{dlmf}.
\end{lemma}
\begin{proof}
\leanok
\end{proof}

\begin{lemma}[Comparison of weights]
\label{lem:U-to-T}
\uses{lem:U-integral,thm:measureT-integrability,thm:measureU-integrability}
\lean{Polynomial.Chebyshev.integral_measureU_eq_classicalMeasureT}\leanok
For $f\in\mathcal I(\mu_U)$, the function $x\mapsto f(x)(1-x^2)$ is $\mu_T$-integrable and
\[\int f\,d\mu_U=\int f(x)(1-x^2)\,d\mu_T(x).\]
Reference: \cite[Table 18.3.1]{dlmf}.
\end{lemma}
\begin{proof}
\leanok
\end{proof}

\begin{lemma}[Product reduction]
\label{lem:U-product}\uses{lem:U-to-T,lem:U-integrable,thm:chebyshev-trig}
\lean{Polynomial.Chebyshev.integral_eval_U_real_mul_eval_U_real_measureU}
\leanok

For integers $n,m$,
\[\int U_n U_m\,d\mu_U=\frac12\left(\int T_{n-m}\,d\mu_T
-\int T_{n+m+2}\,d\mu_T\right).\]

Reference: \cite[Table 18.3.1 and equations 18.5.1--18.5.2]{dlmf}.
\end{lemma}
\begin{proof}
\leanok
\end{proof}

\end{supporting}

\chapter{Chebyshev series: coefficients and completeness}\label{chap:coefficients}

\rvsummary{Defines Chebyshev coefficients and proves that Chebyshev expansions are complete in the weighted $L^2$ space.}

Orthogonality gives coefficients by integration. Finite expansions are recovered exactly; completeness explains what happens for infinite expansions.

\section{Coefficients and finite expansions}\label{sec:textbook-coefficients-coefficients}

The coefficient at index $n$ requires integrability of $fT_n$. We keep the full constant term: $f=a_0+\sum_{n\ge1}a_nT_n$.

\begin{definition}[Inputs for a fixed Chebyshev coefficient]
\label{def:coefficient-input}\uses{thm:T-integrable-transport}
\lean{Polynomial.Chebyshev.CoefficientInput, Polynomial.Chebyshev.CoefficientInput.ofL1, Polynomial.Chebyshev.CoefficientInput.ofContinuousOn, Polynomial.Chebyshev.CoefficientInput.add, Polynomial.Chebyshev.CoefficientInput.constMul}\leanok
For $n\in\N$, define
\[
\mathcal I_n=\{f:\R\to\R: fT_n\text{ is }\mu_T\text{-integrable}\}.
\]
Functions in this domain retain their pointwise values. Addition and real
scalar multiplication are pointwise. The defining condition concerns $fT_n$,
not $f$ alone.
Reference: \cite[Chapter 3]{atap}.
\end{definition}

\begin{theorem}[Characterisation of coefficient inputs]
\label{thm:coefficient-input}\uses{def:coefficient-input,thm:chebyshev-recurrence}
\lean{Polynomial.Chebyshev.CoefficientInput.integrable, Polynomial.Chebyshev.CoefficientInput.ofL1_val, Polynomial.Chebyshev.CoefficientInput.ofContinuousOn_val, Polynomial.Chebyshev.coefficientInput_zero_iff, Polynomial.Chebyshev.coefficientInput_congr_ae}\leanok
For $f\in\mathcal I_n$, the product $fT_n$ is $\mu_T$-integrable.
Every $\mu_T$-integrable real function, and every $f:\R\to\R$ continuous on
$[-1,1]$, belongs to $\mathcal I_n$ for every $n$, with unchanged pointwise values.
The domain $\mathcal I_0$ consists exactly of the $\mu_T$-integrable real functions.
Membership of $\mathcal I_n$ is unchanged by $\mu_T$-almost-everywhere equality.
\end{theorem}
\begin{proof}
\leanok
\end{proof}

\begin{definition}[Chebyshev coefficients]
\label{def:coefficients}\uses{def:coefficient-input}
\lean{Polynomial.Chebyshev.coefficient}
\leanok

At each index $n$, for $f\in\mathcal I_n$, set
\[a_0(f)=\pi^{-1}\int f\,d\mu_T,\qquad
 a_n(f)=2\pi^{-1}\int fT_n\,d\mu_T\quad(n>0).\]
The expansion convention is $a_0+\sum_{n\geq1}a_nT_n$.

Reference: \cite[Chapters 3 and 8]{atap},
with the constant-term convention made explicit here.
\end{definition}

\begin{theorem}[Cosine integral formula]
\label{thm:coefficient-cos}\uses{thm:coefficient-characterisation,thm:chebyshev-trig}
\lean{Polynomial.Chebyshev.coefficient_eq_integral_cos, Polynomial.Chebyshev.coefficient_integrable_cos}
\leanok

For $n\in\N$ and $f\in\mathcal I_n$, the cosine integrand is integrable and
\[a_n(f)=\frac{c_n}{\pi}\int_0^\pi f(\cos\theta)\cos(n\theta)\,d\theta,
\qquad c_0=1,\quad c_n=2\ (n>0).\]

Reference: \cite[Chapters 3 and 8]{atap}.
\end{theorem}
\begin{proof}
\leanok
\end{proof}

\begin{theorem}[Recovery of a finite expansion]
\label{thm:coefficient-finite}\uses{thm:coefficient-T}
\lean{Polynomial.Chebyshev.coefficient_sum_eval_T}
\leanok

For a finite set $S\subset\N$ and real coefficients $b_k$, let
$p(x)=\sum_{k\in S}b_kT_k(x)$. Then $a_n(p)=b_n$ if $n\in S$, and
$a_n(p)=0$ otherwise.

Reference: \cite[Chapter 3]{atap}.
\end{theorem}
\begin{proof}
\leanok
\end{proof}

\begin{theorem}[Uniform coefficient bound]
\label{thm:coefficient-bound}
\lean{Polynomial.Chebyshev.abs_coefficient_le}
\leanok
\uses{thm:coefficient-cos}
For $f\in\mathcal I_n$, if $M\ge0$ and $|f(x)|\le M$ on $[-1,1]$,
then $|a_n(f)|\le M$ at $n=0$ and $|a_n(f)|\le2M$ at $n>0$.

Reference: the coefficient formula in \cite[Chapters 3 and 8]{atap}.
\end{theorem}
\begin{proof}
\leanok
\end{proof}

\begin{supporting}[Coefficient identities and interval dependence]

\begin{theorem}[Classical coefficient formula]
\label{thm:coefficient-characterisation}\uses{def:coefficients}
\lean{Polynomial.Chebyshev.coefficient_apply}\leanok
For $f\in\mathcal I_n$,
\[a_n(f)=\frac{c_n}{\pi}\int fT_n\,d\mu_T,
\qquad c_0=1,\quad c_n=2\quad(n>0).\]
The product $fT_n$ is integrable, so this integral is finite.
\end{theorem}
\begin{proof}
\leanok
\end{proof}

\begin{theorem}[Coefficients of a basis polynomial]
\label{thm:coefficient-T}\uses{thm:T-orthogonal,thm:coefficient-linear,thm:chebyshev-recurrence}
\lean{Polynomial.Chebyshev.coefficient_eval_T, Polynomial.Chebyshev.coefficient_const}
\leanok

For $m,n\in\N$, $a_n(T_m)$ is $1$ when $m=n$ and $0$ otherwise.

References: \cite[Table 18.3.1]{dlmf} and \cite[Chapter 3]{atap}.
\end{theorem}
\begin{proof}
\leanok
\end{proof}

\begin{theorem}[Dependence only on interval values]
\label{thm:coefficient-congr}
\lean{Polynomial.Chebyshev.coefficient_congr, Polynomial.Chebyshev.coefficient_congr_ae}
\leanok
\uses{thm:coefficient-cos}
For $f,g\in\mathcal I_n$ agreeing on $[-1,1]$, or agreeing $\mu_T$-almost
everywhere, $a_n(f)=a_n(g)$.

Reference: \cite[Chapter 3]{atap}.
\end{theorem}
\begin{proof}
\leanok
\end{proof}

\end{supporting}

\section{Completeness in the weighted space}\label{sec:textbook-coefficients-completeness}

In $L^2(\mu_T)$, functions that agree almost everywhere represent the same element. A Hilbert basis is an orthonormal family with dense span; its expansions converge in the norm of this space.

\begin{definition}[Polynomial classes in weighted L2]
\label{def:polynomial-L2}\uses{thm:T-integrable-transport}
\lean{Polynomial.Chebyshev.memLp_eval_measureT, Polynomial.Chebyshev.polynomialToL2}
\leanok

Let $\mathcal P:\R[X]\to L^2(\R,\mu_T;\R)$ send a polynomial $p$ to the class of
$x\mapsto p(x)$.


Reference: \cite[Table 18.3.1]{dlmf}.
\end{definition}

\begin{theorem}[First-kind completeness]
\label{thm:T-basis}
\lean{Polynomial.Chebyshev.hilbertBasisT, Polynomial.Chebyshev.hilbertBasisT_apply}
\leanok
\uses{thm:T-orthonormal,thm:T-dense}
The classes of $T_0/\sqrt\pi$ and $T_n\sqrt{2/\pi}$ for $n>0$ form a
Hilbert basis of $L^2(\R,\mu_T;\R)$.

Reference: \cite[Table 18.3.1]{dlmf},
with polynomial density from the Weierstrass theorem.
\end{theorem}
\begin{proof}
\leanok
\end{proof}

\begin{theorem}[Weighted L2 convergence of Chebyshev expansions]
\label{thm:T-L2-expansion}
\lean{Polynomial.Chebyshev.hasSum_coefficient_polynomialToL2, Polynomial.Chebyshev.tendsto_sum_coefficient_polynomialToL2, Polynomial.Chebyshev.tendsto_norm_sub_sum_coefficient_polynomialToL2}
\leanok
\uses{thm:T-coordinates}
For every $f\in L^2(\R,\mu_T;\R)$,
\[\lim_{N\to\infty}\left\|f-\sum_{n=0}^{N-1}a_n(f)\mathcal P(T_n)\right\|_{L^2(\mu_T)}=0.\]
This asserts convergence in the weighted $L^2$ norm, not pointwise or uniform
convergence.

References: \cite[18.2(viii)]{dlmf}, \cite[Chapter 3]{atap}.
\end{theorem}
\begin{proof}
\leanok
\end{proof}

\begin{theorem}[Chebyshev Parseval identity]
\label{thm:T-parseval}\uses{thm:T-coordinates}
\lean{Polynomial.Chebyshev.hasSum_coefficient_energy, Polynomial.Chebyshev.integral_sq_eq_sum_coefficient_sq}\leanok
For real $f\in L^2(\mu_T)$ with full-constant coefficients $a_n$,
\[
 \int f^2\,d\mu_T=\pi a_0^2+\frac\pi2\sum_{n=1}^\infty a_n^2.
\]
The weighted energy series converges to the integral.

Reference: the Hilbert-space Parseval identity applied to the Chebyshev basis.
\end{theorem}
\begin{proof}
\leanok
\end{proof}

\begin{supporting}[Polynomial classes, density and normalization]

\begin{theorem}[Polynomials in the first-kind weighted L2]
\label{thm:polynomial-L2-basic}\uses{def:polynomial-L2}
\lean{Polynomial.Chebyshev.coe_polynomialToL2, Polynomial.Chebyshev.inner_polynomialToL2}\leanok
Every real polynomial $p$ is square integrable for $\mu_T$. The class
$\mathcal P(p)$ agrees $\mu_T$-almost everywhere with $x\mapsto p(x)$, and
$\langle\mathcal P(p),\mathcal P(q)\rangle=\int pq\,d\mu_T$ for all $p,q$.

Reference: \cite[Table 18.3.1]{dlmf}.
The map is real linear.
\end{theorem}
\begin{proof}
\leanok
\end{proof}

\begin{theorem}[Hilbert coordinates and Chebyshev coefficients]
\label{thm:T-coordinates}\uses{thm:T-basis,def:coefficients,thm:chebyshev-recurrence}
\lean{Polynomial.Chebyshev.inner_polynomialToL2_left, Polynomial.Chebyshev.hilbertBasisT_repr_smul}
\leanok

For $f\in L^2(\R,\mu_T;\R)$ and $n\in\N$,
\[\langle e_n,f\rangle e_n=a_n(f)\mathcal P(T_n).\]
Coefficients are evaluated on any almost-everywhere representative and are independent of that choice.

References: \cite[Chapter 3]{atap},
and \cite[Table 18.3.1]{dlmf}.
\end{theorem}
\begin{proof}
\leanok
\end{proof}

\begin{lemma}[Finite first-kind measure and interval support]
\label{lem:T-finite}\uses{thm:T-integrable-transport}
\lean{Polynomial.Chebyshev.isFiniteMeasure_measureT, Polynomial.Chebyshev.ae_mem_Icc_measureT}
\leanok

The measure $\mu_T$ is finite and almost every point belongs to $[-1,1]$.

Reference: \cite[Table 18.3.1]{dlmf}.
\end{lemma}
\begin{proof}
\leanok
\end{proof}

\begin{theorem}[Density of polynomial classes]
\label{thm:polynomial-dense}\uses{lem:T-finite,thm:polynomial-L2-basic}
\lean{Polynomial.Chebyshev.polynomialToL2_denseRange}
\leanok

The range of $\mathcal P$ is dense in $L^2(\R,\mu_T;\R)$.

Reference: Weierstrass approximation, as in \cite[Chapter 1]{atap}.
\end{theorem}
\begin{proof}
\leanok
\end{proof}

\begin{definition}[Normalised first-kind classes]
\label{def:T-normalised}\uses{def:polynomial-L2}
\lean{Polynomial.Chebyshev.normalizedT}
\leanok

Write $h_0=\pi$ and $h_n=\pi/2$ for $n>0$, and set
$e_n=h_n^{-1/2}\mathcal P(T_n)$.
This is $T_0/\sqrt\pi$ at index zero and $T_n\sqrt{2/\pi}$ otherwise.

Reference: \cite[Table 18.3.1]{dlmf}.
\end{definition}

\begin{theorem}[First-kind orthonormality]
\label{thm:T-orthonormal}\uses{def:T-normalised,thm:T-orthogonal,thm:polynomial-L2-basic}
\lean{Polynomial.Chebyshev.inner_polynomialToL2_T, Polynomial.Chebyshev.orthonormal_normalizedT}
\leanok

The family $(e_n)_{n\in\N}$ is orthonormal.

Reference: \cite[Table 18.3.1]{dlmf}.
\end{theorem}
\begin{proof}
\leanok
\end{proof}

\begin{theorem}[Dense Chebyshev span]
\label{thm:T-dense}
\lean{Polynomial.Chebyshev.span_polynomialToL2_T, Polynomial.Chebyshev.normalizedT_dense_span}
\leanok
\uses{thm:polynomial-dense,def:T-normalised}
The closed real linear span of $(e_n)_{n\in\N}$ is the whole weighted $L^2$ space.

Reference: \cite[Chapters 1 and 3]{atap}.
\end{theorem}
\begin{proof}
\leanok
\end{proof}

\begin{theorem}[Linearity of coefficients]
\label{thm:coefficient-linear}
\lean{Polynomial.Chebyshev.coefficient_add, Polynomial.Chebyshev.coefficient_add_of_continuousOn, Polynomial.Chebyshev.coefficient_const_mul, Polynomial.Chebyshev.coefficient_sum}
\leanok
\uses{def:coefficients}
If $fT_n$ and $gT_n$ are integrable for $\mu_T$, then
$a_n(f+g)=a_n(f)+a_n(g)$. For every real $c$ and $f\in\mathcal I_n$,
$a_n(cf)=c\,a_n(f)$. For every finite family $(f_j)_{j\in J}$ in $\mathcal I_n$, $a_n(\sum_{j\in J}f_j)=\sum_{j\in J}a_n(f_j)$. In particular, additivity holds for real-valued functions continuous
on $[-1,1]$.

Reference: the weighted-integral convention in \cite[Chapter 3]{atap}.
\end{theorem}
\begin{proof}
\leanok
\end{proof}

\end{supporting}

\chapter{Smoothness and algebraic decay}\label{chap:smoothness}

\rvsummary{Turns finite smoothness into algebraic coefficient decay and uniform truncation error bounds.}

The decay of Chebyshev coefficients controls truncation error. Finite smoothness gives algebraic rates, with the last derivative allowed to have bounded variation.

\section{From smoothness to coefficient decay}\label{sec:textbook-smoothness-decay}

Total variation measures the sum of absolute increments over partitions. The derivative relation below can be iterated when the intermediate functions are absolutely continuous; the final derivative may have jumps.

\begin{theorem}[Derivative relation for coefficients]
\label{thm:derivative-coefficients}\uses{thm:cosine-integration-by-parts}
\lean{Polynomial.Chebyshev.absolutelyContinuous_cosine_pullback, Polynomial.Chebyshev.coefficient_eq_sub_derivative_coefficients, Polynomial.Chebyshev.coefficient_one_eq_derivative_coefficients}\leanok
Suppose $f$ is absolutely continuous on $[-1,1]$ with derivative representative
$g$ almost everywhere, and $\theta\mapsto g(\cos\theta)$ is integrable on
$[0,\pi]$. Then $f(\cos\theta)$ is absolutely continuous, with derivative
$-g(\cos\theta)\sin\theta$ almost everywhere. For $n\ge2$,
\[
a_n(f)=\frac{a_{n-1}(g)-a_{n+1}(g)}{2n}.
\]
The exceptional first-degree case is $a_1(f)=a_0(g)-a_2(g)/2$,
with the full constant coefficient convention.
The integrability assumption holds in particular when $g$ has bounded variation.

Reference: the repeated-integration step in \cite[Chapter 7]{atap}.
\end{theorem}
\begin{proof}
\leanok
\end{proof}

\begin{theorem}[Bounded-variation coefficient decay]
\label{thm:bv-coefficient-decay}\uses{thm:bv-interval-bound,thm:coefficient-cos}
\lean{Polynomial.Chebyshev.abs_coefficient_le_of_boundedVariation, Polynomial.Chebyshev.abs_coefficient_le_of_eVariationOn_le}\leanok
If $f$ has bounded variation on $[-1,1]$, then for $n\ge1$,
\[
|a_n(f)|\le\frac{2\operatorname{Var}_{[-1,1]}(f)}{\pi n}.
\]
In particular, if the variation is at most $V\ge0$, then
$|a_n(f)|\le2V/(\pi n)$. Continuity is not assumed.

Reference: \cite[Theorem 7.1]{atap}, case $\nu=0$.
\end{theorem}
\begin{proof}
\leanok
\end{proof}

\begin{theorem}[Finite-smoothness coefficient decay]
\label{thm:smooth-coefficient-decay}\uses{thm:derivative-coefficients,thm:bv-coefficient-decay}
\lean{Polynomial.Chebyshev.abs_coefficient_le_descFactorial_of_derivative_chain, Polynomial.Chebyshev.descFactorial_coefficient_bound_le_pow, Polynomial.Chebyshev.abs_coefficient_le_pow_of_derivative_chain}\leanok
Let $\nu\ge0$. Let $f_0=f,f_1,\ldots,f_\nu$ be real functions on $[-1,1]$.
For $j<\nu$, assume $f_j$ is absolutely continuous with derivative $f_{j+1}$
almost everywhere. Assume $f_\nu$ has total variation at most $V$, where
$V\ge0$. For $k\ge\nu+1$,
\[|a_k(f)|\le \frac{2V}{\pi k(k-1)\cdots(k-\nu)}
\le\frac{2V}{\pi(k-\nu)^{\nu+1}}.\]
The derivative chain specifies representatives at jumps without requiring
classical differentiability there. For $\nu=0$, only bounded variation of $f$
is assumed.

Reference: \cite[Theorem 7.1]{atap}.
\end{theorem}
\begin{proof}
\leanok
\end{proof}

\section{From coefficient decay to approximation error}\label{sec:textbook-smoothness-error}

First establish that the coefficient series represents $f$. Its tail then bounds the truncation error.

\begin{theorem}[Representation by an absolutely summable Chebyshev series]
\label{thm:series-representation}\uses{thm:T-L2-expansion,thm:coefficient-congr}
\lean{Polynomial.Chebyshev.hasSum_coefficient_mul_eval_T}\leanok
If $f$ is continuous on $[-1,1]$ and $\sum_{n\ge0}|a_n(f)|<\infty$, then
\[
f(x)=\sum_{n\ge0}a_n(f)T_n(x)\qquad (x\in[-1,1]).
\]
The endpoints are included.

Reference: \cite[Chapters 4 and 7]{atap}.
\end{theorem}
\begin{proof}
\leanok
\end{proof}

\begin{theorem}[Chebyshev coefficient-tail bound]
\label{thm:coefficient-tail}\uses{def:chebyshev-mathlib}
\lean{Polynomial.Chebyshev.norm_sub_sum_smul_le_tsum_tail, Polynomial.Chebyshev.summable_mul_eval_T, Polynomial.Chebyshev.abs_sub_sum_le_tsum_tail, Polynomial.Chebyshev.abs_sub_sum_le_geometric}\leanok
For real or complex coefficients, if $\sum_{k\ge0}|a_k|<\infty$ and $f(x)=\sum_{k\ge0}a_kT_k(x)$ on $[-1,1]$, then
\[|f(x)-\sum_{k=0}^N a_kT_k(x)|\le\sum_{k>N}|a_k|\quad(x\in[-1,1]).\]
For real coefficients, if $|a_k|\le Cq^k$ with $0\le q<1$, this is at most $Cq^{N+1}/(1-q)$.


Reference: \cite[equation 4.8 and Theorems 7.2 and 8.2]{atap}.
\end{theorem}
\begin{proof}
\leanok
\end{proof}

\begin{theorem}[Algebraic truncation error]
\label{thm:algebraic-error}\uses{thm:smooth-coefficient-decay,thm:coefficient-tail,thm:series-representation,thm:pseries-tail}
\lean{Polynomial.Chebyshev.summable_abs_coefficient_of_derivative_chain, Polynomial.Chebyshev.abs_sub_sum_coefficient_le_of_derivative_chain, Polynomial.Chebyshev.sup_abs_sub_sum_coefficient_le_of_derivative_chain, Polynomial.Chebyshev.bddAbove_abs_sub_sum_coefficient_of_derivative_chain}\leanok
Under the hypotheses of Theorem~\ref{thm:smooth-coefficient-decay}, with
$\nu\ge1$ and $N>\nu$, the degree-$N$
Chebyshev projection satisfies
\[\sup_{x\in[-1,1]}|f(x)-\sum_{k=0}^N a_k(f)T_k(x)|
\le\frac{2V}{\pi\nu(N-\nu)^\nu}.\]
The coefficients are absolutely summable.

Reference: \cite[Theorem 7.2]{atap}.
\end{theorem}
\begin{proof}
\leanok
\end{proof}

\begin{supporting}[Integration by parts and tail estimates]

\begin{theorem}[Cosine integration by parts]
\label{thm:cosine-integration-by-parts}\uses{thm:coefficient-cos,thm:integrable-derivative}
\lean{Polynomial.Chebyshev.coefficient_eq_integral_sin, Polynomial.Chebyshev.abs_coefficient_le_integral_abs_deriv, Polynomial.Chebyshev.coefficient_eq_integral_sin_of_absolutelyContinuous}\leanok
Put $g(\theta)=f(\cos\theta)$. If $g$ is absolutely continuous on $[0,\pi]$,
let $h$ be any integrable derivative representative on that interval. For $n>0$,
\[a_n(f)=-\frac{2}{\pi n}\int_0^\pi h(\theta)\sin(n\theta)\,d\theta.\]
Such representatives exist, and the formula is independent of their choice.
For continuous $g$ differentiable on $(0,\pi)$ with integrable derivative,
we also have $|a_n(f)|\le 2\|g'\|_{L^1(0,\pi)}/(\pi n)$.

Reference: \cite[Chapter 7]{atap}.
\end{theorem}
\begin{proof}
\leanok
\end{proof}

\begin{theorem}[Integration against a signed density]
\label{thm:signed-density-integral}
\lean{MeasureTheory.VectorMeasure.integral_withDensityᵥ_of_bound, MeasureTheory.VectorMeasure.setIntegral_withDensityᵥ_of_bound}\leanok
Let $\mu$ be a measure, $h\in L^1(\mu)$ real-valued, and $u$ an
almost-everywhere strongly measurable, essentially bounded function with
values in a real normed vector space. Integrating $u$ against the signed
measure with density $h$ agrees with $\int h u\,d\mu$. The same identity
holds after restriction to a measurable set.
\end{theorem}
\begin{proof}
\leanok
\end{proof}

\begin{theorem}[Bounded-variation integration by parts]
\label{thm:bv-integration-by-parts}\uses{thm:signed-density-integral,thm:integrable-derivative}
\lean{BoundedVariationOn.vectorMeasure_eq_withDensityᵥ_of_integral_eq_sub, BoundedVariationOn.vectorMeasure_eq_withDensityᵥ_deriv, BoundedVariationOn.integral_mul_eq_sub_vectorMeasure, BoundedVariationOn.abs_integral_mul_le_of_primitive}\leanok
Let $f,g:\mathbb R\to\mathbb R$ have bounded variation on $\mathbb R$,
with $g$ continuous. Suppose $h\in L^1(\mathbb R)$ and
$g(v)-g(u)=\int_u^v h(x)\,dx$ for every $u\le v$.
Then $dg=h\,dx$, and, for $a\le b$,
\[
\int_a^b h(x)f(x)\,dx
 =g(b)f(b+)-g(a)f(a-)-\int_{[a,b]}g\,df.
\]
In particular, if $g(a)=g(b)=0$ and $|g|\le C$ on $[a,b]$, with $C\ge0$,
then
\[
\left|\int_a^b h(x)f(x)\,dx\right|
 \le C\operatorname{Var}_{\mathbb R}(f).
\]
The density identification also holds for a continuous, locally absolutely
continuous function $g$ of bounded variation and any global integrable
derivative representative $h\in\mathcal D_{dx}(g)$. Jumps of $f$ are allowed.

Reference: the integration-by-parts step in \cite[Chapter 7]{atap}.
\end{theorem}
\begin{proof}
\leanok
\end{proof}

\begin{theorem}[Bounded-variation interval estimate]
\label{thm:bv-interval-bound}\uses{thm:bv-integration-by-parts}
\lean{BoundedVariationOn.eVariationOn_comp_projIcc_le, BoundedVariationOn.comp_projIcc, BoundedVariationOn.ae_hasDerivAt_comp_projIcc, BoundedVariationOn.abs_integral_mul_le_of_lipschitz_primitive}\leanok
Let $a\le b$, let $f$ have bounded variation on $[a,b]$, and let
$g:\mathbb R\to\mathbb R$ be Lipschitz. Suppose $g'=h$ on $(a,b)$,
$h$ is integrable on $[a,b]$, $g(a)=g(b)=0$, and $|g|\le C$ there, with $C\ge0$.
Then
\[
\left|\int_a^b h(x)f(x)\,dx\right|
 \le C\operatorname{Var}_{[a,b]}(f).
\]

Reference: the bounded-variation integration step in \cite[Chapter 7]{atap}.
\end{theorem}
\begin{proof}
\leanok
\end{proof}

\begin{theorem}[$p$-series tail estimate]
\label{thm:pseries-tail}
\lean{Real.tsum_one_div_nat_add_pow_le}\leanok
For integers $\nu,M\ge1$,
\[
\sum_{k=M+1}^{\infty}\frac{1}{k^{\nu+1}}
 \le\frac{1}{\nu M^\nu}.
\]

Reference: the integral-test estimate in \cite[Theorem 7.2]{atap}.
\end{theorem}
\begin{proof}
\leanok
\end{proof}

\end{supporting}

\chapter{Interpolation at Chebyshev points}\label{chap:interpolation}

\rvsummary{Interpolation at Chebyshev points: discrete orthogonality, aliasing, the barycentric formula and logarithmic Lebesgue constants.}

Interpolation constructs a polynomial from sampled values. On the Chebyshev extrema grid, a discrete cosine transform gives its coefficients, while the barycentric formula evaluates it directly from the samples.

\section{Interpolation and the cosine grid}\label{sec:textbook-interpolation-grid}

For distinct nodes $x_0,\ldots,x_N$, the cardinal polynomials are $\ell_j(X)=\prod_{i\ne j}(X-x_i)/(x_j-x_i)$. Their values at the nodes are $\ell_j(x_i)=\delta_{ij}$, so $\sum_j f(x_j)\ell_j$ interpolates the data (Theorem~\ref{thm:lagrange-interpolation}). We now choose the cosine grid.

\begin{definition}[Classical Chebyshev cosine nodes]
\label{def:classical-chebyshev-nodes}
\lean{Polynomial.Chebyshev.classicalNode}\leanok
For a positive natural order $N$ and a natural angular index $j$, define
$x_{N,j}=\cos(j\pi/N)$. The extrema grid consists of $0\le j\le N$.
Reference: \cite[Chapter 4]{atap}.
\end{definition}

\begin{definition}[Finite extrema grid]
\label{def:extrema-grid}\uses{def:classical-chebyshev-nodes}
\lean{Polynomial.Chebyshev.extremaGridNode}\leanok
For $N\ge0$, index the grid by $0\le j\le N$. When $N>0$ its points are
$x_{N,j}=\cos(j\pi/N)$. For $N=0$ there is one node, $x_{0,0}=1$.
\end{definition}

\begin{definition}[Classical extrema cardinal polynomials]
\label{def:extrema-cardinal}\uses{def:lagrange-basis,thm:lagrange-input,thm:classical-chebyshev-nodes}
\lean{Polynomial.Chebyshev.extremaCardinal}\leanok
For a finite-grid index $j$, let $\ell_{N,j}$ be its classical Lagrange
cardinal polynomial. At order zero it is the constant polynomial $1$.
Reference: \cite[Chapter 4]{atap}.
\end{definition}

\begin{theorem}[Extrema cardinal-polynomial characterisation]
\label{thm:extrema-cardinal}\uses{def:extrema-cardinal,thm:lagrange-basis,thm:extrema-grid}
\lean{Polynomial.Chebyshev.extremaCardinal_zero,Polynomial.Chebyshev.extremaCardinal_positive,Polynomial.Chebyshev.eval_extremaCardinal}\leanok
For positive order the cardinal polynomials are the classical Lagrange products;
for order zero the single polynomial is $1$. For every order,
$\ell_{N,i}(x_{N,j})=\delta_{ij}$.
\end{theorem}
\begin{proof}
\leanok
\end{proof}

\begin{definition}[Chebyshev-extrema interpolant]
\label{def:extrema-interpolant}\uses{thm:classical-chebyshev-nodes,def:lagrange-interpolant}
\lean{Polynomial.Chebyshev.extremaInterpolant}\leanok
For $N\ge1$, let $I_Nf\in\C[X]$ be the Lagrange interpolant of $f:\R\to\C$ at the
$N+1$ distinct Chebyshev extrema $x_j=\cos(j\pi/N)$, $0\le j\le N$.
For $N=0$, define $I_0f$ to be the constant polynomial $f(1)$.

Reference: \cite[Chapters 4 and 8]{atap}.
\end{definition}

\begin{theorem}[Interpolation and degree]
\label{thm:extrema-interpolant-basic}\uses{thm:extrema-interpolant-definition,thm:lagrange-interpolant}
\lean{Polynomial.Chebyshev.eval_extremaInterpolant_classicalNode, Polynomial.Chebyshev.degree_extremaInterpolant_lt}\leanok
For $N\ge1$ and $0\le j\le N$, $I_Nf(x_j)=f(x_j)$, and $I_Nf$ has degree at
most $N$.

Reference: \cite[Chapter 4]{atap}.
\end{theorem}
\begin{proof}
\leanok
\end{proof}

\begin{supporting}[Lagrange products and interpolation domains]

\begin{theorem}[Classical Chebyshev node characterisation]
\label{thm:classical-chebyshev-nodes}\uses{def:classical-chebyshev-nodes}
\lean{Polynomial.Chebyshev.classicalNode_apply,Polynomial.Chebyshev.strictAntiOn_classicalNode,Polynomial.Chebyshev.classicalNode_mem_Icc,Polynomial.Chebyshev.classicalNode_endpoints}\leanok
The nodes obey the displayed cosine formula and lie in $[-1,1]$.
They are strictly decreasing for $0\le j\le N$, with endpoints $1$ and $-1$.
\end{theorem}
\begin{proof}
\leanok
\end{proof}

\begin{theorem}[Finite extrema-grid characterisation]
\label{thm:extrema-grid}\uses{def:extrema-grid,thm:classical-chebyshev-nodes}
\lean{Polynomial.Chebyshev.extremaGridNode_zero,Polynomial.Chebyshev.extremaGridNode_positive,Polynomial.Chebyshev.extremaGridNode_mem_Icc,Polynomial.Chebyshev.extremaGridNode_injective}\leanok
For $N=0$, $x_{0,0}=1$; for $N>0$, $x_{N,j}=\cos(j\pi/N)$. All grid points lie in
$[-1,1]$ and are pairwise distinct.
\end{theorem}
\begin{proof}
\leanok
\end{proof}

\begin{theorem}[Extrema-interpolant characterisation]
\label{thm:extrema-interpolant-definition}\uses{def:extrema-interpolant}
\lean{Polynomial.Chebyshev.extremaInterpolant_positive,Polynomial.Chebyshev.extremaInterpolant_zero}\leanok
For positive order, $I_Nf$ is classical interpolation on the cosine grid.
For order zero, $I_0f$ is exactly the constant polynomial $f(1)$.
\end{theorem}
\begin{proof}
\leanok
\end{proof}

\begin{definition}[Domain of the Lagrange products]
\label{def:lagrange-input}
\lean{Lagrange.BasisInput}\leanok
For a finite index set $s$ and field-valued nodes $v$, define the admissible
indices to be those $i$ satisfying $v_i\ne v_j$ for every
$j\in s\setminus\{i\}$. An admissible index need not belong to $s$.
\end{definition}

\begin{theorem}[Nonvanishing of Lagrange denominators]
\label{thm:lagrange-input}\uses{def:lagrange-input}
\lean{Lagrange.BasisInput.of_injOn,Lagrange.BasisInput.id}\leanok
If the nodes $v_i$ are distinct on $s$, then every $i\in s$ is admissible:
$v_i-v_j\ne0$ for $j\in s\setminus\{i\}$.
For identity nodes on a finite subset $s$ of a field, every field element $x$
is admissible, since $x-y\ne0$ for $y\in s\setminus\{x\}$; this includes $x\notin s$.
\end{theorem}
\begin{proof}
\leanok
\end{proof}

\begin{definition}[Classical Lagrange polynomials]
\label{def:lagrange-basis}\uses{def:lagrange-input}
\lean{Lagrange.classicalBasis}\leanok
For an admissible index $i$ of Definition~\ref{def:lagrange-input}, define
$\ell_i(X)=\prod_{j\in s\setminus\{i\}}(X-v_j)/(v_i-v_j)$.
When $i\notin s$ this is the normalized nodal product, extending the
cardinal-polynomial construction. Reference: \cite[3.3(i)]{dlmf}.
\end{definition}

\begin{theorem}[Classical Lagrange product characterisation]
\label{thm:lagrange-basis}\uses{def:lagrange-basis}
\lean{Lagrange.classicalBasis_eq_product,Lagrange.eval_classicalBasis_self,Lagrange.eval_classicalBasis_of_ne,Lagrange.classicalBasis_ne_zero}\leanok
The polynomial is the displayed product, has $\ell_i(v_i)=1$, and vanishes
at $v_j$ for $j\in s$, $j\ne i$. In particular it is nonzero.
\end{theorem}
\begin{proof}
\leanok
\end{proof}

\begin{theorem}[Real and complex cardinal products]
\label{thm:lagrange-complex}\uses{def:lagrange-basis}
\lean{Polynomial.Chebyshev.complex_basisInput,Lagrange.eval_classicalBasis_complex}\leanok
For a finite set of natural indices $S$, real nodes $v_i$, and any index $j$
with $v_j\ne v_i$ for every $i\in S\setminus\{j\}$, the normalized cardinal
product has the same value at a real point whether computed over $\mathbb R$
or over $\mathbb C$. The same separation condition holds over $\mathbb C$.
The index $j$ need not belong to $S$.
\end{theorem}
\begin{proof}
\leanok
\end{proof}

\begin{theorem}[Lagrange interpolation]
\label{thm:lagrange-interpolation}\uses{thm:lagrange-input,def:lagrange-basis}
\lean{Lagrange.natDegree_classicalBasis,Lagrange.eq_sum_classicalBasis}\leanok
For a finite set $s$ of indices and distinct field-valued nodes $v_i$, there are
Lagrange polynomials $\ell_i$ with $\ell_i(v_j)=\delta_{ij}$ for
$i,j\in s$. Each $\ell_i$ is nonzero of degree $|s|-1$, and every polynomial
of degree less than $|s|$ is its Lagrange interpolant.

Reference: \cite[3.3(i)]{dlmf}.
\end{theorem}
\begin{proof}
\leanok
\end{proof}

\begin{definition}[Classical nodal weights]
\label{def:lagrange-nodal-weight}\uses{def:lagrange-input}
\lean{Lagrange.classicalNodalWeight}\leanok
For an admissible index $i$ of Definition~\ref{def:lagrange-input}, define
$\omega_i=\prod_{j\in s\setminus\{i\}}(v_i-v_j)^{-1}$.
This also permits $i\notin s$ whenever all displayed denominators are nonzero.
Reference: \cite[3.3(i)]{dlmf}.
\end{definition}

\begin{theorem}[Classical nodal-weight characterisation]
\label{thm:lagrange-nodal-weight}\uses{def:lagrange-nodal-weight,def:lagrange-basis,def:polynomial-linear-maps}
\lean{Lagrange.classicalNodalWeight_apply,Lagrange.classicalNodalWeight_ne_zero,Lagrange.classicalNodalWeight_eq_derivative,Lagrange.derivative_nodal_ne_zero,Lagrange.eval_classicalBasis_not_at_node}\leanok
The weight is the displayed product and is nonzero. Writing
$q(X)=\prod_{j\in s}(X-v_j)$, if $i\in s$ then $q'(v_i)\ne0$ and
$\omega_i=1/q'(v_i)$. For $i\in s$ and $x\ne v_i$,
$\ell_i(x)=q(x)\omega_i/(x-v_i)$.
\end{theorem}
\begin{proof}
\leanok
\end{proof}

\begin{definition}[Finite interpolation configurations]
\label{def:lagrange-interpolation-input}
\lean{Lagrange.InterpolationInput}\leanok
A finite interpolation configuration consists of a finite index set $s$
and field-valued nodes $(v_i)_{i\in s}$ with $v_i\ne v_j$ whenever $i\ne j$.
The set $s$ may be empty.
\end{definition}

\begin{theorem}[Nonvanishing denominators at interpolation nodes]
\label{thm:lagrange-interpolation-input}\uses{def:lagrange-interpolation-input,thm:lagrange-input}
\lean{Lagrange.InterpolationInput.basis}\leanok
For a finite interpolation configuration $(s,v)$, every $i\in s$
satisfies $v_i-v_j\ne0$ for $j\in s\setminus\{i\}$, so the normalized
nodal product $\ell_i$ is defined.
\end{theorem}
\begin{proof}
\leanok
\end{proof}

\begin{definition}[Classical interpolation polynomial]
\label{def:lagrange-interpolant}\uses{def:lagrange-interpolation-input}
\lean{Lagrange.classicalInterpolate}\leanok
For a finite interpolation configuration $(s,v)$ and values $(r_i)_{i\in s}$, define
$I r=\sum_{i\in s}r_i\ell_i$. For the empty set this is the zero polynomial.
Reference: \cite[3.3(i)]{dlmf}.
\end{definition}

\begin{theorem}[Classical interpolation characterisation]
\label{thm:lagrange-interpolant}\uses{def:lagrange-interpolant,thm:lagrange-interpolation-input,def:lagrange-basis}
\lean{Lagrange.classicalInterpolate_apply,Lagrange.eval_classicalInterpolate,Lagrange.degree_classicalInterpolate}\leanok
The interpolant is the displayed cardinal-polynomial sum, attains the data
at every node, and has degree less than $|s|$, with the degree of zero
understood as $-\infty$.
\end{theorem}
\begin{proof}
\leanok
\end{proof}

\begin{definition}[Domain of barycentric evaluation]
\label{def:barycentric-input}\uses{def:lagrange-interpolation-input}
\lean{Lagrange.BarycentricInput}\leanok
A barycentric evaluation configuration consists of a nonempty finite set of distinct nodes and
an evaluation point $z$ different from every node.
\end{definition}

\begin{theorem}[Nonvanishing barycentric denominators]
\label{thm:barycentric-input}\uses{def:barycentric-input,thm:lagrange-nodal-weight,thm:lagrange-interpolation-input}
\lean{Lagrange.BarycentricInput.valid}\leanok
For each node, $\omega_i\ne0$ and $z-v_i\ne0$.
\end{theorem}
\begin{proof}
\leanok
\end{proof}

\end{supporting}

\section{From values to coefficients}\label{sec:textbook-interpolation-dct}

Discrete orthogonality gives a cosine transform of the sampled values. The half weights at the endpoints also affect the first and last coefficient.

\begin{definition}[Half-endpoint discrete sum and DCT-I coefficients]
\label{def:extrema-dct}\uses{def:chebyshev-mathlib,def:classical-chebyshev-nodes}
\lean{Polynomial.Chebyshev.extremaSum, Polynomial.Chebyshev.extremaCoefficient, Polynomial.Chebyshev.dctInterpolant}\leanok
For $N\ge1$, put $x_j=\cos(j\pi/N)$ and
$S_N(g)=\tfrac12g(0)+\sum_{j=1}^{N-1}g(j)+\tfrac12g(N)$.
For real-valued data and $0\le k\le N$, define
\[
 c_k=\frac{2-\mathbf1_{k=0\text{ or }k=N}}{N}
 S_N\bigl(j\mapsto f(x_j)T_k(x_j)\bigr),\qquad
 P_Nf=\sum_{k=0}^N c_kT_k.
\]

For $N=0$, define $P_0f$ to be the constant polynomial $f(1)$.

Reference: \cite[Chapters 4 and 6]{atap}.
\end{definition}

\begin{theorem}[Discrete Chebyshev orthogonality]
\label{thm:discrete-T-orthogonal}\uses{def:extrema-dct,thm:chebyshev-trig,thm:chebyshev-recurrence}
\lean{Polynomial.Chebyshev.extremaSum_mul_eval_T_eq_zero, Polynomial.Chebyshev.extremaSum_sq_eval_T, Polynomial.Chebyshev.extremaCoefficient_eval_T}\leanok
For $N\ge1$ and $0\le j,k\le N$, the discrete inner product
$S_N(\ell\mapsto T_j(x_\ell)T_k(x_\ell))$ vanishes if $j\ne k$.
Its diagonal is $N$ for $k=0,N$ and $N/2$ otherwise.
Consequently $c_k(T_j)=\delta_{jk}$.

Reference: \cite[Chapter 4]{atap}.
\end{theorem}
\begin{proof}
\leanok
\end{proof}

\begin{theorem}[DCT reconstruction]
\label{thm:dct-reconstruction}\uses{thm:discrete-T-orthogonal,thm:coefficient-T,thm:extrema-dct-definition,thm:lagrange-interpolant,thm:classical-chebyshev-nodes}
\lean{Polynomial.Chebyshev.dctInterpolant_eval, Polynomial.Chebyshev.dctInterpolant_eq_lagrange, Polynomial.Chebyshev.eval_dctInterpolant_node, Polynomial.Chebyshev.coefficient_dctInterpolant}\leanok
For $N\ge1$, $P_Nf$ equals the Lagrange interpolant at the extrema and
$P_Nf(x_j)=f(x_j)$ for $0\le j\le N$. It reproduces every real polynomial
of degree at most $N$. Its integral-defined Chebyshev coefficient of index $k$
is $c_k$ for $k\le N$ and zero otherwise.

Reference: \cite[Chapters 4 and 6]{atap}.
\end{theorem}
\begin{proof}
\leanok
\end{proof}

\paragraph{Example (Julia).} This direct implementation evaluates the DCT formula in $O(N^2)$ operations. Sampling $T_3$ at six extrema gives its unit coefficient vector in exact arithmetic. The computation uses floating-point values. An FFT-based implementation can be faster; none is implemented here.

\begin{juliacode}
function chebcoeffs(f, N)            # c_k of Definition (DCT-I)
    @assert N >= 1
    x = [cos(j*pi/N) for j in 0:N]
    w = [j == 0 || j == N ? 1/2 : 1.0 for j in 0:N]
    [(k == 0 || k == N ? 1 : 2)/N *
        sum(w[j+1]*f(x[j+1])*cos(k*j*pi/N) for j in 0:N) for k in 0:N]
end
chebcoeffs(x -> 4x^3 - 3x, 5)        # [0, 0, 0, 1, 0, 0] up to rounding
\end{juliacode}

\begin{supporting}[Discrete-transform identities]

\begin{theorem}[Discrete-transform characterisation]
\label{thm:extrema-dct-definition}\uses{def:extrema-dct,thm:chebyshev-recurrence}
\lean{Polynomial.Chebyshev.extremaCoefficient_positive,Polynomial.Chebyshev.dctInterpolant_zero}\leanok
For positive order the coefficients are the stated half-endpoint cosine sum.
For $N=0$, $P_0f$ is the constant polynomial $f(1)$.
\end{theorem}
\begin{proof}
\leanok
\end{proof}

\end{supporting}

\section{Aliasing and interpolation error}\label{sec:textbook-interpolation-aliasing}

Frequencies above the grid order can have identical samples. The folded index identifies which coefficients become indistinguishable.

\begin{definition}[Folded frequency]
\label{def:extrema-aliasing}
\lean{Polynomial.Chebyshev.aliasIndex}\leanok
For $N\ge1$, fold a frequency $k$ by setting $r=k\bmod 2N$ and
$A_N(k)=\min(r,2N-r)$. For $N=0$, set $A_0(k)=0$.
\end{definition}

\begin{theorem}[Aliasing at Chebyshev extrema]
\label{thm:extrema-aliasing}\uses{thm:chebyshev-trig,def:extrema-aliasing,def:classical-chebyshev-nodes}
\lean{Polynomial.Chebyshev.aliasIndex_zero, Polynomial.Chebyshev.aliasIndex_positive, Polynomial.Chebyshev.aliasIndex_le, Polynomial.Chebyshev.aliasIndex_eq_self, Polynomial.Chebyshev.eval_T_node_aliasIndex}\leanok
For $N\ge1$, set $r=k\bmod 2N$ and $m=\min(r,2N-r)$. Then $0\le m\le N$,
$m=k$ if $k\le N$, and $T_k(\cos(j\pi/N))=T_m(\cos(j\pi/N))$ for every
nonnegative integer $j$. At order zero, $A_0(k)=0$ for every $k$.

Reference: \cite[Chapter 4]{atap}.
\end{theorem}
\begin{proof}
\leanok
\end{proof}

\begin{theorem}[Grouped coefficient aliasing]
\label{thm:dct-coefficient-aliasing}\uses{thm:extrema-aliasing,thm:extrema-grid,thm:discrete-T-orthogonal}
\lean{Polynomial.Chebyshev.extremaCoefficient_eval_T_alias, Polynomial.Chebyshev.extremaCoefficient_sum_eval_T, Polynomial.Chebyshev.hasSum_alias_coefficients}\leanok
Let $N\ge1$, $0\le k\le N$, and let $A_N(m)$ be the folded alias index.
For a finite expansion $f=\sum_{m\in s}a_mT_m$,
\[
 c_k(f)=\sum_{\substack{m\in s\\A_N(m)=k}}a_m.
\]
More generally, if $\sum_{m\ge0}a_mT_m(x_j)$ is absolutely summable with sum
$f(x_j)$ at every extremum, then
$\sum_{m\ge0}\mathbf1_{A_N(m)=k}a_m$ is summable with sum $c_k(f)$.
Summability is unconditional: sums are limits over finite subsets of indices.

Reference: \cite[Theorem 4.2]{atap}.
\end{theorem}
\begin{proof}
\leanok
\end{proof}

\begin{theorem}[Interpolation coefficient-tail bound]
\label{thm:interpolation-tail}\uses{thm:extrema-interpolant-definition,thm:extrema-aliasing,thm:extrema-cardinal}
\lean{Polynomial.Chebyshev.extremaInterpolant_T, Polynomial.Chebyshev.norm_eval_extremaInterpolant_T_le_one, Polynomial.Chebyshev.eval_extremaInterpolant, Polynomial.Chebyshev.hasSum_eval_extremaInterpolant, Polynomial.Chebyshev.tendsto_sum_eval_extremaInterpolant, Polynomial.Chebyshev.norm_sub_extremaInterpolant_le_tsum_tail}\leanok
Suppose $\sum_{k\ge0}|a_k|<\infty$ and $f(x)=\sum_{k\ge0}a_kT_k(x)$ on
$[-1,1]$. For $N\ge1$ and $x\in[-1,1]$,
\[
 |f(x)-I_Nf(x)|\le2\sum_{k>N}|a_k|.
\]
In particular $I_NT_k=T_m$ for the aliased index $m$, so $|I_NT_k(x)|\le1$.

Reference: \cite[Chapter 4]{atap}.

The nodal expansion is
$I_N f(z)=\sum_{j=0}^{N} f(x_{N,j})\ell_{N,j}(z)$ using the finite grid of
Definition~\ref{def:extrema-grid}. Interpolation commutes with any Chebyshev
series whose ordinary partial sums converge at each of these finitely many points; absolute convergence is not required. Both identities also
hold for $N=0$, when the only interpolation node is $x_{0,0}=1$.
\end{theorem}
\begin{proof}
\leanok
\end{proof}

\section{Barycentric evaluation}\label{sec:textbook-interpolation-barycentric}

The barycentric formula evaluates the interpolant directly from its samples. Its weights are $(-1)^j$, halved at the endpoints. At a node, use the prescribed sample value.

\begin{definition}[Extrema nodal polynomial and weights]
\label{def:extrema-barycentric}\uses{def:polynomial-linear-maps,def:chebyshev-mathlib}
\lean{Polynomial.Chebyshev.extremaNodal, Polynomial.Chebyshev.extremaEndpointFactor, Polynomial.Chebyshev.extremaBarycentricWeight}\leanok
For $N\ge1$ and $0\le j\le N$, write $x_j=\cos(j\pi/N)$ for the extrema nodes. Set
$F_N=(X^2-1)T_N'$, and let $e_j=2$ for $j=0$ or $j=N$, and $e_j=1$ otherwise.
The normalised barycentric weights are $w_j=(-1)^j/e_j$.

Reference: \cite[Theorem 5.2]{atap}.
\end{definition}

\begin{theorem}[Chebyshev barycentric formula]
\label{thm:extrema-barycentric}\uses{thm:extrema-nodal,thm:lagrange-complex,thm:extrema-interpolant-definition,def:lagrange-nodal-weight,thm:lagrange-input}
\lean{Polynomial.Chebyshev.extremaEndpointFactor_pos, Polynomial.Chebyshev.nodalWeight_extrema, Polynomial.Chebyshev.extremaBarycentricScale_ne_zero, Polynomial.Chebyshev.nodalWeight_complex_extrema, Polynomial.Chebyshev.extremaBarycentric_denominator_ne_zero, Polynomial.Chebyshev.eval_extremaInterpolant_barycentric, Polynomial.Chebyshev.eval_extremaBasis}\leanok
For $N\ge1$, arbitrary $f:\mathbb R\to\mathbb C$, and $z\in\mathbb C$ distinct
from all nodes,
\[
 I_Nf(z)=\frac{\sum_{j=0}^N w_j f(x_j)/(z-x_j)}
                   {\sum_{j=0}^N w_j/(z-x_j)}.
\]
The denominator is nonzero. At a node the interpolant takes its prescribed value. The order-zero interpolant is the constant $f(1)$.
The product weights equal $\operatorname{lc}(F_N)w_j/N^2$, with a nonzero common
factor. For real $x\ne x_j$, the cardinal polynomial satisfies
\[
 \ell_j(x)=\frac{F_N(x)(-1)^j}{N^2 e_j(x-x_j)}.
\]

Reference: \cite[Theorem 5.2]{atap}.
\end{theorem}
\begin{proof}
\leanok
\end{proof}

\section{Sensitivity and near-best interpolation}\label{sec:textbook-interpolation-lebesgue}

The sum of the absolute cardinal values measures sensitivity to perturbations in the samples. Its maximum over the interval also controls the error relative to a polynomial competitor.

\begin{definition}[Extrema Lebesgue function]
\label{def:extrema-lebesgue}\uses{def:extrema-cardinal}
\lean{Polynomial.Chebyshev.extremaLebesgueFunction}\leanok
Define $\lambda_N(x)=\sum_{j=0}^N|\ell_{N,j}(x)|$ using the classical
finite-grid cardinal polynomials, for every real $x$ and $N\ge0$.
Reference: \cite[Chapter 15]{atap}.
\end{definition}

Here $\sup$ is the least upper bound of a nonempty bounded set; the boundedness assertion is recorded in Theorem~\ref{thm:extrema-lebesgue-function}.

\begin{definition}[Classical extrema Lebesgue constant]
\label{def:extrema-lebesgue-constant}\uses{thm:extrema-lebesgue-function,def:classical-real-supremum}
\lean{Polynomial.Chebyshev.extremaLebesgueSupremumInput,Polynomial.Chebyshev.extremaLebesgueConstant}\leanok

For $N\ge0$, define $\Lambda_N$ as the classical real supremum of $\{\lambda_N(x):x\in[-1,1]\}$.
Reference: \cite[Chapter 15]{atap}.
\end{definition}

\begin{theorem}[Logarithmic extrema Lebesgue bound]
\label{thm:extrema-lebesgue}\uses{def:extrema-lebesgue-constant}
\lean{Polynomial.Chebyshev.extremaLebesgueFunction_le_log, Polynomial.Chebyshev.extremaLebesgueFunction_le_constant, Polynomial.Chebyshev.extremaLebesgueConstant_zero, Polynomial.Chebyshev.extremaLebesgueConstant_le_log}\leanok
The Lebesgue function is bounded above on $[-1,1]$, so the real supremum is
well-founded and bounds every pointwise value. For $N\ge1$,
\[
 \lambda_N(x)\le2\pi+2H_N\le2\pi+2+2\log N,\qquad
 \Lambda_N\le2\pi+2+2\log N.
\]
For $N=0$, $\lambda_0(x)=\Lambda_0=1$.


Reference: \cite[Chapter 15]{atap}; its Theorem 15.2 gives the sharper constant.
\end{theorem}
\begin{proof}
\leanok
\end{proof}

\begin{theorem}[Near-best Chebyshev interpolation]
\label{thm:extrema-near-best}\uses{thm:extrema-lebesgue,thm:interpolation-tail}
\lean{Polynomial.Chebyshev.norm_eval_extremaInterpolant_le, Polynomial.Chebyshev.extremaInterpolant_eq_self, Polynomial.Chebyshev.extremaInterpolant_sub, Polynomial.Chebyshev.norm_sub_extremaInterpolant_le_competitor, Polynomial.Chebyshev.norm_sub_extremaInterpolant_le_lebesgueConstant, Polynomial.Chebyshev.norm_sub_extremaInterpolant_le_log_competitor}\leanok
For arbitrary $f:\mathbb R\to\mathbb C$ and any complex polynomial $p$ of
degree at most $N$, suppose $|f(x)-p(x)|\le\varepsilon$ throughout $[-1,1]$.
Then, on the same interval,
\[
 |f(x)-I_Nf(x)|\le(1+\lambda_N(x))\varepsilon
                  \le(1+\Lambda_N)\varepsilon.
\]
For $N\ge1$, this is at most $(2\pi+3+2\log N)\varepsilon$.
No continuity assumption is needed for this uniform-competitor formulation.

Reference: \cite[Theorem 15.1]{atap}.
\end{theorem}
\begin{proof}
\leanok
\end{proof}

\begin{supporting}[Suprema and harmonic-sum estimates]

\begin{definition}[Classical real supremum]
\label{def:classical-real-supremum}
\lean{Real.SupremumInput,Real.classicalSup}\leanok
For a nonempty set $S\subseteq\R$ that is bounded above, define $\sup S$ to be its least upper bound in $\R$.
\end{definition}

\begin{theorem}[Classical real supremum characterisation]
\label{thm:classical-real-supremum}\uses{def:classical-real-supremum}
\lean{Real.classicalSup_isLUB}\leanok
For a nonempty set $S\subseteq\R$ that is bounded above, $x\le\sup S$ for every $x\in S$. If $x\le b$ for every $x\in S$, then $\sup S\le b$.
\end{theorem}
\begin{proof}
\leanok
\end{proof}

\begin{theorem}[Lebesgue function characterisation and boundedness]
\label{thm:extrema-lebesgue-function}\uses{def:extrema-lebesgue,thm:extrema-cardinal-bounds,thm:harmonic-distance-sum,thm:extrema-cardinal}
\lean{Polynomial.Chebyshev.extremaLebesgueFunction_apply,Polynomial.Chebyshev.extremaLebesgueFunction_zero,Polynomial.Chebyshev.extremaLebesgueFunction_nonneg,Polynomial.Chebyshev.extremaLebesgueFunction_le_harmonic,Polynomial.Chebyshev.bddAbove_extremaLebesgueFunction}\leanok
The function equals the displayed finite sum, is nonnegative, and
$\lambda_0(x)=1$. For $x\in[-1,1]$,
$\lambda_N(x)\le2\pi+2H_N$; in particular its range on that interval is
bounded above. Reference: \cite[Chapter 15]{atap}, with nonsharp constants.
\end{theorem}
\begin{proof}
\leanok
\end{proof}

\begin{theorem}[Lebesgue constant characterisation]
\label{thm:extrema-lebesgue-constant}\uses{def:extrema-lebesgue-constant,thm:classical-real-supremum}
\lean{Polynomial.Chebyshev.extremaLebesgueConstant_isLUB}\leanok
The set $\{\lambda_N(x):x\in[-1,1]\}$ is nonempty and bounded above. The constant $\Lambda_N$ is the least upper bound of
$\{\lambda_N(x):x\in[-1,1]\}$.
\end{theorem}
\begin{proof}
\leanok
\end{proof}

\begin{theorem}[Barycentric interpolation]
\label{thm:lagrange-barycentric}\uses{def:lagrange-interpolant,def:lagrange-nodal-weight,thm:lagrange-interpolation-input,def:barycentric-input}
\lean{Lagrange.classicalBarycentric_denominator_ne_zero,Lagrange.eval_classicalInterpolate_barycentric}\leanok
Define the product weights $\omega_j=\prod_{i\ne j}(x_j-x_i)^{-1}$.
For a nonempty finite set of distinct nodes, away from every node,
\[
 I f(z)=\frac{\sum_j \omega_j f(x_j)/(z-x_j)}
                   {\sum_j \omega_j/(z-x_j)},
\]
and the denominator is nonzero.

Reference: \cite[Theorem 5.1]{atap}.
\end{theorem}
\begin{proof}
\leanok
\end{proof}

\begin{theorem}[Extrema nodal identities]
\label{thm:extrema-nodal}\uses{def:extrema-barycentric,thm:chebyshev-ode,thm:chebyshev-trig}
\lean{Polynomial.Chebyshev.eval_derivative_T_node, Polynomial.Chebyshev.eval_extremaNodal_node, Polynomial.Chebyshev.eval_derivative_extremaNodal, Polynomial.Chebyshev.eval_derivative_extremaNodal_node, Polynomial.Chebyshev.extremaNodal_ne_zero, Polynomial.Chebyshev.natDegree_extremaNodal_le, Polynomial.Chebyshev.extremaNodal_eq_leadingCoeff_mul_nodal, Polynomial.Chebyshev.eval_extremaNodal_cos}\leanok
For $N\ge1$, $F_N$ is a nonzero scalar multiple of the monic nodal polynomial:
\[
 F_N=\operatorname{lc}(F_N)\prod_{j=0}^N(X-x_j),\qquad
 F_N'(x_j)=e_jN^2(-1)^j.
\]
Also $F_N(\cos\theta)=-N\sin\theta\sin(N\theta)$.
The derivative of $T_N$ vanishes at the interior nodes.

Reference: \cite[Chapters 3 and 5]{atap}.
\end{theorem}
\begin{proof}
\leanok
\end{proof}

\begin{theorem}[Cardinal polynomial bounds]
\label{thm:extrema-cardinal-bounds}\uses{thm:extrema-barycentric}
\lean{Polynomial.Chebyshev.sin_mul_abs_sub_le_pi_mul_abs_cos_sub, Polynomial.Chebyshev.node_angle_mem_Icc, Polynomial.Chebyshev.abs_eval_extremaBasis_mul_angle_le, Polynomial.Chebyshev.abs_eval_extremaBasis_cos_le_pi}\leanok
For $N\ge1$, $0\le j\le N$ and $\theta\in[0,\pi]$,
\[
 |\ell_j(\cos\theta)|\le\pi,\qquad
 N\left|\theta-\frac{j\pi}{N}\right|\,|\ell_j(\cos\theta)|
       \le\pi|\sin(N\theta)|.
\]


Reference: \cite[Chapter 15]{atap}; elementary trigonometric estimates.
\end{theorem}
\begin{proof}
\leanok
\end{proof}

\begin{theorem}[Harmonic reciprocal-distance bound]
\label{thm:harmonic-distance-sum}\uses{thm:harmonic-log}
\lean{Finset.sum_le_two_mul_add_harmonic}\leanok
Let $0\le t\le N$, $M\ge0$, and suppose $0\le a_j\le M$ and
$|t-j|a_j\le1$ for $j=0,\ldots,N$. Then
\[
 \sum_{j=0}^N a_j\le2M+2H_N,\qquad H_N=\sum_{j=1}^N\frac1j.
\]

Reference: elementary harmonic-sum estimate used for \cite[Chapter 15]{atap}.
\end{theorem}
\begin{proof}
\leanok
\end{proof}

\begin{theorem}[Harmonic logarithmic bound]
\label{thm:harmonic-log}
\lean{harmonic, harmonic_le_one_add_log_of_pos}\leanok
For $N\ge1$, $H_N\le1+\log N$.
\end{theorem}
\begin{proof}
\leanok
\end{proof}

\end{supporting}

\chapter{Analyticity and geometric convergence}\label{chap:analytic}

\rvsummary{Analyticity in a Bernstein ellipse is equivalent to geometric coefficient decay, which gives geometric convergence.}

Analyticity on a Bernstein ellipse gives geometric coefficient decay and geometric approximation error. The converse theorems recover analyticity from sufficiently fast approximation.

\section{Bernstein ellipses}\label{sec:textbook-analytic-ellipse}

The Joukowski map sends an annulus to an ellipse with foci $-1$ and $1$. This geometry relates complex analyticity to approximation on the interval.

\begin{definition}[Joukowski map]
\label{def:joukowski}
\lean{Complex.joukowski}
\leanok
Define $J:\mathbb C\setminus\{0\}\to\mathbb C$ by
$J(z)=(z+z^{-1})/2$.

Reference: \cite[Chapter 8]{atap}.
\end{definition}

\begin{definition}[Open filled Bernstein ellipse]
\label{def:ellipse}
\lean{Complex.bernsteinEllipse}\leanok
For $\rho>1$, define the open filled ellipse by
\[E_\rho=\{w\in\C:|w-1|+|w+1|<\rho+\rho^{-1}\}.\]

Reference: \cite[Chapter 8]{atap}.
\end{definition}

\begin{theorem}[Ellipse as an annulus image]
\label{thm:ellipse-annulus}\uses{def:ellipse,thm:joukowski-foci,thm:joukowski-surjective}
\lean{Complex.joukowski_mem_bernsteinEllipse_iff, Complex.bernsteinEllipse_eq_image_annulus}\leanok
For $\rho>1$, $E_\rho=J(\{z:\rho^{-1}<|z|<\rho\})$.
For nonzero $z$, membership of $J(z)$ in $E_\rho$ is equivalent to these
annulus inequalities.

Reference: \cite[Chapter 8]{atap}.
\end{theorem}
\begin{proof}
\leanok
\end{proof}

\begin{theorem}[Approximation interval inside the ellipse]
\label{thm:ellipse-interval}\uses{def:ellipse}
\lean{Complex.ofReal_mem_bernsteinEllipse}\leanok
For $\rho>1$, every real $x\in[-1,1]$ belongs to $E_\rho$.

Reference: \cite[Chapter 8]{atap}.
\end{theorem}
\begin{proof}
\leanok
\end{proof}

\begin{supporting}[Ellipse geometry]

\begin{theorem}[Joukowski value and unit-circle characterisation]
\label{thm:joukowski-characterization}\uses{def:joukowski}
\lean{Complex.joukowski_apply, Complex.unitCircle_re_mem, Complex.unitCircle_ne_zero, Complex.joukowski_unitCircle}\leanok
For every nonzero $z$, $J(z)=(z+z^{-1})/2$.
If $|z|=1$, then $z\ne0$, $-1\le\operatorname{Re}z\le1$, and
$J(z)=\operatorname{Re}z$.
Reference: \cite[Chapter 8]{atap}.
\end{theorem}
\begin{proof}
\leanok
\end{proof}

\begin{theorem}[Bernstein ellipse characterisation]
\label{thm:ellipse-characterization}\uses{def:ellipse}
\lean{Complex.mem_bernsteinEllipse_iff}\leanok
For every parameter $\rho>1$ and $w\in\mathbb C$,
\[
w\in E_\rho\quad\Longleftrightarrow\quad
|w-1|+|w+1|<\rho+\rho^{-1}.
\]
Reference: \cite[Chapter 8]{atap}.
\end{theorem}
\begin{proof}
\leanok
\end{proof}

\begin{theorem}[Symmetry of the Joukowski map and openness of the ellipse]
\label{thm:joukowski-ellipse-basic}\uses{def:joukowski,def:ellipse}
\lean{Complex.joukowski_inv, Complex.joukowski_neg, Complex.isOpen_bernsteinEllipse}\leanok
For every $z\ne0$, $J(1/z)=J(z)$ and $J(-z)=-J(z)$. For every $\rho>1$, the set
$E_\rho$ is open.

Reference: \cite[Chapter 8]{atap}.
\end{theorem}
\begin{proof}
\leanok
\end{proof}

\begin{theorem}[Joukowski focal-distance identity]
\label{thm:joukowski-foci}\uses{def:joukowski}
\lean{Complex.joukowski_sub_one, Complex.joukowski_add_one, Complex.norm_joukowski_sub_one_add_norm_joukowski_add_one}
\leanok
For $z\ne0$,
\[|J(z)-1|+|J(z)+1|=|z|+|z|^{-1}.\]

Reference: \cite[Chapter 8]{atap}.
\end{theorem}
\begin{proof}
\leanok
\end{proof}

\begin{theorem}[Nonzero Joukowski preimages]
\label{thm:joukowski-surjective}\uses{def:joukowski}
\lean{Complex.exists_joukowski_eq}\leanok
Every complex $w$ is $J(z)$ for some $z\ne0$.

Reference: \cite[Chapter 8]{atap}.
\end{theorem}
\begin{proof}
\leanok
\end{proof}

\begin{theorem}[Ellipse boundary]
\label{thm:ellipse-boundary}\uses{thm:ellipse-annulus,thm:joukowski-analytic,thm:joukowski-ellipse-basic,thm:open-domain-analytic}
\lean{Complex.image_joukowski_sphere, Complex.joukowski_mem_closure_bernsteinEllipse, Complex.frontier_bernsteinEllipse}\leanok
For $\rho>1$, the boundary of $E_\rho$ is exactly $J(\{|z|=\rho\})$,
which is the level set $|w-1|+|w+1|=\rho+\rho^{-1}$.

Reference: \cite[Chapter 8]{atap}.
\end{theorem}
\begin{proof}
\leanok
\end{proof}

\begin{theorem}[Analyticity of the Joukowski map]
\label{thm:joukowski-analytic}\uses{def:joukowski,def:open-domain-analytic}
\lean{Complex.analyticOnOpen_joukowski, Complex.continuous_joukowski}\leanok
The map $J:\mathbb C\setminus\{0\}\to\mathbb C$ is complex analytic
and continuous on its domain.

Reference: \cite[Chapter 8]{atap}.
\end{theorem}
\begin{proof}
\leanok
\end{proof}

\end{supporting}

\section{Complex coefficients and geometric decay}\label{sec:textbook-analytic-decay}

Complex coefficients use the same full constant convention as real coefficients. The contour formula uses $\mathcal I_{0,1}$, the integral around the unit circle from Definition~\ref{def:circle-integral-input}.

\begin{definition}[Inputs for a fixed complex Chebyshev coefficient]
\label{def:complex-coefficient-input}\uses{thm:complex-integrable-input,def:coefficient-input}
\lean{Polynomial.Chebyshev.ComplexCoefficientInput, Polynomial.Chebyshev.ComplexCoefficientInput.ofL1, Polynomial.Chebyshev.ComplexCoefficientInput.ofContinuousOn, Polynomial.Chebyshev.ComplexCoefficientInput.ofReal, Polynomial.Chebyshev.ComplexCoefficientInput.re, Polynomial.Chebyshev.ComplexCoefficientInput.im}\leanok
For $n\in\N$, define
\[
\mathcal I_n^{\C}=\{f:\R\to\C: fT_n\text{ is }\mu_T\text{-integrable}\},
\]
retaining pointwise function values. Real functions are regarded as complex
functions by scalar inclusion; real and imaginary parts are taken pointwise.
Reference: \cite[Chapter 3]{atap}.
\end{definition}

\begin{definition}[Complex Chebyshev coefficients]
\label{def:complex-coefficients}\uses{def:complex-coefficient-input}
\lean{Polynomial.Chebyshev.complexCoefficient}\leanok
For $f\in\mathcal I_n^{\C}$, define $a_n(f)$ by the weighted integrals of
Definition~\ref{def:coefficients}, with the same full constant coefficient
convention: $a_0(f)=\pi^{-1}\int f\,d\mu_T$ and $a_n(f)=2\pi^{-1}\int fT_n\,d\mu_T$
for $n>0$.

Reference: \cite[Chapters 3 and 8]{atap}.
\end{definition}

\begin{definition}[Laurent integrand of continuous interval data]
\label{def:joukowski-laurent-input}\uses{thm:joukowski-characterization,def:circle-integral-input}
\lean{Complex.joukowskiLaurentInput}\leanok
For $f\in C([-1,1],\mathbb C)$ and $n\ge0$, define the integrable function on the unit circle
\[
L_n[f](z)=\frac{f(\operatorname{Re}z)}{z^{n+1}},\qquad |z|=1.
\]
Its numerator is the Joukowski pullback of the interval data.
Reference: \cite[Theorems 3.1 and 8.1]{atap}.
\end{definition}

\begin{theorem}[Chebyshev coefficient decay]
\label{thm:coefficient-decay}\uses{thm:joukowski-coefficient-bound,thm:ellipse-interval,thm:complex-coefficients-basic,thm:circle-integral-input,thm:joukowski-laurent-input}
\lean{Polynomial.Chebyshev.complexCoefficient_eq_laurent_integral, Polynomial.Chebyshev.norm_complexCoefficient_le_of_analytic}\leanok
For a complex-valued function whose restriction to $[-1,1]$ is continuous,
the coefficient has the contour representation
\[
a_n(f)=\frac{\kappa_n}{2\pi i}\mathcal I_{0,1}(L_n[f|_{[-1,1]}]),
\qquad \kappa_0=1,\quad \kappa_n=2\ (n>0).
\]
Let $\rho>1$. If $f$ is complex analytic on $E_\rho$ and $|f(z)|\le M$ there, then
$|a_n(f)|\le2M\rho^{-n}$ for $n\ge1$, and $|a_0(f)|\le M$.
No boundary continuity is assumed.

Reference: \cite[Theorem 8.1]{atap}.
\end{theorem}
\begin{proof}
\leanok
\end{proof}

\begin{theorem}[Geometric truncation and interpolation error]
\label{thm:geometric-error}\uses{thm:coefficient-decay,thm:coefficient-tail,thm:complex-series-representation,thm:interpolation-tail}
\lean{Polynomial.Chebyshev.norm_sub_sum_complex_le_tsum_tail, Polynomial.Chebyshev.norm_sub_sum_complex_le_geometric, Polynomial.Chebyshev.norm_sub_sum_complexCoefficient_le_of_analytic, Polynomial.Chebyshev.norm_sub_extremaInterpolant_le_of_analytic}\leanok
For $\rho>1$, let $f$ be analytic on the open ellipse $E_\rho$ and satisfy $|f(z)|\le M$ there. Then the degree-$N$ truncated Chebyshev expansion
has uniform error at most $2M\rho^{-N}/(\rho-1)$ on $[-1,1]$.
For $N\ge1$, interpolation at the $N+1$ extrema $\cos(j\pi/N)$ has
error at most $4M\rho^{-N}/(\rho-1)$. No continuity on the ellipse boundary
is assumed.

Reference: \cite[Theorem 8.2]{atap}.
\end{theorem}
\begin{proof}
\leanok
\end{proof}

\paragraph{Example (Julia).} For $f(x)=1/(1+25x^2)$, the poles $\pm i/5$ lie on the ellipse with $\rho=(1+\sqrt{26})/5$. Using \texttt{chebcoeffs} from Chapter~\ref{chap:interpolation}, inspect the geometric decay of the even coefficients. Odd coefficients vanish in exact arithmetic. Rescaling high-degree numerical coefficients also magnifies their rounding errors.

\begin{juliacode}
c = chebcoeffs(x -> 1/(1 + 25x^2), 200)
rho = 0.2 + sqrt(1.04)
[abs(c[n+1]) * rho^n for n in 20:20:80]   # inspect the rescaled even coefficients
\end{juliacode}

\begin{supporting}[Complex coefficient identities and contour estimates]

\begin{definition}[Integrable complex interval data]
\label{def:complex-integrable-input}\uses{def:measureT}
\lean{Polynomial.Chebyshev.ComplexIntegrableInput}\leanok
Let $\mathcal L^1_T(\C)$ be the set of functions $f:\R\to\C$ that are
$\mu_T$-integrable, retaining their pointwise representatives.
Reference: \cite[Chapter 3]{atap}.
\end{definition}

\begin{theorem}[Integrability of complex interval data]
\label{thm:complex-integrable-input}\uses{def:complex-integrable-input,thm:T-integrable-transport}
\lean{Polynomial.Chebyshev.ComplexIntegrableInput.integrable, Polynomial.Chebyshev.integrable_classicalMeasureT_complex}\leanok
Every $f\in\mathcal L^1_T(\C)$ is $\mu_T$-integrable.
If $f:\R\to\C$ is continuous on $[-1,1]$, then $f$ is $\mu_T$-integrable
and hence belongs to $\mathcal L^1_T(\C)$.
\end{theorem}
\begin{proof}
\leanok
\end{proof}

\begin{theorem}[Characterisation of complex coefficient inputs]
\label{thm:complex-coefficient-input}\uses{def:complex-coefficient-input,thm:chebyshev-recurrence}
\lean{Polynomial.Chebyshev.ComplexCoefficientInput.integrable, Polynomial.Chebyshev.ComplexCoefficientInput.ofL1_val, Polynomial.Chebyshev.ComplexCoefficientInput.ofContinuousOn_val, Polynomial.Chebyshev.ComplexCoefficientInput.ofReal_val, Polynomial.Chebyshev.ComplexCoefficientInput.re_val, Polynomial.Chebyshev.ComplexCoefficientInput.im_val, Polynomial.Chebyshev.complexCoefficientInput_zero_iff}\leanok
For $f\in\mathcal I_n^{\C}$, the product $fT_n$ is $\mu_T$-integrable.
Every $\mu_T$-integrable complex function, and every $f:\R\to\C$ continuous
on $[-1,1]$, belongs to $\mathcal I_n^{\C}$ for every $n$.
Scalar inclusion maps $\mathcal I_n$ into $\mathcal I_n^{\C}$, and real and
imaginary parts map $\mathcal I_n^{\C}$ into $\mathcal I_n$.
All these operations preserve the indicated pointwise values.
Membership of $\mathcal I_0^{\C}$ is exactly $\mu_T$-integrability of $f$.
\end{theorem}
\begin{proof}
\leanok
\end{proof}

\begin{theorem}[Classical complex coefficient formula]
\label{thm:complex-coefficient-characterisation}\uses{def:complex-coefficients}
\lean{Polynomial.Chebyshev.complexCoefficient_apply}\leanok
For $f\in\mathcal I_n^{\C}$,
\[a_n(f)=\frac{c_n}{\pi}\int fT_n\,d\mu_T,
\qquad c_0=1,\quad c_n=2\quad(n>0).\]
The product $fT_n$ is integrable, so this integral is finite.
\end{theorem}
\begin{proof}
\leanok
\end{proof}

\begin{theorem}[Real case and cosine formula]
\label{thm:complex-coefficients-basic}\uses{def:coefficients,thm:complex-coefficient-input,thm:chebyshev-trig,def:complex-coefficients}
\lean{Polynomial.Chebyshev.complexCoefficient_ofReal, Polynomial.Chebyshev.complexCoefficient_eq_integral_cos}\leanok
For $f\in\mathcal I_n$, the complex coefficient at index $n$
equals its real coefficient from Definition~\ref{def:coefficients}. For $f$ continuous on $[-1,1]$,
\[a_n(f)=\frac{c_n}{\pi}\int_0^\pi f(\cos\theta)\cos(n\theta)\,d\theta,
\qquad c_0=1,\quad c_n=2\ (n>0).\]

Reference: \cite[Chapters 3 and 8]{atap}.
\end{theorem}
\begin{proof}
\leanok
\end{proof}

\begin{theorem}[Values of the unit-circle Laurent integrand]
\label{thm:joukowski-laurent-input}\uses{def:joukowski-laurent-input}
\lean{Complex.joukowskiLaurentInput_apply}\leanok
For $f\in C([-1,1],\mathbb C)$, $n\ge0$ and $|z|=1$,
$L_n[f](z)=f(\operatorname{Re}z)/z^{n+1}$.
\end{theorem}
\begin{proof}
\leanok
\end{proof}

\begin{theorem}[Complex Chebyshev series representation]
\label{thm:complex-series-representation}\uses{def:complex-coefficients,thm:series-representation,thm:complex-coefficient-input}
\lean{Polynomial.Chebyshev.complexCoefficient_re, Polynomial.Chebyshev.complexCoefficient_im, Polynomial.Chebyshev.hasSum_complexCoefficient_mul_eval_T}\leanok
If $f:[-1,1]\to\mathbb C$ is continuous and $\sum_{n\ge0}|a_n(f)|<\infty$,
then $f(x)=\sum_{n\ge0}a_n(f)T_n(x)$ at every $x\in[-1,1]$, including the
endpoints. The real and imaginary parts of $a_n(f)$ are the corresponding
real Chebyshev coefficients.

Reference: \cite[Chapters 4 and 8]{atap}.
\end{theorem}
\begin{proof}
\leanok
\end{proof}

\begin{theorem}[Open-annulus Laurent coefficient bound]
\label{thm:laurent-bound}\uses{thm:circle-integral-bound}
\lean{Complex.norm_laurent_integral_le, Complex.laurent_integral_eq_of_annulus, Complex.norm_laurent_integral_le_of_open_annulus}\leanok
For $r>0$, let $q_r\in\mathcal C_{0,r}$ satisfy
$|z^{n+1}q_r(z)|\le M$ on $|z|=r$. Then
\[
\left|\frac{\mathcal I_{0,r}(q_r)}{2\pi i}\right|\le M r^{-n}.
\]
If $g$ is analytic on the closed annulus $r\le|z|\le R$, where $0<r\le R$,
and $q_r(z)=g(z)/z^{n+1}$ and $q_R(z)=g(z)/z^{n+1}$ on their respective
circles, then $\mathcal I_{0,r}(q_r)=\mathcal I_{0,R}(q_R)$.
Finally, let $\rho>1$. If $g$ is analytic at every $z$ with
$1\le|z|<\rho$ and $|g(z)|\le M$ there, then every $q\in\mathcal C_{0,1}$ with $q(z)=g(z)/z^{n+1}$ satisfies
\[
\left|\frac{\mathcal I_{0,1}(q)}{2\pi i}\right|\le M\rho^{-n}.
\]
Here $n\ge0$; no analyticity inside the inner circle is required.

Reference: \cite[proof of Theorem 8.1]{atap}.
\end{theorem}
\begin{proof}
\leanok
\end{proof}

\begin{theorem}[Joukowski pullback coefficient bound]
\label{thm:joukowski-coefficient-bound}\uses{thm:laurent-bound,thm:ellipse-annulus,thm:joukowski-analytic,thm:joukowski-characterization,thm:open-domain-analytic}
\lean{Complex.norm_joukowski_laurent_integral_le}\leanok
For $\rho>1$ and $f$ analytic and bounded by $M$ in $E_\rho$, let $q\in\mathcal C_{0,1}$ satisfy $q(z)=f(J(z))/z^{n+1}$, $n\ge0$.
Then
\[
\left|\frac{\mathcal I_{0,1}(q)}{2\pi i}\right|\le M\rho^{-n}.
\]

Reference: \cite[proof of Theorem 8.1]{atap}.
\end{theorem}
\begin{proof}
\leanok
\end{proof}

\end{supporting}

\section{Converses and analytic extensions}\label{sec:textbook-analytic-converse}

An analytic extension gives decay at every smaller ellipse radius. Conversely, those decay estimates determine an analytic extension.

\begin{definition}[Complex Chebyshev series]
\label{def:complex-chebyshev-series}\uses{def:chebyshev-mathlib}
\lean{Polynomial.Chebyshev.ComplexChebyshevSeriesPoint, Polynomial.Chebyshev.complexChebyshevSeries}\leanok
For $a:\mathbb N\to\mathbb C$, let $D_a$ consist of the complex points $z$
where the natural-order partial sums $\sum_{n<N}a_nT_n(z)$ converge to a finite
complex limit. Define $S_a:D_a\to\mathbb C$ to be that limit. This includes
conditionally convergent series; no value is assigned elsewhere.

Reference: \cite[Chapters 3 and 8]{atap}.
\end{definition}

\begin{theorem}[Analytic extension characterized by smaller-radius decay]
\label{thm:analytic-decay-characterization}\uses{thm:coefficient-decay,thm:coefficient-decay-converse}
\lean{Polynomial.Chebyshev.complexCoefficient_congr_of_continuousOn, Polynomial.Chebyshev.coefficient_decay_of_analytic_extension, Polynomial.Chebyshev.exists_analytic_extension_iff_coefficient_decay}\leanok
For continuous complex data $f$ on $[-1,1]$ and $\rho>1$, an analytic
extension to $E_\rho$ exists if and only if, for every $1<r<\rho$, there is
$C_r$ such that $|a_n(f)|\le C_r r^{-n}$ for every $n\ge0$.
The constant may depend on $r$.

Reference: \cite[Theorems 8.1--8.3]{atap}.
\end{theorem}
\begin{proof}
\leanok
\end{proof}

\begin{theorem}[Bernstein converse from polynomial approximation]
\label{thm:bernstein-converse}\uses{thm:complex-coefficient-support,thm:coefficient-decay-converse}
\lean{Polynomial.Chebyshev.continuousOn_of_geometric_polynomial_approximation, Polynomial.Chebyshev.coefficient_decay_of_geometric_polynomial_approximation, Polynomial.Chebyshev.exists_analytic_extension_of_geometric_polynomial_approximation}\leanok
Let $\rho>1$ and suppose complex polynomials $p_n$ of degree at most $n$
satisfy $|f(x)-p_n(x)|\le C\rho^{-n}$ for every $n\ge0$ and $x\in[-1,1]$.
Then $f$ has an analytic extension to the open ellipse $E_\rho$.
Continuity of $f$ need not be assumed.

Reference: \cite[Theorem 8.3]{atap}.
\end{theorem}
\begin{proof}
\leanok
\end{proof}

\begin{supporting}[Construction of analytic extensions]

\begin{theorem}[Classical Chebyshev sum characterisation]
\label{thm:complex-chebyshev-series-value}\uses{def:complex-chebyshev-series}
\lean{Polynomial.Chebyshev.tendsto_complexChebyshevSeries, Polynomial.Chebyshev.complexChebyshevSeries_eq_of_tendsto, Polynomial.Chebyshev.ComplexChebyshevSeriesPoint.ofHasSum, Polynomial.Chebyshev.ComplexChebyshevSeriesPoint.ofSummable, Polynomial.Chebyshev.complexChebyshevSeries_ofHasSum}\leanok
For $z\in D_a$, the natural-order partial sums converge to $S_a(z)$, and any
finite limit of those partial sums equals $S_a(z)$. Unconditional summability
supplies a point of $D_a$; if the unconditional sum is $s$, then $S_a(z)=s$.
\end{theorem}
\begin{proof}
\leanok
\end{proof}

\begin{theorem}[Compact smaller ellipses]
\label{thm:ellipse-exhaustion}\uses{thm:ellipse-annulus}
\lean{Complex.exists_smaller_bernsteinEllipse, Complex.closure_bernsteinEllipse_subset, Complex.isCompact_closure_bernsteinEllipse}\leanok
Every point of $E_\rho$, $\rho>1$, belongs to some $E_r$ with $1<r<\rho$.
For these radii, $\overline{E_r}$ is compact and contained in $E_\rho$.

Reference: \cite[Chapter 8]{atap}.
\end{theorem}
\begin{proof}
\leanok
\end{proof}

\begin{theorem}[Complex Chebyshev growth]
\label{thm:chebyshev-ellipse-growth}\uses{thm:ellipse-annulus,thm:chebyshev-recurrence,def:chebyshev-mathlib}
\lean{Polynomial.Chebyshev.eval_T_joukowski, Polynomial.Chebyshev.norm_eval_T_le_on_bernsteinEllipse}\leanok
For $w\ne0$ and $n\ge0$,
$T_n(J(w))=(w^n+w^{-n})/2$. Consequently $|T_n(z)|\le r^n$
for $r>1$ and $z\in E_r$.

Reference: \cite[Exercise 8.8]{atap}.
\end{theorem}
\begin{proof}
\leanok
\end{proof}

\begin{theorem}[Analyticity from geometric coefficient decay]
\label{thm:analytic-chebyshev-series}\uses{thm:chebyshev-ellipse-growth,thm:ellipse-exhaustion,thm:joukowski-ellipse-basic,thm:complex-chebyshev-series-value,thm:open-domain-analytic}
\lean{Polynomial.Chebyshev.summable_norm_chebyshevSeries, Polynomial.Chebyshev.summable_chebyshevSeries_of_geometric, Polynomial.Chebyshev.ComplexChebyshevSeriesPoint.ofGeometric, Polynomial.Chebyshev.ComplexChebyshevSeriesPoint.ofGeometric_val, Polynomial.Chebyshev.norm_complexChebyshevSeries_le, Polynomial.Chebyshev.analyticOnOpen_complexChebyshevSeries}\leanok
If $\rho>1$ and $|a_n|\le C\rho^{-n}$ for all $n\ge0$, then $S_a$ is
analytic on the open ellipse $E_\rho$. For $1<r<\rho$ and $z\in E_r$, the
series converges absolutely and $|S_a(z)|\le C/(1-r/\rho)$.
Boundedness on all of $E_\rho$ is not claimed.

Reference: \cite[Theorem 8.3]{atap}.
\end{theorem}
\begin{proof}
\leanok
\end{proof}

\begin{theorem}[Coefficient-decay converse]
\label{thm:coefficient-decay-converse}\uses{thm:analytic-chebyshev-series,thm:complex-series-representation}
\lean{Polynomial.Chebyshev.exists_analytic_extension_of_coefficient_decay}\leanok
Let $f:[-1,1]\to\mathbb C$ be continuous and $\rho>1$. If its full-constant
Chebyshev coefficients satisfy $|a_n(f)|\le C\rho^{-n}$, then $f$ has an
analytic extension to $E_\rho$, agreeing at every interval point.

Reference: \cite[Theorems 8.1--8.3]{atap}.
\end{theorem}
\begin{proof}
\leanok
\end{proof}

\begin{theorem}[Complex coefficient bounds and polynomial support]
\label{thm:complex-coefficient-support}\uses{thm:sequence-finite-coordinates,thm:complex-coefficients-basic,thm:coefficient-T}
\lean{Polynomial.Chebyshev.norm_complexCoefficient_le, Polynomial.Chebyshev.complexCoefficient_sub, Polynomial.Chebyshev.complexCoefficient_const_mul, Polynomial.Chebyshev.complexCoefficient_sum, Polynomial.Chebyshev.complexCoefficient_eval_T, Polynomial.Chebyshev.complexCoefficient_polynomial_eq_zero}\leanok
For continuous complex $f$ on $[-1,1]$, $|f|\le M$, $M\ge0$, implies
$|a_0(f)|\le M$ and $|a_n(f)|\le2M$ for $n>0$. Coefficients commute with
subtraction, complex scalar multiplication, and finite sums of continuous
functions. Moreover $a_n(T_m)=\delta_{mn}$, and $\deg p<n$ implies $a_n(p)=0$
for complex polynomials.

Reference: \cite[Chapters 3 and 8]{atap}.
\end{theorem}
\begin{proof}
\leanok
\end{proof}

\end{supporting}

\chapter{Best approximation}\label{chap:best}

\rvsummary{Existence, equioscillation and uniqueness of best uniform approximations, and near-best Chebyshev interpolation.}

Best uniform approximation provides a benchmark for interpolation. For real functions, equioscillation characterizes the unique minimizing polynomial.

\section{Best polynomial error}\label{sec:textbook-best-best}

The best error is the benchmark against which a polynomial approximation can be compared. The first theorem shows that its infimum is attained.

\begin{definition}[Best uniform polynomial error]
\label{def:best-uniform-error}\uses{def:extrema-cardinal,thm:extrema-grid}
\lean{Polynomial.Chebyshev.polynomialOnInterval, Polynomial.Chebyshev.bestUniformError, Polynomial.Chebyshev.interpolationOnInterval}\leanok
For $f\in C([-1,1],\mathbb C)$, define
\[
 E_N(f)=\inf_{\deg p\le N}\|f-p\|_\infty,
\]
the infimum over complex polynomials $p$ of degree at most $N$, restricted to
$[-1,1]$. Here $\|\cdot\|_\infty$ is the norm of $C([-1,1],\mathbb C)$. The interpolant $I_N f$ is the cardinal sum $\sum_{j=0}^N f(x_j)\ell_j$, evaluating $f$ only at its extrema nodes in $[-1,1]$.

Reference: \cite[Chapter 16]{atap}.
\end{definition}

\begin{theorem}[Attainment of best uniform polynomial error]
\label{thm:best-error-attainment}\uses{thm:finite-dimensional-best-approximation,thm:best-uniform-error}
\lean{Polynomial.Chebyshev.polynomialOnIntervalLinear, Polynomial.Chebyshev.exists_bestUniformError, Polynomial.Chebyshev.exists_real_bestApproximation, Polynomial.exists_bestApproximation}\leanok
The complex best-error infimum $E_N(f)$ on $[-1,1]$ is attained by a polynomial
of degree at most $N$. For continuous real data on any compact interval $[a,b]$,
a real degree-at-most-$N$ minimizer exists as well. Existence alone does not
require $a<b$.

Reference: \cite[Theorem 10.1]{atap}.
\end{theorem}
\begin{proof}
\leanok
\end{proof}

\begin{supporting}[Infimum bounds]

\begin{theorem}[Best-error infimum bounds]
\label{thm:best-uniform-error}\uses{def:best-uniform-error}
\lean{Polynomial.Chebyshev.bestUniformError_nonneg, Polynomial.Chebyshev.bestUniformError_le}\leanok
The competitor-error set is nonempty and bounded below by zero. Thus
$0\le E_N(f)\le\|f-p\|_\infty$ for every polynomial of degree at most $N$.


Reference: \cite[Chapter 16]{atap}.
\end{theorem}
\begin{proof}
\leanok
\end{proof}

\end{supporting}

\section{Equioscillation}\label{sec:textbook-best-alternation}

For real data, a pattern of alternating extreme errors characterizes the best approximant. The lower-bound theorem also applies before the extreme errors have equal magnitude.

\begin{definition}[Equioscillating error]
\label{def:equioscillation}
\lean{Polynomial.Equioscillates}\leanok
A real continuous error $e$ on $[a,b]$ equioscillates for degree $N$ if there
are $a\le x_0<\cdots<x_{N+1}\le b$ and $\sigma\in\{-1,1\}$ such that
$e(x_i)=\sigma(-1)^i\|e\|_\infty$. The zero-error case is included.

Reference: \cite[Theorem 10.1]{atap}.
\end{definition}

\begin{theorem}[De la Vall\'ee Poussin bound and sufficient alternation]
\label{thm:alternation-sufficient}\uses{def:equioscillation,thm:alternation-root-count}
\lean{Polynomial.lower_bound_of_alternating_error, Polynomial.Equioscillates.norm_sub_le}\leanok
Let $p$ have degree at most $N$. If $N+2$ increasing points satisfy
$\sigma(-1)^i(f(x_i)-p(x_i))\ge L$, where $|\sigma|=1$, then every
polynomial $q$ of degree at most $N$ satisfies $L\le\|f-q\|_\infty$.
In particular, an equioscillating error certifies a best approximant.

Reference: \cite[Theorem 10.1 and Exercise 10.3]{atap}.
\end{theorem}
\begin{proof}
\leanok
\end{proof}

For example, the error in approximating $x^2$ by $1/2$ takes the values $1/2,-1/2,1/2$ at $-1,0,1$. The alternation criterion certifies this best constant approximation.

\begin{theorem}[Real equioscillation characterization]
\label{thm:equioscillation-characterization}\uses{thm:alternation-sufficient,thm:extreme-moment-certificate,thm:finite-annihilator-alternation}
\lean{Polynomial.equioscillates_of_bestApproximation, Polynomial.equioscillates_iff_bestApproximation}\leanok
Let $a<b$, let $f\in C([a,b],\mathbb R)$, and let $p$ have degree at most $N$.
Then $p$ minimizes the uniform error among all real polynomials of degree at
most $N$ if and only if $f-p$ equioscillates at $N+2$ increasing points.
This includes zero error.

Reference: \cite[Theorem 10.1]{atap}.
\end{theorem}
\begin{proof}
\leanok
\end{proof}

\begin{theorem}[Unique real minimax polynomial]
\label{thm:unique-real-minimax}\uses{thm:best-error-attainment,thm:equioscillation-characterization}
\lean{Polynomial.bestApproximation_unique, Polynomial.existsUnique_bestApproximation, Polynomial.existsUnique_equioscillating_approximation}\leanok
For every $a<b$, $N\ge0$, and continuous real $f$ on $[a,b]$, there is a
unique best polynomial of degree at most $N$. Equivalently, there is a unique
such polynomial whose error equioscillates.

Reference: \cite[Theorem 10.1 and Exercise 10.10]{atap}.
\end{theorem}
\begin{proof}
\leanok
\end{proof}

\rvnote{thm:unique-real-minimax}{No complex uniqueness or complex sign-alternation characterization is asserted here.}

\section{Comparison with interpolation}\label{sec:textbook-best-near-best}

The Lebesgue constant compares Chebyshev interpolation with the best polynomial of the same degree.

\begin{theorem}[Uniform near-best interpolation against the infimum]
\label{thm:infimum-near-best}\uses{thm:best-uniform-error,thm:extrema-near-best}
\lean{Polynomial.Chebyshev.norm_sub_interpolationOnInterval_le_bestUniformError, Polynomial.Chebyshev.norm_sub_interpolationOnInterval_le_log_bestUniformError}\leanok
For $f\in C([-1,1],\mathbb C)$,
\[
 \|f-I_Nf\|_\infty\le(1+\Lambda_N)E_N(f).
\]
For $N\ge1$, the proved logarithmic estimate also gives the factor
$2\pi+3+2\log N$ in place of $1+\Lambda_N$.


Reference: \cite[Theorem 15.1 and Chapter 16]{atap}.
\end{theorem}
\begin{proof}
\leanok
\end{proof}

\begin{supporting}[Best-error converse and convex certificates]

\begin{theorem}[Bernstein converse for the best error]
\label{thm:best-error-bernstein-converse}\uses{thm:bernstein-converse,thm:best-error-attainment}
\lean{Polynomial.Chebyshev.exists_analytic_extension_of_bestUniformError}\leanok
For $f\in C([-1,1],\mathbb C)$ and $\rho>1$, if $E_n(f)\le C\rho^{-n}$
for every $n\ge0$, then $f$ has an analytic extension to $E_\rho$.

Reference: \cite[Theorem 8.3]{atap}.
\end{theorem}
\begin{proof}
\leanok
\end{proof}

\begin{theorem}[Finite-dimensional best approximation]
\label{thm:finite-dimensional-best-approximation}
\lean{Submodule.exists_norm_sub_le_of_finiteDimensional}\leanok
Let $V$ be a finite-dimensional subspace of a normed space over a proper,
nontrivially normed field. For every $f$, there exists $p\in V$ such that
$\|f-p\|\le\|f-q\|$ for every $q\in V$. No inner product or uniqueness
assumption is needed.

Reference: \cite[Theorem 10.1]{atap}, existence argument.
\end{theorem}
\begin{proof}
\leanok
\end{proof}

\begin{theorem}[Alternating signs force polynomial degree]
\label{thm:alternation-root-count}
\lean{Polynomial.exists_eval_eq_zero_between_of_mul_neg, Polynomial.le_natDegree_of_alternating_eval}\leanok
If a real polynomial changes sign across each of $n$ consecutive intervals
with strictly increasing endpoints, then its natural degree is at least $n$.
Each sign change gives a root in the open interval.

Reference: \cite[proof of Theorem 10.1]{atap}.
\end{theorem}
\begin{proof}
\leanok
\end{proof}

\begin{theorem}[Convex separation and Carath\'eodory certificates]
\label{thm:mathlib-convex-certificates}
\lean{geometric_hahn_banach_point_closed, eq_pos_convex_span_of_mem_convexHull}\mathlibok
A point outside a closed convex set in a real locally convex space can be
strictly separated from that set by a continuous linear functional.
A point in a convex hull has a positive convex combination supported on an
affinely independent finite family.
\end{theorem}
\begin{proof}
\mathlibok
\end{proof}

\begin{theorem}[Compact finite-dimensional convex hulls]
\label{thm:compact-convex-hull}\uses{thm:mathlib-convex-certificates}
\lean{mem_convexHull_exists_fin_combination, IsCompact.convexHull_finiteDimensional}\leanok
In dimension $d$, every convex-hull point is a convex combination of at most
$d+1$ points. The convex hull of a compact set in a finite-dimensional real
normed space is compact.
\end{theorem}
\begin{proof}
\leanok
\end{proof}

\begin{theorem}[Strict improvement from the extreme set]
\label{thm:uniform-extreme-improvement}
\lean{ContinuousMap.exists_norm_sub_smul_lt_of_extreme_sign}\leanok
Let $e,q$ be continuous real functions on a compact space, with $\|e\|>0$.
If $e(x)q(x)>0$ at every point where $|e(x)|=\|e\|$, there exists a real
$t>0$ such that $\|e-tq\|<\|e\|$.

Reference: \cite[Theorem 10.1]{atap}.
\end{theorem}
\begin{proof}
\leanok
\end{proof}

\begin{definition}[Extreme moment vectors]
\label{def:extreme-moment}
\lean{Polynomial.extremeMoment}\leanok
For a real error function $e$ on $[a,b]$, set
$v_N(x)=e(x)(1,x,\ldots,x^N)\in\mathbb R^{N+1}$.


Reference: \cite[Theorem 10.1]{atap}, convex-separation formulation.
\end{definition}

\rvnote{def:extreme-moment}{These are the signed evaluation vectors used in the minimax certificate.}

\begin{theorem}[Finite positive certificate for a minimax error]
\label{thm:extreme-moment-certificate}\uses{def:extreme-moment,thm:uniform-extreme-improvement,thm:compact-convex-hull}
\lean{Polynomial.zero_mem_convexHull_extremeMoment, Polynomial.exists_extremeMoment_certificate}\leanok
Suppose $e$ is continuous and real, $\|e\|>0$, and
$\|e\|\le\|e-q\|$ for every polynomial $q$ of degree at most $N$.
Then zero lies in the convex hull of $v_N(x)$ over the extreme set
$\{x:|e(x)|=\|e\|\}$. There are at most $N+2$ distinct extreme points
$x_i$ and positive weights $w_i$ with $\sum_iw_i=1$ and
$\sum_iw_i v_N(x_i)=0$.

Reference: \cite[Theorem 10.1]{atap}.
\end{theorem}
\begin{proof}
\leanok
\end{proof}

\begin{theorem}[Alternation of finite annihilating weights]
\label{thm:finite-annihilator-alternation}
\lean{Polynomial.sum_mul_eval_eq_zero_of_moments, Polynomial.exists_strictMono_reindex, Polynomial.card_ge_of_nonzero_annihilating_weights, Polynomial.mul_weights_neg_of_annihilating}\leanok
Vanishing moments through degree $N$ imply annihilation of all polynomials of
that degree. A nonempty family of distinct nodes carrying nonzero annihilating
weights must have at least $N+2$ nodes. If there are exactly $N+2$ increasing
nodes, adjacent weights have opposite signs. Finite injective node families
can be reindexed in increasing order.

Reference: \cite[Theorem 10.1]{atap}.
\end{theorem}
\begin{proof}
\leanok
\end{proof}

\end{supporting}

\chapter{Calculus and arithmetic on Chebyshev series}\label{chap:calculus}

\rvsummary{Clenshaw evaluation, differentiation, integration and multiplication carried out exactly on Chebyshev coefficients.}

A coefficient representation supports evaluation, differentiation, integration and multiplication. Finite formulas come first; convergence hypotheses justify the operations on infinite series.

\section{Clenshaw evaluation}\label{sec:textbook-calculus-evaluation}

A backward recurrence evaluates an expansion without forming every basis polynomial separately. The state has two entries because the basis satisfies a three-term recurrence.

\begin{definition}[Backward Clenshaw state]
\label{def:clenshaw}
\lean{Polynomial.clenshawState, Polynomial.clenshaw}\leanok
Let $R$ be a ring, not necessarily commutative, and let $a,b,c:\N\to R$.
For $n$ coefficients starting at index $k$, define $S_{0,k}=(0,0)$.
If $S_{n,k+1}=(u,v)$, set
\[
 S_{n+1,k}=(c_k+b_kv,\ u+a_kv).
\]
If $S_{n,0}=(u,v)$, the evaluator returns $u p_0+v p_1$.
Here $n$ is the number of coefficients, not the largest degree.
The initial values $p_0,p_1$ are supplied separately.

References: \cite{clenshaw}; \cite[3.11(ii), 18.40(i)]{dlmf}.
\end{definition}

\begin{theorem}[Clenshaw evaluation]
\label{thm:clenshaw}\uses{def:clenshaw}
\lean{Polynomial.clenshawState_spec, Polynomial.clenshaw_eq_sum, Polynomial.clenshaw_zero, Polynomial.clenshaw_one}\leanok
Let $R$ be a commutative ring. Suppose $p_{k+2}=a_kp_{k+1}+b_kp_k$
for every $k\in\N$.
Then the evaluator returns exactly
\[
 \sum_{i=0}^{n-1}c_i p_i.
\]
More generally, if $S_{n,k}=(u,v)$, then
$u p_k+v p_{k+1}=\sum_{i=0}^{n-1}c_{k+i}p_{k+i}$.
The empty expansion is zero and a one-coefficient expansion is $c_0p_0$.
No division or orthogonality is assumed.

References: \cite{clenshaw}; \cite[18.40(i)]{dlmf}.
\end{theorem}
\begin{proof}
\leanok
\end{proof}

\rvnote{thm:clenshaw}{This is an exact-arithmetic theorem; no floating-point error or stability estimate is asserted.}

\begin{theorem}[Chebyshev Clenshaw evaluation]
\label{thm:chebyshev-clenshaw}\uses{thm:clenshaw,thm:chebyshev-recurrence,def:chebyshev-mathlib}
\lean{Polynomial.Chebyshev.clenshawT, Polynomial.Chebyshev.clenshawT_eq_sum}\leanok
For any commutative ring and any $x$, use $a_k=2x$, $b_k=-1$,
$p_0=1$ and $p_1=x$. The resulting executable evaluator returns
$\sum_{i=0}^{n-1}c_iT_i(x)$.
The coefficient $c_0$ is the full constant term; unlike \cite[3.11(ii)]{dlmf},
it is not halved.

References: \cite{clenshaw}; \cite[3.11(ii)]{dlmf}.
\end{theorem}
\begin{proof}
\leanok
\end{proof}

\paragraph{Example (Julia).} The backward state can be compared with direct cosine summation at an interval point. Both evaluations use the full constant coefficient.

\begin{juliacode}
function clenshaw(c, a, b, p0, p1)  # backward state of the definition
    u, v = zero(p0), zero(p0)
    for k in length(c)-1:-1:0
        u, v = c[k+1] + b(k)*v, u + a(k)*v
    end
    u*p0 + v*p1
end
c = [1/(1 + k^2) for k in 0:20]
x = 0.3
direct = sum(c[k+1]*cos(k*acos(x)) for k in 0:20)
clenshaw(c, k -> 2x, k -> -1.0, 1.0, x) - direct   # compare the two evaluations
\end{juliacode}

\begin{supporting}[Finite-array and stream formulations]

\begin{theorem}[Chebyshev recurrence]
\label{thm:chebyshev-recurrence}
\lean{Polynomial.Chebyshev.T_zero, Polynomial.Chebyshev.T_one, Polynomial.Chebyshev.T_add_two}\mathlibok
Over a commutative ring, $T_0=1$, $T_1=X$ and
$T_{k+2}=2XT_{k+1}-T_k$ for every integer $k$.

Reference: \cite[Table 18.9.1]{dlmf}.
\end{theorem}
\begin{proof}
\mathlibok
\end{proof}

\begin{definition}[Clenshaw evaluation from finite data]
\label{def:finite-clenshaw}
\lean{Polynomial.finiteClenshawState, Polynomial.finiteClenshaw}\leanok
Over a ring, take $n$ coefficients $c_0,\ldots,c_{n-1}$ and recurrence
coefficients $a_k,b_k$ only for $0\le k<n-2$.
Define the states for zero, one and two coefficients by
$(0,0)$, $(c_0,0)$ and $(c_0,c_1)$, respectively.
For $n\ge3$, shift all three finite arrays by one index, compute the suffix
state $(u,v)$ recursively, and return $(c_0+b_0v,u+a_0v)$.
The evaluator applied to the resulting state $(u,v)$ returns $up_0+vp_1$.
Thus no recurrence coefficient is evaluated for an expansion with at most two coefficients.

References: \cite{clenshaw}; \cite[18.40(i)]{dlmf}.
\end{definition}

\begin{theorem}[Finite and stream Clenshaw states agree]
\label{thm:finite-clenshaw}\uses{def:finite-clenshaw,def:clenshaw}
\lean{Polynomial.finiteClenshawState_eq, Polynomial.finiteClenshaw_eq}\leanok
For any ring and streams $a,b,c:\N\to R$, restricting them to the finite
indices used by an $n$-coefficient expansion gives the same state and evaluation
as the backward Clenshaw recurrence. The state identity also holds for every
suffix starting at an arbitrary index $k$.
\end{theorem}
\begin{proof}
\leanok
\end{proof}

\end{supporting}

\section{Integration of coefficients}\label{sec:textbook-calculus-integration}

Primitives of individual modes give primitives of finite expansions. Their endpoint differences give the moments used in Clenshaw--Curtis quadrature.

\begin{definition}[Admissible degrees for Chebyshev primitives]
\label{def:T-primitive-input}
\lean{Polynomial.Chebyshev.PrimitiveInput}\leanok
For a field $K$ and degree $n$, call $n$ an admissible primitive degree when the
denominators in the adjacent-degree formula are nonzero: there is no condition
at $n=0$, one requires $4\ne0$ at $n=1$, and for $n=m+2$ one requires
$2(m+3)\ne0$ and $2(m+1)\ne0$ in $K$.
\end{definition}

\begin{definition}[Chebyshev primitives]
\label{def:T-primitives}\uses{def:chebyshev-mathlib,def:T-primitive-input}
\lean{Polynomial.Chebyshev.primitiveT, Polynomial.Chebyshev.primitiveTAt}\leanok
Over a field $K$, for an admissible primitive degree $n$ as in Definition~\ref{def:T-primitive-input}, define $Q_0=T_1$, $Q_1=T_2/4$, and
\[
 Q_n=\frac{T_{n+1}}{2(n+1)}-\frac{T_{n-1}}{2(n-1)}\quad(n\ge2).
\]
The anchored polynomial is $Q_n-Q_n(a)$.

Reference: \cite[Chapter 19]{atap}.
\end{definition}

\begin{theorem}[Exact coefficient calculus]
\label{thm:T-coefficient-calculus}\uses{def:T-primitives,thm:chebyshev-recurrence,def:polynomial-linear-maps,thm:T-primitive-input}
\lean{Polynomial.Chebyshev.derivative_primitiveT, Polynomial.Chebyshev.derivative_sum_primitiveT, Polynomial.Chebyshev.derivative_primitiveTAt, Polynomial.Chebyshev.eval_primitiveTAt_self, Polynomial.Chebyshev.derivative_T_succ_eval_neg_one}\leanok
Over a characteristic-zero field, $Q_n'=T_n$, so differentiating $\sum_{n\in s}a_nQ_n$ recovers
$\sum_{n\in s}a_nT_n$. Anchored primitives have the same derivative and vanish
at their anchor. The negative endpoint derivative is
$T_{n+1}'(-1)=(-1)^n(n+1)^2$.

Reference: \cite[Chapters 19 and 21]{atap}.
\end{theorem}
\begin{proof}
\leanok
\end{proof}

\begin{definition}[Unweighted Chebyshev moments]
\label{def:integral-moment}
\lean{Polynomial.Chebyshev.integralMoment}\leanok
For a natural index $n$, define $m_1=0$ and, for $n\ne1$,
\[
 m_n=\frac{1+(-1)^n}{1-n^2}.
\]
The exceptional odd mode is defined separately, without division by zero.
\end{definition}

\begin{theorem}[Unweighted integration of Chebyshev coefficients]
\label{thm:T-unweighted-moments}\uses{thm:integral-moment}
\lean{Polynomial.Chebyshev.integral_sum_eval_T}\leanok
With respect to Lebesgue measure,
\[
 \int_{-1}^1T_n(x)\,dx=\begin{cases}2/(1-n^2),&n\text{ even},\\0,&n\text{ odd}.\end{cases}
\]
Integration of a finite Chebyshev expansion is the corresponding weighted sum
of its coefficients.

Reference: \cite[equation (19.6) and Theorem 19.2]{atap}.
\end{theorem}
\begin{proof}
\leanok
\end{proof}

\begin{supporting}[Primitive domains and moment identities]

\begin{theorem}[Admissible primitive degrees in characteristic zero]
\label{thm:T-primitive-input}\uses{def:T-primitive-input}
\lean{Polynomial.Chebyshev.PrimitiveInput.ofCharZero}\leanok
Over a field of characteristic zero, every $n\in\N$ is an admissible
primitive degree: $4\ne0$, and $2(m+3)\ne0$, $2(m+1)\ne0$ for every $m\in\N$.
\end{theorem}
\begin{proof}
\leanok
\end{proof}

\begin{theorem}[Chebyshev moments]
\label{thm:integral-moment}\uses{def:integral-moment,thm:T-coefficient-calculus,thm:chebyshev-ode}
\lean{Polynomial.Chebyshev.integralMoment_denominator_ne_zero,Polynomial.Chebyshev.integralMoment_eq_ite,Polynomial.Chebyshev.integral_eval_T,Polynomial.Chebyshev.integralMoment_eq_integral,Polynomial.Chebyshev.integral_eval_T_eq_primitive,Polynomial.Chebyshev.integral_eval_T_eq_ite}\leanok
For $n\ne1$, the denominator $1-n^2$ is nonzero. For every $n\ge0$,
\[
 m_n=\int_{-1}^1T_n(x)\,dx
 =Q_n(1)-Q_n(-1)
 =\begin{cases}2/(1-n^2),&n\text{ even},\\0,&n\text{ odd}.\end{cases}
\]
Reference: \cite[equation (19.6)]{atap}.
\end{theorem}
\begin{proof}
\leanok
\end{proof}

\end{supporting}

\section{Differentiating and integrating infinite series}\label{sec:textbook-calculus-differentiation}

Write $S_a$ for the convergent Chebyshev series of Definition~\ref{def:complex-chebyshev-series}. Differentiation uses the second-kind basis, also denoted $C_n^{(1)}=U_n$; the precise bridge is Theorem~\ref{thm:chebyshev-gegenbauer-entry}. Geometric coefficient decay justifies differentiation on smaller ellipses.

The derivative is defined on the domain where it exists, as in Definition~\ref{def:open-domain-derivative-value}; geometric decay supplies that existence (Theorem~\ref{thm:chebyshev-series-derivative-input}).

\begin{theorem}[Chebyshev series differentiation into ultraspherical coefficients]
\label{thm:chebyshev-series-differentiation}\uses{thm:chebyshev-series-higher-derivatives,thm:normal-series-differentiation}
\lean{Polynomial.Chebyshev.hasSum_derivative_complexChebyshevSeries, Polynomial.Chebyshev.hasSum_gegenbauer_derivative_complexChebyshevSeries, Polynomial.Chebyshev.tendstoUniformlyOn_derivative_complexChebyshevSeries}\leanok
If $|a_n|\le C\rho^{-n}$, $\rho>1$, then on $E_\rho$,
\[
 S_a'(z)=\sum_{n\ge0}a_nT_n'(z)
        =\sum_{n\ge0}(n+1)a_{n+1}C_n^{(1)}(z).
\]
For $1<r<\rho$, derivatives of partial sums converge uniformly on every
compact subset of $E_r$.

References: \cite[Chapter 21]{atap}; \cite[18.9.21]{dlmf}.
\end{theorem}
\begin{proof}
\leanok
\end{proof}

\begin{theorem}[Uniform geometric derivative error]
\label{thm:chebyshev-series-derivative-error}\uses{thm:chebyshev-series-differentiation}
\lean{Polynomial.Chebyshev.norm_derivative_complexChebyshevSeries_sub_sum_le, Polynomial.Chebyshev.exists_uniform_derivative_series_bound}\leanok
Assume $\rho>1$ and $|a_n|\le C\rho^{-n}$ for every $n$. Let $1<r<\rho$ and $\delta>0$. If $z\in E_r$ and the
radius-$\delta$ circle about $z$ lies in $E_r$, then
\[
 \left|S_a'(z)-\sum_{n<N}a_nT_n'(z)\right|
 \le \frac{C(r/\rho)^N}{\delta(1-r/\rho)}.
\]
Every compact $K\subset E_r$ admits one $\delta>0$ for which this bound holds
for all $z\in K$ and all $N\ge0$, in particular for the real interval.

Reference: \cite[Theorem 21.1]{atap}.
\end{theorem}
\begin{proof}
\leanok
\end{proof}

\begin{theorem}[Anchored integration of absolute Chebyshev expansions]
\label{thm:chebyshev-anchored-series-integral}\uses{thm:absolute-series-coefficients,thm:T-coefficient-calculus}
\lean{Polynomial.Chebyshev.eval_primitiveTAt_eq_integral, Polynomial.Chebyshev.norm_eval_primitiveTAt_le_two, Polynomial.Chebyshev.hasSum_primitiveTAt}\leanok
Write $A_n$ for the existing primitive of $T_n$ anchored by $A_n(-1)=0$.
For $x\in[-1,1]$, $A_n(x)=\int_{-1}^xT_n(t)\,dt$ and $|A_n(x)|\le2$.
If $\sum_n|a_n|<\infty$ and $f$ agrees with $S_a$ on $[-1,1]$, then
$f$ is integrable on every $[-1,x]$, $x\in[-1,1]$, and
\[
 \int_{-1}^x f(t)\,dt=\sum_{n\ge0}a_n A_n(x).
\]


Reference: \cite[Chapters 19 and 21]{atap}.
\end{theorem}
\begin{proof}
\leanok
\end{proof}

\begin{theorem}[Uniform anchored-integration error]
\label{thm:chebyshev-anchored-series-error}\uses{thm:chebyshev-anchored-series-integral,thm:series-operation-errors}
\lean{Polynomial.Chebyshev.norm_integral_sub_sum_primitiveTAt_le, Polynomial.Chebyshev.tendstoUniformlyOn_primitiveTAt, Polynomial.Chebyshev.norm_integral_sub_sum_primitiveTAt_le_geometric}\leanok
Let $f$ agree with $S_a$ on $[-1,1]$, and let $A_n$ be the primitive of $T_n$ anchored at $-1$. If $\sum_n|a_n|<\infty$, $\sum_{n<N}a_n A_n(x)$ converges to $\int_{-1}^x f(t)\,dt$ uniformly on
$[-1,1]$, with error at most
$2\sum_{n\ge N}|a_n|$ after the first $N$ modes. If $|a_n|\le C\rho^{-n}$,
$\rho>1$, this is at most $2C\rho^{-N}/(1-\rho^{-1})$.

Reference: \cite[Chapter 21]{atap}.
\end{theorem}
\begin{proof}
\leanok
\end{proof}

\begin{supporting}[Derivative domains and normal convergence]

\begin{definition}[Derivative inputs of geometric Chebyshev expansions]
\label{def:chebyshev-series-derivative-input}\uses{thm:analytic-chebyshev-series,thm:open-domain-analytic-derivatives}
\lean{Polynomial.Chebyshev.complexChebyshevSeriesDerivInput}\leanok
For $\rho>1$, $|a_n|\le C\rho^{-n}$ and $k\ge0$, regard $S_a|_{E_\rho}$
as an element of $\mathcal D^k(E_\rho)$. Its order-$k$ derivative value is
$D_{E_\rho}^kS_a$.
\end{definition}

\begin{theorem}[Geometric-series derivative input characterisation]
\label{thm:chebyshev-series-derivative-input}\uses{def:chebyshev-series-derivative-input,thm:open-domain-derivative-value}
\lean{Polynomial.Chebyshev.complexChebyshevSeriesDerivInput_spec}\leanok
For $\rho>1$, $|a_n|\le C\rho^{-n}$ and $k\ge0$, the function
$S_a|_{E_\rho}$ admits an order-$k$ derivative chain on $E_\rho$.
The chain starts at $S_a|_{E_\rho}$ and ends at $D_{E_\rho}^kS_a$.
\end{theorem}
\begin{proof}
\leanok
\end{proof}

\begin{theorem}[Normal convergence supplies classical Chebyshev values]
\label{thm:normal-chebyshev-domain}\uses{thm:complex-chebyshev-series-value,thm:open-domain-analytic}
\lean{Polynomial.Chebyshev.ComplexChebyshevSeriesPoint.ofNormal, Polynomial.Chebyshev.ComplexChebyshevSeriesPoint.ofNormal_val, Polynomial.Chebyshev.analyticOnOpen_complexChebyshevSeries_of_normal}\leanok
Let $U\subseteq\mathbb C$ be open. If $\sum_n u_n<\infty$ and
$|a_nT_n(z)|\le u_n$ for every $n$ and $z\in U$, then $U\subseteq D_a$
and $S_a|_U$ is analytic.
\end{theorem}
\begin{proof}
\leanok
\end{proof}

\end{supporting}

\section{Multiplication of series}\label{sec:textbook-calculus-products}

The product formula $T_mT_n=\tfrac12(T_{m+n}+T_{|m-n|})$ groups pairs of modes by their output degree. Both terms contribute when those degrees coincide.

\begin{definition}[Infinite Chebyshev product coefficients]
\label{def:chebyshev-product-coefficients}\uses{def:regrouped-coefficients,thm:chebyshev-product-input}
\lean{Polynomial.Chebyshev.chebyshevProductCoefficients}\leanok
For coefficient sequences $a,b$ such that $(a_mb_n)_{|m-n|=k}$ is summable, define
\[
 c_k=\frac12\left(\sum_{m+n=k}a_mb_n+
                         \sum_{|m-n|=k}a_mb_n\right),\qquad m,n\ge0.
\]
Both sums contribute when their output degrees coincide. This retains the
full-constant convention, including multiplication by $T_0$.

References: \cite[18.18.21]{dlmf}; \cite[Section 2.2, equation (2.7)]{olver-townsend}.
\end{definition}

\begin{theorem}[Multiplication of absolute Chebyshev expansions]
\label{thm:chebyshev-series-product}\uses{def:chebyshev-product-coefficients,thm:series-regrouping,thm:absolute-series-coefficients}
\lean{Polynomial.Chebyshev.summable_norm_chebyshevProductCoefficients, Polynomial.Chebyshev.tsum_norm_chebyshevProductCoefficients_le, Polynomial.Chebyshev.hasSum_chebyshevProductCoefficients, Polynomial.Chebyshev.complexChebyshevSeries_product, Polynomial.Chebyshev.complexCoefficient_product_series}\leanok
If $\sum|a_n|$ and $\sum|b_n|$ are finite, then $c$ is absolutely summable,
\[
 \sum_k|c_k|\le\left(\sum_m|a_m|\right)\left(\sum_n|b_n|\right),
 \qquad S_a(x)S_b(x)=\sum_{k\ge0}c_kT_k(x)\quad(x\in[-1,1]).
\]
Moreover any $f\in\mathcal I_k^{\C}$ agreeing with
$S_aS_b$ on $[-1,1]$ has coefficient $c_k$.

References: \cite[18.18.21]{dlmf}; \cite[Section 2.2]{olver-townsend}.
\end{theorem}
\begin{proof}
\leanok
\end{proof}

\begin{theorem}[Product truncation error]
\label{thm:chebyshev-series-product-error}\uses{thm:chebyshev-series-product,thm:series-operation-errors}
\lean{Polynomial.Chebyshev.norm_product_series_sub_sum_le,Polynomial.Chebyshev.tendsto_chebyshevProductCoefficients_tail}\leanok
If $\sum_n|a_n|<\infty$ and $\sum_n|b_n|<\infty$, with $c$ defined by the product-coefficient formula of Definition~\ref{def:chebyshev-product-coefficients}, uniformly for
$x\in[-1,1]$,
\[
 \left|S_a(x)S_b(x)-\sum_{k<N}c_kT_k(x)\right|\le\sum_{k\ge N}|c_k|.
\]
The right side tends to zero.

Reference: \cite[Chapter 3]{atap}.
\end{theorem}
\begin{proof}
\leanok
\end{proof}

\begin{supporting}[Regrouping and product-coefficient domains]

\begin{definition}[Domain of a collected coefficient]
\label{def:regrouped-coefficient-input}
\lean{Polynomial.RegroupCoefficientInput}\leanok
For an index map $e:I\to\N$ and output index $k$, consider coefficient families $c:I\to\C$ whose restriction to $\{i:e(i)=k\}$ is summable.
Convergence is required only in this fibre.
\end{definition}

\begin{theorem}[Summability on coefficient fibres]
\label{thm:regrouped-coefficient-input}\uses{def:regrouped-coefficient-input}
\lean{Polynomial.RegroupCoefficientInput.ofSummable, Polynomial.RegroupCoefficientInput.ofSummableNorm, Polynomial.summable_mul_mode, Polynomial.RegroupCoefficientInput.ofWeightedBound}\leanok
For an index map $e:I\to\N$, a summable family $(c_i)_{i\in I}$
restricts to a summable family on every fibre $\{i:e(i)=k\}$.
In particular, absolute summability of $(c_i)$ suffices.
More generally, if $|P_k|\le w_k$ and $\sum_i|c_i|w_{e(i)}<\infty$, then
$(c_iP_{e(i)})_{i\in I}$ and each of its fibre restrictions are summable.
This includes fibres with $P_k=w_k=0$.
\end{theorem}
\begin{proof}
\leanok
\end{proof}

\rvnote{thm:regrouped-coefficient-input}{This includes fibres with $P_k=w_k=0$; no convergence of the unweighted coefficient fibre is asserted in this more general case.}

\begin{definition}[Collecting coefficients by output mode]
\label{def:regrouped-coefficients}\uses{def:regrouped-coefficient-input}
\lean{Polynomial.regroupCoefficients}\leanok
For a coefficient family $c:I\to\mathbb C$ summable on $\{i:e(i)=k\}$, define
$c^{[e]}_k=\sum_{i:e(i)=k}c_i$.

Reference: elementary regrouping, used for the aliasing formula \cite[Theorem 4.2]{atap}.
\end{definition}

\begin{theorem}[Characterization of the collected coefficient]
\label{thm:regrouped-coefficient-value}\uses{def:regrouped-coefficients}
\lean{Polynomial.hasSum_regroupCoefficients}\leanok
If $(c_i)_{e(i)=k}$ is summable, its sum is $c^{[e]}_k$.
\end{theorem}
\begin{proof}
\leanok
\end{proof}

\begin{theorem}[Absolute and weighted regrouping]
\label{thm:series-regrouping}\uses{def:regrouped-coefficients,thm:regrouped-coefficient-input}
\lean{Polynomial.summable_norm_regroupCoefficients, Polynomial.tsum_norm_regroupCoefficients_le, Polynomial.hasSum_regrouped_series_of_weighted_bound, Polynomial.hasSum_regrouped_coefficients_of_weighted_bound, Polynomial.hasSum_regrouped_series}\leanok
If $\sum_i|c_i|<\infty$, then $\sum_k|c^{[e]}_k|\le\sum_i|c_i|$.
For mode values $P_k$ with $|P_k|\le w_k$ and
$\sum_i|c_i|w_{e(i)}<\infty$, grouping the contributions preserves the represented value:
\[
 \sum_k\left(\sum_{i:e(i)=k}c_iP_{e(i)}\right)=\sum_i c_iP_{e(i)}.
\]
Every inner contribution family is summable, even when $w_k=P_k=0$.
If each unweighted coefficient fibre is also summable, this becomes
\[
 \sum_k c^{[e]}_kP_k=\sum_i c_iP_{e(i)}.
\]
In particular, absolute coefficient summability suffices for a uniformly bounded
family. Degree-dependent weights allow growing Jacobi mode norms.

Reference: elementary regrouping, used for the aliasing formula \cite[Theorem 4.2]{atap}.
\end{theorem}
\begin{proof}
\leanok
\end{proof}

\begin{definition}[Domain of a Chebyshev product coefficient]
\label{def:chebyshev-product-input}
\lean{Polynomial.Chebyshev.ChebyshevProductInput}\leanok
At output index $k$, consider pairs of coefficient sequences $a,b$ for which the family $(a_mb_n)_{|m-n|=k}$ is summable.
This condition concerns only the chosen output index.
\end{definition}

\begin{theorem}[Summability of Chebyshev product coefficients]
\label{thm:chebyshev-product-input}\uses{def:chebyshev-product-input,def:regrouped-coefficient-input}
\lean{Polynomial.Chebyshev.finite_add_fibre, Polynomial.Chebyshev.ChebyshevProductInput.addInput, Polynomial.Chebyshev.ChebyshevProductInput.ofSummableNorm}\leanok
For $k\in\N$, the set $\{(m,n)\in\N^2:m+n=k\}$ is finite, so
$\sum_{m+n=k}a_mb_n$ exists for any sequences $a,b$.
If both sequences are absolutely summable, then
$(a_mb_n)_{|m-n|=k}$ is summable for every $k$, and $(a,b)$ belongs to
the product-coefficient domain of Definition~\ref{def:chebyshev-product-input}.
\end{theorem}
\begin{proof}
\leanok
\end{proof}

\begin{theorem}[Characterization of a Chebyshev product coefficient]
\label{thm:chebyshev-product-coefficient-formula}\uses{def:chebyshev-product-coefficients}
\lean{Polynomial.Chebyshev.chebyshevProductCoefficients_eq}\leanok
For $k\in\N$ and sequences $a,b$ such that $(a_mb_n)_{|m-n|=k}$ is
summable, the product coefficient is
\[
 c_k=\frac12\left(\sum_{m+n=k}a_mb_n+\sum_{|m-n|=k}a_mb_n\right).
\]
The first sum is finite. The formula includes $k=0$.
\end{theorem}
\begin{proof}
\leanok
\end{proof}

\end{supporting}

\section{Other intervals}\label{sec:textbook-calculus-intervals}

An affine coordinate maps $[-1,1]$ to $[a,b]$ when $a<b$. It preserves approximation errors and supplies the scale factors for differentiation and integration.

\begin{definition}[Affine interval coordinates]
\label{def:affine-chebyshev-interval}
\lean{Polynomial.Chebyshev.intervalMap, Polynomial.Chebyshev.intervalCoord}\leanok
For arbitrary $a,b$ define $A$; for $a\ne b$ define its inverse coordinate $B$:
\[
 A(t)=\frac{b-a}{2}t+\frac{a+b}{2},\qquad
 B(x)=\frac{2x-a-b}{b-a}.
\]


References: Chebfun interval scaling \cite{chebfun}; \cite[3.11.12]{dlmf}.
\end{definition}

\begin{theorem}[Affine transport of approximation and calculus]
\label{thm:affine-chebyshev-calculus}\uses{def:affine-chebyshev-interval,def:interval-integrable-input}
\lean{Polynomial.Chebyshev.intervalMap_neg_one, Polynomial.Chebyshev.intervalMap_one, Polynomial.Chebyshev.intervalMap_intervalCoord, Polynomial.Chebyshev.intervalCoord_intervalMap, Polynomial.Chebyshev.intervalMap_mem_Icc, Polynomial.Chebyshev.intervalCoord_mem_Icc, Polynomial.Chebyshev.hasDerivAt_intervalMap, Polynomial.Chebyshev.hasDerivAt_intervalCoord, Polynomial.Chebyshev.hasDerivAt_comp_intervalCoord, Polynomial.Chebyshev.intervalIntegrable_comp_intervalMap, Polynomial.Chebyshev.integral_comp_intervalMap, Polynomial.Chebyshev.abs_sub_comp_intervalCoord_le, Polynomial.Chebyshev.eval_comp_intervalCoord, Polynomial.Chebyshev.degree_comp_intervalCoord}\leanok
For $a<b$, $A$ and $B$ are inverse maps between $[-1,1]$ and $[a,b]$,
with $A(-1)=a$ and $A(1)=b$. Composition $p\circ B$ is ordinary polynomial
composition and preserves degree. Uniform approximation errors transfer unchanged.
At points where $g$ has derivative $d$ at $B(x)$, $g\circ B$ has derivative
$2d/(b-a)$ at $x$. Moreover,
\[
 \int_a^b f(x)\,dx=\frac{b-a}{2}\int_{-1}^1 f(A(t))\,dt.
\]
The integral identity assumes that $f$ is integrable between $a$ and $b$ and also asserts integrability of $f\circ A$ on $[-1,1]$. It holds for either orientation and
for equal endpoints. The approximation and interval-membership assertions
assume $a<b$; the coordinate $B$ requires distinct endpoints.

Reference: affine change of variables, as in \cite[3.11.12]{dlmf}.
\end{theorem}
\begin{proof}
\leanok
\end{proof}

\begin{supporting}[Oriented integrals and termwise operations]

\begin{definition}[Integrable functions between two endpoints]
\label{def:interval-integrable-input}
\lean{MeasureTheory.IntervalIntegrableInput}\leanok
For real endpoints $a,b$, consider real functions $f$
with finite Lebesgue integral of $|f|$ between the endpoints. Either orientation
is allowed. At equal endpoints every function is admitted.
\end{definition}

\begin{theorem}[Interval-input characterisation]
\label{thm:interval-integrable-input}\uses{def:interval-integrable-input}
\lean{MeasureTheory.IntervalIntegrableInput.integrable}\leanok
If $f$ belongs to the interval-integrability domain for endpoints $a,b$,
then $f$ is integrable on the interval between $a$ and $b$; its oriented
integral $\int_a^b f(x)\,dx$ therefore exists.
\end{theorem}
\begin{proof}
\leanok
\end{proof}

\begin{theorem}[Series tail and product error bounds]
\label{thm:series-operation-errors}
\lean{Polynomial.norm_series_sub_sum_le_tail, Polynomial.norm_mul_sub_mul_le_errors}\leanok
If $\sum f_n=s$, $|f_n|\le v_n$ and $\sum v_n<\infty$, then
$|s-\sum_{n<N}f_n|\le\sum_{n\ge N}v_n$.
For complex numbers with $|f|\le A$, $|q|\le B$, $|f-p|\le\varepsilon$ and
$|g-q|\le\eta$, one has $|fg-pq|\le A\eta+B\varepsilon$.

Reference: elementary estimates, used for the products and calculus of \cite[Chapters 3 and 21]{atap}.
\end{theorem}
\begin{proof}
\leanok
\end{proof}

\begin{theorem}[Differentiation of normally convergent polynomial expansions]
\label{thm:normal-series-differentiation}\uses{thm:series-operation-errors,def:polynomial-linear-maps}
\lean{Polynomial.hasSum_derivative_series, Polynomial.tendstoUniformlyOn_derivative_series, Polynomial.norm_derivative_series_sub_sum_le}\leanok
Let $P_n\in\mathbb C[X]$, $a_n\in\mathbb C$ and $U\subseteq\mathbb C$ be open.
Suppose $|a_nP_n(z)|\le u_n$ on $U$, with $\sum u_n<\infty$.
Then $S(z)=\sum a_nP_n(z)$ satisfies $S'(z)=\sum a_nP_n'(z)$ on $U$,
and derivatives of partial sums converge uniformly on each compact $K\subset U$.
If $z\in U$ and the circle of radius $\delta>0$ about $z$ lies in $U$, then
\[
 \left|S'(z)-\sum_{n<N}a_nP_n'(z)\right|
 \le \delta^{-1}\sum_{n\ge N}u_n.
\]
No orthogonality hypothesis is needed.

Reference: Weierstrass's theorem on locally uniform limits of analytic functions.
\end{theorem}
\begin{proof}
\leanok
\end{proof}

\begin{theorem}[Termwise integration under a uniform majorant]
\label{thm:series-termwise-integration}
\lean{Polynomial.hasSum_integral_series}\leanok
For continuous $F_n:\mathbb R\to\mathbb C$, suppose $|F_n(x)|\le u_n$ between
$a$ and $b$, $\sum u_n<\infty$, and $\sum F_n(x)=f(x)$ there. Then
$\sum_n\int_a^b F_n(x)\,dx=\int_a^b f(x)\,dx$.
Either orientation of the interval is allowed.

Reference: dominated convergence for series.
\end{theorem}
\begin{proof}
\leanok
\end{proof}

\begin{theorem}[Recovering coefficients from absolute expansions]
\label{thm:absolute-series-coefficients}\uses{thm:series-termwise-integration,thm:complex-coefficient-support,thm:complex-chebyshev-series-value}
\lean{Polynomial.Chebyshev.norm_eval_T_complex_le_one, Polynomial.Chebyshev.summable_chebyshevSeries_of_absolute, Polynomial.Chebyshev.ComplexChebyshevSeriesPoint.ofAbsolute, Polynomial.Chebyshev.ComplexChebyshevSeriesPoint.ofAbsolute_val, Polynomial.Chebyshev.norm_complexChebyshevSeries_le_tsum, Polynomial.Chebyshev.continuousOn_complexChebyshevSeries, Polynomial.Chebyshev.complexCoefficient_complexChebyshevSeries}\leanok
If $\sum_n|a_n|<\infty$, then $S_a$ is continuous on $[-1,1]$,
$|S_a(x)|\le\sum_n|a_n|$, and its integral Chebyshev coefficient at index $k$
is exactly $a_k$, including the full constant coefficient. More precisely,
any $f\in\mathcal I_k^{\C}$ agreeing with $S_a$ on $[-1,1]$ has coefficient $a_k$.
All interval points belong to $D_a$.

Reference: \cite[Chapters 3 and 4]{atap}.
\end{theorem}
\begin{proof}
\leanok
\end{proof}

\end{supporting}

\chapter{Clenshaw--Curtis quadrature}\label{chap:cc}

\rvsummary{Clenshaw--Curtis quadrature: polynomial exactness, positive weights and geometric error bounds.}

Integrating the Chebyshev interpolant gives Clenshaw--Curtis quadrature. Its weights and error estimates follow from interpolation, moments and positivity.

\section{Integrating the interpolant}\label{sec:textbook-cc-rule}

Clenshaw--Curtis quadrature integrates the polynomial reconstructed from the cosine-grid samples. This gives both a coefficient formula and a nodal-weight formula.

\begin{definition}[Clenshaw--Curtis quadrature]
\label{def:clenshaw-curtis}\uses{def:extrema-dct,def:integral-moment}
\lean{Polynomial.Chebyshev.clenshawCurtis}\leanok
For $N\ge1$, define $\mathcal Q_Nf=\sum_{k=0}^Nc_k(f)m_k$, where
$m_k=\int_{-1}^1T_k$. Separately, set $\mathcal Q_0f=2f(1)$ for the singleton interpolant.

Reference: \cite[Chapter 19]{atap}.
\end{definition}

\begin{theorem}[Clenshaw--Curtis exactness and baseline error]
\label{thm:clenshaw-curtis-baseline}\uses{def:clenshaw-curtis,thm:dct-reconstruction,thm:geometric-error,thm:T-unweighted-moments,thm:lagrange-complex}
\lean{Polynomial.Chebyshev.clenshawCurtis_eq_integral_dctInterpolant, Polynomial.Chebyshev.clenshawCurtis_eval, Polynomial.Chebyshev.abs_integral_sub_clenshawCurtis_le, Polynomial.Chebyshev.abs_integral_sub_clenshawCurtis_le_of_analytic}\leanok
For $N\ge1$, $\mathcal Q_Nf=\int_{-1}^1P_Nf$, and this quadrature is exact
for polynomials of degree at most $N$. If $f$ is interval integrable and
$|f-P_Nf|\le E$ on $[-1,1]$, its quadrature error is at most $2E$.
If $\rho>1$, $F$ is analytic on the open Bernstein ellipse $E_\rho$, and
$|F|\le M$ there, then
\[
 \left|\int_{-1}^1\Re F(x)\,dx-\mathcal Q_N(\Re F)\right|
 \le\frac{8M}{\rho^N(\rho-1)}.
\]


Reference: \cite[Theorem 19.1 and proof of Theorem 19.3]{atap}.
\end{theorem}
\begin{proof}
\leanok
\end{proof}

\rvnote{thm:clenshaw-curtis-baseline}{Positivity and a direct competitor estimate are given in Theorem~\ref{thm:clenshaw-curtis-positive-error}.

This is the interpolation-derived bound, not the stronger estimate of \cite[Theorem 19.3]{atap}.}

\begin{supporting}[The order-zero rule]

\begin{theorem}[Singleton quadrature]
\label{thm:clenshaw-curtis-zero}\uses{def:clenshaw-curtis,thm:extrema-dct-definition}
\lean{Polynomial.Chebyshev.clenshawCurtis_zero}\leanok
The order-zero rule integrates its constant interpolant: $\mathcal Q_0f=2f(1)$.
\end{theorem}
\begin{proof}
\leanok
\end{proof}

\end{supporting}

\section{Nodal weights and positivity}\label{sec:textbook-cc-weights}

The following identity connects the explicit weights to the coefficient-defined rule. Positivity controls the effect of errors in the samples.

\begin{definition}[Explicit Clenshaw--Curtis nodal weights]
\label{def:clenshaw-curtis-weights}\uses{def:chebyshev-mathlib,def:classical-chebyshev-nodes,def:integral-moment}
\lean{Polynomial.Chebyshev.extremaEndpointWeight, Polynomial.Chebyshev.clenshawCurtisWeight}\leanok
For $N\ge1$ and $0\le j\le N$, write $x_j=\cos(j\pi/N)$, $d_j=1/2$ at $j=0,N$ and $d_j=1$
otherwise, and let $m_k$ be the classical moment from Definition~\ref{def:integral-moment}. Define
\[
 w_j=\frac{d_j}{N}\sum_{k=0}^N b_k m_k T_k(x_j),\qquad
 b_k=\begin{cases}1&k=0\text{ or }k=N,\\2&\text{otherwise}.\end{cases}
\]
The nodal endpoints and highest DCT mode have separate half factors.

Reference: \cite[Chapters 4 and 19]{atap}.
\end{definition}

\begin{theorem}[Nodal and coefficient quadrature agree]
\label{thm:clenshaw-curtis-weights}\uses{thm:clenshaw-curtis-baseline,def:clenshaw-curtis-weights}
\lean{Polynomial.Chebyshev.extremaSum_eq_sum_endpointWeight, Polynomial.Chebyshev.clenshawCurtis_eq_sum_weights, Polynomial.Chebyshev.clenshawCurtisWeight_cosine, Polynomial.Chebyshev.clenshawCurtisWeight_eq_integral_basis}\leanok
For $N\ge1$ and arbitrary real nodal data, the coefficient-defined rule
is $Q_Nf=\sum_{j=0}^Nw_jf(x_j)$. Each weight is the integral of its Lagrange
cardinal polynomial. In the explicit formula, replace $T_k(x_j)$ by
$\cos(kj\pi/N)$ and $m_k$ by $2/(1-k^2)$ for even $k$ and zero for odd $k$.
Thus the highest term contributes when $N$ is even and vanishes when $N$ is odd.

Reference: \cite[Chapters 4 and 19]{atap}.
\end{theorem}
\begin{proof}
\leanok
\end{proof}

\begin{theorem}[Positive Clenshaw--Curtis weights and endpoint formulas]
\label{thm:clenshaw-curtis-positive}\uses{thm:cc-moment-sums,thm:clenshaw-curtis-weights}
\lean{Polynomial.Chebyshev.clenshawCurtis_kernel_pos, Polynomial.Chebyshev.clenshawCurtisWeight_pos, Polynomial.Chebyshev.sum_clenshawCurtisWeight, Polynomial.Chebyshev.clenshawCurtisWeight_last, Polynomial.Chebyshev.clenshawCurtisWeight_zero_even, Polynomial.Chebyshev.clenshawCurtisWeight_zero_odd}\leanok
For $N\ge1$, every weight is strictly positive and $\sum_{j=0}^Nw_j=2$.
The endpoint weights agree and are
\[
 w_0=w_N=\begin{cases}
 1/(N^2-1),&N\text{ even},\\
 1/N^2,&N\text{ odd}.
 \end{cases}
\]
This includes the two-point rule $N=1$. More generally,
$\sum_{k=0}^Nb_km_kT_k(x)>0$ for $x\in[-1,1]$.

References: \cite[Chapter 19]{atap}; \cite{trefethen-gauss-cc}.
\end{theorem}
\begin{proof}
\leanok
\end{proof}

\paragraph{Example (Julia).} The explicit weight formula uses the unweighted Chebyshev moments. For a positive grid order, the following code computes the weights and displays their positivity, sum and endpoint value.

\begin{juliacode}
m(k) = iseven(k) ? 2/(1 - k^2) : 0.0          # integral of T_k
function ccweights(N)
    @assert N >= 1
    [sum((k == 0 || k == N ? 1 : 2)/N *
    (j == 0 || j == N ? 1/2 : 1.0) * cos(k*j*pi/N) * m(k) for k in 0:N)
    for j in 0:N]
end
w = ccweights(8)
(minimum(w) > 0, sum(w), w[1])    # positivity, total weight and endpoint weight
\end{juliacode}

\section{Error estimates and interval scaling}\label{sec:textbook-cc-error}

A polynomial competitor bounds the quadrature error. The sharper analytic estimate uses the even coefficient tail.

\begin{theorem}[Positive-rule Clenshaw--Curtis error]
\label{thm:clenshaw-curtis-positive-error}\uses{thm:clenshaw-curtis-positive}
\lean{Polynomial.Chebyshev.abs_clenshawCurtis_sub_le, Polynomial.Chebyshev.abs_integral_sub_clenshawCurtis_le_of_competitor}\leanok
If $|f(x_j)-g(x_j)|\le\varepsilon$ at all nodes, then
$|Q_Nf-Q_Ng|\le2\varepsilon$. If $f$ is interval-integrable and a real polynomial
$p$ of degree at most $N$ satisfies $|f(x)-p(x)|\le\varepsilon$ throughout
$[-1,1]$, then
\[
 \left|\int_{-1}^1f(x)\,dx-Q_Nf\right|\le4\varepsilon.
\]
Here $N\ge1$.

Reference: \cite[proof of Theorem 19.3]{atap}.
\end{theorem}
\begin{proof}
\leanok
\end{proof}

\begin{theorem}[Sharper analytic Clenshaw--Curtis bound]
\label{thm:cc-sharp-analytic}\uses{thm:cc-even-tail}
\lean{Polynomial.Chebyshev.tsum_even_tail_le_geometric, Polynomial.Chebyshev.abs_integral_sub_clenshawCurtis_le_geometric_even, Polynomial.Chebyshev.abs_integral_sub_clenshawCurtis_le_of_analytic_cutoff, Polynomial.Chebyshev.abs_integral_sub_clenshawCurtis_le_sharp_of_analytic, Polynomial.Chebyshev.abs_integral_sub_clenshawCurtis_le_sharp_of_analytic_even}\leanok
Let $\rho>1$, let $F$ be analytic on the open Bernstein ellipse $E_\rho$,
and suppose $|F(z)|\le M$ there. For $N\ge2$,
\[
 \left|\int_{-1}^1\Re F(x)\,dx-Q_N(\Re F)\right|
 \le \frac{64M}{15\rho^{N-1}(\rho^2-1)}.
\]
If $N$ is even, replace $\rho^{N-1}$ by $\rho^N$.
More generally, any $1\le m$ with $2m\le N$ gives the denominator
$15\rho^{2m}(\rho^2-1)$. The underlying series bound uses
$|a_k|\le Cq^k$, $0\le q<1$, to give
$\frac{32}{15}Cq^{2(m+1)}/(1-q^2)$.

Reference: \cite[Theorem 19.3, equation (19.7)]{atap}.
\end{theorem}
\begin{proof}
\leanok
\end{proof}

\begin{theorem}[Interpolation and Clenshaw--Curtis on a physical interval]
\label{thm:affine-chebyshev-quadrature}\uses{thm:affine-chebyshev-calculus,thm:clenshaw-curtis-baseline}
\lean{Polynomial.Chebyshev.eval_dctInterpolant_interval_node, Polynomial.Chebyshev.integral_dctInterpolant_interval, Polynomial.Chebyshev.clenshawCurtis_interval_eval, Polynomial.Chebyshev.abs_integral_sub_clenshawCurtis_interval_le, Polynomial.Chebyshev.abs_integral_sub_clenshawCurtis_interval_le_of_analytic}\leanok
For $a<b$ and $N\ge1$, $P_N(f\circ A)\circ B$ interpolates $f$ at the mapped
extrema $A(x_j)$. Its integral is the scaled rule
$\frac{b-a}{2}\mathcal Q_N(f\circ A)$, which is exact for every physical
polynomial of degree at most $N$.
If $f\circ A$ is interval integrable and its canonical interpolation error is
at most $E$ on $[-1,1]$, the physical quadrature error is at most $(b-a)E$.
In particular, let $\rho>1$, let $F$ be analytic on the open ellipse $E_\rho$,
and let $|F|\le M$ there. For $f(x)=\Re F(B(x))$ the error is at most
\[
 \frac{4(b-a)M}{\rho^N(\rho-1)}.
\]

Reference: \cite[Chapters 8 and 19]{atap}, with affine scaling.
\end{theorem}
\begin{proof}
\leanok
\end{proof}

\begin{supporting}[Even tails, moment sums and aliasing]

\begin{theorem}[Clenshaw--Curtis error as an even coefficient tail]
\label{thm:cc-even-tail}\uses{thm:cc-mode-aliasing,thm:cc-mode-error}
\lean{Polynomial.Chebyshev.hasSum_integralMoment, Polynomial.Chebyshev.hasSum_clenshawCurtis, Polynomial.Chebyshev.hasSum_clenshawCurtis_error_even_tail, Polynomial.Chebyshev.abs_integral_sub_clenshawCurtis_le_even_tail}\leanok
Suppose $\sum_{k\ge0}|a_k|<\infty$ and $f(x)=\sum_{k\ge0}a_kT_k(x)$
on $[-1,1]$. Integration and quadrature commute with this series.
If $N\ge1$ and $2m\le N$, then the error has the absolutely convergent expansion
\[
 \int_{-1}^1 f-Q_Nf
 =\sum_{j\ge0} a_{2(j+m+1)}\bigl(m_{2(j+m+1)}-Q_NT_{2(j+m+1)}\bigr).
\]
If also $m\ge1$, its absolute value is at most
$\frac{32}{15}\sum_{j\ge0}|a_{2(j+m+1)}|$.

Reference: \cite[proof of Theorem 19.3]{atap}.
\end{theorem}
\begin{proof}
\leanok
\end{proof}

\begin{theorem}[Telescoping Chebyshev moments]
\label{thm:cc-moment-sums}\uses{thm:integral-moment}
\lean{Polynomial.Chebyshev.integralMoment_odd, Polynomial.Chebyshev.integralMoment_even, Polynomial.Chebyshev.sum_abs_integralMoment_even, Polynomial.Chebyshev.sum_abs_integralMoment_lt_three, Polynomial.Chebyshev.integralMoment_nonpos, Polynomial.Chebyshev.sum_integralMoment, Polynomial.Chebyshev.sum_endpoint_modes}\leanok
The odd moments vanish, and every positive-degree moment is nonpositive.
For every $m\ge0$,
\[
 \sum_{k=0}^{2m}|m_k|=3-\frac1{2m+1},\qquad
 \sum_{k=0}^Nm_k=4-\sum_{k=0}^N|m_k|.
\]
In particular, every finite absolute moment sum is strictly below $3$.
For $N\ge1$, the endpoint-mode sum satisfies
$\sum_{k=0}^Nb_km_k=2\sum_{k=0}^Nm_k-2-m_N$.

Reference: \cite[equation (19.6)]{atap}.
\end{theorem}
\begin{proof}
\leanok
\end{proof}

\begin{theorem}[Parity-preserving quadrature aliasing]
\label{thm:cc-mode-aliasing}\uses{thm:dct-coefficient-aliasing,thm:cc-moment-sums,def:clenshaw-curtis}
\lean{Polynomial.Chebyshev.aliasIndex_mod_two, Polynomial.Chebyshev.clenshawCurtis_eval_T_alias, Polynomial.Chebyshev.clenshawCurtis_eval_T_odd}\leanok
For $N\ge1$, the folded index $A_N(k)$ has the same parity as $k$, and
$Q_NT_k=m_{A_N(k)}$. In particular, $Q_NT_{2k+1}=0$ for every $k\ge0$,
even beyond the polynomial exactness degree.

Reference: \cite[Chapters 4 and 19]{atap}.
\end{theorem}
\begin{proof}
\leanok
\end{proof}

\begin{theorem}[Individual Clenshaw--Curtis mode errors]
\label{thm:cc-mode-error}\uses{thm:clenshaw-curtis-positive-error}
\lean{Polynomial.Chebyshev.abs_clenshawCurtis_eval_T_le_two, Polynomial.Chebyshev.abs_integralMoment_le_two_fifteenths, Polynomial.Chebyshev.abs_integralMoment_sub_clenshawCurtis_le}\leanok
For $N\ge1$ and every $k\ge0$, $|Q_NT_k|\le2$.
For $k\ge4$, $|m_k|\le2/15$, and hence
$|m_k-Q_NT_k|\le32/15$.

Reference: \cite[proof of Theorem 19.3]{atap}.
\end{theorem}
\begin{proof}
\leanok
\end{proof}

\end{supporting}

\chapter{Rootfinding with colleague matrices}\label{chap:roots}

\rvsummary{Roots of a Chebyshev expansion as the eigenvalues of its colleague matrix.}

The roots of a polynomial can be recovered as eigenvalues of a finite multiplication matrix. Choosing the Chebyshev basis gives the colleague matrix.

\section{Multiplication modulo a polynomial}\label{sec:textbook-roots-construction}

Let $K$ be a field and let $P_j$ have degree $j$. Every polynomial of degree below $n$ has coordinates in $P_0,\ldots,P_{n-1}$. In $K[x]/(g)$, multiplication by $x$ is followed by taking the remainder modulo $g$. The matrix columns are the coordinates of these basis images.

\begin{definition}[Comrade matrix in a degree-graded basis]
\label{def:comrade-matrix}\uses{thm:sequence-finite-coordinates,thm:mathlib-quotient-charpoly}
\lean{Polynomial.Sequence.quotientBasis, Polynomial.Sequence.comradeMatrix}\leanok
For monic $g$ of degree $n$, give $K[x]/(g)$ the basis whose coordinates are
the sequence coordinates $[r]_0,\ldots,[r]_{n-1}$ of the remainder $r$ modulo
$g$; by Theorem~\ref{thm:comrade-basis} its vectors are the images of
$P_0,\ldots,P_{n-1}$. Define $C_g$ to be the matrix of multiplication by $x$ in
this basis, with basis images as columns.

References: \cite{barnett-comrade}; \cite[Chapter 18]{atap}.
\end{definition}

\begin{theorem}[Comrade characteristic polynomial and final-column correction]
\label{thm:comrade-matrix}\uses{thm:comrade-basis}
\lean{Polynomial.Sequence.comradeMatrix_apply, Polynomial.Sequence.charpoly_comradeMatrix, Polynomial.Sequence.mem_spectrum_comradeMatrix_iff, Polynomial.Sequence.comradeMatrix_rank_one_entry, Polynomial.Sequence.comradeMatrix_apply_of_lt}\leanok
The characteristic polynomial of $C_g$ is $g$, and its spectrum over $K$ is the
set of roots of $g$ in $K$. Write $[q]_i$ for sequence coordinates and $[x^n]q$
for a monomial coefficient. Its entries are
\[
 (C_g)_{ij}=[xP_j]_i-([x^n](xP_j))[g]_i,\qquad 0\le i,j<n.
\]
The correction vanishes for $j+1<n$, so only the final column changes.
This includes arbitrary monic input polynomials, with repeated roots allowed.

References: \cite{barnett-comrade}; \cite[Chapter 18]{atap}.
\end{theorem}
\begin{proof}
\leanok
\end{proof}

\begin{supporting}[Quotient bases and finite coordinates]

\begin{theorem}[Comrade basis vectors and coordinates]
\label{thm:comrade-basis}\uses{def:comrade-matrix}
\lean{Polynomial.Sequence.quotientBasis_apply, Polynomial.Sequence.quotientBasis_repr}\leanok
The $i$th basis vector of $K[x]/(g)$ is the image of $P_i$, and the
coordinates of an element are the sequence coordinates of its remainder
modulo $g$.
\end{theorem}
\begin{proof}
\leanok
\end{proof}

\end{supporting}

\section{Chebyshev colleague matrices}\label{sec:textbook-roots-colleague}

In the Chebyshev basis the multiplication recurrence determines all but a correction in the final column. The resulting eigenvalue problem finds the roots of the polynomial.

\begin{theorem}[Chebyshev colleague rootfinding]
\label{thm:chebyshev-colleague}\uses{thm:comrade-matrix,def:chebyshev-polynomial-basis,def:chebyshev-mathlib}
\lean{Polynomial.Chebyshev.colleagueMatrix, Polynomial.Chebyshev.charpoly_colleagueMatrix, Polynomial.Chebyshev.mem_spectrum_colleagueMatrix_iff, Polynomial.Chebyshev.colleagueMatrix_size, Polynomial.Chebyshev.colleagueMatrix_entry, Polynomial.Chebyshev.mem_complex_spectrum_colleagueMatrix_iff}\leanok
For every nonzero real polynomial $p$, define its colleague matrix by applying
the comrade construction in the $T$ basis to $p/\operatorname{lc}(p)$.
Its size is $\deg p$ and its characteristic polynomial is $p/\operatorname{lc}(p)$.
The entry formula is the final-column correction above, in the full-constant
Chebyshev convention. Its real spectrum is the real roots of $p$; after
embedding its entries in $\mathbb C$, its complex spectrum is all complex roots
of $p$. Repeated roots are allowed.

References: \cite{good-colleague}; \cite[Theorem 18.1]{atap}.
\end{theorem}
\begin{proof}
\leanok
\end{proof}

\paragraph{Example (Julia).} For degree $n\ge2$, the final-column correction in the full-constant convention is $c_i/(2c_n)$. The example forms that matrix and checks the polynomial residual at its computed eigenvalues. The recurrence starts with $xT_0=T_1$.

\begin{juliacode}
using LinearAlgebra
function colleague(c)               # p = sum c[k+1] T_k of degree n >= 2
    n = length(c) - 1
    @assert n >= 2 && !iszero(c[end])
    C = zeros(n, n)
    C[2, 1] = 1                     # x T_0 = T_1
    for j in 1:n-1, i in (j, j+2)   # x T_j = (T_{j-1} + T_{j+1})/2
        i <= n && (C[i, j+1] = 1/2)
    end
    C[:, n] .-= c[1:n] ./ (2c[n+1]) # final-column correction
    C
end
function T(k, z)                    # T_k(z) by the recurrence
    a, b = one(z), z
    for _ in 1:k; a, b = b, 2z*b - a; end
    a
end
c = [0.2, -1.0, 0.5, 0.3, 0.0, 1.0] # T_5 + 0.3 T_3 + ... + 0.2
r = eigvals(colleague(c))
maximum(abs(sum(c[k+1]*T(k, z) for k in 0:5)) for z in r)  # polynomial residual at the computed eigenvalues
count(z -> abs(imag(z)) < 1e-12 && abs(real(z)) <= 1, r)   # number of computed real roots in the interval
\end{juliacode}

\begin{supporting}[Characteristic-polynomial and coordinate proofs]

\begin{theorem}[Quotient power basis and characteristic polynomial]
\label{thm:mathlib-quotient-charpoly}
\lean{AdjoinRoot.powerBasis, AdjoinRoot.powerBasis', AdjoinRoot.minpoly_powerBasis_gen_of_monic, charpoly_leftMulMatrix, LinearMap.charpoly_toMatrix, Matrix.mem_spectrum_iff_isRoot_charpoly}\mathlibok
For a monic polynomial $g$ over a field, $K[x]/(g)$ has the power basis
$1,x,\ldots,x^{\deg g-1}$. Multiplication by its generating root has
characteristic polynomial $g$. Characteristic polynomials are invariant under
choice of basis, and their roots in the base field are exactly the matrix spectrum.
\end{theorem}
\begin{proof}
\mathlibok
\end{proof}

\begin{theorem}[Finite coordinates of a polynomial sequence]
\label{thm:sequence-finite-coordinates}
\lean{Polynomial.Sequence.basis_repr_eq_zero_of_degree_lt, Polynomial.Sequence.sum_fin_repr, Polynomial.Sequence.degree_sum_fin_lt}\leanok
Let $P_j$ have degree $j$ over a field. A polynomial of degree below $n$ has
zero coordinates at every index $j\ge n$ and reconstructs from its first $n$
coordinates. Conversely, a combination of those $n$ basis members has degree
below $n$. No orthogonality hypothesis is required.

Reference: elementary linear algebra used in the comrade construction \cite{barnett-comrade}.
\end{theorem}
\begin{proof}
\leanok
\end{proof}

\end{supporting}

\part{Orthogonal polynomials and the classical families}\label{part:op}
\rvsummary{General orthogonal polynomials and the Jacobi family, which supply the many weighted bases that ApproxFun works with.}

\chapter{Orthogonal polynomials for a general measure}\label{chap:general}

\rvsummary{Recurrence and zero structure for general measures, followed by weighted expansions under a separate polynomial-density hypothesis.}

The recurrence, kernel and zero structure of orthogonal polynomials follow from a general measure and its polynomial inner product. Completeness requires a separate density assumption.

\section{Measures and orthogonality}\label{sec:textbook-general-measure}

A degree-graded sequence has one polynomial $P_n$ of degree $n$ for each $n\ge0$. Orthogonality is a separate condition on the integrals of its products. Squared norms need not be positive until a nondegeneracy hypothesis is imposed.

\begin{definition}[Orthogonal polynomial sequence]
\label{def:orthogonal-sequence}
\lean{Polynomial.Sequence.IsOrthogonal}
\leanok
For $P:\operatorname{Polynomial.Sequence}\R$ and a measure $\mu$, require
integrability of each product $P_mP_n$ and vanishing integrals when $m\ne n$.
This predicate alone imposes no normalisation or positivity of squared norms.

References: \cite[18.2(i)]{dlmf}, \cite[Chapter 3]{szego}.
\end{definition}

\begin{theorem}[First-kind Chebyshev sequence is orthogonal]
\label{thm:T-sequence}\uses{def:orthogonal-sequence,thm:T-orthogonal}
\lean{Polynomial.Chebyshev.isOrthogonal_chebyshevTsequence}\leanok
The sequence $(T_n)_{n\in\N}$ is an orthogonal polynomial sequence for $\mu_T$
in the sense of Definition~\ref{def:orthogonal-sequence}.

Reference: \cite[Table 18.3.1]{dlmf}.
\end{theorem}
\begin{proof}
\leanok
\end{proof}

\begin{definition}[Finite polynomial moments]
\label{def:polynomial-moments}
\lean{Polynomial.HasFiniteMoments}\leanok
A Borel measure $\mu$ on $\R$ has finite polynomial moments when
$x\mapsto p(x)^2$ is $\mu$-integrable for every $p\in\R[X]$.
No interval-support or connectedness restriction is imposed.

References: \cite[18.2(i)]{dlmf}, \cite[Chapter 3]{szego}.
\end{definition}

\begin{definition}[Positive-definite polynomial integral]
\label{def:polynomial-positive}\uses{def:polynomial-moments}
\lean{Polynomial.HasPositiveDefiniteIntegral}\leanok
Require every polynomial square to be integrable and $\int p(x)^2\,d\mu(x)>0$ for every nonzero $p\in\R[X]$.

References: \cite[18.2(i)]{dlmf}, \cite[Chapter 3]{szego}.
\end{definition}

\begin{definition}[Integrable polynomial pairs]
\label{def:polynomial-integrable-pair}
\lean{Polynomial.IntegrablePair, Polynomial.integralInner}\leanok
For a Borel measure $\mu$ on $\R$, an integrable polynomial pair consists of
$p,q\in\R[X]$ together with the condition $pq\in L^1(\mu)$.
Its pairing is the finite integral $\langle p,q\rangle_\mu=\int pq\,d\mu$.
This pairing is defined only on integrable pairs; no assumption on other
polynomial moments is required.
\end{definition}

\begin{theorem}[Polynomial integral inner product]
\label{thm:polynomial-inner}\uses{def:polynomial-positive,thm:polynomial-integrable-pair}
\lean{Polynomial.HasFiniteMoments.integrable_mul, Polynomial.HasFiniteMoments.isFiniteMeasure, Polynomial.HasFiniteMoments.pair, Polynomial.integralInnerProductCore}\leanok
Under the moment and positivity hypotheses, $\mu$ is finite, every polynomial
product is integrable, and $\langle p,q\rangle_\mu=\int p(x)q(x)\,d\mu(x)$
defines an inner product on $\R[X]$.

References: \cite[18.2(i)]{dlmf}, \cite[Chapter 3]{szego}.
\end{theorem}
\begin{proof}
\leanok
\end{proof}

\begin{theorem}[Polynomial moments on a finite interval]
\label{thm:compact-polynomial-moments}\uses{def:polynomial-moments}
\lean{Polynomial.hasFiniteMoments_of_ae_mem_Icc}\leanok
A finite measure on $\R$ supported almost everywhere on a closed bounded
interval has finite polynomial moments.
\end{theorem}
\begin{proof}
\leanok
\end{proof}

\begin{supporting}[Orthogonality, moments and pairing identities]

\begin{theorem}[Basic properties of orthogonality]
\label{thm:orthogonal-sequence-basic}\uses{def:orthogonal-sequence}
\lean{Polynomial.Sequence.isOrthogonal_zero, Polynomial.Sequence.IsOrthogonal.integrable_sq, Polynomial.Sequence.IsOrthogonal.congr_ae}\leanok
Every polynomial sequence is orthogonal for the zero measure. If $P$ is
orthogonal for $\mu$, each $P_n^2$ is $\mu$-integrable. If $P$ is orthogonal for
$\mu$ and, for every $n$, $P_n$ and $Q_n$ agree $\mu$-almost everywhere as
functions, then $Q$ is orthogonal for $\mu$.

Reference: \cite[18.2(i)]{dlmf}.
\end{theorem}
\begin{proof}
\leanok
\end{proof}

\begin{theorem}[Characterisation of positive polynomial integrals]
\label{thm:polynomial-positive-characterisation}\uses{def:polynomial-positive}
\lean{Polynomial.HasPositiveDefiniteIntegral.hasFiniteMoments,Polynomial.HasPositiveDefiniteIntegral.integral_sq_pos}\leanok
A positive-definite polynomial integral has all polynomial moments finite.
Every nonzero polynomial has strictly positive squared integral.
\end{theorem}
\begin{proof}
\leanok
\end{proof}

\begin{theorem}[Characterisation of the integrable pairing]
\label{thm:polynomial-integrable-pair}\uses{def:polynomial-integrable-pair}
\lean{Polynomial.integralInner_eq_integral, Polynomial.integralInner_comm, Polynomial.integralInner_self}\leanok
The pairing of an integrable pair $(p,q)$ is its integral $\int pq\,d\mu$.
It is symmetric. For an integrable diagonal pair, it equals $\int p^2\,d\mu$.
\end{theorem}
\begin{proof}
\leanok
\end{proof}

\begin{theorem}[Moments supplied by an orthogonal polynomial sequence]
\label{thm:orthogonal-moments}\uses{def:polynomial-moments,thm:orthogonal-sequence-basic}
\lean{Polynomial.Sequence.IsOrthogonal.hasFiniteMoments}\leanok
If a polynomial sequence is orthogonal for $\mu$, including integrability of
all its pairwise products, then $\mu$ has finite polynomial moments.
No positivity assumption is needed.
\end{theorem}
\begin{proof}
\leanok
\end{proof}

\end{supporting}

\section{Existence and polynomial coefficients}\label{sec:textbook-general-existence}

The integral inner product gives existence and uniqueness up to scaling. Orthogonality to lower-degree polynomials then identifies expansion coefficients by integration.

\begin{theorem}[Existence of an orthogonal sequence]
\label{thm:general-existence}\uses{def:orthogonal-sequence,thm:polynomial-inner}
\lean{Polynomial.exists_isOrthogonal}\leanok
If $\mu$ has finite polynomial moments and a positive definite polynomial pairing, there exists
$P:\operatorname{Polynomial.Sequence}\R$ with $P.\operatorname{IsOrthogonal}(\mu)$.
In particular $\deg P_n=n$. Measures on unbounded sets or disjoint intervals are
included when they satisfy the hypotheses.

References: \cite[18.2(i)]{dlmf}, \cite[Chapter 3]{szego}.
\end{theorem}
\begin{proof}
\leanok
\end{proof}

\begin{theorem}[Uniqueness up to scaling]
\label{thm:general-uniqueness}\uses{thm:orthogonal-lower-degree}
\lean{Polynomial.Sequence.IsOrthogonal.eq_leadingCoeff_div_smul, Polynomial.Sequence.IsOrthogonal.exists_smul}\leanok
If $\mu$ has finite polynomial moments and a positive definite polynomial pairing, any two orthogonal sequences $P,Q$
for the same $\mu$ satisfy $Q_n=c_nP_n$, where
\[c_n=\frac{\operatorname{lc}(Q_n)}{\operatorname{lc}(P_n)}\ne0.\]


References: \cite[18.2(i)]{dlmf}, \cite[Chapter 3]{szego}.
\end{theorem}
\begin{proof}
\leanok
\end{proof}

\rvnote{thm:general-uniqueness}{This concerns a fixed measure; it does not assert moment determinacy or $L^2$ completeness.}

\begin{theorem}[Orthogonality to lower-degree polynomials]
\label{thm:orthogonal-lower-degree}\uses{def:orthogonal-sequence,thm:polynomial-inner}
\lean{Polynomial.Sequence.IsOrthogonal.integralInner_eq_zero_of_degree_lt}\leanok
If $P$ is orthogonal for a measure with finite polynomial moments and
$\deg p<n$, then $\langle p,P_n\rangle_\mu=0$.

References: \cite[18.2(i)]{dlmf}, \cite[Chapter 3]{szego}.
\end{theorem}
\begin{proof}
\leanok
\end{proof}

\begin{theorem}[Integral projection coordinates]
\label{thm:general-coordinates}\uses{def:orthogonal-sequence,thm:polynomial-inner}
\lean{Polynomial.Sequence.IsOrthogonal.integralInner_eq_basis_repr_mul, Polynomial.Sequence.IsOrthogonal.basis_repr_eq_integralInner_div}\leanok
For an orthogonal sequence $P$ under the moment and positivity hypotheses,
the coefficient of $P_n$ in a polynomial $p$ is
\[\frac{\langle p,P_n\rangle_\mu}{\langle P_n,P_n\rangle_\mu}.\]

References: \cite[18.2(i)]{dlmf}, \cite[Chapter 3]{szego}.
\end{theorem}
\begin{proof}
\leanok
\end{proof}

\section{Recurrence coefficients and kernels}\label{sec:textbook-general-recurrence}

Write $h_n$ for the squared integral norm. The recurrence coefficients express multiplication by $x$ in the polynomial basis; the finite kernel sums the corresponding normalized products.

\begin{definition}[Integrable polynomial squares]
\label{def:integrable-polynomial-square}
\lean{Polynomial.Sequence.IntegrableSquare, Polynomial.Sequence.IntegrableSquare.of}\leanok
For a measure $\mu$ on $\R$, define the square-integrability domain
\[
\mathcal Q_\mu=\{p\in\R[X]:x\mapsto p(x)^2\text{ is }\mu\text{-integrable}\}.
\]
No condition is imposed on other polynomial moments.
\end{definition}

\begin{definition}[Classical squared polynomial norms]
\label{def:general-squared-norm}\uses{def:integrable-polynomial-square,def:polynomial-integrable-pair}
\lean{Polynomial.Sequence.squaredNorm}\leanok
For a degree-indexed polynomial sequence $P$, a measure $\mu$, an index $n$,
with $P_n\in\mathcal Q_\mu$, define $h_n=\int P_n(x)^2\,d\mu(x)$.
Zero squared norms are permitted.
\end{definition}

\begin{definition}[Classical recurrence coefficient inputs]
\label{def:recurrence-coefficient-input}\uses{def:general-squared-norm}
\lean{Polynomial.Sequence.RecurrenceBInput, Polynomial.Sequence.RecurrenceCInput, Polynomial.Sequence.RecurrenceBInput.of, Polynomial.Sequence.RecurrenceCInput.of}\leanok
The domain of the diagonal coefficient at $n$ requires integrability of $P_n^2$ and
$xP_n^2$, and $h_n\ne0$. The domain of the backward coefficient at $n$ requires
integrability of $P_n^2$ and $P_{n+1}^2$, and $h_n\ne0$.
No nonzero condition is imposed on $h_{n+1}$, and no moments at other
indices are required.
\end{definition}

\begin{definition}[Classical recurrence coefficients]
\label{def:general-recurrence-coefficients}\uses{def:recurrence-coefficient-input}
\lean{Polynomial.Sequence.recurrenceA, Polynomial.Sequence.recurrenceB, Polynomial.Sequence.recurrenceC}\leanok
Write $k_n=\operatorname{lc}(P_n)$. For $B_n$ and $C_n$, respectively, use the domains of Definition~\ref{def:recurrence-coefficient-input}. Define
\[
 A_n=\frac{k_n}{k_{n+1}},\quad
 B_n=\frac{\langle XP_n,P_n\rangle_\mu}{h_n},\quad
 C_n=A_n\frac{h_{n+1}}{h_n}.
\]

References: \cite[18.2(iv)]{dlmf}; \cite[Chapter 3]{szego}.
\end{definition}

\begin{theorem}[Explicit recurrence and coefficient signs]
\label{thm:general-coefficient-signs}\uses{thm:general-three-term,def:general-recurrence-coefficients,thm:general-squared-norm}
\lean{Polynomial.Sequence.recurrenceA_ne_zero, Polynomial.Sequence.recurrenceC_ne_zero, Polynomial.Sequence.recurrenceA_mul_recurrenceC_pos, Polynomial.Sequence.recurrenceA_eq_one_of_monic, Polynomial.Sequence.recurrenceC_pos_of_monic, Polynomial.Sequence.IsOrthogonal.three_term, Polynomial.Sequence.IsOrthogonal.two_term_zero}\leanok
Assume finite polynomial moments, positive definiteness and orthogonality. Then
\[
 XP_0=A_0P_1+B_0P_0,\qquad
 XP_{n+1}=A_{n+1}P_{n+2}+B_{n+1}P_{n+1}+C_nP_n.
\]
Moreover, $h_n>0$, $A_n\ne0$, $C_n\ne0$, and $A_nC_n>0$.
If every $P_n$ is monic, then $A_n=1$ and $C_n>0$.

Reference: \cite[18.2(iv)]{dlmf}.
\end{theorem}
\begin{proof}
\leanok
\end{proof}

\begin{definition}[Classical finite kernel inputs]
\label{def:finite-kernel-input}\uses{def:general-squared-norm}
\lean{Polynomial.Sequence.KernelInput, Polynomial.Sequence.KernelInput.of}\leanok
For a polynomial sequence $P$, a measure $\mu$, and a cutoff $n$, the finite-kernel domain requires integrability of $P_j^2$ and $h_j\ne0$ for every $0\le j\le n$.
No conditions are imposed on higher indices.
\end{definition}

\begin{definition}[Finite polynomial kernel]
\label{def:general-kernel}\uses{def:finite-kernel-input}
\lean{Polynomial.Sequence.kernel}\leanok
For a cutoff $n$ in the domain of Definition~\ref{def:finite-kernel-input} and arbitrary real $x,y$, set $K_n(x,y)=\sum_{j=0}^n P_j(x)P_j(y)/h_j$.

Reference: \cite[18.2.12]{dlmf}.
\end{definition}

\begin{theorem}[Christoffel--Darboux]
\label{thm:general-cd}\uses{thm:general-coefficient-signs,def:general-kernel}
\lean{Polynomial.Sequence.IsOrthogonal.christoffelDarboux}\leanok
Let $P$ be an orthogonal polynomial sequence for a measure $\mu$ with finite polynomial moments and a positive definite polynomial pairing. For every $n\ge0$ and $x,y\in\mathbb R$,
\[
 (x-y)K_n(x,y)=\frac{A_n}{h_n}
  \bigl(P_{n+1}(x)P_n(y)-P_n(x)P_{n+1}(y)\bigr).
\]
The cleared form also holds at $x=y$.

Reference: \cite[18.2.12]{dlmf}.
\end{theorem}
\begin{proof}
\leanok
\end{proof}

\begin{theorem}[Diagonal kernel and simplicity]
\label{thm:general-cd-diag}\uses{thm:general-cd,def:polynomial-linear-maps}
\lean{Polynomial.Sequence.IsOrthogonal.christoffelDarboux_diag, Polynomial.Sequence.kernel_diag_pos, Polynomial.Sequence.IsOrthogonal.wronskian_ne_zero, Polynomial.Sequence.IsOrthogonal.not_eval_eq_zero_and_succ, Polynomial.Sequence.IsOrthogonal.derivative_eval_ne_zero}\leanok
Let $P$ be an orthogonal polynomial sequence for a measure $\mu$ with finite polynomial moments and a positive definite polynomial pairing.
The diagonal identity is
\[
 0<K_n(x,x)=\frac{A_n}{h_n}
   \bigl(P'_{n+1}(x)P_n(x)-P'_n(x)P_{n+1}(x)\bigr).
\]
Thus the Wronskian never vanishes, consecutive members have no common real
root, and $P_n(x)=0$ implies $P'_n(x)\ne0$.

References: \cite[18.2.13]{dlmf}; \cite[Chapter 3]{szego}.
\end{theorem}
\begin{proof}
\leanok
\end{proof}

\begin{supporting}[Norm, coefficient and kernel domains]

\begin{theorem}[Characterisation of square data]
\label{thm:integrable-polynomial-square}\uses{def:integrable-polynomial-square}
\lean{Polynomial.Sequence.IntegrableSquare.integrable_square}\leanok
For $p\in\mathcal Q_\mu$, the function $x\mapsto p(x)^2$ is
$\mu$-integrable.
\end{theorem}
\begin{proof}
\leanok
\end{proof}

\begin{theorem}[Squared-norm formula and positivity]
\label{thm:general-squared-norm}\uses{def:general-squared-norm,thm:polynomial-positive-characterisation,thm:polynomial-inner}
\lean{Polynomial.Sequence.squaredNorm_apply, Polynomial.Sequence.squaredNorm_pos}\leanok
The squared norm equals the integral of $P_n^2$.
Under positive-definite polynomial integration it is strictly positive.
\end{theorem}
\begin{proof}
\leanok
\end{proof}

\begin{theorem}[Nonzero recurrence denominators]
\label{thm:recurrence-coefficient-input}\uses{def:recurrence-coefficient-input}
\lean{Polynomial.Sequence.RecurrenceBInput.norm_ne, Polynomial.Sequence.RecurrenceCInput.norm_ne}\leanok
On either recurrence-coefficient domain of
Definition~\ref{def:recurrence-coefficient-input},
$h_n=\int P_n(x)^2\,d\mu(x)$ is nonzero.
\end{theorem}
\begin{proof}
\leanok
\end{proof}

\begin{theorem}[Classical recurrence coefficient formulas]
\label{thm:recurrence-coefficient-characterisation}\uses{def:general-recurrence-coefficients}
\lean{Polynomial.Sequence.recurrenceB_apply, Polynomial.Sequence.recurrenceC_apply}\leanok
On the diagonal and backward coefficient domains of Definition~\ref{def:recurrence-coefficient-input}, respectively,
$B_n=\langle XP_n,P_n\rangle_\mu/h_n$ and
$C_n=A_nh_{n+1}/h_n$, with integrable numerators and nonzero denominators.
\end{theorem}
\begin{proof}
\leanok
\end{proof}

\begin{theorem}[Characterisation of finite kernel inputs]
\label{thm:finite-kernel-input}\uses{def:finite-kernel-input,thm:integrable-polynomial-square}
\lean{Polynomial.Sequence.KernelInput.valid}\leanok
For a cutoff $n$ in the finite-kernel domain, every $0\le j\le n$
satisfies $P_j\in\mathcal Q_\mu$ and $h_j\ne0$.
If these conditions hold for all $j\in\N$, every cutoff belongs to this domain.
\end{theorem}
\begin{proof}
\leanok
\end{proof}

\begin{theorem}[Classical finite kernel formula]
\label{thm:finite-kernel-characterisation}\uses{def:general-kernel}
\lean{Polynomial.Sequence.kernel_apply}\leanok
For a cutoff $n$ in the domain of Definition~\ref{def:finite-kernel-input},
$K_n(x,y)=\sum_{j=0}^n P_j(x)P_j(y)/h_j$.
Each denominator is a finite nonzero squared integral.
\end{theorem}
\begin{proof}
\leanok
\end{proof}

\end{supporting}

\section{Zeros and parity}\label{sec:textbook-general-zeros}

The zeros are real and simple, and consecutive members interlace. These results do not require bounded or connected support.

\begin{theorem}[Real simple zeros]
\label{thm:general-real-zeros}\uses{thm:general-cd-diag}
\lean{Polynomial.Sequence.IsOrthogonal.natDegree_eq_one_of_mem_normalizedFactors, Polynomial.Sequence.IsOrthogonal.splits, Polynomial.Sequence.IsOrthogonal.card_roots, Polynomial.Sequence.IsOrthogonal.roots_nodup, Polynomial.Sequence.IsOrthogonal.separable}\leanok
Let $P$ be an orthogonal polynomial sequence for a measure $\mu$ with finite polynomial moments and a positive definite polynomial pairing.
Each $P_n$ splits over $\mathbb R$, is separable, and has exactly $n$ distinct
real roots. No boundedness or connectedness of the support is required.

References: \cite[18.2(vi)]{dlmf}; \cite[Chapter 3]{szego}.
\end{theorem}
\begin{proof}
\leanok
\end{proof}

\begin{theorem}[Strict root separation]
\label{thm:general-separation}\uses{thm:general-cd-diag}
\lean{Polynomial.exists_eval_eq_zero_between_of_wronskian_ne_zero, Polynomial.Sequence.IsOrthogonal.exists_root_succ_between, Polynomial.Sequence.IsOrthogonal.exists_root_between_succ}\leanok
Between any two distinct real zeros of $P_n$ lies a zero of $P_{n+1}$, and
between any two distinct real zeros of $P_{n+1}$ lies a zero of $P_n$.
Let $P$ be an orthogonal polynomial sequence for a measure $\mu$ with finite polynomial moments and a positive definite polynomial pairing.

References: \cite[18.2(vi)]{dlmf}; \cite[Chapter 3]{szego}.
\end{theorem}
\begin{proof}
\leanok
\end{proof}

\begin{theorem}[Outermost roots of consecutive members]
\label{thm:general-outer-roots}\uses{thm:general-real-zeros,thm:general-separation}
\lean{Polynomial.Sequence.IsOrthogonal.exists_outer_roots}\leanok
Under the finite-moment and positive-definite-integral hypotheses, for $n\ge1$
there are distinct roots $a<b$ of $P_{n+1}$ such that every root $x$ of $P_n$
satisfies $a<x<b$. Thus $P_{n+1}$ has a root below the smallest root of $P_n$
and another above its largest root.

References: \cite[Chapter 3]{szego}; \cite[18.2(vi)]{dlmf}.
\end{theorem}
\begin{proof}
\leanok
\end{proof}

\begin{theorem}[Positive normalization and even-measure parity]
\label{thm:even-positive-sequences}\uses{thm:general-existence,thm:general-uniqueness,thm:jacobi-coefficient-parity,thm:jacobi-weight-basic}
\lean{Polynomial.exists_positive_orthonormal_sequence, Polynomial.Sequence.IsOrthogonal.comp_neg_X, Polynomial.Sequence.IsOrthogonal.coeff_eq_zero_of_parity, MeasureTheory.isNegInvariant_withDensity, Polynomial.isNegInvariant_jacobiMeasure}\leanok
A real polynomial measure with finite moments and positive definite polynomial
pairing admits a positive-leading orthonormal sequence. If the measure is even,
every orthogonal sequence satisfies $p_n(-x)=(-1)^n p_n(x)$.
Even densities preserve reflection invariance.

Reference: \cite[18.2(iii) and 18.2.21]{dlmf}.
\end{theorem}
\begin{proof}
\leanok
\end{proof}

\section{Completeness and expansions}\label{sec:textbook-general-completeness}

Completeness needs polynomial density in $L^2(\mu)$ in addition to moments and positivity. Compact support supplies that density; the Hilbert-basis theorem also applies whenever it is known by another argument.

\begin{definition}[Polynomial classes in general real L2]
\label{def:general-polynomial-L2}\uses{def:polynomial-moments}
\lean{Polynomial.memLp_eval_of_hasFiniteMoments, Polynomial.polynomialToL2}\leanok
For a real measure $\mu$ with finite polynomial moments, let
$J_\mu:\R[X]\to L^2(\mu)$ send $p$ to the class of $x\mapsto p(x)$.
\end{definition}

\begin{theorem}[Density for compactly supported measures]
\label{thm:compact-polynomial-dense}\uses{thm:general-polynomial-L2-basic,thm:polynomial-inner}
\lean{Polynomial.polynomialToL2_denseRange}\leanok
If additionally $x\in[a,b]$ for $\mu$-almost every $x$, then $J_\mu$ has dense range. No positivity or connected-support hypothesis is required.

Reference: \cite[18.2(viii)]{dlmf}.
\end{theorem}
\begin{proof}
\leanok
\end{proof}

\begin{definition}[Classical normalized coefficient inputs]
\label{def:integral-coefficient-input}\uses{thm:L2-polynomial-product-integrable,thm:polynomial-inner}
\lean{Polynomial.IntegralCoefficientInput, Polynomial.IntegralCoefficientInput.ofL2, Polynomial.IntegralCoefficientInput.ofPolynomial, Polynomial.IntegralCoefficientInput.ofReal, Polynomial.IntegralCoefficientInput.integrable}\leanok
For $\mathbb K=\R$ or $\C$, a real polynomial $p$ and a measure $\mu$, define
\[
\mathcal A_{\mu,p}^{\mathbb K}
 =\{f:\R\to\mathbb K: p^2,fp\text{ are }\mu\text{-integrable and }
                    \int p^2\,d\mu>0\}.
\]
Functions retain their pointwise representatives. These conditions concern
only $p$, not other polynomial moments.
Reference: \cite[18.2(viii)]{dlmf}.
\end{definition}

\begin{definition}[Normalised polynomial classes]
\label{def:general-normalised-L2}\uses{def:integral-coefficient-input,def:general-squared-norm}
\lean{Polynomial.Sequence.normalizedToL2, Polynomial.Sequence.integralCoefficient}\leanok
For a degree-indexed sequence $P$ and a measure with finite polynomial moments, a normalization index is a natural number $n$ together with $h_n>0$; define $e_n=h_n^{-1/2}J_\mu P_n$ on this domain. For $f\in\mathcal A_{\mu,P_n}^{\R}$, define the original-normalisation coefficient by $a_n(f)=h_n^{-1}\int fP_n\,d\mu$.

Reference: \cite[18.2(viii)]{dlmf}.
\end{definition}

\begin{theorem}[Hilbert basis from polynomial density]
\label{thm:general-hilbert-basis}\uses{def:general-normalised-L2,thm:general-polynomial-L2-basic,thm:general-squared-norm,def:orthogonal-sequence}
\lean{Polynomial.Sequence.IsOrthogonal.orthonormal_normalizedToL2, Polynomial.Sequence.span_polynomialToL2, Polynomial.Sequence.normalizedToL2_dense_span, Polynomial.Sequence.IsOrthogonal.hilbertBasis, Polynomial.Sequence.IsOrthogonal.hilbertBasis_apply}\leanok
Assume finite polynomial moments, positive definiteness, and orthogonality of $P$.
If polynomial classes have dense range in $L^2(\mu)$, the normalised family is
orthonormal with dense span and defines a Hilbert basis. No support restriction
is required here.

Reference: \cite[18.2(viii)]{dlmf}.
\end{theorem}
\begin{proof}
\leanok
\end{proof}

\begin{theorem}[General polynomial L2 expansion]
\label{thm:general-L2-expansion}\uses{thm:general-hilbert-basis}
\lean{Polynomial.Sequence.IsOrthogonal.hilbertBasis_repr_smul, Polynomial.Sequence.IsOrthogonal.hasSum_integralCoefficient}\leanok
Let $P$ be orthogonal for $\mu$, assume finite polynomial moments and a positive definite polynomial pairing, and suppose polynomial classes are dense in $L^2(\mu;\R)$. Then the Hilbert-coordinate summands are $a_n(f)J_\mu P_n$, and their series sums to every $f\in L^2(\mu)$.

Reference: \cite[18.2(viii)]{dlmf}.
\end{theorem}
\begin{proof}
\leanok
\end{proof}

\begin{supporting}[Coefficient domains, normalization and multiplication proofs]

\begin{theorem}[Polynomials in general weighted L2]
\label{thm:general-polynomial-L2-basic}\uses{def:general-polynomial-L2}
\lean{Polynomial.coe_polynomialToL2, Polynomial.inner_polynomialToL2, Polynomial.inner_polynomialToL2_left}\leanok
If $\mu$ has finite polynomial moments, every real polynomial is square
integrable for $\mu$, and $J_\mu p$ agrees $\mu$-almost everywhere with
$x\mapsto p(x)$. Moreover $\langle J_\mu p,J_\mu q\rangle=\int pq\,d\mu$ and
$\langle J_\mu p,f\rangle=\int fp\,d\mu$ for every $f\in L^2(\mu)$.
The map is linear.
\end{theorem}
\begin{proof}
\leanok
\end{proof}

\begin{theorem}[Integrability of polynomial products with square-integrable functions]
\label{thm:L2-polynomial-product-integrable}\uses{def:general-polynomial-L2}
\lean{Polynomial.integrable_L2_mul_eval}\leanok
If $\mu$ has all polynomial moments finite, $p$ is a real polynomial, and
$f\in L^2(\mu;\mathbb K)$ with $\mathbb K=\R$ or $\C$, then $fp$ is integrable.
\end{theorem}
\begin{proof}
\leanok
\end{proof}

\begin{theorem}[Characterisation of normalized coefficient inputs]
\label{thm:integral-coefficient-input}\uses{def:integral-coefficient-input}
\lean{Polynomial.IntegralCoefficientInput.normalization, Polynomial.IntegralCoefficientInput.ofL2_val, Polynomial.IntegralCoefficientInput.ofPolynomial_val, Polynomial.IntegralCoefficientInput.ofReal_val}\leanok
Let $\mathbb K=\mathbb R$ or $\mathbb C$.
For $f\in\mathcal A_{\mu,p}^{\mathbb K}$, the product $fp$ is integrable
and $\int p^2\,d\mu$ is finite and positive.
If $\mu$ has finite polynomial moments, every $L^2(\mu;\mathbb K)$ function
has an integrable product with each real polynomial $p$.
If the polynomial pairing is also positive definite and $p\ne0$, every such
$L^2$ function belongs to $\mathcal A_{\mu,p}^{\mathbb K}$; real polynomials
do as well. These identifications retain the pointwise representatives.
Scalar inclusion sends $\mathcal A_{\mu,p}^{\R}$ into
$\mathcal A_{\mu,p}^{\C}$ by $f(x)\mapsto f(x)+0i$.
\end{theorem}
\begin{proof}
\leanok
\end{proof}

\begin{definition}[Finite classical coefficient families]
\label{def:integral-coefficient-family-input}\uses{def:integral-coefficient-input}
\lean{Polynomial.IntegralCoefficientFamilyInput, Polynomial.IntegralCoefficientFamilyInput.at, Polynomial.IntegralCoefficientFamilyInput.ofL2}\leanok
For a polynomial family $P$, a cutoff $N$, and $\mathbb K=\R$ or $\C$,
define the simultaneous coefficient domain
\[
\mathcal A_{\mu,P,\le N}^{\mathbb K}
 =\bigcap_{n=0}^N\mathcal A_{\mu,P_n}^{\mathbb K}.
\]
Thus a member is one pointwise function $f:\R\to\mathbb K$ for which each
$P_n^2$ and $fP_n$ is integrable and $\int P_n^2\,d\mu>0$, for $0\le n\le N$.
No condition is imposed at higher indices.
\end{definition}

\begin{theorem}[Characterisation of finite coefficient families]
\label{thm:integral-coefficient-family-input}\uses{def:integral-coefficient-family-input}
\lean{Polynomial.IntegralCoefficientFamilyInput.at_val, Polynomial.IntegralCoefficientFamilyInput.ofL2_val, Polynomial.IntegralCoefficientFamilyInput.ofL2_at}\leanok
Let $\mathbb K=\mathbb R$ or $\mathbb C$.
For $f\in\mathcal A_{\mu,P,\le N}^{\mathbb K}$ and $n\le N$, its
membership in $\mathcal A_{\mu,P_n}^{\mathbb K}$ uses the same function $f$.
Under finite moments and positive definiteness, every $L^2(\mu;\mathbb K)$
function belongs to this simultaneous domain for a degree-indexed polynomial
sequence $P$. Its representative at each index is the representative used
for the individual coefficient domain.
\end{theorem}
\begin{proof}
\leanok
\end{proof}

\begin{theorem}[Positive-norm normalization formula]
\label{thm:general-normalization-characterisation}\uses{def:general-normalised-L2}
\lean{Polynomial.Sequence.normalizedToL2_apply}\leanok
For a positive-norm index, $e_n=h_n^{-1/2}J_\mu P_n$.
The square root and inverse are taken at a strictly positive squared norm.
\end{theorem}
\begin{proof}
\leanok
\end{proof}

\begin{theorem}[Classical integral coefficient formula]
\label{thm:general-coefficient-characterisation}\uses{def:general-normalised-L2}
\lean{Polynomial.Sequence.integralCoefficient_apply}\leanok
For $f\in\mathcal A_{\mu,P_n}^{\R}$,
$a_n(f)=(\int P_n^2\,d\mu)^{-1}\int fP_n\,d\mu$.
The numerator is integrable and the denominator is finite and positive.
\end{theorem}
\begin{proof}
\leanok
\end{proof}

\begin{theorem}[Tridiagonal multiplication]
\label{thm:general-tridiagonal}\uses{thm:general-coordinates,thm:orthogonal-lower-degree}
\lean{Polynomial.integralInner_X_mul, Polynomial.Sequence.IsOrthogonal.integralInner_X_mul_eq_zero, Polynomial.Sequence.IsOrthogonal.basis_repr_X_mul_eq_zero}\leanok
For an orthogonal sequence under the moment and positivity hypotheses,
the coefficient of $P_m$ in $XP_n$ is zero if $n+1<m$ or $m+1<n$.
The integral identity $\langle XP_n,P_m\rangle_\mu=0$ needs only finite moments
and orthogonality under the same index condition.

References: \cite[18.2(iv)]{dlmf}, \cite[Chapter 3]{szego}.
\end{theorem}
\begin{proof}
\leanok
\end{proof}

\begin{theorem}[Three-term multiplication expansion]
\label{thm:general-three-term}\uses{thm:general-tridiagonal}
\lean{Polynomial.Sequence.IsOrthogonal.exists_three_term, Polynomial.Sequence.IsOrthogonal.exists_two_term_zero}\leanok
For each $n\ge0$, there are real $a,b,c$ with
\[XP_{n+1}=aP_{n+2}+bP_{n+1}+cP_n.\]
At index zero there are real $a,b$ with $XP_0=aP_1+bP_0$.
The hypotheses are finite polynomial moments, positive definiteness, and
orthogonality.

References: \cite[18.2(iv)]{dlmf}, \cite[Chapter 3]{szego}.
\end{theorem}
\begin{proof}
\leanok
\end{proof}

\end{supporting}

\chapter{Jacobi polynomials}\label{chap:jacobi}

\rvsummary{The Jacobi family: explicit formulas, derivatives, recurrences, orthogonality, exact norms, zeros and completeness.}

Jacobi polynomials organize the classical families on $[-1,1]$. Algebraic formulas allow arbitrary parameters; orthogonality and weighted expansions use $\alpha,\beta>-1$.

\section{Definition and normalization}\label{sec:textbook-jacobi-definition}

The finite-sum definition is polynomial in the parameters and works over any commutative $\Q$-algebra, including $\Q$, $\R$ and $\C$. Some parameter choices lower the degree; the leading-coefficient theorem states exactly when this happens.

\begin{definition}[Jacobi polynomials]
\label{def:jacobi}
\lean{Polynomial.jacobi}\leanok
For a commutative $\Q$-algebra $R$, $\alpha,\beta\in R$, and $n\in\N$, define
\[P_n^{(\alpha,\beta)}(X)=\sum_{\ell=0}^n
\frac{(n+\alpha+\beta+1)_\ell(\alpha+\ell+1)_{n-\ell}}
 {\ell!(n-\ell)!2^\ell}(X-1)^\ell.\]

Reference: \cite[18.5.7]{dlmf}.
\end{definition}

\begin{theorem}[Degree and leading term]
\label{thm:jacobi-degree}\uses{def:jacobi}
\lean{Polynomial.natDegree_jacobi_le, Polynomial.degree_jacobi_le, Polynomial.coeff_jacobi_self, Polynomial.degree_jacobi_eq, Polynomial.degree_jacobi_eq_iff, Polynomial.leadingCoeff_jacobi}\leanok
Let $R$ be a commutative $\Q$-algebra, let $\alpha,\beta\in R$, and let $n\in\N$.
The Jacobi polynomial $P_n^{(\alpha,\beta)}(X)$ has degree at most $n$, and its coefficient of $X^n$ is
\[
 c_n=\frac{(n+\alpha+\beta+1)_n}{2^n n!}.
\]
It has degree exactly $n$ if and only if $c_n\ne0$. In that case, its leading coefficient is $c_n$.

Reference: \cite[Table 18.3.1]{dlmf}.
\end{theorem}
\begin{proof}
\leanok
\end{proof}

\begin{theorem}[Jacobi value at one]
\label{thm:jacobi-value-one}\uses{def:jacobi}
\lean{Polynomial.eval_jacobi_one, Polynomial.succ_mul_eval_jacobi_one}\leanok
Let $R$ be a commutative $\Q$-algebra, let $\alpha,\beta\in R$, and let $n\in\N$. Then
$P_n^{(\alpha,\beta)}(1)=(\alpha+1)_n/n!$. In particular,
$(n+1)P_{n+1}^{(\alpha,\beta)}(1)=(n+\alpha+1)P_n^{(\alpha,\beta)}(1)$.

Reference: \cite[Table 18.6.1]{dlmf}.
\end{theorem}
\begin{proof}
\leanok
\end{proof}

\begin{theorem}[Special values and symmetry]
\label{thm:jacobi-values}\uses{thm:jacobi-ode}
\lean{Polynomial.jacobi_comp_neg_X, Polynomial.eval_jacobi_neg, Polynomial.eval_jacobi_neg_one}\leanok
Let $R$ be a commutative $\Q$-algebra, let $\alpha,\beta\in R$, and let $n\in\N$.
The Jacobi polynomials satisfy
\[P_n^{(\alpha,\beta)}(-X)=(-1)^nP_n^{(\beta,\alpha)}(X),\qquad
P_n^{(\alpha,\beta)}(-1)=(-1)^n\frac{(\beta+1)_n}{n!}.\]

Reference: \cite[Table 18.6.1]{dlmf}.
\end{theorem}
\begin{proof}
\leanok
\end{proof}

\begin{supporting}[Parameter substitution and low degrees]

\begin{theorem}[Low degrees and parameter substitution]
\label{thm:jacobi-low-degree}\uses{def:jacobi}
\lean{Polynomial.jacobi_zero, Polynomial.jacobi_one, Polynomial.map_jacobi}\leanok
For a commutative $\Q$-algebra $R$ and $\alpha,\beta\in R$,
\[P_0^{(\alpha,\beta)}=1,\qquad
P_1^{(\alpha,\beta)}=(\alpha+1)+\tfrac12(\alpha+\beta+2)(X-1).\]
For every $\Q$-algebra homomorphism $f:R\to S$ and $n\in\N$, applying $f$ to the
coefficients of $P_n^{(\alpha,\beta)}$ gives $P_n^{(f(\alpha),f(\beta))}$.

Reference: \cite[18.5.7]{dlmf}.
\end{theorem}
\begin{proof}
\leanok
\end{proof}

\end{supporting}

\section{Polynomial and weighted identities}\label{sec:textbook-jacobi-identities}

The derivative shifts both parameters. The recurrence, differential equation, Rodrigues formula and terminating Gauss series give complementary descriptions of the same polynomials.

\begin{theorem}[Jacobi derivative]
\label{thm:jacobi-derivative}\uses{def:jacobi,def:polynomial-linear-maps}
\lean{Polynomial.derivative_jacobi_succ}\leanok
Let $R$ be a commutative $\Q$-algebra, let $\alpha,\beta\in R$, and let $n\in\N$. Then
\[(P_{n+1}^{(\alpha,\beta)})'=\tfrac12(n+\alpha+\beta+2)
 P_n^{(\alpha+1,\beta+1)}.\]

Reference: \cite[18.9.15]{dlmf}.
\end{theorem}
\begin{proof}
\leanok
\end{proof}

\begin{theorem}[Cleared Jacobi recurrence]
\label{thm:jacobi-recurrence}\uses{thm:jacobi-ode-unique,thm:jacobi-ode}
\lean{Polynomial.jacobi_recurrence}\leanok
Let $R$ be a commutative $\Q$-algebra, let $\alpha,\beta\in R$, and let $n\in\N$. Put $s=\alpha+\beta$. Then
\begin{align*}
&2(n+2)(n+s+2)(2n+s+2)P_{n+2}^{(\alpha,\beta)}\\
&\quad=(2n+s+3)\bigl[(2n+s+2)(2n+s+4)X+\alpha^2-\beta^2\bigr]
 P_{n+1}^{(\alpha,\beta)}\\
&\qquad\quad-2(n+\alpha+1)(n+\beta+1)(2n+s+4)P_n^{(\alpha,\beta)}.
\end{align*}
No factor depending on parameters is inverted, including at degenerate specialisations.

Reference: \cite[18.9.1--18.9.2]{dlmf}.
\end{theorem}
\begin{proof}
\leanok
\end{proof}

\begin{theorem}[Jacobi differential equation]
\label{thm:jacobi-ode}\uses{thm:jacobi-derivative,thm:jacobi-value-one,thm:jacobi-low-degree}
\lean{Polynomial.jacobi_differentialEquation}\leanok
Let $R$ be a commutative $\Q$-algebra, let $\alpha,\beta\in R$, and let $n\in\N$.
The Jacobi polynomial $y=P_n^{(\alpha,\beta)}(X)$ satisfies
\[(1-X^2)y''+[\beta-\alpha-(\alpha+\beta+2)X]y'
+n(n+\alpha+\beta+1)y=0.\]

Reference: \cite[Table 18.8.1]{dlmf}.
\end{theorem}
\begin{proof}
\leanok
\end{proof}

\begin{definition}[Jacobi weight and measure]
\label{def:jacobi-measure}\uses{def:open-interval-weight-measure}
\lean{Polynomial.jacobiWeight, Polynomial.jacobiMeasure}\leanok
For real $\alpha,\beta$, define the continuous map
$w_{\alpha,\beta}:(-1,1)\to\R_{\ge0}$ by $w_{\alpha,\beta}(x)=(1-x)^\alpha(1+x)^\beta$.
Let $\mu_{\alpha,\beta}=\mu_{w_{\alpha,\beta}}$ be the measure of
Definition~\ref{def:open-interval-weight-measure} for this weight on $(-1,1)$.

Reference: \cite[Table 18.3.1]{dlmf}.
\end{definition}

\begin{theorem}[Jacobi integral characterisation]
\label{thm:jacobi-measure-characterisation}
\uses{def:jacobi-measure,thm:open-interval-weight-measure}
\lean{Polynomial.integrable_jacobiMeasure_iff, Polynomial.integral_jacobiMeasure}\leanok
A real function $f$ is $\mu_{\alpha,\beta}$-integrable if and only if $x\mapsto w_{\alpha,\beta}(x)f(x)$ is Lebesgue-integrable on $(-1,1)$. When these equivalent integrability conditions hold,
\[\int f\,d\mu_{\alpha,\beta}=\int_{x\in(-1,1)}w_{\alpha,\beta}(x)f(x)\,dx.\]
\end{theorem}
\begin{proof}
\leanok
\end{proof}

\begin{theorem}[Weighted Jacobi Rodrigues formula]
\label{thm:jacobi-rodrigues}\uses{thm:jacobi-rodrigues-step,thm:jacobi-weight-calculus,def:open-domain-iterated-derivative,thm:open-domain-derivative,thm:jacobi-weight-basic}
\lean{Polynomial.hasDerivAt_jacobiWeight_mul_jacobi, Polynomial.hasIteratedDeriv_jacobiWeight, Polynomial.eval_jacobi_eq_rodrigues}\leanok
For all real $\alpha,\beta$, $n\ge0$, and $-1<x<1$, write
$w_{\alpha,\beta}(x)=(1-x)^\alpha(1+x)^\beta$. Then
\[\frac{d}{dx}\left(w_{\alpha+1,\beta+1}(x)P_n^{(\alpha+1,\beta+1)}(x)\right)
 =-2(n+1)w_{\alpha,\beta}(x)P_{n+1}^{(\alpha,\beta)}(x),\]
\[\frac{d^n}{dx^n}w_{\alpha+n,\beta+n}(x)
 =(-2)^n n!\,w_{\alpha,\beta}(x)P_n^{(\alpha,\beta)}(x).\]
The weight is strictly positive in the open interval, so division gives the
usual Rodrigues representation. No integrability or endpoint regularity
hypothesis is imposed.

Reference: \cite[18.9.16, 18.5.5 and Table 18.5.1]{dlmf}.
\end{theorem}
\begin{proof}
\leanok
\end{proof}

\begin{theorem}[Jacobi hypergeometric representation]
\label{thm:jacobi-representations}\uses{def:jacobi,thm:ordinary-hypergeometric}
\lean{Polynomial.eval_jacobi_eq_hypergeometric_sum, Polynomial.eval_jacobi_eq_ordinaryHypergeometric}\leanok
Let $R$ be a characteristic-zero field with a rational algebra structure, let $\alpha,\beta,x\in R$, and let $n\in\N$. Assume $(\alpha+1)_n\ne0$. Then
\[P_n^{(\alpha,\beta)}(x)=\frac{(\alpha+1)_n}{n!}
 \sum_{k=0}^n\frac{(-n)_k(n+\alpha+\beta+1)_k}{(\alpha+1)_k k!}
 \left(\frac{1-x}{2}\right)^k.\]
For real or complex scalars this equals
\[\frac{(\alpha+1)_n}{n!}\,
 {}_2F_1\left(-n,n+\alpha+\beta+1;\alpha+1;\frac{1-x}{2}\right),\]
where ${}_2F_1$ is the Gauss series of Definition~\ref{def:ordinary-hypergeometric}. The finite nonvanishing condition excludes
zeros among exactly the denominator factors needed through degree $n$;
$x$ is unrestricted.

Reference: \cite[18.5.7]{dlmf}.
\end{theorem}
\begin{proof}
\leanok
\end{proof}

\begin{supporting}[Rodrigues proofs, Beta and Euler integrals]

\begin{theorem}[Jacobi density characterisation]
\label{thm:jacobi-weight-characterisation}
\uses{def:jacobi-measure}
\lean{Polynomial.jacobiWeight_apply}\leanok
For every $x\in(-1,1)$, $w_{\alpha,\beta}(x)=(1-x)^\alpha(1+x)^\beta$.
\end{theorem}
\begin{proof}
\leanok
\end{proof}

\begin{theorem}[Algebraic Jacobi Rodrigues step]
\label{thm:jacobi-rodrigues-step}\uses{thm:jacobi-ode}
\lean{Polynomial.jacobi_rodrigues_step}\leanok
Let $R$ be a commutative $\Q$-algebra, let $\alpha,\beta\in R$, and let $n\in\N$. Put $Q=P_n^{(\alpha+1,\beta+1)}$. Then
\[(1-X^2)Q'+[\beta-\alpha-(\alpha+\beta+2)X]Q
  =-2(n+1)P_{n+1}^{(\alpha,\beta)}.\]
No parameter factor is cancelled in the final identity.

Reference: \cite[18.5.5 and 18.9.16]{dlmf}.
\end{theorem}
\begin{proof}
\leanok
\end{proof}

\begin{theorem}[Uniqueness at the Jacobi endpoint]
\label{thm:jacobi-ode-unique}\uses{def:polynomial-linear-maps}
\lean{Polynomial.eq_zero_of_jacobi_ode}\leanok
Let $R$ be a characteristic-zero commutative domain and a $\Q$-algebra.
Suppose $\alpha+k+1\ne0$ for every $k\in\N$. For any $\beta,\gamma\in R$,
a polynomial $p$ satisfying
\[(1-X^2)p''+[\beta-\alpha-(\alpha+\beta+2)X]p'+\gamma p=0,\qquad p(1)=0\]
is zero.

Reference: \cite[Table 18.8.1]{dlmf}.
\end{theorem}
\begin{proof}
\leanok
\end{proof}

\begin{theorem}[Pochhammer binomial series]
\label{thm:hypergeometric-binomial}
\lean{Complex.choose_add_nat_sub_one_eq,Complex.hasSum_pochhammer_mul_pow,Complex.summable_norm_pochhammer_mul_pow}\leanok
For $a,z\in\C$ with $|z|<1$,
\[
 (1-z)^{-a}=\sum_{n=0}^\infty\frac{(a)_n}{n!}z^n,
\]
with absolute convergence and the principal complex power.
Reference: the binomial expansion used in \cite[15.6.1]{dlmf}.
\end{theorem}
\begin{proof}
\leanok
\end{proof}

\begin{definition}[Euler Beta integral]
\label{def:euler-beta}
\lean{Complex.BetaParameter,Complex.betaIntegralOn}\leanok
For $a,b\in\C$ with $\Re a,\Re b>0$, define
\[
 B(a,b)=\int_{(0,1)} t^{a-1}(1-t)^{b-1}\,dt.
\]
The powers are principal complex powers of positive real bases; the integration
variable belongs to the open interval. Reference: \cite[5.12.1]{dlmf}.
\end{definition}

\begin{theorem}[Characterisation of the Euler Beta integral]
\label{thm:euler-beta}\uses{def:euler-beta,thm:gamma-function}
\lean{Complex.integrable_betaIntegralOn,Complex.betaIntegralOn_eq_Gamma,Complex.betaIntegralOn_ne_zero}\leanok
For every pair in the preceding domain the integrand is integrable, and
\[
 B(a,b)=\frac{\Gamma(a)\Gamma(b)}{\Gamma(a+b)}\ne0.
\]
All three Gamma arguments lie in the open right half-plane.
Reference: \cite[5.12.1]{dlmf}.
\end{theorem}
\begin{proof}
\leanok
\end{proof}

\begin{theorem}[Euler Beta moments]
\label{thm:hypergeometric-beta-moments}\uses{thm:euler-beta,thm:gauss-coefficients}
\lean{Complex.betaIntegralOn_add_nat_left,Complex.integral_betaWeight_mul_powOn,Complex.eulerGaussParameters,Complex.integral_eulerHypergeometric_termOn}\leanok
For $\Re b>0$, $\Re(c-b)>0$, and $n\in\N$,
\[
 B(b+n,c-b)=B(b,c-b)\frac{(b)_n}{(c)_n}.
\]
Thus integration of $t^n$ against $t^{b-1}(1-t)^{c-b-1}$, followed by
multiplication by $(a)_n/n!$, gives $B(b,c-b)$ times the $n$th ordinary
Gauss-series coefficient.
These integrands are integrable on $(0,1)$. More generally, the shifted
moment $\int_{(0,1)}t^{u-n-1}(1-t)^{v-1}t^n\,dt=B(u,v)$ is integrable for
$\Re u,\Re v>0$, even when the unshifted exponent would not be integrable.
The Euler parameters $(a,b,b+d)$ are admissible whenever $\Re b,\Re d>0$.
Reference: \cite[5.12.1, 15.6.1]{dlmf}.
\end{theorem}
\begin{proof}
\leanok
\end{proof}

\begin{theorem}[Euler integral for the Gauss function]
\label{thm:hypergeometric-euler-integral}\uses{thm:hypergeometric-binomial,thm:hypergeometric-beta-moments,thm:regularized-gauss,thm:hypergeometric-calculus}
\lean{Complex.hasSum_integral_binomialOn,Complex.integral_eulerHypergeometricOn,Complex.integral_eulerRegularizedHypergeometricOn,Complex.ofReal_ascPochhammer_eval,Real.eulerGaussParameters,Real.integral_eulerHypergeometricOn}\leanok
Let $a,b,c,z\in\C$, $\Re b>0$, $\Re(c-b)>0$, and $|z|<1$. Then
\[
 \int_0^1\frac{t^{b-1}(1-t)^{c-b-1}}{(1-zt)^a}\,dt
 = B(b,c-b)\,{}_2F_1(a,b;c;z)
 = \Gamma(b)\Gamma(c-b)\,\mathbf F(a,b;c;z).
\]
All integrals here are over $(0,1)$ and their integrability is asserted.
The real and complex definitions agree on real parameters and arguments, and the same integral formula holds over the reals. The dominated-convergence
lemma takes any complex $L^1(0,1)$ weight and asserts integrability of each term
and of the summed integrand, together with the sum identity.
Reference: \cite[15.6.1]{dlmf}, which prints the regularized normalization.
\end{theorem}
\begin{proof}
\leanok
\end{proof}

\begin{theorem}[Symmetric Jacobi coefficient parity]
\label{thm:jacobi-coefficient-parity}\uses{thm:jacobi-values}
\lean{Polynomial.coeff_comp_neg_X, Polynomial.coeff_jacobi_symmetric_eq_zero_of_parity}\leanok
The coefficient of $x^j$ in $P_n^{(\beta,\beta)}$ vanishes if $j$ and $n$ have
different parity, for every real $\beta$, including degenerate parameters.

Reference: \cite[Table 18.6.1]{dlmf}.
\end{theorem}
\begin{proof}
\leanok
\end{proof}

\end{supporting}

\section{Orthogonality and zeros}\label{sec:textbook-jacobi-orthogonality}

For $\alpha,\beta>-1$, the Jacobi weight is integrable and the polynomials are orthogonal. The degree-zero norm has a separate formula, valid also when $\alpha+\beta+1=0$.

\begin{theorem}[Jacobi weight integrability]
\label{thm:jacobi-integrable}\uses{thm:mathlib-power-integral,thm:jacobi-measure-characterisation}
\lean{Polynomial.integrable_jacobiWeight, Polynomial.integrable_jacobiMeasure}\leanok
For $\alpha,\beta>-1$, the weight is integrable on $(-1,1)$ and every
real-valued function continuous on $[-1,1]$ is $\mu_{\alpha,\beta}$-integrable.


Reference: \cite[Table 18.3.1]{dlmf}.
\end{theorem}
\begin{proof}
\leanok
\end{proof}

\begin{theorem}[Jacobi pairwise orthogonality]
\label{thm:jacobi-pairwise}\uses{thm:jacobi-weight-calculus,thm:jacobi-moments,thm:jacobi-sequence,thm:jacobi-ode,thm:jacobi-weight-closed,thm:open-domain-derivative,def:orthogonal-sequence}
\lean{Polynomial.integral_mul_jacobi_eq_zero, Polynomial.isOrthogonal_jacobiSequence}\leanok
For $\alpha,\beta>-1$ and distinct $m,n$,
$\int P_m^{(\alpha,\beta)}P_n^{(\alpha,\beta)}\,d\mu_{\alpha,\beta}=0$.
The sequence satisfies the general orthogonality predicate.

Reference: \cite[Table 18.3.1]{dlmf}.
\end{theorem}
\begin{proof}
\leanok
\end{proof}

Write $h_n=\int (P_n^{(\alpha,\beta)})^2\,d\mu_{\alpha,\beta}$ for the squared norm.

\begin{theorem}[Exact Jacobi squared norms]
\label{thm:jacobi-norm}\uses{thm:jacobi-energy,thm:jacobi-mass}
\lean{Polynomial.integral_sq_jacobi_zero, Polynomial.integral_sq_jacobi_succ}\leanok
For $\alpha,\beta>-1$,
\[h_0=\frac{2^{\alpha+\beta+1}\Gamma(\alpha+1)\Gamma(\beta+1)}{\Gamma(\alpha+\beta+2)},\]
and for $n\ge1$,
\[h_n=\frac{2^{\alpha+\beta+1}\Gamma(n+\alpha+1)\Gamma(n+\beta+1)}
 {(2n+\alpha+\beta+1)\Gamma(n+\alpha+\beta+1)n!}.\]


Reference: \cite[Table 18.3.1]{dlmf}.
\end{theorem}
\begin{proof}
\leanok
\end{proof}

\begin{theorem}[Real simple Jacobi zeros]
\label{thm:jacobi-real-zeros}\uses{thm:jacobi-pairwise,thm:jacobi-positive,thm:general-real-zeros}
\lean{Polynomial.splits_jacobi, Polynomial.card_roots_jacobi, Polynomial.separable_jacobi}\leanok
For $\alpha,\beta>-1$, $P_n^{(\alpha,\beta)}$ splits over $\mathbb{R}$, is separable, and has $n$ roots counted with multiplicity, hence $n$ distinct real zeros.

Reference: \cite[18.2(vi)]{dlmf}.
\end{theorem}
\begin{proof}
\leanok
\end{proof}

\begin{theorem}[Jacobi root location]
\label{thm:jacobi-root-location}\uses{thm:jacobi-pairwise,thm:jacobi-positive,thm:orthogonal-lower-degree,thm:jacobi-values}
\lean{Polynomial.integral_mul_jacobi_eq_zero_of_degree_lt, Polynomial.root_jacobi_mem_Ioo}\leanok
For $\alpha,\beta>-1$, $P_n$ is orthogonal to every polynomial of degree below $n$, and every real root of $P_n$ lies in $(-1,1)$.

Reference: \cite[18.2(i), (vi)]{dlmf}.
\end{theorem}
\begin{proof}
\leanok
\end{proof}

\begin{theorem}[Jacobi strict interlacing]
\label{thm:jacobi-interlacing}\uses{thm:general-outer-roots,thm:jacobi-positive,thm:jacobi-pairwise}
\lean{Polynomial.not_eval_jacobi_eq_zero_and_succ, Polynomial.exists_root_jacobi_succ_between, Polynomial.exists_root_jacobi_between_succ, Polynomial.exists_outer_roots_jacobi}\leanok
For $\alpha,\beta>-1$, consecutive members have no common zero and mutually separate each other's real zeros. For $n\ge1$, $P_{n+1}$ has zeros $a<b$ with every zero of $P_n$ strictly between them.

Reference: \cite[18.2(vi)]{dlmf}.
\end{theorem}
\begin{proof}
\leanok
\end{proof}

\begin{supporting}[Weight properties and finite moments]

\begin{theorem}[Jacobi weight positivity and support]
\label{thm:jacobi-weight-basic}\uses{def:jacobi-measure}
\lean{Polynomial.measurable_jacobiWeight, Polynomial.jacobiWeight_pos, Polynomial.jacobiWeight_nonneg, Polynomial.ae_mem_Ioo_jacobiMeasure}\leanok
For all real $\alpha,\beta$, the weight $w_{\alpha,\beta}$ is measurable
and strictly positive on its domain $(-1,1)$, and
$\mu_{\alpha,\beta}$-almost every point lies in $(-1,1)$.

Reference: \cite[Table 18.3.1]{dlmf}.
\end{theorem}
\begin{proof}
\leanok
\end{proof}

\begin{theorem}[Jacobi total mass]
\label{thm:jacobi-mass}\uses{thm:unit-jacobi-measure-characterisation,thm:euler-beta,thm:jacobi-integrable,thm:unit-jacobi-weight,thm:integrable-input}
\lean{Polynomial.integral_jacobiWeight, Polynomial.integral_one_jacobiMeasure}\leanok
For $\alpha,\beta>-1$,
\[\int_{-1}^1 w_{\alpha,\beta}(x)\,dx=\int1\,d\mu_{\alpha,\beta}
 =\frac{2^{\alpha+\beta+1}\Gamma(\alpha+1)\Gamma(\beta+1)}{\Gamma(\alpha+\beta+2)}.\]
This is nonsingular also when $\alpha+\beta+1=0$.

Reference: \cite[Table 18.3.1]{dlmf}.
\end{theorem}
\begin{proof}
\leanok
\end{proof}

\end{supporting}

\section{Jacobi expansions}\label{sec:textbook-jacobi-expansion}

The integral coefficients retain the polynomial normalization used above. Every function in the weighted $L^2$ space admits the resulting norm-convergent expansion.

\begin{definition}[Jacobi coefficients]
\label{def:jacobi-coefficients}\uses{thm:jacobi-positive,thm:jacobi-moments,thm:jacobi-sequence,def:integral-coefficient-input}
\lean{Polynomial.jacobiCoefficient, Polynomial.IntegralCoefficientInput.ofJacobiL2}\leanok
For $f\in\mathcal A_{\mu_{\alpha,\beta},P_n^{(\alpha,\beta)}}^{\R}$, define $a_n(f)=h_n^{-1}\int fP_n\,d\mu_{\alpha,\beta}$, where
$h_n=\int P_n^2\,d\mu_{\alpha,\beta}$.

Reference: \cite[18.2(viii)]{dlmf}.
\end{definition}

\begin{theorem}[Jacobi coefficient characterisation]
\label{thm:jacobi-coefficient-characterisation}\uses{def:jacobi-coefficients}
\lean{Polynomial.jacobiCoefficient_apply}\leanok
For $f\in\mathcal A_{\mu_{\alpha,\beta},P_n^{(\alpha,\beta)}}^{\R}$, $a_n(f)=h_n^{-1}\int fP_n\,d\mu_{\alpha,\beta}$, with an integrable numerator
and finite positive $h_n$. Under $\alpha,\beta>-1$, every $L^2(\mu_{\alpha,\beta})$ function belongs to this coefficient domain at every index, in the original polynomial normalization.
\end{theorem}
\begin{proof}
\leanok
\end{proof}

The normalized polynomials are $P_n^{(\alpha,\beta)}/\sqrt{h_n}$, viewed as weighted $L^2$ classes. The map $J_\mu$ denotes polynomial evaluation followed by passage to this class (Definition~\ref{def:general-polynomial-L2}).

\begin{theorem}[Jacobi completeness]
\label{thm:jacobi-basis}\uses{def:jacobi-normalised,thm:jacobi-pairwise,thm:compact-polynomial-dense,thm:general-hilbert-basis}
\lean{Polynomial.polynomialToL2_jacobi_denseRange, Polynomial.orthonormal_normalizedJacobi, Polynomial.hilbertBasisJacobi, Polynomial.hilbertBasisJacobi_apply}\leanok
For $\alpha,\beta>-1$, the normalised Jacobi polynomials form a Hilbert basis of real $L^2(\mu_{\alpha,\beta})$.

Reference: \cite[18.2(viii)]{dlmf}.
\end{theorem}
\begin{proof}
\leanok
\end{proof}

\begin{theorem}[Jacobi L2 expansions]
\label{thm:jacobi-L2-expansion}\uses{thm:jacobi-basis,thm:general-L2-expansion,def:jacobi-coefficients}
\lean{Polynomial.hasSum_jacobiCoefficient_polynomialToL2, Polynomial.tendsto_sum_jacobiCoefficient_polynomialToL2}\leanok
For $\alpha,\beta>-1$ and every $f\in L^2(\mu_{\alpha,\beta})$, the series $\sum_{n\ge0}a_n(f)J_\mu P_n$ sums to $f$. Its natural-order partial sums converge in $L^2$.

Reference: \cite[18.2(viii)]{dlmf}.
\end{theorem}
\begin{proof}
\leanok
\end{proof}

\begin{supporting}[Coefficient identification and weight calculus]

\begin{theorem}[Jacobi coefficients as general coefficients]
\label{thm:jacobi-coefficients-general}\uses{def:general-normalised-L2,def:jacobi-coefficients}
\lean{Polynomial.integralCoefficient_jacobiSequence}\leanok
For $\alpha,\beta>-1$, the integral coefficients of the Jacobi sequence for
$\mu_{\alpha,\beta}$ (Definition~\ref{def:general-normalised-L2}) are the
Jacobi coefficients $a_n(f)$, for $f\in\mathcal A_{\mu_{\alpha,\beta},P_n^{(\alpha,\beta)}}^{\R}$.

Reference: \cite[18.2(viii)]{dlmf}.
\end{theorem}
\begin{proof}
\leanok
\end{proof}

\begin{theorem}[Finite Jacobi moments]
\label{thm:jacobi-moments}\uses{thm:jacobi-integrable,thm:polynomial-inner}
\lean{Polynomial.hasFiniteMoments_jacobiMeasure, Polynomial.isFiniteMeasure_jacobiMeasure}\leanok
For $\alpha,\beta>-1$, $\mu_{\alpha,\beta}$ has finite polynomial moments and finite total mass.

Reference: \cite[Table 18.3.1]{dlmf}.
\end{theorem}
\begin{proof}
\leanok
\end{proof}

\begin{theorem}[Positive Jacobi pairing]
\label{thm:jacobi-positive}\uses{thm:jacobi-integrable,def:polynomial-positive,thm:jacobi-weight-basic}
\lean{Polynomial.hasPositiveDefiniteIntegral_jacobiMeasure}\leanok
For $\alpha,\beta>-1$ and $p\ne0$, $\int p^2\,d\mu_{\alpha,\beta}>0$.

Reference: \cite[18.2(i)]{dlmf}.
\end{theorem}
\begin{proof}
\leanok
\end{proof}

\begin{theorem}[Jacobi polynomial sequence]
\label{thm:jacobi-sequence}\uses{thm:jacobi-degree,thm:jacobi-low-degree}
\lean{Polynomial.coeff_jacobi_self_pos, Polynomial.degree_jacobi_of_neg_one_lt, Polynomial.jacobiSequence, Polynomial.jacobiSequence_apply}\leanok
For $\alpha,\beta>-1$, $P_n^{(\alpha,\beta)}$ has positive top coefficient and exact degree $n$. These polynomials define a degree-indexed sequence.

Reference: \cite[Table 18.3.1]{dlmf}.
\end{theorem}
\begin{proof}
\leanok
\end{proof}

\begin{definition}[Continuous closed-interval Jacobi weight]
\label{def:jacobi-weight-closed}
\lean{Polynomial.jacobiWeightClosed}\leanok
For $\alpha,\beta\ge0$, define the continuous map
$\overline w_{\alpha,\beta}:[-1,1]\to\R_{\ge0}$ by
$\overline w_{\alpha,\beta}(x)=(1-x)^\alpha(1+x)^\beta$, where $0^0=1$.

Reference: \cite[Table 18.3.1]{dlmf}.
\end{definition}

\begin{theorem}[Closed Jacobi weight characterisation]
\label{thm:jacobi-weight-closed}\uses{def:jacobi-weight-closed,def:jacobi-measure}
\lean{Polynomial.jacobiWeightClosed_restrict, Polynomial.jacobiWeightClosed_one, Polynomial.jacobiWeightClosed_neg_one}\leanok
For $\alpha,\beta\ge0$, the weight $\overline w_{\alpha,\beta}$ agrees with
$w_{\alpha,\beta}$ on $(-1,1)$, and
\[\overline w_{\alpha,\beta}(1)=\begin{cases}2^\beta&\alpha=0,\\0&\alpha>0,\end{cases}
\qquad
\overline w_{\alpha,\beta}(-1)=\begin{cases}2^\alpha&\beta=0,\\0&\beta>0.\end{cases}\]
\end{theorem}
\begin{proof}
\leanok
\end{proof}

\begin{theorem}[Jacobi weight calculus]
\label{thm:jacobi-weight-calculus}\uses{def:jacobi-measure,def:open-domain-derivative}
\lean{Polynomial.jacobiWeight_add_one, Polynomial.hasDerivAt_jacobiWeight_add_one, Polynomial.tendsto_jacobiWeight_add_one_one, Polynomial.tendsto_jacobiWeight_add_one_neg_one}\leanok
On $(-1,1)$, for arbitrary real $\alpha,\beta$,
\[w_{\alpha+1,\beta+1}=(1-x^2)w_{\alpha,\beta},\qquad
w_{\alpha+1,\beta+1}'=[\beta-\alpha-(\alpha+\beta+2)x]w_{\alpha,\beta}.\]
If $\alpha>-1$, the right endpoint limit of the shifted weight is zero,
with no restriction on $\beta$. If $\beta>-1$, its left endpoint limit is zero,
with no restriction on $\alpha$.

Reference: \cite[Tables 18.3.1 and 18.8.1]{dlmf}.
\end{theorem}
\begin{proof}
\leanok
\end{proof}

\begin{theorem}[Jacobi energy and norm shift]
\label{thm:jacobi-energy}\uses{thm:jacobi-weight-calculus,thm:jacobi-integrable,thm:jacobi-ode,thm:jacobi-weight-closed,thm:open-domain-derivative}
\lean{Polynomial.integral_sq_derivative_jacobi, Polynomial.integral_sq_jacobi_succ_shift}\leanok
For $\alpha,\beta>-1$, set $h_n^{(\alpha,\beta)}=\int(P_n^{(\alpha,\beta)})^2\,d\mu_{\alpha,\beta}$.
Then
\[\int ((P_n^{(\alpha,\beta)})')^2\,d\mu_{\alpha+1,\beta+1}
 =n(n+\alpha+\beta+1)h_n^{(\alpha,\beta)},\]
\[h_{n+1}^{(\alpha,\beta)}=\frac{n+\alpha+\beta+2}{4(n+1)}h_n^{(\alpha+1,\beta+1)}.\]

Reference: \cite[18.9.15, Tables 18.3.1 and 18.8.1]{dlmf}.
\end{theorem}
\begin{proof}
\leanok
\end{proof}

\begin{definition}[Normalised Jacobi classes]
\label{def:jacobi-normalised}\uses{thm:jacobi-sequence,thm:jacobi-moments,def:general-normalised-L2,thm:jacobi-positive,thm:general-squared-norm}
\lean{Polynomial.normalizedJacobi}\leanok
For $\alpha,\beta>-1$, set $e_n=h_n^{-1/2}J_{\mu_{\alpha,\beta}}P_n^{(\alpha,\beta)}$, where $h_n$ is the squared integral norm.

Reference: \cite[Table 18.3.1]{dlmf}.
\end{definition}

\end{supporting}

\chapter{Ultraspherical, Legendre and Chebyshev as Jacobi cases}\label{chap:classical}

\rvsummary{Ultraspherical, Legendre and Chebyshev polynomials as Jacobi cases, with exact normalising bridges.}

The Jacobi family includes Gegenbauer, Legendre and both Chebyshev families after the normalizations described here. These identifications connect their formulas and weighted expansions.

\section{Gegenbauer polynomials}\label{sec:textbook-classical-gegenbauer}

The ultraspherical family is a rescaling of the symmetric Jacobi family. Its polynomial definition remains meaningful at degenerate parameters, including $\lambda=0$.

\begin{definition}[Gegenbauer polynomials]
\label{def:gegenbauer}
\lean{Polynomial.gegenbauer}\leanok
Gegenbauer polynomials are also called \emph{ultraspherical polynomials};
see \cite[\S18.1(ii)]{dlmf}.
Let $R$ be a commutative $\Q$-algebra, let $\lambda\in R$, and let $n\in\N$. Define
\[C_n^{(\lambda)}(X)=\sum_{k=0}^{\lfloor n/2\rfloor}
 \frac{(-1)^k(\lambda)_{n-k}}{k!(n-2k)!}(2X)^{n-2k}.\]

Reference: \cite[18.5.10]{dlmf}.
\end{definition}

\begin{theorem}[Gegenbauer degree and leading term]
\label{thm:gegenbauer-degree}\uses{def:gegenbauer}
\lean{Polynomial.natDegree_gegenbauer_le, Polynomial.degree_gegenbauer_le, Polynomial.coeff_gegenbauer_self, Polynomial.degree_gegenbauer_eq, Polynomial.degree_gegenbauer_eq_iff, Polynomial.leadingCoeff_gegenbauer}\leanok
Let $R$ be a commutative $\Q$-algebra, let $\lambda\in R$, and let $n\in\N$.
The Gegenbauer polynomial $C_n^{(\lambda)}(X)$ has degree at most $n$, and its coefficient of $X^n$ is
\[
 c_n=\frac{2^n(\lambda)_n}{n!}.
\]
It has degree exactly $n$ if and only if $c_n\ne0$. In that case, its leading coefficient is $c_n$.

Reference: \cite[Table 18.3.1]{dlmf}.
\end{theorem}
\begin{proof}
\leanok
\end{proof}

\begin{theorem}[Zero Gegenbauer parameter]
\label{thm:gegenbauer-zero-parameter}\uses{def:gegenbauer}
\lean{Polynomial.gegenbauer_zero_parameter}\leanok
For a commutative $\Q$-algebra $R$, $C_n^{(0)}=0$ for every $n\ge1$.
The renormalised family is
treated after the recurrence (Theorem~\ref{thm:gegenbauer-parameter-limit}).

Reference: \cite[18.5.10]{dlmf}.
\end{theorem}
\begin{proof}
\leanok
\end{proof}

\begin{theorem}[Gegenbauer derivative]
\label{thm:gegenbauer-derivative}\uses{def:gegenbauer,def:polynomial-linear-maps}
\lean{Polynomial.derivative_gegenbauer_succ}\leanok
Let $R$ be a commutative $\Q$-algebra, let $\lambda\in R$, and let $n\in\N$. Then
\[(C_{n+1}^{(\lambda)})'=2\lambda C_n^{(\lambda+1)}.\]

Reference: \cite[18.9.19]{dlmf}.
\end{theorem}
\begin{proof}
\leanok
\end{proof}

\begin{theorem}[Cleared Gegenbauer parameter conversion]
\label{thm:gegenbauer-conversion}\uses{thm:gegenbauer-derivative,thm:gegenbauer-values-zero,thm:gegenbauer-low-degree}
\lean{Polynomial.gegenbauer_conversion_zero, Polynomial.gegenbauer_conversion_one, Polynomial.gegenbauer_conversion}\leanok
Let $R$ be a commutative $\Q$-algebra, let $\lambda\in R$, and let $n\in\N$. Then
\[(n+\lambda+2)C_{n+2}^{(\lambda)}
=\lambda\bigl(C_{n+2}^{(\lambda+1)}-C_n^{(\lambda+1)}\bigr).\]
The initial identities are $\lambda C_0^{(\lambda)}=\lambda C_0^{(\lambda+1)}$
and $(\lambda+1)C_1^{(\lambda)}=\lambda C_1^{(\lambda+1)}$.
No parameter is divided by.

Reference: \cite[18.9.7]{dlmf}.
\end{theorem}
\begin{proof}
\leanok
\end{proof}

\begin{theorem}[Cleared Gegenbauer recurrence]
\label{thm:gegenbauer-recurrence}\uses{thm:gegenbauer-conversion}
\lean{Polynomial.gegenbauer_recurrence}\leanok
Let $R$ be a commutative $\Q$-algebra, let $\lambda\in R$, and let $n\in\N$. Then
\[(n+2)C_{n+2}^{(\lambda)}-2(n+\lambda+1)X C_{n+1}^{(\lambda)}
 +(n+2\lambda)C_n^{(\lambda)}=0.\]
No parameter is divided by.

Reference: \cite[Table 18.9.1]{dlmf}.
\end{theorem}
\begin{proof}
\leanok
\end{proof}

\begin{theorem}[Gegenbauer differential equation]
\label{thm:gegenbauer-ode}\uses{thm:gegenbauer-derivative,thm:gegenbauer-values-zero,thm:gegenbauer-low-degree}
\lean{Polynomial.gegenbauer_differentialEquation}\leanok
Let $R$ be a commutative $\Q$-algebra, let $\lambda\in R$, and let $n\in\N$. The Gegenbauer polynomial $y=C_n^{(\lambda)}(X)$ satisfies
\[(1-X^2)y''-(2\lambda+1)Xy'+n(n+2\lambda)y=0.\]


Reference: \cite[Table 18.8.1]{dlmf}.
\end{theorem}
\begin{proof}
\leanok
\end{proof}

\begin{theorem}[Gegenbauer endpoint value]
\label{thm:gegenbauer-value-one}\uses{thm:gegenbauer-recurrence,thm:jacobi-value-one,thm:jacobi-low-degree}
\lean{Polynomial.eval_gegenbauer_one}\leanok
Let $R$ be a commutative $\Q$-algebra, let $\lambda\in R$, and let $n\in\N$. Then
$C_n^{(\lambda)}(1)=(2\lambda)_n/n!$.

Reference: \cite[Table 18.6.1]{dlmf}.
\end{theorem}
\begin{proof}
\leanok
\end{proof}

\begin{theorem}[Gegenbauer parity]
\label{thm:gegenbauer-parity}\uses{def:gegenbauer}
\lean{Polynomial.gegenbauer_comp_neg_X, Polynomial.eval_gegenbauer_neg}\leanok
Let $R$ be a commutative $\Q$-algebra, let $\lambda\in R$, and let $n\in\N$. Then
$C_n^{(\lambda)}(-X)=(-1)^n C_n^{(\lambda)}(X)$.

Reference: \cite[Table 18.6.1]{dlmf}.
\end{theorem}
\begin{proof}
\leanok
\end{proof}

\begin{theorem}[Gegenbauer--Jacobi bridge]
\label{thm:gegenbauer-jacobi}\uses{thm:gegenbauer-ode,thm:gegenbauer-value-one,thm:jacobi-ode-unique,thm:jacobi-ode}
\lean{Polynomial.gegenbauer_jacobi}\leanok
Let $R$ be a commutative $\Q$-algebra, let $\lambda\in R$, and let $n\in\N$. Then
\[(\lambda+\tfrac12)_n C_n^{(\lambda)}
 =(2\lambda)_n P_n^{(\lambda-1/2,\lambda-1/2)}.\]
No parameter factor is cancelled, so the identity includes singular
specialisations.

Reference: \cite[18.7.1]{dlmf}.
\end{theorem}
\begin{proof}
\leanok
\end{proof}

\begin{theorem}[Gegenbauer hypergeometric representations]
\label{thm:gegenbauer-hypergeometric}\uses{thm:gegenbauer-jacobi,thm:jacobi-representations,thm:hypergeometric-inverse-square}
\lean{Polynomial.eval_gegenbauer_eq_ordinaryHypergeometric, Polynomial.eval_gegenbauer_eq_ordinaryHypergeometric_inv_sq}\leanok
Let $\lambda$ and $x$ be real or complex and $n\in\N$. If
$(\lambda+\tfrac12)_n\ne0$, then
\[C_n^{(\lambda)}(x)=\frac{(2\lambda)_n}{n!}\,
 {}_2F_1\!\left(-n,n+2\lambda;\lambda+\tfrac12;\frac{1-x}2\right).\]
If $x\ne0$ and $(1-\lambda-n)_{\lfloor n/2\rfloor}\ne0$, then
\[C_n^{(\lambda)}(x)=(2x)^n\frac{(\lambda)_n}{n!}\,
 {}_2F_1\!\left(-\tfrac n2,\tfrac{1-n}2;1-\lambda-n;\frac1{x^2}\right).\]
Here ${}_2F_1$ is the Gauss series of Definition~\ref{def:ordinary-hypergeometric}. The conditions exclude exactly the
denominator factors that occur. The second condition is satisfied at
$\lambda=0$ and $\lambda=1$ for every $n$.

Reference: \cite[18.5.9 and 18.5.10]{dlmf}.
\end{theorem}
\begin{proof}
\leanok
\end{proof}

\begin{supporting}[Low degrees, parameter limits and hypergeometric reindexing]

\begin{theorem}[Low degrees and parameter substitution]
\label{thm:gegenbauer-low-degree}\uses{def:gegenbauer}
\lean{Polynomial.gegenbauer_zero, Polynomial.gegenbauer_one, Polynomial.map_gegenbauer}\leanok
For a commutative $\Q$-algebra $R$ and $\lambda\in R$,
$C_0^{(\lambda)}=1$ and $C_1^{(\lambda)}=2\lambda X$. For every $\Q$-algebra
homomorphism $f:R\to S$ and $n\in\N$, applying $f$ to the coefficients of
$C_n^{(\lambda)}$ gives $C_n^{(f(\lambda))}$.

Reference: \cite[18.5.10]{dlmf}.
\end{theorem}
\begin{proof}
\leanok
\end{proof}

\begin{theorem}[Gegenbauer values at zero]
\label{thm:gegenbauer-values-zero}\uses{def:gegenbauer}
\lean{Polynomial.eval_gegenbauer_even_zero, Polynomial.eval_gegenbauer_odd_zero}\leanok
Let $R$ be a commutative $\Q$-algebra, let $\lambda\in R$, and let $m\in\N$. Then
\[C_{2m}^{(\lambda)}(0)=(-1)^m\frac{(\lambda)_m}{m!},\qquad
C_{2m+1}^{(\lambda)}(0)=0.\]


Reference: \cite[Table 18.6.1]{dlmf}.
\end{theorem}
\begin{proof}
\leanok
\end{proof}

\begin{definition}[Reduced Gegenbauer polynomials]
\label{def:gegenbauer-reduced}
\lean{Polynomial.gegenbauerReduced}\leanok
For a commutative $\Q$-algebra $R$ and $\lambda\in R$, normalize
$R_0^{(\lambda)}=1$. For $n\ge1$, define
\[R_n^{(\lambda)}(X)=\sum_{k=0}^{\lfloor n/2\rfloor}
 \frac{(-1)^k(\lambda+1)_{n-k-1}}{k!(n-2k)!}(2X)^{n-2k}.\]

Reference: \cite[18.5.10]{dlmf}, \cite[\S18.7(iii)]{dlmf}.
\end{definition}

\begin{theorem}[Factoring out the parameter]
\label{thm:gegenbauer-reduced-factor}\uses{def:gegenbauer, def:gegenbauer-reduced}
\lean{Polynomial.gegenbauerReduced_zero, Polynomial.gegenbauer_eq_C_mul_gegenbauerReduced, Polynomial.map_gegenbauerReduced}\leanok
For a commutative $\Q$-algebra $R$ and $\lambda\in R$,
$R_0^{(\lambda)}=1$, and $C_n^{(\lambda)}=\lambda R_n^{(\lambda)}$ for $n\ge1$.
For every $n\ge0$ and every $\Q$-algebra homomorphism $f:R\to S$, applying $f$ to the coefficients of $R_n^{(\lambda)}$ gives
$R_n^{(f(\lambda))}$.

Reference: \cite[18.5.10]{dlmf}.
\end{theorem}
\begin{proof}
\leanok
\end{proof}

\begin{theorem}[Reduced Gegenbauer recurrence]
\label{thm:gegenbauer-reduced-recurrence}\uses{thm:gegenbauer-recurrence,thm:gegenbauer-reduced-factor}
\lean{Polynomial.gegenbauerReduced_recurrence}\leanok
For a commutative $\Q$-algebra $R$, $\lambda\in R$ and $n\ge1$,
\[(n+2)R_{n+2}^{(\lambda)}-2(n+\lambda+1)XR_{n+1}^{(\lambda)}
 +(n+2\lambda)R_n^{(\lambda)}=0.\]

Reference: \cite[Table 18.9.1]{dlmf}.
\end{theorem}
\begin{proof}
\leanok
\end{proof}

\begin{theorem}[Chebyshev T at parameter zero]
\label{thm:gegenbauer-reduced-chebyshev}\uses{def:chebyshev-mathlib,thm:gegenbauer-reduced-recurrence,thm:chebyshev-recurrence}
\lean{Polynomial.natCast_mul_gegenbauerReduced_zero}\leanok
For a commutative $\Q$-algebra $R$ and $n\ge1$, $nR_n^{(0)}=2T_n$. When $n$ is
invertible, $R_n^{(0)}=(2/n)T_n$.

Reference: \cite[18.1.1]{dlmf}.
\end{theorem}
\begin{proof}
\leanok
\end{proof}

\begin{theorem}[Gegenbauer polynomials divided by the parameter]
\label{thm:gegenbauer-parameter-limit}\uses{thm:gegenbauer-reduced-chebyshev}
\lean{Polynomial.continuous_eval_gegenbauerReduced, Polynomial.tendsto_eval_gegenbauer_div_nhdsNE_zero}\leanok
For $n\ge1$ and $x\in\R$, $\lambda\mapsto R_n^{(\lambda)}(x)$ is continuous on
$\R$, and
\[\lim_{\lambda\to0,\ \lambda\ne0}\frac{C_n^{(\lambda)}(x)}{\lambda}
 =\frac2n\,T_n(x).\]

Reference: \cite[\S18.7(iii)]{dlmf}.
\end{theorem}
\begin{proof}
\leanok
\end{proof}

\begin{theorem}[Quadratic Gauss reindexing]
\label{thm:hypergeometric-inverse-square}\uses{thm:ordinary-hypergeometric}
\lean{Polynomial.four_pow_mul_ascPochhammer_eval_mul_eval_add_half, Polynomial.two_mul_pow_mul_ordinaryHypergeometric_inv_sq}\leanok
Over a characteristic-zero field, $4^k(a)_k(a+\tfrac12)_k=(2a)_{2k}$ for all $a$
and $k\in\N$. For real or complex $c$, $x\ne0$, and $n\in\N$ with
$(c)_{\lfloor n/2\rfloor}\ne0$,
\[(2x)^n\,{}_2F_1\!\left(-\tfrac n2,\tfrac{1-n}2;c;\frac1{x^2}\right)
 =\sum_{k=0}^{\lfloor n/2\rfloor}\frac{n!}{k!\,(n-2k)!\,(c)_k}(2x)^{n-2k},\]
where ${}_2F_1$ is the Gauss series of Definition~\ref{def:ordinary-hypergeometric}. Every denominator factor in the finite sum is nonzero.

Reference: \cite[5.2.8 and 15.2.1]{dlmf}.
\end{theorem}
\begin{proof}
\leanok
\end{proof}

\end{supporting}

\section{Gegenbauer orthogonality and expansions}\label{sec:textbook-classical-gegenbauer-integrals}

The weight is $(1-x^2)^{\lambda-1/2}$ on $(-1,1)$. Orthogonal expansions require $\lambda>-1/2$ and $\lambda\ne0$; the measure itself is still meaningful at zero.

\begin{definition}[Gegenbauer measure]
\label{def:gegenbauer-measure}\uses{def:jacobi-measure}
\lean{Polynomial.gegenbauerWeight, Polynomial.gegenbauerMeasure}\leanok
For real $\lambda$, let $w_\lambda=w_{\lambda-1/2,\lambda-1/2}:(-1,1)\to\R_{\ge0}$
be the Jacobi weight of Definition~\ref{def:jacobi-measure}, and let
$\mu_\lambda=\mu_{w_\lambda}$ be the measure of
Definition~\ref{def:open-interval-weight-measure}.

Reference: \cite[Table 18.3.1]{dlmf}.
\end{definition}

\begin{theorem}[Gegenbauer density and integral characterisation]
\label{thm:gegenbauer-measure-characterisation}\uses{def:gegenbauer-measure,thm:open-interval-weight-measure,thm:jacobi-weight-characterisation}
\lean{Polynomial.gegenbauerWeight_apply,Polynomial.integrable_gegenbauerMeasure_iff,Polynomial.integral_gegenbauerMeasure}\leanok
For every $x\in(-1,1)$, $w_\lambda(x)=(1-x^2)^{\lambda-1/2}$. A real function $f$ is $\mu_\lambda$-integrable if and only if
$w_\lambda f$ is integrable on the open interval. When these equivalent integrability conditions hold,
\[
 \int f\,d\mu_\lambda=\int_{x\in(-1,1)}w_\lambda(x)f(x)\,dx.
\]
Reference: \cite[Table 18.3.1]{dlmf}.
\end{theorem}
\begin{proof}
\leanok
\end{proof}

\begin{theorem}[Gegenbauer orthogonality]
\label{thm:gegenbauer-orthogonal}\uses{thm:gegenbauer-moments,thm:gegenbauer-sequence,thm:jacobi-pairwise}
\lean{Polynomial.integral_mul_gegenbauer_eq_zero, Polynomial.isOrthogonal_gegenbauerSequence}\leanok
For $\lambda>-1/2$ and $m\ne n$,
\[\int C_m^{(\lambda)}C_n^{(\lambda)}\,d\mu_\lambda=0.\]
For $\lambda\ne0$ the sequence satisfies the general orthogonality predicate.
The pairwise integral identity itself also holds at $\lambda=0$.

Reference: \cite[Table 18.3.1]{dlmf}.
\end{theorem}
\begin{proof}
\leanok
\end{proof}

\begin{theorem}[Exact Gegenbauer squared norms]
\label{thm:gegenbauer-norm}\uses{thm:gegenbauer-norm-transfer,thm:jacobi-norm}
\lean{Polynomial.integral_sq_gegenbauer_zero, Polynomial.integral_sq_gegenbauer_succ}\leanok
For $\lambda>-1/2$, set $h_n=\int(C_n^{(\lambda)})^2\,d\mu_\lambda$. Then
\[h_0=\frac{\sqrt\pi\,\Gamma(\lambda+1/2)}{\Gamma(\lambda+1)}.\]
For $n\ge1$ and $\lambda\ne0$,
\[h_n=\frac{2^{1-2\lambda}\pi\,\Gamma(n+2\lambda)}
 {(n+\lambda)\Gamma(\lambda)^2n!}.\]
The separate degree-zero formula also covers $\lambda=0$.

Reference: \cite[Table 18.3.1]{dlmf}.
\end{theorem}
\begin{proof}
\leanok
\end{proof}

\begin{definition}[Gegenbauer coefficients]
\label{def:gegenbauer-coefficients}\uses{thm:gegenbauer-moments,thm:gegenbauer-sequence,def:integral-coefficient-input}
\lean{Polynomial.gegenbauerCoefficient, Polynomial.IntegralCoefficientInput.ofGegenbauerL2}\leanok
For $f\in\mathcal A_{\mu_\lambda,C_n^{(\lambda)}}^{\R}$, define $a_n(f)=h_n^{-1}\int fC_n^{(\lambda)}\,d\mu_\lambda$, where
$h_n=\int(C_n^{(\lambda)})^2\,d\mu_\lambda$.

Reference: \cite[18.2(viii)]{dlmf}.
\end{definition}

\begin{theorem}[Gegenbauer coefficient characterisation]
\label{thm:gegenbauer-coefficient-characterisation}\uses{def:gegenbauer-coefficients}
\lean{Polynomial.gegenbauerCoefficient_apply}\leanok
For $f\in\mathcal A_{\mu_\lambda,C_n^{(\lambda)}}^{\R}$, $a_n(f)=h_n^{-1}\int fC_n^{(\lambda)}\,d\mu_\lambda$, with an integrable numerator
and finite positive $h_n$. Under $\lambda>-1/2$, $\lambda\ne0$, every $L^2(\mu_\lambda)$ function belongs to this coefficient domain at every index, in the original polynomial normalization.
\end{theorem}
\begin{proof}
\leanok
\end{proof}

Normalization means dividing $C_n^{(\lambda)}$ by $\sqrt{h_n}$; $J_\mu$ denotes its weighted $L^2$ class.

\begin{theorem}[Gegenbauer completeness]
\label{thm:gegenbauer-basis}\uses{def:gegenbauer-normalised,thm:gegenbauer-orthogonal,thm:compact-polynomial-dense,thm:general-hilbert-basis}
\lean{Polynomial.polynomialToL2_gegenbauer_denseRange, Polynomial.orthonormal_normalizedGegenbauer, Polynomial.hilbertBasisGegenbauer, Polynomial.hilbertBasisGegenbauer_apply}\leanok
For $\lambda>-1/2$, $\lambda\ne0$, the normalised Gegenbauer polynomials
form a Hilbert basis of real $L^2(\mu_\lambda)$.

Reference: \cite[18.2(viii)]{dlmf}.
\end{theorem}
\begin{proof}
\leanok
\end{proof}

\begin{theorem}[Gegenbauer L2 expansions]
\label{thm:gegenbauer-L2-expansion}\uses{thm:gegenbauer-basis,thm:general-L2-expansion,def:gegenbauer-coefficients}
\lean{Polynomial.hasSum_gegenbauerCoefficient_polynomialToL2, Polynomial.tendsto_sum_gegenbauerCoefficient_polynomialToL2}\leanok
For $\lambda>-1/2$, $\lambda\ne0$, and every $f\in L^2(\mu_\lambda)$,
the series $\sum_{n\ge0}a_n(f)J_{\mu_\lambda}C_n^{(\lambda)}$ sums to $f$.
Its natural-order partial sums converge in $L^2$.

Reference: \cite[18.2(viii)]{dlmf}.
\end{theorem}
\begin{proof}
\leanok
\end{proof}

\begin{theorem}[Gegenbauer root location]
\label{thm:gegenbauer-root-location}\uses{thm:jacobi-root-location,thm:gegenbauer-orthogonal}
\lean{Polynomial.integral_mul_gegenbauer_eq_zero_of_degree_lt, Polynomial.root_gegenbauer_mem_Ioo}\leanok
For $\lambda>-1/2$, $\lambda\ne0$, $C_n^{(\lambda)}$ is orthogonal to every
polynomial of degree below $n$, and every real root lies in $(-1,1)$.

Reference: \cite[18.2(i), (vi)]{dlmf}.
\end{theorem}
\begin{proof}
\leanok
\end{proof}

\begin{theorem}[Gegenbauer strict interlacing]
\label{thm:gegenbauer-interlacing}\uses{thm:general-outer-roots,thm:gegenbauer-orthogonal}
\lean{Polynomial.not_eval_gegenbauer_eq_zero_and_succ, Polynomial.exists_root_gegenbauer_succ_between, Polynomial.exists_root_gegenbauer_between_succ, Polynomial.exists_outer_roots_gegenbauer}\leanok
For $\lambda>-1/2$, $\lambda\ne0$, consecutive Gegenbauer polynomials have
no common zero and mutually separate each other's real zeros.
For $n\ge1$, $C_{n+1}^{(\lambda)}$ has roots $a<b$ with every root of
$C_n^{(\lambda)}$ strictly between them.

Reference: \cite[18.2(vi)]{dlmf}.
\end{theorem}
\begin{proof}
\leanok
\end{proof}

\begin{supporting}[Measure bridges, norm transfer and normalization]

\begin{theorem}[Gegenbauer coefficients as general coefficients]
\label{thm:gegenbauer-coefficients-general}\uses{def:general-normalised-L2,def:gegenbauer-coefficients}
\lean{Polynomial.integralCoefficient_gegenbauerSequence}\leanok
For $\lambda>-1/2$, $\lambda\ne0$, the integral coefficients of the Gegenbauer
sequence for $\mu_\lambda$ (Definition~\ref{def:general-normalised-L2}) are the
Gegenbauer coefficients $a_n(f)$, for $f\in\mathcal A_{\mu_\lambda,C_n^{(\lambda)}}^{\R}$.

Reference: \cite[18.2(viii)]{dlmf}.
\end{theorem}
\begin{proof}
\leanok
\end{proof}

\begin{theorem}[Gegenbauer--Jacobi measure bridge]
\label{thm:gegenbauer-measure-bridge}\uses{def:gegenbauer-measure,thm:jacobi-weight-basic}
\lean{Polynomial.gegenbauerWeight_eq_jacobiWeight, Polynomial.gegenbauerMeasure_eq_jacobiMeasure, Polynomial.ae_mem_Ioo_gegenbauerMeasure}\leanok
For every real $\lambda$, the Gegenbauer weight equals
$w_{\lambda-1/2,\lambda-1/2}$ on $(-1,1)$, and
$\mu_\lambda=\mu_{\lambda-1/2,\lambda-1/2}$.
The measure assigns zero mass to $\mathbb{R}\setminus(-1,1)$.

Reference: \cite[Table 18.3.1]{dlmf}.
\end{theorem}
\begin{proof}
\leanok
\end{proof}

\begin{theorem}[Gegenbauer moments and positivity]
\label{thm:gegenbauer-moments}\uses{thm:gegenbauer-measure-bridge,thm:jacobi-moments,thm:jacobi-positive}
\lean{Polynomial.integrable_gegenbauerWeight, Polynomial.integrable_gegenbauerMeasure, Polynomial.hasFiniteMoments_gegenbauerMeasure, Polynomial.isFiniteMeasure_gegenbauerMeasure, Polynomial.hasPositiveDefiniteIntegral_gegenbauerMeasure}\leanok
For $\lambda>-1/2$, the weight is integrable, every real-valued function continuous on
$[-1,1]$ is $\mu_\lambda$-integrable, all polynomial moments are finite, and
$\int p^2\,d\mu_\lambda>0$ for every nonzero polynomial $p$.
The measure has finite total mass. These assertions include $\lambda=0$.

Reference: \cite[Table 18.3.1]{dlmf}.
\end{theorem}
\begin{proof}
\leanok
\end{proof}

\begin{theorem}[Divided Gegenbauer--Jacobi bridge]
\label{thm:gegenbauer-divided-bridge}\uses{thm:gegenbauer-jacobi}
\lean{Polynomial.gegenbauer_eq_jacobi_mul, Polynomial.gegenbauer_jacobi_factor_ne_zero}\leanok
For $\lambda>-1/2$, write $c_n=(2\lambda)_n/(\lambda+1/2)_n$. Then
\[C_n^{(\lambda)}=c_nP_n^{(\lambda-1/2,\lambda-1/2)}.\]
If additionally $\lambda\ne0$, every $c_n$ is nonzero, including when
$-1/2<\lambda<0$.

Reference: \cite[18.7.1]{dlmf}.
\end{theorem}
\begin{proof}
\leanok
\end{proof}

\begin{theorem}[Gegenbauer polynomial sequence]
\label{thm:gegenbauer-sequence}\uses{thm:gegenbauer-divided-bridge,thm:jacobi-sequence}
\lean{Polynomial.degree_gegenbauer_of_neg_half_lt, Polynomial.gegenbauerSequence, Polynomial.gegenbauerSequence_apply}\leanok
For $\lambda>-1/2$, $\lambda\ne0$, $C_n^{(\lambda)}$ has exact degree $n$
and these polynomials define a degree-indexed sequence.

Reference: \cite[Table 18.3.1]{dlmf}.
\end{theorem}
\begin{proof}
\leanok
\end{proof}

\begin{theorem}[Gegenbauer norm transfer]
\label{thm:gegenbauer-norm-transfer}\uses{thm:gegenbauer-divided-bridge,thm:gegenbauer-measure-bridge}
\lean{Polynomial.integral_sq_gegenbauer_eq_jacobi}\leanok
For $\lambda>-1/2$, the squared norm of $C_n^{(\lambda)}$ is $c_n^2$ times
the squared norm of $P_n^{(\lambda-1/2,\lambda-1/2)}$, with $c_n$ as in the
divided bridge.

Reference: \cite[18.7.1]{dlmf}.
\end{theorem}
\begin{proof}
\leanok
\end{proof}

\begin{theorem}[Real simple Gegenbauer zeros]
\label{thm:gegenbauer-real-zeros}\uses{thm:gegenbauer-orthogonal,thm:general-real-zeros}
\lean{Polynomial.splits_gegenbauer, Polynomial.card_roots_gegenbauer, Polynomial.separable_gegenbauer}\leanok
For $\lambda>-1/2$, $\lambda\ne0$, $C_n^{(\lambda)}$ splits over $\mathbb{R}$,
is separable, and has $n$ roots counted with multiplicity, hence $n$ distinct
real zeros.

Reference: \cite[18.2(vi)]{dlmf}.
\end{theorem}
\begin{proof}
\leanok
\end{proof}

\begin{definition}[Normalised Gegenbauer classes]
\label{def:gegenbauer-normalised}\uses{thm:gegenbauer-sequence,thm:gegenbauer-moments,def:general-normalised-L2,thm:general-squared-norm}
\lean{Polynomial.normalizedGegenbauer}\leanok
For $\lambda>-1/2$, $\lambda\ne0$, let
$e_n=h_n^{-1/2}J_{\mu_\lambda}C_n^{(\lambda)}$, using the positive square root
of the squared norm.

Reference: \cite[Table 18.3.1]{dlmf}.
\end{definition}

\end{supporting}

\section{Legendre polynomials}\label{sec:textbook-classical-legendre}

Legendre polynomials are the Jacobi case $\alpha=\beta=0$, normalized by $P_n(1)=1$. Their orthogonality measure is ordinary Lebesgue measure on the interval.

\begin{definition}[Legendre polynomials]
\label{def:legendre}\uses{def:polynomial-linear-maps}
\lean{Polynomial.legendre}\leanok
Let $R$ be a commutative $\Q$-algebra and let $n\in\N$. Define (Rodrigues' formula)
\[P_n(X)=\frac{1}{2^n n!}\left(\frac{d}{dX}\right)^n(X^2-1)^n.\]

Reference: \cite[18.5.5, Table 18.5.1 and Table 18.3.1]{dlmf}.
\end{definition}

\begin{theorem}[Legendre--Jacobi bridge]
\label{thm:legendre-jacobi}\uses{thm:legendre-shifted,def:jacobi}
\lean{Polynomial.legendre_eq_jacobi}\leanok
Let $R$ be a commutative $\Q$-algebra and let $n\in\N$. Then $P_n=P_n^{(0,0)}$.

Reference: \cite[18.7.9]{dlmf}.
\end{theorem}
\begin{proof}
\leanok
\end{proof}

\begin{theorem}[Legendre endpoint values]
\label{thm:legendre-values}\uses{thm:legendre-jacobi,thm:jacobi-value-one,thm:jacobi-low-degree}
\lean{Polynomial.legendre_zero, Polynomial.legendre_one, Polynomial.eval_legendre_one, Polynomial.eval_legendre_neg_one}\leanok
Let $R$ be a commutative $\Q$-algebra and let $n\in\N$. The Legendre polynomials satisfy $P_0=1$, $P_1=X$, $P_n(1)=1$, and $P_n(-1)=(-1)^n$.

Reference: \cite[Table 18.6.1]{dlmf}.
\end{theorem}
\begin{proof}
\leanok
\end{proof}

\begin{theorem}[Legendre three-term recurrence]
\label{thm:legendre-recurrence}\uses{thm:legendre-recurrence-coefficients,thm:general-coefficient-signs}
\lean{Polynomial.legendre_recurrence}\leanok
Let $R$ be a commutative $\Q$-algebra and let $n\in\N$. Then\[(n+2)P_{n+2}=(2n+3)XP_{n+1}-(n+1)P_n.\]

Reference: \cite[18.9.1]{dlmf}.
\end{theorem}
\begin{proof}
\leanok
\end{proof}

\begin{theorem}[Legendre parity]
\label{thm:legendre-parity}\uses{def:legendre}
\lean{Polynomial.legendre_comp_neg_X, Polynomial.eval_legendre_neg}\leanok
Let $R$ be a commutative $\Q$-algebra and let $n\in\N$. Then $P_n(-X)=(-1)^nP_n(X)$.

Reference: \cite[Table 18.6.1]{dlmf}.
\end{theorem}
\begin{proof}
\leanok
\end{proof}

\begin{theorem}[Legendre degree and leading coefficient]
\label{thm:legendre-degree}\uses{thm:legendre-jacobi,thm:jacobi-degree}
\lean{Polynomial.degree_legendre_le, Polynomial.coeff_legendre_self, Polynomial.degree_legendre, Polynomial.natDegree_legendre, Polynomial.leadingCoeff_legendre, Polynomial.legendreSequence, Polynomial.legendreSequence_apply}\leanok
Let $R$ be a commutative $\Q$-algebra and let $n\in\N$.
The Legendre polynomial $P_n(X)\in R[X]$ has degree at most $n$, and its coefficient of $X^n$ is
$2^{-n}\binom{2n}{n}$.
If $R$ is nontrivial, this coefficient is nonzero, $P_n$ has degree exactly $n$, and its leading coefficient is $2^{-n}\binom{2n}{n}$.
For nontrivial $R$, the Legendre polynomials therefore form a degree-indexed polynomial sequence.

Reference: \cite[Table 18.3.1]{dlmf}.
\end{theorem}
\begin{proof}
\leanok
\end{proof}

\begin{theorem}[Legendre derivative]
\label{thm:legendre-derivative}\uses{thm:legendre-jacobi,thm:jacobi-derivative}
\lean{Polynomial.derivative_legendre_succ}\leanok
Let $R$ be a commutative $\Q$-algebra and let $n\in\N$. Then $P_{n+1}'=(n+2)P_n^{(1,1)}/2$.

Reference: \cite[18.9.15]{dlmf}.
\end{theorem}
\begin{proof}
\leanok
\end{proof}

\begin{theorem}[Legendre differential equation]
\label{thm:legendre-ode}\uses{thm:legendre-jacobi,thm:jacobi-ode}
\lean{Polynomial.legendre_differentialEquation}\leanok
Let $R$ be a commutative $\Q$-algebra and let $n\in\N$. The Legendre polynomials satisfy
\[(1-X^2)P_n''-2XP_n'+n(n+1)P_n=0\]


Reference: \cite[Table 18.8.1]{dlmf}.
\end{theorem}
\begin{proof}
\leanok
\end{proof}

\begin{definition}[Legendre measure]
\label{def:legendre-measure}
\lean{Polynomial.legendreMeasure}\leanok
Let $\mu_L$ be Lebesgue measure restricted to $(-1,1]$, the Legendre weight one.

Reference: \cite[Table 18.3.1]{dlmf}.
\end{definition}

\begin{theorem}[Legendre squared norms and integral pairing]
\label{thm:legendre-norm}\uses{thm:legendre-low-degree,thm:legendre-values,thm:legendre-derivative-leading-coefficient}
\lean{Polynomial.integral_sq_legendre, Polynomial.integral_sq_legendre_measure, Polynomial.integral_mul_legendre, Polynomial.integral_mul_legendre_measure}\leanok
For every $n\ge0$,
\[\int_{-1}^{1}P_n(x)^2\,dx=\frac{2}{2n+1}.\]
Consequently $\int P_mP_n\,d\mu_L=\delta_{mn}\,2/(2n+1)$, with the same
formula for the interval integral.

Reference: \cite[Table 18.3.1]{dlmf}.
\end{theorem}
\begin{proof}
\leanok
\end{proof}

\begin{theorem}[Legendre root location]
\label{thm:legendre-root-location}\uses{thm:legendre-low-degree}
\lean{Polynomial.root_legendre_mem_Ioo}\leanok
Every real zero of $P_n$ lies strictly in $(-1,1)$.

Reference: \cite[18.2(vi)]{dlmf}.
\end{theorem}
\begin{proof}
\leanok
\end{proof}

\begin{theorem}[Legendre strict interlacing]
\label{thm:legendre-interlacing}\uses{thm:general-outer-roots,thm:legendre-positive,thm:legendre-orthogonal}
\lean{Polynomial.not_eval_legendre_eq_zero_and_succ, Polynomial.exists_root_legendre_succ_between, Polynomial.exists_root_legendre_between_succ, Polynomial.exists_outer_roots_legendre}\leanok
Consecutive Legendre polynomials have no common zero. Between any two distinct zeros of either member lies a zero of the other. For $n\ge1$, $P_{n+1}$ has two zeros $a<b$ such that every zero of $P_n$ lies strictly between $a$ and $b$.

Reference: \cite[18.2(vi)]{dlmf}.
\end{theorem}
\begin{proof}
\leanok
\end{proof}

\begin{definition}[Legendre coefficients]
\label{def:legendre-coefficients}\uses{def:legendre-measure,def:legendre,def:integral-coefficient-input}
\lean{Polynomial.legendreCoefficient}\leanok
For fixed $n$, let $f:\R\to\R$ satisfy $fP_n\in L^1(\mu_L)$, with pointwise values retained. Define $a_n(f)=\frac{2n+1}{2}\int f(x)P_n(x)\,d\mu_L(x)$.

Reference: \cite[18.2(viii)]{dlmf}.
\end{definition}

The map $J_{\mu_L}$ sends a polynomial to its $L^2(\mu_L)$ class. Convergence below is in this norm.

\begin{theorem}[Legendre L2 expansion]
\label{thm:legendre-L2-expansion}\uses{thm:legendre-basis,thm:legendre-coefficient-characterization}
\lean{Polynomial.hilbertBasisLegendre_repr_smul, Polynomial.hasSum_legendreCoefficient_polynomialToL2, Polynomial.tendsto_sum_legendreCoefficient_polynomialToL2}\leanok
For every $f\in L^2(\mu_L)$ the Hilbert-coordinate summands equal $a_n(f)J_{\mu_L}P_n$, their series sums to $f$, and $\sum_{n=0}^{N-1}a_n(f)J_{\mu_L}P_n\to f$ in $L^2$ as $N\to\infty$.

Reference: \cite[18.2(viii)]{dlmf}.
\end{theorem}
\begin{proof}
\leanok
\end{proof}

\begin{supporting}[Coefficient maps, shifted polynomials and Legendre structure]

\begin{theorem}[Legendre polynomials under coefficient maps]
\label{thm:legendre-maps}\uses{def:legendre}
\lean{Polynomial.map_legendre, Polynomial.map_legendre_rat}\leanok
Let $R,S$ be commutative $\Q$-algebras, let $n\in\N$, and let $f:R\to S$ be a $\Q$-algebra homomorphism. Applying $f$ to the
coefficients of $P_n$ over $R$ gives $P_n$ over $S$. In particular, $P_n$ over
$R$ is the image of the rational polynomial $P_n$.

Reference: \cite[18.5.5]{dlmf}.
\end{theorem}
\begin{proof}
\leanok
\end{proof}

\begin{theorem}[Legendre hypergeometric representation]
\label{thm:legendre-hypergeometric}\uses{thm:legendre-jacobi, thm:jacobi-representations}
\lean{Polynomial.eval_legendre_eq_ordinaryHypergeometric}\leanok
For real or complex $x$ and $n\in\N$,
\[P_n(x)={}_2F_1\!\left(-n,n+1;1;\frac{1-x}2\right),\]
where ${}_2F_1$ is the Gauss series of Definition~\ref{def:ordinary-hypergeometric}. No hypothesis is needed.

Reference: \cite[18.5.7]{dlmf}, \cite[Table 18.3.1]{dlmf}.
\end{theorem}
\begin{proof}
\leanok
\end{proof}

\begin{theorem}[Legendre and the reflected shifted family]
\label{thm:legendre-shifted}\uses{thm:shifted-legendre,thm:legendre-parity,thm:legendre-maps}
\lean{Polynomial.legendre_comp_one_sub_two_mul_X, Polynomial.legendre_eq_shiftedLegendre_comp, Polynomial.legendre_comp_two_mul_X_sub_one}\leanok
Let $R$ be a commutative $\Q$-algebra and let $n\in\N$. After mapping the integer coefficients of the shifted Legendre polynomial $S_n$ into $R$,
\[P_n(1-2X)=S_n(X),\qquad P_n(X)=S_n((1-X)/2).\]
Equivalently, $P_n(2X-1)=(-1)^nS_n(X)$. The parity factor arises because $S_n$ is
oriented by $S_1(X)=1-2X$.

Reference: \cite[18.7.10 and Table 18.6.1]{dlmf}.
\end{theorem}
\begin{proof}
\leanok
\end{proof}

\begin{theorem}[Leading coefficient after Legendre differentiation]
\label{thm:legendre-derivative-leading-coefficient}\uses{thm:legendre-degree}
\lean{Polynomial.coeff_derivative_legendre_succ}\leanok
Let $R$ be a commutative $\Q$-algebra and let $n\in\N$. Then
\[ [X^n]P_{n+1}'=(2n+1)[X^n]P_n. \]

Reference: \cite[18.3 and 18.9]{dlmf}.
\end{theorem}
\begin{proof}
\leanok
\end{proof}

\begin{theorem}[Shifted Legendre family]
\label{thm:shifted-legendre}
\lean{Polynomial.shiftedLegendre, Polynomial.coeff_shiftedLegendre, Polynomial.degree_shiftedLegendre, Polynomial.factorial_mul_shiftedLegendre_eq, Polynomial.neg_one_pow_mul_shiftedLegendre_comp_one_sub_X_eq}\mathlibok
Define the integer polynomial
\[S_n(X)=\sum_{k=0}^n(-1)^k\binom nk\binom{n+k}{n}X^k.\]
It has degree $n$, satisfies $n!S_n=(d/dX)^n[X^n(1-X)^n]$, and
$(-1)^nS_n(1-X)=S_n(X)$. In particular $S_1(X)=1-2X$: its orientation is
opposite to the frequently used shifted convention $P_n(2X-1)$.

Reference: \cite[18.7.10]{dlmf}, with the reflected orientation $x\mapsto1-x$.
\end{theorem}
\begin{proof}
\mathlibok
\end{proof}

\begin{theorem}[Finite Legendre moments]
\label{thm:legendre-moments}\uses{def:legendre-measure,def:polynomial-moments}
\lean{Polynomial.hasFiniteMoments_legendreMeasure}\leanok
Every real polynomial has integrable square with respect to $\mu_L$.

Reference: \cite[Table 18.3.1]{dlmf}.
\end{theorem}
\begin{proof}
\leanok
\end{proof}

\begin{theorem}[Legendre orthogonality]
\label{thm:legendre-orthogonal}\uses{thm:classical-jacobi-measures,thm:legendre-moments,thm:legendre-degree,thm:jacobi-pairwise}
\lean{Polynomial.integral_mul_legendre_eq_zero, Polynomial.isOrthogonal_legendreSequence}\leanok
For distinct $m,n\ge0$,
\[\int_{-1}^{1}P_m(x)P_n(x)\,dx=0.\]
The Legendre polynomial sequence is orthogonal for $\mu_L$ in the general
\texttt{IsOrthogonal} API, including integrability of every pairwise product.

Reference: \cite[Table 18.3.1]{dlmf}.
\end{theorem}
\begin{proof}
\leanok
\end{proof}

\begin{theorem}[Legendre orthogonality to lower degrees]
\label{thm:legendre-low-degree}\uses{thm:legendre-orthogonal,thm:orthogonal-lower-degree}
\lean{Polynomial.integral_mul_legendre_eq_zero_of_degree_lt}\leanok
For every real polynomial $q$ of degree less than $n$,
$\int_{-1}^{1}q(x)P_n(x)\,dx=0$.

Reference: \cite[18.2(i)]{dlmf}.
\end{theorem}
\begin{proof}
\leanok
\end{proof}

\begin{theorem}[Positive Legendre measure]
\label{thm:legendre-positive}\uses{thm:legendre-moments,thm:polynomial-inner}
\lean{Polynomial.hasPositiveDefiniteIntegral_legendreMeasure, Polynomial.isFiniteMeasure_legendreMeasure}\leanok
The measure $\mu_L$ is finite and $\int p^2\,d\mu_L>0$ for every nonzero real polynomial $p$.
\end{theorem}
\begin{proof}
\leanok
\end{proof}

\begin{theorem}[Legendre recurrence coefficients]
\label{thm:legendre-recurrence-coefficients}\uses{thm:legendre-norm,thm:legendre-positive,def:general-recurrence-coefficients,thm:general-squared-norm}
\lean{Polynomial.squaredNorm_legendreSequence, Polynomial.recurrenceA_legendreSequence, Polynomial.recurrenceB_legendreSequence, Polynomial.recurrenceC_legendreSequence}\leanok
In the general recurrence convention, $h_n=2/(2n+1)$, $A_n=(n+1)/(2n+1)$, $B_n=0$, and $C_n=(n+1)/(2n+3)$.

Reference: \cite[18.9.1 and Table 18.9.1]{dlmf}.
\end{theorem}
\begin{proof}
\leanok
\end{proof}

\begin{theorem}[Legendre Christoffel--Darboux identities]
\label{thm:legendre-cd}\uses{thm:legendre-recurrence-coefficients,thm:general-cd-diag}
\lean{Polynomial.legendre_christoffelDarboux, Polynomial.legendre_christoffelDarboux_diag}\leanok
Set $K_n(x,y)=\sum_{k=0}^n\frac{2k+1}{2}P_k(x)P_k(y)$. For real $x,y$,
\[(x-y)K_n(x,y)=\frac{n+1}{2}\bigl(P_{n+1}(x)P_n(y)-P_n(x)P_{n+1}(y)\bigr),\]
\[K_n(x,x)=\frac{n+1}{2}\bigl(P_{n+1}'(x)P_n(x)-P_n'(x)P_{n+1}(x)\bigr).\]

Reference: \cite[18.2.12--13]{dlmf}.
\end{theorem}
\begin{proof}
\leanok
\end{proof}

\begin{theorem}[Real simple Legendre zeros]
\label{thm:legendre-real-zeros}\uses{thm:legendre-positive,thm:legendre-orthogonal,thm:general-real-zeros}
\lean{Polynomial.splits_legendre, Polynomial.card_roots_legendre, Polynomial.separable_legendre}\leanok
$P_n$ splits over $\R$, is separable, and has exactly $n$ roots counted with multiplicity, hence $n$ distinct real zeros.

Reference: \cite[18.2(vi)]{dlmf}.
\end{theorem}
\begin{proof}
\leanok
\end{proof}

\begin{definition}[Normalised Legendre classes]
\label{def:legendre-normalised}\uses{def:general-polynomial-L2,thm:legendre-moments,def:legendre}
\lean{Polynomial.normalizedLegendre}\leanok
Define $e_n=\sqrt{(2n+1)/2}\,J_{\mu_L}P_n$.

Reference: \cite[Table 18.3.1]{dlmf}.
\end{definition}

\begin{theorem}[Legendre Hilbert basis]
\label{thm:legendre-basis}\uses{def:legendre-normalised,thm:compact-polynomial-dense,thm:legendre-norm}
\lean{Polynomial.ae_mem_Icc_legendreMeasure, Polynomial.polynomialToL2_legendre_denseRange, Polynomial.inner_polynomialToL2_legendre, Polynomial.orthonormal_normalizedLegendre, Polynomial.span_polynomialToL2_legendre, Polynomial.normalizedLegendre_dense_span, Polynomial.hilbertBasisLegendre, Polynomial.hilbertBasisLegendre_apply}\leanok
The family $(e_n)_{n\ge0}$ is orthonormal with dense span in real $L^2(\mu_L)$ and defines a Hilbert basis.

Reference: \cite[18.2(viii)]{dlmf}.
\end{theorem}
\begin{proof}
\leanok
\end{proof}

\begin{theorem}[Legendre coefficient characterization]
\label{thm:legendre-coefficient-characterization}\uses{def:legendre-coefficients,thm:legendre-positive,thm:legendre-degree}
\lean{Polynomial.IntegralCoefficientInput.ofLegendreL2, Polynomial.legendreCoefficient_apply}\leanok
For every $f\in L^2(\mu_L;\mathbb R)$ and $n\in\N$, the product $fP_n$ is
$\mu_L$-integrable. More generally, for every real-valued function $f:\mathbb R\to\mathbb R$ with
$fP_n$ integrable, its Legendre coefficient is
\[
 a_n(f)=\frac{2n+1}{2}\int f(x)P_n(x)\,d\mu_L(x).
\]
\end{theorem}
\begin{proof}
\leanok
\end{proof}

\end{supporting}

\section{Chebyshev as Jacobi cases}\label{sec:textbook-classical-chebyshev}

The first- and second-kind Chebyshev polynomials correspond to Jacobi parameters $-1/2$ and $1/2$. Their polynomial normalizations differ by explicit factors, while their measures agree exactly.

\begin{theorem}[Chebyshev differential equations and endpoint values]
\label{thm:chebyshev-ode}
\lean{Polynomial.Chebyshev.one_sub_X_sq_mul_derivative_derivative_T_eq_poly_in_T, Polynomial.Chebyshev.one_sub_X_sq_mul_derivative_derivative_U_eq_poly_in_U, Polynomial.Chebyshev.T_eval_one, Polynomial.Chebyshev.U_eval_one}\mathlibok
The polynomials $T_n$ and $U_n$ satisfy
\[(1-X^2)T_n''-XT_n'+n^2T_n=0,\quad T_n(1)=1,\]
\[(1-X^2)U_n''-3XU_n'+n(n+2)U_n=0,\quad U_n(1)=n+1.\]
These hold over every commutative ring and for every integer $n$.

Reference: \cite[Tables 18.8.1 and 18.6.1]{dlmf}.
\end{theorem}
\begin{proof}
\mathlibok
\end{proof}

\begin{theorem}[Chebyshev--Jacobi bridges]
\label{thm:chebyshev-jacobi}\uses{thm:chebyshev-ode,thm:jacobi-ode-unique,def:chebyshev-mathlib,thm:jacobi-ode}
\lean{Polynomial.jacobi_eq_chebyshev_T, Polynomial.jacobi_eq_chebyshev_U}\leanok
Let $R$ be a commutative $\Q$-algebra and let $n\in\N$. Then
\[P_n^{(-1/2,-1/2)}=\frac{(1/2)_n}{n!}T_n,\qquad
(n+1)P_n^{(1/2,1/2)}=\frac{(3/2)_n}{n!}U_n.\]

Reference: \cite[18.7.3--4]{dlmf}.
\end{theorem}
\begin{proof}
\leanok
\end{proof}

\begin{theorem}[Classical Jacobi measure bridges]
\label{thm:classical-jacobi-measures}\uses{def:legendre-measure,thm:classical-chebyshev-weights,thm:jacobi-weight-characterisation,def:measureT,def:measureU}
\lean{Polynomial.legendreMeasure_eq_jacobiMeasure, Polynomial.measureT_eq_jacobiMeasure, Polynomial.measureU_eq_jacobiMeasure}\leanok
The existing measures satisfy
\[
 \mu_L=\mu_{0,0},\qquad \mu_T=\mu_{-1/2,-1/2},\qquad
 \mu_U=\mu_{1/2,1/2}.
\]
These are exact measure equalities, with no scalar normalization factor.
Endpoints have zero mass, and restricting the measures to the open or half-open interval gives the same measure.

Reference: \cite[Table 18.3.1]{dlmf}.
\end{theorem}
\begin{proof}
\leanok
\end{proof}

\begin{supporting}[Hypergeometric representations and norm compatibility]

\begin{theorem}[Chebyshev hypergeometric representations]
\label{thm:chebyshev-hypergeometric}\uses{thm:chebyshev-jacobi,thm:jacobi-representations}
\lean{Polynomial.eval_chebyshev_T_eq_ordinaryHypergeometric, Polynomial.eval_chebyshev_U_eq_ordinaryHypergeometric}\leanok
For real or complex $x$ and $n\in\N$,
\[T_n(x)={}_2F_1\!\left(-n,n;\tfrac12;\frac{1-x}2\right),\qquad
 U_n(x)=(n+1)\,{}_2F_1\!\left(-n,n+2;\tfrac32;\frac{1-x}2\right),\]
where ${}_2F_1$ is the Gauss series of Definition~\ref{def:ordinary-hypergeometric}. No hypothesis is needed.

Reference: \cite[18.5.11\_2 and 18.5.11\_4]{dlmf}.
\end{theorem}
\begin{proof}
\leanok
\end{proof}

\begin{theorem}[Classical norm compatibility]
\label{thm:classical-jacobi-norms}\uses{thm:legendre-norm,thm:U-norm}
\lean{Polynomial.integral_sq_jacobi_zero_zero, Polynomial.integral_sq_chebyshev_T_jacobiMeasure, Polynomial.integral_sq_chebyshev_U_jacobiMeasure}\leanok
In the corresponding Jacobi measures, the squared norms of $P_n^{(0,0)}$, $T_n$
and $U_n$ are respectively $2/(2n+1)$, $\pi$ for $n=0$ and $\pi/2$ otherwise,
and $\pi/2$.

Reference: \cite[Table 18.3.1]{dlmf}.
\end{theorem}
\begin{proof}
\leanok
\end{proof}

\end{supporting}

\chapter{Connection coefficients}\label{chap:connections}

\rvsummary{Changing Jacobi parameters without changing the represented polynomial: the connection formulas behind basis conversion.}

Changing parameters changes the basis, not the polynomial being represented. The connection coefficients give the corresponding finite change of coordinates.

\section{One-step parameter changes}\label{sec:textbook-connections-raising}

A connection formula rewrites the same polynomial in a different basis. Raising one Jacobi parameter by one gives a two-term relation.

\begin{theorem}[Jacobi parameter raising]
\label{thm:jacobi-parameter-raising}\uses{thm:jacobi-values}
\lean{Polynomial.jacobi_raise_beta, Polynomial.jacobi_raise_alpha}\leanok
For $n\ge0$, over every commutative $\Q$-algebra, write $s=\alpha+\beta$. Then
\begin{align*}
(2n+s+3)P_{n+1}^{(\alpha,\beta)}
 &= (n+s+2)P_{n+1}^{(\alpha,\beta+1)}+(n+\alpha+1)P_n^{(\alpha,\beta+1)},\\
(2n+s+3)P_{n+1}^{(\alpha,\beta)}
 &= (n+s+2)P_{n+1}^{(\alpha+1,\beta)}-(n+\beta+1)P_n^{(\alpha+1,\beta)}.
\end{align*}
No parameter factor is cancelled. The constant members are all one.

Reference: \cite[18.9.5]{dlmf} and its reflection via Table 18.6.1.
\end{theorem}
\begin{proof}
\leanok
\end{proof}

\section{Finite connection formulas}\label{sec:textbook-connections-connections}

The quotient $A_{nk}/d_{nk}$ gives a one-parameter connection coefficient in the orthogonality range. The cleared expressions $B_{nk}$ and $D_n$ also describe changes at degenerate parameters, without division.

\begin{definition}[Explicit Jacobi connection data]
\label{def:jacobi-connection-data}
\lean{Polynomial.jacobiAlphaConnectionNumerator, Polynomial.jacobiAlphaConnectionDenominator, Polynomial.jacobiAlphaConnectionScale, Polynomial.jacobiAlphaConnectionCleared}\leanok
For $0\le k\le n$ define polynomial expressions
\begin{align*}
 A_{nk}(\alpha,\beta,\gamma)
   &=(\beta+k+1)_{n-k}(n+\gamma+\beta+1)_k(\gamma-\alpha)_{n-k},\\
 d_{nk}(\alpha,\beta)
   &=(k+\alpha+\beta+1)_k(\alpha+\beta+2k+2)_{n-k}(n-k)!,\\
 D_n(\alpha,\beta)&=\prod_{k=0}^n d_{nk}(\alpha,\beta),\\
 B_{nk}(\alpha,\beta,\gamma)
   &=A_{nk}(\alpha,\beta,\gamma)\prod_{0\le j\le n,\ j\ne k}d_{nj}(\alpha,\beta).
\end{align*}
These are defined over every commutative ring.

Reference: \cite[18.18.14]{dlmf}.
\end{definition}

\begin{theorem}[One-parameter Jacobi connections in the orthogonality range]
\label{thm:jacobi-alpha-connection}\uses{thm:jacobi-mixed-moments,def:jacobi-connection-data,thm:jacobi-finite-connections,def:jacobi-coefficients,thm:jacobi-pairwise,thm:general-coordinates,thm:integral-coefficient-input}
\lean{Polynomial.ascPochhammer_jacobi_leading_pos, Polynomial.jacobiBasis_repr_eq_coefficient, Polynomial.jacobiCoefficient_jacobi_add, Polynomial.jacobiAlphaConnectionDenominator_pos, Polynomial.jacobiFiniteConnectionCoefficient_eq_alpha, Polynomial.jacobi_eq_sum_alphaConnection, Polynomial.jacobiAlphaConnectionScale_smul_of_neg_one_lt}\leanok
For $\alpha,\beta>-1$, every $d_{nk}(\alpha,\beta)>0$, and for arbitrary real $\gamma$,
\[ P_n^{(\gamma,\beta)}=
 \sum_{k=0}^n \frac{A_{nk}(\alpha,\beta,\gamma)}{d_{nk}(\alpha,\beta)}P_k^{(\alpha,\beta)}. \]
The coefficients equal both the integral projection coefficients and the finite
determinant coefficients. Degree zero and $\alpha+\beta+1=0$ are included.

Reference: \cite[18.18.14]{dlmf}.
\end{theorem}
\begin{proof}
\leanok
\end{proof}

\begin{theorem}[General Jacobi connection coefficients]
\label{thm:jacobi-connections}\uses{thm:jacobi-universal-connection}\leanok
\lean{Polynomial.jacobi_generalConnection_cleared, Polynomial.jacobi_eq_sum_generalConnection, Polynomial.jacobi_generalConnection_scale_pos}
Let
\[ L_n=D_n(\alpha,\delta)\prod_{i=0}^nD_i(\beta,\alpha),\qquad
 Q_{ni}=\prod_{\substack{0\le l\le n\\l\ne i}}D_l(\beta,\alpha). \]
Over every commutative rational algebra, for arbitrary source and target parameters,
\[ L_nP_n^{(\gamma,\delta)}=
 \sum_{i=0}^n\sum_{j=0}^i
 B_{ni}(\alpha,\delta,\gamma)Q_{ni}(-1)^{i+j}
 B_{ij}(\beta,\alpha,\delta)P_j^{(\alpha,\beta)}. \]
Over a field with $L_n\ne0$, divide by $L_n$ to obtain the explicit finite
connection expansion. Thus the degree-$j$ coefficient is the displayed finite
convolution over $i\ge j$. The common scale is positive over $\R$ when
$\alpha,\beta,\delta>-1$, with unrestricted $\gamma$; the universal cleared
identity also includes all degenerate parameter choices.

Reference: \cite[18.9(ii) and 18.18(iv)]{dlmf}.
\end{theorem}
\begin{proof}
\leanok
\end{proof}

\begin{supporting}[Universal identities, determinant coefficients and mixed moments]

\begin{theorem}[Connection data under ring homomorphisms]
\label{thm:jacobi-connection-data-maps}\uses{def:jacobi-connection-data}
\lean{Polynomial.map_jacobiAlphaConnectionNumerator, Polynomial.map_jacobiAlphaConnectionDenominator, Polynomial.map_jacobiAlphaConnectionScale, Polynomial.map_jacobiAlphaConnectionCleared}\leanok
For every ring homomorphism $f:R\to S$, applying $f$ to $A_{nk}$, $d_{nk}$,
$D_n$ and $B_{nk}$ gives the same expressions with parameters
$f(\alpha),f(\beta),f(\gamma)$.
\end{theorem}
\begin{proof}
\leanok
\end{proof}

\begin{theorem}[Universal one-parameter Jacobi connections]
\label{thm:jacobi-universal-connection}\uses{thm:jacobi-alpha-connection,thm:jacobi-values,thm:jacobi-connection-data-maps}
\lean{Polynomial.eq_of_map_eval₂_eq_on, Polynomial.jacobiAlphaConnectionScale_smul, Polynomial.jacobiBetaConnectionScale_smul}\leanok
For arbitrary parameters in any commutative rational algebra,
\begin{align*}
 D_n(\alpha,\beta)P_n^{(\gamma,\beta)}
   &=\sum_{k=0}^n B_{nk}(\alpha,\beta,\gamma)P_k^{(\alpha,\beta)},\\
 D_n(\beta,\alpha)P_n^{(\alpha,\delta)}
   &=\sum_{k=0}^n(-1)^{n+k}B_{nk}(\beta,\alpha,\delta)P_k^{(\alpha,\beta)}.
\end{align*}
No nondegeneracy assumption is needed for either cleared identity.

Reference: \cite[18.18.14 and Table 18.6.1]{dlmf}.
\end{theorem}
\begin{proof}
\leanok
\end{proof}

\begin{definition}[Finite determinant connection coefficients]
\label{def:finite-connections}\uses{def:jacobi}
\lean{Polynomial.connectionMatrix, Polynomial.connectionNumerator, Polynomial.FiniteConnectionInput, Polynomial.finiteConnectionCoefficient, Polynomial.jacobiConnectionMatrix, Polynomial.jacobiConnectionNumerator, Polynomial.jacobiFiniteConnectionCoefficient}\leanok
For a finite family $p_0,\ldots,p_{N-1}$, let $A_{ij}=[X^i]p_j$.
For $q$ let $D_j(q)$ be the determinant obtained by replacing column $j$ of $A$
with the first $N$ coefficients of $q$. Over a field, on the domain $\det A\ne0$, put $c_j(q)=D_j(q)/\det A$.
The Jacobi versions use target parameters $(\alpha,\beta)$, source parameters
$(\gamma,\delta)$, and target degrees $0,\ldots,n$ for a source of degree $n$.

\end{definition}

\begin{theorem}[Characterisation of divided finite connection coefficients]
\label{thm:finite-connection-characterisation}\uses{def:finite-connections}
\lean{Polynomial.finiteConnectionCoefficient_mul_det, Polynomial.jacobiFiniteConnectionCoefficient_mul_det}\leanok
For every nonsingular target family, multiplying each divided connection
coefficient by the target determinant recovers its Cramer numerator.
The same identity holds for the Jacobi specialization, with unrestricted
source parameters and a nonsingular target matrix.
\end{theorem}
\begin{proof}
\leanok
\end{proof}

\begin{theorem}[Universal finite connection identity]
\label{thm:finite-connections}\uses{def:finite-connections}
\lean{Polynomial.connectionNumerator_eq_det, Polynomial.det_connectionMatrix, Polynomial.det_smul_eq_sum_connectionNumerator, Polynomial.eq_sum_finiteConnectionCoefficient}\leanok
Over any commutative ring, if $\deg p_j<N$ and $\deg q<N$, then
\[ (\det A)q=\sum_{j<N}D_j(q)p_j. \]
If $\deg p_j\le j$, then $\det A=\prod_{j<N}[X^j]p_j$.
Over a field, $\det A\ne0$ gives $q=\sum_{j<N}c_j(q)p_j$.

Reference: Cramer's rule.
\end{theorem}
\begin{proof}
\leanok
\end{proof}

\begin{theorem}[Finite Jacobi connections throughout the orthogonality range]
\label{thm:jacobi-finite-connections}\uses{thm:finite-connections,thm:coordinate-basis-vectors}
\lean{Polynomial.det_jacobiConnectionMatrix, Polynomial.det_smul_jacobi_eq_sum_connectionNumerator, Polynomial.jacobi_eq_sum_finiteConnectionCoefficient, Polynomial.det_jacobiConnectionMatrix_pos, Polynomial.jacobi_eq_sum_finiteConnectionCoefficient_of_neg_one_lt, Polynomial.jacobiFiniteConnectionCoefficient_eq_repr}\leanok
The cleared determinant identity applies to Jacobi families over every
commutative rational algebra, without restrictions on any of the four parameters.
Its determinant is the product of the degree-$j$ coefficients $[X^j]P_j$
of the target polynomials. These are their leading coefficients when the degrees
do not drop. Over a field,
a nonzero target determinant gives the divided expansion.
Over $\R$, target parameters $\alpha,\beta>-1$ make the determinant strictly positive;
source parameters $\gamma,\delta$ remain unrestricted. The divided coefficients
agree with the existing target polynomial-basis coordinates.
\end{theorem}
\begin{proof}
\leanok
\end{proof}

\begin{theorem}[Mixed Jacobi moments and weight shifts]
\label{thm:jacobi-mixed-moments}\uses{thm:jacobi-rodrigues,thm:gamma-polynomial-growth,thm:jacobi-weight-closed,thm:jacobi-mass}
\lean{Polynomial.hasDerivAt_jacobiWeight_mul_jacobi_mixed, Polynomial.integral_jacobi_mixed_succ, Polynomial.integral_jacobi_mixed, Polynomial.integral_mul_jacobi_succ, Polynomial.integral_jacobi_mul_jacobi_shift, Polynomial.integral_one_jacobiMeasure_beta_add_nat}\leanok
Write $m(\alpha,\beta)=\int1\,d\mu_{\alpha,\beta}$.
For $\alpha,\beta>-1$ and arbitrary $\gamma\in\R$,
\begin{align*}
 2^n n!\int P_n^{(\gamma,\beta)}\,d\mu_{\alpha,\beta}
   &=(\gamma-\alpha)_n m(\alpha,\beta+n),\\
 (\alpha+\beta+2)_m m(\alpha,\beta+m)
   &=2^m(\beta+1)_m m(\alpha,\beta),\\
 4^k k!\int P_{k+m}^{(\gamma,\beta)}P_k^{(\alpha,\beta)}\,d\mu_{\alpha,\beta}
   &=(k+m+\gamma+\beta+1)_k
     \int P_m^{(\gamma+k,\beta+k)}\,d\mu_{\alpha+k,\beta+k}.
\end{align*}
The underlying integration-by-parts identity holds against any polynomial.

Reference: \cite[18.18.14]{dlmf}.
\end{theorem}
\begin{proof}
\leanok
\end{proof}

\end{supporting}

\chapter{Approximation in Jacobi and ultraspherical series}\label{chap:jacobi-approx}

\rvsummary{Weighted $L^2$ and uniform geometric convergence of Jacobi and ultraspherical expansions.}

Chebyshev approximation bounds also control projections in other weighted polynomial bases. Kernel estimates then turn weighted norm bounds into uniform bounds for Jacobi and Gegenbauer expansions.

\section{Partial sums and weighted best approximation}\label{sec:textbook-jacobi-approx-projection}

The finite coefficient domain $\mathcal A_{\mu,P,\le N}^{\R}$ consists of real functions for which every product $fP_n$ and square $P_n^2$ is integrable and $h_n>0$, for $n\le N$ (Definition~\ref{def:integral-coefficient-family-input}). The partial sum is the best polynomial approximation in the weighted $L^2$ norm.

\begin{definition}[Polynomial partial sums and kernel vectors]
\label{def:weighted-partial-sum}\uses{def:general-normalised-L2,def:integral-coefficient-family-input}
\lean{Polynomial.Sequence.partialSum, Polynomial.Sequence.kernelVector}\leanok
A kernel cutoff is a natural number $N$ together with $h_n>0$ for every $0\le n\le N$. For $f\in\mathcal A_{\mu,P,\le N}^{\R}$, write $S_Nf=\sum_{n=0}^Na_n(f)P_n$ as a polynomial, and let
$v_{N,x}=\sum_{n=0}^NP_n(x)h_n^{-1/2}q_n$ be the kernel vector in weighted L2,
where $q_n=P_n/\sqrt{h_n}$ and the kernel vector is defined only at a kernel cutoff.
\end{definition}

\begin{theorem}[Weighted polynomial best approximation]
\label{thm:weighted-best-approximation}\uses{thm:finite-best-approximation,thm:general-hilbert-basis}
\lean{Polynomial.Sequence.inner_normalizedToL2_smul, Polynomial.Sequence.polynomialToL2_mem_span_normalizedToL2, Polynomial.Sequence.IsOrthogonal.norm_sub_sum_integralCoefficient_le, Polynomial.Sequence.IsOrthogonal.norm_sub_sum_integralCoefficient_le_of_ae_bound}\leanok
Let $P_n$ be orthogonal for $\mu$, with finite polynomial moments and positive
definite polynomial integral. For $f\in L^2(\mu;\R)$, let
$S_Nf=\sum_{n=0}^Na_n(f)P_n$ use the integral coefficients in the original
normalisation. Every real polynomial $p$ of degree at most $N$ satisfies
\[ \|f-S_Nf\|_{L^2(\mu)}\le\|f-p\|_{L^2(\mu)}. \]
If $|f-p|\le\varepsilon$ almost everywhere, with $\varepsilon\ge0$, the error is
at most $\sqrt{\mu(\R)}\,\varepsilon$. No density or support condition is assumed.

Reference: \cite[18.2(viii)]{dlmf}.
\end{theorem}
\begin{proof}
\leanok
\end{proof}

\begin{theorem}[Geometric weighted L2 approximation]
\label{thm:weighted-analytic-approximation}\uses{thm:weighted-best-approximation,thm:real-analytic-competitor}
\lean{Polynomial.Sequence.IsOrthogonal.norm_sub_sum_integralCoefficient_le_of_analytic}\leanok
Under the weighted best-approximation hypotheses, assume additionally that
$\mu$ is supported almost everywhere on $[-1,1]$. For $\rho>1$, let $F$ be
analytic and bounded by $M$ in $E_\rho$, and let $f\in L^2(\mu;\R)$ agree almost
everywhere with $\operatorname{Re}F$. Then
\[ \|f-S_Nf\|_{L^2(\mu)}\le
 \sqrt{\mu(\R)}\,\frac{2M}{\rho^N(\rho-1)}. \]
The support may be disconnected.

Reference: the competitor of \cite[Theorem 8.2]{atap}.
\end{theorem}
\begin{proof}
\leanok
\end{proof}

\begin{theorem}[Geometric Jacobi L2 approximation]
\label{thm:jacobi-analytic-approximation}\uses{thm:weighted-analytic-approximation,thm:jacobi-pairwise,def:jacobi-coefficients}
\lean{Polynomial.norm_sub_sum_jacobiCoefficient_le, Polynomial.norm_sub_sum_jacobiCoefficient_le_of_analytic}\leanok
Let $\alpha,\beta>-1$, $f\in L^2(\mu_{\alpha,\beta};\R)$, and
$S_Nf=\sum_{n=0}^N a_n(f)P_n^{(\alpha,\beta)}$.
For every real polynomial $p$ of degree at most $N$,
\[
 \|f-S_Nf\|_{L^2(\mu_{\alpha,\beta})}\le
 \|f-p\|_{L^2(\mu_{\alpha,\beta})}.
\]
If $\rho>1$, $F$ is analytic on $E_\rho$, $|F|\le M$ there, and
$f=\operatorname{Re}F$ almost everywhere for $\mu_{\alpha,\beta}$, then
\[
 \|f-S_Nf\|_{L^2(\mu_{\alpha,\beta})}\le
 \sqrt{\mu_{\alpha,\beta}(\R)}\frac{2M}{\rho^N(\rho-1)}.
\]
The original Jacobi measure is used, including singular endpoint weights.
\end{theorem}
\begin{proof}
\leanok
\end{proof}

\paragraph{Example (Julia).} For $F(z)=1/(1+25z^2)$ choose $\rho=1.2$. Put $b=(\rho-\rho^{-1})/2<1/5$. On $E_\rho$, $\Re(1+25z^2)\ge1-25b^2>0$, so $M=(1-25b^2)^{-1}$ is an upper bound for $|F|$. The Gauss--Legendre rule of Chapter~\ref{chap:gauss} estimates the coefficients and weighted errors below; these computed norms are numerical estimates, not certified integrals.

\begin{juliacode}
using LinearAlgebra
n = 200                             # Gauss-Legendre nodes
b = [k / sqrt(4k^2 - 1) for k in 1:n-1]
E = eigen(SymTridiagonal(zeros(n), b))
x, w = E.values, 2 .* E.vectors[1, :].^2
P = ones(n, 62); P[:, 2] = x        # Legendre P_k at the nodes, k <= 61
for k in 1:60
    P[:, k+2] = ((2k+1) .* x .* P[:, k+1] .- k .* P[:, k]) ./ (k+1)
end
f = 1 ./ (1 .+ 25x.^2)
a = [(2k+1)/2 * dot(w, f .* P[:, k+1]) for k in 0:61]
L2err(N) = sqrt(dot(w, (f .- P[:, 1:N+1]*a[1:N+1]).^2))
rho = 1.2                           # poles lie on the ellipse rho = 1.2198...
bminor = (rho - 1/rho) / 2
M = 1 / (1 - 25bminor^2)             # analytic upper bound on the ellipse
[L2err(N) / (sqrt(2)*2M/(rho^N*(rho - 1))) for N in 10:10:60]
                                    # compare estimated errors with the analytic bound
\end{juliacode}

\begin{supporting}[Partial-sum and kernel-vector formulas]

\begin{theorem}[Classical finite kernel-vector formula]
\label{thm:kernel-vector-characterisation}\uses{def:weighted-partial-sum}
\lean{Polynomial.Sequence.kernelVector_apply}\leanok
For a kernel cutoff $N$ and $x\in\R$,
$v_{N,x}=\sum_{n=0}^N P_n(x)h_n^{-1/2}e_n$.
Every summand uses its positive-norm index.
\end{theorem}
\begin{proof}
\leanok
\end{proof}

\begin{theorem}[Classical finite partial sum formula]
\label{thm:weighted-partial-sum-characterisation}\uses{def:weighted-partial-sum}
\lean{Polynomial.Sequence.partialSum_apply}\leanok
For $f\in\mathcal A_{\mu,P,\le N}^{\R}$,
\[
 S_Nf=\sum_{n=0}^N\frac{\int fP_n\,d\mu}{h_n}P_n,
 \qquad h_n=\int P_n^2\,d\mu>0.
\]
Each term uses the same function $f$ in its index-$n$ coefficient domain.
\end{theorem}
\begin{proof}
\leanok
\end{proof}

\end{supporting}

\section{From weighted error to uniform error}\label{sec:textbook-jacobi-approx-uniform}

A weighted norm does not determine a point value. The Christoffel kernel bounds evaluation of a finite polynomial projection and supplies the extra factor needed for a uniform estimate.

The hypotheses here are finite polynomial moments, a positive-definite polynomial pairing and orthogonality of $P$, as in the weighted best-approximation theorem. The kernel is $K_N(x,y)=\sum_{n=0}^N P_n(x)P_n(y)/h_n$ from Definition~\ref{def:general-kernel}.

\begin{theorem}[Christoffel kernel evaluation bounds]
\label{thm:kernel-evaluation-bound}\uses{thm:weighted-projection,def:general-kernel}
\lean{Polynomial.Sequence.inner_kernelVector, Polynomial.Sequence.IsOrthogonal.norm_kernelVector_sq, Polynomial.Sequence.IsOrthogonal.norm_kernelVector, Polynomial.Sequence.IsOrthogonal.norm_eval_partialSum_le_kernel, Polynomial.Sequence.IsOrthogonal.norm_eval_le_kernel_mul_norm, Polynomial.Sequence.IsOrthogonal.norm_sub_eval_partialSum_le_kernel}\leanok
The existing Christoffel kernel satisfies
$\langle v_{N,x},f\rangle=(S_Nf)(x)$ and $\|v_{N,x}\|_2^2=K_N(x,x)$.
Consequently, for $\deg p\le N$,
$|p(x)|\le\sqrt{K_N(x,x)}\|p\|_2$.
If $|f-p|\le\varepsilon$ almost everywhere and $|y-p(x)|\le\varepsilon$, then
\[ |y-(S_Nf)(x)|\le
 (1+\sqrt{K_N(x,x)}\sqrt{\mu(\R)})\varepsilon. \]
Here $f$ is a real L2 class and $y$ specifies the desired pointwise value;
no point evaluation of an L2 equivalence class is assumed.

Reference: \cite[18.2(v), (viii)]{dlmf}.
\end{theorem}
\begin{proof}
\leanok
\end{proof}

\begin{definition}[Explicit Jacobi growth constants]
\label{def:jacobi-growth-constants}\uses{def:jacobi,def:jacobi-measure,thm:gamma-function}
\lean{Polynomial.jacobiEndpointBound, Polynomial.jacobiReciprocalNormBound, Polynomial.jacobiKernelBound}\leanok
For $a+1\notin\{0,-1,-2,\ldots\}$, $\alpha,\beta>-1$, and an integer
$k\ge0$, define
\begin{align*}
 E_k(a)&=1+|a+1|+(k+1)^k/\Gamma(a+1),\\
 B_k(\alpha,\beta)&=1+h_0^{-1}+h_1^{-1}+h_2^{-1}
 +\frac{(2+|\alpha+\beta+3|)(k+1)^k}{2^{\alpha+\beta+1}},\\
 C_k(\alpha,\beta)&=[1+2(E_k(\alpha+1)+E_k(\beta))]^2B_k(\alpha,\beta),
\end{align*}
where $h_n=\int(P_n^{(\alpha,\beta)})^2\,d\mu_{\alpha,\beta}$.
The constants depend on $k$ and the parameters, not on $n$.
\end{definition}

\begin{theorem}[Uniform analytic Jacobi approximation and convergence]
\label{thm:jacobi-uniform-approximation}\uses{thm:analytic-kernel-error,thm:jacobi-kernel-growth,thm:jacobi-pairwise,def:jacobi-coefficients}
\lean{Polynomial.norm_sub_sum_jacobiCoefficient_le_uniform_of_analytic, tendsto_natCast_add_one_pow_div_const_pow, tendstoUniformlyOn_of_polynomial_geometric_bound, Polynomial.tendstoUniformlyOn_sum_jacobiCoefficient_of_analytic}\leanok
For $\alpha,\beta>-1$, $\rho>1$, $F$ analytic and bounded by $M$ on $E_\rho$,
and a real weighted L2 class $f=\operatorname{Re}F$ almost everywhere,
\[ |\operatorname{Re}F(x)-S_Nf(x)|\le
 [1+\sqrt{C_k(\alpha,\beta)}(N+1)^{3k+3}\sqrt{\mu_{\alpha,\beta}(\R)}]
 \frac{2M}{\rho^N(\rho-1)} \]
for every $x\in[-1,1]$ and any integer
$k\ge\max(\alpha+1,\beta,\alpha+\beta+1)$.
Such an integer always exists, and the Jacobi partial sums converge uniformly
on the closed interval, including both endpoints. No boundary regularity of
$F$ on the ellipse is assumed.
\end{theorem}
\begin{proof}
\leanok
\end{proof}

\begin{theorem}[Gegenbauer kernel and analytic approximation]
\label{thm:gegenbauer-analytic-approximation}\uses{thm:gegenbauer-norm-transfer,thm:jacobi-uniform-approximation,thm:gegenbauer-orthogonal,def:gegenbauer-coefficients,thm:weighted-analytic-approximation}
\lean{Polynomial.gegenbauer_kernel_eq_jacobi, Polynomial.norm_sub_sum_gegenbauerCoefficient_le_of_analytic, Polynomial.norm_sub_sum_gegenbauerCoefficient_le_uniform_of_analytic, Polynomial.tendstoUniformlyOn_sum_gegenbauerCoefficient_of_analytic}\leanok
For $\lambda>-1/2$, $\lambda\ne0$, the Gegenbauer kernel equals the Jacobi kernel
at $\alpha=\beta=\lambda-1/2$. Let $\rho>1$, let $F$ be analytic and bounded by
$M$ on $E_\rho$, and let $f\in L^2(\mu_\lambda;\R)$ agree almost everywhere with
$\operatorname{Re}F$. Its Gegenbauer partial sums satisfy
\[ \|f-S_Nf\|_2\le\sqrt{\mu_\lambda(\R)}\frac{2M}{\rho^N(\rho-1)}. \]
The preceding uniform Jacobi error bound also holds, with
$C_k(\lambda-1/2,\lambda-1/2)$ and the Gegenbauer mass, whenever
$k\ge\max(\lambda+1/2,2\lambda)$ is an integer. Consequently the expansions
converge uniformly on $[-1,1]$. Negative admissible parameters are included;
$\lambda=0$ is excluded because it does not give a polynomial basis.
\end{theorem}
\begin{proof}
\leanok
\end{proof}

\begin{supporting}[Positivity of explicit growth constants]

\begin{theorem}[Classical Jacobi growth constants]
\label{thm:jacobi-growth-constants}\uses{def:jacobi-growth-constants,thm:jacobi-positive,thm:jacobi-sequence}
\lean{Polynomial.jacobiEndpointBound_pos, Polynomial.integral_sq_jacobi_pos, Polynomial.jacobiReciprocalNormBound_pos, Polynomial.jacobiKernelBound_nonneg}\leanok
For $a,\alpha,\beta>-1$, every squared norm $h_n$ is strictly positive,
$E_k(a)>0$, $B_k(\alpha,\beta)>0$, and $C_k(\alpha,\beta)\ge0$.
Thus the reciprocals appearing in these constants are reciprocals of genuine
nonzero values, and the Gamma argument $a+1$ lies in the positive half-line.
\end{theorem}
\begin{proof}
\leanok
\end{proof}

\end{supporting}

\section{Complex data and differentiation}\label{sec:textbook-jacobi-approx-complex}

Complex coefficients use the same real orthogonal basis. Differentiating Jacobi expansions shifts both parameters; normal convergence justifies the termwise formulas.

The domains $\mathcal A_{\mu,p}^{\mathbb K}$ and $\mathcal A_{\mu,P,\le N}^{\mathbb K}$ require an integrable numerator and a finite positive squared norm for each coefficient being taken (Definitions~\ref{def:integral-coefficient-input} and~\ref{def:integral-coefficient-family-input}).

\begin{definition}[Scalar extensions and complex classical coefficients]
\label{def:scalar-orthogonal-coefficients}\uses{def:weighted-partial-sum,def:jacobi,def:gegenbauer,def:gegenbauer-measure}
\lean{Polynomial.polynomialToL2Scalar, Polynomial.Sequence.normalizedToL2Scalar, Polynomial.Sequence.integralCoefficientScalar, Polynomial.Sequence.partialSumValueScalar, Polynomial.Sequence.kernelVectorScalar, Polynomial.jacobiComplexCoefficient, Polynomial.gegenbauerComplexCoefficient}\leanok
For $\mathbb K=\R$ or $\C$, embed the existing real polynomial functions and
normalized vectors in $L^2(\mu;\mathbb K)$ at positive-norm indices. The scalar kernel vector is the real kernel vector embedded into this space, at a cutoff with positive squared norms throughout. For $f\in\mathcal A_{\mu,P_n}^{\mathbb K}$ and $f\in\mathcal A_{\mu,P,\le N}^{\mathbb K}$, respectively, define
\[ a_n^{\mathbb K}(f)=h_n^{-1}\int f(x)P_n(x)\,d\mu(x),\qquad
 S_N^{\mathbb K}f(x)=\sum_{n=0}^N a_n^{\mathbb K}(f)P_n(x). \]
For $f$ in the respective domains $\mathcal A_{\mu_{\alpha,\beta},P_n^{(\alpha,\beta)}}^{\C}$ and $\mathcal A_{\mu_\lambda,C_n^{(\lambda)}}^{\C}$, define the complex Jacobi and Gegenbauer coefficients
$h_n^{-1}\int fP_n^{(\alpha,\beta)}\,d\mu_{\alpha,\beta}$ and
$h_n^{-1}\int fC_n^{(\lambda)}\,d\mu_\lambda$, with $h_n$ the corresponding
squared norms, in the original polynomial normalizations.

Reference: \cite[18.18]{dlmf}.
\end{definition}

\begin{theorem}[Complex weighted and kernel approximation bounds]
\label{thm:complex-weighted-approximation}\uses{thm:scalar-orthogonal-projection,thm:geometric-error,thm:complex-chebyshev-competitor}
\lean{Polynomial.Sequence.IsOrthogonal.norm_sub_sum_integralCoefficientScalar_le_of_analytic, Polynomial.Sequence.IsOrthogonal.norm_sub_partialSumValueScalar_le_kernel_of_analytic}\leanok
Let $P$ be orthogonal for a positive-definite measure with finite moments and
almost-everywhere support in $[-1,1]$. Let $\rho>1$, let $F$ be analytic on the
open ellipse $E_\rho$ with $|F|\le M$, and let $f\in L^2(\mu;\C)$ equal $F$
almost everywhere on the real axis. With $\varepsilon_N=2M/[\rho^N(\rho-1)]$,
\[ \|f-S_N^{\C}f\|_2\le\sqrt{\mu(\R)}\varepsilon_N,\qquad
 |F(x)-S_N^{\C}f(x)|\le
 (1+\sqrt{K_N(x,x)}\sqrt{\mu(\R)})\varepsilon_N \]
for $x\in[-1,1]$. No connected-support or density hypothesis is needed.

Reference: \cite[Theorem 8.2]{atap}.
\end{theorem}
\begin{proof}
\leanok
\end{proof}

\begin{theorem}[Jacobi series differentiation with parameter shift]
\label{thm:jacobi-series-differentiation}\uses{thm:normal-series-differentiation,thm:jacobi-derivative,thm:jacobi-low-degree}
\lean{Polynomial.hasSum_derivative_jacobi_series}\leanok
For complex $\alpha,\beta$, suppose $a_nP_n^{(\alpha,\beta)}$ has a summable
uniform majorant on an open set $U$. Then throughout $U$,
\[
 \frac{d}{dz}\sum_{n\ge0}a_nP_n^{(\alpha,\beta)}(z)
 =\sum_{n\ge0}\frac{n+\alpha+\beta+2}{2}\,a_{n+1}
       P_n^{(\alpha+1,\beta+1)}(z).
\]
No real-weight parameter restrictions are required for this identity.

Reference: \cite[18.9.15]{dlmf}.
\end{theorem}
\begin{proof}
\leanok
\end{proof}

\begin{theorem}[Higher derivatives of Jacobi series]
\label{thm:jacobi-series-higher-derivatives}\uses{thm:jacobi-iterated-derivative,thm:normal-series-higher-derivatives}
\lean{Polynomial.hasSum_iterated_derivative_jacobi_series}\leanok
For complex $\alpha,\beta$, let $U\subset\C$ be open and suppose $\sum_n u_n<\infty$ with $|a_nP_n^{(\alpha,\beta)}(z)|\le u_n$ for every $z\in U$. For $S(z)=\sum_n a_nP_n^{(\alpha,\beta)}(z)$,
\[
 S^{(k)}(z)=\sum_{n\ge0}a_{n+k}
 \left[\prod_{j=0}^{k-1}\frac{n+k+\alpha+\beta+j+1}{2}\right]
 P_n^{(\alpha+k,\beta+k)}(z).
\]

Reference: \cite[18.9.15]{dlmf}, iterated.
\end{theorem}
\begin{proof}
\leanok
\end{proof}

\begin{supporting}[Projection constructions, kernel estimates and growth proofs]

\begin{theorem}[Scalar normalization and kernel-vector formulas]
\label{thm:scalar-normalization-characterisation}\uses{def:scalar-orthogonal-coefficients}
\lean{Polynomial.Sequence.normalizedToL2Scalar_eq, Polynomial.Sequence.kernelVectorScalar_apply}\leanok
Let $\mathbb K=\mathbb R$ or $\mathbb C$.
At a positive-norm index, the scalar normalized vector is
$h_n^{-1/2}J_\mu^{\mathbb K}P_n$. At a kernel cutoff, the scalar kernel
vector is the scalar embedding of $v_{N,x}$.
\end{theorem}
\begin{proof}
\leanok
\end{proof}

\begin{theorem}[Classical scalar coefficient and partial sum formulas]
\label{thm:scalar-coefficient-characterisation}\uses{def:scalar-orthogonal-coefficients}
\lean{Polynomial.Sequence.integralCoefficientScalar_apply, Polynomial.Sequence.partialSumValueScalar_apply}\leanok
Let $\mathbb K=\mathbb R$ or $\mathbb C$.
For $f\in\mathcal A_{\mu,P_n}^{\mathbb K}$,
$a_n^{\mathbb K}(f)=(\int P_n^2\,d\mu)^{-1}\int fP_n\,d\mu$,
with the real denominator embedded in $\mathbb K$.
For $f\in\mathcal A_{\mu,P,\le N}^{\mathbb K}$,
$S_N^{\mathbb K}f(x)=\sum_{n=0}^N a_n^{\mathbb K}(f)P_n(x)$.
\end{theorem}
\begin{proof}
\leanok
\end{proof}

\begin{theorem}[Classical complex Jacobi and Gegenbauer coefficient formulas]
\label{thm:complex-family-coefficient-characterisation}\uses{def:scalar-orthogonal-coefficients,def:jacobi-coefficients,def:gegenbauer-coefficients}
\lean{Polynomial.jacobiComplexCoefficient_apply, Polynomial.gegenbauerComplexCoefficient_apply, Polynomial.jacobiComplexCoefficient_ofReal, Polynomial.gegenbauerComplexCoefficient_ofReal}\leanok
For $f\in\mathcal A_{\mu_{\alpha,\beta},P_n^{(\alpha,\beta)}}^{\C}$, the
complex Jacobi coefficient is
\[
 \frac{\int fP_n^{(\alpha,\beta)}\,d\mu_{\alpha,\beta}}
 {\int (P_n^{(\alpha,\beta)})^2\,d\mu_{\alpha,\beta}}.
\]
For $f\in\mathcal A_{\mu_\lambda,C_n^{(\lambda)}}^{\C}$, the complex
Gegenbauer coefficient is
\[
 \frac{\int fC_n^{(\lambda)}\,d\mu_\lambda}
 {\int (C_n^{(\lambda)})^2\,d\mu_\lambda}.
\]
Both numerators are integrals of integrable functions and both denominators
are finite and positive. For real-valued $f$ these equal the real coefficients
viewed in $\C$.
\end{theorem}
\begin{proof}
\leanok
\end{proof}

\begin{definition}[Complex Chebyshev truncation in a weighted space]
\label{def:complex-chebyshev-competitor}\uses{def:scalar-orthogonal-coefficients,def:complex-coefficients}
\lean{Polynomial.complexChebyshevToL2}\leanok
Let $\mu$ have all moments finite and let $f\in\mathcal L^1_T(\C)$.
The degree-$N$ Chebyshev truncation in $L^2(\mu;\C)$ is
\[q_N=\sum_{n=0}^N a_n(f)[T_n],\]
where $[T_n]$ denotes the polynomial's $L^2(\mu;\C)$ class.
Reference: \cite[Theorem 8.2]{atap}.
\end{definition}

\begin{theorem}[Weighted Chebyshev truncation representation]
\label{thm:complex-chebyshev-competitor}\uses{def:complex-chebyshev-competitor,thm:general-polynomial-L2-basic}
\lean{Polynomial.coe_complexChebyshevToL2}\leanok
For $f\in\mathcal L^1_T(\C)$, the weighted truncation is represented
$\mu$-almost everywhere by $\sum_{n=0}^N a_n(f)T_n(x)$.
\end{theorem}
\begin{proof}
\leanok
\end{proof}

\begin{theorem}[Best finite orthonormal approximation]
\label{thm:finite-best-approximation}
\lean{Orthonormal.inner_sub_sum_inner_sum_eq_zero, Orthonormal.norm_sub_sum_inner_le, Orthonormal.norm_sub_sum_inner_le_norm, Orthonormal.norm_sub_sum_inner_le_of_mem_span}\leanok
For an orthonormal family $v_i$ in a real or complex inner-product space and a
finite index set $s$, the sum $Q_s f=\sum_{i\in s}\langle v_i,f\rangle v_i$
minimises $\|f-y\|$ over the span of the selected vectors. Its residual is
orthogonal to that span. No completeness or Hilbert-basis hypothesis is needed.

Reference: the Hilbert-space best-approximation theorem.
\end{theorem}
\begin{proof}
\leanok
\end{proof}

\begin{theorem}[Real analytic polynomial competitors]
\label{thm:real-analytic-competitor}\uses{thm:geometric-error}
\lean{Polynomial.Chebyshev.exists_real_polynomial_approximation_of_analytic}\leanok
If $\rho>1$ and $F$ is analytic and bounded in norm by $M$ in the open ellipse
$E_\rho$, there is a real polynomial $p_N$ of degree at most $N$ satisfying
\[ |\operatorname{Re}F(x)-p_N(x)|\le\frac{2M}{\rho^N(\rho-1)}
 \qquad(-1\le x\le1). \]
No boundary regularity is needed.

Reference: \cite[Theorem 8.2]{atap}.
\end{theorem}
\begin{proof}
\leanok
\end{proof}

\begin{theorem}[Finite projection and reproduction]
\label{thm:weighted-projection}\uses{def:weighted-partial-sum,thm:weighted-best-approximation}
\lean{Polynomial.polynomialToL2_injective, Polynomial.Sequence.degree_partialSum_le, Polynomial.Sequence.polynomialToL2_partialSum, Polynomial.Sequence.IsOrthogonal.partialSum_polynomialToL2, Polynomial.Sequence.partialSum_sub, Polynomial.Sequence.IsOrthogonal.norm_polynomialToL2_partialSum_le}\leanok
Under finite moments, positive definiteness, and orthogonality, the polynomial
map into L2 is injective, $\deg S_Nf\le N$, and $S_N$ reproduces every polynomial
of degree at most $N$. It respects subtraction and has weighted L2 norm at most
one. These statements impose no support or density hypothesis.
\end{theorem}
\begin{proof}
\leanok
\end{proof}

\begin{theorem}[Analytic projection error with kernel factor]
\label{thm:analytic-kernel-error}\uses{thm:kernel-evaluation-bound,thm:real-analytic-competitor}
\lean{Polynomial.Sequence.IsOrthogonal.norm_sub_eval_partialSum_le_kernel_of_analytic}\leanok
Let $P$ be orthogonal for a measure $\mu$ with finite polynomial moments and positive definite polynomial pairing, supported almost everywhere on $[-1,1]$. Let $\rho>1$, let $F$ be analytic on $E_\rho$ with $|F|\le M$ there, and let $f\in L^2(\mu;\R)$ agree almost everywhere with $\operatorname{Re}F$. Then for every
$x\in[-1,1]$,
\[ |\operatorname{Re}F(x)-(S_Nf)(x)|\le
 (1+\sqrt{K_N(x,x)}\sqrt{\mu(\R)})
 \frac{2M}{\rho^N(\rho-1)}. \]

\end{theorem}
\begin{proof}
\leanok
\end{proof}

\rvnote{thm:analytic-kernel-error}{A growth estimate on the kernel is still needed to infer uniform convergence.}

\begin{theorem}[Elementary Gamma and rising-factorial growth bounds]
\label{thm:gamma-polynomial-growth}\uses{thm:gamma-function}
\lean{Real.Gamma_add_nat_eq_ascPochhammer_mul, Real.Gamma_nat_add_le_factorial_mul_pow, Real.factorial_le_Gamma_nat_add, Real.ascPochhammer_div_factorial_le_polynomial}\leanok
For $a>0$, $\Gamma(a+n)=(a)_n\Gamma(a)$. If $a\le k+1$ and $n+a\ge2$,
$\Gamma(n+a)\le n!(n+k)^k$. If $n\ge1$ and $a\ge1$,
$\Gamma(n+a)\ge n!$. For $a>0$ and $a\le k+1$, all $n\ge0$ satisfy
\[ \frac{(a)_n}{n!}\le
 \left(1+|a|+\frac{(k+1)^k}{\Gamma(a)}\right)(n+1)^k. \]
These bounds are nonsharp.

Reference: \cite[5.5.1, 5.6]{dlmf}.
\end{theorem}
\begin{proof}
\leanok
\end{proof}

\begin{theorem}[Jacobi endpoint control from the differential equation]
\label{thm:jacobi-sonin-bound}\uses{thm:jacobi-parameter-raising}
\lean{Polynomial.hasDerivAt_jacobiSoninEnergy, Polynomial.abs_eval_jacobi_le_max_endpoints, Polynomial.abs_eval_jacobi_succ_le_shifted_endpoints}\leanok
For $p=P_n^{(\alpha,\beta)}$ and $c=n(n+\alpha+\beta+1)$, the energy
$E(x)=cp(x)^2+(1-x^2)p'(x)^2$ satisfies
$E'(x)=2[\alpha-\beta+(\alpha+\beta+1)x]p'(x)^2$.
If $\alpha+\beta+1\ge0$, then
$|p(x)|\le\max(|p(-1)|,|p(1)|)$ on $[-1,1]$.
For every $\alpha,\beta>-1$, parameter raising instead bounds
$|P_{n+1}^{(\alpha,\beta)}(x)|$ by the sum of the endpoint maxima of
$P_{n+1}^{(\alpha+1,\beta)}$ and $P_n^{(\alpha+1,\beta)}$.
Thus the final estimate covers the whole orthogonality range.

References: \cite[Table 18.8.1 and 18.9(ii)]{dlmf}; cf. \cite[18.14(i)]{dlmf}.
\end{theorem}
\begin{proof}
\leanok
\end{proof}

\begin{theorem}[Global polynomial growth of Jacobi values and reciprocal norms]
\label{thm:jacobi-polynomial-growth}\uses{thm:gamma-polynomial-growth,thm:jacobi-sonin-bound,thm:jacobi-growth-constants,thm:jacobi-norm}
\lean{Polynomial.abs_eval_jacobi_one_le_polynomial, Polynomial.abs_eval_jacobi_neg_one_le_polynomial, Polynomial.max_abs_eval_jacobi_endpoints_le_polynomial, Polynomial.abs_eval_jacobi_le_polynomial, Polynomial.inv_integral_sq_jacobi_succ_le, Polynomial.inv_integral_sq_jacobi_succ_le_polynomial, Polynomial.inv_integral_sq_jacobi_le_polynomial}\leanok
Let $\alpha,\beta>-1$ and choose an integer
$k\ge\max(\alpha+1,\beta,\alpha+\beta+1)$. For every $n\ge0$ and $x\in[-1,1]$,
\begin{align*}
 |P_n^{(\alpha,\beta)}(x)|
 &\le[1+2(E_k(\alpha+1)+E_k(\beta))](n+1)^k,\\
 h_n^{-1}&\le B_k(\alpha,\beta)(n+1)^{k+2}.
\end{align*}
The norms and the constants $E_k$ and $B_k$ are positive in this range.
\end{theorem}
\begin{proof}
\leanok
\end{proof}

\begin{theorem}[Uniform Jacobi kernel growth]
\label{thm:jacobi-kernel-growth}\uses{thm:jacobi-polynomial-growth,thm:jacobi-moments,def:general-kernel,thm:general-squared-norm}
\lean{Polynomial.jacobi_kernel_term_le_polynomial, Polynomial.jacobi_kernel_le_polynomial, Polynomial.sqrt_jacobi_kernel_le_polynomial}\leanok
For $\alpha,\beta>-1$ and an integer $k\ge\max(\alpha+1,\beta,\alpha+\beta+1)$, with $C_k$ as in Theorem~\ref{thm:jacobi-growth-constants},
\[ K_N(x,x)\le C_k(\alpha,\beta)(N+1)^{3k+3},\qquad
 \sqrt{K_N(x,x)}\le\sqrt{C_k(\alpha,\beta)}(N+1)^{3k+3} \]
uniformly for $x\in[-1,1]$.
\end{theorem}
\begin{proof}
\leanok
\end{proof}

\begin{theorem}[Scalar orthogonality and kernel reproduction]
\label{thm:scalar-orthogonal-projection}\uses{thm:kernel-evaluation-bound,thm:scalar-normalization-characterisation,thm:gegenbauer-sequence,thm:gegenbauer-moments}
\lean{inner_ofReal_compLp, norm_ofReal_compLp, Polynomial.inner_polynomialToL2Scalar_left, Polynomial.Sequence.IsOrthogonal.orthonormal_normalizedToL2Scalar, Polynomial.Sequence.polynomialToL2Scalar_mem_span, Polynomial.Sequence.inner_normalizedToL2Scalar_smul, Polynomial.Sequence.inner_kernelVectorScalar, Polynomial.Sequence.IsOrthogonal.norm_kernelVectorScalar, Polynomial.Sequence.IsOrthogonal.inner_kernelVectorScalar_polynomialToL2Scalar, Polynomial.integralCoefficientScalar_jacobiSequence, Polynomial.integralCoefficientScalar_gegenbauerSequence}\leanok
Let $\mathbb K=\mathbb R$ or $\mathbb C$.
The real orthonormal vectors remain orthonormal after scalar extension.
Degree-at-most-$N$ real polynomials belong to their first $N+1$-dimensional
$\mathbb K$-span. The lifted kernel vector represents $S_N^{\mathbb K}$ and has
norm $\sqrt{K_N(x,x)}$. It reproduces degree-at-most-$N$ real polynomials.
The named complex classical coefficients agree with the real coefficients on
admissible real inputs, and for $\alpha,\beta>-1$, respectively $\lambda>-1/2$, $\lambda\ne0$,
they are the scalar coefficients $a_n^{\mathbb C}$ of the Jacobi, respectively
Gegenbauer, sequence, on the coefficient domains of Definition~\ref{def:scalar-orthogonal-coefficients}.

Reference: \cite[18.2(viii), 18.18]{dlmf}.
\end{theorem}
\begin{proof}
\leanok
\end{proof}

\begin{theorem}[Iterated Jacobi differentiation]
\label{thm:jacobi-iterated-derivative}\uses{thm:jacobi-derivative,thm:jacobi-degree}
\lean{Polynomial.iterate_derivative_jacobi, Polynomial.iterate_derivative_jacobi_eq_zero}\leanok
Over a commutative $\mathbb Q$-algebra, for $n,k\ge0$,
\[
 D^kP_{n+k}^{(\alpha,\beta)}=
 \left[\prod_{j=0}^{k-1}\frac{n+k+\alpha+\beta+j+1}{2}\right]
 P_n^{(\alpha+k,\beta+k)}.
\]
Also $D^kP_n^{(\alpha,\beta)}=0$ when $n<k$.
There are no positivity restrictions on the parameters.

Reference: \cite[18.9.15]{dlmf}, iterated.
\end{theorem}
\begin{proof}
\leanok
\end{proof}

\begin{theorem}[Approximation in other classical families]
\label{thm:beyond-chebyshev}\uses{thm:gegenbauer-analytic-approximation,thm:complex-weighted-approximation}
\lean{Polynomial.norm_sub_sum_jacobiComplexCoefficient_le_of_analytic, Polynomial.norm_sub_sum_jacobiComplexCoefficient_le_uniform_of_analytic, Polynomial.tendstoUniformlyOn_sum_jacobiComplexCoefficient_of_analytic, Polynomial.norm_sub_sum_gegenbauerComplexCoefficient_le_of_analytic, Polynomial.norm_sub_sum_gegenbauerComplexCoefficient_le_uniform_of_analytic, Polynomial.tendstoUniformlyOn_sum_gegenbauerComplexCoefficient_of_analytic}\leanok
Both real and complex analytic Jacobi and Gegenbauer expansions satisfy weighted
L2 error bounds and converge uniformly on $[-1,1]$. The complex bounds have
exactly the constants of the real-valued theorems above, with $|F-S_Nf|$ replacing
$|\operatorname{Re}F-S_Nf|$. They hold for every $\alpha,\beta>-1$ and every
$\lambda>-1/2$, $\lambda\ne0$, respectively, under the open-ellipse hypotheses
$\rho>1$, analyticity and $|F|\le M$.
The uniform polynomial factors are explicit and nonsharp; no boundary regularity
on the ellipse is assumed.

References: \cite[18.18]{dlmf} and \cite[Chapter 8]{atap}.
\end{theorem}
\begin{proof}
\leanok
\end{proof}

\end{supporting}

\chapter{Gaussian quadrature and the Jacobi matrix}\label{chap:gauss}

\rvsummary{Gaussian quadrature for a general measure and its Golub--Welsch eigenvalue characterisation.}

Gaussian quadrature combines orthogonality with interpolation. The spectral description of its nodes and weights gives the Golub--Welsch construction.

\section{Exactness and positive weights}\label{sec:textbook-gauss-gauss}

Integrating the cardinal polynomials gives an interpolatory rule. Choosing the zeros of an orthogonal polynomial as nodes raises its exactness from degree below $n$ to degree below $2n$.

\begin{definition}[Classical quadrature inputs]
\label{def:quadrature-input}\uses{def:lagrange-basis}
\lean{Lagrange.QuadratureInput}\leanok
For a measure $\mu$, finite nodes $v$ indexed by $s$, and an index $i$,
the quadrature-weight domain requires $v_i\ne v_j$ for
$j\in s\setminus\{i\}$ and $\mu$-integrability of
$\ell_i(X)=\prod_{j\in s\setminus\{i\}}(X-v_j)/(v_i-v_j)$.
Only this individual polynomial is required to be integrable.
\end{definition}

\begin{definition}[Interpolatory and Gaussian weights]
\label{def:quadrature-weights}\uses{def:quadrature-input}
\lean{Lagrange.quadratureWeight, Polynomial.Sequence.gaussWeight}\leanok
For a measure $\mu$ on $\mathbb R$ and a finite node set, define
$w_i=\int\ell_i\,d\mu$ on the domain of Definition~\ref{def:quadrature-input}, where $\ell_i$ is the Lagrange basis polynomial of the nodes
(Theorem~\ref{thm:lagrange-interpolation}).
For a polynomial sequence $P$, its degree-$n$ Gaussian weights use the finite
set $s_n$ of distinct real roots of $P_n$, with the identity node map.
No ordering of the roots is chosen. The normalized-product integral also
defines a weight at an off-root point whenever it is integrable; the identity
node map has no vanishing denominator at any real point.

Reference: \cite[3.5(iv)--(v)]{dlmf}.
\end{definition}

\begin{theorem}[Interpolatory quadrature exactness]
\label{thm:interpolatory-quadrature}\uses{def:quadrature-weights,thm:lagrange-interpolation,thm:quadrature-input}
\lean{Lagrange.integral_eq_sum_quadratureWeight_mul_eval}\leanok
Assume finite polynomial moments. For distinct nodes indexed by $s$ and
$\deg p<|s|$,
\[
 \int p\,d\mu=\sum_{i\in s}w_i p(v_i).
\]

Reference: \cite[3.5(iv)]{dlmf}.
\end{theorem}
\begin{proof}
\leanok
\end{proof}

\begin{theorem}[General Gaussian quadrature exactness]
\label{thm:general-gauss-exact}\uses{thm:interpolatory-quadrature,thm:general-real-zeros}
\lean{Polynomial.Sequence.IsOrthogonal.card_gaussNodes, Polynomial.Sequence.IsOrthogonal.integral_eq_sum_gaussWeight_mul_eval}\leanok
Let $P$ be an orthogonal polynomial sequence for $\mu$, under the finite-moment
and positive-definiteness hypotheses. Its root set $s_n$ has cardinality $n$.
For every polynomial $p$ with $\deg p<2n$,
\[
 \int p\,d\mu=\sum_{x\in s_n}w_xp(x).
\]
This includes the zero polynomial and requires no boundedness or connectedness
of the support.

References: \cite[3.5(v)]{dlmf}; \cite[Section 3.4]{szego}.
\end{theorem}
\begin{proof}
\leanok
\end{proof}

\begin{theorem}[Positive Gaussian weights]
\label{thm:general-gauss-positive}\uses{thm:general-gauss-exact}
\lean{Polynomial.Sequence.IsOrthogonal.gaussWeight_eq_integral_sq, Polynomial.Sequence.IsOrthogonal.gaussWeight_pos}\leanok
Let $P$ be an orthogonal polynomial sequence for a measure $\mu$ with finite polynomial moments and a positive definite polynomial pairing. Let $s_n$ be the distinct real roots of $P_n$ and $w_x$ their Gaussian weights. Then for every $x\in s_n$,
\[
 w_x=\int\ell_x^2\,d\mu>0.
\]

Reference: \cite[3.5(v)]{dlmf}.
\end{theorem}
\begin{proof}
\leanok
\end{proof}

\begin{theorem}[Total Gaussian weight]
\label{thm:general-gauss-mass}\uses{thm:general-gauss-exact}
\lean{Polynomial.Sequence.IsOrthogonal.sum_gaussWeight}\leanok
Let $P$ be an orthogonal polynomial sequence for a measure $\mu$ with finite polynomial moments and a positive definite polynomial pairing. Let $s_n$ be the distinct real roots of $P_n$ and $w_x$ their Gaussian weights. Then for $n\ge1$,
$\sum_{x\in s_n}w_x=\mu(\mathbb R)$, with the finite mass viewed as a real number.

Reference: \cite[3.5(v)]{dlmf}.
\end{theorem}
\begin{proof}
\leanok
\end{proof}

\begin{supporting}[Quadrature-weight domains and formulas]

\begin{theorem}[Classical quadrature input characterisation]
\label{thm:quadrature-input}\uses{def:quadrature-input,thm:polynomial-inner}
\lean{Lagrange.QuadratureInput.ofMoments,Lagrange.QuadratureInput.integrable_eval}\leanok
On the quadrature-weight domain, $\ell_i$ is $\mu$-integrable and
$\int\ell_i\,d\mu$ is finite.
If $\mu$ has finite polynomial moments, this integrability holds whenever
$v_i-v_j\ne0$ for all $j\in s\setminus\{i\}$, including admissible $i\notin s$.
\end{theorem}
\begin{proof}
\leanok
\end{proof}

\begin{theorem}[Classical quadrature weight formulas]
\label{thm:quadrature-weight-formulas}\uses{def:quadrature-weights}
\lean{Lagrange.quadratureWeight_apply,Polynomial.Sequence.gaussWeight_apply}\leanok
For each index $i$ in the quadrature-weight domain of Definition~\ref{def:quadrature-input}, the interpolatory weight is $\int\ell_i\,d\mu$.
For a polynomial sequence, its Gaussian weight is the same integral for
the identity nodes given by the distinct real roots.
\end{theorem}
\begin{proof}
\leanok
\end{proof}

\end{supporting}

\section{Gauss--Jacobi quadrature and error}\label{sec:textbook-gauss-jacobi}

For Jacobi weights, the nodes lie in $(-1,1)$. Positivity transfers a uniform polynomial approximation bound to a quadrature error bound.

\begin{definition}[Gauss--Jacobi weights]
\label{def:jacobi-gauss-weights}\uses{def:quadrature-weights,def:jacobi-measure,def:jacobi}
\lean{Polynomial.jacobiGaussWeight}\leanok
Use the roots of $P_n^{(\alpha,\beta)}$ as nodes and integrate their Lagrange
basis polynomials against the original Jacobi measure, on the domain where the normalized nodal product is integrable, as in Definition~\ref{def:quadrature-input}. These weights do not
renormalise the measure to have total mass one.

Reference: \cite[3.5(v)]{dlmf}.
\end{definition}

\begin{theorem}[Gauss--Jacobi quadrature]
\label{thm:quadrature}\uses{def:jacobi-gauss-weights,thm:general-gauss-positive,thm:general-gauss-mass,thm:jacobi-root-location}
\lean{Polynomial.card_jacobiGaussNodes, Polynomial.jacobiGaussNodes_subset_Ioo, Polynomial.integral_eq_sum_jacobiGaussWeight_mul_eval, Polynomial.jacobiGaussWeight_pos, Polynomial.sum_jacobiGaussWeight}\leanok
For $\alpha,\beta>-1$ and $n\ge1$, the $n$ distinct zeros of
$P_n^{(\alpha,\beta)}$ lie in $(-1,1)$ and carry strictly positive quadrature
weights. The resulting rule is exact for all polynomials of degree at most
$2n-1$ (stated in Lean as $\deg p<2n$), and its weights sum to the total Jacobi
mass.

Reference: \cite[3.5(v)]{dlmf}.
\end{theorem}
\begin{proof}
\leanok
\end{proof}

\begin{theorem}[Gaussian quadrature competitor bound]
\label{thm:general-gauss-error}\uses{thm:positive-quadrature-error,thm:general-gauss-positive,thm:general-gauss-mass}
\lean{Polynomial.Sequence.IsOrthogonal.abs_integral_sub_gauss_le_of_competitor}\leanok
Under the general Gaussian moment, positivity, and orthogonality hypotheses,
let $n\ge1$ and let $f$ be integrable. Any real polynomial $p$ with
$\deg p<2n$ and $|f-p|\le\varepsilon$ both almost everywhere and at all
Gaussian nodes gives error at most $2\mu(\R)\varepsilon$.
No interval or connected-support hypothesis is needed.

Reference: \cite[proof of Theorem 19.3]{atap}.
\end{theorem}
\begin{proof}
\leanok
\end{proof}

\begin{theorem}[Geometric Gauss--Jacobi quadrature error]
\label{thm:jacobi-gauss-geometric-error}\uses{thm:general-gauss-error,thm:quadrature,thm:real-analytic-competitor}
\lean{Polynomial.abs_integral_sub_jacobiGauss_le_of_competitor, Polynomial.abs_integral_sub_jacobiGauss_le_of_analytic}\leanok
Let $\alpha,\beta>-1$ and $n\ge1$. A uniform polynomial competitor on $[-1,1]$
of degree below $2n$ and error at most $\varepsilon$ bounds Gauss--Jacobi
quadrature error by $2\mu_{\alpha,\beta}(\R)\varepsilon$, for an integrable $f$.
In particular, if $\rho>1$, $F$ is analytic on the open ellipse $E_\rho$,
and $|F|\le M$ there, then for $f(x)=\Re F(x)$,
\[
 \left|\int f\,d\mu_{\alpha,\beta}-Q_n(f)\right|
 \le\frac{4\mu_{\alpha,\beta}(\R)M}{\rho^{2n-1}(\rho-1)}.
\]
The original Jacobi mass is retained. This is a baseline geometric estimate,
not the sharper Legendre-specific constant of \cite[Theorem 19.3]{atap}.

Reference: \cite[proof of Theorem 19.3]{atap}.
\end{theorem}
\begin{proof}
\leanok
\end{proof}

\begin{supporting}[Gauss--Jacobi weight formulas]

\begin{theorem}[Gauss--Jacobi weight formula]
\label{thm:jacobi-gauss-weight-formula}\uses{def:jacobi-gauss-weights}
\lean{Polynomial.jacobiGaussWeight_apply}\leanok
Let $s$ be the distinct real roots of $P_n^{(\alpha,\beta)}$ and let $x$
lie in the quadrature-weight domain of Definition~\ref{def:quadrature-input}
for the identity nodes on $s$ and measure $\mu_{\alpha,\beta}$. Then
\[
 w_x=\int\ell_x(t)\,d\mu_{\alpha,\beta}(t),\qquad
 \ell_x(t)=\prod_{y\in s\setminus\{x\}}\frac{t-y}{x-y}.
\]
\end{theorem}
\begin{proof}
\leanok
\end{proof}

\end{supporting}

\section{The Jacobi matrix and Golub--Welsch}\label{sec:textbook-gauss-spectrum}

The recurrence matrix gives the nodes as eigenvalues. After normalization it is symmetric; the first components of its unit eigenvectors give the weights.

\begin{definition}[Finite recurrence multiplication section]
\label{def:finite-jacobi-matrix}
\lean{Polynomial.Sequence.jacobiMatrix}\leanok
For a degree-graded sequence, define $A_n$ by $(A_n)_{ij}=[xP_j]_i$ for
$0\le i,j<n$. For an orthogonal sequence we call it the finite Jacobi matrix
in the original polynomial normalization.

Reference: \cite[3.5(vi)]{dlmf}.
\end{definition}

\begin{theorem}[Finite Jacobi characteristic polynomial]
\label{thm:finite-jacobi-charpoly}\uses{def:finite-jacobi-matrix,thm:comrade-matrix}
\lean{Polynomial.Sequence.repr_X_mul_mod_next, Polynomial.Sequence.charpoly_jacobiMatrix, Polynomial.Sequence.mem_spectrum_jacobiMatrix_iff, Polynomial.Sequence.sum_jacobiMatrix_mul_eval}\leanok
For any degree-graded sequence over a field,
\[
 \det(tI-A_n)=\frac{P_n(t)}{\operatorname{lc}(P_n)}.
\]
Thus the spectrum of $A_n$ consists precisely of the roots of $P_n$ in the
base field. At a root $x$, the evaluation vector satisfies
$\sum_{j<n}(A_n)_{ji}P_j(x)=xP_i(x)$.

References: \cite[3.5(vi)]{dlmf}; \cite[Chapter 18]{atap}.
\end{theorem}
\begin{proof}
\leanok
\end{proof}

\begin{definition}[Symmetric normalized Jacobi section]
\label{def:normalized-jacobi-matrix}\uses{def:general-squared-norm}
\lean{Polynomial.Sequence.NormalizedJacobiMatrixInput, Polynomial.Sequence.normalizedJacobiMatrix}\leanok
For a real polynomial sequence, a measure and a finite size $n$, the matrix domain requires integrability of every mixed product $xP_jP_i$ and every square $P_i^2$ for $i,j<n$, together with $h_i=\int P_i^2\,d\mu>0$. On this domain set
\[
 (J_n)_{ij}=\frac{\int xP_j(x)P_i(x)\,d\mu(x)}{\sqrt{h_i}\sqrt{h_j}}.
\]


Reference: \cite[3.5(vi)]{dlmf}.
\end{definition}

\begin{theorem}[Symmetric Jacobi spectrum and nonzero adjacent entries]
\label{thm:normalized-jacobi-spectrum}\uses{thm:finite-jacobi-charpoly,thm:general-coefficient-signs,thm:normalized-jacobi-characterisation}
\lean{Polynomial.Sequence.IsOrthogonal.normalizedJacobiMatrix_eq_zero, Polynomial.Sequence.IsOrthogonal.normalizedJacobiMatrix_entry, Polynomial.Sequence.IsOrthogonal.charpoly_normalizedJacobiMatrix, Polynomial.Sequence.IsOrthogonal.mem_spectrum_normalizedJacobiMatrix_iff, Polynomial.Sequence.IsOrthogonal.normalizedJacobiMatrix_superdiag_ne_zero}\leanok
Assume finite moments, positive-definite polynomial integration and orthogonality.
Then $J_n$ is symmetric tridiagonal with nonzero adjacent entries and
\[
 (J_n)_{ij}=\sqrt{h_i}(A_n)_{ij}/\sqrt{h_j},\qquad
 \det(tI-J_n)=P_n(t)/\operatorname{lc}(P_n).
\]
Its spectrum is exactly the Gaussian nodes. Signs of leading coefficients are
unrestricted; adjacent entries need not be positive in this convention.

Reference: \cite[3.5(vi)]{dlmf}.
\end{theorem}
\begin{proof}
\leanok
\end{proof}

\begin{definition}[Classical Gaussian evaluation-vector inputs]
\label{def:gauss-vector-input}\uses{def:general-squared-norm,def:quadrature-weights}
\lean{Polynomial.Sequence.GaussEigenvectorInput}\leanok
For a point $x$ and each coordinate $i<n$, the Gauss-vector domain requires integrability of the Gaussian
weight's normalized nodal product $\ell_x$, nonnegativity of $w_x=\int\ell_x\,d\mu$, and an
integrable polynomial square with strictly positive squared norm $h_i$.
No conditions are required when $n=0$. A point need not be a Gaussian node if these
local conditions hold.
\end{definition}

\begin{definition}[Classical Gaussian evaluation vector]
\label{def:gauss-vector}\uses{def:gauss-vector-input}
\lean{Polynomial.Sequence.gaussEigenvector}\leanok
On the domain of Definition~\ref{def:gauss-vector-input}, define $u_i(x)=\sqrt{w_x}P_i(x)/\sqrt{h_i}$,
$i<n$. The square roots have nonnegative arguments and their denominators
are positive. Reference: \cite[3.5(vi)]{dlmf}.
\end{definition}

\begin{theorem}[Golub--Welsch nodes and weights]
\label{thm:golub-welsch}\uses{thm:normalized-jacobi-spectrum,thm:gauss-evaluation-vector,thm:hessenberg-eigenvector}
\lean{Polynomial.Sequence.IsOrthogonal.normalizedJacobiMatrix_mulVec_gaussEigenvector, Polynomial.Sequence.IsOrthogonal.golubWelsch, Polynomial.Sequence.IsOrthogonal.gaussWeight_eq_of_unit_eigenvector}\leanok
Under finite moments, positive-definite polynomial integration and orthogonality,
the canonical vector $u(x)$ at each Gaussian node is a unit eigenvector of $J_n$
with eigenvalue $x$. For \emph{every} real unit eigenvector $v$ for that node,
\[
 w_x=\mu(\mathbb R)v_0^2.
\]
Thus the formula is independent of eigenvector sign. These statements concern
exact matrices and vectors, without a floating-point eigensolver error claim.

References: \cite[3.5(vi)]{dlmf}; \cite{golub-welsch}.
\end{theorem}
\begin{proof}
\leanok
\end{proof}

\begin{theorem}[Jacobi and Chebyshev spectral quadrature]
\label{thm:classical-spectral-quadrature}\uses{thm:golub-welsch,thm:T-sequence,thm:classical-jacobi-measures,thm:jacobi-positive,thm:jacobi-pairwise,def:jacobi-gauss-weights}
\lean{Polynomial.mem_spectrum_jacobi_iff, Polynomial.jacobi_golubWelsch, Polynomial.jacobiGaussWeight_eq_of_unit_eigenvector, Polynomial.Chebyshev.hasFiniteMoments_measureT, Polynomial.Chebyshev.hasPositiveDefiniteIntegral_measureT, Polynomial.Chebyshev.mem_spectrum_normalizedJacobiMatrix_iff, Polynomial.Chebyshev.gaussWeight_eq_of_unit_eigenvector}\leanok
Let $P_n$ be either $P_n^{(\alpha,\beta)}$ with $\alpha,\beta>-1$, or
$T_n$, and let $\mu$ be its original orthogonality measure.
The eigenvalues of the normalized Jacobi matrix $J_n$ are exactly the roots
of $P_n$. For $n\ge1$ and every real unit eigenvector $v$ of $J_n$ with
eigenvalue $x$, the Gaussian weight is
\[
 w_x=\mu(\R)v_0^2.
\]
Reference: \cite[3.5(vi)]{dlmf}.
\end{theorem}
\begin{proof}
\leanok
\end{proof}

\paragraph{Example (Julia).} For the Legendre measure, the total mass is $2$. Golub--Welsch therefore gives weights equal to twice the squared first eigenvector components. The example compares the computed rule with exact monomial integrals through degree $2n-1$ and at the next even degree.

\begin{juliacode}
using LinearAlgebra
n = 5
b = [k / sqrt(4k^2 - 1) for k in 1:n-1]        # orthonormal recurrence
E = eigen(SymTridiagonal(zeros(n), b))
x, w = E.values, 2 .* E.vectors[1, :].^2
maximum(abs(dot(w, x.^k) - (iseven(k) ? 2/(k+1) : 0)) for k in 0:2n-1)
                                               # residual through the exactness degree
dot(w, x.^(2n)) - 2/(2n+1)                     # error at the next even degree
\end{juliacode}

\begin{supporting}[Normalization, evaluation vectors and eigenvector uniqueness]

\begin{theorem}[Classical normalized-section characterization]
\label{thm:normalized-jacobi-characterisation}\uses{def:normalized-jacobi-matrix,thm:general-tridiagonal,thm:general-squared-norm}
\lean{Polynomial.Sequence.NormalizedJacobiMatrixInput.ofMoments, Polynomial.Sequence.NormalizedJacobiMatrixInput.denominator_pos, Polynomial.Sequence.normalizedJacobiMatrix_mul_norm, Polynomial.Sequence.normalizedJacobiMatrix_isSymm}\leanok
Every normalizing denominator $\sqrt{h_i}\sqrt{h_j}$ is positive, and
$(J_n)_{ij}\sqrt{h_i}\sqrt{h_j}=\int xP_jP_i\,d\mu$.
The resulting matrix is symmetric. Finite polynomial moments and a positive
definite polynomial pairing ensure these integrability and positivity conditions for every finite size.
The size-zero domain has no integrability or positivity obligations.
\end{theorem}
\begin{proof}
\leanok
\end{proof}

\begin{theorem}[Classical Gaussian vector input characterisation]
\label{thm:gauss-vector-input}\uses{def:gauss-vector-input,thm:general-gauss-positive}
\lean{Polynomial.Sequence.GaussEigenvectorInput.valid,Polynomial.Sequence.GaussEigenvectorInput.empty,Polynomial.Sequence.IsOrthogonal.gaussEigenvectorInput}\leanok
On the Gauss-vector domain, each $i<n$ satisfies $w_x\ge0$ and
$h_i=\int P_i^2\,d\mu>0$. For $n=0$ the domain imposes no conditions.
If $\mu$ has finite polynomial moments and a positive definite polynomial
pairing and $P$ is orthogonal for $\mu$, every Gaussian node belongs to this domain.
\end{theorem}
\begin{proof}
\leanok
\end{proof}

\begin{theorem}[Gaussian evaluation-vector coordinate formula]
\label{thm:gauss-vector-formula}\uses{def:gauss-vector}
\lean{Polynomial.Sequence.gaussEigenvector_apply}\leanok
Every coordinate of the vector equals $\sqrt{w_x}P_i(x)/\sqrt{h_i}$.
\end{theorem}
\begin{proof}
\leanok
\end{proof}

\begin{theorem}[Positive quadrature competitor bound]
\label{thm:positive-quadrature-error}
\lean{MeasureTheory.abs_integral_sub_sum_le_of_nonneg_weights}\leanok
Let $\mu$ be any finite measure and $Q(f)=\sum_{i\in s}w_i f(v_i)$ a finite
rule with $w_i\ge0$ and $\sum_iw_i=\mu(\Omega)$. If $f,g$ are integrable,
$Q(g)=\int g\,d\mu$, and $|f-g|\le\varepsilon$ both almost everywhere and
at every sampled node, then
\[
 \left|\int f\,d\mu-Q(f)\right|\le2\mu(\Omega)\varepsilon.
\]
The measured space need not be an interval or even a subset of $\R$.

Reference: the positive-weight argument in the proof of \cite[Theorem 19.3]{atap}.
\end{theorem}
\begin{proof}
\leanok
\end{proof}

\begin{theorem}[Normalized Gaussian evaluation vectors]
\label{thm:gauss-evaluation-vector}\uses{thm:sequence-finite-coordinates,thm:gauss-vector-input,def:gauss-vector}
\lean{Polynomial.Sequence.IsOrthogonal.gauss_lagrange_repr, Polynomial.Sequence.IsOrthogonal.gaussWeight_mul_sum_eval_sq, Polynomial.Sequence.IsOrthogonal.sum_sq_gaussEigenvector, Polynomial.Sequence.squaredNorm_zero_eq_mass, Polynomial.Sequence.IsOrthogonal.gaussWeight_eq_mass_mul_first_sq}\leanok
Under the same moment, positivity and orthogonality hypotheses, let $x$ be a
Gaussian node of order $n\ge1$, with weight $w_x$. Its cardinal polynomial
has coordinates $w_xP_i(x)/h_i$ for $i<n$. Consequently
\[
 w_x\sum_{i<n}\frac{P_i(x)^2}{h_i}=1,\qquad
 u_i(x)=\sqrt{w_x}\frac{P_i(x)}{\sqrt{h_i}},\qquad
 \sum_{i<n}u_i(x)^2=1.
\]
Moreover $w_x=\mu(\mathbb R)u_0(x)^2$.

References: \cite[3.5(vi)]{dlmf}; \cite{golub-welsch}.
\end{theorem}
\begin{proof}
\leanok
\end{proof}

\begin{theorem}[Eigenvectors determined by their first component]
\label{thm:hessenberg-eigenvector}
\lean{Matrix.eq_zero_of_hessenberg_eigenvector, Matrix.hessenberg_eigenvector_eq_smul, Matrix.hessenberg_eigenvector_first_sq_eq}\leanok
Let $A$ over a field have size $n+1$, vanish at $A_{ij}$ when $j>i+1$, and
have nonzero superdiagonal entries. An eigenvector with first component zero
is zero. Two eigenvectors for the same eigenvalue, one with nonzero first
component, are proportional with ratio determined by those components.
If both squared-coordinate sums equal one, their first-component squares agree.
For real vectors these sums are their Euclidean squared norms.

Reference: \cite{golub-van-loan}, eigenvalue methods.
\end{theorem}
\begin{proof}
\leanok
\end{proof}

\end{supporting}

\part{Operators: the ultraspherical spectral method}\label{part:operators}
\rvsummary{Banded operators on coefficient sequences and the ultraspherical spectral method for linear differential equations.}

\chapter{Coefficients, bases and operators}\label{chap:operators}
\rvsummary{Linear maps on polynomials as matrices acting on coefficient sequences, the language of ApproxFun.}

ApproxFun \cite{approxfun-jl,approxfun} represents operators by infinite banded matrices acting on
coefficient vectors. We first define operators as linear maps on polynomials
and take their matrices in specified source and target bases. On finitely
supported coefficients the action involves only finite sums. Bandwidth
records which entries vanish. Analytic extensions to convergent series or
$\ell^2$ require the additional hypotheses stated later.

\begin{supporting}[Reference details: Polynomial maps and basis matrices]
\begin{definition}[Polynomial maps and basis matrices]
\label{def:polynomial-linear-maps}
\lean{LinearMap.mulLeft, Polynomial.derivative, LinearMap.id, Module.Basis.toMatrix}\mathlibok
Consider the linear maps $p\mapsto Xp$, $p\mapsto p'$, and $p\mapsto p$.
For a target basis $c$ and any family $v$, the basis matrix is
$(c.\mathrm{toMatrix}(v))_{ij}=(c.\mathrm{repr}(v_j))_i$, the $i$th coordinate of $v_j$.
This construction allows infinite index sets.

Reference: \cite[\S2]{olver-townsend}.
\end{definition}
\end{supporting}


\begin{definition}[Operators between polynomial bases]
\label{def:operators}\uses{def:polynomial-linear-maps}
\lean{Module.Basis.coefficientMap, Matrix.IsBanded}\leanok
For source and target bases $b,c$ and a linear map $L$, define the coefficient
map $c.\mathrm{repr}\circ L\circ b.\mathrm{repr}^{-1}$ on finitely supported
coefficients. Its matrix is $A_{ij}=(c.\mathrm{repr}(L(b_j)))_i$.
For natural indices, lower bandwidth $\ell$ and upper bandwidth $u$ mean
$A_{ij}=0$ whenever $j+\ell<i$ or $i+u<j$.
Parameter conversion uses the identity map with different source and target bases.

Reference: \cite[\S2]{olver-townsend}.
\end{definition}

\begin{supporting}[Reference details: Enlarging bandwidths]
\begin{theorem}[Enlarging bandwidths]
\label{thm:bandwidth-mono}\uses{def:operators}
\lean{Matrix.IsBanded.mono}\leanok
If a matrix has bandwidths $(\ell,u)$, $\ell\le\ell'$ and $u\le u'$, then it
has bandwidths $(\ell',u')$.
\end{theorem}
\begin{proof}
\leanok
\end{proof}
\end{supporting}


\begin{theorem}[Exact finite coefficient action]
\label{thm:operator-coefficients}\uses{def:operators}
\lean{Module.Basis.repr_symm_coefficientMap, Module.Basis.coefficientMap_single, Module.Basis.coefficientMap_apply, Module.Basis.coefficientMap_comp}\leanok
The coefficient map sends the unit column at $j$ to $c.\mathrm{repr}(L(b_j))$.
For finitely supported $v$, its $i$th output is
$\sum_{j\in\operatorname{supp}(v)} A_{ij}v_j$; reconstruction in $c$ equals
$L$ applied to reconstruction in $b$. Composition agrees exactly with
composition of the underlying linear maps. No infinite matrix sum is required.

Reference: \cite[\S2]{olver-townsend}.
\end{theorem}
\begin{proof}
\leanok
\end{proof}

\begin{theorem}[Exact rectangular coefficient sections]
\label{thm:finite-operator-sections}\uses{thm:operator-coefficients}
\lean{Module.Basis.coefficientMap_apply_range, Module.Basis.coefficientMap_eq_zero_of_bandwidth}\leanok
Let $L$ be a linear map with natural-indexed source and target bases and matrix
$A$. For an input coefficient vector $v$ supported on indices $j<N$,
\[
 (Lv)_i=\sum_{j=0}^{N-1}A_{ij}v_j
\]
for every output index $i$. If $A$ has bandwidths $(\ell,u)$, all output
coefficients with $i\ge N+\ell$ vanish. Thus $N+\ell$ rows suffice to retain
the entire result; a square truncation is not assumed exact.
Here $Lv$ denotes the coefficient map, and the statements hold over a
commutative ring without finite-dimensional basis assumptions.

Reference: \cite[Section 3]{olver-townsend}.
\end{theorem}
\begin{proof}
\leanok
\end{proof}

Matrix entries depend on the choice and normalisation of each basis.
The following definitions fix the Chebyshev, Jacobi and Gegenbauer coordinates
used below.

\begin{supporting}[Reference details: Chebyshev polynomial coordinate basis and related results]
\begin{definition}[Chebyshev polynomial coordinate basis]
\label{def:chebyshev-polynomial-basis}
\lean{Polynomial.chebyshevBasis}\leanok
Let the Chebyshev coordinate basis of $\R[X]$ be $(T_n)_{n\in\N}$, in the
polynomial normalisation $T_n(1)=1$. It differs from the normalised weighted
Hilbert basis.

Reference: \cite[Sections 2--3]{olver-townsend}.
\end{definition}

\begin{definition}[Jacobi polynomial basis]
\label{def:jacobi-polynomial-basis}\uses{thm:jacobi-sequence}
\lean{Polynomial.jacobiBasis}\leanok
For real $\alpha,\beta>-1$, let the Jacobi basis of $\R[X]$ be the basis built
from the Jacobi sequence, in the original normalisation. This is an algebraic
polynomial basis, distinct from the weighted L2 Hilbert basis.

Reference: \cite[Table 18.3.1]{dlmf}.
\end{definition}

\begin{definition}[Gegenbauer polynomial basis]
\label{def:gegenbauer-polynomial-basis}\uses{thm:gegenbauer-sequence}
\lean{Polynomial.gegenbauerBasis}\leanok
For real $\lambda>-1/2$, $\lambda\ne0$, let the Gegenbauer basis of $\R[X]$ be
the basis built from the sequence $C_n^{(\lambda)}$ in its original
normalisation. Theorem~\ref{thm:coordinate-basis-vectors} identifies the vectors
of all three bases.

Reference: \cite[Table 18.3.1]{dlmf}.
\end{definition}

\begin{theorem}[Basis vectors of the coordinate bases]
\label{thm:coordinate-basis-vectors}\uses{def:chebyshev-polynomial-basis,def:jacobi-polynomial-basis,def:gegenbauer-polynomial-basis,def:chebyshev-mathlib}
\lean{Polynomial.chebyshevBasis_apply, Polynomial.jacobiBasis_apply, Polynomial.gegenbauerBasis_apply}\leanok
For every $n$, the $n$th vectors of the Chebyshev, Jacobi ($\alpha,\beta>-1$)
and Gegenbauer ($\lambda>-1/2$, $\lambda\ne0$) coordinate bases are $T_n$,
$P_n^{(\alpha,\beta)}$ and $C_n^{(\lambda)}$.
\end{theorem}
\begin{proof}
\leanok
\end{proof}

Multiplication by $x$ in an orthogonal basis is the tridiagonal Jacobi operator
$J$ of Chapter~\ref{chap:general}. Multiplication by any polynomial $a$ is then
$a(J)$, and it can be evaluated by running Clenshaw's recurrence in the operator
algebra.

\begin{theorem}[General multiplication matrix]
\label{thm:general-multiplication-matrix}\uses{def:operators,thm:general-coefficient-signs}
\lean{Polynomial.Sequence.IsOrthogonal.mul_matrix_succ, Polynomial.Sequence.IsOrthogonal.mul_matrix_zero, Polynomial.Sequence.IsOrthogonal.mul_matrix_isBanded}\leanok
For an orthogonal polynomial sequence with finite moments and positive definite
integral, the multiplication matrix is tridiagonal. Column $n+1$ has entries
$A_{n+1},B_{n+1},C_n$ in rows $n+2,n+1,n$, with the recurrence coefficients
already defined. Column zero has entries $A_0,B_0$ in rows one and zero.
No interval-support hypothesis is imposed.

Reference: \cite[18.2(iv)]{dlmf}.
\end{theorem}
\begin{proof}
\leanok
\end{proof}
\end{supporting}


\begin{theorem}[Multiplication as polynomial functional calculus]
\label{thm:polynomial-functional-calculus}\uses{thm:operator-coefficients}
\lean{Polynomial.aeval_mulLeft_X, Polynomial.coefficientMap_eq_conj, Polynomial.coefficientMap_mulLeft_eq_aeval, Polynomial.repr_symm_aeval_mulLeft}\leanok
Over any commutative semiring, let $J$ be multiplication by $X$ on polynomials.
For every polynomial $p$, multiplication by $p$ is exactly $p(J)$.
In any polynomial basis $b$, the coefficient maps satisfy the same identity
$M[p]=p(J_b)$, and reconstruct $p\,b^{-1}(v)$ for every finitely supported
coefficient vector $v$. Here $p(J)$ is the evaluation of $p$ at $J$ in the
algebra of linear endomorphisms.


Reference: \cite[Sections 2--3]{olver-townsend}.
\end{theorem}
\begin{proof}
\leanok
\end{proof}

\rvnote{thm:polynomial-functional-calculus}{No finite-section matrix functional calculus is asserted.}


\begin{theorem}[Operator Clenshaw]
\label{thm:operator-clenshaw}\uses{thm:polynomial-functional-calculus,def:chebyshev-mathlib,thm:chebyshev-recurrence,thm:clenshaw}
\lean{Polynomial.clenshawState_map, Polynomial.clenshaw_map, Polynomial.Chebyshev.clenshaw_polynomial_T, Polynomial.Chebyshev.clenshaw_aeval_T, Polynomial.Chebyshev.clenshaw_mulLeft_T, Polynomial.Chebyshev.clenshaw_coefficientMap_T}\leanok
Clenshaw state and evaluation commute with ring homomorphisms; the recurrence
itself is defined for rings without requiring commutativity.
Let $R$ be commutative, $A$ any $R$-algebra, and $J\in A$. Apply the backward
recurrence with $a_k=2J$, $b_k=-1$, scalar coefficients $c_k1_A$ and initial
values $1_A,J$. The result is
\[
 \left(\sum_{k=0}^{n-1}c_kT_k\right)(J).
\]
In particular it computes the exact multiplication operator, or its coefficient
map in any polynomial basis. The full constant coefficient is retained.


References: \cite{clenshaw}; \cite[Sections 2--3]{olver-townsend}.
\end{theorem}
\begin{proof}
\leanok
\end{proof}

\rvnote{thm:operator-clenshaw}{
No commutativity of the entire endomorphism algebra is required: the proof transports a polynomial identity by evaluation at a single $J$.
}

For convergent expansions the same matrices act on infinite coefficient
sequences, provided the reconstruction converges absolutely.

\begin{supporting}[Reference details: Coefficient maps on infinite sequences and related results]
\begin{definition}[Coefficient maps on infinite sequences]
\label{def:series-coefficient-map}
\lean{Module.Basis.SeriesCoefficientInput, Module.Basis.seriesCoefficientMap}\leanok
Let $b,c$ be polynomial bases over a commutative ring $R$, let
$L:R[x]\to R[x]$ be $R$-linear, and let $\sigma:R\to\mathbb C$ be a ring map.
Write $M_{ij}=[Lb_j]_{c,i}$. At index $i$, an admissible input is a complex
sequence $a$ for which $\sum_j a_j\sigma(M_{ij})$ converges unconditionally.
Define $\beta_i$ only for such an input, as this sum.

Reference: \cite[Sections 2--3]{olver-townsend}.
\end{definition}

\begin{theorem}[Classical matrix-row sum]
\label{thm:series-coefficient-map-value}\uses{def:series-coefficient-map}
\lean{Module.Basis.hasSum_seriesCoefficientMap_row}\leanok
For every admissible input at row $i$, the series $\sum_j a_j\sigma(M_{ij})$
converges unconditionally to $\beta_i$.
\end{theorem}
\begin{proof}
\leanok
\end{proof}
\end{supporting}



\begin{theorem}[Absolutely convergent matrix reconstruction]
\label{thm:series-matrix-reconstruction}\uses{def:series-coefficient-map,thm:regrouped-coefficient-input,def:regrouped-coefficients}
\lean{Module.Basis.hasSum_repr_eval₂, Module.Basis.hasSum_seriesCoefficientMap, Module.Basis.hasSum_seriesCoefficientMap_of_rows, Module.Basis.summable_norm_seriesCoefficientMap_terms}\leanok
At a complex point $z$, suppose
$\sum_{j,i}|a_j\sigma(M_{ij})c_i(z)|<\infty$, with evaluation through $\sigma$.
The convergent weighted contributions
$\gamma_i=\sum_j a_j\sigma(M_{ij})c_i(z)$ then satisfy
\[
 \sum_i\gamma_i=\sum_j a_j(Lb_j)(z).
\]
If $a$ also supplies an admissible coefficient input at every row, then
$\gamma_i=\beta_i c_i(z)$ and the left side is $\sum_i\beta_i c_i(z)$.
A sufficient condition for weighted reconstruction is $|c_i(z)|\le w_i$ and
\[
 \sum_j |a_j|\sum_{i\in\operatorname{supp}[Lb_j]_c}|\sigma(M_{ij})|w_i<\infty.
\]
Each inner sum is finite.
\end{theorem}
\begin{proof}
\leanok
\end{proof}

\rvnote{thm:series-matrix-reconstruction}{Weighted convergence alone does not assert convergence of a coefficient row whose output mode vanishes at $z$.}


\begin{supporting}[Supporting results: finite rows of banded series maps]
\begin{theorem}[Finite rows and compatibility with finite coefficients]
\label{thm:banded-series-coefficient-map}\uses{def:series-coefficient-map,thm:operator-coefficients}
\lean{Module.Basis.SeriesCoefficientInput.ofBanded, Module.Basis.SeriesCoefficientInput.ofFinsupp, Module.Basis.seriesCoefficientMap_eq_sum_of_isBanded, Module.Basis.seriesCoefficientMap_finsupp}\leanok
If $M$ has bandwidths $(\ell,u)$, each row sum $\sum_j a_j\sigma(M_{ij})$ is finite for every sequence $a$, and
$\beta_i=\sum_{j<i+u+1}a_j\sigma(M_{ij})$ for every sequence $a$.
On finitely supported sequences embedded from $R$, this agrees exactly with the
existing finite coefficient map. Neither coordinate identity needs an
analytic convergence hypothesis.
\end{theorem}
\begin{proof}
\leanok
\end{proof}

\end{supporting}

We next compute these matrices in the classical bases, including their
nonzero bands.

\chapter{Banded multiplication, differentiation and conversion}\label{chap:banded}
\rvsummary{The banded multiplication, differentiation and conversion matrices of the Jacobi and ultraspherical bases.}

In Jacobi and ultraspherical bases, multiplication by $x$ is tridiagonal.
Differentiation raises the parameters and has a single superdiagonal.
One-step conversion has bandwidth $(0,1)$ for Jacobi and $(0,2)$ for
ultraspherical bases. These sparse matrices are those of the ultraspherical spectral
method \cite{olver-townsend}, used by ApproxFun \cite{approxfun-jl}. The
identities with denominators cleared hold over every commutative
$\Q$-algebra. Divided entries and bandwidths require nonvanishing denominators;
these conditions hold in the real orthogonality ranges.

\begin{theorem}[Jacobi multiplication matrix]
\label{thm:jacobi-multiplication-matrix}\uses{def:operators,thm:jacobi-recurrence,thm:coordinate-basis-vectors}
\lean{Polynomial.jacobi_mul_matrix_succ, Polynomial.jacobi_mul_matrix_zero, Polynomial.jacobi_mul_matrix_succ_div, Polynomial.jacobi_mul_matrix_zero_div, Polynomial.jacobi_mul_matrix_isBanded, Polynomial.jacobiBasis_mul_isBanded}\leanok
In a Jacobi target basis, write $s=\alpha+\beta$. For input degree $n\ge1$,
only rows $n+1,n,n-1$ can be nonzero, with respective entries
\[\frac{2(n+1)(n+s+1)}{(2n+s+1)(2n+s+2)},\qquad
\frac{\beta^2-\alpha^2}{(2n+s)(2n+s+2)},\qquad
\frac{2(n+\alpha)(n+\beta)}{(2n+s)(2n+s+1)}.\]
Column zero has entries $2/(s+2)$ in row one and $(\beta-\alpha)/(s+2)$ in row zero.
The divided formulas hold over fields with a rational algebra structure when
the displayed denominator factors are nonzero; cleared identities hold over
commutative rational algebras. For real $\alpha,\beta>-1$ all factors are
nonzero and the matrix has bandwidths $(1,1)$.

Reference: \cite[18.9.1--18.9.2]{dlmf}.
\end{theorem}
\begin{proof}
\leanok
\end{proof}

\begin{theorem}[Classical differentiation matrices]
\label{thm:classical-derivative-matrices}\uses{def:operators,thm:coordinate-basis-vectors}
\lean{Polynomial.jacobi_derivative_matrix_succ, Polynomial.jacobi_derivative_matrix_zero, Polynomial.jacobi_derivative_matrix_isBanded, Polynomial.jacobiBasis_derivative_isBanded, Polynomial.gegenbauer_derivative_matrix_succ, Polynomial.gegenbauer_derivative_matrix_zero, Polynomial.gegenbauer_derivative_matrix_isBanded, Polynomial.gegenbauerBasis_derivative_isBanded}\leanok
Jacobi differentiation from $(\alpha,\beta)$ to $(\alpha+1,\beta+1)$ has
entry $(n+\alpha+\beta+2)/2$ in row $n$, column $n+1$.
Gegenbauer differentiation from $\lambda$ to $\lambda+1$ has entry $2\lambda$
in row $n$, column $n+1$. All other entries vanish, including column zero.
These are identities over commutative rational algebras whenever the specified
target bases exist; both matrices have bandwidths $(0,1)$.
The concrete real bases above supply them throughout their orthogonality ranges.

References: \cite[18.9.15, 18.9.19]{dlmf}; \cite{olver-townsend}.
\end{theorem}
\begin{proof}
\leanok
\end{proof}

\begin{theorem}[Jacobi parameter-conversion matrices]
\label{thm:jacobi-conversion-matrices}\uses{def:operators,thm:jacobi-parameter-raising,thm:coordinate-basis-vectors}
\lean{Polynomial.jacobi_raise_beta_matrix_succ, Polynomial.jacobi_raise_alpha_matrix_succ, Polynomial.jacobi_conversion_matrix_zero, Polynomial.jacobi_raise_beta_matrix_succ_div, Polynomial.jacobi_raise_alpha_matrix_succ_div, Polynomial.jacobi_raise_beta_matrix_isBanded, Polynomial.jacobi_raise_alpha_matrix_isBanded, Polynomial.jacobiBasis_raise_alpha_isBanded, Polynomial.jacobiBasis_raise_beta_isBanded}\leanok
Let $s=\alpha+\beta$. For input degree $n\ge1$, conversion raising $\alpha$
has entries $(n+s+1)/(2n+s+1)$ in row $n$ and
$-(n+\beta)/(2n+s+1)$ in row $n-1$.
Conversion raising $\beta$ has the same diagonal entry and
$(n+\alpha)/(2n+s+1)$ in row $n-1$.
Column zero is the unit column at zero. All other entries vanish.
The divided formulas require $2n+s+1\ne0$ over a field; the cleared identities
hold over every commutative rational algebra with the specified target basis.
For real $\alpha,\beta>-1$ both matrices have bandwidths $(0,1)$ with no
additional denominator hypotheses.

Reference: \cite[18.9.5]{dlmf} and its reflection.
\end{theorem}
\begin{proof}
\leanok
\end{proof}

\begin{supporting}[Reference details: Cleared Jacobi derivative lowering]
\begin{theorem}[Cleared Jacobi derivative lowering]
\label{thm:jacobi-derivative-lowering}\uses{thm:jacobi-ode-unique,thm:jacobi-ode}
\lean{Polynomial.jacobi_derivative_lowering}\leanok
For $n\ge0$ over every commutative $\Q$-algebra, put $s=\alpha+\beta$. Then
\begin{align*}
 (2n+s+2)(1-X^2)(P_{n+1}^{(\alpha,\beta)})'
 &= (n+1)[\alpha-\beta-(2n+s+2)X]P_{n+1}^{(\alpha,\beta)}\\
 &\quad+2(n+\alpha+1)(n+\beta+1)P_n^{(\alpha,\beta)}.
\end{align*}
This includes degenerate parameters, without division by parameter factors.

Reference: \cite[18.9.17]{dlmf}, with its degree shifted by one.
\end{theorem}
\begin{proof}
\leanok
\end{proof}
\end{supporting}


\begin{theorem}[Gegenbauer multiplication matrix]
\label{thm:gegenbauer-multiplication-matrix}\uses{def:operators,thm:coordinate-basis-vectors}
\lean{Polynomial.gegenbauer_mul_matrix_succ, Polynomial.gegenbauer_mul_matrix_zero, Polynomial.gegenbauer_mul_matrix_succ_div, Polynomial.gegenbauer_mul_matrix_zero_div, Polynomial.gegenbauer_mul_matrix_isBanded, Polynomial.gegenbauerBasis_mul_isBanded}\leanok
In a Gegenbauer basis, column $n\ge1$ of multiplication has entries
$(n+1)/(2(n+\lambda))$ in row $n+1$ and
$(n+2\lambda-1)/(2(n+\lambda))$ in row $n-1$; the diagonal is zero.
Column zero has entry $1/(2\lambda)$ in row one.
Over a field the divided formulas require their denominators nonzero;
cleared identities hold over commutative rational algebras.
For real $\lambda>-1/2$, $\lambda\ne0$, the matrix has bandwidths $(1,1)$.

Reference: \cite[18.9.1, Table 18.9.1]{dlmf}.
\end{theorem}
\begin{proof}
\leanok
\end{proof}

\begin{theorem}[Gegenbauer conversion matrix]
\label{thm:gegenbauer-conversion-matrix}\uses{def:operators,thm:coordinate-basis-vectors}
\lean{Polynomial.gegenbauer_conversion_matrix_zero, Polynomial.gegenbauer_conversion_matrix_one, Polynomial.gegenbauer_conversion_matrix_add_two, Polynomial.gegenbauer_conversion_matrix_one_div, Polynomial.gegenbauer_conversion_matrix_add_two_div, Polynomial.gegenbauer_conversion_matrix_isBanded, Polynomial.gegenbauerBasis_conversion_isBanded}\leanok
Conversion from parameter $\lambda$ to $\lambda+1$ has entries
$\lambda/(n+\lambda)$ in row $n$ and $-\lambda/(n+\lambda)$ in row $n-2$
of column $n\ge2$. Column one has entry $\lambda/(\lambda+1)$ in row one;
column zero is the unit column at zero. All other entries vanish.
Over a field the divided formulas require their denominators nonzero;
cleared identities hold over commutative rational algebras.
For real $\lambda>-1/2$ the target basis exists and the matrix has bandwidths
$(0,2)$, also at the degenerate source parameter zero.

References: \cite[18.9.7]{dlmf}; \cite{olver-townsend}.
\end{theorem}
\begin{proof}
\leanok
\end{proof}

\paragraph{Example (Julia).} The conversion matrix of
Theorem~\ref{thm:gegenbauer-conversion-matrix} for $\lambda=0.7$ reproduces
each $C_n^{(\lambda)}(x)$ from values of $C^{(\lambda+1)}$, both computed by the
recurrence of Theorem~\ref{thm:gegenbauer-recurrence}. Its 40 nonzeros lie on
the diagonal and the second superdiagonal, the bandwidths $(0,2)$.
\begin{verbatim}
function gegenbauer(n, lam, x)      # three-term recurrence
    n == 0 && return 1.0
    a, b = 1.0, 2lam*x
    for k in 1:n-1
        a, b = b, (2(k + lam)*x*b - (k + 2lam - 1)*a)/(k + 1)
    end
    b
end
N, lam = 20, 0.7
S = zeros(N+1, N+1)                 # conversion C^(lam) -> C^(lam+1)
S[1, 1] = 1
S[2, 2] = lam/(lam + 1)
for n in 2:N
    S[n+1, n+1] = lam/(n + lam)
    S[n-1, n+1] = -lam/(n + lam)
end
x = 0.37
maximum(abs(gegenbauer(n, lam, x) -
    sum(S[i+1, n+1]*gegenbauer(i, lam + 1, x) for i in 0:N)) for n in 0:N)
                                    # about 1e-16
band = all(S[i, j] == 0 for i in 1:N+1, j in 1:N+1 if !(0 <= j - i <= 2))
(count(!iszero, S), band)           # (40, true)
\end{verbatim}

\begin{theorem}[Polynomial multiplication bandwidth in Gegenbauer bases]
\label{thm:gegenbauer-polynomial-multiplication}\uses{thm:univariate-modification-operators,def:gegenbauer-polynomial-basis,thm:gegenbauer-orthogonal}
\lean{Polynomial.gegenbauerBasis_polynomial_mul_isBanded}\leanok
For real Gegenbauer bases with $\lambda>-1/2$, $\lambda\ne0$, a polynomial
multiplier $p$ has bandwidths $(d,d)$, where $d=\operatorname{natDegree}p$.

Reference: \cite[Sections 2--3]{olver-townsend}.
\end{theorem}
\begin{proof}
\leanok
\end{proof}

\begin{supporting}[Reference details: Finite Gegenbauer recurrence domain and related results]
\begin{definition}[Finite Gegenbauer recurrence domain]
\label{def:gegenbauer-clenshaw-input}
\lean{Polynomial.GegenbauerClenshawInput}\leanok
For a field $K$ and a coefficient count $n$, the Gegenbauer Clenshaw domain requires $(k+2)\ne0$ in $K$ for every $0\le k<n-2$.
There is no condition when $n\le2$.
\end{definition}

\begin{theorem}[Rational-algebra Gegenbauer inputs]
\label{thm:gegenbauer-clenshaw-input}\uses{def:gegenbauer-clenshaw-input}
\lean{Polynomial.GegenbauerClenshawInput.ofRatAlgebra}\leanok
Over a field with a rational algebra structure, $k+2\ne0$ for every
$k\in\N$. Thus every coefficient count $n$ satisfies the Gegenbauer Clenshaw
denominator conditions of Definition~\ref{def:gegenbauer-clenshaw-input}.
\end{theorem}
\begin{proof}
\leanok
\end{proof}

\begin{definition}[Finite Gegenbauer Clenshaw evaluator]
\label{def:gegenbauer-clenshaw}\uses{def:finite-clenshaw,def:gegenbauer-clenshaw-input}
\lean{Polynomial.clenshawGegenbauer}\leanok
For a field $K$, parameters $\lambda,x\in K$, $n$ coefficients $c_0,\ldots,c_{n-1}$, and the denominator conditions of Definition~\ref{def:gegenbauer-clenshaw-input}, use finite Clenshaw evaluation with
\[
 a_k=\frac{2(k+\lambda+1)x}{k+2},\quad
 b_k=-\frac{k+2\lambda}{k+2}\quad(0\le k<n-2),\qquad
 p_0=1,\quad p_1=2\lambda x.
\]
References: \cite{clenshaw}; \cite[Table 18.9.1]{dlmf}.
\end{definition}

\begin{theorem}[Gegenbauer Clenshaw evaluation]
\label{thm:gegenbauer-clenshaw}\uses{thm:clenshaw,thm:gegenbauer-recurrence,thm:finite-clenshaw,def:gegenbauer-clenshaw,thm:gegenbauer-clenshaw-input}
\lean{Polynomial.clenshawGegenbauer_eq_sum}\leanok
Over a field with a rational algebra structure, for every parameter $\lambda$,
use
\[
 a_k=\frac{2(k+\lambda+1)x}{k+2},\qquad
 b_k=-\frac{k+2\lambda}{k+2},\qquad p_0=1,\quad p_1=2\lambda x.
\]
The evaluator returns $\sum_{i=0}^{n-1}c_iC_i^{(\lambda)}(x)$.
Only positive natural denominators occur, so this includes $\lambda=0$.

References: \cite{clenshaw}; \cite[Table 18.9.1]{dlmf}.
\end{theorem}
\begin{proof}
\leanok
\end{proof}

\begin{definition}[Finite Jacobi recurrence domain]
\label{def:jacobi-clenshaw-input}
\lean{Polynomial.JacobiClenshawInput}\leanok
For a field $K$, parameters $\alpha,\beta\in K$ and a coefficient count $n$,
the Jacobi Clenshaw domain requires $2\ne0$ when $n\ge2$, and
\[
 d_k=2(k+2)(k+\alpha+\beta+2)(2k+\alpha+\beta+2)\ne0
 \qquad(0\le k<n-2).
\]
Empty and constant expansions impose no denominator conditions.
\end{definition}

\begin{theorem}[Constructing finite Jacobi inputs]
\label{thm:jacobi-clenshaw-input}\uses{def:jacobi-clenshaw-input}
\lean{Polynomial.JacobiClenshawInput.ofDenominators, Polynomial.JacobiClenshawInput.of_neg_one_lt}\leanok
In a field with a rational algebra structure, if $d_k\ne0$ for every $k\ge0$,
the denominator conditions of Definition~\ref{def:jacobi-clenshaw-input} hold for every coefficient count $n$.
For real $\alpha,\beta>-1$, these denominator conditions hold without further assumptions.
\end{theorem}
\begin{proof}
\leanok
\end{proof}

\begin{definition}[Finite Jacobi Clenshaw evaluator]
\label{def:jacobi-clenshaw}\uses{def:finite-clenshaw,def:jacobi-clenshaw-input}
\lean{Polynomial.clenshawJacobi}\leanok
For a field $K$, parameters $\alpha,\beta,x\in K$, $n$ coefficients $c_0,\ldots,c_{n-1}$, and the denominator conditions of Definition~\ref{def:jacobi-clenshaw-input}, return $0$ if $n=0$ and $c_0$ if $n=1$.
For $n\ge2$, use finite Clenshaw evaluation with
\[
 a_k=\frac{(2k+\alpha+\beta+3)
 ((2k+\alpha+\beta+2)(2k+\alpha+\beta+4)x+\alpha^2-\beta^2)}{d_k},
 \qquad
 b_k=-\frac{2(k+\alpha+1)(k+\beta+1)(2k+\alpha+\beta+4)}{d_k}
 \quad(0\le k<n-2),
\]
where $d_k=2(k+2)(k+\alpha+\beta+2)(2k+\alpha+\beta+2)$, and
$p_0=1$, $p_1=\alpha+1+\tfrac12(\alpha+\beta+2)(x-1)$.
References: \cite{clenshaw}; \cite[18.9.1--2]{dlmf}.
\end{definition}

\begin{theorem}[Jacobi Clenshaw evaluation]
\label{thm:jacobi-clenshaw}\uses{thm:clenshaw,thm:jacobi-recurrence,thm:finite-clenshaw,def:jacobi-clenshaw,thm:jacobi-clenshaw-input}
\lean{Polynomial.clenshawJacobi_eq_sum, Polynomial.clenshawJacobi_eq_sum_of_neg_one_lt}\leanok
Over a field with a rational algebra structure, write
\[
 d_k=2(k+2)(k+\alpha+\beta+2)(2k+\alpha+\beta+2).
\]
Assume $d_k\ne0$ for every $k\in\N$. Use
\[
\begin{aligned}
 a_k&=\frac{(2k+\alpha+\beta+3)
 ((2k+\alpha+\beta+2)(2k+\alpha+\beta+4)x+\alpha^2-\beta^2)}{d_k},\\
 b_k&=-\frac{2(k+\alpha+1)(k+\beta+1)(2k+\alpha+\beta+4)}{d_k}.
\end{aligned}
\]
with $p_0=1$ and $p_1=\alpha+1+\tfrac12(\alpha+\beta+2)(x-1)$.
The evaluator returns $\sum_{i=0}^{n-1}c_iP_i^{(\alpha,\beta)}(x)$.
For real $\alpha,\beta>-1$, all $d_k$ are positive and no extra hypotheses
are needed. In particular, $\alpha+\beta=-1$ is included.

References: \cite{clenshaw}; \cite[18.9.1--2]{dlmf}.
\end{theorem}
\begin{proof}
\leanok
\end{proof}
\end{supporting}


Chebyshev expansions enter the ultraspherical bases through $C^{(1)}=U$.
Conversion $T\to C^{(1)}$ has bandwidth $(0,2)$, and the $k$th derivative maps
$T$ to $C^{(k)}$ with a single nonzero diagonal $k$ places above the main one.
Multiplication in the $T$ basis has a Toeplitz-plus-Hankel form. The entry
formula below includes the adjustments required by the full constant term.

\begin{theorem}[Entry into the ultraspherical bases]
\label{thm:chebyshev-gegenbauer-entry}\uses{thm:gegenbauer-recurrence,def:chebyshev-mathlib}
\lean{Polynomial.gegenbauer_one_eq_chebyshev_U, Polynomial.two_mul_chebyshev_T_eq_gegenbauer_sub, Polynomial.derivative_chebyshev_T_succ, Polynomial.iterate_derivative_chebyshev_T, Polynomial.iterate_derivative_chebyshev_T_eq_zero}\leanok
Over a commutative rational algebra, $C_n^{(1)}=U_n$ and
$2T_{n+2}=C_{n+2}^{(1)}-C_n^{(1)}$. For $k,n\ge0$,
\[
 D^{k+1}T_{n+k+1}=2^k k!(n+k+1)C_n^{(k+1)}.
\]
Over a characteristic-zero field, $D^kT_n=0$ when $n<k$.
These identities retain the full constant coefficient convention.

References: \cite[18.7.4 and 18.9.9]{dlmf}; \cite[\S2]{olver-townsend}.
\end{theorem}
\begin{proof}
\leanok
\end{proof}

\begin{supporting}[Reference details: Chebyshev representations in $1/x^2$]
\begin{theorem}[Chebyshev representations in $1/x^2$]
\label{thm:chebyshev-hypergeometric-inverse-square}\uses{thm:chebyshev-gegenbauer-entry,thm:gegenbauer-hypergeometric,thm:gegenbauer-reduced-chebyshev}
\lean{Polynomial.eval_chebyshev_U_eq_ordinaryHypergeometric_inv_sq, Polynomial.eval_chebyshev_T_eq_ordinaryHypergeometric_inv_sq}\leanok
For real or complex $x\ne0$ and $n\in\N$,
\[U_n(x)=(2x)^n\,{}_2F_1\!\left(-\tfrac n2,\tfrac{1-n}2;-n;\frac1{x^2}\right),\]
and for $n\ge1$,
\[T_n(x)=2^{n-1}x^n\,{}_2F_1\!\left(-\tfrac n2,\tfrac{1-n}2;1-n;\frac1{x^2}\right),\]
where ${}_2F_1$ is the Gauss series of Definition~\ref{def:ordinary-hypergeometric}.

Reference: \cite[18.5.11\_1 and 18.5.11\_3]{dlmf}.
\end{theorem}
\begin{proof}
\leanok
\end{proof}
\end{supporting}


\begin{theorem}[Chebyshev-to-ultraspherical conversion matrix]
\label{thm:chebyshev-ultraspherical-conversion-matrix}\uses{thm:chebyshev-gegenbauer-entry,def:operators,thm:chebyshev-recurrence,thm:coordinate-basis-vectors}
\lean{Polynomial.chebyshev_gegenbauer_conversion_zero, Polynomial.chebyshev_gegenbauer_conversion_one, Polynomial.chebyshev_gegenbauer_conversion_add_two, Polynomial.chebyshev_gegenbauer_conversion_isBanded, Polynomial.chebyshevBasis_conversion_isBanded}\leanok
Over a field with a rational algebra structure, take a target polynomial basis
$c_i=C_i^{(1)}$. With rows indexing output coefficients and columns input
coefficients, conversion $S$ from $T$ has columns
\[
 S_{i0}=\delta_{i0},\qquad S_{i1}=\tfrac12\delta_{i1},\qquad
 S_{i,n+2}=\tfrac12(\delta_{i,n+2}-\delta_{in}).
\]
Thus its bandwidths are $(0,2)$. The constant column is not halved.
The result is also instantiated on the concrete real polynomial bases.

Reference: \cite[equation (2.8)]{olver-townsend}.
\end{theorem}
\begin{proof}
\leanok
\end{proof}

\begin{theorem}[Higher Chebyshev derivative matrices]
\label{thm:chebyshev-ultraspherical-derivative-matrix}\uses{thm:chebyshev-gegenbauer-entry,def:operators,thm:coordinate-basis-vectors}
\lean{Polynomial.chebyshev_iterate_derivative_matrix, Polynomial.chebyshev_iterate_derivative_matrix_zero, Polynomial.chebyshev_iterate_derivative_isBanded, Polynomial.chebyshevBasis_iterate_derivative_isBanded}\leanok
For $k\ge0$, differentiation of order $k+1$ from $T$ into a target basis
$C^{(k+1)}$ has entries
\[
 (D_{k+1})_{i,n+k+1}=2^k k!(n+k+1)\delta_{in},
\]
with columns below $k+1$ zero. Its bandwidths are $(0,k+1)$.
These identities hold over a field with a rational algebra structure and are
instantiated on the real polynomial bases.

Reference: \cite[Section 3]{olver-townsend}.
\end{theorem}
\begin{proof}
\leanok
\end{proof}

\begin{theorem}[Chebyshev multiplication entries]
\label{thm:chebyshev-multiplication-entries}\uses{thm:coordinate-basis-vectors,thm:chebyshev-recurrence}
\lean{Polynomial.two_mul_chebyshevBasis_repr_mul, Polynomial.chebyshev_mul_matrix_zero_column, Polynomial.chebyshev_mul_matrix_zero_row, Polynomial.chebyshev_mul_matrix_positive}\leanok
Write $p=\sum_m a_mT_m$ with a full constant coefficient, and let $M$ represent
multiplication by $p$ in the $T$ basis. Then
\[
 M_{i0}=a_i,\qquad M_{0j}=\tfrac12a_j\quad(j>0),\qquad
 M_{ij}=\frac{(1+\delta_{ij})a_{|i-j|}+a_{i+j}}2\quad(i,j>0).
\]
In particular the positive diagonal is $a_0+a_{2i}/2$.
These are the Toeplitz-plus-Hankel entries with the exceptional row, column and
constant coefficient accounted for.

References: \cite[Section 2.2, equation (2.7)]{olver-townsend}; \cite[18.18.21]{dlmf}.
\end{theorem}
\begin{proof}
\leanok
\end{proof}

\begin{definition}[Integer ultraspherical basis ladder]
\label{def:ultraspherical-ladder}\uses{def:chebyshev-polynomial-basis,def:gegenbauer-polynomial-basis}
\lean{Polynomial.ultrasphericalBasis}\leanok
Let $B_0$ be the Chebyshev coordinate basis and, for integers $k\ge1$, let $B_k$
be the Gegenbauer basis with $\lambda=k$.

Reference: \cite[Section 3]{olver-townsend}.
\end{definition}

\begin{supporting}[Reference details: Rungs of the ultraspherical ladder]
\begin{theorem}[Rungs of the ultraspherical ladder]
\label{thm:ultraspherical-ladder-vectors}\uses{def:ultraspherical-ladder,thm:coordinate-basis-vectors}
\lean{Polynomial.ultrasphericalBasis_zero, Polynomial.ultrasphericalBasis_succ_apply}\leanok
$B_0$ is the Chebyshev coordinate basis, and for $k,j\ge0$ the $j$th vector of
$B_{k+1}$ is $C_j^{(k+1)}$.
\end{theorem}
\begin{proof}
\leanok
\end{proof}
\end{supporting}


\begin{theorem}[Common-target ultraspherical conversion and differentiation]
\label{thm:ultraspherical-ladder}\uses{thm:chebyshev-ultraspherical-conversion-matrix,thm:chebyshev-ultraspherical-derivative-matrix,thm:gegenbauer-conversion-matrix,thm:conversion-composition-support,thm:ultraspherical-ladder-vectors}
\lean{Polynomial.ultrasphericalBasis_step_isBanded, Polynomial.ultrasphericalBasis_raise_isBanded, Polynomial.ultrasphericalBasis_conversion_isBanded, Polynomial.ultrasphericalBasis_derivative_isBanded, Polynomial.ultrasphericalBasis_derivative_to_level_isBanded}\leanok
Conversion $B_k\to B_{k+r}$ has bandwidths $(0,2r)$.
Differentiation $D^k$ from $B_0$ to $B_k$ has bandwidths $(0,k)$, including $k=0$.
Consequently, for $0\le k\le N$, differentiation represented from $B_0$ to
$B_N$ has bandwidths $(0,2N-k)$.

Reference: \cite[Section 3]{olver-townsend}.
\end{theorem}
\begin{proof}
\leanok
\end{proof}

\begin{theorem}[Convergent ultraspherical basis conversion]
\label{thm:ultraspherical-series-conversion}\uses{thm:series-matrix-reconstruction,thm:banded-series-coefficient-map,thm:ultraspherical-ladder}
\lean{Polynomial.ultraspherical_series_conversion_eq_sum, Polynomial.hasSum_ultraspherical_series_conversion, Polynomial.hasSum_ultraspherical_series_conversion_of_rows, Polynomial.hasSum_ultraspherical_series_conversion_of_le}\leanok
Use the existing real bases $B_0=T$, $B_k=C^{(k)}$ for $k\ge1$, evaluated
at complex arguments. For $k\le N$, conversion from $B_k$ to $B_N$ computes
coordinate $i$ from indices $j<i+2(N-k)+1$. Under the weighted summability
condition of matrix reconstruction, it preserves the represented function.
In either direction, the weighted summability condition gives reconstruction
by the contributions $\gamma_i$ of Theorem~\ref{thm:series-matrix-reconstruction}.
For reverse conversion, these may be factored into coefficients and modes
when each coefficient row converges. Forward conversion has finite rows,
so this additional row convergence is automatic.
\end{theorem}
\begin{proof}
\leanok
\end{proof}

Banded classical operators: see Theorems~\ref{thm:operator-coefficients}, \ref{thm:jacobi-multiplication-matrix}, \ref{thm:classical-derivative-matrices}, \ref{thm:jacobi-conversion-matrices}, \ref{thm:gegenbauer-multiplication-matrix} and~\ref{thm:gegenbauer-conversion-matrix}.


A differential equation combines these maps. Each differentiated term first
lands in its own basis; conversion brings the terms to a common target so
that their coefficient matrices can be added.

\chapter{Linear differential equations as almost-banded systems}\label{chap:odes}
\rvsummary{Linear differential equations with polynomial coefficients as exact almost-banded systems in ultraspherical bases.}

To solve $\sum_k a_k(x)u^{(k)}=f$ with boundary conditions, the ultraspherical spectral
method \cite{olver-townsend}, as implemented in ApproxFun \cite{approxfun-jl}, represents
$u$ in $T$ and the residual in $C^{(N)}$. Each term $a_kD^k$ becomes
multiplication composed with differentiation and conversion, so the residual
matrix is banded with bandwidths fixed by $\deg a_k$ and $N$. Boundary
conditions add a few dense rows. This chapter proves the exact algebra of that
assembly and the truncation errors for analytic data. Solver convergence is
separate: Chapter~\ref{chap:solvers} treats first-order problems with polynomial
coefficients, assuming invertibility of the infinite preconditioned operator.

\begin{definition}[Polynomial differential operator]
\label{def:polynomial-differential-operator}\uses{def:polynomial-linear-maps}
\lean{Polynomial.differentialOperator}\leanok
For polynomial coefficients $a_k$ over a commutative ring, assemble the
endomorphism $L=\sum_{k=0}^N M[a_k]D^k$ on the polynomial module, where $D$ is
polynomial differentiation.

Reference: \cite[Section 3]{olver-townsend}.
\end{definition}

\begin{theorem}[Exact polynomial differential assembly]
\label{thm:polynomial-differential-operator}\uses{def:polynomial-differential-operator,thm:operator-coefficients}
\lean{Polynomial.differentialOperator_apply, Polynomial.differentialOperator_one_apply, Polynomial.repr_symm_differentialOperator, Polynomial.differentialOperator_isBanded}\leanok
The assembled operator satisfies $Lp=\sum_{k=0}^Na_kD^kp$ for every polynomial
$p$. Its coefficient map between any source and target polynomial bases
reconstructs that exact residual. If each summand has bandwidths bounded by
$(\ell,u)$ in those common bases, then their sum has bandwidths $(\ell,u)$.
This is an algebraic assembly theorem with explicit term-bandwidth hypotheses,
not a general all-order ultraspherical bandwidth or ODE solvability theorem.

Reference: \cite[Section 3]{olver-townsend}.
\end{theorem}
\begin{proof}
\leanok
\end{proof}

\begin{theorem}[All-order ultraspherical residual bandwidth]
\label{thm:all-order-ultraspherical}\uses{thm:ultraspherical-ladder,thm:gegenbauer-polynomial-multiplication,thm:polynomial-differential-operator,thm:bandwidth-mono,thm:finite-operator-sections}
\lean{Polynomial.ultrasphericalBasis_mul_isBanded, Polynomial.ultrasphericalBasis_differential_term_isBanded, Polynomial.ultraspherical_differentialOperator_isBanded, Polynomial.ultraspherical_differentialOperator_isBanded_of_degree_le, Polynomial.ultraspherical_differentialOperator_coeff_eq_zero}\leanok
Let $N\ge1$ and $L=\sum_{k=0}^N a_kD^k$, with real polynomial coefficients and
$d_k=\operatorname{natDegree}a_k$. In Chebyshev input and $B_N$ output,
each term has bandwidths $(d_k,d_k+2N-k)$. Hence $L$ has bandwidths $(\ell,u)$
whenever $d_k\le\ell$ and $d_k+2N-k\le u$ for every $k\le N$.
A common degree bound $d$ gives $(d,d+2N)$. For input supported on the first $m$
coefficients, all residual coefficients at indices $i\ge m+d$ vanish.
These are exact polynomial statements.

Reference: \cite[Section 3]{olver-townsend}.
\end{theorem}
\begin{proof}
\leanok
\end{proof}

\begin{theorem}[Chebyshev endpoint boundary rows]
\label{thm:chebyshev-boundary-rows}\uses{thm:T-coefficient-calculus,thm:coordinate-basis-vectors,thm:chebyshev-ode}
\lean{Polynomial.chebyshevBasis_functional_apply, Polynomial.chebyshevBasis_eval_one, Polynomial.chebyshevBasis_eval_neg_one, Polynomial.chebyshevBasis_derivative_eval_one, Polynomial.chebyshevBasis_derivative_eval_neg_one}\leanok
For a finite expansion $p=\sum_jv_jT_j$, the four boundary rows are
\[
 p(1)=\sum_jv_j,\quad p(-1)=\sum_j(-1)^jv_j,\quad
 p'(1)=\sum_jj^2v_j,\quad p'(-1)=\sum_j(-1)^{j+1}j^2v_j.
\]
All sums are over the finite support. More generally, any real linear functional
acts by the row of its values on the basis polynomials.

Reference: \cite[Sections 2.4 and 3.2]{olver-townsend}.
\end{theorem}
\begin{proof}
\leanok
\end{proof}

\begin{definition}[Boundary-bordered coefficient map]
\label{def:bordered-coefficient-map}\uses{def:operators}
\lean{Module.Basis.borderedCoefficientMap}\leanok
For source and target bases $b,c$, a linear residual map $L$ and a linear map
$B$ of $q$ boundary constraints, the bordered coefficient map sends finite
coefficients $v$ to $(Bp,\operatorname{repr}_c(Lp))$, where
$p=\operatorname{repr}_b^{-1}v$.

Reference: \cite[Sections 2.4 and 3.2]{olver-townsend}.
\end{definition}

\begin{theorem}[Exact boundary and residual assembly]
\label{thm:bordered-coefficient-map}\uses{def:bordered-coefficient-map}
\lean{Module.Basis.borderedCoefficientMap_apply, Module.Basis.borderedCoefficientMap_eq_iff}\leanok
The bordered equation with right-hand side $(\beta,\operatorname{repr}_c f)$
is equivalent to $Bp=\beta$ and $Lp=f$ for the reconstructed finite expansion.
This holds for module bases over a commutative semiring. In particular it applies
to polynomial differential operators and their boundary functionals.
It asserts neither invertibility nor convergence of a truncated ODE solver.

Reference: \cite[Sections 2.4 and 3.2]{olver-townsend}.
\end{theorem}
\begin{proof}
\leanok
\end{proof}

\paragraph{Example (Julia).} Solve $u'+xu=0$, $u(-1)=1$, whose solution is
$e^{(1-x^2)/2}$. Each matrix below is one stated in Chapter~\ref{chap:banded}:
differentiation $T\to U$, conversion $T\to U$, and multiplication by $x$ in the
$U$ basis. The residual block is banded, and the boundary row is
$u(-1)=\sum_j(-1)^jv_j$.
\begin{verbatim}
N = 30
D = zeros(N+1, N); for n in 1:N-1; D[n, n+1] = n; end   # T_n' = n U_{n-1}
S = zeros(N+1, N); S[1,1] = 1; S[2,2] = 1/2            # T -> U
for n in 0:N-3; S[n+3, n+3] += 1/2; S[n+1, n+3] -= 1/2; end
M = zeros(N+1, N+1)                     # x U_n = (U_{n+1}+U_{n-1})/2
for i in 1:N; M[i+1, i] = 1/2; M[i, i+1] = 1/2; end
L = D + M*S
v = [[(-1)^j for j in 0:N-1]'; L[1:N-1, :]] \ [1.0; zeros(N-1)]
u(x) = sum(v[j+1]*cos(j*acos(x)) for j in 0:N-1)
err(x) = abs(u(x) - exp((1 - x^2)/2))
maximum(err, range(-1, 1, length=201))          # about 1e-15
\end{verbatim}
The residual block \texttt{L[1:N-1,:]} has 84 nonzeros out of 870. Exactness of
the assembled operator on polynomials is established by the preceding results.
The numerical error illustrates the solve; it is not a floating-point error bound.

For analytic solutions the same matrices act on the full coefficient sequence.
Higher derivatives of Chebyshev series land in $C^{(k)}$ with geometric errors,
and truncating a represented solution leaves geometrically small residual and
boundary defects.

\begin{supporting}[Reference details: Classical derivatives of normally convergent polynomial series and related results]
\begin{theorem}[Classical derivatives of normally convergent polynomial series]
\label{thm:normal-series-classical-domain}\uses{thm:open-domain-analytic-derivatives}
\lean{Polynomial.analyticOnOpen_of_hasSum_polynomial, Polynomial.exists_iteratedDerivInput_of_polynomialSeries}\leanok
On an open $U\subseteq\mathbb C$, suppose $\sum_nu_n<\infty$ and
$|a_nP_n(z)|\le u_n$ for every $n$ and $z\in U$. Every function $f:U\to\mathbb C$
represented by this series is analytic. For every $m\ge0$, the sum belongs to $\mathcal D^m(U)$: it admits an order-$m$ derivative chain on $U$, with initial function $z\mapsto\sum_n a_nP_n(z)$.
\end{theorem}
\begin{proof}
\leanok
\end{proof}

\begin{definition}[Differential action on complex functions]
\label{def:function-differential-operator}\uses{def:open-domain-derivative-jet}
\lean{Polynomial.differentialOperatorOnFunctions}\leanok
Let $U\subseteq\mathbb C$ be open and $f:U\to\mathbb C$. For complex polynomial
coefficients $A_0,\ldots,A_m$ and a point $z\in U$ at which all derivatives of orders $0,\ldots,m$ exist in the sense of Definition~\ref{def:open-domain-derivative-jet}, define
$\mathcal Lf(z)=\sum_{k=0}^m A_k(z)J_kf(z)$. The highest derivative need
only exist at $z$; no value is assigned at points without the required jet.

Reference: \cite[Section 3]{olver-townsend}.
\end{definition}

\begin{theorem}[Classical differential-action characterisation]
\label{thm:function-differential-operator}\uses{def:function-differential-operator,thm:open-domain-derivative-jet,thm:open-domain-analytic-derivatives,thm:normal-series-higher-derivatives,thm:polynomial-differential-operator}
\lean{Polynomial.differentialOperatorOnFunctions_ofIterated, Polynomial.iteratedDerivInputOnOpen, Polynomial.iteratedDerivInputOnOpen_val, Polynomial.differentialOperatorOnFunctions_eval}\leanok
For $f\in\mathcal D^m(U)$, the action equals
$\sum_{k=0}^m A_k(z)D_U^kf(z)$, where $D_U^kf$ is obtained from the order-$k$ truncation of the derivative chain of $f$. Every polynomial $p$ restricted to $U$ belongs to $\mathcal D^m(U)$, and
$\mathcal Lp(z)=(Lp)(z)$, where $L$ is polynomial differential assembly.
\end{theorem}
\begin{proof}
\leanok
\end{proof}
\end{supporting}



\begin{theorem}[Consistency of analytic and coefficient differential operators]
\label{thm:series-differential-operator}\uses{thm:series-matrix-reconstruction,thm:function-differential-operator}
\lean{Polynomial.differentialOperator_map, Polynomial.hasSum_differentialOperator_series, Polynomial.hasSum_seriesCoefficientMap_differentialOperator, Polynomial.hasSum_seriesCoefficientMap_differentialOperator_eval₂, Polynomial.hasSum_seriesCoefficientMap_differentialOperator_eval₂_of_rows}\leanok
The polynomial and function actions agree on polynomials, and polynomial
assembly commutes with changing the coefficient ring.
If $U\subset\C$ is open and $|a_nP_n(z)|\le u_n$ on $U$ with $\sum_n u_n<\infty$, then
$\mathcal LS=\sum_n a_n LP_n$ on $U$.
If the coefficient matrix reconstruction is absolutely convergent at $z$,
its weighted contribution series represents $\mathcal LS(z)$. If each
coefficient row converges, this is also the coefficient-times-mode series.
This includes bases over a ring $R$ evaluated through $\sigma:R\to\mathbb C$.

\end{theorem}
\begin{proof}
\leanok
\end{proof}

\rvnote{thm:series-differential-operator}{
No ODE solvability assertion is made.
}

\begin{supporting}[Reference details: Higher Chebyshev derivatives in ultraspherical bases and related results]
\begin{theorem}[Higher Chebyshev derivatives in ultraspherical bases]
\label{thm:chebyshev-series-higher-derivatives}\uses{thm:normal-series-higher-derivatives,thm:chebyshev-gegenbauer-entry,def:chebyshev-series-derivative-input,thm:open-domain-derivative-value}
\lean{Polynomial.Chebyshev.norm_chebyshev_term_le_geometric, Polynomial.Chebyshev.hasSum_iterated_derivative_complexChebyshevSeries, Polynomial.Chebyshev.hasSum_gegenbauer_iterated_derivative_series}\leanok
If $\rho>1$ and $|a_n|\le C\rho^{-n}$, the complex Chebyshev sum
$S_a=\sum_n a_nT_n$ can be differentiated termwise to every order on $E_\rho$.
For $k\ge0$ its positive-order derivative is
\[
 S_a^{(k+1)}(z)=\sum_{n\ge0}a_{n+k+1}2^k k!(n+k+1)C_n^{(k+1)}(z).
\]

Reference: \cite[Section 3]{olver-townsend}.
\end{theorem}
\begin{proof}
\leanok
\end{proof}

\begin{theorem}[Uniform all-order geometric derivative errors]
\label{thm:chebyshev-series-higher-error}\uses{thm:chebyshev-series-higher-derivatives}
\lean{Polynomial.Chebyshev.norm_iterated_derivative_complexChebyshevSeries_sub_sum_le, Polynomial.Chebyshev.exists_uniform_iterated_derivative_series_bound}\leanok
Assume $\rho>1$ and $|a_n|\le C\rho^{-n}$ for every $n$; choose $1<r<\rho$, $z\in E_r$,
and $\delta>0$ with its circle about $z$ contained in $E_r$. Then
\[
 \left|S_a^{(k)}(z)-\sum_{n<N}a_nT_n^{(k)}(z)\right|
 \le \frac{k!}{\delta^k}\frac{C(r/\rho)^N}{1-r/\rho}.
\]
For every compact $K\subset E_r$ a single positive $\delta$ works for all
$z\in K$, all derivative orders $k$, and all cutoffs $N$.

Reference: \cite[Theorem 21.1]{atap}.
\end{theorem}
\begin{proof}
\leanok
\end{proof}

\begin{theorem}[Finite rows of polynomial ultraspherical residuals]
\label{thm:ultraspherical-series-row-input}\uses{thm:all-order-ultraspherical,thm:chebyshev-multiplication-entries,thm:banded-series-coefficient-map}
\lean{Polynomial.exists_isBanded_chebyshev_mul, Polynomial.exists_isBanded_ultraspherical_differentialOperator, Polynomial.ultrasphericalDifferentialSeriesInput, Polynomial.ultrasphericalDifferentialSeriesInput_val}\leanok
Polynomial multiplication in the Chebyshev basis has finite bandwidth.
For every $m\ge0$ and real polynomial coefficients $A_k$, the matrix from
$T$ to $B_m$ for $\sum_{k=0}^mA_kD^k$ has finite bandwidth. Thus for every complex coefficient sequence $a$, each row series $\sum_j M_{ij}a_j$ has only finitely many nonzero terms and converges,
including for the order-zero multiplication operator.
\end{theorem}
\begin{proof}
\leanok
\end{proof}
\end{supporting}


\begin{theorem}[Analytic ultraspherical residual assembly]
\label{thm:analytic-ultraspherical-residual}\uses{thm:series-differential-operator,thm:ultraspherical-series-row-input}
\lean{Polynomial.ultraspherical_series_differentialOperator_eq_sum, Polynomial.hasSum_ultraspherical_differentialOperator_series}\leanok
Let $L=\sum_{k=0}^m A_kD^k$ have real polynomial coefficients.
On a normally convergent complex Chebyshev expansion, the existing matrix
from $T$ to $B_m$, under the weighted reconstruction condition of Theorem~\ref{thm:series-differential-operator},
represents the analytic differential expression $\mathcal LS$.
For $m\ge1$ and $\operatorname{natDegree} A_k\le d$, each output coefficient
uses only indices $j<i+d+2m+1$. This is the same matrix as finite polynomial
assembly, without introducing a second complex basis definition.
\end{theorem}
\begin{proof}
\leanok
\end{proof}

\begin{theorem}[Geometric residual and boundary errors]
\label{thm:chebyshev-series-residual-boundary-errors}\uses{thm:series-residual-boundary-errors,thm:chebyshev-series-higher-error}
\lean{Polynomial.Chebyshev.norm_differentialOperator_complexChebyshevSeries_sub_sum_le, Polynomial.Chebyshev.norm_boundary_complexChebyshevSeries_sub_sum_le}\leanok
For $\rho>1$, $|a_n|\le C\rho^{-n}$ and $1<r<\rho$, the residual and boundary inequalities of Theorem~\ref{thm:series-residual-boundary-errors} hold for the Chebyshev truncations with
\[
 t_N=\frac{C(r/\rho)^N}{1-r/\rho},
\]
provided the evaluation points and the specified positive-radius circles lie
in $E_r$.
\end{theorem}
\begin{proof}
\leanok
\end{proof}

\rvnote{thm:chebyshev-series-residual-boundary-errors}{
These estimates concern truncation of a represented function, not convergence or invertibility of a numerical ODE solver.
}

Both real endpoints can be used with sufficiently small radii, since
$\pm1\in E_r$ and $E_r$ is open; this step is not formalised.

\begin{supporting}[Supporting results: all-order series differentiation, generic residual bounds]
\begin{theorem}[All-order differentiation of polynomial series]
\label{thm:normal-series-higher-derivatives}\uses{def:polynomial-linear-maps,thm:series-operation-errors}
\lean{Polynomial.iteratedDeriv_eval, Polynomial.tendstoLocallyUniformlyOn_iteratedDeriv, Polynomial.tendstoLocallyUniformlyOn_iterated_derivative_series, Polynomial.tendstoUniformlyOn_iterated_derivative_series, Polynomial.hasSum_iterated_derivative_series, Polynomial.norm_iterated_derivative_series_sub_sum_le}\leanok
Let $U\subset\mathbb C$ be open and $P_n\in\mathbb C[x]$. Suppose
$\sum_n u_n<\infty$ and $|a_nP_n(z)|\le u_n$ for every $z\in U$.
Then $S=\sum_n a_nP_n$ satisfies, for every integer $k\ge0$,
\[
 S^{(k)}(z)=\sum_n a_n P_n^{(k)}(z),
\]
with locally uniform convergence of the differentiated finite sums, hence
uniform convergence on every compact subset of $U$. If $z\in U$, $\delta>0$,
and the circle of radius $\delta$ about $z$ lies in $U$, then
\[
 \left|S^{(k)}(z)-\sum_{n<N}a_nP_n^{(k)}(z)\right|
 \le\frac{k!}{\delta^k}\sum_{n\ge N}u_n.
\]
No new majorant for each derivative order is assumed.

Reference: Weierstrass's theorem on locally uniform limits of analytic functions.
\end{theorem}
\begin{proof}
\leanok
\end{proof}

\begin{theorem}[Residual and boundary truncation errors]
\label{thm:series-residual-boundary-errors}\uses{thm:normal-series-classical-domain,thm:function-differential-operator}
\lean{Polynomial.norm_differentialOperator_series_sub_sum_le, Polynomial.norm_boundary_series_sub_sum_le}\leanok
Let $p_N=\sum_{n<N}a_nP_n$ and $t_N=\sum_{n\ge N}u_n$.
For $z\in U$ and a radius-$\delta$ circle contained in $U$, with $\delta>0$,
\[
 |\mathcal LS(z)-(Lp_N)(z)|\le
 \left(\sum_{k=0}^m |A_k(z)|\frac{k!}{\delta^k}\right)t_N.
\]
For a finite boundary functional $Bf=\sum_i c_i f^{(k_i)}(z_i)$, where
$z_i\in U$ and positive radii $\delta_i$ have their circles in $U$,
\[
 |BS-Bp_N|\le\left(\sum_i|c_i|\frac{k_i!}{\delta_i^{k_i}}\right)t_N.
\]
This includes value, derivative, and Robin conditions.
\end{theorem}
\begin{proof}
\leanok
\end{proof}

\end{supporting}

\chapter{From exact operators to convergent solvers}\label{chap:solvers}
\rvsummary{Convergent finite solves for polynomial differential operators, under explicit bounded-boundary and infinite-operator invertibility assumptions.}

For convergence of truncated solves, we precondition the bordered operator
to obtain a compact perturbation of the identity $I+K$. Stable finite sections
then converge, with error controlled by the projection error of the solution.
For first-order problems $u'+pu$ with polynomial $p$, the preconditioner is
$R=\operatorname{diag}(1,1,1/2,1/3,\ldots)$. Its decreasing diagonal compensates
for the growing differentiation entries, giving a bounded operator on
coefficient space.

If the infinite preconditioned operator is invertible, large finite sections are
uniquely solvable, and convergence of the preconditioned coefficients in $\ell^2$
implies uniform convergence of the reconstructed function, including the
endpoints. For $p=0$ we prove invertibility.

\begin{supporting}[Reference details: Inversion on invertible indices and related results]
\begin{definition}[Inversion on invertible indices]
\label{def:invertible-index}
\lean{Ring.InvertibleIndex, Ring.inverseAt}\leanok
For a family $A:I\to R$ in a monoid, let $I_A=\{i\in I:A_i\text{ is invertible}\}$. For $i\in I_A$, define $A_i^{-1}$ by inversion in the group of units of $R$. There is no inverse value at an index outside $I_A$.
\end{definition}

\begin{theorem}[Indexed inverse identities]
\label{thm:invertible-index}\uses{def:invertible-index}
\lean{Ring.inverseAt_mul}\leanok
For each $i\in I_A$, the indexed inverse satisfies $A_i^{-1}A_i=A_iA_i^{-1}=1$.
\end{theorem}
\begin{proof}
\leanok
\end{proof}

\begin{definition}[Limiting filter on invertible indices]
\label{def:invertible-index-filter}\uses{def:invertible-index}
\lean{Ring.invertibleIndexFilter}\leanok
For a filter $\mathcal F$ on $I$, the induced filter on $I_A$ is its pullback along the inclusion $I_A\hookrightarrow I$. For $I=\mathbb N$ and $\mathcal F$ the filter of tails, this means tending to infinity through invertible indices.
\end{definition}

\begin{theorem}[Preservation of the limiting process]
\label{thm:invertible-index-filter}\uses{def:invertible-index-filter}
\lean{Ring.map_invertibleIndexFilter, Ring.invertibleIndexFilter_neBot}\leanok
If $A_i$ is invertible eventually along $\mathcal F$, pushing the induced filter forward to $I$ recovers $\mathcal F$. If $\mathcal F$ is nontrivial, its induced filter on $I_A$ is also nontrivial. In particular, an eventually invertible sequence has a cofinal domain of invertible indices.
\end{theorem}
\begin{proof}
\leanok
\end{proof}

\begin{theorem}[Continuity of indexed inverses]
\label{thm:invertible-index-continuity}\uses{def:invertible-index-filter}
\lean{Ring.tendsto_inverseAt}\leanok
In a complete normed ring, if $A_i\to U$ along $\mathcal F$ and $U$ is a unit, then $A_i^{-1}\to U^{-1}$ along the induced filter on invertible indices.
\end{theorem}
\begin{proof}
\leanok
\end{proof}
\end{supporting}


\begin{theorem}[Stable finite sections of an invertible compact perturbation]
\label{thm:compact-finite-section-convergence}\uses{thm:compact-projection-approximation,thm:invertible-index-filter}
\lean{ContinuousLinearMap.liftedFiniteSection, ContinuousLinearMap.invertible_liftedFiniteSection_filter_neBot, ContinuousLinearMap.eventually_stable_liftedFiniteSection, ContinuousLinearMap.existsUnique_projected_solution, ContinuousLinearMap.eventually_norm_liftedFiniteSection_error_le, ContinuousLinearMap.tendsto_liftedFiniteSection_solution}\leanok
Let $E$ be a real or complex Banach space, $K$ compact, and $T=I+K$ invertible. Let uniformly bounded projections $P_n$ converge strongly to $I$. For sufficiently large $n$, $A_n=I+P_nK$ is invertible with uniformly bounded inverses. For $f=Tx$, the unique projected solution $x_n\in\operatorname{ran}P_n$ of $P_nTx_n=P_nf$ is $A_n^{-1}P_nf$ and satisfies $\|x_n-x\|\le(\|T^{-1}\|+1)\|P_nx-x\|$. The solution $x_n$ is defined for every index $n$ at which $A_n$ is invertible. These indices include every sufficiently large $n$, and $x_n\to x$ as $n\to\infty$ through them.

Reference: \cite[Section 4]{olver-townsend}.
\end{theorem}
\begin{proof}
\leanok
\end{proof}

\begin{theorem}[First-order compact preconditioning]
\label{thm:first-order-compact-preconditioner}\uses{thm:lp-compact-diagonal}
\lean{Polynomial.Chebyshev.firstOrderPreconditioner, Polynomial.Chebyshev.isCompactOperator_firstOrderPreconditioner, Polynomial.Chebyshev.firstOrderBoundary, Polynomial.Chebyshev.firstOrderCompactRemainder, Polynomial.Chebyshev.isCompactOperator_firstOrderCompactRemainder, Polynomial.Chebyshev.firstOrder_bordered_row_zero, Polynomial.Chebyshev.firstOrder_bordered_row_succ}\leanok
Let $R=\operatorname{diag}(1,1,1/2,1/3,\ldots)$ and $Bv=\sum_j(-1)^j(Rv)_j$. The endpoint functional $B$ is bounded and $R$ is compact. For any bounded lower-order map $A$, the bordered operator has row zero $Bv$ and rows $v_{j+1}+(ARv)_j$. It is $I+K$ with $K$ compact.

Reference: \cite[Section 4.1, Lemma 4.3]{olver-townsend}.
\end{theorem}
\begin{proof}
\leanok
\end{proof}

\begin{theorem}[Uniform reconstruction from preconditioned coefficients]
\label{thm:first-order-reconstruction}\uses{thm:first-order-compact-preconditioner,thm:absolute-series-coefficients}
\lean{Polynomial.Chebyshev.hasSum_firstOrderBoundary, Polynomial.Chebyshev.ComplexChebyshevSeriesPoint.ofFirstOrder, Polynomial.Chebyshev.ComplexChebyshevSeriesPoint.ofFirstOrder_val, Polynomial.Chebyshev.complexChebyshevSeries_firstOrderPreconditioner_neg_one, Polynomial.Chebyshev.summable_norm_firstOrderPreconditioner, Polynomial.Chebyshev.norm_firstOrder_reconstruction_sub_le, Polynomial.Chebyshev.tendstoUniformlyOn_firstOrder_reconstruction}\leanok
Every $v\in\ell^2$ gives absolutely summable coefficients $Rv$. The Chebyshev series with coefficients $Rv$ takes the value $Bv$ at $-1$. Reconstruction errors are bounded uniformly by $\|((-1)^jR_{jj})_j\|_2\|v-w\|_2$. Thus convergence in the preconditioned coefficient norm implies uniform convergence including endpoints.
\end{theorem}
\begin{proof}
\leanok
\end{proof}

The reconstructed series is continuous on $[-1,1]$, as a uniform limit of
polynomials; this is not stated as a separate formal result.

\begin{theorem}[Polynomial first-order finite-section convergence]
\label{thm:first-order-polynomial-convergence}\uses{thm:lp-finite-systems,thm:first-order-reconstruction,thm:bounded-chebyshev-operators}
\lean{Polynomial.Chebyshev.firstOrderPolynomialTerm, Polynomial.Chebyshev.firstOrderPolynomialTerm_coefficientLp, Polynomial.Chebyshev.isCompactOperator_firstOrderPolynomial, Polynomial.Chebyshev.eventually_existsUnique_firstOrderPolynomial_system, Polynomial.Chebyshev.tendstoUniformlyOn_firstOrderPolynomial_solution}\leanok
For real polynomial $p$, take $A=S p(J)$ in the bordered assembly for $u\mapsto (u(-1),u\prime+pu)$. This extends exactly the existing polynomial lower-order map. If the infinite preconditioned bordered operator is invertible, sufficiently large finite systems are uniquely solvable. Indexed by the
invertible finite sections tending to infinity, their actual solutions
reconstruct functions converging uniformly on $[-1,1]$.

Reference: \cite[Section 4.2]{olver-townsend}.
\end{theorem}
\begin{proof}
\leanok
\end{proof}

\rvnote{thm:first-order-polynomial-convergence}{
This is a first-order specialization; higher-order solver convergence is not asserted here.
}

\begin{supporting}[Reference details: Invertibility of the bordered derivative and related results]
\begin{theorem}[Invertibility of the bordered derivative]
\label{thm:first-order-derivative-invertibility}\uses{thm:first-order-compact-preconditioner}
\lean{Polynomial.Chebyshev.firstOrderBoundaryCorrection_sq, Polynomial.Chebyshev.isUnit_firstOrderDerivative}\leanok
The boundary correction $K_B$ satisfies $K_B^2=0$, and hence the bordered
derivative $I+K_B$ is invertible with inverse $I-K_B$.
\end{theorem}
\begin{proof}
\leanok
\end{proof}

\begin{definition}[Solution of the bordered derivative system]
\label{def:first-order-derivative-solution}\uses{thm:first-order-derivative-invertibility,def:invertible-index}
\lean{Polynomial.Chebyshev.firstOrderDerivativeSolution}\leanok
For endpoint value $c\in\mathbb C$ and derivative coordinates $f\in\ell^2$,
define $v=(I+K_B)^{-1}(c,f_0,f_1,\ldots)$, using the inverse of the
invertible bordered derivative. The reconstructed solution has coefficients $Rv$.
\end{definition}
\end{supporting}


\begin{theorem}[A boundary problem with proved invertibility]
\label{thm:first-order-derivative-solve}\uses{thm:first-order-reconstruction,thm:chebyshev-series-differentiation,def:first-order-derivative-solution,thm:compact-finite-section-convergence}
\lean{Polynomial.Chebyshev.firstOrderDerivativeSolution_equation, Polynomial.Chebyshev.firstOrderDerivativeSolution_boundary, Polynomial.Chebyshev.firstOrderDerivativeSolution_succ, Polynomial.Chebyshev.firstOrderDerivativeSolution_deriv, Polynomial.Chebyshev.tendstoUniformlyOn_firstOrder_derivative_solution}\leanok
For $p=0$, the boundary correction $K_B$ satisfies $K_B^2=0$, so the bordered derivative has explicit inverse $I-K_B$. Given endpoint $c$ and $U$ coefficients $f\in\ell^2$, its solution has $u(-1)=c$ and preconditioned coordinates $v_{j+1}=f_j$. The actual finite solves, indexed by invertible sections tending to
infinity, converge uniformly to the reconstructed function without an
additional invertibility assumption. If the actual solution coefficients $Rv$ satisfy a geometric bound with $\rho>1$, the classical derivative on $E_\rho$ is $u\prime(z)=\sum_j f_jU_j(z)$.
\end{theorem}
\begin{proof}
\leanok
\end{proof}

\begin{supporting}[Supporting results: coordinate projections, compact diagonals, bounded operators]
\begin{theorem}[Coordinate projections on coefficient space]
\label{thm:lp-coordinate-projections}\uses{thm:mathlib-hilbert-sums}
\lean{lp.coefficientProjection, lp.coefficientProjection_apply, lp.norm_coefficientProjection_le, lp.coefficientProjection_idempotent, lp.tendsto_coefficientProjection, lp.isCompactOperator_coefficientProjection}\leanok
On real or complex $\ell^2$, $P_n$ retains precisely coordinates $0,\ldots,n-1$. It is an idempotent compact operator of norm at most one, and $P_nx\to x$ for every $x$.
\end{theorem}
\begin{proof}
\leanok
\end{proof}

\begin{theorem}[Compact diagonal coefficient operators]
\label{thm:lp-compact-diagonal}\uses{thm:lp-coordinate-projections}
\lean{lp.coefficientDiagonal, lp.coefficientDiagonal_apply, lp.norm_coefficientDiagonal_le, lp.tendsto_coefficientDiagonal_sections, lp.norm_coefficientDiagonal_tail_le, lp.isCompactOperator_coefficientDiagonal}\leanok
A bounded real or complex scalar sequence $d_j\to0$ defines a compact diagonal operator $D$, with $\|D\|\le C$ whenever $|d_j|\le C$. The finite diagonal sections converge to $D$ in operator norm. If $|d_j|\le\varepsilon$ for $j\ge n$, then $\|D-P_nD\|\le\varepsilon$.
\end{theorem}
\begin{proof}
\leanok
\end{proof}

\begin{theorem}[Strong projections approximate compact operators]
\label{thm:compact-projection-approximation}
\lean{ContinuousLinearMap.tendstoUniformlyOn_of_strongly_tendsto_id, ContinuousLinearMap.tendsto_compact_comp_of_strongly_tendsto_id}\leanok
On a real or complex normed space, uniformly bounded continuous linear maps $P_n$ converging strongly to the identity converge uniformly on compact sets. For compact $K$, $P_nK\to K$ in operator norm.
\end{theorem}
\begin{proof}
\leanok
\end{proof}

\begin{theorem}[Finite coefficient systems]
\label{thm:lp-finite-systems}\uses{thm:compact-finite-section-convergence,thm:lp-coordinate-projections}
\lean{lp.coefficientProjection_eq_self_iff, lp.coefficientProjection_eq_iff, lp.eventually_existsUnique_coefficient_system, lp.eventually_coefficient_system_error_le, lp.tendsto_coefficient_system_solution, lp.coefficient_system_filter_neBot}\leanok
On complex $\ell^2$, $P_nx=x$ exactly when $x_j=0$ for $j\ge n$, and
$P_nx=P_ny$ exactly when the first $n$ coordinates agree.
For compact $K$ on $\ell^2$ with $I+K$ invertible, the first $n$ equations of $(I+K)x_n=f$ with $x_{n,j}=0$ for $j\ge n$ have a unique solution for every sufficiently large $n$. For every sufficiently large $n$, the lifted finite section has an actual
inverse and its solution obeys the preceding tail error bound. The inverse norm
constant is $\|(I+K)^{-1}\|+1$. On the cofinal set of invertible cutoffs,
these actual solves converge in $\ell^2$; the corresponding limit filter is
nontrivial. No inverse is evaluated at a singular cutoff.

Reference: \cite[Section 4.2]{olver-townsend}.
\end{theorem}
\begin{proof}
\leanok
\end{proof}

\begin{theorem}[Bounded multiplication and entry conversion]
\label{thm:bounded-chebyshev-operators}\uses{thm:chebyshev-ultraspherical-conversion-matrix,thm:lp-coordinate-projections,thm:ultraspherical-ladder-vectors}
\lean{Module.Basis.coefficientLp, Polynomial.Chebyshev.boundedMulX, Polynomial.Chebyshev.boundedConversion, Polynomial.Chebyshev.boundedPolynomialMultiplier, Polynomial.Chebyshev.boundedMulX_coefficientLp, Polynomial.Chebyshev.boundedConversion_coefficientLp, Polynomial.Chebyshev.boundedPolynomialMultiplier_coefficientLp}\leanok
The full-constant Chebyshev multiplication matrix extends to a bounded operator $J$ on complex $\ell^2$. For a real polynomial $p$, $p(J)$ agrees with multiplication by $p$ on finite polynomial coordinates. The bounded map $S$ with $(Sa)_0=a_0-a_2/2$ and $(Sa)_j=(a_j-a_{j+2})/2$ for $j>0$ agrees with conversion from $T$ to $U=C^{(1)}$.
\end{theorem}
\begin{proof}
\leanok
\end{proof}

\begin{theorem}[First-order ultraspherical residual]
\label{thm:first-order-ultraspherical}\uses{thm:polynomial-differential-operator,thm:chebyshev-ultraspherical-conversion-matrix,thm:chebyshev-ultraspherical-derivative-matrix,thm:gegenbauer-polynomial-multiplication,thm:conversion-composition-support,thm:bandwidth-mono}
\lean{Polynomial.chebyshev_firstOrder_isBanded}\leanok
For arbitrary real polynomial coefficients $a_0,a_1$, represent
$Lu=a_0u+a_1u'$ with Chebyshev input and $C^{(1)}$ output. Put
$d_j=\operatorname{natDegree}a_j$. The residual matrix has bandwidths
\[
 \bigl(\max(d_0,d_1),\ \max(d_0+2,d_1+1)\bigr).
\]
The constant column follows the full-constant convention. Boundary constraints
and invertibility are not included in this residual operator.

Reference: \cite[Section 2]{olver-townsend}.
\end{theorem}
\begin{proof}
\leanok
\end{proof}

\end{supporting}

\section{Higher-order ultraspherical convergence}\label{sec:higher-order-ultraspherical}

For a normalized differential equation of order $N$, differentiation is sparse
when its output is expressed in $C^{(N)}$. Polynomial multiplication is bounded,
and conversion from a lower positive ultraspherical level is compact. These
facts give a compact remainder after right preconditioning. Boundary rows must
be bounded in the chosen coefficient norm; invertibility of the infinite
bordered operator remains an assumption.

\begin{supporting}[Weighted spaces and coefficient operators]
\begin{definition}[Degree-weighted coefficient space]
\label{def:degree-weighted-l2}
\lean{lp.degreeWeightedL2Submodule, lp.DegreeWeightedL2, lp.degreeWeightedL2Equiv, lp.degreeWeightedL2Isometry}\leanok
For $s\in\mathbb N$, let $W_s$ consist of complex sequences $v$ for which
$((n+1)^s v_n)_{n\ge0}\in\ell^2$, with norm
$\|v\|_s=\|((n+1)^s v_n)_n\|_{\ell^2}$.
Write $U_sv=((n+1)^sv_n)_n$.
\end{definition}

\begin{theorem}[Weighted coefficient-space structure]
\label{thm:degree-weighted-l2}\uses{def:degree-weighted-l2}
\lean{lp.norm_degreeWeightedL2, lp.completeSpace_degreeWeightedL2}\leanok
The map $U_s:W_s\to\ell^2$ is a linear isometric equivalence, with
$(U_s^{-1}w)_n=w_n/(n+1)^s$. In particular $W_s$ is a complex Banach space.
\end{theorem}
\begin{proof}
\leanok
\end{proof}

\begin{definition}[Weighted coordinate operators]
\label{def:degree-weighted-operators}\uses{thm:lp-compact-diagonal,def:degree-weighted-l2}
\lean{lp.degreeWeightedTransport, lp.degreeWeightedDiagonal, lp.degreeWeightLeftRatio, lp.degreeWeightRightRatio, lp.degreeWeightedShiftLeft, lp.degreeWeightedShiftRight, lp.degreeWeightedProjection, lp.degreeWeightedInclusion}\leanok
On $W_s$, define the diagonal $D_dv=(d_nv_n)_n$ for a bounded sequence $d$,
the shifts $(Lv)_n=v_{n+1}$ and $(Rv)_0=0$, $(Rv)_{n+1}=v_n$,
and the truncation $(P_mv)_n=v_n$ for $n<m$ and $0$ otherwise.
Let $\iota_s:W_s\to\ell^2$ be the coordinate inclusion. For a bounded linear operator $A$ on
$\ell^2$, its weighted transport is $U_s^{-1}AU_s$.
\end{definition}

\begin{theorem}[Weighted coordinate-operator properties]
\label{thm:degree-weighted-operators}\uses{def:degree-weighted-operators}
\lean{lp.degreeWeightedDiagonal_apply, lp.degreeWeightedShiftLeft_apply, lp.degreeWeightedShiftRight_zero, lp.degreeWeightedShiftRight_succ, lp.degreeWeightedProjection_apply, lp.norm_degreeWeightedProjection_le, lp.tendsto_degreeWeightedProjection, lp.degreeWeightedInclusion_apply, lp.degreeWeightedInclusion_injective, lp.isCompactOperator_degreeWeightedDiagonal, lp.isCompactOperator_degreeWeightedProjection, lp.degreeWeightedShiftLeft_pow_apply, lp.degreeWeightedShiftRight_pow_apply, lp.degreeWeightedShiftRight_mul_left}\leanok
These operators are bounded. The inclusion is injective; $\|P_m\|\le1$,
$P_mv\to v$ in $W_s$, and $P_m$ is compact. A bounded diagonal tending to zero
is compact on $W_s$. Moreover $(L^Nv)_n=v_{n+N}$ and $R^NL^N=I-P_N$.
\end{theorem}
\begin{proof}
\leanok
\end{proof}

\begin{definition}[Positive-level coefficient operators]
\label{def:positive-ultraspherical-operators}\uses{thm:lp-compact-diagonal}
\lean{Polynomial.Chebyshev.positiveConversionDiagonal, Polynomial.Chebyshev.positiveConversionScaling, Polynomial.Chebyshev.boundedPositiveConversion, Polynomial.Chebyshev.positiveMulXUp, Polynomial.Chebyshev.positiveMulXDown, Polynomial.Chebyshev.boundedPositiveMulX, Polynomial.Chebyshev.boundedPositivePolynomialMultiplier}\leanok
For $k\in\mathbb N$ and $\lambda=k+1$, define $d_n=\lambda/(n+\lambda)$,
$u_n=(n+1)/(2(n+\lambda))$, and
$\ell_n=(n+2\lambda-1)/(2(n+\lambda))$.
On complex $\ell^2$ set $S_\lambda=(I-L^2)D_d$,
$J_\lambda=RD_u+LD_\ell$, and $M_\lambda[a]=a(J_\lambda)$ for $a\in\mathbb R[x]$.
\end{definition}

\begin{theorem}[Positive ultraspherical conversion and multiplication]
\label{thm:positive-ultraspherical-operators}\uses{def:positive-ultraspherical-operators,thm:bounded-chebyshev-operators,thm:gegenbauer-conversion-matrix,thm:gegenbauer-multiplication-matrix}
\lean{Polynomial.Chebyshev.isCompactOperator_boundedPositiveConversion, Polynomial.Chebyshev.boundedPositiveConversion_coefficientLp, Polynomial.Chebyshev.boundedPositiveMulX_coefficientLp, Polynomial.Chebyshev.boundedPositivePolynomialMultiplier_coefficientLp}\leanok
The map $S_\lambda$ is compact on $\ell^2$ and converts polynomial
$C^{(\lambda)}$ coordinates to $C^{(\lambda+1)}$ coordinates.
The bounded maps $J_\lambda$ and $M_\lambda[a]$ represent multiplication by
$x$ and $a(x)$, respectively, in the $C^{(\lambda)}$ basis.
\end{theorem}
\begin{proof}
\leanok
\end{proof}

\begin{definition}[Weighted polynomial coordinates]
\label{def:weighted-polynomial-coordinates}\uses{def:degree-weighted-l2}
\lean{Module.Basis.degreeWeightedCoefficient}\leanok
For a real polynomial basis $b=(b_n)$ and $p\in\mathbb R[x]$, let
$c_b(p)\in W_s$ be its finite coordinate sequence, regarded as complex coefficients.
\end{definition}

\begin{theorem}[Weighted polynomial coordinate identities]
\label{thm:weighted-polynomial-coordinates}\uses{def:weighted-polynomial-coordinates,thm:degree-weighted-operators,thm:bounded-chebyshev-operators}
\lean{Module.Basis.degreeWeightedCoefficient_inclusion, Module.Basis.degreeWeightedCoefficient_add, Module.Basis.degreeWeightedCoefficient_smul}\leanok
Weighted polynomial coordinates preserve addition and real scalar multiplication.
Their inclusion into $\ell^2$ is the ordinary finite polynomial coefficient sequence.
\end{theorem}
\begin{proof}
\leanok
\end{proof}

\begin{definition}[Weighted ultraspherical operators]
\label{def:weighted-ultraspherical-operators}\uses{def:positive-ultraspherical-operators,def:degree-weighted-operators}
\lean{Polynomial.Chebyshev.weightedChebyshevConversion, Polynomial.Chebyshev.weightedChebyshevMulX, Polynomial.Chebyshev.weightedChebyshevPolynomialMultiplier, Polynomial.Chebyshev.weightedPositiveConversion, Polynomial.Chebyshev.weightedPositiveMulX, Polynomial.Chebyshev.weightedPositivePolynomialMultiplier, Polynomial.Chebyshev.weightedPositiveRaise}\leanok
On $W_s$, use the same coordinate formulas for $S_\lambda$, $J_\lambda$ and
$M_\lambda[a]$ at integer levels $\lambda\ge1$. At level zero set
$S_0=\tfrac12(I-L^2+P_1)$,
$J_0=\tfrac12(R(I+P_1)+L)$ and $M_0[a]=a(J_0)$.
Write $S_{j\to N}=S_{N-1}\cdots S_j$ for successive conversion, with
$S_{j\to j}=I$.
\end{definition}

\begin{theorem}[Weighted ultraspherical polynomial agreement]
\label{thm:weighted-ultraspherical-operators}\uses{def:weighted-ultraspherical-operators,thm:positive-ultraspherical-operators,thm:weighted-polynomial-coordinates}
\lean{Polynomial.Chebyshev.weightedChebyshevConversion_coefficient, Polynomial.Chebyshev.weightedChebyshevMulX_coefficient, Polynomial.Chebyshev.weightedChebyshevPolynomialMultiplier_coefficient, Polynomial.Chebyshev.weightedPositiveConversion_coefficient, Polynomial.Chebyshev.weightedPositiveMulX_coefficient, Polynomial.Chebyshev.weightedPositivePolynomialMultiplier_coefficient, Polynomial.Chebyshev.weightedPositiveRaise_coefficient, Polynomial.Chebyshev.isCompactOperator_weightedPositiveConversion, Polynomial.Chebyshev.isCompactOperator_weightedPositiveRaise}\leanok
On $W_s$, multiplication by $x$ and by a polynomial, conversion from the
Chebyshev basis to the first ultraspherical basis, and conversion between
positive ultraspherical parameters are bounded operators. On finite
polynomial coefficient vectors, their coordinates equal those of the
polynomial multiplication and conversion maps in
Definition~\ref{def:weighted-ultraspherical-operators}.
Every conversion $S_{j\to N}$ with $1\le j<N$ is compact on $W_s$.
\end{theorem}
\begin{proof}
\leanok
\end{proof}

\begin{definition}[Higher-order right preconditioning]
\label{def:weighted-ultraspherical-preconditioner}\uses{def:degree-weighted-operators}
\lean{Polynomial.Chebyshev.weightedFirstOrderPreconditioner, Polynomial.Chebyshev.ultrasphericalDerivativeScale, Polynomial.Chebyshev.weightedUltrasphericalPreconditioner, Polynomial.Chebyshev.weightedPreconditionedDerivative}\leanok
Let $N=k+1$ and $c_k=2^k k!$. Define
$\mathcal R_N=c_k^{-1}\operatorname{diag}(1,1,1/2,1/3,\ldots)$ on $W_s$.
For $j\ge0$, define the preconditioned order-$(j+1)$ derivative by
$E_{N,j+1}=(c_j/c_k)L^{j+1}$.
\end{definition}

\begin{theorem}[Preconditioned polynomial derivatives]
\label{thm:weighted-ultraspherical-preconditioner}\uses{def:weighted-ultraspherical-preconditioner,thm:chebyshev-ultraspherical-derivative-matrix,thm:degree-weighted-operators,thm:ultraspherical-ladder-vectors,def:weighted-polynomial-coordinates}
\lean{Polynomial.Chebyshev.ultrasphericalDerivativeScale_ne_zero, Polynomial.Chebyshev.isCompactOperator_weightedUltrasphericalPreconditioner, Polynomial.Chebyshev.weightedPreconditionedDerivative_top, Polynomial.ultraspherical_derivative_coordinate, Polynomial.Chebyshev.weightedPreconditionedDerivative_coefficient}\leanok
The map $\mathcal R_N$ is compact and $E_{N,N}=L^N$.
If $\mathcal R_Nv=c_T(p)$ for a polynomial $p$, then
$E_{N,j+1}v=c_{C^{(j+1)}}(p^{(j+1)})$ for every $j\ge0$.
\end{theorem}
\begin{proof}
\leanok
\end{proof}

The diagonal above differs from the usual order-$N$ preconditioner only in
finitely many modes below degree $N$. These modes are annihilated by the
highest derivative; the chosen normalization also makes every positive
lower derivative a scalar multiple of a shift.

\begin{definition}[Bounded boundary rows]
\label{def:weighted-boundary-rows}\uses{def:degree-weighted-operators}
\lean{lp.degreeWeightedBoundaryCorrection, lp.degreeWeightedBorderedRemainder}\leanok
For bounded maps $B,A:W_s\to W_s$, define
$K_{N,B,A}=P_N(B-I)+R^N A$. The first $N$ coordinates of $Bv$ specify
the boundary data; coordinates of $Bv$ with index at least $N$ are unused.
\end{definition}

\begin{theorem}[Bordering a compact interior remainder]
\label{thm:weighted-boundary-rows}\uses{def:weighted-boundary-rows,thm:degree-weighted-operators}
\lean{lp.isCompactOperator_degreeWeightedBoundaryCorrection, lp.isCompactOperator_degreeWeightedBorderedRemainder, lp.degreeWeightedBorderedRemainder_rows}\leanok
If $A$ is compact, so is $K_{N,B,A}$. The first $N$ rows of $I+K_{N,B,A}$
are $B$, and its row $N+i$ is $v_{N+i}+(Av)_i$.
\end{theorem}
\begin{proof}
\leanok
\end{proof}

\end{supporting}

\begin{definition}[Preconditioned polynomial differential system]
\label{def:higher-order-ultraspherical-system}\uses{def:weighted-boundary-rows,def:weighted-ultraspherical-operators,def:weighted-ultraspherical-preconditioner}
\lean{Polynomial.Chebyshev.weightedUltrasphericalZeroTerm, Polynomial.Chebyshev.weightedUltrasphericalPositiveTerm, Polynomial.Chebyshev.weightedUltrasphericalLowerTerms, Polynomial.Chebyshev.weightedUltrasphericalRemainder}\leanok
Let $N=k+1$, $a_0,\ldots,a_k\in\mathbb R[x]$, and $B:W_s\to W_s$ be bounded.
Set
\[
 A=S_{0\to N}M_0[a_0]\mathcal R_N+
   \sum_{j=1}^{N-1}S_{j\to N}M_j[a_j]E_{N,j},\qquad
 K=P_N(B-I)+R^N A,\qquad T=I+K.
\]
\end{definition}

\begin{theorem}[Higher-order compact preconditioning]
\label{thm:higher-order-ultraspherical-compact}\uses{def:higher-order-ultraspherical-system,thm:weighted-boundary-rows,thm:weighted-ultraspherical-preconditioner,thm:weighted-ultraspherical-operators}
\lean{Polynomial.Chebyshev.isCompactOperator_weightedUltrasphericalZeroTerm, Polynomial.Chebyshev.isCompactOperator_weightedUltrasphericalPositiveTerm, Polynomial.Chebyshev.isCompactOperator_weightedUltrasphericalLowerTerms, Polynomial.Chebyshev.isCompactOperator_weightedUltrasphericalRemainder, Polynomial.Chebyshev.weightedUltraspherical_polynomial_equation, Polynomial.Chebyshev.weightedUltraspherical_boundary_rows, Polynomial.Chebyshev.weightedUltraspherical_interior_polynomial_rows}\leanok
The operators $A$ and $K$ are compact on $W_s$.
If $\mathcal R_Nv=c_T(p)$ for a polynomial $p$, the first $N$ rows of $Tv$
are $Bv$, and the remaining rows are the $C^{(N)}$ coordinates of
\[
 p^{(N)}+\sum_{j=0}^{N-1}a_jp^{(j)}.
\]
Reference: \cite[Section 4.1]{olver-townsend}.
\end{theorem}
\begin{proof}
\leanok
\end{proof}

\begin{supporting}[Finite-section estimates and reconstruction]
\begin{theorem}[Finite systems in weighted coefficient space]
\label{thm:weighted-finite-systems}\uses{thm:degree-weighted-operators,thm:compact-finite-section-convergence}
\lean{lp.degreeWeightedProjection_idempotent, lp.degreeWeightedProjection_eq_self_iff, lp.degreeWeightedProjection_eq_iff, lp.eventually_existsUnique_degreeWeighted_system, lp.eventually_degreeWeighted_system_error_le, lp.tendsto_degreeWeighted_system_solution, lp.degreeWeighted_system_filter_neBot}\leanok
Let $K$ be compact on $W_s$ and $T=I+K$ invertible. For all sufficiently large
$n$, the equations $P_nTx_n=P_nf$, $P_nx_n=x_n$ have a unique solution.
For $f=Tx$, along the cofinal set of invertible lifted sections,
$x_n\to x$ and, eventually,
\[
 \|x_n-x\|_s\le(\|T^{-1}\|+1)\|P_nx-x\|_s.
\]
\end{theorem}
\begin{proof}
\leanok
\end{proof}

\begin{definition}[Reconstruction after right preconditioning]
\label{def:weighted-ultraspherical-reconstruction}\uses{thm:first-order-reconstruction,def:weighted-ultraspherical-preconditioner}
\lean{Polynomial.Chebyshev.weightedUltrasphericalReconstruction, Polynomial.Chebyshev.weightedUltrasphericalFunction}\leanok
Let $\widehat{\mathcal R}_N:W_s\to W_{s+1}$ have the same coordinate formula
as $\mathcal R_N$. For $x\in[-1,1]$, define
$\mathcal F_N(v)(x)=\sum_{n\ge0}(\mathcal R_Nv)_nT_n(x)$ using the convergent
Chebyshev series.
\end{definition}

\begin{theorem}[Gain of a degree weight and uniform reconstruction]
\label{thm:weighted-ultraspherical-reconstruction}\uses{def:weighted-ultraspherical-reconstruction,thm:degree-weighted-operators}
\lean{Polynomial.Chebyshev.hasSum_weightedUltrasphericalFunction, Polynomial.Chebyshev.weightedUltrasphericalReconstruction_apply, Polynomial.Chebyshev.weightedUltrasphericalReconstruction_eq, Polynomial.Chebyshev.tendstoUniformly_weightedUltrasphericalFunction}\leanok
The map $\widehat{\mathcal R}_N$ is bounded and has precisely the
right-preconditioned coefficients. If $v_m\to v$ in $W_s$, then
$\widehat{\mathcal R}_Nv_m\to\widehat{\mathcal R}_Nv$ in $W_{s+1}$ and
$\mathcal F_N(v_m)\to\mathcal F_N(v)$ uniformly on $[-1,1]$.
\end{theorem}
\begin{proof}
\leanok
\end{proof}

\begin{definition}[Higher-order finite solves]
\label{def:higher-order-ultraspherical-solves}\uses{thm:compact-finite-section-convergence,def:higher-order-ultraspherical-system}
\lean{Polynomial.Chebyshev.weightedUltrasphericalFiniteSection, Polynomial.Chebyshev.weightedUltrasphericalFiniteSolution}\leanok
For the preceding polynomial differential system, let $T_n=I+P_nK$.
At a cutoff $n$ for which $T_n$ is invertible, define $v_n=T_n^{-1}P_nf$.
\end{definition}

\end{supporting}

\begin{theorem}[Higher-order ultraspherical finite-section convergence]
\label{thm:higher-order-ultraspherical-convergence}\uses{def:higher-order-ultraspherical-solves,thm:weighted-finite-systems,thm:weighted-ultraspherical-reconstruction,thm:higher-order-ultraspherical-compact}
\lean{Polynomial.Chebyshev.eventually_existsUnique_weightedUltraspherical_system, Polynomial.Chebyshev.tendsto_weightedUltraspherical_system_solution, Polynomial.Chebyshev.weightedUltraspherical_filter_neBot, Polynomial.Chebyshev.eventually_weightedUltraspherical_error_le, Polynomial.Chebyshev.eventually_weightedUltraspherical_reconstruction_error_le, Polynomial.Chebyshev.tendsto_weightedUltraspherical_reconstruction, Polynomial.Chebyshev.tendstoUniformly_weightedUltraspherical_solution}\leanok
Let $s\in\mathbb N$, $N\ge1$, and let $T$ be the preconditioned polynomial
differential system above, with bounded boundary rows on $W_s$.
Assume $T$ is invertible. For all sufficiently large cutoffs the finite
systems are uniquely solvable. For $f=Tv$, their solutions satisfy
\[
 \|v_n-v\|_s\le(\|T^{-1}\|+1)\|P_nv-v\|_s.
\]
Writing $u_n=\widehat{\mathcal R}_Nv_n$ and $u=\widehat{\mathcal R}_Nv$,
the original coefficients also satisfy
\[
 \|u_n-u\|_{s+1}\le\|\widehat{\mathcal R}_N\|(\|T^{-1}\|+1)\|P_nv-v\|_s.
\]
Thus $v_n\to v$ in $W_s$, $u_n\to u$ in $W_{s+1}$, and the
reconstructed functions converge uniformly on $[-1,1]$.
All limits are through the cofinal set of invertible cutoffs.

Reference: \cite[Section 4.1]{olver-townsend}.
\end{theorem}
\begin{proof}
\leanok
\end{proof}

Part~\ref{part:beyond} treats measure modification and several variables.

\part{Measure modification and several variables}\label{part:beyond}
\rvsummary{New bases from old: measure modification, semiclassical Jacobi families, and classical bases on triangles, disks and annuli.}

\chapter{Polynomial measure modification in one variable}\label{chap:modification}
\rvsummary{Polynomial measure modification, with connection matrices as banded Cholesky factors of Gram matrices.}

Multiplying a weight by a nonnegative polynomial $u$ changes the orthogonal
polynomials. The Gram matrix records the new inner products in the old basis,
and a triangular factor describes the corresponding change of coordinates.
This is the Gram/Cholesky viewpoint on measure modification.

The connection from the old basis to the
new one is upper triangular with bandwidth $\deg u$, and multiplication by $u$
has bandwidths $(\deg u,\deg u)$. Leading sections of the connection are the
Cholesky factors of the new Gram matrix, which is itself a section of
multiplication by $u$, with the orthonormal and positive-diagonal conventions
specified below. This gives a structured way to build new bases from old ones
\cite{gautschi,gutleb-olver-slevinsky}.

\begin{theorem}[Univariate polynomial-modification support]
\label{thm:univariate-modification-support}\uses{thm:general-coordinates,thm:even-positive-sequences,def:polynomial-density-measure}
\lean{Polynomial.HasFiniteMoments.withDensity_polynomial, Polynomial.HasPositiveDefiniteIntegral.withDensity_polynomial, Polynomial.Sequence.IsOrthogonal.basis_repr_eq_zero_of_polynomial_modification, Polynomial.Sequence.IsOrthogonal.basis_repr_eq_zero_of_parity, Polynomial.Sequence.IsOrthogonal.integralInner_pos_of_leadingCoeff_pos}\leanok
A nonzero nonnegative polynomial density $u$ preserves finite moments and
positive definiteness. The connection from an old degree-$n$ orthogonal
polynomial into the modified sequence has zero coefficient at index $m$ if
$m>n$ or $m+\deg u<n$. For even original and modified measures, opposite
parities also vanish. Equal-degree pairings of positive-leading polynomials
with the target orthogonal sequence are positive.

Reference: \cite[Theorem 1.3 and Lemma 2.8]{gutleb-olver-slevinsky}.
\end{theorem}
\begin{proof}
\leanok
\end{proof}

\begin{theorem}[Polynomial multiplication and conversion operators]
\label{thm:univariate-modification-operators}\uses{thm:univariate-modification-support,thm:operator-coefficients}
\lean{Polynomial.Sequence.IsOrthogonal.polynomial_mul_isBanded, Polynomial.Sequence.IsOrthogonal.conversion_isBanded, Polynomial.Sequence.conversion_intertwines_multiplication}\leanok
For a real measure with finite moments and positive definite polynomial pairing,
multiplication by a polynomial $u$ in an orthogonal basis has lower and upper
bandwidth at most $\operatorname{natDegree}u$ (zero for the zero polynomial).
For nonzero $u\ge0$ almost everywhere, conversion
from an original orthogonal sequence to one for $u\,d\mu$ has bandwidths
$(0,\deg u)$. If $C$ denotes conversion and $J_P,J_Q$ multiplication by $x$
in source and target coordinates, then $CJ_P=J_QC$ on finitely supported
coefficients. This is not a same-size finite-section similarity assertion.

Reference: \cite[Section 2]{gutleb-olver-slevinsky}.
\end{theorem}
\begin{proof}
\leanok
\end{proof}

\begin{theorem}[Univariate connection sections and Cholesky]
\label{thm:univariate-connection-gram}\uses{thm:univariate-modification-support,thm:cholesky-lu,thm:positive-gram-factor}
\lean{Polynomial.Sequence.IsOrthogonal.sum_range_repr, Polynomial.Sequence.IsOrthogonal.connection_upperTriangular, Polynomial.Sequence.IsOrthogonal.connection_diag_pos, Polynomial.Sequence.IsOrthogonal.connection_gram, Polynomial.Sequence.IsOrthogonal.connection_eq_of_gram, Polynomial.Sequence.IsOrthogonal.modification_gram, Polynomial.Sequence.IsOrthogonal.modification_lu_support}\leanok
Let $Q$ be an orthonormal sequence for a measure $\nu$ with finite moments and
positive definite polynomial pairing, and let $P$ be any degree-indexed sequence.
For $0\le i,j<N$, set $C_{ij}=[Q_i]P_j$. Then
\[
 [\langle P_i,P_j\rangle_\nu]_{i,j<N}=C^{\mathsf T}C.
\]
The matrix $C$ is upper triangular. If both sequences have positive leading
coefficients, its diagonal is positive and it is the unique such Gram factor.
When $P$ is orthonormal for $\mu$ and $\nu=u\,d\mu$ for a nonzero nonnegative
polynomial $u$, this Gram section equals the leading section of full
multiplication by $u$ in the $P$ basis. If both sequences have positive leading
coefficients, every exact unpivoted LU factorisation with unit lower diagonal
has the same upper-factor zeros as $C$.

Reference: \cite[Lemma 2.8]{gutleb-olver-slevinsky}.
\end{theorem}
\begin{proof}
\leanok
\end{proof}

\begin{theorem}[Composition of polynomial conversions]
\label{thm:conversion-composition-support}\uses{thm:univariate-connection-gram,thm:polynomial-density-measure,thm:operator-coefficients}
\lean{Polynomial.withDensity_polynomial_mul, Module.Basis.toMatrix_comp_isBanded, Polynomial.Sequence.conversion_isBanded_trans, Polynomial.Sequence.IsOrthogonal.connection_comp, Matrix.HasOffsetSupport, Matrix.HasOffsetSupport.mul}\leanok
Successive nonnegative polynomial measure changes multiply their densities.
Leading univariate connection sections compose exactly through an intermediate
orthogonal sequence, without padding. Their upper bandwidths add; more generally,
composition adds lower and upper bandwidths of coefficient maps separately.
For finite matrices with compatible labels in an additive commutative group,
if source-minus-target differences lie in finite sets $S$ and $T$, those of
the product lie in $S+T$. This applies also to pairs of degree labels.
The support bound permits additional cancellations and assumes the same
intermediate ordered basis for both factors.
\end{theorem}
\begin{proof}
\leanok
\end{proof}

\begin{theorem}[Diagonal normalization and unpivoted LU]
\label{thm:cholesky-lu}
\lean{Matrix.exists_lu_of_diagonal_gram_factor, Matrix.unitLower_upper_factorization_unique, Matrix.upper_eq_of_diagonal_gram_factor}\leanok
Suppose $C$ is upper triangular with nonzero diagonal, $D$ is invertible diagonal,
and $C^TC=DWD^{-1}$. For any exact unpivoted factorization $W=LU$ with unit lower
triangular $L$ and upper triangular $U$,
\[
 U_{ij}=D_i^{-1}C_{ii}C_{ij}D_j.
\]
Such a factorization exists and is unique, and $U$ and $C$ have exactly the same
zero entries.
\end{theorem}
\begin{proof}
\leanok
\end{proof}

\rvnote{thm:cholesky-lu}{
This makes no assertion about floating-point thresholding.
}

\paragraph{Example (Julia).} Modify the Legendre weight by $u=1+x^2$. In the
orthonormal Legendre basis, the Gram matrix of the new measure is the leading
section of $u(J)$. To get that section exactly, compute $J^2$ from a padded $J$
and then truncate. Its Cholesky factor is the connection matrix: upper
triangular with bandwidth $\deg u=2$ and a positive diagonal.
\begin{verbatim}
using LinearAlgebra
N = 8
b = [k / sqrt(4k^2 - 1) for k in 1:N+1]   # orthonormal Legendre, padded
J = SymTridiagonal(zeros(N+2), b)
W = (I + J^2)[1:N, 1:N]                   # Gram of P for (1+x^2) dx
R = cholesky(Symmetric(W)).U
maximum(abs(R[i, j]) for i in 1:N, j in 1:N if j > i + 2)   # 0.0
\end{verbatim}

\begin{supporting}[Supporting results: triangular Gram factors]
\begin{theorem}[Gram matrices in orthonormal coordinates]
\label{thm:mathlib-gram-factor}
\lean{Matrix.gram, Matrix.gram_eq_conjTranspose_mul}\mathlibok
For vectors $P_j$ in a finite-dimensional real inner-product space and an
orthonormal basis $Q$, set $R_{ij}=\langle Q_i,P_j\rangle$. Then
$G(P)=R^\top R$. This identity is independent of polynomial variables or ordering.
\end{theorem}
\begin{proof}
\mathlibok
\end{proof}

\begin{theorem}[Positive triangular Gram factors]
\label{thm:positive-gram-factor}\uses{thm:mathlib-gram-factor,def:polynomial-linear-maps}
\lean{InnerProductSpace.inner_gramSchmidtNormed_self, InnerProductSpace.exists_orthonormalBasis_pos_upperTriangular, InnerProductSpace.gram_eq_transpose_toMatrix_mul, InnerProductSpace.exists_pos_upperTriangular_gram_factor, Matrix.eq_of_pos_upperTriangular_gram_eq}\leanok
Every ordered basis of a finite-dimensional real inner-product space admits
orthonormalization with upper triangular coordinate matrix $R$ and strictly
positive diagonal. It satisfies $G(P)=R^\top R$. Among upper triangular matrices
with positive diagonal, this factor is unique.

Reference: \cite[Section 4.2]{golub-van-loan}.
\end{theorem}
\begin{proof}
\leanok
\end{proof}

\end{supporting}

Repeated polynomial modifications correspond to compositions of the
connection maps. In the next chapter, integer parameter shifts of the
semiclassical Jacobi family give examples of these modifications, together
with sparse differentiation identities.

\chapter{Semiclassical Jacobi polynomials}\label{chap:semiclassical}
\rvsummary{Semiclassical Jacobi families for $x^a(1-x)^b(t-x)^c$, with sparse conversion, differentiation, and Cholesky/QR hierarchy construction.}

The weight $x^a(1-x)^b(t-x)^c$ on $(0,1)$, with $t>1$, adds a third factor to
the Jacobi weight, with the extra root $x=t$ outside the interval. The family
retains sparse conversion and differentiation.
Integer parameter shifts are polynomial modifications, so conversion is banded
by the previous chapter. A Pearson identity $(\sigma w)'=\tau w$ makes
differentiation sparse. For simultaneously raised parameters and positive
orthonormal scaling, weighted differentiation is the negative transpose of the
ordinary derivative matrix. The partially weighted variant has its own
adjoint identity and support pattern, stated separately below.
Cholesky and QR factorizations construct the parameter hierarchy, with exact
finite Jacobi updates in Section~\ref{sec:semiclassical-hierarchy}.
These families are used in spectral methods on annuli (Chapter~\ref{chap:annulus}).

\begin{definition}[Semiclassical Jacobi weights and measures]
\label{def:semiclassical-jacobi-measure}\uses{def:open-interval-weight-measure}
\lean{Polynomial.semiclassicalJacobiWeight, Polynomial.semiclassicalJacobiMeasure}\leanok
For $t\in(1,\infty)$ and arbitrary real $a,b,c$, define the continuous
nonnegative weight $w_{t,a,b,c}(x)=x^a(1-x)^b(t-x)^c$ on $(0,1)$.
Define $\mu_{t,a,b,c}$ by this density on the open interval, pushed forward
to $\R$ by inclusion. The measure need not be finite.

Reference: \cite[Section 3]{pgso}.
\end{definition}



\begin{supporting}[Reference details: Semiclassical density characterisation and related results]
\begin{theorem}[Semiclassical density characterisation]
\label{thm:semiclassical-jacobi-weight-characterisation}
\uses{def:semiclassical-jacobi-measure}
\lean{Polynomial.semiclassicalJacobiWeight_apply}\leanok
At every $x\in(0,1)$, $w(x)=x^\alpha(1-x)^\beta(t-x)^\gamma$, with all three power bases positive.
\end{theorem}
\begin{proof}
\leanok
\end{proof}



\begin{theorem}[Semiclassical integral characterisation]
\label{thm:semiclassical-jacobi-measure-characterisation}
\uses{def:semiclassical-jacobi-measure,thm:open-interval-weight-measure}
\lean{Polynomial.integrable_semiclassicalJacobiMeasure_iff, Polynomial.integral_semiclassicalJacobiMeasure}\leanok
A real function $f$ is $\mu_{t,\alpha,\beta,\gamma}$-integrable if and only if $x\mapsto w(x)f(x)$ is Lebesgue-integrable on $(0,1)$. When these equivalent integrability conditions hold,
\[\int f\,d\mu_{t,\alpha,\beta,\gamma}=\int_{x\in(0,1)}w(x)f(x)\,dx.\]
\end{theorem}
\begin{proof}
\leanok
\end{proof}

\begin{theorem}[Semiclassical Jacobi measure foundation]
\label{thm:semiclassical-jacobi-measure}\uses{thm:semiclassical-jacobi-measure-characterisation,thm:general-existence,thm:mathlib-power-integral}

\lean{Polynomial.integrable_semiclassicalJacobiWeight, Polynomial.integrable_semiclassicalJacobiMeasure, Polynomial.semiclassicalJacobiMeasure_foundation}\leanok
For $t\in(1,\infty)$, $\alpha,\beta>-1$ and $\gamma\in\R$, the classical density is integrable on $(0,1)$. Every real-valued function continuous on $[0,1]$ is integrable for its measure. The measure has finite mass and all polynomial moments, and the polynomial pairing is positive definite. There exists an orthogonal polynomial sequence of exact degrees. No normalization is fixed by this existence assertion.
Reference: \cite[Section 3]{pgso}.
\end{theorem}
\begin{proof}
\leanok
\end{proof}
\end{supporting}


\begin{definition}[Positive orthonormal semiclassical Jacobi coordinates]
\label{def:semiclassical-jacobi-sequence}\uses{thm:semiclassical-jacobi-measure,thm:even-positive-sequences}
\lean{Polynomial.semiclassicalJacobiSequence}\leanok
For $t>1$, $\alpha,\beta>-1$ and arbitrary real $\gamma$, let $(Q_n)$ be a
degree-indexed polynomial sequence chosen to be orthonormal for
$x^\alpha(1-x)^\beta(t-x)^\gamma\,dx$ on $(0,1)$, with positive leading
coefficients; such a sequence exists by Theorem~\ref{thm:semiclassical-jacobi-measure}.

Reference: \cite[Section 3]{pgso}.
\end{definition}

\begin{theorem}[Positive orthonormal semiclassical Jacobi sequence]
\label{thm:semiclassical-jacobi-orthonormality}\uses{def:semiclassical-jacobi-sequence}
\lean{Polynomial.semiclassicalJacobiSequence_spec}\leanok
For $t>1$, $\alpha,\beta>-1$ and $\gamma\in\R$, the chosen sequence $(Q_n)$
is orthogonal for the semiclassical Jacobi measure, has positive leading
coefficients, and satisfies $\int Q_n^2\,d\mu=1$ for every $n$.

Reference: \cite[Section 3]{pgso}.
\end{theorem}
\begin{proof}
\leanok
\end{proof}

\begin{supporting}[Reference details: Semiclassical Jacobi polynomial basis and related results]
\begin{definition}[Semiclassical Jacobi polynomial basis]
\label{def:semiclassical-jacobi-basis}\uses{thm:semiclassical-jacobi-orthonormality}
\lean{Polynomial.semiclassicalJacobiBasis}\leanok
For $t>1$, $\alpha,\beta>-1$ and $\gamma\in\R$, let $B$ be the polynomial
basis whose $n$th member is $Q_n$.

Reference: \cite[Section 3]{pgso}.
\end{definition}

\begin{theorem}[Orthonormality and coordinates]
\label{thm:semiclassical-jacobi-sequence-basic}\uses{def:semiclassical-jacobi-basis,thm:general-coordinates}
\lean{Polynomial.semiclassicalJacobiBasis_repr}\leanok
The sequence $(Q_n)$ is orthogonal for the semiclassical Jacobi measure, its
members have positive leading coefficients, and each has unit squared norm.
The $n$th coordinate of a polynomial $p$ in the basis $B$ is the pairing
$\langle p,Q_n\rangle$ for this measure.

Reference: \cite[Section 3]{pgso}.
\end{theorem}
\begin{proof}
\leanok
\end{proof}
\end{supporting}


\begin{theorem}[Semiclassical Jacobi parameter raising]
\label{thm:semiclassical-jacobi-shifts}\uses{def:semiclassical-jacobi-measure}
\lean{Polynomial.semiclassicalJacobiWeight_add_nat, Polynomial.semiclassicalJacobiMeasure_add_nat, Polynomial.semiclassicalJacobiWeight_add_one}\leanok
For $t>1$ and nonnegative integers $k,l,m$, increasing $(a,b,c)$ by $(k,l,m)$
is exactly multiplication of the measure by $x^k(1-x)^l(t-x)^m$.
No integrability restrictions on the exponents are needed for this measure identity.

Reference: \cite[Section 3]{pgso}.
\end{theorem}
\begin{proof}
\leanok
\end{proof}

\begin{theorem}[Semiclassical integer conversion and Gram factors]
\label{thm:semiclassical-conversion}\uses{thm:semiclassical-jacobi-shifts,thm:semiclassical-jacobi-sequence-basic,thm:univariate-connection-gram}
\lean{Polynomial.semiclassicalJacobiConnection, Polynomial.semiclassicalJacobiConnection_eq_zero, Polynomial.semiclassicalJacobiConnection_upperTriangular, Polynomial.semiclassicalJacobiConnection_diag_pos, Polynomial.semiclassicalJacobiConnection_gram}\leanok
Raising $(\alpha,\beta,\gamma)$ by nonnegative integers $(a,b,c)$ gives an upper
connection of bandwidth at most $a+b+c$. Its leading finite sections have positive
diagonal and factor the target-measure Gram matrix as $R^{\mathsf T}R$.
All coefficients are coordinates in the actual normalized polynomial bases.

References: \cite[Theorem 3.2]{pgso}; \cite[Lemma 2.8]{gutleb-olver-slevinsky}.
\end{theorem}
\begin{proof}
\leanok
\end{proof}

\begin{theorem}[Semiclassical Jacobi Pearson identity]
\label{thm:semiclassical-jacobi-pearson}\uses{thm:semiclassical-jacobi-measure,thm:semiclassical-jacobi-shifts,def:polynomial-linear-maps,thm:open-domain-derivative}
\lean{Polynomial.semiclassicalJacobiSigma, Polynomial.semiclassicalJacobiTau, Polynomial.hasDerivAt_semiclassicalJacobiWeight_add_one, Polynomial.exists_continuous_semiclassicalJacobiWeight_add_one, Polynomial.tendsto_semiclassicalJacobiWeight_add_one_zero, Polynomial.tendsto_semiclassicalJacobiWeight_add_one_one, Polynomial.integral_semiclassicalJacobi_pearson}\leanok
Set $\sigma=x(1-x)(t-x)$ and
$\tau=(a+1)(1-x)(t-x)-(b+1)x(t-x)-(c+1)x(1-x)$.
On $(0,1)$, $(\sigma w)'=\tau w$. For $t>1$, $a,b>-1$, $c\in\R$,
the raised weight extends continuously to $[0,1]$, with value zero at both
endpoints. For the one-sided limit at zero only $a>-1$ is needed, while
for the limit at one only $b>-1$ is needed (always with $t>1$).
Consequently every polynomial $p$ satisfies
$\int(\sigma p'+\tau p)\,d\mu_{t,a,b,c}=0$.
These polynomials form a Pearson pair.
\end{theorem}
\begin{proof}
\leanok
\end{proof}

\rvnote{thm:semiclassical-jacobi-pearson}{This is a Pearson pair; no claim of minimal semiclassical class is made.}


\begin{theorem}[Sparse semiclassical differentiation]
\label{thm:semiclassical-differentiation}\uses{thm:semiclassical-jacobi-pearson,thm:semiclassical-jacobi-sequence-basic,thm:operator-coefficients}
\lean{Polynomial.integralInner_semiclassicalJacobi_derivative, Polynomial.semiclassicalJacobiDerivative, Polynomial.semiclassicalJacobiDerivative_eq_zero, Polynomial.semiclassicalJacobi_derivative_coefficientMap_exact, Polynomial.semiclassicalJacobi_derivative_coefficientMap_apply}\leanok
Differentiation from parameters $(\alpha,\beta,\gamma)$ to
$(\alpha+1,\beta+1,\gamma+1)$ has entries only when $j<n\le j+2$, for row $j$
and column $n$. Its action on finitely supported coefficients reconstructs the
ordinary polynomial derivative exactly. The underlying adjoint identity is
\[
 \langle p',q\rangle_{\mu_+}
 =-\langle p,\sigma q'+\tau q\rangle_\mu.
\]
No endpoint regularity beyond $\alpha,\beta>-1$ is imposed.

Reference: \cite[Theorem 3.4]{pgso}.
\end{theorem}
\begin{proof}
\leanok
\end{proof}

\begin{theorem}[Weighted semiclassical differentiation]
\label{thm:semiclassical-weighted-differentiation}\uses{thm:semiclassical-differentiation}
\lean{Polynomial.semiclassicalJacobiWeightedDerivativeMap, Polynomial.semiclassicalJacobiWeightedDerivative, Polynomial.semiclassicalJacobiWeightedDerivative_eq_neg_transpose, Polynomial.semiclassicalJacobiWeightedDerivative_eq_zero, Polynomial.semiclassicalJacobi_weightedDerivative_coefficientMap_exact}\leanok
For $t>1$, $\alpha,\beta>-1$, $\gamma\in\mathbb R$, weighted differentiation
from the simultaneously raised family to the original family is the polynomial
map $q\mapsto\sigma q'+\tau q$. With the positive orthonormal scaling, its
matrix is $-D^{\mathsf T}$, where $D$ is the unweighted derivative matrix.
Column $j$ is supported only in rows $j+1,j+2$. Its coefficient action is exact
on all finite expansions.

Reference: \cite[Section 3.4]{pgso}.
\end{theorem}
\begin{proof}
\leanok
\end{proof}

\begin{supporting}[Reference details: Partial weighted semiclassical differentiation]
\begin{theorem}[Partial weighted semiclassical differentiation]
\label{thm:semiclassical-partial-weighted-differentiation}\uses{thm:semiclassical-differentiation}
\lean{Polynomial.semiclassicalJacobiCWeightedDerivativeMap, Polynomial.integralInner_semiclassicalJacobi_cWeightedDerivative, Polynomial.semiclassicalJacobiCWeightedDerivative, Polynomial.semiclassicalJacobiCWeightedDerivative_eq_zero, Polynomial.semiclassicalJacobi_cWeightedDerivative_coefficientMap_exact}\leanok
For $t>1$, $a,b>-1$, and real $\gamma$, the map
$D_c p=(t-x)p'-(\gamma+1)p$ from parameters $(a,b,\gamma+1)$ to
$(a+1,b+1,\gamma)$ has support only in rows $n-1,n$ of column $n$.
Its coefficient action is exact on finite polynomial expansions. Its adjoint identity is
\[
 \langle D_cp,z\rangle_{a+1,b+1,\gamma}
 =-\langle p,x(1-x)z'+[(a+1)(1-x)-(b+1)x]z\rangle_{a,b,\gamma+1}.
\]

Reference: \cite[Section 3.4]{pgso}.
\end{theorem}
\begin{proof}
\leanok
\end{proof}
\end{supporting}


One-dimensional families also give separated bases in several variables.
The next chapter treats product domains, triangles and disks. On the annulus,
the radial factor will be a semiclassical Jacobi polynomial.

\section{Constructing the parameter hierarchy}\label{sec:semiclassical-hierarchy}
Fix $t>1$, $\alpha,\beta>-1$ and $\gamma\in\R$. Write $p_j^{(\gamma)}$ for the
positive-leading orthonormal semiclassical Jacobi polynomial with these parameters.
Changing $\gamma$ gives two ways to construct the next family: Cholesky for a
unit shift and QR for a two-step shift. The finite formulas below include the
extra mode produced by multiplication by $x$.

\begin{definition}[Jacobi matrices and raising sections]
\label{def:semiclassical-hierarchy-matrices}\uses{thm:semiclassical-conversion}
\lean{Polynomial.semiclassicalJacobiMatrix, Polynomial.semiclassicalJacobiRaisingSection}\leanok
Let $X^{(\gamma)}_{ij}=[p_i^{(\gamma)}](x p_j^{(\gamma)})$.
For $c,N\in\N$, let $R_{c,N}$ be the leading $N\times N$ section of the
connection from $p^{(\gamma)}$ to $p^{(\gamma+c)}$:
$(R_{c,N})_{ij}=[p_i^{(\gamma+c)}]p_j^{(\gamma)}$, $0\le i,j<N$.
Here $[p_i]q$ denotes the coordinate of $q$ in the indicated polynomial basis.
\end{definition}

\begin{supporting}[Structure and Gram identities]
\begin{theorem}[Structure of the hierarchy matrices]
\label{thm:semiclassical-hierarchy-matrices}\uses{def:semiclassical-hierarchy-matrices,thm:general-multiplication-matrix}
\lean{Polynomial.semiclassicalJacobiMatrix_isBanded, Polynomial.semiclassicalJacobiMatrix_symmetric, Polynomial.semiclassicalJacobiRaisingSection_triangular, Polynomial.semiclassicalJacobiRaisingSection_isUnit, Polynomial.semiclassicalJacobiRaisingSection_gram}\leanok
$X^{(\gamma)}$ is symmetric and tridiagonal. Each $R_{c,N}$ is upper triangular,
with positive diagonal, and invertible. Moreover,
\[
 [\,[p_i^{(\gamma)}]((t-x)^c p_j^{(\gamma)})\,]_{i,j<N}
   =R_{c,N}^{\mathsf T}R_{c,N}.
\]
The multiplication on the left is performed on polynomials before taking the section.
\end{theorem}
\begin{proof}
\leanok
\end{proof}
\end{supporting}

\begin{theorem}[Cholesky construction of a unit shift]
\label{thm:semiclassical-cholesky-shift}\uses{thm:semiclassical-hierarchy-matrices}
\lean{Polynomial.semiclassicalJacobiRaisingSection_cholesky, Polynomial.semiclassicalJacobiRaisingSection_cholesky_unique}\leanok
For every $N\in\N$,
\[
 tI_N-X^{(\gamma)}_N=R_{1,N}^{\mathsf T}R_{1,N}.
\]
Thus $R_{1,N}$ is the unique upper triangular factor with positive diagonal.
Reference: \cite[Section 3.1]{pgso}.
\end{theorem}
\begin{proof}
\leanok
\end{proof}

\begin{theorem}[Padded Jacobi update]
\label{thm:semiclassical-padded-update}\uses{thm:semiclassical-hierarchy-matrices}
\lean{Polynomial.semiclassicalJacobi_padded_intertwining, Polynomial.semiclassicalJacobi_padded_update}\leanok
For any other real parameter $\gamma'$, write $R$ for the connection from
$p^{(\gamma)}$ to $p^{(\gamma')}$. With indices starting at zero,
\[
 R_{[0,N),[0,N+1)}X^{(\gamma)}_{[0,N+1),[0,N)}
   =X^{(\gamma')}_N R_{[0,N),[0,N)}.
\]
In particular, when $\gamma'=\gamma+c$, $c\in\N$, there is a matrix $S$ with
$R_{c,N}S=SR_{c,N}=I_N$, and
\[
 X^{(\gamma+c)}_N=
 R_{[0,N),[0,N+1)}X^{(\gamma)}_{[0,N+1),[0,N)}S.
\]
Reference: \cite[Sections 3.2 and 3.5]{pgso}.
\end{theorem}
\begin{proof}
\leanok
\end{proof}

\begin{definition}[Rectangular factors for a two-step shift]
\label{def:semiclassical-qr}\uses{def:semiclassical-jacobi-basis}
\lean{Polynomial.semiclassicalJacobiQRSection, Polynomial.semiclassicalJacobiShiftSection}\leanok
For $0\le i<N+1$ and $0\le j<N$, define
\[
 (A_N)_{ij}=[p_i^{(\gamma)}]((t-x)p_j^{(\gamma)}),\qquad
 (Q_N)_{ij}=[p_i^{(\gamma)}]((t-x)p_j^{(\gamma+2)}).
\]
\end{definition}

\begin{theorem}[Semiclassical Jacobi hierarchy construction]
\label{thm:semiclassical-hierarchies}\uses{def:semiclassical-qr,def:semiclassical-hierarchy-matrices}
\lean{Polynomial.semiclassicalJacobiShiftSection_eq_QR, Polynomial.semiclassicalJacobiQRSection_orthogonal, Polynomial.semiclassicalJacobiQRSection_update}\leanok
The two-step raising section is the positive-diagonal triangular factor in the
rectangular QR factorization
\[
 A_N=Q_NR_{2,N},\qquad Q_N^{\mathsf T}Q_N=I_N.
\]
The twice-raised Jacobi matrix is obtained by the exact padded compression
\[
 X^{(\gamma+2)}_N=Q_N^{\mathsf T}X^{(\gamma)}_{N+1}Q_N.
\]
Together with the unit-shift Cholesky factorization, these identities construct
successive levels of the semiclassical Jacobi parameter hierarchy.
Reference: \cite[Sections 3.1, 3.2 and 3.5]{pgso}.
\end{theorem}
\begin{proof}
\leanok
\end{proof}


\chapter{Classical orthogonal polynomials in several variables}\label{chap:multivariate}
\rvsummary{Classical orthogonal polynomials on squares, triangles and disks, built by separation of variables.}

On squares, triangles and disks, separation of variables constructs
orthogonal polynomials from one-dimensional families. The resulting bases
have explicit norms and span each total-degree space, with sparse
multiplication and parameter conversion. We begin with the multivariate
version of the polynomial inner product from Chapter~\ref{chap:general}.

\begin{definition}[Finite multivariate polynomial moments]
\label{def:multivariate-moments}
\lean{MvPolynomial.HasFiniteMoments}\leanok
A measure $\mu$ on $\R^\sigma$ has finite polynomial moments when every
polynomial square is $\mu$-integrable.
\end{definition}

\begin{definition}[Positive-definite multivariate polynomial integral]
\label{def:multivariate-positive}\uses{def:multivariate-moments}
\lean{MvPolynomial.HasPositiveDefiniteIntegral}\leanok
The polynomial integral is positive definite when every polynomial square
is integrable and every nonzero polynomial has strictly positive squared integral.
\end{definition}

\begin{supporting}[Reference details: Characterisation of positive multivariate polynomial integrals and related results]
\begin{theorem}[Characterisation of positive multivariate polynomial integrals]
\label{thm:multivariate-positive-characterisation}\uses{def:multivariate-positive}
\lean{MvPolynomial.HasPositiveDefiniteIntegral.hasFiniteMoments,MvPolynomial.HasPositiveDefiniteIntegral.integral_sq_pos}\leanok
A positive-definite polynomial integral has all polynomial moments finite.
Every nonzero polynomial has strictly positive squared integral.
\end{theorem}
\begin{proof}
\leanok
\end{proof}

\begin{definition}[Integrable multivariate polynomial pairs]
\label{def:multivariate-integrable-pair}
\lean{MvPolynomial.IntegrablePair, MvPolynomial.integralInner}\leanok
For a measure $\mu$ on $\R^\sigma$, an integrable polynomial pair consists of
polynomials $p,q$ together with $pq\in L^1(\mu)$. Its pairing is the finite
integral $\langle p,q\rangle_\mu=\int pq\,d\mu$, defined on these pairs only.
\end{definition}

\begin{theorem}[Characterisation of the multivariate integrable pairing]
\label{thm:multivariate-integrable-pair}\uses{def:multivariate-integrable-pair}
\lean{MvPolynomial.integralInner_eq_integral, MvPolynomial.integralInner_comm, MvPolynomial.integralInner_self}\leanok
An integrable pair has pairing $\int pq\,d\mu$. The pairing is symmetric;
on an integrable diagonal pair it equals $\int p^2\,d\mu$.
\end{theorem}
\begin{proof}
\leanok
\end{proof}

\begin{theorem}[Integrable products and finite mass]
\label{thm:multivariate-moments-basic}\uses{def:multivariate-moments}
\lean{MvPolynomial.HasFiniteMoments.integrable_mul, MvPolynomial.HasFiniteMoments.isFiniteMeasure}\leanok
If every polynomial square is $\mu$-integrable, then $\mu$ is a finite measure,
and, for finite $\sigma$, every product $pq$ of polynomials is
$\mu$-integrable.
\end{theorem}
\begin{proof}
\leanok
\end{proof}

\begin{definition}[Multivariate polynomial measure inner products]
\label{def:multivariate-degree-spaces}\uses{thm:multivariate-moments-basic,def:multivariate-positive,thm:multivariate-integrable-pair}
\lean{MvPolynomial.HasFiniteMoments.pair, MvPolynomial.integralInnerProductCore}\leanok
For finite $\sigma$, assume explicitly that $\mu$ has finite polynomial moments
and positive-definite polynomial integral. Equip the polynomial space with the inner product $\langle p,q\rangle_\mu=\int pq\,d\mu$. There is no connectedness or product-support restriction.
\end{definition}
\end{supporting}


\begin{supporting}[Supporting results: total-degree spaces]
\begin{definition}[Finite total-degree spaces]
\label{def:mathlib-total-degree}
\lean{MvPolynomial.restrictTotalDegree}\mathlibok
For a variable set $\sigma$ and $N\in\N$, consider the submodule of polynomials
of total degree at most $N$.
\end{definition}

\begin{theorem}[Total-degree spaces are finite-dimensional]
\label{thm:mathlib-total-degree-finite}\uses{def:mathlib-total-degree}
\lean{MvPolynomial.instFiniteSubtypeMemSubmoduleRestrictTotalDegreeOfFinite}\mathlibok
For a finite variable set $\sigma$, the space of real polynomials of total
degree at most $N$ is finite-dimensional.
\end{theorem}
\begin{proof}
\mathlibok
\end{proof}

\end{supporting}

\section{Tensor-product domains}\label{sec:tensor}
\rvsummary{Product bases on tensor-product domains and their coordinate-multiplication structure.}

Products of univariate orthogonal families are orthogonal for product measures.
Multiplication by a coordinate acts on one factor only, so the univariate
recurrences determine the multivariate multiplication matrices. Total degree
organises the products into finite polynomial spaces, while the product
measure makes their pairings factor into one-dimensional pairings.

\begin{supporting}[Reference details: Product moments, bases and strict positivity]
\begin{theorem}[Product moments, bases and strict positivity]
\label{thm:tensor-product-bases}\uses{thm:tensor-product-pairings,thm:general-existence,thm:mathlib-total-degree-finite,def:multivariate-degree-spaces}
\lean{MvPolynomial.hasFiniteMoments_pi, MvPolynomial.tensorProductPolynomial_linearIndependent, MvPolynomial.tensorProductDegreeBasis, MvPolynomial.hasPositiveDefiniteIntegral_pi}\leanok
Finite polynomial moments are preserved under finite products. Orthogonal coordinate sequences with positive squared norms form a basis of every complete multivariate total-degree space. Strict positivity of every coordinate polynomial pairing implies strict positivity of the product pairing. These results impose no interval or compact-support assumption.
\end{theorem}
\begin{proof}
\leanok
\end{proof}
\end{supporting}


\begin{theorem}[Coordinate multiplication on tensor bases]
\label{thm:tensor-product-operators}\uses{thm:orthogonal-moments,thm:tensor-product-bases,thm:general-tridiagonal}
\lean{MvPolynomial.tensorProductPolynomial_update, MvPolynomial.X_mul_tensorProductPolynomial, MvPolynomial.integralInner_X_mul_tensorProductPolynomial, MvPolynomial.integralInner_X_mul_tensorProductPolynomial_eq_zero_of_other_index, MvPolynomial.integralInner_X_mul_tensorProductPolynomial_eq_zero_of_gap}\leanok
Multiplication by coordinate $x_i$ is exactly multiplication by $x$ in factor $i$. For $\sigma$-finite factor measures, pairing $x_i$ times a product with another product gives the product of the univariate pairings, with $x$ acting in factor $i$. If each coordinate sequence is orthogonal for its factor measure, these pairing entries vanish when any other coordinate degree changes; if each factor measure also has finite polynomial moments, they vanish when the $i$th degree changes by more than one.
\end{theorem}
\begin{proof}
\leanok
\end{proof}

Orthogonal polynomials on tensor-product domains: see Theorems~\ref{thm:tensor-product-bases}, \ref{thm:tensor-product-operators} and~\ref{thm:tensor-product-coordinates}.

\begin{supporting}[Supporting results: factored pairings and coordinates]
\begin{theorem}[Coordinate products and factored pairings]
\label{thm:tensor-product-pairings}\uses{def:orthogonal-sequence,def:polynomial-integrable-pair,thm:disk-polynomial-lift,def:multivariate-integrable-pair}
\lean{MvPolynomial.tensorProductPolynomial, MvPolynomial.eval_tensorProductPolynomial, MvPolynomial.totalDegree_tensorProductPolynomial_le, MvPolynomial.integralInner_tensorProductPolynomial, MvPolynomial.integralInner_tensorProductPolynomial_eq_zero, MvPolynomial.integralInner_tensorProductPolynomial_self}\leanok
For a finite coordinate set, the product of univariate polynomial factors is an actual multivariate polynomial of degree at most the sum of their degrees. Under a product of sigma-finite measures, its pairing with another coordinate product is the product of the coordinate pairings. Orthogonal coordinate sequences give orthogonal multiindices; normalized factors give normalized products.
\end{theorem}
\begin{proof}
\leanok
\end{proof}

\begin{theorem}[Normalized tensor coordinates]
\label{thm:tensor-product-coordinates}\uses{thm:tensor-product-bases,thm:measure-conversion-cholesky}
\lean{MvPolynomial.integralGram_tensorProductDegreeBasis, MvPolynomial.tensorProductDegreeBasis_repr}\leanok
If each coordinate family is normalized, every complete tensor degree basis has identity Gram matrix and its coordinates are the exact polynomial integral projections.
\end{theorem}
\begin{proof}
\leanok
\end{proof}

\end{supporting}

\section{The triangle: Proriol polynomials}\label{sec:triangle}
\rvsummary{Proriol polynomials on the triangle: orthogonality, complete bases and sparse operators.}

On the triangle the Proriol polynomials couple an outer Jacobi
polynomial in $x$ to an inner one in $y/x$, scaled by $x^k$ so the result is a
polynomial. Coordinate multiplication couples only neighbouring degrees, and
raising any weight parameter by an integer gives a sparse conversion with banded
inner blocks.

\begin{definition}[Proriol polynomials on the triangle]
\label{def:proriol-triangle}\uses{def:scaled-koornwinder-lift,def:unit-jacobi-measure}
\lean{MvPolynomial.proriol}\leanok
For real $\alpha,\beta,\gamma$, write $\widetilde P_j^{(a,b)}(x)=P_j^{(a,b)}(2x-1)$,
and for a polynomial $q(z)=\sum_j c_jz^j$ of degree at most $k$ let
$x^kq(y/x)$ denote the polynomial $\sum_j c_jy^jx^{k-j}$. Define
\[
 Q_{n,k}^{\alpha,\beta,\gamma}(x,y)=
 \widetilde P_{n-k}^{(\alpha,\beta+\gamma+2k+1)}(x)\;
 x^k\widetilde P_k^{(\beta,\gamma)}(y/x),
\]
used for $0\le k\le n$.

References: \cite[18.37.7]{dlmf}; \cite[Section 2.4]{dunkl-xu}.
\end{definition}

\begin{theorem}[Separated form and degree]
\label{thm:proriol-triangle-basic}\uses{def:proriol-triangle,thm:scaled-lift-basic,thm:jacobi-degree}
\lean{MvPolynomial.eval_proriol, MvPolynomial.eval_proriol_scaled, MvPolynomial.totalDegree_proriol_le}\leanok
For $x\ne0$,
\[
 Q_{n,k}^{\alpha,\beta,\gamma}(x,y)=
 P_{n-k}^{(\alpha,\beta+\gamma+2k+1)}(2x-1)\,x^kP_k^{(\beta,\gamma)}(2y/x-1),
\]
the triangle polynomial of \cite[18.37.7]{dlmf} on $0<y<x<1$. For all real
$x,z$, including $x=0$,
$Q_{n,k}^{\alpha,\beta,\gamma}(x,xz)=\widetilde P_{n-k}^{(\alpha,\beta+\gamma+2k+1)}(x)\,
x^k\,\widetilde P_k^{(\beta,\gamma)}(z)$. If $k\le n$, the polynomial has total
degree at most $n$.

References: \cite[18.37.7]{dlmf}; \cite[Section 2.4]{dunkl-xu}.
\end{theorem}
\begin{proof}
\leanok
\end{proof}

\begin{theorem}[Proriol orthogonality and squared norms]
\label{thm:proriol-triangle-orthogonality}\uses{thm:proriol-triangle-measure,thm:proriol-triangle-basic,thm:unit-jacobi-measure,thm:multivariate-integrable-pair,thm:jacobi-growth-constants}
\lean{MvPolynomial.integralInner_proriolMeasure_eq_zero, MvPolynomial.integral_sq_proriolMeasure, MvPolynomial.integral_sq_proriolMeasure_pos}\leanok
For $\alpha,\beta,\gamma>-1$, distinct valid triangular indices give orthogonal
Proriol polynomials for the standard triangle measure. Write $h_j^{(a,b)}$
for the standard Jacobi squared norm. Then
\[
 \|Q_{n,k}^{\alpha,\beta,\gamma}\|^2
 =2^{-c-2k}h_{n-k}^{(\alpha,\beta+\gamma+2k+1)}h_k^{(\beta,\gamma)}>0,
 \qquad c=\alpha+2\beta+2\gamma+3.
\]
The product formula includes degree zero without introducing singular Gamma quotients.

Reference: \cite[18.37.8]{dlmf}.
\end{theorem}
\begin{proof}
\leanok
\end{proof}

\begin{theorem}[Complete Proriol bases and integral coordinates]
\label{thm:proriol-triangle-bases}\uses{thm:proriol-triangle-orthogonality,thm:mathlib-total-degree-finite,def:multivariate-degree-spaces}
\lean{MvPolynomial.proriolDegreeBasis, MvPolynomial.proriolDegreeBasis_apply, MvPolynomial.proriolDegreeBasis_repr, MvPolynomial.hasPositiveDefiniteIntegral_proriolMeasure}\leanok
For $\alpha,\beta,\gamma>-1$, the polynomials $Q_{n,k}$ with $0\le k\le n\le N$
form a basis of the complete total-degree-$N$ space. Coordinates are
$\langle p,Q_{n,k}\rangle/\|Q_{n,k}\|^2$ in the standard triangle measure.
This measure has finite polynomial moments and a strictly positive polynomial pairing.

Reference: \cite[Section 2.4]{dunkl-xu}.
\end{theorem}
\begin{proof}
\leanok
\end{proof}

\rvnote{thm:proriol-triangle-bases}{
This measure has finite polynomial moments and a strictly positive polynomial pairing, supplying the hypotheses of the general Gram/Cholesky construction.
}

\begin{supporting}[Reference details: Classical Proriol multiplication entry and related results]
\begin{definition}[Classical Proriol multiplication entry]
\label{def:proriol-multiplication-entry}\uses{def:multivariate-integrable-pair,def:proriol-triangle,thm:proriol-triangle-measure}
\lean{MvPolynomial.ProriolMultiplicationInput, MvPolynomial.proriolMultiplicationEntry}\leanok
For real parameters and natural indices, the multiplication-entry domain requires integrability of
$x_\ell Q_{n,k}Q_{m,j}$ and $Q_{m,j}^2$ for the Proriol measure, together with
a strictly positive target squared norm $h_{m,j}=\int Q_{m,j}^2\,d\mu$.
Define the multiplication entry by
$M^\ell_{(m,j),(n,k)}=(\int x_\ell Q_{n,k}Q_{m,j}\,d\mu)/h_{m,j}$.
This local domain does not require positivity for other polynomial pairings.
\end{definition}

\begin{theorem}[Proriol multiplication-entry characterization]
\label{thm:proriol-multiplication-entry}\uses{def:proriol-multiplication-entry,thm:proriol-triangle-orthogonality}
\lean{MvPolynomial.ProriolMultiplicationInput.of_neg_one_lt, MvPolynomial.proriolMultiplicationEntry_mul_norm}\leanok
On the multiplication-entry domain,
\[
 h_{m,j}M^\ell_{(m,j),(n,k)}=\int x_\ell Q_{n,k}Q_{m,j}\,d\mu.
\]
The integrand is integrable and $h_{m,j}>0$. For
$\alpha,\beta,\gamma>-1$, these conditions hold for every choice of natural
indices and either coordinate.
\end{theorem}
\begin{proof}
\leanok
\end{proof}
\end{supporting}


\begin{theorem}[Sparse Proriol coordinate multiplication]
\label{thm:proriol-triangle-multiplication}\uses{thm:proriol-multiplication-entry,thm:proriol-triangle-bases,thm:general-tridiagonal}
\lean{MvPolynomial.proriolMultiplicationEntry_eq_repr, MvPolynomial.proriolMultiplicationEntry_eq_zero}\leanok
For valid indices and $\alpha,\beta,\gamma>-1$, the coefficient of $Q_{m,j}$
in either $xQ_{n,k}$ or $yQ_{n,k}$ vanishes when $|m-n|>1$.
For $x$ it also vanishes unless $j=k$; for $y$ it vanishes when $|j-k|>1$.
These are actual coordinates in any containing complete degree space, in the
normalization of \cite{dlmf}: both matrices are block tridiagonal, with diagonal inner
blocks for $x$ and tridiagonal inner blocks for $y$.

Reference: \cite[Section 2.2]{gutleb-olver}.
\end{theorem}
\begin{proof}
\leanok
\end{proof}

\begin{supporting}[Reference details: All integer Proriol parameter shifts]
\begin{theorem}[All integer Proriol parameter shifts]
\label{thm:proriol-all-shift-measures}\uses{thm:proriol-triangle-measure,def:multivariate-integrable-pair}
\lean{MvPolynomial.proriolShiftDensity, MvPolynomial.totalDegree_proriolShiftDensity_le, MvPolynomial.integral_proriolMeasure_add_nat, MvPolynomial.integralInner_proriolMeasure_add_nat}\leanok
For admissible $\alpha,\beta,\gamma$ and nonnegative integers $a,b,c$, raising
these parameters multiplies the weight by $u=(1-x)^a(x-y)^b y^c$.
This holds for every polynomial integral and Gram pairing; $\deg u\le a+b+c$.
\end{theorem}
\begin{proof}
\leanok
\end{proof}
\end{supporting}


\begin{theorem}[Sparse outer-parameter modification on the triangle]
\label{thm:proriol-alpha-conversion}\uses{thm:proriol-triangle-bases}
\lean{MvPolynomial.integral_proriolMeasure_alpha_add_nat, MvPolynomial.integralInner_proriolMeasure_alpha_add_nat, MvPolynomial.proriol_alpha_connection_eq_zero}\leanok
For $\alpha,\beta,\gamma>-1$ and an integer $r\ge0$, replacing $\alpha$ by
$\alpha+r$ multiplies the triangle weight by $(1-x)^r$ at the level of every
polynomial integral and Gram pairing. In the expansion of a source $Q_{n,k}$
in the new Proriol basis, the coefficient of the target $Q_{m,j}$ is zero unless
$j=k$ and $m\le n\le m+r$.
\end{theorem}
\begin{proof}
\leanok
\end{proof}

\begin{theorem}[Banded inner blocks for Proriol conversion]
\label{thm:proriol-all-shift-support}\uses{thm:proriol-all-shift-measures,thm:proriol-triangle-bases,thm:orthogonal-lower-degree}
\lean{Polynomial.integralInner_shiftedJacobi_add_nat_eq_zero, MvPolynomial.integralInner_proriol_shift_eq_zero, MvPolynomial.proriol_connection_eq_zero}\leanok
For nonnegative integer parameter shifts $(a,b,c)$ as in Theorem~\ref{thm:proriol-all-shift-measures}, the coefficient of the target $Q_{m,j}$ in a source
$Q_{n,k}$ vanishes unless
\[
 m\le n\le m+a+b+c,\qquad j\le k\le j+b+c.
\]
The statement concerns actual coordinates in every complete finite degree space.
\end{theorem}
\begin{proof}
\leanok
\end{proof}

\begin{supporting}[Supporting results: affine Jacobi measures and the triangle measure]




\begin{definition}[Classical affine Jacobi density]
\label{def:unit-jacobi-weight}
\uses{def:jacobi-measure}
\lean{Polynomial.unitJacobiWeight}\leanok
On the open interval $(0,1)$, set $\widehat w_{\alpha,\beta}(x)=2w_{\alpha,\beta}(2x-1)$.
\end{definition}

\begin{theorem}[Affine Jacobi density characterisation]
\label{thm:unit-jacobi-weight}\uses{def:unit-jacobi-weight}

\lean{Polynomial.unitJacobiWeight_apply}\leanok
For every $x\in(0,1)$, $\widehat w_{\alpha,\beta}(x)=2^{\alpha+\beta+1}(1-x)^\alpha x^\beta$.
\end{theorem}
\begin{proof}
\leanok
\end{proof}

\begin{definition}[Classical affine Jacobi measure and polynomials]
\label{def:unit-jacobi-measure}
\uses{def:jacobi-measure,def:jacobi}
\lean{Polynomial.unitJacobiMeasure, Polynomial.shiftedJacobi}\leanok
Let $\widehat\mu_{\alpha,\beta}$ be the pushforward of $\mu_{\alpha,\beta}$ by $t\mapsto(t+1)/2$. Set $\widehat P_n^{(\alpha,\beta)}(x)=P_n^{(\alpha,\beta)}(2x-1)$.
\end{definition}

\begin{theorem}[Affine Jacobi integral characterisation]
\label{thm:unit-jacobi-measure-characterisation}\uses{def:unit-jacobi-measure,thm:jacobi-measure-characterisation,def:unit-jacobi-weight}

\lean{Polynomial.integrable_unitJacobiWeight_mul, Polynomial.integral_unitJacobiMeasure_subtype}\leanok
For a measurable $f\in\mathcal I(\widehat\mu_{\alpha,\beta})$, the product $\widehat w_{\alpha,\beta}f$ is integrable on the open interval $(0,1)$, and
\[\int f\,d\widehat\mu_{\alpha,\beta}=\int_{x\in(0,1)}\widehat w_{\alpha,\beta}(x)f(x)\,dx.\]
\end{theorem}
\begin{proof}
\leanok
\end{proof}

\begin{theorem}[Affine Jacobi measures and radial shifts]
\label{thm:unit-jacobi-measure}\uses{def:unit-jacobi-measure,thm:integrable-input,thm:jacobi-pairwise}

\lean{Polynomial.integral_pow_unitJacobiMeasure_integrable, Polynomial.isOrthogonal_shiftedJacobiSequence}\leanok
The shifted Jacobi polynomials are orthogonal for $\widehat\mu_{\alpha,\beta}$ when $\alpha,\beta>-1$. For arbitrary real $\alpha,\beta$, $r\in\N$ and measurable $f\in\mathcal I(\widehat\mu_{\alpha,\beta+r})$, the function $x^rf(x)$ is integrable for $\widehat\mu_{\alpha,\beta}$ and
\[\int x^rf(x)\,d\widehat\mu_{\alpha,\beta}=2^{-r}\int f(x)\,d\widehat\mu_{\alpha,\beta+r}.\]
\end{theorem}
\begin{proof}
\leanok
\end{proof}

\begin{theorem}[The standard Proriol triangle measure]
\label{thm:proriol-triangle-measure}\uses{thm:koornwinder-moments,thm:multivariate-moments-basic,thm:unit-jacobi-measure-characterisation,thm:jacobi-moments,thm:unit-jacobi-weight}
\lean{MvPolynomial.proriolMeasure, MvPolynomial.integral_proriolProductMeasure_density, MvPolynomial.integral_proriolMeasure_density, MvPolynomial.hasFiniteMoments_proriolMeasure}\leanok
For $\alpha,\beta,\gamma>-1$, transport
$\widehat\mu_{\alpha,\beta+\gamma+1}\otimes\widehat\mu_{\beta,\gamma}$
by $(x,z)\mapsto(x,xz)$ and multiply by $2^{-c}$, where
$c=\alpha+2\beta+2\gamma+3$. Polynomial integrals in the resulting measure are exactly
\[
 \int_{x\in(0,1)}(1-x)^\alpha\!\int_{y\in(0,x)}(x-y)^\beta y^\gamma p(x,y)\,dy\,dx.
\]
Thus the auxiliary affine mass has been removed from the weight of \cite[18.37.8]{dlmf}.

The standard triangle measure has all polynomial moments.

Reference: \cite[18.37.8]{dlmf}.
\end{theorem}
\begin{proof}
\leanok
\end{proof}

\end{supporting}

\section{The disk: Dunkl--Xu and Zernike polynomials}\label{sec:disk}
\rvsummary{Cartesian Dunkl--Xu bases and their parameter conversions, followed by polar Zernike bases for angular separation.}

The disk has two classical bases. The Cartesian Dunkl--Xu polynomials separate
$x$ from $y/\sqrt{1-x^2}$; a parity lift keeps them polynomial, even on the
boundary. Raising the weight parameter multiplies the measure by $1-x^2-y^2$
and gives a four-term connection with explicit coefficients. The Zernike
polynomials instead separate radius and angle, which makes them natural for
the Laplacian and for annuli.

\begin{supporting}[Reference details: Polynomial lift for Cartesian disk families]
\begin{definition}[Polynomial lift for Cartesian disk families]
\label{def:disk-polynomial-lift}
\lean{MvPolynomial.diskLift}\leanok
For $p(x)=\sum_j c_jx^j$ define the bivariate polynomial
$L_np(x,y)=\sum_{j=0}^n c_j y^j(1-x^2)^{\lfloor(n-j)/2\rfloor}$.

\end{definition}

\rvnote{def:disk-polynomial-lift}{
Its evaluation on the disk, for $p$ of degree at most $n$ and parity $n$, is Theorem~\ref{thm:disk-polynomial-lift}.

The definition uses no division or square roots.
}
\end{supporting}


\begin{definition}[Cartesian Dunkl--Xu polynomials]
\label{def:dunkl-xu-polynomials}\uses{def:disk-polynomial-lift,def:jacobi}
\lean{MvPolynomial.dunklXuDisk}\leanok
Set $D_{n,k}^{\beta}(x,y)=P_{n-k}^{(\beta+k+1/2,\beta+k+1/2)}(x)
L_kP_k^{(\beta,\beta)}(x,y)$, using $0\le k\le n$ in the degree-$n$ block.
This is the Cartesian Jacobi normalization, with increasing $k$ within a block;
it is not yet orthonormally scaled.

Reference: \cite[Section 2.3]{dunkl-xu}.
\end{definition}

\begin{theorem}[Dunkl--Xu separated formula]
\label{thm:dunkl-xu-separated}\uses{def:dunkl-xu-polynomials,thm:disk-polynomial-lift,thm:jacobi-coefficient-parity,thm:jacobi-degree}
\lean{MvPolynomial.eval_dunklXuDisk, MvPolynomial.totalDegree_dunklXuDisk_le, MvPolynomial.dunklXuDisk_zero_zero, MvPolynomial.dunklXuDisk_zero_inner}\leanok
For $x^2<1$, put $s=\sqrt{1-x^2}$. Then
$D_{n,k}^{\beta}(x,y)=P_{n-k}^{(\beta+k+1/2,\beta+k+1/2)}(x)s^kP_k^{(\beta,\beta)}(y/s)$.
The defining polynomial has total degree at most $n$ when $k\le n$ and is
well defined on the boundary too. Its constant member is one; the $k=0$
subfamily is the symmetric Jacobi family in $x$.

Reference: \cite[Section 2.3]{dunkl-xu}.
\end{theorem}
\begin{proof}
\leanok
\end{proof}

\begin{supporting}[Reference details: Measurability domain of the product measure and related results]
\begin{definition}[Measurability domain of the product measure]
\label{def:product-measure-input}
\lean{MeasureTheory.ProductMeasureInput,MeasureTheory.ProductMeasureInput.ofSFinite}\leanok
For an outer measure $\mu$, the admissible inner measures $\nu$ are those
for which the measure-valued map $x\mapsto(y\mapsto(x,y))_*\nu$ is
$\mu$-almost-everywhere measurable. This is the measurability condition used
to form the product measure $\mu\otimes\nu$.
\end{definition}
\begin{theorem}[Product measure on measurable sets]
\label{thm:product-measure-input}\uses{def:product-measure-input}
\lean{MeasureTheory.ProductMeasureInput.prod_apply}\leanok
For every measurable set $A$, the product measure satisfies
$(\mu\otimes\nu)(A)=\int \nu(\{y:(x,y)\in A\})\,d\mu(x)$,
where the integral is nonnegative and may be infinite.
\end{theorem}
\begin{proof}
\leanok
\end{proof}

\begin{definition}[Polynomial density measure]
\label{def:polynomial-density-measure}
\lean{Polynomial.polynomialDensityMeasure}\leanok
For a real polynomial $p$, write $D_p=\{x\in\mathbb R:p(x)\ge0\}$.
For a measure $\eta$ on $D_p$, its polynomial density measure is the real-coordinate
pushforward of $p\,d\eta$. The density has values in $[0,\infty)$ on its entire domain.
\end{definition}

\begin{theorem}[Polynomial density measure characterization]
\label{thm:polynomial-density-measure}\uses{def:polynomial-density-measure}
\lean{Polynomial.polynomialDensityMeasure_apply}\leanok
For every measurable $E\subseteq\mathbb R$, the polynomial density measure assigns
$\int_{\iota^{-1}(E)}p(x)\,d\eta(x)$, interpreted as a nonnegative integral, where
$\iota:D_p\hookrightarrow\mathbb R$ is the coordinate inclusion.
\end{theorem}
\begin{proof}
\leanok
\end{proof}

\begin{definition}[Degree-dependent disk outer measure]
\label{def:disk-outer-measure}\uses{def:polynomial-density-measure}
\lean{Polynomial.diskOuterMeasure}\leanok
For $k\in\mathbb N$, let $D_k=\{x\in\mathbb R:(1-x^2)^k\ge0\}$.
For a measure $\eta$ on $D_k$, the shifted disk outer measure is the real-coordinate
pushforward of $(1-x^2)^k\,d\eta(x)$. The domain is all of $\mathbb R$ for even
$k$ and $[-1,1]$ for odd $k$.
\end{definition}

\begin{theorem}[Disk outer measure characterization]
\label{thm:disk-outer-measure}\uses{def:disk-outer-measure,thm:polynomial-density-measure}
\lean{Polynomial.diskOuterMeasure_apply, Polynomial.diskOuterMeasure_zero}\leanok
For every measurable $E\subseteq\mathbb R$, the shifted measure assigns the
nonnegative integral $\int_{\iota^{-1}(E)}(1-x^2)^k\,d\eta(x)$.
At degree zero, a real measure transported to $D_0=\mathbb R$ is recovered unchanged.
\end{theorem}
\begin{proof}
\leanok
\end{proof}

\begin{theorem}[Jacobi realization of the disk outer shift]
\label{thm:jacobi-disk-outer-shift}\uses{def:disk-outer-measure,thm:jacobi-symmetric-measure-shift}
\lean{Polynomial.diskOuterMeasure_jacobi}\leanok
For every real $\alpha$ and $k\in\mathbb N$, transport the symmetric Jacobi
measure $\mu_{\alpha,\alpha}$ to $D_k=\{x:(1-x^2)^k\ge0\}$.
Its outer degree-$k$ shift equals $\mu_{\alpha+k,\alpha+k}$.
\end{theorem}
\begin{proof}
\leanok
\end{proof}

\begin{definition}[Polynomially modified disk outer measure]
\label{def:disk-outer-modification}\uses{def:polynomial-density-measure}
\lean{Polynomial.diskModificationWeight, Polynomial.diskOuterModification}\leanok
For $h\in\mathbb R[X]$ and $l\in\mathbb N$, put $w(x)=(1-x^2)^l h(x)$.
A measure $\eta$ on $\{x:w(x)\ge0\}$ has modified outer measure given by the
real-coordinate pushforward of $w\,d\eta$. All source points lie in the
nonnegative-density domain.
\end{definition}

\begin{theorem}[Modified disk outer measure characterization]
\label{thm:disk-outer-modification}\uses{def:disk-outer-modification,thm:polynomial-density-measure}
\lean{Polynomial.diskOuterModification_apply}\leanok
For every measurable $E\subseteq\mathbb R$, the modified measure assigns
$\int_{\iota^{-1}(E)}(1-x^2)^l h(x)\,d\eta(x)$ as a nonnegative integral.
\end{theorem}
\begin{proof}
\leanok
\end{proof}
\end{supporting}


\begin{definition}[Cartesian disk measure]
\label{def:disk-measure}\uses{def:product-measure-input,thm:jacobi-weight-basic}
\lean{MvPolynomial.diskCoordinates, MvPolynomial.cartesianDiskMeasure, MvPolynomial.dunklXuDiskMeasure}\leanok
The Cartesian chart $(x,z)\mapsto(x,\sqrt{1-x^2}z)$ has domain
$[-1,1]\times\R$. For a measure $\mu$ on $[-1,1]$ and a real inner measure $\nu$ with the product-kernel measurability condition of Definition~\ref{def:product-measure-input},
its pushforward of $\mu\otimes\nu$ is the Cartesian measure.
Real outer measures supported on $[-1,1]$ are transported to this domain.
The disk measure uses the Jacobi factors
$\mu_{\beta+1/2,\beta+1/2}$ and $\mu_{\beta,\beta}$ in this construction.
The disk orthogonality range is $\beta>-1$.
\end{definition}

\begin{supporting}[Reference details: Characterisation of the Cartesian chart and pushforward and related results]
\begin{theorem}[Characterisation of the Cartesian chart and pushforward]
\label{thm:disk-coordinate-characterisation}\uses{def:disk-measure,thm:integrable-input}
\lean{MvPolynomial.continuous_diskCoordinates,MvPolynomial.diskCoordinates_sq_nonneg,MvPolynomial.cartesianDiskMeasure_integral}\leanok
The Cartesian disk chart
$\Phi(x,z)=(x,\sqrt{1-x^2}\,z)$ on $[-1,1]\times\R$ is continuous,
including at $x=\pm1$, and $1-x^2\ge0$ on this domain.
If $f$ is integrable for the Cartesian measure
$\omega=\Phi_*(\mu\otimes\nu)$, then $f\circ\Phi$ is integrable for
$\mu\otimes\nu$, and
\[
 \int f\,d\omega=\int f\circ\Phi\,d(\mu\otimes\nu).
\]
\end{theorem}
\begin{proof}
\leanok
\end{proof}

\begin{theorem}[Disk integrability and moments]
\label{thm:disk-integrability}\uses{thm:jacobi-moments,thm:disk-coordinate-characterisation,def:multivariate-moments}
\lean{MvPolynomial.integrable_dunklXuDiskMeasure,MvPolynomial.hasFiniteMoments_dunklXuDiskMeasure,MvPolynomial.isFiniteMeasure_dunklXuDiskMeasure}\leanok
For $\beta>-1$, every continuous real-valued function on the plane is integrable for the
Cartesian disk measure. All polynomial moments exist, and the measure is finite.
\end{theorem}
\begin{proof}
\leanok
\end{proof}

\begin{theorem}[Jacobian identity for the disk weights]
\label{thm:disk-weight-jacobian}\uses{def:jacobi-measure}
\lean{MvPolynomial.diskCoordinates_weight}\leanok
For $\beta\in\R$ and $x,z\in(-1,1)$, put $r_x=\sqrt{1-x^2}$.
Writing $w_{a,b}$ for the Jacobi weight,
\[
 w_{\beta+1/2,\beta+1/2}(x)\,w_{\beta,\beta}(z)
 =r_x\bigl(1-x^2-(r_xz)^2\bigr)^\beta.
\]
\end{theorem}
\begin{proof}
\leanok
\end{proof}

\begin{definition}[Classical vertical disk-slice input]
\label{def:disk-slice-input}\uses{thm:jacobi-integrable}
\lean{MvPolynomial.diskSliceInput}\leanok
For $\beta>-1$, continuous $f:\R^2\to\R$, and $x\in(-1,1)$, put
$r_x=\sqrt{1-x^2}$. The function
\[
 y\longmapsto (1-x^2-y^2)^\beta f(x,y),\qquad y\in(-r_x,r_x),
\]
is denoted by $F_{\beta,f,x}$. Its domain is the open slice $(-r_x,r_x)$.
\end{definition}
\begin{theorem}[Classical slice domain and integrability]
\label{thm:disk-slice-input}\uses{def:disk-slice-input,thm:integrable-input}
\lean{MvPolynomial.diskSlice_radical_pos,MvPolynomial.diskSliceInput_integrable}\leanok
For $\beta>-1$, continuous $f:\R^2\to\R$, $x\in(-1,1)$ and
$y\in(-\sqrt{1-x^2},\sqrt{1-x^2})$, the base $1-x^2-y^2$ is positive.
The function $F_{\beta,f,x}(y)=(1-x^2-y^2)^\beta f(x,y)$ is Lebesgue-integrable
on that open slice, including when $-1<\beta<0$.
\end{theorem}
\begin{proof}
\leanok
\end{proof}

\begin{definition}[Classical outer disk-slice input]
\label{def:disk-slices-input}\uses{def:disk-slice-input,thm:disk-integrability,thm:disk-weight-jacobian}
\lean{MvPolynomial.diskSlicesInput}\leanok
For $\beta>-1$ and continuous $f:\R^2\to\R$, set $r_x=\sqrt{1-x^2}$
for $x\in(-1,1)$ and define
\[
 G_{\beta,f}(x)=\int_{y\in(-r_x,r_x)}(1-x^2-y^2)^\beta f(x,y)\,dy.
\]
The integrand is the slice function of Definition~\ref{def:disk-slice-input}.
\end{definition}
\begin{theorem}[Integrability of the outer disk-slice input]
\label{thm:disk-slices-input}\uses{def:disk-slices-input}
\lean{MvPolynomial.diskSlicesInput_integrable}\leanok
For $\beta>-1$ and continuous $f:\R^2\to\R$, the function
\[
 x\longmapsto G_{\beta,f}(x)
 =\int_{y\in(-\sqrt{1-x^2},\sqrt{1-x^2})}(1-x^2-y^2)^\beta f(x,y)\,dy
\]
is Lebesgue-integrable on $(-1,1)$.
\end{theorem}
\begin{proof}
\leanok
\end{proof}

\begin{theorem}[Disk density formula]
\label{thm:disk-density}\uses{def:disk-slices-input}
\lean{MvPolynomial.integral_dunklXuDiskMeasure_eq_slices}\leanok
For $\beta>-1$ and continuous $f$, integration against the disk measure is
\[
\int_{-1}^1\int_{-\sqrt{1-x^2}}^{\sqrt{1-x^2}}
 (1-x^2-y^2)^\beta f(x,y)\,dy\,dx.
\]
Both integrals use the open-slice inputs of Definitions~\ref{def:disk-slice-input} and~\ref{def:disk-slices-input}; all their integrands are classically integrable.
There is no additional normalization constant.
\end{theorem}
\begin{proof}
\leanok
\end{proof}

\rvnote{thm:disk-density}{
Integrability and finiteness of moments are given by Theorem~\ref{thm:disk-integrability}.
}
\end{supporting}


\begin{theorem}[Disk orthogonality and product norms]
\label{thm:disk-orthogonality}\uses{thm:jacobi-symmetric-measure-shift,thm:jacobi-positive,thm:cartesian-disk-bases,def:dunkl-xu-polynomials,thm:jacobi-pairwise}
\lean{MvPolynomial.integral_mul_dunklXuDisk_eq_zero, MvPolynomial.integral_sq_dunklXuDisk, MvPolynomial.integral_sq_dunklXuDisk_pos}\leanok
For $\beta>-1$, distinct valid pairs $(n,k)$ are orthogonal. Writing
$h_r^{\alpha,\alpha}$ for the usual Jacobi squared norm,
\[
 \|D_{n,k}^\beta\|^2
 =h_{n-k}^{\beta+k+1/2,\beta+k+1/2}h_k^{\beta,\beta}>0.
\]

Reference: \cite[Section 2.3]{dunkl-xu}.
\end{theorem}
\begin{proof}
\leanok
\end{proof}

\begin{theorem}[Disk polynomial bases and positivity]
\label{thm:disk-basis}\uses{thm:disk-orthogonality,thm:disk-integrability,thm:dunkl-xu-separated}
\lean{MvPolynomial.linearIndependent_dunklXuDisk, MvPolynomial.dunklXuDegreeBasis, MvPolynomial.dunklXuDegreeBasis_apply, MvPolynomial.span_dunklXuDisk, MvPolynomial.dunklXuBasis, MvPolynomial.dunklXuBasis_apply, MvPolynomial.hasPositiveDefiniteIntegral_dunklXuDiskMeasure}\leanok
The disk polynomials through degree $N$ form a basis of the full bivariate
polynomial space of degree at most $N$. The full family, indexed by $n\ge0$
and $0\le k\le n$, is a basis of the polynomial algebra. Every nonzero
polynomial has strictly positive squared integral for $\beta>-1$.
\end{theorem}
\begin{proof}
\leanok
\end{proof}

\begin{supporting}[Reference details: Classical Dunkl--Xu polynomial norm and related results]
\begin{definition}[Classical Dunkl--Xu polynomial norm]
\label{def:dunkl-xu-norm}\uses{def:disk-measure,def:dunkl-xu-polynomials}
\lean{MvPolynomial.dunklXuNorm}\leanok
For every parameter and index choice for which $(D_{n,k}^{\beta})^2$ is
integrable with respect to $\mu_\beta$, define
\[
 \|D_{n,k}^{\beta}\|_{\mu_\beta}=\sqrt{\int (D_{n,k}^{\beta})^2\,d\mu_\beta}.
\]
The domain consists of all such square-integrable polynomials; no additional
parameter restrictions are imposed.
\end{definition}
\begin{theorem}[Characterisation of the Dunkl--Xu norm]
\label{thm:dunkl-xu-norm}\uses{def:dunkl-xu-norm}
\lean{MvPolynomial.dunklXuNorm_nonneg,MvPolynomial.dunklXuNorm_sq}\leanok
Whenever $(D_{n,k}^{\beta})^2$ is $\mu_\beta$-integrable,
\[
 \|D_{n,k}^{\beta}\|_{\mu_\beta}\ge0,\qquad
 \|D_{n,k}^{\beta}\|_{\mu_\beta}^2=\int (D_{n,k}^{\beta})^2\,d\mu_\beta.
\]
\end{theorem}
\begin{proof}
\leanok
\end{proof}
\end{supporting}


\begin{theorem}[Normalized disk bases]
\label{thm:disk-normalization}\uses{thm:disk-basis,def:dunkl-xu-norm}
\lean{MvPolynomial.normalizedDunklXuBasis, MvPolynomial.integralInner_normalizedDunklXuBasis, MvPolynomial.normalizedDunklXuBasis_repr, MvPolynomial.normalizedDunklXuBasis_repr_eq_integralInner, MvPolynomial.normalizedDunklXuDegreeBasis}\leanok
Dividing by the positive square root of the product norm gives orthonormal
polynomial bases, both globally and in each finite degree space. Coordinates
are integral pairings; normalization multiplies raw coordinates by positive
norm factors.
\end{theorem}
\begin{proof}
\leanok
\end{proof}

The next results describe how raising the disk weight changes coordinates:
they give the support of the conversion matrix and the values and signs of
its coefficients.

\begin{supporting}[Reference details: Disk measure raising and related results]
\begin{theorem}[Disk measure raising]
\label{thm:disk-measure-raising}\uses{thm:jacobi-disk-outer-shift,thm:cartesian-separated-modification,thm:disk-integrability}
\lean{MvPolynomial.diskBoundary, MvPolynomial.integral_dunklXuDiskMeasure_add_one, MvPolynomial.integralInner_dunklXuDiskMeasure_add_one}\leanok
For $\beta>-1$, increasing $\beta$ by one multiplies the disk measure by
$1-x^2-y^2$, as an identity against every continuous real-valued test function on $\mathbb R^2$ and in particular for all
polynomial pairings.
\end{theorem}
\begin{proof}
\leanok
\end{proof}

\begin{definition}[Symmetric equal-degree raising coefficient]
\label{def:symmetric-domain-raise-a}
\lean{Polynomial.JacobiSymmetricRaiseAInput, Polynomial.jacobiSymmetricRaiseA}\leanok
For $\beta\in\mathbb R$ and natural degrees, define
\[
 A_0(\beta)=1,\quad A_1(\beta)=\frac{\beta+1}{\beta+2},\quad
 A_{m+2}(\beta)=\frac{(2m+2\beta+4)(m+2\beta+3)(m+2\beta+4)}
 {(2m+2\beta+5)(2m+2\beta+6)(2m+2\beta+4)}.
\]
The domain at each degree requires exactly the displayed denominator to be
nonzero; degree zero admits every real parameter.
\end{definition}

\begin{theorem}[Symmetric equal-degree raising coefficient characterization]
\label{thm:symmetric-domain-raise-a}\uses{def:symmetric-domain-raise-a}
\lean{Polynomial.JacobiSymmetricRaiseAInput.of_neg_one_lt, Polynomial.jacobiSymmetricRaiseA_apply}\leanok
The coefficient $A_n(\beta)$ equals the expression in Definition~\ref{def:symmetric-domain-raise-a} whenever its denominator conditions hold, including the separate low-degree cases specified there.
Every $\beta>-1$ lies in this domain for every natural index.
\end{theorem}
\begin{proof}
\leanok
\end{proof}

\begin{definition}[Symmetric lower-degree raising coefficient]
\label{def:symmetric-domain-raise-b}
\lean{Polynomial.JacobiSymmetricRaiseBInput, Polynomial.jacobiSymmetricRaiseB}\leanok
Define $B_0(\beta)=B_1(\beta)=0$ for every real $\beta$, and
\[
 B_{m+2}(\beta)=-\frac{(2m+2\beta+6)(m+\beta+2)^2}
 {(2m+2\beta+5)(2m+2\beta+6)(2m+2\beta+4)}
\]
when its displayed denominator is nonzero.
\end{definition}

\begin{theorem}[Symmetric lower-degree raising coefficient characterization]
\label{thm:symmetric-domain-raise-b}\uses{def:symmetric-domain-raise-b}
\lean{Polynomial.JacobiSymmetricRaiseBInput.of_neg_one_lt, Polynomial.jacobiSymmetricRaiseB_apply}\leanok
The coefficient $B_n(\beta)$ equals the expression in Definition~\ref{def:symmetric-domain-raise-b} whenever its denominator conditions hold, including the separate low-degree cases specified there.
Every $\beta>-1$ lies in this domain for every natural index.
\end{theorem}
\begin{proof}
\leanok
\end{proof}

\begin{definition}[Symmetric degree-plus-two lowering coefficient]
\label{def:symmetric-domain-lower-a}
\lean{Polynomial.JacobiSymmetricLowerInput, Polynomial.jacobiSymmetricLowerA}\leanok
Write $D_n(\beta)=(2n+2\beta+2)(n+2\beta+2)(2n+2\beta+3)(2n+2\beta+4)$. For $D_n(\beta)\ne0$, define
\[
 E_n(\beta)=-\frac{4(n+1)(n+2)(n+2\beta+2)(2n+2\beta+2)}{D_n(\beta)}.
\]
\end{definition}

\begin{theorem}[Symmetric degree-plus-two lowering coefficient characterization]
\label{thm:symmetric-domain-lower-a}\uses{def:symmetric-domain-lower-a}
\lean{Polynomial.JacobiSymmetricLowerInput.of_neg_one_lt, Polynomial.jacobiSymmetricLowerA_apply}\leanok
The coefficient $E_n(\beta)$ equals the expression in Definition~\ref{def:symmetric-domain-lower-a} whenever its denominator conditions hold, including the separate low-degree cases specified there.
Every $\beta>-1$ lies in this domain for every natural index.
\end{theorem}
\begin{proof}
\leanok
\end{proof}

\begin{definition}[Symmetric equal-degree lowering coefficient]
\label{def:symmetric-domain-lower-b}\uses{def:symmetric-domain-lower-a}
\lean{Polynomial.jacobiSymmetricLowerB}\leanok
With $D_n(\beta)=(2n+2\beta+2)(n+2\beta+2)(2n+2\beta+3)(2n+2\beta+4)$, define
\[
 F_n(\beta)=\frac{4(n+\beta+1)^2(2n+2\beta+4)(n+2\beta+2)}{D_n(\beta)}
\]
on the same domain $D_n(\beta)\ne0$.
\end{definition}

\begin{theorem}[Symmetric equal-degree lowering coefficient characterization]
\label{thm:symmetric-domain-lower-b}\uses{def:symmetric-domain-lower-b}
\lean{Polynomial.jacobiSymmetricLowerB_apply}\leanok
The coefficient $F_n(\beta)$ equals the expression in Definition~\ref{def:symmetric-domain-lower-b} whenever its denominator conditions hold, including the separate low-degree cases specified there.
The domain is the same as that of the degree-plus-two lowering coefficient.
\end{theorem}
\begin{proof}
\leanok
\end{proof}

\begin{definition}[Double symmetric raising coefficient A]
\label{def:symmetric-domain-double-a}\uses{def:symmetric-domain-raise-a}
\lean{Polynomial.JacobiSymmetricDoubleRaiseAInput, Polynomial.jacobiSymmetricDoubleRaiseA}\leanok
For $n\in\mathbb N$ and $\beta\in\mathbb R$, when every constituent coefficient below is defined, set $\widehat A_n(\beta)=A_n(\beta)A_n(\beta+1)$.
\end{definition}

\begin{theorem}[Double symmetric raising coefficient A characterization]
\label{thm:symmetric-domain-double-a}\uses{def:symmetric-domain-double-a,thm:symmetric-domain-raise-a}
\lean{Polynomial.JacobiSymmetricDoubleRaiseAInput.of_neg_one_lt, Polynomial.jacobiSymmetricDoubleRaiseA_apply}\leanok
The coefficient $\widehat A_n(\beta)$ equals the expression in Definition~\ref{def:symmetric-domain-double-a} whenever its denominator conditions hold, including the separate low-degree cases specified there.
Every $\beta>-1$ lies in this domain for every natural index.
\end{theorem}
\begin{proof}
\leanok
\end{proof}

\begin{definition}[Double symmetric raising coefficient B]
\label{def:symmetric-domain-double-b}\uses{def:symmetric-domain-raise-b,def:symmetric-domain-raise-a}
\lean{Polynomial.JacobiSymmetricDoubleRaiseBInput, Polynomial.jacobiSymmetricDoubleRaiseB}\leanok
For $n\in\mathbb N$ and $\beta\in\mathbb R$, when every constituent coefficient below is defined, set $\widehat B_n(\beta)=A_n(\beta)B_n(\beta+1)+B_n(\beta)A_{\max(n-2,0)}(\beta+1)$.
\end{definition}

\begin{theorem}[Double symmetric raising coefficient B characterization]
\label{thm:symmetric-domain-double-b}\uses{def:symmetric-domain-double-b,thm:symmetric-domain-raise-b,thm:symmetric-domain-raise-a}
\lean{Polynomial.JacobiSymmetricDoubleRaiseBInput.of_neg_one_lt, Polynomial.jacobiSymmetricDoubleRaiseB_apply}\leanok
The coefficient $\widehat B_n(\beta)$ equals the expression in Definition~\ref{def:symmetric-domain-double-b} whenever its denominator conditions hold, including the separate low-degree cases specified there.
Every $\beta>-1$ lies in this domain for every natural index.
\end{theorem}
\begin{proof}
\leanok
\end{proof}

\begin{definition}[Double symmetric raising coefficient C]
\label{def:symmetric-domain-double-c}\uses{def:symmetric-domain-raise-b}
\lean{Polynomial.JacobiSymmetricDoubleRaiseCInput, Polynomial.jacobiSymmetricDoubleRaiseC}\leanok
For $n\in\mathbb N$ and $\beta\in\mathbb R$, when every constituent coefficient below is defined, set $\widehat C_n(\beta)=B_n(\beta)B_{\max(n-2,0)}(\beta+1)$.
\end{definition}

\begin{theorem}[Double symmetric raising coefficient C characterization]
\label{thm:symmetric-domain-double-c}\uses{def:symmetric-domain-double-c,thm:symmetric-domain-raise-b}
\lean{Polynomial.JacobiSymmetricDoubleRaiseCInput.of_neg_one_lt, Polynomial.jacobiSymmetricDoubleRaiseC_apply}\leanok
The coefficient $\widehat C_n(\beta)$ equals the expression in Definition~\ref{def:symmetric-domain-double-c} whenever its denominator conditions hold, including the separate low-degree cases specified there.
Every $\beta>-1$ lies in this domain for every natural index.
\end{theorem}
\begin{proof}
\leanok
\end{proof}

\begin{definition}[Cartesian disk raising coefficient A]
\label{def:symmetric-domain-disk-a}\uses{def:symmetric-domain-raise-a}
\lean{MvPolynomial.DunklXuDiskRaiseAInput, MvPolynomial.dunklXuDiskRaiseA}\leanok
For $n,k\in\mathbb N$ and $\beta\in\mathbb R$, when every constituent coefficient is defined, set $\mathcal A_{n,k}(\beta)=A_k(\beta)A_{\max(n-k,0)}(\beta+k+1/2)$.
\end{definition}

\begin{theorem}[Cartesian disk raising coefficient A characterization]
\label{thm:symmetric-domain-disk-a}\uses{def:symmetric-domain-disk-a,thm:symmetric-domain-raise-a}
\lean{MvPolynomial.DunklXuDiskRaiseAInput.of_neg_one_lt, MvPolynomial.dunklXuDiskRaiseA_apply}\leanok
The coefficient $\mathcal A_{n,k}(\beta)$ equals the expression in Definition~\ref{def:symmetric-domain-disk-a} whenever its denominator conditions hold, including the separate low-degree cases specified there.
Every $\beta>-1$ lies in this domain for every natural index.
\end{theorem}
\begin{proof}
\leanok
\end{proof}

\begin{definition}[Cartesian disk raising coefficient B]
\label{def:symmetric-domain-disk-b}\uses{def:symmetric-domain-raise-b,def:symmetric-domain-raise-a}
\lean{MvPolynomial.DunklXuDiskRaiseBInput, MvPolynomial.dunklXuDiskRaiseB}\leanok
For $n,k\in\mathbb N$ and $\beta\in\mathbb R$, when every constituent coefficient is defined, set $\mathcal B_{n,k}(\beta)=A_k(\beta)B_{\max(n-k,0)}(\beta+k+1/2)$.
\end{definition}

\begin{theorem}[Cartesian disk raising coefficient B characterization]
\label{thm:symmetric-domain-disk-b}\uses{def:symmetric-domain-disk-b,thm:symmetric-domain-raise-b,thm:symmetric-domain-raise-a}
\lean{MvPolynomial.DunklXuDiskRaiseBInput.of_neg_one_lt, MvPolynomial.dunklXuDiskRaiseB_apply}\leanok
The coefficient $\mathcal B_{n,k}(\beta)$ equals the expression in Definition~\ref{def:symmetric-domain-disk-b} whenever its denominator conditions hold, including the separate low-degree cases specified there.
Every $\beta>-1$ lies in this domain for every natural index.
\end{theorem}
\begin{proof}
\leanok
\end{proof}

\begin{definition}[Cartesian disk raising coefficient C]
\label{def:symmetric-domain-disk-c}\uses{def:symmetric-domain-raise-b,def:symmetric-domain-lower-a}
\lean{MvPolynomial.DunklXuDiskRaiseCInput, MvPolynomial.dunklXuDiskRaiseC}\leanok
Let $n,k\in\mathbb N$. For $k<2$ the coefficient is zero for every real $\beta$. For $k\ge2$, when every constituent coefficient is defined, set $\mathcal C_{n,k}(\beta)=B_k(\beta)E_{\max(n-k,0)}(\beta+1+(k-2)+1/2)$.
\end{definition}

\begin{theorem}[Cartesian disk raising coefficient C characterization]
\label{thm:symmetric-domain-disk-c}\uses{def:symmetric-domain-disk-c,thm:symmetric-domain-raise-b,thm:symmetric-domain-lower-a}
\lean{MvPolynomial.DunklXuDiskRaiseCInput.of_neg_one_lt, MvPolynomial.dunklXuDiskRaiseC_apply}\leanok
The coefficient $\mathcal C_{n,k}(\beta)$ equals the expression in Definition~\ref{def:symmetric-domain-disk-c} whenever its denominator conditions hold, including the separate low-degree cases specified there.
Every $\beta>-1$ lies in this domain for every natural index.
\end{theorem}
\begin{proof}
\leanok
\end{proof}

\begin{definition}[Cartesian disk raising coefficient D]
\label{def:symmetric-domain-disk-d}\uses{def:symmetric-domain-raise-b,def:symmetric-domain-lower-b}
\lean{MvPolynomial.DunklXuDiskRaiseDInput, MvPolynomial.dunklXuDiskRaiseD}\leanok
Let $n,k\in\mathbb N$. For $k<2$ the coefficient is zero for every real $\beta$. For $k\ge2$, when every constituent coefficient is defined, set $\mathcal D_{n,k}(\beta)=B_k(\beta)F_{\max(n-k,0)}(\beta+1+(k-2)+1/2)$.
\end{definition}

\begin{theorem}[Cartesian disk raising coefficient D characterization]
\label{thm:symmetric-domain-disk-d}\uses{def:symmetric-domain-disk-d,thm:symmetric-domain-raise-b,thm:symmetric-domain-lower-a}
\lean{MvPolynomial.DunklXuDiskRaiseDInput.of_neg_one_lt, MvPolynomial.dunklXuDiskRaiseD_apply}\leanok
The coefficient $\mathcal D_{n,k}(\beta)$ equals the expression in Definition~\ref{def:symmetric-domain-disk-d} whenever its denominator conditions hold, including the separate low-degree cases specified there.
Every $\beta>-1$ lies in this domain for every natural index.
\end{theorem}
\begin{proof}
\leanok
\end{proof}
\end{supporting}


\begin{theorem}[Four-term Dunkl--Xu parameter raising]
\label{thm:dunkl-xu-raising}\uses{thm:jacobi-symmetric-coefficients,thm:disk-basis,thm:symmetric-domain-disk-a,thm:symmetric-domain-disk-b,thm:symmetric-domain-disk-c,thm:symmetric-domain-disk-d,thm:jacobi-symmetric-conversion}
\lean{MvPolynomial.dunklXuDisk_raise, MvPolynomial.exists_dunklXuDisk_raise, MvPolynomial.dunklXuBasis_raise_repr_eq_zero, MvPolynomial.dunklXuBasis_raise_repr_self_pos}\leanok
For $\beta>-1$ and $0\le k\le n$, the polynomial $D_{n,k}^{\beta}$ belongs to
the real span of $D_{m,j}^{\beta+1}$ with
\[
 (m,j)\in\{(n,k),(n-2,k),(n,k-2),(n-2,k-2)\},\qquad 0\le j\le m.
\]
Omit indices that become negative; subtraction here is integer subtraction.
For $\beta>-1$ and natural indices $k\le n$, put $q=n-k$ and
$\gamma=\beta+k+1/2$. Define the four coefficients, in the displayed index order, by
\[
 a=A_k(\beta)A_q(\gamma),\qquad b=A_k(\beta)B_q(\gamma),\qquad
 c=\begin{cases}0&k<2,\\B_k(\beta)E_q(\gamma-1)&k\ge2,\end{cases}\qquad
 d=\begin{cases}0&k<2,\\B_k(\beta)F_q(\gamma-1)&k\ge2.\end{cases}
\]
Here $A,B,E,F$ are the symmetric Jacobi coefficients on the domains specified in Theorem~\ref{thm:jacobi-symmetric-conversion}.
When $\beta>-1$, the unchanged-index coefficient is strictly positive.
This is the ordinary unit parameter shift of the Cartesian disk family.

References: \cite[Section 2.3]{dunkl-xu}; the Cartesian disk raising operator in \texttt{MultivariateOrthogonalPolynomials.jl}~\cite{mop-jl}.
\end{theorem}
\begin{proof}
\leanok
\end{proof}

\begin{supporting}[Reference details: Explicit disk raising coefficients and signs and related results]
\begin{theorem}[Explicit disk raising coefficients and signs]
\label{thm:dunkl-xu-raising-coefficients}\uses{thm:symmetric-domain-disk-a,thm:jacobi-symmetric-coefficients,thm:symmetric-domain-disk-b,thm:symmetric-domain-disk-c,thm:symmetric-domain-disk-d}
\lean{MvPolynomial.dunklXuDiskRaiseA_pos, MvPolynomial.dunklXuDiskRaiseB_nonpos, MvPolynomial.dunklXuDiskRaiseB_neg, MvPolynomial.dunklXuDiskRaiseC_pos, MvPolynomial.dunklXuDiskRaiseD_neg}\leanok
For $\beta>-1$ and $0\le k\le n$, let $q=n-k$ and $\gamma=\beta+k+1/2$.
In the four-term connection, the coefficients at $(n,k),(n-2,k),(n,k-2),(n-2,k-2)$ are
\[
 A_k(\beta)A_q(\gamma),\quad A_k(\beta)B_q(\gamma),\quad
 B_k(\beta)E_q(\gamma-1),\quad B_k(\beta)F_q(\gamma-1).
\]
The last two are defined to be zero when $k<2$. The first is strictly positive;
the second is zero for $q<2$ and strictly negative for $q\ge2$.
When $k\ge2$, the third is strictly positive and the fourth strictly negative.

Reference: \cite[Section 2.3]{dunkl-xu}.
\end{theorem}
\begin{proof}
\leanok
\end{proof}

Zernike polynomials separate radial factors from sine and cosine angular
modes. The Laplacian preserves these angular sectors, which we use in the
annular construction of the next chapter.

\begin{definition}[Polar coordinates on nonnegative squared radii]
\label{def:zernike-coordinates}
\lean{MvPolynomial.zernikeCoordinates}\leanok
For $s\in[0,\infty)$ and $\theta\in\mathbb R$, set $\Phi(s,\theta)=\sqrt{s}(\cos\theta,\sin\theta)$.
\end{definition}

\begin{theorem}[Continuity of polar coordinates]
\label{thm:zernike-coordinates}\uses{def:zernike-coordinates}
\lean{MvPolynomial.continuous_zernikeCoordinates}\leanok
The map $\Phi:[0,\infty)\times\mathbb R\to\mathbb R^2$ is continuous, including at $s=0$.
\end{theorem}
\begin{proof}
\leanok
\end{proof}

\begin{definition}[Auxiliary polar product measure]
\label{def:zernike-product-measure}\uses{def:unit-jacobi-measure,def:product-measure-input,def:zernike-coordinates}
\lean{MvPolynomial.zernikeProductMeasure}\leanok
For real $a,b$, let $\eta_{b,a}$ be the affine Jacobi measure on $[0,1]$, regarded as a measure on $[0,\infty)$. Define $\omega^{\rm aux}_{a,b}=\Phi_*(\eta_{b,a}\otimes d\theta|_{(-\pi,\pi)})$.
\end{definition}

\begin{theorem}[Integrability for the polar product measure]
\label{thm:zernike-product-measure}\uses{def:zernike-product-measure,thm:jacobi-moments,def:multivariate-moments,thm:jacobi-weight-basic,thm:zernike-coordinates}
\lean{MvPolynomial.integrable_zernikeProductMeasure, MvPolynomial.hasFiniteMoments_zernikeProductMeasure}\leanok
If $a,b>-1$, every continuous real-valued function on $\mathbb R^2$ is integrable for $\omega^{\rm aux}_{a,b}$. In particular, all polynomial moments are finite.
\end{theorem}
\begin{proof}
\leanok
\end{proof}
\end{supporting}


\begin{definition}[Standard weighted disk measure]
\label{def:zernike-measure}\uses{def:zernike-product-measure}
\lean{MvPolynomial.zernikeMeasure}\leanok
For real $a,b$, define $\omega_{a,b}=2^{-(a+b+2)}\omega^{\rm aux}_{a,b}$.
\end{definition}

\begin{supporting}[Formal details: Finite moments of the weighted disk measure and related results]
\begin{theorem}[Finite moments of the weighted disk measure]
\label{thm:zernike-measure-characterisation}\uses{def:zernike-measure,thm:zernike-product-measure}
\lean{MvPolynomial.hasFiniteMoments_zernikeMeasure}\leanok
For $a,b>-1$, the measure $\omega_{a,b}$ has all polynomial moments finite.
\end{theorem}
\begin{proof}
\leanok
\end{proof}

\begin{definition}[Integrable angular polynomial input]
\label{def:zernike-angular-input}\uses{def:integrable-input,thm:zernike-coordinates}
\lean{MvPolynomial.zernikeAngularInput}\leanok
For a Cartesian polynomial $p$ and $s\ge0$, let $A_p(s)$ be the function $\theta\mapsto p(\Phi(s,\theta))$ on $(-\pi,\pi)$. Write $J_p(s)=\int_{-\pi}^{\pi}A_p(s)(\theta)\,d\theta$.
\end{definition}

\begin{theorem}[Angular integrability]
\label{thm:zernike-angular-input}\uses{def:zernike-angular-input,thm:integrable-input}
\lean{MvPolynomial.zernikeAngularInput_integrable}\leanok
For every Cartesian polynomial $p$ and $s\ge0$, the function
$\theta\mapsto A_p(s)(\theta)=p(\Phi(s,\theta))$ is Lebesgue-integrable
on $(-\pi,\pi)$.
\end{theorem}
\begin{proof}
\leanok
\end{proof}

\begin{definition}[Integrable squared-radius polynomial input]
\label{def:zernike-squared-radius-input}\uses{thm:zernike-angular-input,thm:zernike-product-measure,thm:unit-jacobi-measure-characterisation,thm:unit-jacobi-weight}
\lean{MvPolynomial.zernikeSquaredRadiusInput}\leanok
For $a,b>-1$ and a Cartesian polynomial $p$, define the function $S_{a,b,p}(s)=(1-s)^bs^aJ_p(s)$ on $s\in(0,1)$.
\end{definition}

\begin{theorem}[Squared-radius integrability]
\label{thm:zernike-squared-radius-input}\uses{def:zernike-squared-radius-input}
\lean{MvPolynomial.zernikeSquaredRadiusInput_integrable}\leanok
The function $S_{a,b,p}$ is integrable on $(0,1)$, including when its weight has endpoint singularities.
\end{theorem}
\begin{proof}
\leanok
\end{proof}

\begin{definition}[Integrable polar polynomial input]
\label{def:zernike-polar-input}\uses{thm:zernike-squared-radius-input}
\lean{MvPolynomial.zernikePolarInput}\leanok
For $a,b>-1$ and a Cartesian polynomial $p$, define the function $R_{a,b,p}(r)=r(r^2)^a(1-r^2)^bJ_p(r^2)$ on $r\in(0,1)$.
\end{definition}

\begin{theorem}[Polar integrability]
\label{thm:zernike-polar-input}\uses{def:zernike-polar-input}
\lean{MvPolynomial.zernikePolarInput_integrable}\leanok
The function $R_{a,b,p}$ is integrable on $(0,1)$, including admissible endpoint singularities.
\end{theorem}
\begin{proof}
\leanok
\end{proof}
\end{supporting}


\begin{definition}[Real Zernike polynomials]
\label{def:zernike-real}\uses{def:unit-jacobi-measure}
\lean{MvPolynomial.diskHarmonics, MvPolynomial.zernikeCos, MvPolynomial.zernikeSin, MvPolynomial.zernike}\leanok
Let $H^c_m$ and $H^s_m$ be the real polynomials with $H^c_m+iH^s_m=(x+iy)^m$.
For real $a,b$ and the shifted Jacobi polynomials
$\widetilde P_k^{(a,b)}(s)=P_k^{(a,b)}(2s-1)$, set
\[
 Z^c_{k,m}=\widetilde P_k^{(b,a+m)}(x^2+y^2)\,H^c_m,\qquad
 Z^s_{k,m}=\widetilde P_k^{(b,a+m)}(x^2+y^2)\,H^s_m.
\]
In total degree $n$, the $j$th polynomial, $0\le j\le n$, is the sine mode with
frequency $m=j+1$ when $n+j$ is odd and the cosine mode with frequency $m=j$
otherwise, with radial degree $k=(n-m)/2$.

Reference: \cite[18.37(i)]{dlmf}.
\end{definition}

\begin{theorem}[Polar form, degrees and ordering]
\label{thm:zernike-real-basic}\uses{def:zernike-real,thm:jacobi-degree}
\lean{MvPolynomial.eval_zernikeCos_polar, MvPolynomial.eval_zernikeSin_polar, MvPolynomial.totalDegree_zernike_le, MvPolynomial.zernike_index_spec}\leanok
In polar coordinates,
\[
 Z^c_{k,m}=P_k^{(b,a+m)}(2r^2-1)\,r^m\cos(m\theta),\qquad
 Z^s_{k,m}=P_k^{(b,a+m)}(2r^2-1)\,r^m\sin(m\theta),
\]
and both have total degree at most $2k+m$. For $0\le j\le n$, the indices of the
$j$th polynomial in degree $n$ satisfy $2k+m=n$, and a sine mode has $m\ne0$.

Reference: \cite[18.37(i)]{dlmf}.
\end{theorem}
\begin{proof}
\leanok
\end{proof}

Within each degree the modes are thus ordered $1,\sin(2\theta),\cos(2\theta),\ldots$
for even degree and $\sin\theta,\cos\theta,\sin(3\theta),\cos(3\theta),\ldots$ for odd
degree, the real ordering in \texttt{MultivariateOrthogonalPolynomials.jl}~\cite{mop-jl}.
For $a=0$, \cite[18.37.1]{dlmf} combines these real modes into a complex mode and
divides by $P_k^{(b,m)}(1)$; these raw real polynomials do not include that
endpoint-normalizing factor. Positive $L^2$ normalization is a separate operation.

\begin{supporting}[Reference details: Classical Zernike polynomial norm and related results]
\begin{definition}[Classical Zernike polynomial norm]
\label{def:zernike-norm}\uses{def:zernike-measure,def:zernike-real}
\lean{MvPolynomial.zernikeNorm}\leanok
For every parameter and index choice for which $(Z_{n,j}^{a,b})^2$ is
integrable with respect to $\omega_{a,b}$, define
\[
 \|Z_{n,j}^{a,b}\|_{\omega_{a,b}}=\sqrt{\int (Z_{n,j}^{a,b})^2\,d\omega_{a,b}}.
\]
The domain consists of all such square-integrable polynomials; no additional
parameter restrictions are imposed.
\end{definition}
\begin{theorem}[Characterisation of the Zernike norm]
\label{thm:zernike-norm}\uses{def:zernike-norm}
\lean{MvPolynomial.zernikeNorm_nonneg,MvPolynomial.zernikeNorm_sq}\leanok
Whenever $(Z_{n,j}^{a,b})^2$ is $\omega_{a,b}$-integrable,
\[
 \|Z_{n,j}^{a,b}\|_{\omega_{a,b}}\ge0,\qquad
 \|Z_{n,j}^{a,b}\|_{\omega_{a,b}}^2=\int (Z_{n,j}^{a,b})^2\,d\omega_{a,b}.
\]
\end{theorem}
\begin{proof}
\leanok
\end{proof}
\end{supporting}


\begin{theorem}[Zernike families on disks]
\label{thm:zernike-polynomials}\uses{thm:measure-conversion-cholesky,thm:zernike-real-basic,thm:mathlib-total-degree-finite,def:zernike-norm,thm:zernike-measure-characterisation,thm:unit-jacobi-measure,thm:jacobi-growth-constants}
\lean{MvPolynomial.integralInner_zernikeMode_eq_zero, MvPolynomial.integral_sq_zernikeMode, MvPolynomial.integral_sq_zernikeMode_pos, MvPolynomial.zernikeDegreeBasis, MvPolynomial.normalizedZernikeDegreeBasis, MvPolynomial.integralGram_normalizedZernikeDegreeBasis, MvPolynomial.diskLaplacian_radial_mul, MvPolynomial.diskLaplacian_zernikeCos, MvPolynomial.diskLaplacian_zernikeSin}\leanok
For $a,b>-1$, the retained real Zernike modes are mutually orthogonal for the
standard radial disk measure and form complete bases of every finite total-degree
polynomial space. Writing $h_k^{(b,a+m)}$ for the standard Jacobi squared norm,
their squared norms are
\[
 2^{-(a+b+2)}2^{-m}h_k^{(b,a+m)}\,c_m,
 \qquad c_0=2\pi,\quad c_m=\pi\ (m>0).
\]
Dividing by the positive square roots gives identity Gram matrices. If $H_m$
is either solid angular harmonic, the Cartesian polynomial Laplacian satisfies
\[
 \Delta\bigl(p(r^2)H_m\bigr)
 =4\bigl(r^2p''(r^2)+(m+1)p'(r^2)\bigr)H_m.
\]
In particular the Laplacian preserves each angular sector, including at the origin.

Reference: \cite[18.37(i)]{dlmf}.
\end{theorem}
\begin{proof}
\leanok
\end{proof}

\begin{supporting}[Reference details: Classical real-mode Zernike normalization and related results]
\begin{definition}[Classical real-mode Zernike normalization]
\label{def:zernike-mode-normalization}\uses{def:zernike-measure,def:zernike-real}
\lean{MvPolynomial.ZernikeModeNormInput, MvPolynomial.ZernikeModeNormalizationInput, MvPolynomial.zernikeModeNorm, MvPolynomial.normalizedZernikeMode}\leanok
For an individual real Zernike mode $Z$, with $Z^2$ integrable against $\omega_{a,b}$, define its norm by $\sqrt{\int Z^2\,d\omega_{a,b}}$.
The normalized polynomial $Z/\|Z\|$ additionally requires a strictly positive
squared integral.
\end{definition}

\begin{theorem}[Real-mode Zernike norm characterization]
\label{thm:zernike-mode-normalization}\uses{def:zernike-mode-normalization,thm:zernike-polynomials}
\lean{MvPolynomial.ZernikeModeNormInput.of_parameters, MvPolynomial.ZernikeModeNormalizationInput.of_parameters, MvPolynomial.zernikeModeNorm_sq, MvPolynomial.integral_sq_normalizedZernikeMode}\leanok
The square of the norm equals the convergent squared integral.
For $a,b>-1$, every mode has a finite norm and every retained mode has positive
norm; its normalized polynomial has squared integral one.
\end{theorem}
\begin{proof}
\leanok
\end{proof}
\end{supporting}




\begin{supporting}[Supporting results: parity lifts, symmetric Jacobi conversions, radial measure]

\begin{theorem}[Disk lift evaluation and degree]
\label{thm:disk-polynomial-lift}\uses{def:disk-polynomial-lift}
\lean{MvPolynomial.eval_diskLift_scaled, MvPolynomial.totalDegree_diskLift_le, MvPolynomial.totalDegree_eval₂_C_X_le, MvPolynomial.diskLift_add, MvPolynomial.diskLift_smul, MvPolynomial.diskLift_add_two}\leanok
The lift is linear in $p$ and has total degree at most $n$. If $\deg p\le n$
and $p$ has parity $n$, then $L_np(x,sz)=s^np(z)$ whenever $s^2=1-x^2$.
This is a polynomial identity and does not require $s\ne0$.
Also $L_{n+2}p=(1-x^2)L_np$ when $\deg p\le n$.
\end{theorem}
\begin{proof}
\leanok
\end{proof}

\begin{theorem}[Two-term symmetric Jacobi conversions]
\label{thm:jacobi-symmetric-conversion}\uses{thm:jacobi-parameter-raising,thm:jacobi-derivative-lowering,thm:jacobi-recurrence,thm:symmetric-domain-raise-b,thm:symmetric-domain-raise-a,def:symmetric-domain-lower-b,thm:symmetric-domain-lower-a}
\lean{Polynomial.exists_jacobi_symmetric_raise, Polynomial.exists_jacobi_symmetric_weight_lower, Polynomial.jacobi_symmetric_raise, Polynomial.jacobi_symmetric_weight_lower, Polynomial.jacobi_symmetric_raise_all}\leanok
Let $A_j,B_j,E_j,F_j$ be the coefficients on the degree-dependent domains of
Definitions~\ref{def:symmetric-domain-raise-a},
\ref{def:symmetric-domain-raise-b}, \ref{def:symmetric-domain-lower-a}, and
\ref{def:symmetric-domain-lower-b}. Every $\beta>-1$ belongs to these domains.
For real $\beta>-1$ and $n\ge0$, there exist real coefficients giving
\[
P_{n+2}^{(\beta,\beta)}\in
\operatorname{span}\{P_{n+2}^{(\beta+1,\beta+1)},P_n^{(\beta+1,\beta+1)}\},
\qquad
(1-X^2)P_n^{(\beta+1,\beta+1)}\in
\operatorname{span}\{P_{n+2}^{(\beta,\beta)},P_n^{(\beta,\beta)}\}.
\]
The range includes $\beta=-1/2$. The raising coefficients are $A_{n+2},B_{n+2}$ and the weighted-lowering coefficients are $E_n,F_n$, respectively.

Reference: \cite[18.9(ii)]{dlmf}.
\end{theorem}
\begin{proof}
\leanok
\end{proof}

\rvnote{thm:jacobi-symmetric-conversion}{
The signs of these coefficients are stated in Theorem~\ref{thm:jacobi-symmetric-coefficients}.
}

\begin{theorem}[Explicit symmetric connection coefficients and signs]
\label{thm:jacobi-symmetric-coefficients}\uses{thm:symmetric-domain-raise-a,thm:symmetric-domain-raise-b,thm:symmetric-domain-lower-a,def:symmetric-domain-lower-b}
\lean{Polynomial.jacobiSymmetricRaiseA_pos, Polynomial.jacobiSymmetricRaiseB_neg, Polynomial.jacobiSymmetricRaiseB_nonpos, Polynomial.jacobiSymmetricLowerA_neg, Polynomial.jacobiSymmetricLowerB_pos}\leanok
Write $A_m,B_m$ for the coefficients at degrees $m,m-2$ in symmetric raising,
and $E_n,F_n$ for the coefficients at degrees $n+2,n$ in weighted lowering.
For $\beta>-1$, put $d_j=2n+2\beta+j$. The explicit formulas are
\begin{align*}
 A_{n+2}&=\frac{d_4(n+2\beta+3)(n+2\beta+4)}{d_5d_6d_4}>0,&
 B_{n+2}&=-\frac{d_6(n+\beta+2)^2}{d_5d_6d_4}<0,\\
 E_n&=-\frac{4(n+1)(n+2)(n+2\beta+2)d_2}{d_2(n+2\beta+2)d_3d_4}<0,&
 F_n&=\frac{4(n+\beta+1)^2d_4(n+2\beta+2)}{d_2(n+2\beta+2)d_3d_4}>0.
\end{align*}
The initial values are $A_0=1$, $A_1=(\beta+1)/(\beta+2)$ and $B_0=B_1=0$.
In particular, no removable singularity is introduced at $\beta=-1/2$.

Reference: \cite[18.9(ii)]{dlmf}.
\end{theorem}
\begin{proof}
\leanok
\end{proof}

\begin{theorem}[Integer Jacobi weight shifts]
\label{thm:jacobi-symmetric-measure-shift}\uses{thm:jacobi-measure-characterisation,thm:jacobi-weight-calculus}

\lean{Polynomial.jacobiWeight_add_nat, Polynomial.integral_jacobiMeasure_add_nat_integrable}\leanok
For $\alpha,\beta\in\R$, $k\in\N$ and $x\in(-1,1)$,
\[w_{\alpha+k,\beta+k}(x)=(1-x^2)^k w_{\alpha,\beta}(x).\]
For $f\in\mathcal I(\mu_{\alpha+k,\beta+k})$, the function $(1-x^2)^kf(x)$ is $\mu_{\alpha,\beta}$-integrable and
\[\int f\,d\mu_{\alpha+k,\beta+k}=\int (1-x^2)^kf(x)\,d\mu_{\alpha,\beta}.\]
\end{theorem}
\begin{proof}
\leanok
\end{proof}

\begin{theorem}[Standard radial disk measure]
\label{thm:zernike-disk-measure}\uses{def:zernike-measure,def:zernike-polar-input}
\lean{MvPolynomial.integral_zernikeMeasure_squaredRadius, MvPolynomial.integral_zernikeMeasure_polar}\leanok
For $a,b>-1$, polynomial integration against the constructed disk measure is
\[
 \int_0^1 r(r^2)^a(1-r^2)^b
   \int_{-\pi}^{\pi}p(r\cos\theta,r\sin\theta)\,d\theta\,dr.
\]
All polynomial moments are finite. In squared-radius coordinates, the area
factor is $1/2$; relative to the affine Jacobi product measure, the mass factor
is $2^{-(a+b+2)}$.
\end{theorem}
\begin{proof}
\leanok
\end{proof}

\end{supporting}

The disk provides the angular description. Removing an inner disk changes
the radial measure and brings in the semiclassical Jacobi family of
Chapter~\ref{chap:semiclassical}.

\chapter{The annulus: Zernike-type bases and a convergent solver}\label{chap:annulus}
\rvsummary{Zernike-type bases on the annulus, a sparse Laplacian, and a convergent Galerkin solver for the Dirichlet problem.}

On the annulus $q<x^2+y^2<1$ the radial factor of a Zernike-type basis is a
semiclassical Jacobi polynomial in $t(1-r^2)$, with $t=(1-q)^{-1}$. The radial family is the one from
Chapter~\ref{chap:semiclassical}. Conversion preserves each angular
mode. The Laplacian becomes banded in the radial index within each mode. A
boundary-vanishing basis turns Dirichlet problems into symmetric tridiagonal
systems.

We first construct the radial basis and sparse operators, then the finite
Galerkin systems and their energy spaces. We prove convergence as both angular
and radial resolution increase, and identify the completed space with physical
functions and Cartesian gradient classes on the annulus.

\section{Geometry, measure and bases}\label{sec:annular-geometry}

\begin{supporting}[Reference details: Nondegenerate annular parameter and related results]
\begin{definition}[Nondegenerate annular parameter]
\label{def:annulus-parameter}
\lean{Polynomial.annulusParameter}\leanok
For $q\in\mathbb R\setminus\{1\}$, define $t(q)=(1-q)^{-1}$.
\end{definition}

\begin{theorem}[Annular parameter identities]
\label{thm:annulus-parameter}\uses{def:annulus-parameter}
\lean{Polynomial.annulusParameter_ne_zero, Polynomial.annulusParameter_mul, Polynomial.one_lt_annulusParameter}\leanok
For $q\ne1$, $t(q)\ne0$ and $t(q)(1-q)=1$. If $0<q<1$, then $t(q)>1$.
\end{theorem}
\begin{proof}
\leanok
\end{proof}

\begin{definition}[Classical annular coordinate map]
\label{def:zernike-annulus-coordinates}\uses{def:zernike-coordinates,def:annulus-parameter}
\lean{MvPolynomial.zernikeAnnulusCoordinates}\leanok
For $q\ne1$, define $\Psi_q(x,\theta)=\Phi(1-x/t(q),\theta)$ on the domain $1-x/t(q)\ge0$.
\end{definition}

\begin{theorem}[Annular chart continuity and admissible parameters]
\label{thm:zernike-annulus-coordinates}\uses{def:zernike-annulus-coordinates,thm:zernike-coordinates}
\lean{MvPolynomial.continuous_zernikeAnnulusCoordinates, MvPolynomial.zernikeAnnulusCoordinates_radical_nonneg}\leanok
The map $\Psi_q$ is continuous on its domain. If $q\ge0$, $q\ne1$, and $0\le x\le1$, then $1-x/t(q)\ge0$.
\end{theorem}
\begin{proof}
\leanok
\end{proof}

\begin{definition}[Auxiliary annular product measure]
\label{def:zernike-annulus-product-measure}\uses{thm:zernike-annulus-coordinates,def:semiclassical-jacobi-measure,def:product-measure-input}
\lean{MvPolynomial.zernikeAnnulusProductMeasure}\leanok
For $q\ge0$, $q\ne1$, and real $a,b$, define $\omega^{\rm aux}_{q,a,b}$ as the pushforward by $\Psi_q$ of the semiclassical measure with parameters $(2,a,b,0)$, regarded as a measure on $[0,1]$, times angular Lebesgue measure on $(-\pi,\pi)$.
\end{definition}

\begin{theorem}[Auxiliary annular integrability]
\label{thm:zernike-annulus-product-measure}\uses{def:zernike-annulus-product-measure,thm:semiclassical-jacobi-measure,def:multivariate-moments}
\lean{MvPolynomial.integrable_zernikeAnnulusProductMeasure, MvPolynomial.hasFiniteMoments_zernikeAnnulusProductMeasure}\leanok
For $0<q<1$ and $a,b>-1$, every continuous real-valued function on $\mathbb R^2$ is integrable for $\omega^{\rm aux}_{q,a,b}$, and all polynomial moments are finite.
\end{theorem}
\begin{proof}
\leanok
\end{proof}
\end{supporting}


\begin{definition}[Standard weighted annular measure]
\label{def:zernike-annulus-measure}\uses{def:zernike-annulus-product-measure}
\lean{MvPolynomial.zernikeAnnulusMeasure}\leanok
For $0\le q<1$ and real $a,b$, define $\omega_{q,a,b}=\frac12t(q)^{-(a+b+1)}\omega^{\rm aux}_{q,a,b}$. The value $q=0$ includes the disk endpoint.
\end{definition}

\begin{supporting}[Formal details: Finite annular polynomial moments]
\begin{theorem}[Finite annular polynomial moments]
\label{thm:zernike-annulus-measure-characterisation}\uses{def:zernike-annulus-measure,thm:zernike-annulus-product-measure}
\lean{MvPolynomial.hasFiniteMoments_zernikeAnnulusMeasure}\leanok
For $0<q<1$ and $a,b>-1$, all polynomial moments of $\omega_{q,a,b}$ are finite.
\end{theorem}
\begin{proof}
\leanok
\end{proof}
\end{supporting}


\begin{theorem}[Annular radial reduction to the semiclassical hierarchy]
\label{thm:zernike-annulus-radial}\uses{thm:semiclassical-jacobi-orthonormality,thm:annulus-parameter}
\lean{Polynomial.semiclassicalJacobiWeight_annulus, Polynomial.integrable_semiclassicalJacobiMeasure_annulus, Polynomial.integral_semiclassicalJacobiMeasure_annulus, Polynomial.zernikeAnnulusRadial, Polynomial.eval_zernikeAnnulusRadial, Polynomial.integral_zernikeAnnulusRadial_mul}\leanok
Let $q=\rho^2\in(0,1)$, $t=(1-q)^{-1}$, $a,b>-1$, and let $Q_k$ denote
the positive orthonormal semiclassical family with parameters $(t,a,b,m)$.
The radial annular polynomials $A_{k,m}(s)=Q_k(t(1-s))$ satisfy
\[
 \int_q^1(1-s)^a(s-q)^b s^m A_{k,m}(s)A_{\ell,m}(s)\,ds
 = t^{-(a+b+m+1)}\delta_{k\ell}.
\]
Thus the angular frequency is the third semiclassical parameter. Here $s=r^2$;
the full planar pairing additionally has the area factor $1/2$ and the angular
pairing. The exact change-of-variables theorem also applies to every integrable real
input for the semiclassical measure. The transported weighted integrand is
integrable on $(q,1)$, and its integral is multiplied by $t^{a+b+m+1}$.
This includes integrable endpoint singularities, separately from the polynomial
orthogonality specialization.

Reference: \cite[Section 4]{pgso}.
\end{theorem}
\begin{proof}
\leanok
\end{proof}

\begin{supporting}[Formal details: Annular squared-radius integrand and related results]
\begin{definition}[Annular squared-radius integrand]
\label{def:zernike-annulus-radial-input}\uses{thm:zernike-annulus-radial,thm:zernike-annulus-product-measure,def:zernike-angular-input}
\lean{MvPolynomial.zernikeAnnulusSquaredRadiusInput}\leanok
For $0<q<1$, $a,b>-1$ and a Cartesian polynomial $p$, define the function $B_{q,a,b,p}(s)=(1-s)^a(s-q)^bJ_p(s)$ on $s\in(q,1)$.
\end{definition}

\begin{theorem}[Annular radial integrability]
\label{thm:zernike-annulus-radial-input}\uses{def:zernike-annulus-radial-input,thm:integrable-input}
\lean{MvPolynomial.zernikeAnnulusSquaredRadiusInput_integrable}\leanok
The function $B_{q,a,b,p}$ is integrable on $(q,1)$, including admissible singularities at either boundary.
\end{theorem}
\begin{proof}
\leanok
\end{proof}
\end{supporting}


\begin{theorem}[Annular measure and complete real polynomial bases]
\label{thm:zernike-annulus-basis}\uses{def:zernike-annulus-radial-input,thm:zernike-annulus-measure-characterisation,thm:zernike-real-basic,thm:mathlib-total-degree-finite,thm:measure-conversion-cholesky}
\lean{MvPolynomial.integral_zernikeAnnulusMeasure_squaredRadius, MvPolynomial.zernikeAnnulusMode, MvPolynomial.integralInner_zernikeAnnulusMode_same_angular, MvPolynomial.integralInner_zernikeAnnulusMode_eq_zero, MvPolynomial.integral_sq_zernikeAnnulusMode_pos, MvPolynomial.zernikeAnnulusDegreeBasis, MvPolynomial.normalizedZernikeAnnulusDegreeBasis, MvPolynomial.integralGram_normalizedZernikeAnnulusDegreeBasis}\leanok
For $q\in(0,1)$ and $a,b>-1$, the standard annular measure satisfies, for every
Cartesian polynomial $p$,
\[
 \int p\,d\mu=\frac12\int_q^1(1-s)^a(s-q)^b
 \int_{-\pi}^{\pi}p(\sqrt{s}\cos\theta,\sqrt{s}\sin\theta)\,d\theta\,ds.
\]
The real modes $Q_k^{(t,a,b,m)}(t(1-r^2))H_m$, with $t=(1-q)^{-1}$ and
$H_m=r^m\cos(m\theta)$ or $r^m\sin(m\theta)$, are mutually orthogonal.
Exclude the zero sine mode. Their squared norms are
$\frac12t^{-(a+b+m+1)}c_m$, with $c_0=2\pi$ and $c_m=\pi$ for $m>0$.
In the established disk ordering, modes of total degree at most $N$ form a
complete basis of that polynomial space. Positive normalization gives identity
Gram matrices.

Reference: \cite[Definition 4.1 and Proposition 4.2]{pgso}.
\end{theorem}
\begin{proof}
\leanok
\end{proof}

\begin{theorem}[Annular area measure]
\label{thm:annular-area-measure}\uses{def:zernike-annulus-measure,thm:annulus-parameter}
\lean{MvPolynomial.zernikeAnnulusMeasure_zero_zero_eq_volume}\leanok
For $0<q<1$, let $\Omega_q=\{(x,y):q<x^2+y^2<1\}$.
The unweighted annular measure is exactly planar Lebesgue measure restricted
to $\Omega_q$:
\[
 \mu_{q,0,0}=\left.\mathrm{vol}_2\right|_{\Omega_q}.
\]
This is equality of measures, hence applies to arbitrary measurable functions
and their $L^2$ classes, not only to polynomial integrands.
\end{theorem}
\begin{proof}
\leanok
\end{proof}

\begin{supporting}[Reference details: Annular polynomial lift and related results]
\begin{definition}[Annular polynomial lift]
\label{def:annulus-lift}\uses{def:zernike-real,def:annulus-parameter}
\lean{MvPolynomial.annulusLift}\leanok
For $q\ne1$ and a real angular harmonic $H_m$, the linear lift sends $p$ to $L_mp=p(t(q)(1-x^2-y^2))H_m(x,y)$.
\end{definition}

\begin{theorem}[Lift of the semiclassical basis]
\label{thm:annulus-lift}\uses{def:annulus-lift,thm:zernike-annulus-basis,def:semiclassical-jacobi-basis}
\lean{MvPolynomial.annulusLift_basis}\leanok
Let $0<q<1$, $a,b>-1$, and let $p_n$ be the semiclassical basis polynomial
with parameters $(t(q),a,b,m)$. For either the cosine or sine angular harmonic $H_m$, its lift is
\[
 L_mp_n(x,y)=p_n\bigl(t(q)(1-x^2-y^2)\bigr)H_m(x,y),
\]
which is the annular mode with radial index $n$ and angular harmonic $H_m$.
\end{theorem}
\begin{proof}
\leanok
\end{proof}
\end{supporting}


\section{Sparse operators and boundary conditions}\label{sec:annular-operators}

\begin{theorem}[Exact annular parameter conversion]
\label{thm:zernike-annulus-conversion}\uses{thm:semiclassical-conversion,thm:operator-coefficients,def:annulus-lift,thm:annulus-parameter}
\lean{MvPolynomial.annulus_conversion_coefficientMap_exact, MvPolynomial.annulus_conversion_eq_zero}\leanok
Within any fixed real angular sector, the semiclassical connection from $(a,b,m)$
to $(a',b',m)$ reconstructs exactly the same Cartesian annular polynomial from
finite radial coefficients. Thus conversion preserves angular phase and frequency.
For nonnegative integer endpoint shifts $(r,s)$, its entry in radial row $\ell$
and column $k$ vanishes unless $\ell\le k\le\ell+r+s$.
The modes use the raw annular normalization with orthonormal semiclassical radial factors.

Reference: \cite[Section 4.1]{pgso}.
\end{theorem}
\begin{proof}
\leanok
\end{proof}

\rvnote{thm:zernike-annulus-conversion}{
The convention here uses the raw annular modes with orthonormal semiclassical radial factors; their planar norms are given above.
}

\begin{theorem}[Banded Cartesian annular Laplacian]
\label{thm:zernike-annulus-laplacian}\uses{thm:semiclassical-partial-weighted-differentiation,thm:zernike-annulus-basis,thm:zernike-polynomials,def:annulus-lift}
\lean{Polynomial.semiclassicalJacobiAnnularLaplacianMap, Polynomial.semiclassicalJacobiAnnularLaplacian_eq_zero, Polynomial.annularLaplacian_comp, MvPolynomial.diskLaplacian_zernikeAnnulusMode, MvPolynomial.diskLaplacian_annulusLift, MvPolynomial.annulus_laplacian_coefficientMap_exact}\leanok
For each angular frequency $m$, the Cartesian annular Laplacian is represented by
\[
 L_{t,m}p=4t\bigl((t-x)p''-(m+1)p'\bigr)=4tD_cDp,
\]
from radial parameters $(a,b,m)$ to $(a+2,b+2,m)$. It preserves angular phase
and frequency. Column $k$ has support only in radial rows $k-1,k-2,k-3$;
negative rows are absent. Acting on finite coefficients and then reconstructing
the Cartesian polynomial agrees exactly with applying the Cartesian Laplacian.


Reference: \cite[Proposition 4.5, equation (4.13)]{pgso}.
\end{theorem}
\begin{proof}
\leanok
\end{proof}

\rvnote{thm:zernike-annulus-laplacian}{
This is an algebraic operator identity; no boundary-value well-posedness or solver convergence is asserted here.
}

For the Dirichlet problem, the finite systems are Galerkin systems for an energy
inner product. They are definite, hence uniquely solvable at every size, and they
give the best approximation in energy. Completing the polynomials in that energy
turns this into convergence of the actual coefficient solves, within each angular
sector.

\begin{supporting}[Reference details: Nondegenerate annular Dirichlet weight and related results]
\begin{definition}[Nondegenerate annular Dirichlet weight]
\label{def:annular-dirichlet-weight}
\lean{Polynomial.annularDirichletWeight}\leanok
For $t\in\mathbb R\setminus\{0\}$, define the polynomial
$W_t(X)=t^{-2}X(1-X)$.
\end{definition}

\begin{theorem}[Evaluation of the annular Dirichlet weight]
\label{thm:annular-dirichlet-weight}\uses{def:annular-dirichlet-weight}
\lean{Polynomial.eval_annularDirichletWeight}\leanok
For $t\ne0$ and every real $x$, $W_t(x)=t^{-2}x(1-x)$.
\end{theorem}
\begin{proof}
\leanok
\end{proof}

\begin{definition}[Nondegenerate radial Dirichlet operator]
\label{def:annular-dirichlet-operator}\uses{thm:semiclassical-partial-weighted-differentiation}
\lean{Polynomial.semiclassicalJacobiABWeightedDerivativeMap, Polynomial.semiclassicalJacobiDirichletMap}\leanok
For $t\in\mathbb R\setminus\{0\}$ and $\gamma\in\mathbb R$, set
$Ep=(X(1-X)p)'$ and
$A_{t,\gamma}p=(4/t)((t-X)(Ep)'-(\gamma+1)Ep)$.
These are linear maps on real polynomials.
\end{definition}

\begin{theorem}[Radial Dirichlet operator characterization]
\label{thm:annular-dirichlet-operator}\uses{def:annular-dirichlet-operator,thm:zernike-annulus-laplacian,def:annular-dirichlet-weight}
\lean{Polynomial.semiclassicalJacobiABWeightedDerivativeMap_apply, Polynomial.semiclassicalJacobiDirichletMap_eq}\leanok
For $t\ne0$, $Ep=X(1-X)p'+(1-2X)p$, and
$A_{t,\gamma}p=L_{t,\gamma}(W_tp)$, where $L_{t,\gamma}$ is the radial Laplacian.
\end{theorem}
\begin{proof}
\leanok
\end{proof}
\end{supporting}


\begin{theorem}[Annular Dirichlet energy identity]
\label{thm:annular-dirichlet-form}\uses{thm:annular-dirichlet-operator}
\lean{Polynomial.integralInner_semiclassicalJacobi_dirichlet, Polynomial.integralInner_semiclassicalJacobi_dirichlet_symm, Polynomial.integralInner_semiclassicalJacobi_dirichlet_neg, Polynomial.semiclassicalJacobiDirichletMap_injective}\leanok
Let $t>1$ and $\gamma\in\mathbb R$. Set $S(x)=x(1-x)$, $Ep=(Sp)'$, and
\[
 A p=\frac4t\bigl((t-x)(Ep)'-(\gamma+1)Ep\bigr).
\]
Then $A p=L_{t,\gamma}(t^{-2}Sp)$, with the radial Laplacian expression above.
In the semiclassical pairings,
\[
 \langle Ap,z\rangle_{1,1,\gamma}
 =-\frac4t\langle Ep,Ez\rangle_{0,0,\gamma+1}.
\]
Thus this form is symmetric, and $\langle Ap,p\rangle_{1,1,\gamma}<0$ for
$p\ne0$. In particular, $A$ is injective on polynomials.
\end{theorem}
\begin{proof}
\leanok
\end{proof}

\begin{theorem}[Symmetric annular Dirichlet and mass matrices]
\label{thm:annular-dirichlet-matrices}\uses{thm:annular-dirichlet-form}
\lean{Polynomial.semiclassicalJacobiDirichletMatrix, Polynomial.semiclassicalJacobiDirichletMatrix_symm, Polynomial.semiclassicalJacobiDirichletMatrix_eq_zero, Polynomial.semiclassicalJacobi_dirichlet_coefficientMap_exact, Polynomial.semiclassicalJacobiDirichletMass, Polynomial.semiclassicalJacobiDirichletMass_symm, Polynomial.semiclassicalJacobiDirichletMass_eq_zero, Polynomial.integralInner_annularDirichletWeight_pos}\leanok
In the positive orthonormal semiclassical basis with parameters $(t,1,1,\gamma)$,
the actual coefficient matrix of $A$ is symmetric and tridiagonal: entry $(i,j)$
vanishes if $j+1<i$ or $i+1<j$. Its coefficient action is exact on finite expansions.
Multiplication by the physical boundary factor $t^{-2}x(1-x)$ has a symmetric
pentadiagonal coefficient matrix: entry $(i,j)$ vanishes if $j+2<i$ or $i+2<j$.
The mass pairing is strictly positive on nonzero polynomials.

\end{theorem}
\begin{proof}
\leanok
\end{proof}

\rvnote{thm:annular-dirichlet-matrices}{
These support bounds do not assert that every permitted entry is nonzero.
}

\begin{supporting}[Reference details: Boundary-vanishing annular lift and related results]
\begin{definition}[Boundary-vanishing annular lift]
\label{def:annulus-dirichlet-lift}\uses{def:annulus-lift,def:polynomial-linear-maps}
\lean{MvPolynomial.annulusBoundaryWeight, MvPolynomial.annulusDirichletLift}\leanok
For $q\ne1$, write $B_q=(1-x^2-y^2)(x^2+y^2-q)$. Define the boundary-weighted lift by $p\mapsto B_qL_mp$.
\end{definition}

\begin{theorem}[Boundary values of the weighted lift]
\label{thm:annulus-dirichlet-lift}\uses{def:annulus-dirichlet-lift}
\lean{MvPolynomial.eval_annulusDirichletLift_eq_zero}\leanok
Every lifted polynomial $B_qL_mp$ vanishes at each real point satisfying $x^2+y^2=q$ or $x^2+y^2=1$.
\end{theorem}
\begin{proof}
\leanok
\end{proof}
\end{supporting}


\begin{theorem}[Exact boundary-vanishing Cartesian annular operators]
\label{thm:zernike-annulus-dirichlet}\uses{def:annulus-dirichlet-lift,thm:annular-dirichlet-operator}
\lean{Polynomial.annularDirichletWeight_comp, MvPolynomial.diskLaplacian_annulusDirichletLift, MvPolynomial.annulus_dirichlet_coefficientMap_exact, MvPolynomial.annulus_dirichlet_mass_coefficientMap_exact}\leanok
Let $q\in(0,1)$ and $t=(1-q)^{-1}$. For every radial polynomial $p$ and either
real solid harmonic $H_m$, define
\[
 W_p=(1-r^2)(r^2-q)\,p(t(1-r^2))H_m.
\]
This polynomial vanishes on both circles $r^2=q$ and $r^2=1$, and
\[
 \Delta W_p=(Ap)(t(1-r^2))H_m,\qquad \gamma=m.
\]
Hence each angular sector has the symmetric tridiagonal Dirichlet block and
pentadiagonal mass block above. The corresponding finite-coefficient operations
reconstruct these Cartesian polynomial identities exactly, with no coupling of
angular frequency or phase.

Reference: \cite[Proposition 4.6, equation (4.15)]{pgso}.
\end{theorem}
\begin{proof}
\leanok
\end{proof}

\rvnote{thm:zernike-annulus-dirichlet}{
Finite-section solvability and energy convergence are given in Theorems~\ref{thm:annular-dirichlet-sections} and~\ref{thm:annular-energy-convergence}.
}

\section{Finite systems and Galerkin approximation}\label{sec:annular-galerkin}

\begin{theorem}[Definiteness and invertibility of annular Dirichlet sections]
\label{thm:annular-dirichlet-sections}\uses{thm:annular-dirichlet-matrices,thm:finite-polynomial-galerkin,thm:mathlib-positive-invertible}
\lean{Polynomial.semiclassicalJacobiDirichletSection, Polynomial.semiclassicalJacobiDirichletMassSection, Polynomial.neg_semiclassicalJacobiDirichletSection_posDef, Polynomial.semiclassicalJacobiDirichletMassSection_posDef, Polynomial.isUnit_semiclassicalJacobiDirichletSection, Polynomial.isUnit_semiclassicalJacobiDirichletMassSection, Polynomial.existsUnique_semiclassicalJacobiDirichletSection}\leanok
For $t>1$, real $\gamma$, and every $N\ge0$, the leading $N\times N$ Dirichlet
section $A_N$ is negative definite, meaning $-A_N$ is positive definite.
The corresponding mass section is positive definite. Both are invertible;
in particular, $A_Nv=g$ has exactly one solution for every finite right-hand side.
\end{theorem}
\begin{proof}
\leanok
\end{proof}

\paragraph{Example (Julia).} For $t=2$, that is $q=1/2$, and angular mode
$\gamma=3$, generate the orthonormal family for $x(1-x)(t-x)^3$ on $(0,1)$ by
the Stieltjes procedure on Gauss--Legendre nodes, together with its first two
derivatives. The $10\times10$ matrix of $A$ is symmetric and tridiagonal up to
rounding, and negative definite, as
Theorems~\ref{thm:annular-dirichlet-matrices}
and~\ref{thm:annular-dirichlet-sections} state.
\begin{verbatim}
using LinearAlgebra
t, g, N = 2.0, 3.0, 10              # annulus q = 1/2, angular mode 3
n = 120                             # Gauss-Legendre nodes on (0,1)
E = eigen(SymTridiagonal(zeros(n), [k/sqrt(4k^2 - 1) for k in 1:n-1]))
x = (E.values .+ 1) ./ 2
w = E.vectors[1, :].^2 .* x .* (1 .- x) .* (t .- x).^g
# Stieltjes: orthonormal Q_k for x(1-x)(t-x)^g dx, with Q_k', Q_k''
Q, D1, D2, beta = [ones(n) / sqrt(sum(w))], [zeros(n)], [zeros(n)], [0.0]
for k in 0:N
    a = dot(w, x .* Q[end].^2)
    r  = (x .- a) .* Q[end]
    r1 = Q[end] .+ (x .- a) .* D1[end]
    r2 = 2 .* D1[end] .+ (x .- a) .* D2[end]
    if k > 0
        r -= beta[end] .* Q[end-1]
        r1 -= beta[end] .* D1[end-1]
        r2 -= beta[end] .* D2[end-1]
    end
    push!(beta, sqrt(dot(w, r.^2)))
    push!(Q, r / beta[end])
    push!(D1, r1 / beta[end])
    push!(D2, r2 / beta[end])
end
S, S1 = x .* (1 .- x), 1 .- 2x      # S = x(1-x), S'' = -2
E1(j) = S1 .* Q[j] .+ S .* D1[j]    # Ep = (Sp)'
E2(j) = -2 .* Q[j] .+ 2 .* S1 .* D1[j] .+ S .* D2[j]   # (Ep)'
Ap(j) = (4/t) .* ((t .- x) .* E2(j) .- (g + 1) .* E1(j))
A = [dot(w, Ap(j) .* Q[i]) for i in 1:N, j in 1:N]
off = maximum(abs(A[i, j]) for i in 1:N, j in 1:N if abs(i - j) > 1)
(norm(A - A') / norm(A), off / norm(A), maximum(eigvals(Symmetric(A))))
                                    # (~1e-14, ~1e-15, negative)
\end{verbatim}

\begin{theorem}[Polynomial Galerkin interpretation of annular sections]
\label{thm:annular-polynomial-galerkin}\uses{thm:annular-dirichlet-sections}
\lean{Polynomial.semiclassicalJacobiDirichletTrial, Polynomial.semiclassicalJacobiDirichletLoad, Polynomial.semiclassicalJacobiDirichletSection_eq_iff_galerkin, Polynomial.existsUnique_semiclassicalJacobiDirichletGalerkin, Polynomial.semiclassicalJacobiDirichletGalerkin_exact_on_trial}\leanok
Let $V_N$ be the span of the first $N$ semiclassical polynomials with parameters
$(t,1,1,\gamma)$. For polynomial forcing $f$, the square-section system is exactly
\[
 \langle Ap_N-f,z\rangle_{1,1,\gamma}=0\qquad(z\in V_N).
\]
There is exactly one coefficient vector solving these equations. If the exact
radial solution already belongs to $V_N$, the finite system recovers its coefficients.
This is residual orthogonality, not an assertion that the full polynomial
residual vanishes for arbitrary forcing.
\end{theorem}
\begin{proof}
\leanok
\end{proof}

\begin{supporting}[Reference details: Polynomial energy best approximation]
\begin{theorem}[Polynomial energy best approximation]
\label{thm:annular-polynomial-best}\uses{thm:annular-polynomial-galerkin}
\lean{Polynomial.semiclassicalJacobiDirichletEnergy, Polynomial.semiclassicalJacobiDirichletEnergy_nonneg, Polynomial.semiclassicalJacobiDirichletEnergy_pythagoras, Polynomial.semiclassicalJacobiDirichletGalerkin_bestApproximation}\leanok
For $t>1$ and real $\gamma$, write
$E(p)=-\langle Ap,p\rangle_{1,1,\gamma}$; finite moments ensure this pairing exists.
If $u$ is a polynomial and
$u_N\in V_N$ solves the Galerkin system with forcing $Au$, then
\[
 E(u-u_N)\le E(u-z)\qquad(z\in V_N).
\]
Energy-orthogonal polynomial vectors satisfy the exact Pythagorean identity.
\end{theorem}
\begin{proof}
\leanok
\end{proof}
\end{supporting}


\section{Energy convergence and physical realization}\label{sec:annular-energy}

\begin{supporting}[Reference details: Annular radial energy completion and related results]
\begin{definition}[Annular radial energy completion]
\label{def:annular-energy-completion}\uses{thm:annular-polynomial-best}
\lean{Polynomial.SemiclassicalJacobiEnergy.PolynomialSpace, Polynomial.SemiclassicalJacobiEnergy.innerCore, Polynomial.SemiclassicalJacobiEnergy.Space, Polynomial.SemiclassicalJacobiEnergy.embed}\leanok
For fixed $t>1$ and real $\gamma$, equip radial polynomials with the inner
product $(p,z)_E=-\langle Ap,z\rangle_{1,1,\gamma}$ of
Theorem~\ref{thm:annular-dirichlet-form}, and let $H_{t,\gamma}$ be its Hilbert
completion, with the canonical isometric polynomial embedding $\iota$. For an
annular angular sector, set $\gamma=m$.
\end{definition}

\begin{theorem}[Energy norm, pairing and density]
\label{thm:annular-energy-completion-basic}\uses{def:annular-energy-completion}
\lean{Polynomial.SemiclassicalJacobiEnergy.norm_sq_eq_energy, Polynomial.SemiclassicalJacobiEnergy.inner_embed, Polynomial.SemiclassicalJacobiEnergy.denseRange_embed}\leanok
The energy norm satisfies $\|p\|_E^2=E(p)$. The embedding satisfies
$(\iota p,\iota z)=-\langle Ap,z\rangle_{1,1,\gamma}$, and its range is dense in
$H_{t,\gamma}$.
\end{theorem}
\begin{proof}
\leanok
\end{proof}
\end{supporting}


The planar Sobolev identification follows after assembling the angular sectors;
see Theorem~\ref{thm:annular-sobolev-identification}.

\begin{theorem}[Convergence of actual annular coefficient solves in energy]
\label{thm:annular-energy-convergence}\uses{thm:annular-energy-completion-basic,thm:mathlib-hilbert-projection-limit}
\lean{Polynomial.SemiclassicalJacobiEnergy.trialSpace, Polynomial.SemiclassicalJacobiEnergy.trialSpace_eq_span, Polynomial.SemiclassicalJacobiEnergy.trialSpace_closure_iSup, Polynomial.SemiclassicalJacobiEnergy.solution, Polynomial.SemiclassicalJacobiEnergy.existsUnique_solution, Polynomial.SemiclassicalJacobiEnergy.coefficients, Polynomial.SemiclassicalJacobiEnergy.coefficients_spec, Polynomial.SemiclassicalJacobiEnergy.trialEmbed_coefficients_eq_projection, Polynomial.SemiclassicalJacobiEnergy.coefficients_bestApproximation, Polynomial.SemiclassicalJacobiEnergy.norm_trialEmbed_coefficients_le, Polynomial.SemiclassicalJacobiEnergy.tendsto_trialEmbed_coefficients}\leanok
Let $\ell$ be any continuous linear functional on $H_{t,\gamma}$, and let $u$ be
the unique energy solution $(u,z)_E=\ell(z)$ for every $z\in H_{t,\gamma}$.
Let $Q_j$ denote the embedded $(t,1,1,\gamma)$ semiclassical trial polynomials.
Solve the actual finite Dirichlet system
\[
 A_Nv^{(N)}=\bigl(-\ell(Q_i)\bigr)_{i<N},\qquad
 u_N=\sum_{j<N}v_j^{(N)}Q_j.
\]
Then $u_N$ is exactly the energy orthogonal projection of $u$ onto $V_N$.
Consequently
\[
 \|u-u_N\|_E\le\|u-z\|_E\quad(z\in V_N),\qquad
 \|u_N\|_E\le\|\ell\|,\qquad u_N\longrightarrow u\text{ in }H_{t,\gamma}.
\]
This is radial convergence in a fixed angular sector when $\gamma=m$.

\end{theorem}
\begin{proof}
\leanok
\end{proof}

\rvnote{thm:annular-energy-convergence}{
No regularity-dependent convergence rate is asserted here.
}

Distinct angular sectors are orthogonal in the physical energy, so the sector
problems combine. Summing the sector energy spaces gives a completed planar
energy space. In it, solving the tridiagonal radial system of
Theorem~\ref{thm:annular-dirichlet-matrices} in each
retained sector yields the orthogonal projection onto the joint trial space, and
these solves converge as the retained sectors and radial degrees grow together.

\begin{supporting}[Reference details: Annular sector energy scale and related results]
\begin{definition}[Annular sector energy scale]
\label{def:annulus-energy-scale}\uses{def:annulus-parameter}
\lean{MvPolynomial.annulusEnergyScale}\leanok
For $q\ne1$ and an angular mode, set $\kappa_m=c_m/(2t(q)^{m+3})$, using the natural-number power. Here $c_m$ is the squared angular norm, with $c_m=0$ for the zero sine mode.
\end{definition}

\begin{theorem}[Positivity of the sector energy scale]
\label{thm:annulus-energy-scale}\uses{def:annulus-energy-scale,thm:annulus-parameter}
\lean{MvPolynomial.annulusEnergyScale_pos}\leanok
For $0<q<1$, the scale $\kappa_m$ is positive for every nonzero angular mode.
\end{theorem}
\begin{proof}
\leanok
\end{proof}

\begin{theorem}[Physical annular sector energy]
\label{thm:annular-physical-energy}\uses{thm:zernike-annulus-dirichlet,thm:annular-energy-completion-basic,def:annulus-energy-scale}
\lean{MvPolynomial.integralInner_annulusLift_dirichletLift, MvPolynomial.integralInner_annulusDirichletLaplacian_eq_zero, MvPolynomial.neg_integralInner_annulusDirichletLaplacian}\leanok
Let $0<q<1$, $t=(1-q)^{-1}$, and $B=(1-r^2)(r^2-q)$.
For a valid real angular harmonic $H_m$ (excluding the zero sine mode), set
$L_m p=p(t(1-r^2))H_m$ and
\[
 \kappa_m=\frac{c_m}{2t^{m+3}}>0,
 \qquad c_0=2\pi\text{ for the constant cosine mode},\quad c_m=\pi\text{ otherwise}.
\]
With the unweighted annular measure $\mu_{q,0,0}$,
\[
 -\int \Delta(BL_m p)\,BL_m z\,d\mu_{q,0,0}
   =\kappa_m(p,z)_{H_{t,m}}.
\]
Pairings between distinct real angular sectors vanish.
\end{theorem}
\begin{proof}
\leanok
\end{proof}

\begin{definition}[Completed angular annular energy space]
\label{def:annular-angular-energy}\uses{def:annular-energy-completion,def:annulus-energy-scale}
\lean{MvPolynomial.AnnularEnergy.Sector, MvPolynomial.AnnularEnergy.Space, MvPolynomial.AnnularEnergy.polynomialEmbed}\leanok
Let $\mathcal H_q=\bigoplus_{s}^{\ell^2}H_{t,m(s)}$, indexed by all real cosine
and nonzero sine sectors $s$. A physical polynomial $BL_{m(s)}p$ embeds into
sector $s$ as $\sqrt{\kappa_{m(s)}}\,p$.
\end{definition}

\begin{theorem}[Pairing, orthogonality and density]
\label{thm:annular-angular-energy-basic}\uses{def:annular-angular-energy,thm:annular-energy-convergence,thm:annulus-energy-scale,thm:annular-physical-energy,thm:hilbert-sum-galerkin}
\lean{MvPolynomial.AnnularEnergy.inner_polynomialEmbed, MvPolynomial.AnnularEnergy.inner_polynomialEmbed_eq_zero, MvPolynomial.AnnularEnergy.trialSpace_closure_iSup}\leanok
Write $U_p=BL_{m(s)}p$. The embedding of sector $s$ preserves the planar
negative-Laplacian pairing, $(\iota_sp,\iota_sz)=-\int\Delta U_p\,U_z\,d\mu$,
where $\mu$ is the annular measure with $a=b=0$. Embeddings of distinct sectors
are orthogonal. The finite angular and radial trial spaces have dense union in
$\mathcal H_q$.
\end{theorem}
\begin{proof}
\leanok
\end{proof}
\end{supporting}


The following theorem gives convergence of the assembled radial solves as
both angular and radial resolution increase.

\begin{theorem}[Joint angular and radial convergence of annular solves]
\label{thm:annular-angular-convergence}\uses{thm:annular-angular-energy-basic}
\lean{MvPolynomial.AnnularEnergy.solution, MvPolynomial.AnnularEnergy.existsUnique_solution, MvPolynomial.AnnularEnergy.solution_apply, MvPolynomial.AnnularEnergy.finiteSolve, MvPolynomial.AnnularEnergy.finiteSolve_eq_projection, MvPolynomial.AnnularEnergy.finiteSolve_bestApproximation, MvPolynomial.AnnularEnergy.norm_finiteSolve_le, MvPolynomial.AnnularEnergy.tendsto_finiteSolve, MvPolynomial.AnnularEnergy.mem_frequencyCutoff, MvPolynomial.AnnularEnergy.tendsto_finiteSolve_frequencyCutoff}\leanok
For any continuous real linear load $\ell$ on $\mathcal H_q$, there is a unique
$u$ satisfying $(u,z)=\ell(z)$ for all $z\in\mathcal H_q$.
Restrict the load to each normalized sector and solve its existing tridiagonal
radial Dirichlet system. The sum $u_{S,N}$ of the retained sector reconstructions
is exactly the orthogonal projection of $u$ onto $V_{S,N}$. Hence
\[
 \|u-u_{S,N}\|\le\|u-z\|\quad(z\in V_{S,N}),\qquad
 \|u_{S,N}\|\le\|\ell\|,\qquad u_{S,N}\longrightarrow u.
\]
In particular, taking every valid angular frequency $m<N$ and radial degree
below $N$ gives convergence as $N\to\infty$. The coordinates use the normalization
in Definition~\ref{def:annular-angular-energy}.
\end{theorem}
\begin{proof}
\leanok
\end{proof}

\rvnote{thm:annular-angular-convergence}{No regularity rate or boundedness of arbitrary physical forcing as an energy load is asserted.}


Green's identity identifies the energy with the Cartesian Dirichlet integral
$\int\nabla U\cdot\nabla V\,dA$ of boundary-weighted lifts, sector by sector.
A weighted derivative embeds the completed radial energy isometrically into
an $L^2$ space, where the derivatives of the computed solutions converge.
Theorem~\ref{thm:annular-sobolev-identification} below identifies the completed
planar space with physical $H^1_0$ functions.

\begin{supporting}[Reference details: Polynomial boundary flux and Green identity on an annulus and related results]
\begin{theorem}[Polynomial boundary flux and Green identity on an annulus]
\label{thm:annular-polynomial-green}\uses{thm:annular-physical-energy,thm:annulus-lift}
\lean{MvPolynomial.eq_top_of_annulusLift_mem, MvPolynomial.diskLaplacian_mul, MvPolynomial.integral_diskLaplacian_diskRadialLift, MvPolynomial.integral_diskLaplacian_boundary_sq, MvPolynomial.integralInner_annulus_boundary_laplacian_symm, MvPolynomial.integralInner_annulus_gradient_eq_neg_laplacian}\leanok
Let $0<q<1$, let $dA$ be the unweighted annular measure, and put
$B_q=(1-x^2-y^2)(x^2+y^2-q)$. For arbitrary Cartesian polynomials $p,z$,
\[
 \int \Delta(B_q^2p)\,dA=0,\qquad
 \int \nabla(B_qp)\cdot\nabla(B_qz)\,dA
 =-\int \Delta(B_qp)\,B_qz\,dA.
\]
The Laplacian pairing on this boundary-weighted polynomial space is symmetric.
More explicitly, the integrated Laplacian of a radial lift with radial polynomial
$F(s)$ is zero in every nonconstant angular sector, and is
$4\pi(F'(1)-qF'(q))$ in the constant sector.
\end{theorem}
\begin{proof}
\leanok
\end{proof}

\begin{theorem}[Cartesian gradient interpretation of annular sector energy]
\label{thm:annular-cartesian-energy}\uses{thm:annular-polynomial-green}
\lean{MvPolynomial.integralInner_annulusDirichlet_gradient, MvPolynomial.integralInner_annulusDirichlet_gradient_eq_zero}\leanok
For boundary-weighted lifts $U_p,U_z$ in a real angular sector $(u,m)$,
\[
 \int \nabla U_p\cdot\nabla U_z\,dA
 =\kappa_{q,u,m}\langle\iota p,\iota z\rangle_{H_{t,m}},
 \qquad t=(1-q)^{-1}.
\]
Distinct angular sectors are orthogonal for this Cartesian gradient pairing.

\end{theorem}
\begin{proof}
\leanok
\end{proof}

\rvnote{thm:annular-cartesian-energy}{
The identification of completed energy vectors with Sobolev functions and their weak gradients is Theorem~\ref{thm:annular-sobolev-identification}.
}

\begin{theorem}[Concrete weighted derivative realization of radial energy]
\label{thm:annular-derivative-realization}\uses{thm:annular-energy-convergence,thm:mathlib-isometry-completion,thm:general-polynomial-L2-basic}
\lean{Polynomial.SemiclassicalJacobiEnergy.DerivativeSpace, Polynomial.SemiclassicalJacobiEnergy.polynomialDerivative, Polynomial.SemiclassicalJacobiEnergy.inner_polynomialDerivative, Polynomial.SemiclassicalJacobiEnergy.derivativeRealization, Polynomial.SemiclassicalJacobiEnergy.derivativeRealization_embed, Polynomial.SemiclassicalJacobiEnergy.derivativeRealization_injective, Polynomial.SemiclassicalJacobiEnergy.range_derivativeRealization, Polynomial.SemiclassicalJacobiEnergy.inner_derivativeRealization, Polynomial.SemiclassicalJacobiEnergy.derivativeLoad_apply, Polynomial.SemiclassicalJacobiEnergy.norm_derivativeLoad_le, Polynomial.SemiclassicalJacobiEnergy.norm_coefficients_derivativeLoad_le, Polynomial.SemiclassicalJacobiEnergy.tendsto_derivativeRealization_coefficients}\leanok
For $t>1$ and real $\gamma$, the map
\[
 Dp=\sqrt{4/t}\,\bigl(x(1-x)p\bigr)'
\]
from radial energy polynomials to $L^2(\mu_{t,0,0,\gamma+1})$ extends to a linear
isometry $D:H_{t,\gamma}\to L^2(\mu_{t,0,0,\gamma+1})$. Its range is exactly the
closed $L^2$ closure of these polynomial derivative classes; no surjectivity onto
all of $L^2$ is claimed. For every $g$ in this $L^2$ space,
$\ell_g(u)=\langle g,Du\rangle$ is a continuous energy load of norm at most
$\|g\|_{L^2}$, and its finite coefficient reconstructions are uniformly bounded
by this norm. The existing energy convergence also gives strong $L^2$ convergence
of the realized derivatives.
\end{theorem}
\begin{proof}
\leanok
\end{proof}

\begin{theorem}[Concrete derivative convergence in all angular sectors]
\label{thm:annular-angular-derivatives}\uses{thm:annular-derivative-realization,thm:annular-angular-convergence}
\lean{MvPolynomial.AnnularEnergy.DerivativeSpace, MvPolynomial.AnnularEnergy.derivativeRealization, MvPolynomial.AnnularEnergy.derivativeRealization_apply, MvPolynomial.AnnularEnergy.derivativeRealization_injective, MvPolynomial.AnnularEnergy.isClosed_range_derivativeRealization, MvPolynomial.AnnularEnergy.norm_sub_derivativeRealization_finiteSolve, MvPolynomial.AnnularEnergy.tendsto_derivativeRealization_finiteSolve}\leanok
Applying $D$ in each normalized real angular sector embeds $\mathcal H_q$
isometrically into the Hilbert sum of the sector derivative $L^2$ spaces, with
closed range. The error of an actual finite annular solve in this concrete space
equals its energy error. Thus the frequency/radial-degree truncations converge
strongly as radial derivative classes as well.
\end{theorem}
\begin{proof}
\leanok
\end{proof}

To apply a forcing function $g\in L^2(dA)$, we reconstruct a function class
on the annulus from each abstract energy vector. The following maps give this
reconstruction and the load functional. Injectivity ensures that distinct
energy vectors represent distinct function classes.

\begin{theorem}[Physical $L^2$ realization and body forcing in a completed sector]
\label{thm:annular-radial-L2}\uses{thm:annular-poincare,thm:multivariate-polynomial-L2,thm:annular-energy-convergence,def:annular-angular-energy,thm:annulus-energy-scale}
\lean{MvPolynomial.AnnularEnergy.PhysicalL2, MvPolynomial.AnnularEnergy.norm_radialPolynomialL2_le, MvPolynomial.AnnularEnergy.radialL2, MvPolynomial.AnnularEnergy.radialL2_polynomialEmbed, MvPolynomial.AnnularEnergy.norm_radialL2_le, MvPolynomial.AnnularEnergy.radialBodyLoad, MvPolynomial.AnnularEnergy.radialBodyLoad_polynomialEmbed, MvPolynomial.AnnularEnergy.norm_radialBodyLoad_le, MvPolynomial.AnnularEnergy.radialPoissonLoad, MvPolynomial.AnnularEnergy.radialPoisson_solution_weak, MvPolynomial.AnnularEnergy.tendsto_radialL2_coefficients}\leanok
Fix a valid real angular sector $s=(u,m)$ and let $\kappa_s$ be its physical
energy scale. Polynomial lifting extends to a continuous linear map
$T_s:H_{t,m}\to L^2(dA)$ with
\[
 \|T_s v\|_{L^2}\leq\sqrt{\kappa_s}\|v\|_{H_{t,m}}.
\]
Every $g\in L^2(dA)$ defines a body load
$b_g(v)=\langle g,T_sv\rangle_{L^2}$, of norm at most
$\sqrt{\kappa_s}\|g\|_{L^2}$. On polynomial trials this is the physical integral
$\int gU_p\,dA$. Feeding the unscaled radial solver the load
$\kappa_s^{-1}b_g$ gives a solution satisfying
\[
 \kappa_s\langle v,z\rangle_{H_{t,m}}=\langle g,T_sz\rangle_{L^2}
 \quad\text{for all }z\in H_{t,m}.
\]
The existing radial coefficient reconstructions converge after applying $T_s$,
for any continuous radial load.
\end{theorem}
\begin{proof}
\leanok
\end{proof}

\begin{theorem}[Faithfulness of the radial physical realization]
\label{thm:annular-radial-L2-injective}\uses{thm:annular-radial-L2}
\lean{MvPolynomial.AnnularEnergy.inner_radialL2_polynomialLaplacian, MvPolynomial.AnnularEnergy.radialL2_injective}\leanok
For a valid sector $s=(u,m)$, the completed map $T_s:H_{t,m}\to L^2(dA)$ is
injective. More precisely, for every completed energy vector $x$ and radial
polynomial $p$,
\[
 \kappa_s\langle x,\iota p\rangle_{H_{t,m}}
 =-\langle T_sx,\Delta U_p\rangle_{L^2(dA)}.
\]
\end{theorem}
\begin{proof}
\leanok
\end{proof}

\begin{theorem}[Faithful physical $L^2$ realization of all annular sectors]
\label{thm:annular-full-L2}\uses{thm:orthogonal-contraction-synthesis,thm:annular-radial-L2-injective}
\lean{MvPolynomial.AnnularEnergy.inner_radialL2_eq_zero, MvPolynomial.AnnularEnergy.sectorL2, MvPolynomial.AnnularEnergy.norm_sectorL2_le, MvPolynomial.AnnularEnergy.physicalL2, MvPolynomial.AnnularEnergy.physicalL2_injective, MvPolynomial.AnnularEnergy.hasSum_physicalL2, MvPolynomial.AnnularEnergy.norm_physicalL2_le, MvPolynomial.AnnularEnergy.physicalL2_polynomialEmbed}\leanok
The full annular energy Hilbert sum $\mathcal H_q$ has an injective continuous
linear realization $T:\mathcal H_q\to L^2(dA)$ given by
\[
 Tx=\sum_s\kappa_s^{-1/2}T_sx_s,\qquad
 \|Tx\|_{L^2}\leq\|x\|_{\mathcal H_q}.
\]
The series converges unconditionally in physical $L^2$. Each physically scaled
polynomial energy vector realizes exactly its boundary-weighted polynomial
$U_p$.
\end{theorem}
\begin{proof}
\leanok
\end{proof}

A physical $L^2$ body force determines a unique solution in the energy
completion. The finite reconstructions converge to it in physical $L^2$,
as the following theorem shows.

\begin{theorem}[Physical body forcing and convergence of assembled annular solves]
\label{thm:annular-body-forcing}\uses{thm:annular-full-L2,thm:annular-angular-convergence}
\lean{MvPolynomial.AnnularEnergy.bodyLoad, MvPolynomial.AnnularEnergy.norm_bodyLoad_le, MvPolynomial.AnnularEnergy.bodyLoad_polynomialEmbed, MvPolynomial.AnnularEnergy.sectorLoad_bodyLoad, MvPolynomial.AnnularEnergy.existsUnique_body_solution, MvPolynomial.AnnularEnergy.norm_solution_bodyLoad_le, MvPolynomial.AnnularEnergy.norm_finiteSolve_bodyLoad_le, MvPolynomial.AnnularEnergy.norm_physicalL2_finiteSolve_bodyLoad_le, MvPolynomial.AnnularEnergy.physicalL2_finiteSolve, MvPolynomial.AnnularEnergy.physicalL2_finiteSolve_bestApproximation, MvPolynomial.AnnularEnergy.tendsto_physicalL2_finiteSolve, MvPolynomial.AnnularEnergy.tendsto_physicalL2_finiteSolve_frequencyCutoff}\leanok
Every physical datum $g\in L^2(dA)$ gives a continuous energy load
$\ell_g(x)=\langle g,Tx\rangle_{L^2}$ with $\|\ell_g\|\leq\|g\|_{L^2}$.
Its value on a polynomial trial is $\int gU_p\,dA$. There is a unique
$x\in\mathcal H_q$ satisfying
\[
 \langle x,z\rangle_{\mathcal H_q}=\langle g,Tz\rangle_{L^2}
 \quad\text{for all }z\in\mathcal H_q.
\]
The full solution and every assembled finite coefficient solve have energy norm
at most $\|g\|_{L^2}$; the finite physical realizations have the same bound.
The physical reconstruction is precisely the finite sum of computed radial
reconstructions with their inverse square-root sector scales. For any continuous
energy load, these reconstructions converge to $Tx$ in physical $L^2$ as both
cutoffs increase, also for the common frequency/radial-degree cutoff. Their
physical error is bounded by the best energy approximation error in the finite
trial space.
\end{theorem}
\begin{proof}
\leanok
\end{proof}

\begin{theorem}[Completed Cartesian gradient in each annular sector]
\label{thm:annular-sector-gradient}\uses{thm:annular-radial-L2,thm:mathlib-isometry-completion,thm:mathlib-hilbert-sums,thm:annular-cartesian-energy}
\lean{MvPolynomial.AnnularEnergy.GradientL2, MvPolynomial.AnnularEnergy.polynomialGradient, MvPolynomial.AnnularEnergy.polynomialGradient_apply, MvPolynomial.AnnularEnergy.inner_polynomialGradient, MvPolynomial.AnnularEnergy.inner_radialPolynomialGradient, MvPolynomial.AnnularEnergy.inner_normalizedPolynomialGradient, MvPolynomial.AnnularEnergy.sectorGradient, MvPolynomial.AnnularEnergy.sectorGradient_polynomialEmbed}\leanok
Let $\mathcal G_q=L^2(dA)\oplus^{\ell^2}L^2(dA)$ be the physical space of two
Cartesian gradient components. Polynomial differentiation defines
$\nabla p=(\partial_xp,\partial_yp)$ in this space. For each valid angular
sector $s$, the map
\[
 \iota p\longmapsto\kappa_s^{-1/2}\nabla U_p
\]
extends to a linear isometry $G_s:H_{t,m}\to\mathcal G_q$. Its two components
on polynomial inputs are exactly the corresponding partial-derivative $L^2$
classes, with this scale.
\end{theorem}
\begin{proof}
\leanok
\end{proof}

The next result controls the function and both Cartesian gradient components
simultaneously. Gradient error is exactly energy error. The following density
results identify this closed image with the zero-boundary Sobolev graph.

\begin{theorem}[Full Cartesian gradient realization and simultaneous convergence]
\label{thm:annular-full-gradient}\uses{thm:annular-sector-gradient,thm:annular-body-forcing}
\lean{MvPolynomial.AnnularEnergy.orthogonalFamily_sectorGradient, MvPolynomial.AnnularEnergy.gradientL2, MvPolynomial.AnnularEnergy.hasSum_gradientL2, MvPolynomial.AnnularEnergy.inner_gradientL2, MvPolynomial.AnnularEnergy.gradientL2_injective, MvPolynomial.AnnularEnergy.isClosed_range_gradientL2, MvPolynomial.AnnularEnergy.gradientL2_polynomialEmbed, MvPolynomial.AnnularEnergy.functionGradient, MvPolynomial.AnnularEnergy.isClosed_range_functionGradient, MvPolynomial.AnnularEnergy.gradient_body_solution_weak, MvPolynomial.AnnularEnergy.norm_gradientL2_error, MvPolynomial.AnnularEnergy.tendsto_functionGradient_finiteSolve}\leanok
The completed Cartesian gradient sectors are mutually orthogonal. Their
synthesis gives a linear isometry $G:\mathcal H_q\to\mathcal G_q$ with
unconditionally convergent sector series and closed range. On physically
scaled polynomial energy vectors it is exactly $\nabla U_p$.
The simultaneous map
\[
 x\longmapsto(Tx,Gx)\in L^2(dA)\times\mathcal G_q
\]
is an isometry for the product maximum norm and has closed range. In particular,
the physical body-force solution satisfies
\[
 \langle Gx,Gz\rangle_{\mathcal G_q}=\langle g,Tz\rangle_{L^2}
 \quad\text{for all }z\in\mathcal H_q.
\]
The Cartesian-gradient error of every actual finite solve equals its energy
error. The common angular/radial cutoff reconstructions converge simultaneously
as physical functions and Cartesian gradient classes.
\end{theorem}
\begin{proof}
\leanok
\end{proof}
\end{supporting}


\section{Weak derivatives and the physical Dirichlet problem}\label{sec:annular-sobolev}

\begin{supporting}[Reference details: The open annular domain and smooth tests and related results]
\begin{definition}[The open annular domain and smooth tests]
\label{def:annular-smooth-tests}\uses{thm:annular-area-measure,thm:annular-radial-L2}
\lean{MvPolynomial.AnnularEnergy.domain, MvPolynomial.AnnularEnergy.test_memLp, MvPolynomial.AnnularEnergy.physicalTestL2}\leanok
For $0<q<1$, the open annular domain is
$\Omega_q=\{(x,y):q<x^2+y^2<1\}$.
A test $\phi\in C_c^\infty(\Omega_q)$ determines its almost-everywhere
class in the physical $L^2(dA_q)$ space.
\end{definition}

\begin{theorem}[Weak Cartesian derivatives of annular energy vectors]
\label{thm:annular-weak-gradient}\uses{def:annular-smooth-tests,thm:annular-full-gradient,thm:polynomial-test-calculus}
\lean{MvPolynomial.AnnularEnergy.coe_physicalTestL2, MvPolynomial.AnnularEnergy.inner_polynomial_physicalTestL2, MvPolynomial.AnnularEnergy.polynomial_weakDerivative, MvPolynomial.AnnularEnergy.gradientL2_weakDerivative}\leanok
For every completed annular energy vector $x\in\mathcal H_q$, its physical
function $Tx$ and Cartesian gradient $Gx$ satisfy
\[
 \langle Tx,D_i\phi\rangle_{L^2(dA_q)}
 =-\langle (Gx)_i,\phi\rangle_{L^2(dA_q)}
 \quad(\phi\in C_c^\infty(\Omega_q),\ i=1,2).
\]
The physical test class agrees almost everywhere with $\phi$, and pairing it
with a polynomial is the ordinary integral of their product.
\end{theorem}
\begin{proof}
\leanok
\end{proof}


\begin{definition}[Smooth annular cutoff trials]
\label{def:annular-cutoff-trials}\uses{thm:smooth-boundary-factor,def:annular-smooth-tests,thm:polynomial-test-calculus}
\lean{MvPolynomial.AnnularEnergy.cutoffTrial, MvPolynomial.AnnularEnergy.cutoffTrial_contDiff, MvPolynomial.AnnularEnergy.cutoffTrial_support, MvPolynomial.AnnularEnergy.cutoffTest}\leanok
Let $0<q<1$, $r^2=x^2+y^2$, and let $p$ be a real Cartesian polynomial.
With the smooth boundary factor $F_a$ of Theorem~\ref{thm:smooth-boundary-factor}, set
\[
 u_a(x,y)=F_a(r^2-q)F_a(1-r^2)p(x,y),\qquad \phi_n=u_{n+1}.
\]

\end{definition}

\begin{theorem}[Smoothness and support of annular cutoff trials]
\label{thm:annular-cutoff-trials}\uses{def:annular-cutoff-trials}
\lean{MvPolynomial.AnnularEnergy.cutoffTrial_contDiff, MvPolynomial.AnnularEnergy.cutoffTrial_support, MvPolynomial.AnnularEnergy.cutoffTest_eventuallyEq}\leanok
For $0<q<1$ and a real Cartesian polynomial $p$, Every $u_a$ is smooth; for $a>0$ it has compact support contained in $\Omega_q$. At each point of $\Omega_q$, the functions $\phi_n$ eventually agree on a neighbourhood with $B_qp$, where $B_q=(1-r^2)(r^2-q)$.
\end{theorem}
\begin{proof}
\leanok
\end{proof}

\begin{theorem}[Smooth approximation of annular polynomial trials]
\label{thm:annular-cutoff-convergence}\uses{thm:annular-cutoff-trials,thm:l2-dominated-convergence}
\lean{MvPolynomial.AnnularEnergy.tendsto_cutoffTest_L2, MvPolynomial.AnnularEnergy.tendsto_cutoffTest_derivative_L2}\leanok
For every real Cartesian polynomial $p$, the compactly supported smooth
cutoff trials satisfy
\[
 \phi_n\longrightarrow B_qp,\qquad
 D_i\phi_n\longrightarrow\partial_i(B_qp)
 \quad\text{in }L^2(dA_q),\quad i=1,2.
\]
\end{theorem}
\begin{proof}
\leanok
\end{proof}


\begin{theorem}[Annular energy vectors have zero Sobolev boundary values]
\label{thm:annular-sobolev-inclusion}\uses{thm:annular-cutoff-convergence,thm:annular-full-gradient,thm:l2-graph-measure-congr}
\lean{MvPolynomial.AnnularEnergy.areaGraphEquiv, MvPolynomial.AnnularEnergy.physicalTestGradient, MvPolynomial.AnnularEnergy.areaGraphEquiv_testGraph, MvPolynomial.AnnularEnergy.tendsto_cutoffTest_gradient_L2, MvPolynomial.AnnularEnergy.polynomialGraph_mem_dirichlet, MvPolynomial.AnnularEnergy.functionGradient_mem_dirichlet}\leanok
The equality of polar and area measures induces an isometric identification
$A_q$ of their function/gradient product spaces. For a smooth test $\phi$,
its physical gradient has coordinates $[D_i\phi]$, and $A_q$ carries its
physical function/gradient pair to its ordinary area-measure graph pair.
The cutoff gradients converge in the finite Hilbert sum. Consequently,
\[
 A_q([B_qp],[\nabla(B_qp)])\in\mathcal H_0^1(\Omega_q)
\]
for every Cartesian polynomial $p$, and every completed energy vector satisfies
\[
 A_q(Tx,Gx)\in\mathcal H_0^1(\Omega_q)\qquad(x\in\mathcal H_q).
\]
\end{theorem}
\begin{proof}
\leanok
\end{proof}


\begin{theorem}[Cartesian polynomial trials in the completed energy graph]
\label{thm:annular-polynomial-graph}\uses{thm:annular-full-gradient}
\lean{MvPolynomial.AnnularEnergy.exists_functionGradient_polynomial}\leanok
For every real Cartesian polynomial $p$, there exists $x\in\mathcal H_q$
with
\[
 (Tx,Gx)=([B_qp],[\nabla(B_qp)]).
\]
Thus the completed energy image contains all boundary-weighted Cartesian
polynomial graph pairs.
\end{theorem}
\begin{proof}
\leanok
\end{proof}

\begin{theorem}[Boundary-preserving approximation of smooth annular tests]
\label{thm:annular-polynomial-density}\uses{thm:annular-sobolev-inclusion,thm:annular-polynomial-graph,thm:multivariate-C1-approximation}
\lean{MvPolynomial.AnnularEnergy.exists_smooth_boundary_quotient, MvPolynomial.AnnularEnergy.exists_boundary_polynomial_approximation, MvPolynomial.AnnularEnergy.physicalTestGraph_mem_range}\leanok
Every $\phi\in C_c^\infty(\Omega_q)$, viewed globally, has a smooth
factorisation $\phi=B_qg$ with $g\in C^\infty(\mathbb R^2)$.
There is a sequence of real Cartesian polynomials $p_n$ such that
\[
 ([B_qp_n],[\nabla(B_qp_n)])\longrightarrow([\phi],[\nabla\phi])
 \quad\text{in }L^2(dA_q)\times\bigoplus_{i=1}^2{}^{\ell^2}L^2(dA_q).
\]
In particular, the physical function/gradient pair of every such test belongs
to the image of $x\mapsto(Tx,Gx)$ on $\mathcal H_q$.
\end{theorem}
\begin{proof}
\leanok
\end{proof}
\end{supporting}


\begin{theorem}[Identification with the annular zero-boundary Sobolev space]
\label{thm:annular-sobolev-identification}\uses{thm:annular-polynomial-density}
\lean{MvPolynomial.AnnularEnergy.dirichletEmbedding, MvPolynomial.AnnularEnergy.dirichletEmbedding_surjective, MvPolynomial.AnnularEnergy.dirichletEquiv, MvPolynomial.AnnularEnergy.coe_dirichletEquiv, MvPolynomial.AnnularEnergy.inner_dirichletEquiv_gradient}\leanok
For $0<q<1$, the map
\[
 E_q:\mathcal H_q\longrightarrow\mathcal H_0^1(\Omega_q),\qquad
 E_qx=A_q(Tx,Gx),
\]
is a surjective linear isometry, where the Sobolev graph carries the function/
gradient product maximum norm. Subscripts $f$ and $\nabla$ denote the
function and gradient components. Moreover,
\[
 \langle(E_qx)_\nabla,(E_qy)_\nabla\rangle_{L^2(\Omega_q)}
 =\langle x,y\rangle_{\mathcal H_q}.
\]
Thus the completed polynomial energy is exactly the zero-boundary Sobolev
space, with its actual function and Cartesian gradient components.
\end{theorem}
\begin{proof}
\leanok
\end{proof}

\begin{theorem}[Weak Dirichlet solvability and spectral convergence]
\label{thm:annular-sobolev-solve}\uses{thm:annular-sobolev-identification}
\lean{MvPolynomial.AnnularEnergy.dirichletSolution, MvPolynomial.AnnularEnergy.dirichletSolution_weak, MvPolynomial.AnnularEnergy.existsUnique_dirichletSolution, MvPolynomial.AnnularEnergy.tendsto_dirichlet_finiteSolve}\leanok
Let $L$ be a continuous real linear functional on $\mathcal H_0^1(\Omega_q)$,
and put $\ell=L\circ E_q$. There is exactly one $u_L\in\mathcal H_0^1(\Omega_q)$
such that
\[
 \langle(u_L)_\nabla,v_\nabla\rangle_{L^2(\Omega_q)}=L(v)
 \qquad(v\in\mathcal H_0^1(\Omega_q)).
\]
It is the image under $E_q$ of the energy Riesz solution for $\ell$.
The finite annular solves with all valid angular frequencies $m<N$ and radial
degree below $N$, transported by $E_q$, converge to $u_L$ in the Sobolev graph
norm as $N\to\infty$.
\end{theorem}
\begin{proof}
\leanok
\end{proof}

\begin{theorem}[Ordinary area forcing and the annular Poisson problem]
\label{thm:annular-sobolev-body-forcing}\uses{thm:annular-sobolev-solve}
\lean{MvPolynomial.AnnularEnergy.dirichletBodyLoad, MvPolynomial.AnnularEnergy.dirichletBodyLoad_apply, MvPolynomial.AnnularEnergy.dirichletBodySolution_weak, MvPolynomial.AnnularEnergy.dirichletBodyLoad_transport}\leanok
For $g\in L^2(\Omega_q)$, define the continuous load
$L_g(v)=\langle g,v_f\rangle_{L^2(\Omega_q)}$. Its weak solution satisfies
\[
 \langle(u_{L_g})_\nabla,v_\nabla\rangle_{L^2(\Omega_q)}
 =\langle g,v_f\rangle_{L^2(\Omega_q)}
 \qquad(v\in\mathcal H_0^1(\Omega_q)).
\]
For a physical force $g_p\in L^2(dA_q)$, put
$g=(A_q(g_p,0))_f$. Then the pulled-back load is precisely
\[
 (L_g\circ E_q)(x)=\langle g_p,Tx\rangle_{L^2(dA_q)}.
\]
\end{theorem}
\begin{proof}
\leanok
\end{proof}

\begin{supporting}[Supporting results: the zero-boundary Sobolev graph]


\begin{theorem}[Uniform polynomial approximation on compact planar sets]
\label{thm:multivariate-uniform-approximation}
\lean{MvPolynomial.exists_polynomial_near_on_compact}\leanok
For every compact $K\subseteq\mathbb R^2$, every continuous real function
$f$ on $K$, and every $\varepsilon>0$, there is a real bivariate polynomial
$p$ such that $|p(x)-f(x)|<\varepsilon$ for every $x\in K$.
\end{theorem}
\begin{proof}
\leanok
\end{proof}

\begin{theorem}[Simultaneous polynomial approximation of values and first derivatives]
\label{thm:multivariate-C1-approximation}\uses{thm:multivariate-uniform-approximation,thm:polynomial-test-calculus,def:polynomial-linear-maps}
\lean{MvPolynomial.exists_polynomial_C1_near_on_square}\leanok
Let $f:\mathbb R^2\to\mathbb R$ be smooth. For every $\varepsilon>0$ there
is a real bivariate polynomial $p$ such that, for all $x\in[-1,1]^2$,
\[
 |p(x)-f(x)|<\varepsilon,\qquad
 |\partial_i p(x)-D_i f(x)|<\varepsilon\quad(i=1,2).
\]
\end{theorem}
\begin{proof}
\leanok
\end{proof}

\begin{theorem}[Smooth cutoff of a vanishing boundary factor]
\label{thm:smooth-boundary-factor}
\lean{Real.exists_transition_deriv_bound, Real.boundaryFactor, Real.boundaryFactor_contDiff, Real.hasDerivAt_boundaryFactor, Real.abs_boundaryFactor_le, Real.abs_deriv_boundaryFactor_le, Real.boundaryFactor_zero, Real.boundaryFactor_eq, Real.deriv_boundaryFactor_eq}\leanok
Let $\eta$ be the smooth transition from $0$ to $1$, with $0\leq\eta\leq1$,
$\eta(t)=0$ for $t\leq0$, and $\eta(t)=1$ for $t\geq1$.
There is $M\geq0$ such that $|\eta'(t)|\leq M$ and
$|t\eta'(t-1)|\leq2M$ for all real $t$. Define
\[
 F_a(r)=r\eta(ar-1).
\]
For every real $a$, this function is smooth and
\[
 F_a'(r)=\eta(ar-1)+ar\eta'(ar-1),\quad
 |F_a(r)|\leq|r|,\quad |F_a'(r)|\leq1+2M.
\]
It vanishes when $ar\leq1$, equals $r$ when $ar\geq2$, and has derivative
$1$ when $ar>2$.
\end{theorem}
\begin{proof}
\leanok
\end{proof}


\begin{theorem}[Transport of function/gradient pairs across equal measures]
\label{thm:l2-graph-measure-congr}\uses{thm:mathlib-hilbert-sums}
\lean{MeasureTheory.L2.graphMeasureCongr, MeasureTheory.L2.graphMeasureCongr_toLp, MeasureTheory.L2.graphMeasureCongr_inner_fst, MeasureTheory.L2.graphMeasureCongr_inner_snd}\leanok
If $\mu=\nu$ are measures on the same measurable space, then
\[
 L^2(\mu)\times\bigoplus_{i=1}^{n}{}^{\ell^2}L^2(\mu)
 \simeq L^2(\nu)\times\bigoplus_{i=1}^{n}{}^{\ell^2}L^2(\nu)
\]
isometrically and linearly. This identification preserves the classes of a
function and of any finite list of square-integrable component functions.
Scalar and gradient inner products are preserved separately.
\end{theorem}
\begin{proof}
\leanok
\end{proof}

\begin{definition}[The zero-boundary Sobolev graph]
\label{def:dirichlet-sobolev-graph}
\lean{MeasureTheory.DirichletSpace.FunctionL2, MeasureTheory.DirichletSpace.GradientL2, MeasureTheory.DirichletSpace.testL2, MeasureTheory.DirichletSpace.testDerivative, MeasureTheory.DirichletSpace.testGradient, MeasureTheory.DirichletSpace.testGraph, MeasureTheory.DirichletSpace.graph, MeasureTheory.DirichletSpace.Space}\leanok
For an open set $\Omega\subseteq\mathbb R^n$, let
$L=L^2(\Omega)$ and $\mathcal G=\bigoplus_{i=1}^n{}^{\ell^2}L$.
Write $D_i\phi$ for the Cartesian derivative of
$\phi\in C_c^\infty(\Omega)$, viewed globally with support compactly contained
in $\Omega$. The zero-boundary Sobolev graph is
\[
 \mathcal H_0^1(\Omega)=
 \overline{\{(\phi,(D_i\phi)_i):\phi\in C_c^\infty(\Omega)\}}
 \subseteq L\times\mathcal G.
\]
The closure uses the product maximum norm. Elements are represented by
function/gradient pairs.
\end{definition}

\begin{theorem}[Cartesian polynomial calculus against smooth tests]
\label{thm:polynomial-test-calculus}\uses{thm:dirichlet-graph-weak-derivatives}
\lean{MvPolynomial.contDiff_eval, MvPolynomial.fderiv_eval_single, MvPolynomial.integral_eval_mul_testDerivative}\leanok
Real multivariate polynomial evaluation is smooth. Its Fr\'echet derivative
in coordinate direction $e_i$ is the evaluation of its algebraic partial
derivative $\partial_i p$. For every $\phi\in C_c^\infty(\Omega)$,
\[
 \int_{\mathbb R^n} p\,D_i\phi
 =-\int_{\mathbb R^n}(\partial_i p)\,\phi.
\]
\end{theorem}
\begin{proof}
\leanok
\end{proof}

\begin{theorem}[Weak derivatives in the zero-boundary Sobolev graph]
\label{thm:dirichlet-graph-weak-derivatives}\uses{def:dirichlet-sobolev-graph}
\lean{MeasureTheory.DirichletSpace.testGradient_apply, MeasureTheory.DirichletSpace.testGraph_mem, MeasureTheory.DirichletSpace.graph_isClosed, MeasureTheory.DirichletSpace.coe_testL2, MeasureTheory.DirichletSpace.testDerivative_apply, MeasureTheory.DirichletSpace.inner_testL2, MeasureTheory.DirichletSpace.inner_testDerivative, MeasureTheory.DirichletSpace.graph_weakDerivative}\leanok
The graph is closed and contains every smooth compactly supported
function/gradient pair. Every $(u,g)\in\mathcal H_0^1(\Omega)$ satisfies
\[
 \langle u,D_i\phi\rangle_{L^2(\Omega)}
 =-\langle g_i,\phi\rangle_{L^2(\Omega)}
 \quad(\phi\in C_c^\infty(\Omega),\ 1\le i\le n).
\]
Thus its gradient coordinates are weak derivatives. On smooth test functions,
the coordinates agree almost everywhere with their classical derivatives,
and their $L^2$ pairings are the ordinary integrals of products.
\end{theorem}
\begin{proof}
\leanok
\end{proof}
\end{supporting}

\begin{supporting}[Supporting results: Hilbert-space tools and physical evaluation]
\begin{theorem}[Uniform polynomial Poincar\'e bound on annuli]
\label{thm:annular-poincare}\uses{thm:annular-polynomial-green}
\lean{MvPolynomial.ae_sq_radius_le_one_zernikeAnnulusMeasure, MvPolynomial.diskEuler_diskRadialLift, MvPolynomial.integral_diskRadialDivergence_diskRadialLift, MvPolynomial.integral_diskRadialDivergence_boundary_sq, MvPolynomial.integralInner_annulus_boundary_euler, MvPolynomial.integralInner_annulus_boundary_le_gradient, MvPolynomial.integralInner_annulus_boundary_le_neg_laplacian}\leanok
For $0<q<1$ and every Cartesian polynomial $p$, put
$u=B_qp$, where $B_q=(1-x^2-y^2)(x^2+y^2-q)$. Then
\[
 \int u^2\,dA\leq\int |\nabla u|^2\,dA
 =-\int u\Delta u\,dA.
\]
The constant is independent of degree, angular frequency, and $q$. Writing $Eu=x\partial_xu+y\partial_yu$, the auxiliary identity is
$\int uEu\,dA=-\int u^2\,dA$.
\end{theorem}
\begin{proof}
\leanok
\end{proof}

\rvnote{thm:annular-poincare}{The constant is independent of degree, angular frequency, and $q$; sharpness is not asserted.}


\begin{theorem}[Multivariate polynomial evaluation in $L^2$]
\label{thm:multivariate-polynomial-L2}\uses{def:multivariate-moments}
\lean{MvPolynomial.memLp_eval_of_hasFiniteMoments, MvPolynomial.polynomialToL2, MvPolynomial.coe_polynomialToL2, MvPolynomial.inner_polynomialToL2}\leanok
For finitely many variables and a measure with finite polynomial moments,
polynomial evaluation defines a linear map to $L^2$. Its representatives agree
almost everywhere with evaluation, and its inner product equals the polynomial
integral pairing.
\end{theorem}
\begin{proof}
\leanok
\end{proof}

\begin{theorem}[Orthogonal-family synthesis]
\label{thm:mathlib-orthogonal-synthesis}\uses{thm:mathlib-hilbert-sums}
\lean{OrthogonalFamily.linearIsometry, OrthogonalFamily.hasSum_linearIsometry, OrthogonalFamily.linearIsometry_apply_single}\mathlibok
A family of linear isometries into a complete inner-product space, with mutually
orthogonal images, induces a linear isometry from the corresponding Hilbert sum.
Its value is the unconditionally convergent sum of the coordinate images.
\end{theorem}
\begin{proof}
\mathlibok
\end{proof}

\begin{theorem}[Synthesis of orthogonal contractions]
\label{thm:orthogonal-contraction-synthesis}\uses{thm:mathlib-orthogonal-synthesis}
\lean{lp.orthogonalFamily_range, lp.rangeCoordinates, lp.orthogonalSynthesis, lp.hasSum_orthogonalSynthesis, lp.norm_orthogonalSynthesis_le, lp.orthogonalSynthesis_single, lp.orthogonalSynthesis_injective}\leanok
Let $T_s:E_s\to H$ be contractions from real inner-product spaces to a real
Hilbert space, with mutually orthogonal images. Then
\[
 T:\bigoplus_s^{\ell^2} E_s\longrightarrow H,\qquad
 Tx=\sum_s T_sx_s
\]
is a continuous linear contraction, and this series converges unconditionally.
It agrees with $T_s$ on each single-coordinate input. If each $T_s$ is injective,
then $T$ is injective. Closedness of the individual ranges is not required.
\end{theorem}
\begin{proof}
\leanok
\end{proof}

\begin{theorem}[Extending isometries to completions]
\label{thm:mathlib-isometry-completion}
\lean{LinearIsometry.fromCompletion, LinearIsometry.fromCompletion_apply_coe, Isometry.isClosedEmbedding}\mathlibok
A linear isometry into a complete normed space extends to the completion of its
domain. An isometry from a complete metric space has closed range.
\end{theorem}
\begin{proof}
\mathlibok
\end{proof}

\begin{theorem}[Hilbert sums]
\label{thm:mathlib-hilbert-sums}
\lean{lp.instInnerProductSpace, lp.inner_eq_tsum, lp.inner_single_right, lp.hasSum_single}\mathlibok
The $\ell^2$ sum of real Hilbert spaces has the sum inner product. Single-sector
injections preserve their inner products, and the finite sums of the coordinates
of any vector converge in the Hilbert-sum norm to that vector.
\end{theorem}
\begin{proof}
\mathlibok
\end{proof}

\begin{theorem}[Joint finite-section convergence in a Hilbert sum]
\label{thm:hilbert-sum-galerkin}\uses{thm:mathlib-hilbert-sums,thm:mathlib-hilbert-projection-limit}
\lean{lp.HilbertSumGalerkin.trialSpace, lp.HilbertSumGalerkin.trialSpace_mono, lp.HilbertSumGalerkin.sum_projection_eq, lp.HilbertSumGalerkin.trialSpace_closure_iSup, lp.HilbertSumGalerkin.tendsto_sum_projection}\leanok
In each real Hilbert sector $E_s$, let $W_{s,N}$ be increasing finite-dimensional
subspaces with dense union. For a finite sector set $S$ define
$V_{S,N}=\bigoplus_{s\in S}W_{s,N}$ inside $\bigoplus_s^{\ell^2}E_s$.
Its orthogonal projection is the sum of the retained sector projections.
As $(S,N)$ grows in the product order (set inclusion and integer order), these
projections converge strongly to the identity.
\end{theorem}
\begin{proof}
\leanok
\end{proof}

\end{supporting}

\begin{supporting}[Supporting results: Galerkin matrices, positive definiteness, Hilbert projections]
\begin{theorem}[Positive finite polynomial Galerkin matrices]
\label{thm:finite-polynomial-galerkin}\uses{thm:polynomial-inner}
\lean{Polynomial.integralInner_operator_sum, Polynomial.posDef_integralInner_operator, Polynomial.integralInner_operator_sum_apply}\leanok
Let $\mu$ have finite polynomial moments, let $L$ be a real linear polynomial
operator symmetric in the integral pairing, and suppose
$\langle Lp,p\rangle_\mu>0$ for every nonzero polynomial $p$.
For any finite linearly independent trial family $f_i$, the matrix
$M_{ij}=\langle Lf_j,f_i\rangle_\mu$ is positive definite. Its quadratic form is
exactly the energy of $\sum_i v_i f_i$, and its rows are the corresponding tested
operator values.
\end{theorem}
\begin{proof}
\leanok
\end{proof}

\begin{theorem}[Invertibility of finite positive definite matrices]
\label{thm:mathlib-positive-invertible}
\lean{Matrix.PosDef.isUnit, Matrix.mulVec_injective_iff_isUnit, Matrix.mulVec_surjective_iff_isUnit}\mathlibok
A finite real positive definite matrix is invertible. Multiplication by an
invertible square matrix is bijective on coefficient vectors.
\end{theorem}
\begin{proof}
\mathlibok
\end{proof}

\begin{theorem}[Hilbert representation and increasing projections]
\label{thm:mathlib-hilbert-projection-limit}
\lean{InnerProductSpace.toDual, InnerProductSpace.toDual_symm_apply, Submodule.starProjection_tendsto_self, Submodule.norm_starProjection_apply_le}\mathlibok
Every continuous real linear functional on a real Hilbert space has a unique
inner-product representative. Orthogonal projections onto increasing closed
subspaces with dense union converge strongly to the identity and are norm
non-increasing.
\end{theorem}
\begin{proof}
\mathlibok
\end{proof}

\end{supporting}

Zernike families on annuli: see Theorems~\ref{thm:zernike-annulus-basis}, \ref{thm:zernike-annulus-conversion} and~\ref{thm:zernike-annulus-laplacian}.

\chapter{Multivariate measure conversion}\label{chap:multivariate-modification}
\rvsummary{Gram factors, polynomial modification and degree-block bandwidth.}
In an ordered polynomial basis, measure conversion can be computed by factoring
a Gram matrix. Degree-block bandwidth describes which blocks vanish. Sparsity
within a block requires further structure, studied in the following chapters.

\begin{supporting}[Reference details: Flags and triangular conversion and related results]
\begin{theorem}[Flags and triangular conversion]
\label{thm:conversion-flags}\uses{def:mathlib-total-degree,def:polynomial-linear-maps}
\lean{Module.Basis.toMatrix_blockTriangular_of_mem_span, Module.Basis.toMatrix_isUpperTriangular_of_mem_span}\leanok
If each $P_j$ belongs to the span of $Q_i$ with degree label $d_i\le d_j$, then
$Q^{-1}P$ is block upper triangular with respect to $d$. A shared scalar flag
$P_j\in\operatorname{span}\{Q_i:i\le j\}$ gives scalar upper triangularity.
Total degree alone does not specify an ordering within each block.
\end{theorem}
\begin{proof}
\leanok
\end{proof}

\begin{theorem}[Exact leading sections of Gram factors]
\label{thm:gram-factor-sections}\uses{thm:positive-gram-factor}
\lean{Matrix.IsUpperTriangular.transpose_mul_self_principal, Matrix.eq_principal_of_pos_upperTriangular_gram_eq}\leanok
For an upper triangular matrix $R$ of size $n+m$, the leading size-$n$ section
of $R^\top R$ equals $R_n^\top R_n$. Positive-diagonal triangular factors of
Gram matrices agreeing on this section consequently agree there too.
\end{theorem}
\begin{proof}
\leanok
\end{proof}
\end{supporting}


\begin{theorem}[Finite Gram matrices and Cholesky measure conversion]
\label{thm:measure-conversion-cholesky}\uses{def:multivariate-degree-spaces,thm:positive-gram-factor}
\lean{MvPolynomial.integralGram, MvPolynomial.existsUnique_pos_upperTriangular_integralGram_factor, MvPolynomial.exists_orthonormal_pos_upperTriangular_basis, MvPolynomial.integralGram_eq_transpose_toMatrix_mul}\leanok
Let $\nu$ have finite polynomial moments and a positive-definite polynomial
pairing, in any finite number of variables. Every ordered basis $P$ of a finite
polynomial subspace has a unique positive-diagonal upper triangular Gram factor.
There exists an orthonormal basis $Q$ for which $P=QR$ with this factor.
For any orthonormal basis $Q$ of the same space, its coordinate matrix satisfies
$G_\nu(P)=R^\top R$; triangularity follows from the appropriate flag hypothesis.
The preceding leading-section result supplies compatibility whenever the ordered
basis vectors in the smaller section are unchanged.

Reference: \cite[Lemma 2.8]{gutleb-olver-slevinsky} and its references.
\end{theorem}
\begin{proof}
\leanok
\end{proof}

\begin{theorem}[Polynomial densities and degree-block bandwidth]
\label{thm:polynomial-density-gram}\uses{thm:measure-conversion-cholesky,def:multivariate-polynomial-density-measure}
\lean{MvPolynomial.HasFiniteMoments.withDensity_polynomial, MvPolynomial.HasPositiveDefiniteIntegral.withDensity_polynomial, MvPolynomial.integralInner_withDensity_polynomial, MvPolynomial.integralGram_withDensity_polynomial, MvPolynomial.integralInner_withDensity_eq_zero_of_degree}\leanok
Let $u$ be a polynomial nonnegative $\mu$-almost everywhere, and set $d\nu=u\,d\mu$.
Finite moments are preserved. If $u\ne0$ and the original polynomial pairing is
positive definite, so is the modified pairing. Moreover
$\langle p,q\rangle_\nu=\langle up,q\rangle_\mu$.
Suppose $\deg P_i\le d_i$ and $P_i$ is orthogonal to every polynomial of total
degree less than $d_i$. Then $\langle P_i,P_j\rangle_\nu=0$ whenever
$d_j+\deg u<d_i$ or $d_i+\deg u<d_j$.
This asserts degree-block bandwidth, without an assertion about sparsity inside blocks.
\end{theorem}
\begin{proof}
\leanok
\end{proof}

\begin{supporting}[Reference details: Exact conversion intertwines multiplication and related results]
\begin{theorem}[Exact conversion intertwines multiplication]
\label{thm:multivariate-intertwining}\uses{thm:operator-coefficients}
\lean{MvPolynomial.coefficientMap_mul_intertwining}\leanok
For any two bases of the multivariate polynomial ring, conversion $C$ and
multiplication by a polynomial $u$ satisfy $C M_P(u)=M_Q(u)C$ as linear maps
on finitely supported coefficients. Orthogonality and boundedness on $\ell^2$
are not required.
\end{theorem}
\begin{proof}
\leanok
\end{proof}

\begin{theorem}[Exact degree padding for truncated products]
\label{thm:polynomial-section-padding}\uses{def:mathlib-total-degree}
\lean{MvPolynomial.iterate_truncated_mul_eq_pow_mul}\leanok
Let a truncation $T$ fix every polynomial of total degree at most $N$.
If $\deg p+k\deg u\le N$, then $k$ applications of $q\mapsto T(uq)$ to $p$
give exactly $u^kp$. The bound is on degree, not coefficient count.

\end{theorem}
\begin{proof}
\leanok
\end{proof}

\rvnote{thm:polynomial-section-padding}{
There is no assertion that applying a polynomial to an unpadded multiplication section gives the corresponding section of the full polynomial operator.
}
\end{supporting}


\begin{theorem}[Degree-block bandwidth of Cholesky factors]
\label{thm:structured-measure-modification}\uses{thm:polynomial-density-gram}
\lean{Matrix.IsUpperTriangular.eq_zero_of_transpose_mul_self_degree_band, MvPolynomial.gram_factor_eq_zero_of_degree}\leanok
Let the scalar basis order have nondecreasing degree labels $d_i$.
An upper triangular factor $R$ with nonzero diagonal inherits upper
block bandwidth $d$ from $R^\top R$: if its Gram entries vanish whenever
$d_i+d<d_j$, then so do $R_{ij}$. In particular, the positive Cholesky factor
for a nonnegative polynomial modification $u$ of a degree-orthogonal family
vanishes when $d_i+\deg u<d_j$.
For one polynomial per degree this is ordinary bandwidth. In higher dimensions
it bounds degree blocks, without determining sparsity inside them.

Reference: \cite[Lemma 2.8 and Section 5]{gutleb-olver-slevinsky}.
\end{theorem}
\begin{proof}
\leanok
\end{proof}

\begin{theorem}[Separated tensor measure conversion]
\label{thm:tensor-product-modification}\uses{thm:polynomial-density-gram,thm:univariate-modification-support,thm:tensor-product-bases}
\lean{MvPolynomial.tensorProductPolynomial_mul, MvPolynomial.integralInner_tensorProductModification, MvPolynomial.integralInner_tensorProductModification_eq_zero_of_gap}\leanok
For $\sigma$-finite factor measures and a nonnegative separated polynomial density $\prod_i h_i(x_i)$, assume each weighted factor product $h_i p_i q_i$ belongs to $L^1(\mu_i)$. The modified pairing of the tensor products of $p_i$ and $q_i$ then factors into the pairings for the modified coordinate measures. If each factor measure has finite polynomial moments and the source and target coordinate sequences are orthogonal for their respective measures, the mixed pairing is zero whenever any coordinate lies outside $0\le n_i-m_i\le\deg h_i$.
\end{theorem}
\begin{proof}
\leanok
\end{proof}

\begin{supporting}[Reference details: Classical Proriol polynomial norm and related results]
\begin{definition}[Classical Proriol polynomial norm]
\label{def:proriol-norm}\uses{thm:proriol-triangle-measure,def:proriol-triangle}
\lean{MvPolynomial.proriolNorm}\leanok
For every parameter and index choice for which $(P_{n,k}^{\alpha,\beta,\gamma})^2$ is
integrable with respect to $\omega_{\alpha,\beta,\gamma}$, define
\[
 \|P_{n,k}^{\alpha,\beta,\gamma}\|_{\omega_{\alpha,\beta,\gamma}}=\sqrt{\int (P_{n,k}^{\alpha,\beta,\gamma})^2\,d\omega_{\alpha,\beta,\gamma}}.
\]
The domain consists of all such square-integrable polynomials; no additional
parameter restrictions are imposed.
\end{definition}
\begin{theorem}[Characterisation of the Proriol norm]
\label{thm:proriol-norm}\uses{def:proriol-norm}
\lean{MvPolynomial.proriolNorm_nonneg,MvPolynomial.proriolNorm_sq}\leanok
Whenever $(P_{n,k}^{\alpha,\beta,\gamma})^2$ is $\omega_{\alpha,\beta,\gamma}$-integrable,
\[
 \|P_{n,k}^{\alpha,\beta,\gamma}\|_{\omega_{\alpha,\beta,\gamma}}\ge0,\qquad
 \|P_{n,k}^{\alpha,\beta,\gamma}\|_{\omega_{\alpha,\beta,\gamma}}^2=\int (P_{n,k}^{\alpha,\beta,\gamma})^2\,d\omega_{\alpha,\beta,\gamma}.
\]
\end{theorem}
\begin{proof}
\leanok
\end{proof}
\end{supporting}


\begin{theorem}[Canonical Proriol Cholesky factors]
\label{thm:proriol-cholesky}\uses{thm:proriol-all-shift-support,thm:canonical-polynomial-gram-factor,thm:ordered-cartesian-even-sector,def:proriol-norm}
\lean{MvPolynomial.normalizedProriolDegreeBasis, MvPolynomial.orderedProriolBasis, MvPolynomial.integralGramFactor_proriol_eq_connection, MvPolynomial.integralGramFactor_proriol_eq_zero}\leanok
Normalize Proriol polynomials by their positive norms and order a complete finite
section by total degree, then inner degree. Under nonnegative integer parameter
raising, the actual normalized connection is the canonical positive upper
triangular Cholesky factor of the target-measure Gram matrix of the source basis.
Both preceding support bounds therefore hold for this canonical factor itself.
\end{theorem}
\begin{proof}
\leanok
\end{proof}

\begin{theorem}[Disk connection as Cholesky factor]
\label{thm:disk-gram-factor}\uses{thm:disk-normalization,thm:dunkl-xu-raising}
\lean{MvPolynomial.dunklXuRaisingMatrix, MvPolynomial.integralGram_eq_dunklXuRaisingMatrix_transpose_mul, MvPolynomial.dunklXuRaisingMatrix_eq_zero, MvPolynomial.dunklXuRaisingMatrix_diag_pos, MvPolynomial.dunklXuRaisingMatrix_eq_zero_of_order}\leanok
On every complete finite degree space, the connection from the normalized
$\beta$ basis to the normalized $\beta+1$ basis has positive diagonal, is upper
triangular in degree then inner-index order, and its transpose product is the
raised-measure Gram matrix. It has the four-position support of parameter
raising. Thus it is the positive Cholesky factor in that order.
\end{theorem}
\begin{proof}
\leanok
\end{proof}

\begin{theorem}[Outer conversion with an unchanged inner factor]
\label{thm:separated-outer-expansion}\uses{thm:univariate-connection-gram}
\lean{Polynomial.Sequence.IsOrthogonal.separated_expansion, Polynomial.Sequence.IsOrthogonal.separated_modification_expansion}\leanok
Let $P,Q$ be orthogonal sequences before and after a nonzero nonnegative
polynomial modification $u$, under the finite-moment and positivity hypotheses.
For any real commutative algebra $A$, algebra homomorphism
$\phi:\mathbb R[x]\to A$, and fixed $v\in A$, one has
\[
 \phi(P_n)v=\sum_{\substack{0\le m\le n\\n\le m+\deg u}}
 [Q_m]P_n\;\phi(Q_m)v.
\]
Applied separately at each inner index, this keeps the inner lift unchanged
and bounds the outer conversion bandwidth independently of the lift's degree.

\end{theorem}
\begin{proof}
\leanok
\end{proof}

\rvnote{thm:separated-outer-expansion}{
This algebraic expansion does not itself assert that a proposed collection of lifts is a basis of a total-degree space.
}

\chapter{Separated Koornwinder systems}\label{chap:koornwinder}
\rvsummary{Separated bases, family-preserving modifications and their connection support.}
Separation of variables reduces parts of the conversion problem to
one-dimensional polynomial families. The results specify the scale function
and basis ordering on which the connection support depends.

\begin{supporting}[Reference details: Vertical coefficient truncation and related results]
\begin{definition}[Vertical coefficient truncation]
\label{def:vertical-truncation}
\lean{MvPolynomial.verticalTruncation}\leanok
For $p\in R[X]$ over a commutative ring and $n\in\N$, define
\[
 V_n(p)(x,y)=\sum_{j=0}^n p_j y^j.
\]

\end{definition}

\rvnote{def:vertical-truncation}{
This construction has no scaling parameter.
}

\begin{theorem}[Evaluation and degree of vertical truncation]
\label{thm:vertical-truncation}\uses{def:vertical-truncation}
\lean{MvPolynomial.eval_verticalTruncation, MvPolynomial.eval_verticalTruncation_homogeneous, MvPolynomial.totalDegree_verticalTruncation_le}\leanok
The total degree of $V_n(p)$ is at most $n$. If $\deg p\le n$, then
$V_n(p)(x,y)=p(y)$. If in addition every nonzero coefficient $p_j$ with
$j\le n$ has $j=n$, then
$V_n(p)(x,\rho z)=\rho^n p(z)$ for every $\rho,z$, including $\rho=0$.
\end{theorem}
\begin{proof}
\leanok
\end{proof}

\begin{definition}[Polynomial lifts for separated scales]
\label{def:scaled-koornwinder-lift}
\lean{MvPolynomial.scaledLift}\leanok
For $b\in R[x]$, a step $s\ge1$ and $n\in\N$, define
\[
 L_{s,n}(q)=\sum_{j=0}^n q_j y^j b(x)^{\lfloor(n-j)/s\rfloor}.
\]
\end{definition}

\begin{theorem}[Evaluation, degree and the disk case]
\label{thm:scaled-lift-basic}\uses{def:scaled-koornwinder-lift,thm:disk-polynomial-lift}
\lean{MvPolynomial.eval_scaledLift, MvPolynomial.totalDegree_scaledLift_le, MvPolynomial.scaledLift_disk}\leanok
If $\deg q\le n$, every nonzero $q_j$ has $s\mid n-j$, $b(x)=\rho^s$ and
$y=\rho z$, then $L_{s,n}(q)(x,y)=\rho^nq(z)$; no condition $\rho\ne0$ is
needed. If $\deg b\le s$, then $L_{s,n}(q)$ has total degree at most $n$. For
$s=2$ and $b=1-x^2$ it is the disk lift of
Definition~\ref{def:disk-polynomial-lift}.
\end{theorem}
\begin{proof}
\leanok
\end{proof}

\begin{definition}[Degree-dependent outer measure]
\label{def:koornwinder-outer-measure}\uses{def:polynomial-density-measure}
\lean{Polynomial.koornwinderOuterMeasure}\leanok
For a real polynomial $a$ and $k\in\mathbb N$, let
$D_{a,k}=\{x\in\mathbb R:a(x)^k\ge0\}$.
For a source measure $\eta$ on this domain, the $k$th outer measure is the
polynomial density measure $a^k\,d\eta$, transported to real coordinates.
For even $k$, the source domain is all of $\mathbb R$.
\end{definition}

\begin{theorem}[Degree-dependent outer measure characterization]
\label{thm:koornwinder-outer-measure}\uses{def:koornwinder-outer-measure,thm:polynomial-density-measure}
\lean{Polynomial.koornwinderOuterMeasure_apply, Polynomial.koornwinderOuterMeasure_zero}\leanok
For a measurable set $E$, the $k$th outer measure assigns
$\int_{\iota^{-1}(E)}a(x)^k\,d\eta(x)$ as a nonnegative integral.
At degree zero, transporting any real source measure to $D_{a,0}=\mathbb R$
and applying this construction recovers that measure.
\end{theorem}
\begin{proof}
\leanok
\end{proof}

\begin{definition}[Separated pushforward measures]
\label{def:koornwinder-measure}\uses{def:product-measure-input}
\lean{MvPolynomial.koornwinderCoordinates, MvPolynomial.KoornwinderMeasureInput, MvPolynomial.koornwinderMeasure}\leanok
Let $D\subseteq\mathbb R$ be an outer domain, $\mu$ a measure on $D$,
$\nu$ a measure on $\mathbb R$, $\rho:D\to\mathbb R$ a scale and $a$ a
polynomial. Write $\bar\mu=\iota_*\mu$, where $\iota:D\hookrightarrow\mathbb R$
is the real coordinate map. When $a\ge0$ almost everywhere for $\bar\mu$,
transport $\bar\mu$ to $D_{a,k}$ and let $\mu_k$ be its degree-dependent
outer measure from Definition~\ref{def:koornwinder-outer-measure}. Let $\omega$ be
the pushforward of $\mu\otimes\nu$ by $(x,z)\mapsto(\iota(x),\rho(x)z)$.
The input includes almost-everywhere measurability certificates for both
the product kernel $x\mapsto(z\mapsto(x,z))_*\nu$ and the coordinate map
$(x,z)\mapsto(x,\rho(x)z)$ under $\mu\otimes\nu$.
The scale may be discontinuous, and no interval support condition is imposed.
Global scales use $D=\mathbb R$. A real measure supported almost everywhere
in a measurable $D$ is represented by its restriction to that domain; its
coordinate pushforward recovers the original measure.

\end{definition}

\rvnote{def:koornwinder-measure}{
Theorem~\ref{thm:koornwinder-moments} assumes $\rho$ continuous with $\rho^2=a$ almost everywhere for $\mu$.
}

\begin{theorem}[Characterisation of separated pushforward measures]
\label{thm:koornwinder-measure-characterisation}\uses{def:koornwinder-measure,thm:integrable-input,thm:polynomial-inner}
\lean{MvPolynomial.KoornwinderMeasureInput.ofContinuous,MvPolynomial.KoornwinderMeasureInput.ofMoments,MvPolynomial.koornwinderMeasure_apply,MvPolynomial.koornwinderMeasure_integral}\leanok
Write $\Phi(x,z)=(\iota(x),\rho(x)z)$ for the separated coordinate map
and $\omega=\Phi_*(\mu\otimes\nu)$ for the measure of
Definition~\ref{def:koornwinder-measure}.
An s-finite inner measure makes the product kernel almost-everywhere measurable;
if $\rho$ is continuous, $\Phi$ is also almost-everywhere measurable.
In particular, finite polynomial moments of the inner measure suffice.
For every measurable set $A$,
$\omega(A)=(\mu\otimes\nu)(\Phi^{-1}(A))$.
If $f$ is $\omega$-integrable, then $f\circ\Phi$ is $\mu\otimes\nu$-integrable and
\[
 \int f\,d\omega=\int f\circ\Phi\,d(\mu\otimes\nu).
\]
\end{theorem}
\begin{proof}
\leanok
\end{proof}

\begin{theorem}[Finite moments for separated measures]
\label{thm:koornwinder-moments}\uses{thm:koornwinder-measure-characterisation,def:multivariate-moments}
\lean{MvPolynomial.hasFiniteMoments_koornwinderMeasure}\leanok
If $\bar\mu$ and $\nu$ have finite polynomial moments, so does $\omega$.
The scale is continuous and its square agrees $\mu$-almost everywhere with
the polynomial $a$.
\end{theorem}
\begin{proof}
\leanok
\end{proof}
\end{supporting}


\begin{definition}[Generalised Koornwinder systems]
\label{def:koornwinder-system}\uses{def:koornwinder-measure,thm:univariate-modification-support,thm:scaled-lift-basic,def:koornwinder-outer-measure}
\lean{MvPolynomial.KoornwinderSystem, MvPolynomial.koornwinderSystemOfLift, MvPolynomial.polynomialScaleKoornwinderSystem, MvPolynomial.sqrtScaleKoornwinderSystem, MvPolynomial.KoornwinderSystem.polynomial}\leanok
Require $a(\iota(x))\ge0$ almost everywhere for $\mu$.
Choose positive-leading orthonormal sequences $q_k$ for $\nu$ and
$p_r^{[k]}$ for $\mu_k$, together with polynomial lifts $L_k$ satisfying
$L_k(x,\rho(x)z)=\rho(x)^kq_k(z)$ almost everywhere in $x$.
Set $P_{n,k}=p_{n-k}^{[k]}L_k$ for $0\le k\le n$.
Two constructions of such systems are defined, under finite moments and
positive definite polynomial pairings for $\bar\mu$ and $\nu$: one for a
nonzero polynomial scale $c$ restricted to $D$ ($a=c^2$), and one for
$\rho(x)=\sqrt{a(\iota(x))}$ on $D=\{x\in\mathbb R:a(x)\ge0\}$ when
$a\ne0$ and $\nu$ is even. The square-root scale has no values outside $D$.
The label $n$ is the separated index.
\end{definition}

\rvnote{def:koornwinder-system}{The separated label $n$ is not asserted to equal total degree in these general cases.}


\begin{supporting}[Reference details: Separated evaluation of a Koornwinder system and related results]
\begin{theorem}[Separated evaluation of a Koornwinder system]
\label{thm:koornwinder-system-characterisation}\uses{def:koornwinder-system}
\lean{MvPolynomial.KoornwinderSystem.eval_polynomial}\leanok
For every pair of natural indices $n,k$, almost every $x\in D$ and every
$z\in\mathbb R$,
\[
P_{n,k}(\iota(x),\rho(x)z)=p_{n-k}^{[k]}(\iota(x))\rho(x)^kq_k(z),
\]
where $n-k$ is natural-number subtraction. In particular this is the
separated expression for every valid index $0\le k\le n$.
\end{theorem}
\begin{proof}
\leanok
\end{proof}

\begin{theorem}[Moments supplied by a separated system]
\label{thm:koornwinder-system-moments}\uses{def:koornwinder-system,thm:koornwinder-moments,thm:orthogonal-moments,thm:koornwinder-outer-measure}
\lean{MvPolynomial.KoornwinderSystem.hasFiniteMoments_inner, MvPolynomial.KoornwinderSystem.hasFiniteMoments_outer, MvPolynomial.KoornwinderSystem.hasFiniteMoments_shifted_outer, MvPolynomial.KoornwinderSystem.hasFiniteMoments_polynomial_scale, MvPolynomial.KoornwinderSystem.hasFiniteMoments}\leanok
A generalised Koornwinder system supplies all moments of its inner measure
and each shifted outer measure. If its scale is continuous and $\rho^2=a$
almost everywhere, the separated measure also has all polynomial moments.
In particular this holds for a polynomial scale $\rho=c$ with $a=c^2$.
\end{theorem}
\begin{proof}
\leanok
\end{proof}
\end{supporting}


\begin{theorem}[Orthonormality of separated systems]
\label{thm:koornwinder-orthonormality}\uses{thm:koornwinder-system-characterisation,def:multivariate-integrable-pair,thm:koornwinder-measure-characterisation}
\lean{MvPolynomial.KoornwinderSystem.integralInner_polynomial, MvPolynomial.KoornwinderSystem.integralInner_polynomial_eq_zero, MvPolynomial.KoornwinderSystem.integralInner_polynomial_self}\leanok
For a Koornwinder system of Definition~\ref{def:koornwinder-system}, with continuous scale $\rho$ satisfying $\rho^2=a$ almost everywhere for the outer measure,
$\langle P_{n,k},P_{m,j}\rangle_\omega=\delta_{nm}\delta_{kj}$
for valid separated indices.
\end{theorem}
\begin{proof}
\leanok
\end{proof}

\begin{supporting}[Reference details: Bases of separated finite spans]
\begin{theorem}[Bases of separated finite spans]
\label{thm:koornwinder-separated-spans}\uses{thm:koornwinder-orthonormality,thm:measure-conversion-cholesky,thm:koornwinder-system-moments}
\lean{MvPolynomial.KoornwinderIndex, MvPolynomial.KoornwinderSystem.sectionSpan, MvPolynomial.KoornwinderSystem.sectionBasis, MvPolynomial.KoornwinderSystem.integralGram_sectionBasis, MvPolynomial.KoornwinderSystem.sectionBasis_repr}\leanok
Let $V_N=\operatorname{span}\{P_{n,k}:0\le k\le n\le N\}$.
These indexed polynomials form an orthonormal basis of $V_N$; the coordinate
of $p\in V_N$ at $(n,k)$ is $\langle p,P_{n,k}\rangle_\omega$.

\end{theorem}
\begin{proof}
\leanok
\end{proof}

\rvnote{thm:koornwinder-separated-spans}{
This statement does not identify $V_N$ with a total-degree space.
}
\end{supporting}


\begin{theorem}[Ordinary systems span total-degree spaces]
\label{thm:koornwinder-ordinary-spans}\uses{thm:koornwinder-separated-spans,thm:mathlib-total-degree-finite}
\lean{MvPolynomial.separatedIndexEquiv, MvPolynomial.KoornwinderSystem.sectionSpan_eq_restrictTotalDegree, MvPolynomial.totalDegree_polynomialScaleKoornwinderSystem_lift, MvPolynomial.totalDegree_sqrtScaleKoornwinderSystem_lift}\leanok
If every $L_k$ has total degree at most $k$, then $V_N$ is exactly the
space of bivariate polynomials of total degree at most $N$.
The lift bound holds for affine polynomial scales and for square-root
quadratic scales with even inner measure.
\end{theorem}
\begin{proof}
\leanok
\end{proof}

\begin{supporting}[Formal details: Polynomial density within an outer domain and related results]
\begin{definition}[Polynomial density within an outer domain]
\label{def:subtype-polynomial-density-measure}
\lean{Polynomial.subtypeDensityMeasure}\leanok
Let $D\subseteq\mathbb R$ and $p\in\mathbb R[X]$. A measure $\eta$ on
$\{x\in D:p(x)\ge0\}$ defines a measure on $D$ by weighting by $p$ and
pushing forward along the inclusion into $D$.
\end{definition}

\begin{theorem}[Outer-domain polynomial density characterization]
\label{thm:subtype-polynomial-density-measure}\uses{def:subtype-polynomial-density-measure}
\lean{Polynomial.subtypeDensityMeasure_apply}\leanok
For every measurable subset $E$ of $D$, the resulting measure assigns the
nonnegative integral $\int_{\iota^{-1}(E)}p(x)\,d\eta(x)$.
\end{theorem}
\begin{proof}
\leanok
\end{proof}
\end{supporting}


\begin{theorem}[Type 1 outer modifications]
\label{thm:koornwinder-type1}\uses{thm:koornwinder-separated-spans,thm:separated-outer-expansion,def:subtype-polynomial-density-measure}
\lean{MvPolynomial.KoornwinderSystem.outerModification, MvPolynomial.KoornwinderSystem.sectionSpan_eq_of_lift_eq, MvPolynomial.KoornwinderSystem.connectionEntry_eq_zero_of_outer_modification}\leanok
Under finite moments and positive definite pairings, and for a nonzero scale
polynomial $a\ge0$ almost everywhere, modify $\mu$ by a nonzero polynomial
$u\ge0$ almost everywhere, using the density measure of Definition~\ref{def:subtype-polynomial-density-measure} after transporting the source to $\{u\ge0\}$, and keeping the inner sequence and lifts fixed. The old
and new separated spans $V_N$ coincide. If moreover $\nu$ has finite moments,
$\rho$ is continuous and $\rho^2=a$ almost everywhere for $\mu$, the connection
coefficient from $(n,k)$ to $(m,j)$ vanishes unless $j=k$ and
$0\le n-m\le\deg u$. No degree restriction on the scale is needed.
\end{theorem}
\begin{proof}
\leanok
\end{proof}

\begin{supporting}[Reference details: Separated modification polynomial and related results]
\begin{definition}[Separated modification polynomial]
\label{def:koornwinder-modification-polynomial}\uses{def:scaled-koornwinder-lift}
\lean{MvPolynomial.koornwinderModificationPolynomial}\leanok
For real polynomials $b,h,g$, a positive integer $s$ and $\ell\in\N$, define
\[
 u(x,y)=h(x)L_{s,s\ell}(g)(x,y).
\]
Reference: \cite[\S2.6.1]{dunkl-xu}.
\end{definition}

\begin{theorem}[Zero-step product as coefficient truncation]
\label{thm:koornwinder-zero-step}\uses{thm:vertical-truncation,thm:disk-polynomial-lift}
\lean{MvPolynomial.eval_zeroStepKoornwinderModification, MvPolynomial.totalDegree_zeroStepKoornwinderModification_le}\leanok
For real polynomials $b,h,g$, let $u(x,y)=h(x)V_0(g)(x,y)=h(x)g_0$.
Its total degree is at most $\deg h$. If $\deg g\le0$ and $b(x)=1$, then
$u(x,y)=b(x)^\ell h(x)g(z)$ for every $y,z$ and $\ell\in\N$.
\end{theorem}
\begin{proof}
\leanok
\end{proof}

\begin{theorem}[Physical densities for separated modifications]
\label{thm:koornwinder-separated-density}\uses{thm:koornwinder-measure-characterisation,thm:scaled-lift-basic,def:koornwinder-modification-polynomial,def:subtype-polynomial-density-measure,def:polynomial-density-measure}
\lean{MvPolynomial.eval_koornwinderModificationPolynomial, MvPolynomial.integral_eval_koornwinder_modification, MvPolynomial.koornwinderModificationPolynomial_affine, MvPolynomial.koornwinderModificationPolynomial_quadratic}\leanok
For a positive integer $s$, suppose $\rho^s=b$ and $\deg g\le s\ell$, with nonzero coefficients of
$g$ only at indices congruent to $s\ell$ modulo $s$.
The polynomial $u=h(x)L_{s,s\ell}(g)$ pulls back to
$b(x)^\ell h(x)g(z)$. If $\rho$ is continuous, $\nu$ is s-finite and both
factor densities are nonnegative almost everywhere, replacing $(\mu,\nu)$ by $(b^\ell h\mu,g\nu)$ gives the same
polynomial integrals as multiplying $\omega$ by $u$: for a polynomial $p$ with
$up\in L^1(\omega)$, the modified integral of $p$ exists and equals
$\int up\,d\omega$.
For $\ell=1$, $h=1$, the affine factor is $\alpha c(x)+\beta y$;
the square-root factor is $\alpha a(x)+\beta y^2$.
\end{theorem}
\begin{proof}
\leanok
\end{proof}
\end{supporting}


\begin{theorem}[Type 2 connection support]
\label{thm:koornwinder-type2-support}\uses{thm:koornwinder-system-moments,thm:koornwinder-system-characterisation,def:multivariate-integrable-pair,def:subtype-polynomial-density-measure,thm:multivariate-moments-basic}
\lean{MvPolynomial.KoornwinderSystem.connectionEntry_affine_support, MvPolynomial.KoornwinderSystem.connectionEntry_quadratic_support}\leanok
Assume finite moments and positive-leading normalised source and target systems.
For an affine scale $c\ge0$, modify the factor measures by $c^\ell h$ and
$g$, where $h,g\ge0$ almost everywhere and $\deg g\le\ell$.
A nonzero connection coefficient requires
\[
 (n-m,k-j)=(s,r),\qquad 0\le s\le\deg h+\ell,\quad 0\le r\le\ell.
\]
For a continuous scale with square $a$, $\deg a\le2$, assume additionally
that $\nu$ is even with positive definite polynomial pairing, and $g$ is
nonzero, even, nonnegative and of degree at most $2\ell$.
Modifying by $a^\ell h$ and $g$ permits only
\[
 (n-m,k-j)=(s,2r),\qquad 0\le s\le\deg h+2\ell,\quad 0\le r\le\ell.
\]
These are containing supports: individual permitted entries may also vanish.
\end{theorem}
\begin{proof}
\leanok
\end{proof}

\begin{theorem}[Separated conversion as a Cholesky factor]
\label{thm:koornwinder-conversion-factors}\uses{thm:koornwinder-ordinary-spans,thm:cholesky-lu}
\lean{MvPolynomial.KoornwinderSystem.connectionEntry_eq_sectionBasis_repr, MvPolynomial.KoornwinderSystem.connectionEntry_gram, MvPolynomial.KoornwinderSystem.connectionEntry_gram_of_ordinary, MvPolynomial.KoornwinderSystem.connectionEntry_diag_pos, MvPolynomial.KoornwinderSystem.connectionEntry_upperTriangular_of_support, MvPolynomial.KoornwinderSystem.connectionEntry_eq_of_gram_factor, MvPolynomial.KoornwinderSystem.connectionEntry_lu_support}\leanok
Suppose the source section lies in the target section and the target
factor measures are absolutely continuous with respect to the source ones.
Mixed pairings are the actual conversion coordinates $C$, and the source
Gram matrix in the target measure is $C^{\mathsf T}C$.
Section compatibility holds for ordinary systems, and for Type 1 modifications.
Using a degree-then-inner ordered enumeration, the preceding support bounds
make $C$ upper triangular; its diagonal is positive. It is the unique positive
upper triangular Gram factor. Any exact unpivoted unit-lower LU factorisation
of that Gram matrix has upper factor with exactly the same zero entries as $C$.
\end{theorem}
\begin{proof}
\leanok
\end{proof}

\begin{theorem}[Composition of separated conversions]
\label{thm:koornwinder-composition}\uses{thm:koornwinder-conversion-factors}
\lean{MvPolynomial.KoornwinderSystem.connectionEntry_comp}\leanok
For successive systems $S,T,U$, if the source section lies in the
intermediate section, mixed-pairing conversion matrices satisfy
$C_{SU}=C_{TU}C_{ST}$. Under compatible finite spans these are actual
coordinate conversions. Their containing offset supports add by Minkowski sum.

\end{theorem}
\begin{proof}
\leanok
\end{proof}

\rvnote{thm:koornwinder-composition}{
No additional sparsity or construction-cost assertion is implied.
}

\begin{supporting}[Reference details: Coordinate multiplication support]
\begin{theorem}[Coordinate multiplication support]
\label{thm:koornwinder-multiplication-support}\uses{thm:koornwinder-system-moments,def:multivariate-integrable-pair,thm:koornwinder-system-characterisation,thm:multivariate-moments-basic,thm:general-tridiagonal}
\lean{MvPolynomial.KoornwinderSystem.multiplicationEntry, MvPolynomial.KoornwinderSystem.multiplicationEntry_x_eq_zero, MvPolynomial.KoornwinderSystem.multiplicationEntry_y_polynomial_support, MvPolynomial.KoornwinderSystem.multiplicationEntry_y_even_support}\leanok
Let multiplication entries be $\langle x_iP_{n,k},P_{m,j}\rangle_\omega$.
For continuous $\rho$ with $\rho^2=a$ almost everywhere for $\mu$,
multiplication by $x$ requires $j=k$ and $|n-m|\le1$.
For polynomial scale $c$, multiplication by $y$ requires
$|k-j|\le1$, $|n-m|\le B$ whenever $B\ge1$, $\deg c\le B$ and
$\deg(c^2)\le B+1$. For an even inner measure with positive definite pairing,
the same bounds hold for continuous scale $\rho$ with $\rho^2=a$ and
$\deg a\le B+1$. Thus ordinary affine and square-root-quadratic cases have
$B=1$. Finite moments are assumed throughout.
For unrestricted scales these statements concern pairing entries, without
asserting completeness in the entire polynomial algebra.
\end{theorem}
\begin{proof}
\leanok
\end{proof}
\end{supporting}


\begin{theorem}[Exact coordinate action for ordinary systems]
\label{thm:koornwinder-operators}\uses{thm:koornwinder-ordinary-spans,thm:koornwinder-multiplication-support}
\lean{MvPolynomial.KoornwinderSystem.coordinate_mul_mem_sectionSpan, MvPolynomial.KoornwinderSystem.coordinate_mul_eq_sum}\leanok
If $\deg L_k\le k$, multiplying a member of the section $V_N$ by either
coordinate lies in $V_{N+1}$ and has the exact finite expansion with the
multiplication entries above as coefficients. Consequently the support bounds
describe the actual coordinate multiplication operators for ordinary systems.
\end{theorem}
\begin{proof}
\leanok
\end{proof}

\begin{theorem}[Compatibility with the Cartesian disk construction]
\label{thm:disk-koornwinder-bridge}\uses{def:cartesian-disk-system,def:koornwinder-system}
\lean{MvPolynomial.CartesianDiskSystem.toKoornwinderSystem, MvPolynomial.CartesianDiskSystem.toKoornwinderSystem_polynomial}\leanok
Every existing Cartesian disk system whose outer measure $\mu$ is supported
in $[-1,1]$ almost everywhere supplies a separated Koornwinder system with
$a=1-x^2$ and $\rho=\sqrt{1-x^2}$. Its polynomials are exactly
the existing disk polynomials, with the same indices and normalisation.
\end{theorem}
\begin{proof}
\leanok
\end{proof}

\chapter{Sparsity and obstructions on the disk}\label{chap:disk-sparsity}
\rvsummary{Exact zeros, parity sectors and obstructions to uniform inner bandwidth.}
The results below give both sparsity patterns and entries that cannot vanish.
The classification depends on the polynomial family and on whether a bandwidth
bound must hold uniformly; sparsity for one example alone does not establish
such a bound.

\begin{supporting}[Reference details: Canonical finite Gram factor and support and related results]
\begin{theorem}[Canonical finite Gram factor and support]
\label{thm:canonical-polynomial-gram-factor}\uses{thm:conversion-composition-support,thm:measure-conversion-cholesky}
\lean{MvPolynomial.integralGramFactor, MvPolynomial.integralGramFactor_spec, MvPolynomial.integralGramFactor_eq, MvPolynomial.integralGramFactor_hasOffsetSupport_iff}\leanok
For any ordered basis of a finite polynomial subspace and any measure with
finite moments and positive definite polynomial pairing, choose the unique
positive upper triangular Gram factor. Its offset support can be checked using
any positive upper triangular factor of the same Gram matrix. No separated
target orthogonal system is assumed.
\end{theorem}
\begin{proof}
\leanok
\end{proof}

\begin{theorem}[Algebraic product-separation test]
\label{thm:product-separation-test}
\lean{Function.IsProductSeparated, Function.isProductSeparated_iff}\leanok
A function $F:X\times Z\to K$ over a field is of the form $a(x)b(z)$
if and only if every rectangle identity
$F(x,z)F(x',z')=F(x,z')F(x',z)$ holds.
\end{theorem}
\begin{proof}
\leanok
\end{proof}

\rvnote{thm:product-separation-test}{
This characterizes product separation; it does not assert a converse from sparse Gram factors.
}
\end{supporting}


\begin{theorem}[Quadratic disk density separation]
\label{thm:quadratic-disk-separation}\uses{thm:product-separation-test}
\lean{Polynomial.quadraticDiskPullback_isProductSeparated_iff, Polynomial.quadraticDiskPullback_rectangle_minor, Polynomial.not_isProductSeparated_one_add_disk_y_sq, Polynomial.quadraticDisk_nonneg_iff}\leanok
The pullback $A+Bx^2+C(1-x^2)z^2$ separates as a product of functions of
$x$ and $z$ exactly when $C=0$ or $A+B=0$.
The physical density $A+Bx^2+Cy^2$ is nonnegative on the unit disk exactly
when $A\ge0$, $A+B\ge0$ and $A+C\ge0$.
In particular, positivity does not imply product separation: $1+y^2$ is an example.
\end{theorem}
\begin{proof}
\leanok
\end{proof}

\begin{supporting}[Reference details: Elimination zeros force vanishing minors and related results]
\begin{theorem}[Elimination zeros force vanishing minors]
\label{thm:elimination-minor-obstruction}
\lean{Matrix.lowerTriangular_mul_submatrix, Matrix.det_submatrix_eq_zero_of_lower_factorization}\leanok
Let $W=LU$ with $L$ lower triangular, and select the first $n$ rows and any
$n$ columns. If some row of the selected upper-factor submatrix vanishes,
the corresponding determinant of $W$ is zero. Thus a nonzero bordered minor
obstructs proposed upper-factor zeros. This is a finite necessary condition,
not an all-degree sparsity classification.
\end{theorem}
\begin{proof}
\leanok
\end{proof}

\begin{theorem}[Double Gegenbauer parameter conversion]
\label{thm:gegenbauer-double-conversion}\uses{thm:gegenbauer-conversion}
\lean{Polynomial.gegenbauer_double_conversion, Polynomial.gegenbauer_double_conversion_two, Polynomial.gegenbauer_double_conversion_three}\leanok
Put $v=n+\lambda+4$. With denominators cleared, two parameter raisings give
\[
 v(v+1)(v-1)C_{n+4}^{(\lambda)}
 =\lambda(\lambda+1)(v-1)C_{n+4}^{(\lambda+2)}
 -2\lambda(\lambda+1)v C_{n+2}^{(\lambda+2)}
 +\lambda(\lambda+1)(v+1)C_n^{(\lambda+2)}.
\]
The identities hold over every commutative rational algebra. Degree two and
three have the corresponding two-term identities with denominators cleared.
\end{theorem}
\begin{proof}
\leanok
\end{proof}

\begin{theorem}[Matrix powers at the maximal level increase]
\label{thm:maximal-level-paths}
\lean{Matrix.maximalStep, Matrix.pow_apply_eq_zero_of_level_gap, Matrix.pow_apply_eq_maximalStep, Matrix.pow_apply_pos_of_positive_path, Matrix.pow_apply_sign_conjugate}\leanok
Let $M$ be a finite real matrix, with integer index levels $\ell(i)$, and suppose
$M_{ij}=0$ when $\ell(j)>\ell(i)+1$. Then $(M^r)_{ij}=0$ when
$\ell(j)>\ell(i)+r$. At equality, $M^r$ has the same entry as the power of
the matrix retaining only entries with $\ell(j)=\ell(i)+1$.
For an entrywise nonnegative matrix, a positive path of length $r$ gives a
strictly positive entry of its $r$-th power. Diagonal signs $s_i^2=1$ telescope:
$(\operatorname{diag}(s)M\operatorname{diag}(s))^r_{ij}=s_i(M^r)_{ij}s_j$.
\end{theorem}
\begin{proof}
\leanok
\end{proof}

\begin{theorem}[Finite inverse expansion and first surviving pairing]
\label{thm:inverse-path-factorization}
\lean{eq_geom_sum_add_pow_mul_of_one_sub_mul_eq_one, LinearMap.apply_rightInverse_eq_of_powers_vanish, LinearMap.apply_rightInverse_ne_zero_of_powers_vanish}\leanok
In any unital ring, $(1-x)y=1$ implies the exact identity
$y=\sum_{q<r}x^q+x^ry$.
In a real algebra, let $f$ be a real linear functional and $(1+tK)Y=1$.
If $f(K^q)=0$ for $q<r$, then
\[
 f(Y)=(-t)^r f(K^rY).
\]
In particular, this pairing is nonzero if $t\ne0$ and the last factor is nonzero.
No infinite series or convergence hypothesis is used.
\end{theorem}
\begin{proof}
\leanok
\end{proof}

\begin{theorem}[Local nonvanishing of inverse pairings]
\label{thm:inverse-path-nonvanishing}\uses{thm:inverse-path-factorization}
\lean{ContinuousLinearMap.eventually_apply_rightInverse_ne_zero}\leanok
Let $K$ lie in a normed real algebra, $f$ be a continuous real linear functional,
and $Y(t)$ be continuous at zero with $Y(0)=1$ and $(1+tK)Y(t)=1$ near zero.
If $f(K^q)=0$ for $q<r$ but $f(K^r)\ne0$, then $f(Y(t))\ne0$ throughout
some punctured neighbourhood of zero. Completeness and analyticity are not assumed.
\end{theorem}
\begin{proof}
\leanok
\end{proof}

\begin{theorem}[Finite-matrix inverse path criterion]
\label{thm:finite-inverse-path}\uses{thm:inverse-path-factorization}
\lean{Matrix.eventually_apply_inv_one_add_smul_ne_zero}\leanok
Let $K$ be a finite real square matrix and $f$ a real linear functional on
matrices. If $f(K^q)=0$ for every $q<r$ and $f(K^r)\ne0$, then
$f((I+tK)^{-1})\ne0$ for all sufficiently small nonzero real $t$.
\end{theorem}
\begin{proof}
\leanok
\end{proof}

\begin{theorem}[Schur path criterion for Cholesky nonvanishing]
\label{thm:cholesky-schur-path}\uses{thm:finite-inverse-path,thm:mathlib-gram-factor}
\lean{Matrix.schur_gram_upper_blocks, Matrix.IsUpperTriangular.transpose_mul_self_first_row, Matrix.schur_gram_first_row_eq_of_pencil, Matrix.sandwichEntry, Matrix.eventually_cholesky_first_row_ne_zero}\leanok
For an upper block factor $R=\begin{pmatrix}A&B\\0&D\end{pmatrix}$ with $A$
invertible, the Schur complement of the leading block of $R^\mathsf{T}R$ is
$D^\mathsf{T}D$. Suppose $D$ is upper triangular and
\[
 A^\mathsf{T}A=I+tK,\qquad A^\mathsf{T}B=tE,\qquad B^\mathsf{T}A=tF,
 \qquad (B^\mathsf{T}B+D^\mathsf{T}D)_{0j}=0.
\]
Then
\[
 D_{00}D_{0j}=-t^2\bigl(F(I+tK)^{-1}E\bigr)_{0j}.
\]
If $(FK^qE)_{0j}=0$ for $q<r$ and $(FK^rE)_{0j}\ne0$, then for every
sufficiently small nonzero $t$, every such factorization has $D_{0j}\ne0$.
Continuity of the selected factors is not required.
\end{theorem}
\begin{proof}
\leanok
\end{proof}

\begin{definition}[Classical deformation branch]
\label{def:disk-deformation-branch}
\lean{Polynomial.diskDeformation}\leanok
For $\delta\in\mathbb R$, let $D_\delta=\{t\in\mathbb R:1-\delta+\delta t^2\ge0\}$.
The positive branch $\phi_\delta:D_\delta\to\mathbb R$ is
$\phi_\delta(t)=\sqrt{1-\delta+\delta t^2}$.
\end{definition}

\begin{theorem}[Deformation branch characterization]
\label{thm:disk-deformation-branch}\uses{def:disk-deformation-branch}
\lean{Polynomial.continuous_diskDeformation, Polynomial.diskDeformation_sq, Polynomial.diskDeformation_radical_nonneg, Polynomial.diskDeformation_mem_Icc}\leanok
The branch is continuous and satisfies $\phi_\delta(t)^2=1-\delta+\delta t^2$ on $D_\delta$.
If $0\le\delta\le1$, then $D_\delta=\mathbb R$, and the branch maps $[-1,1]$ into $[-1,1]$.
\end{theorem}
\begin{proof}
\leanok
\end{proof}

\begin{definition}[Symmetric outer-measure deformation]
\label{def:disk-deformation}\uses{def:disk-deformation-branch}
\lean{Polynomial.diskDeformedMeasure}\leanok
For a measure $\mu$ on $D_\delta$, define
\[
 \mu_\delta=\tfrac12\bigl((\phi_\delta)_*\mu+(-\phi_\delta)_*\mu\bigr).
\]
For $0\le\delta\le1$, a real source measure is identified with its measure on $D_\delta=\mathbb R$.
In the Cartesian disk application, the scale remains $\sqrt{1-x^2}$ and the inner factor measure is unchanged.
\end{definition}

\begin{theorem}[Symmetric deformation integral characterization]
\label{thm:disk-deformation-characterisation}\uses{def:disk-deformation,thm:disk-deformation-branch}
\lean{Polynomial.integral_diskDeformedMeasure, Polynomial.integral_diskDeformedMeasure_even, Polynomial.diskDeformedMeasure_ne_zero, Polynomial.finite_fiber_diskDeformation}\leanok
For a finite measure $\mu$ on $D_\delta$ supported on $[-1,1]$ and continuous $f$,
\[
\int f\,d\mu_\delta=\tfrac12\int f(\phi_\delta(t))\,d\mu(t)
+\tfrac12\int f(-\phi_\delta(t))\,d\mu(t).
\]
For even $f$ this equals the single positive-branch integral. A nonzero source gives a nonzero deformed measure.
If $\delta\ne0$, the branch has finite fibers.
\end{theorem}
\begin{proof}
\leanok
\end{proof}

\begin{theorem}[Admissibility and shifted deformation integrals]
\label{thm:disk-deformation-admissible}\uses{thm:disk-deformation-characterisation,thm:compact-polynomial-moments,def:integrable-input,def:disk-outer-measure,thm:univariate-modification-support}
\lean{Polynomial.ae_mem_Icc_diskDeformedMeasure, Polynomial.hasFiniteMoments_diskDeformedMeasure, Polynomial.nullSingletonClass_diskDeformedMeasure, Polynomial.hasPositiveDefiniteIntegral_diskDeformedMeasure, Polynomial.integral_diskOuterMeasure_diskDeformedMeasure}\leanok
Let $\mu$ be a nonzero finite measure supported on $[-1,1]$, giving zero mass to
singletons. For $0<\delta\le1$, $\mu_\delta$ is finite, even, supported on $[-1,1]$,
has all polynomial moments, and has a strictly positive polynomial pairing.
Consequently every shifted outer measure $(1-x^2)^k\,d\mu_\delta$ is admissible.
For continuous $f$, writing $d\mu^{[k]}=(1-t^2)^k\,d\mu(t)$,
\[
 \int f(x)(1-x^2)^k\,d\mu_\delta(x)
 =\frac{\delta^k}{2}\left(\int f(\phi_\delta(t))\,d\mu^{[k]}(t)
 +\int f(-\phi_\delta(t))\,d\mu^{[k]}(t)\right).
\]
The integral identity needs only finiteness and support, and also holds at $\delta=0$.
\end{theorem}
\begin{proof}
\leanok
\end{proof}

\begin{definition}[Nondegenerate even-polynomial transport]
\label{def:even-disk-transport}
\lean{Polynomial.evenDiskTransport}\leanok
For $\delta\in\mathbb R\setminus\{0\}$ and a real polynomial $p$, let
$E_p(s)=\sum_j[X^{2j}]p\,s^j$. Define
$\mathcal T_\delta p=E_p((X^2-1+\delta)/\delta)$.

\end{definition}

\begin{theorem}[Even-polynomial transport characterization]
\label{thm:even-disk-transport-characterisation}\uses{def:even-disk-transport}
\lean{Polynomial.evenDiskTransport_comp_sq, Polynomial.even_eval_evenDiskTransport, Polynomial.evenDiskTransport_eq_comp_sq}\leanok
The transported polynomial is even. For every polynomial $q$ and $\delta\ne0$,
$\mathcal T_\delta(q(X^2))=q((X^2-1+\delta)/\delta)$.
More generally, $\mathcal T_\delta p$ is the composition of
$E_p((X-1+\delta)/\delta)$ with $X^2$.
\end{theorem}
\begin{proof}
\leanok
\end{proof}

\begin{theorem}[Invertible transport of even polynomials]
\label{thm:even-disk-transport}\uses{thm:even-disk-transport-characterisation}
\lean{Polynomial.comp_sq_contract_of_even_coeff, Polynomial.eval_evenDiskTransport_sq, Polynomial.natDegree_evenDiskTransport, Polynomial.leadingCoeff_evenDiskTransport_pos, Polynomial.evenDiskTransport_inv}\leanok
Every even real polynomial is $p(x)=q(x^2)$. For $\delta\ne0$, its transport is
\[
 (\mathcal T_\delta p)(x)=q\!\left(\frac{x^2-1+\delta}{\delta}\right).
\]
For $\delta\ne0$ this preserves degree and has inverse $\mathcal T_{1/\delta}$
on even polynomials. Positive $\delta$ preserves a positive leading coefficient.
Whenever $x^2=1-\delta+\delta t^2$,
$(\mathcal T_\delta p)(x)=p(t)$.
With $E_p(s)=\sum_j[X^{2j}]p\,s^j$, every even $p$ satisfies $p(X)=E_p(X^2)$.
\end{theorem}
\begin{proof}
\leanok
\end{proof}

\begin{theorem}[Canonical even outer polynomials under deformation]
\label{thm:deformed-outer-normalization}\uses{thm:disk-deformation-admissible,thm:even-disk-transport,def:cartesian-disk-system}
\lean{Polynomial.integralInner_evenDiskTransport, Polynomial.integralInner_normalized_evenDiskTransport, Polynomial.normalized_evenDiskTransport_eq, MvPolynomial.CartesianDiskSystem.outer_even_diskDeformation_eq}\leanok
Let $\mu$ be a finite even measure supported in $[-1,1]$ almost everywhere,
with finite polynomial moments, and assume that $\mu$ and $\mu_\delta$ have
positive definite polynomial pairings. For $0<\delta\le1$, transport multiplies
the shifted pairing of even polynomials by $\delta^k$.
If $p_n^{[k]}$ and $p_{\delta,n}^{[k]}$ are the positive-leading orthonormal
outer sequences, then for $a,b\ge0$ and $0<\delta\le1$,
\[
 p_{\delta,2a}^{[2b]}=\delta^{-b}\mathcal T_\delta p_{2a}^{[2b]}.
\]
\end{theorem}
\begin{proof}
\leanok
\end{proof}

\begin{theorem}[Canonical disk deformation and quadratic Gram pullback]
\label{thm:deformed-disk-gram}\uses{thm:deformed-outer-normalization,thm:disk-polynomial-lift,thm:disk-coordinate-characterisation}
\lean{MvPolynomial.eval_cartesianDiskPolynomial_evenDiskTransport, MvPolynomial.CartesianDiskSystem.inner_eq, MvPolynomial.CartesianDiskSystem.eval_polynomial_diskDeformation, MvPolynomial.CartesianDiskSystem.integral_quadratic_polynomial_diskDeformation}\leanok
Let $S,T$ be the canonical positive-leading Cartesian systems for
$(\mu,\nu)$ and $(\mu_\delta,\nu)$ under the preceding hypotheses, with finite
factor measures supported on $[-1,1]$ and positive polynomial pairings.
For $n=2a+2b$, $k=2b$, $t\in[-1,1]$, and $z\in\mathbb R$,
\[
 T_{n,k}(\pm\phi_\delta(t),\sqrt{1-\phi_\delta(t)^2}\,z)
 =S_{n,k}(t,\sqrt{1-t^2}\,z).
\]
Thus for any two even-even indices $i,j$ and any $A,B,C\in\mathbb R$,
\[
 \int (A+Bx^2+Cy^2)T_iT_j\,d\sigma_\delta
 =\int \bigl[(A+B)+\delta\{Cy^2-B(1-x^2)\}\bigr]S_iS_j\,d\sigma,
\]
where $\sigma,\sigma_\delta$ are the corresponding Cartesian disk measures.

\end{theorem}
\begin{proof}
\leanok
\end{proof}

\rvnote{thm:deformed-disk-gram}{
This identifies the even-even quadratic Gram entries as an affine function of $\delta$; it does not assert nonvanishing of Cholesky entries.
}

\begin{theorem}[Leading-degree pairing and double Jacobi raising]
\label{thm:double-jacobi-pairing}\uses{thm:jacobi-symmetric-coefficients,thm:jacobi-pairwise,thm:jacobi-positive,thm:symmetric-domain-double-c,thm:symmetric-domain-double-a,thm:symmetric-domain-double-b,thm:orthogonal-lower-degree,thm:jacobi-symmetric-conversion}
\lean{Polynomial.Sequence.IsOrthogonal.integralInner_eq_leadingCoeff_div_mul, Polynomial.Sequence.IsOrthogonal.integralInner_pos_of_degree_eq, Polynomial.jacobi_symmetric_double_raise, Polynomial.jacobiSymmetricDoubleRaiseA_pos, Polynomial.jacobiSymmetricDoubleRaiseB_neg, Polynomial.jacobiSymmetricDoubleRaiseC_pos, Polynomial.integralInner_jacobi_double_raise, Polynomial.integralInner_jacobi_double_raise_top_pos, Polynomial.integralInner_jacobi_double_raise_middle_neg, Polynomial.integralInner_jacobi_double_raise_bottom_pos}\leanok
For an orthogonal degree-indexed sequence with finite moments, a polynomial of degree $n$ pairs with its $n$th member through the ratio of leading coefficients. Two symmetric Jacobi parameter raisings express $P_n^{(\beta,\beta)}$ in target degrees $n,n-2,n-4$, with strictly positive, negative, positive coefficients when present, throughout $\beta>-1$. Missing initial-degree terms have zero coefficients. The corresponding pairings in the target weight have the same strict signs.
\end{theorem}
\begin{proof}
\leanok
\end{proof}

\begin{theorem}[Maximal inner-step disk multiplication signs]
\label{thm:disk-quadratic-step-signs}\uses{thm:double-jacobi-pairing,def:multivariate-integrable-pair,thm:multivariate-moments-basic,thm:dunkl-xu-separated,thm:disk-integrability,thm:jacobi-symmetric-measure-shift}
\lean{MvPolynomial.integralInner_y_sq_dunklXuDisk_separated, MvPolynomial.integralInner_y_sq_dunklXuDisk_inner_add_two, MvPolynomial.integralInner_X_sq_jacobi_add_two_pos, MvPolynomial.integralInner_y_sq_dunklXuDisk_up_pos, MvPolynomial.integralInner_y_sq_dunklXuDisk_same_degree_neg, MvPolynomial.integralInner_y_sq_dunklXuDisk_down_pos, MvPolynomial.integralInner_y_sq_dunklXuDisk_inner_add_two_eq_zero}\leanok
In the original Jacobi-normalized Dunkl--Xu basis, multiplication by $y^2$ with inner-index increase two has total-degree changes $+2,0,-2$ with strict signs $+,-,+$. Invalid lower degrees are omitted; the top coefficient remains positive at outer degrees zero and one. These are signs of actual weighted pairings, valid for every disk parameter $\beta>-1$.
\end{theorem}
\begin{proof}
\leanok
\end{proof}

\begin{theorem}[Normalized quadratic perturbation and exact inner zeros]
\label{thm:disk-quadratic-matrix-support}\uses{thm:disk-quadratic-step-signs,thm:disk-normalization}
\lean{MvPolynomial.dunklXuQuadraticMatrix, MvPolynomial.integralInner_outer_mul_dunklXuDisk_eq_zero, MvPolynomial.integralInner_mul_normalizedDunklXuBasis, MvPolynomial.dunklXuQuadraticMatrix_symm, MvPolynomial.dunklXuQuadraticMatrix_of_inner_ne, MvPolynomial.dunklXuQuadraticMatrix_eq_zero_of_inner_gap, MvPolynomial.dunklXuQuadraticMatrix_eq_zero_of_inner_parity}\leanok
In the normalized basis, the perturbation $Cy^2-B(1-x^2)$ has symmetric multiplication matrix. Every polynomial in $x$ preserves the inner index, so off-inner-diagonal entries depend only on $C$. Inner gaps larger than two and opposite inner parities vanish exactly. Positive norm scalings preserve the preceding strict signs.
\end{theorem}
\begin{proof}
\leanok
\end{proof}

\begin{theorem}[Boundary-to-boundary maximal paths]
\label{thm:maximal-sandwich-paths}\uses{thm:maximal-level-paths}
\lean{Matrix.sum_mul_pow_mul_eq_zero_of_level_gap, Matrix.sum_mul_pow_mul_eq_maximalStep, Matrix.sum_mul_pow_mul_sign_conjugate, Matrix.sum_mul_pow_mul_pos_of_positive_path}\leanok
For a boundary row, $q$ interior matrix steps and a boundary column, a level gap exceeding $q+2$ forces the sandwich to vanish if each edge increases level by at most one. At gap exactly $q+2$, only maximal steps contribute. Diagonal signs telescope through both boundary edges. Nonnegative signed entries and one positive path imply a strictly positive signed sandwich.
\end{theorem}
\begin{proof}
\leanok
\end{proof}

\begin{theorem}[A noncancelling disk path in every degree]
\label{thm:disk-all-degree-path}\uses{thm:disk-quadratic-matrix-support,thm:maximal-sandwich-paths}
\lean{MvPolynomial.dunklXuInteriorFunctional, MvPolynomial.dunklXuEvenIndex, MvPolynomial.DunklXuEvenInterior, MvPolynomial.dunklXuInteriorPath, MvPolynomial.dunklXuInteriorPath_degree, MvPolynomial.dunklXuInteriorPath_step, MvPolynomial.dunklXuEvenY2Entry_signed_nonneg, MvPolynomial.dunklXuInterior_sandwich_eq_zero, MvPolynomial.dunklXuInterior_first_coefficient_ne_zero}\leanok
Reindex the even-even sector by $(n,k)=(2a+2b,2b)$. For every $r\ge2$, let $I=\{(a,b):a+b<r\}$, with boundary indices $(r,0)$ and $(0,r)$. For the perturbation $y^2-B(1-x^2)$, write its interior block as $K$ and boundary row and column as $F,G$. For any interior matrix $M$, define the real linear functional
$\Phi(M)=\sum_{i,j\in I}F_iM_{ij}G_j=FMG$.
Then $FK^qG=0$ for $q<r-2$, but $FK^{r-2}G\ne0$. The diagonal sign $(-1)^a$ makes all maximal steps nonnegative. The strictly positive path enters at $(r-2,1)$, follows $(r-3,2),\ldots,(0,r-1)$, and exits to $(0,r)$.
\end{theorem}
\begin{proof}
\leanok
\end{proof}

\begin{theorem}[Nonzero finite disk inverse sandwich]
\label{thm:disk-inverse-sandwich}\uses{thm:disk-all-degree-path,thm:finite-inverse-path}
\lean{MvPolynomial.dunklXuQuadraticMatrix_scale, MvPolynomial.dunklXuInteriorFunctional_pow_scale, MvPolynomial.dunklXuInterior_first_coefficient_ne_zero_of_vertical_ne_zero, MvPolynomial.eventually_dunklXuInterior_inverse_ne_zero_of_vertical_ne_zero}\leanok
For any $C\ne0$, any $B$, and any $r\ge2$, the quadratic perturbation $Cy^2-B(1-x^2)$, with interior block $K$ and boundary factors $F,G$ as in Theorem~\ref{thm:disk-all-degree-path}, has nonzero first allowed sandwich coefficient. Scaling contributes $C^{q+2}$ to the $q$th power sandwich. Consequently $F(I+tK)^{-1}G\ne0$ for every sufficiently small nonzero real $t$, where the neighborhood may depend on $r,B,C$ and the disk parameter.
\end{theorem}
\begin{proof}
\leanok
\end{proof}

\begin{theorem}[Forced distant factor entry for the even disk pencil]
\label{thm:disk-quadratic-cholesky-path}\uses{thm:disk-inverse-sandwich,thm:cholesky-schur-path}
\lean{MvPolynomial.dunklXuQuadraticLeftBlock, MvPolynomial.dunklXuQuadraticRightBlock, MvPolynomial.dunklXuQuadraticBoundaryBlock, MvPolynomial.dunklXuInteriorFunctional_eq_sandwich, MvPolynomial.dunklXuQuadraticBoundaryBlock_first_last, MvPolynomial.eventually_dunklXuQuadratic_cholesky_first_last_ne_zero}\leanok
For the even-sector pencil $I+tK_{\rm full}$ of $Cy^2-B(1-x^2)$, $C\ne0$, and $r\ge2$, split at degree $2r$. Any compatible upper block Gram factor $\left(\begin{smallmatrix}A&B_0\\0&D\end{smallmatrix}\right)$ with invertible $A$ and upper triangular $D$ has $D_{0,r}\ne0$ for all sufficiently small nonzero $t$. Compatibility means its four transpose-product blocks equal those of the displayed pencil.
\end{theorem}
\begin{proof}
\leanok
\end{proof}

\rvnote{thm:disk-quadratic-cholesky-path}{
The canonical ordered-factor conclusion is given in Theorem~\ref{thm:canonical-quadratic-distant-entry}.
}

\begin{theorem}[Parity sectors of a complete Gram factor]
\label{thm:gram-factor-sectors}
\lean{Matrix.IsUpperTriangular.apply_eq_zero_of_gram_sectors, Matrix.IsUpperTriangular.transpose_mul_self_sector, Matrix.IsUpperTriangular.gram_sector_blocks}\leanok
An upper triangular real factor with nonzero diagonal preserves every sector decomposition of its transpose-product Gram matrix: all cross-sector factor entries vanish. Restricting to a complete sector therefore commutes with forming the Gram matrix. Splitting such a sector into an earlier interior and a later boundary gives the usual four block Gram identities, even when other sectors are interleaved in the full ordering.
\end{theorem}
\begin{proof}
\leanok
\end{proof}
\end{supporting}


\begin{theorem}[Reflection sectors of Cartesian disk quadratics]
\label{thm:cartesian-quadratic-parity}\uses{thm:gram-factor-sectors,thm:canonical-polynomial-gram-factor,thm:polynomial-density-gram,thm:disk-polynomial-lift,thm:cartesian-disk-system-moments}
\lean{MeasureTheory.map_prod_neg_left, MeasureTheory.map_prod_neg_right, MvPolynomial.CartesianDiskSystem.eval_polynomial_neg_inner, MvPolynomial.CartesianDiskSystem.eval_polynomial_neg_outer, MvPolynomial.CartesianDiskSystem.integralInner_coordinateQuadratic_eq_zero_of_parity, MvPolynomial.CartesianDiskSystem.integralGramFactor_coordinateQuadratic_eq_zero_of_parity}\leanok
For factor measures $\mu,\nu$ on $\mathbb R$ with $\nu$ s-finite, evenness of either factor makes $\mu\times\nu$ invariant under reflection of its corresponding coordinate.
For even admissible outer and inner factor measures, the Cartesian polynomial indexed by $(n,k)$ has reflection signs $(-1)^{n-k}$ and $(-1)^k$. A diagonal quadratic density preserves these reflections. Its Gram matrix and canonical positive upper triangular factor have exact zeros between different parity pairs $((n-k)\bmod2,k\bmod2)$. The canonical-factor statement assumes a nonnegative density and the moments and positivity needed to define that factor.
\end{theorem}
\begin{proof}
\leanok
\end{proof}

\begin{supporting}[Reference details: Actual quadratic Gram pencil and full-factor implication and related results]
\begin{theorem}[Actual quadratic Gram pencil and full-factor implication]
\label{thm:deformed-quadratic-pencil}\uses{thm:disk-quadratic-cholesky-path,thm:gram-factor-sectors,thm:polynomial-density-gram,thm:dunkl-xu-cartesian-normalization,thm:deformed-disk-gram}
\lean{MvPolynomial.integral_quadratic_deformed_eq_dunklXu_pencil, MvPolynomial.integralInner_deformed_quadratic_eq_evenGramPencil, MvPolynomial.eventually_dunklXu_full_factor_entry_ne_zero}\leanok
Under $x^2=1-\delta+\delta t^2$, the actual modified even-sector Gram matrix is $(A+B)I+\delta K$, where $K$ is the normalized disk multiplication matrix of $Cy^2-B(1-x^2)$. For $A+B=1$, $C\ne0$, every complete-sector embedding into an upper triangular full factor inherits the distant-entry nonvanishing for all sufficiently small nonzero $\delta$. The embedding includes the entire sector and respects interior-before-boundary and increasing boundary order.
\end{theorem}
\begin{proof}
\leanok
\end{proof}

\begin{theorem}[Admissible deformed Jacobi factors and positive rescaling]
\label{thm:quadratic-witness-admissible}\uses{thm:disk-deformation-admissible,thm:quadratic-disk-separation,thm:polynomial-density-gram,thm:canonical-polynomial-gram-factor,thm:jacobi-positive,thm:cartesian-disk-bases,thm:jacobi-moments,thm:cartesian-disk-system-existence}
\lean{Polynomial.hasPositiveDefiniteIntegral_diskDeformed_jacobiMeasure, MvPolynomial.dunklXuDeformedCartesianSystem, MvPolynomial.ae_coordinateQuadratic_nonneg_cartesianDiskMeasure, MvPolynomial.CartesianDiskSystem.hasPositiveDefiniteIntegral_coordinateQuadratic, MvPolynomial.integralGram_withDensity_smul, MvPolynomial.integralGramFactor_eq_sqrt_smul_of_gram_eq, MvPolynomial.integralGramFactor_eq_zero_iff_of_gram_scale}\leanok
For $\beta>-1$ and $0<\delta\le1$, the deformed outer Jacobi factor defines an actual positive Cartesian system. The scalar conditions $A\ge0$, $A+B>0$, $A+C\ge0$ make the diagonal quadratic modification nonnegative and preserve finite moments and strict polynomial positivity. A positive scalar Gram rescaling multiplies its canonical factor by the positive square root, preserving every zero.
\end{theorem}
\begin{proof}
\leanok
\end{proof}

\begin{theorem}[Complete degree and even-sector ordering]
\label{thm:ordered-cartesian-even-sector}\uses{thm:disk-all-degree-path}
\lean{MvPolynomial.diskDegreeOrder, MvPolynomial.CartesianDiskSystem.orderedDegreeBasis, MvPolynomial.dunklXuEvenSectorEmbedding_injective, MvPolynomial.dunklXuEvenSectorEmbedding_covers, MvPolynomial.dunklXuEvenSectorEmbedding_before, MvPolynomial.dunklXuEvenSectorEmbedding_strictMono}\leanok
The complete Cartesian basis through degree $2r$ admits an explicit increasing-total-degree, increasing-inner-index enumeration. Its entire even-even sector is the interior $a+b<r$ followed by the top block $(r-j,j)$, $0\le j\le r$, in half-degree coordinates. The embedding is injective, covers every even-even index, and has the ordering required by sector elimination.
\end{theorem}
\begin{proof}
\leanok
\end{proof}

\begin{theorem}[Canonical factor nonvanishing without quadratic normalization]
\label{thm:canonical-quadratic-distant-entry}\uses{thm:deformed-quadratic-pencil,thm:cartesian-quadratic-parity,thm:quadratic-witness-admissible}
\lean{MvPolynomial.eventually_integralGramFactor_deformed_quadratic_ne_zero, MvPolynomial.eventually_integralGramFactor_admissible_quadratic_ne_zero}\leanok
Let $A\ge0$, $A+B>0$, $A+C\ge0$, $C\ne0$ and $r\ge2$. For all sufficiently small positive $\delta$, the actual canonical Gram factor for the quadratically modified deformed measure has a nonzero first-to-last top-even-sector entry.  No $A+B=1$ normalization remains.
The conclusion holds for any complete ordered source basis with the stated sector enumeration; the intermediate measure satisfies the required admissibility conditions.
\end{theorem}
\begin{proof}
\leanok
\end{proof}

\begin{definition}[Multivariate polynomial density measure]
\label{def:multivariate-polynomial-density-measure}
\lean{MvPolynomial.polynomialDensityMeasure}\leanok
For a real polynomial $p$ in finitely many variables, put
$D_p=\{x:p(x)\ge0\}$. For a measure $\eta$ on $D_p$, define its polynomial
density measure as the coordinate pushforward of $p\,d\eta$.
\end{definition}

\begin{theorem}[Multivariate polynomial density characterization]
\label{thm:multivariate-polynomial-density-measure}\uses{def:multivariate-polynomial-density-measure}
\lean{MvPolynomial.polynomialDensityMeasure_apply}\leanok
For every measurable set $E$, the density measure assigns
$\int_{\iota^{-1}(E)}p(x)\,d\eta(x)$ as a nonnegative integral.
\end{theorem}
\begin{proof}
\leanok
\end{proof}

\begin{definition}[Deformed quadratic disk measure]
\label{def:deformed-quadratic-disk-measure}\uses{def:disk-deformation,def:disk-measure,def:multivariate-polynomial-density-measure}
\lean{MvPolynomial.DunklXuDeformedQuadraticInput, MvPolynomial.dunklXuDeformedQuadraticMeasure}\leanok
Let $\beta\in\mathbb R$ and $0\le\delta\le1$. Form the Cartesian disk measure
$\sigma_{\beta,\delta}$ from the symmetrically deformed outer Jacobi measure
with parameters $(\beta+1/2,\beta+1/2)$ and inner Jacobi measure with parameters
$(\beta,\beta)$. For $r(x,y)=A+Bx^2+Cy^2$, require $r\ge0$ almost everywhere for $\sigma_{\beta,\delta}$.
Transport the source to $\{r\ge0\}$ and take its polynomial density measure;
this is the deformed quadratic disk measure.
\end{definition}

\begin{theorem}[Deformed quadratic disk measure characterization]
\label{thm:deformed-quadratic-disk-measure}\uses{def:deformed-quadratic-disk-measure,thm:quadratic-witness-admissible,thm:multivariate-polynomial-density-measure}
\lean{MvPolynomial.DunklXuDeformedQuadraticInput.of_nonneg, MvPolynomial.dunklXuDeformedQuadraticMeasure_apply}\leanok
The inequalities $A\ge0$, $A+B\ge0$, $A+C\ge0$ imply that $A+Bx^2+Cy^2$ is nonnegative almost everywhere for $\sigma_{\beta,\delta}$. For every measurable set $E$, the modified measure assigns the
nonnegative integral of $A+Bx^2+Cy^2$ over $E$ against the transported source
on $\{A+Bx^2+Cy^2\ge0\}$.
\end{theorem}
\begin{proof}
\leanok
\end{proof}

\begin{theorem}[Varying-measure quadratic obstruction in every degree]
\label{thm:quadratic-varying-measure-obstruction}\uses{thm:canonical-quadratic-distant-entry,thm:ordered-cartesian-even-sector,thm:deformed-quadratic-disk-measure}
\lean{MvPolynomial.dunklXuDeformedDegreeBasis, MvPolynomial.dunklXuDeformedQuadraticMeasure_admissible, MvPolynomial.exists_integralGramFactor_quadratic_distant_entry}\leanok
Fix $\beta>-1$ and a diagonal quadratic $A+Bx^2+Cy^2$ with $A\ge0$, $A+B>0$, $A+C\ge0$, and $C\ne0$. For every $r\ge2$, there exists $0<\delta\le1$ such that the canonical factor on the complete polynomial space through degree $2r$, in degree/inner order, has a nonzero entry from $(n,k)=(2r,0)$ to $(2r,2r)$. The measures, bases, positivity and ordering are all constructed. Thus no finite inner bandwidth can hold uniformly over these admissible factor measures. The measure may depend on $r$.
\end{theorem}
\begin{proof}
\leanok
\end{proof}

\rvnote{thm:quadratic-varying-measure-obstruction}{
The chosen measure may depend on $r$; this does not assert unbounded fill for the original uniform disk or for one fixed deformation.
}

\begin{theorem}[Canonical factors of ordered separated connections]
\label{thm:ordered-separated-canonical}\uses{thm:ordered-cartesian-even-sector,thm:canonical-polynomial-gram-factor,thm:polynomial-density-gram,thm:cartesian-connection-support}
\lean{MvPolynomial.CartesianDiskSystem.orderedConnectionMatrix, MvPolynomial.CartesianDiskSystem.integralGram_eq_orderedConnectionMatrix_transpose_mul, MvPolynomial.CartesianDiskSystem.orderedConnectionMatrix_isUpperTriangular, MvPolynomial.CartesianDiskSystem.orderedConnectionMatrix_diag_pos, MvPolynomial.CartesianDiskSystem.integralGramFactor_eq_orderedConnectionMatrix, MvPolynomial.CartesianDiskSystem.integralGramFactor_eq_zero_of_separated_inner_gap}\leanok
For factor measures $\mu,\nu$ supported in $[-1,1]$, in total-degree/inner-index order, a separated polynomial modification has a positive upper triangular connection matrix. It equals the canonical Gram factor whenever the target polynomial pairings agree. For the actual polynomial-density measure, the factor vanishes whenever the absolute inner-index gap exceeds the degree of the inner modifying polynomial.
\end{theorem}
\begin{proof}
\leanok
\end{proof}
\end{supporting}


\begin{definition}[Uniform Cartesian inner bandwidth]
\label{def:uniform-cartesian-inner-bandwidth}\uses{thm:ordered-cartesian-even-sector,thm:canonical-polynomial-gram-factor,def:multivariate-polynomial-density-measure}
\lean{MvPolynomial.HasUniformCartesianInnerBandwidth}\leanok
A polynomial density has uniform Cartesian inner bandwidth if there exists one finite bound $w$ such that, for every admissible pair of factor measures, every Cartesian source system and every complete degree section, the canonical modified Gram factor has zero entries whenever the absolute inner-index gap exceeds $w$. The factor measures have finite moments, strictly positive polynomial pairings and support in $[-1,1]$; the inner measure is even. Only densities nonnegative almost everywhere for the Cartesian source are tested, using their polynomial density measures on the nonnegative locus. The modified pairing must also have the moments and positivity required to define its canonical factor. The bound is chosen before both the measures and the section.
\end{definition}

\begin{theorem}[Classification of uniform diagonal-quadratic sparsity]
\label{thm:uniform-quadratic-classification}\uses{def:uniform-cartesian-inner-bandwidth,thm:ordered-separated-canonical,thm:quadratic-varying-measure-obstruction}
\lean{MvPolynomial.hasUniformCartesianInnerBandwidth_separated, MvPolynomial.not_hasUniformCartesianInnerBandwidth_coordinateQuadratic, MvPolynomial.coordinateQuadratic_separated_of_uniform_inner_bandwidth, MvPolynomial.hasUniformCartesianInnerBandwidth_coordinateQuadratic_iff}\leanok
For a nonzero diagonal quadratic $u=A+Bx^2+Cy^2$ nonnegative on the disk, uniform finite Cartesian inner bandwidth holds if and only if $C(A+B)=0$. Nonnegativity is stated equivalently as $A\ge0$, $A+B\ge0$, $A+C\ge0$. This is a classification over all admissible factor measures, with fixed Cartesian orientation and degree/inner ordering.
\end{theorem}
\begin{proof}
\leanok
\end{proof}

\begin{supporting}[Reference details: Polynomial matrix pencil and adjugate numerator and related results]
\begin{definition}[Polynomial matrix pencil and adjugate numerator]
\label{def:matrix-polynomial-pencil}
\lean{Matrix.polynomialPencil, Matrix.inversePencilNumerator}\leanok
For a finite real square matrix $K$, let $P_K(X)=I+XK$ be its polynomial matrix pencil. For compatible real matrices $F,E$, define the polynomial $q_{ij}(X)=(F\operatorname{adj}(P_K(X))E)_{ij}$.
\end{definition}

\begin{theorem}[Evaluation of the polynomial pencil and numerator]
\label{thm:matrix-polynomial-pencil}\uses{def:matrix-polynomial-pencil}
\lean{Matrix.eval_polynomialPencil, Matrix.eval_inversePencilNumerator}\leanok
For every real $t$, $P_K(t)=I+tK$ and $q_{ij}(t)=(F\operatorname{adj}(I+tK)E)_{ij}$.
\end{theorem}
\begin{proof}
\leanok
\end{proof}

\begin{definition}[Invertible pencil parameters]
\label{def:invertible-pencil-parameter}
\lean{Matrix.InvertiblePencilParameter}\leanok
For a finite real square matrix $K$, let $D_K=\{t\in\mathbb R: I+tK\text{ is invertible}\}$. An invertible pencil parameter is an element of $D_K$.
\end{definition}

\begin{theorem}[Characterization of invertible pencil parameters]
\label{thm:invertible-pencil-parameter}\uses{thm:matrix-polynomial-pencil}
\lean{Matrix.isUnit_pencil_iff_det_ne_zero, Matrix.eventually_cofinite_isUnit_pencil}\leanok
A real parameter belongs to $D_K$ if and only if $\det P_K(t)\ne0$. The complement of $D_K$ is finite.
\end{theorem}
\begin{proof}
\leanok
\end{proof}

\begin{definition}[Classical inverse of an affine pencil]
\label{def:inverse-pencil}\uses{def:invertible-pencil-parameter}
\lean{Matrix.inversePencil}\leanok
For $t\in D_K$, define $R_K(t)$ as the inverse of $I+tK$.
\end{definition}

\begin{theorem}[Inverse-pencil identities]
\label{thm:inverse-pencil}\uses{def:inverse-pencil}
\lean{Matrix.inversePencil_mul}\leanok
For every $t\in D_K$, $R_K(t)(I+tK)=(I+tK)R_K(t)=I$.
\end{theorem}
\begin{proof}
\leanok
\end{proof}

\begin{theorem}[Finite exceptional parameters for inverse pencils]
\label{thm:finite-pencil-exceptions}\uses{thm:invertible-pencil-parameter,thm:inverse-pencil}
\lean{Matrix.inversePencil_sandwich_eq, Matrix.eventually_cofinite_inversePencil_sandwich_ne_zero}\leanok
For fixed finite real matrices $F,K,E$ and $t\in D_K$, $(FR_K(t)E)_{ij}=q_{ij}(t)/\det P_K(t)$, with nonzero denominator. If this entry is nonzero at one parameter in $D_K$, then outside a finite exceptional set every real parameter belongs to $D_K$ and its entry is nonzero. All singular parameters belong to the exceptional set.
\end{theorem}
\begin{proof}
\leanok
\end{proof}

\begin{theorem}[Cofinite nonvanishing of the quadratic canonical factor]
\label{thm:cofinite-quadratic-factors}\uses{thm:finite-pencil-exceptions,thm:quadratic-varying-measure-obstruction}
\lean{MvPolynomial.dunklXu_full_factor_entry_ne_zero_of_inverse, MvPolynomial.eventually_cofinite_dunklXu_full_factor_entry_ne_zero, MvPolynomial.eventually_cofinite_integralGramFactor_admissible_quadratic_ne_zero, MvPolynomial.eventually_cofinite_integralGramFactor_quadratic_distant_entry}\leanok
At an invertible interior-pencil parameter, a nonzero inverse sandwich forces the corresponding distant factor entry to be nonzero. For each fixed $r\ge2$ and admissible nonseparated quadratic, only finitely many deformation parameters can fail the distant-entry implication. In particular, every admissible $0<\delta\le1$ outside that finite set has a nonzero $(2r,0)$ to $(2r,2r)$ entry in its actual canonical factor. The finite exceptional set can depend on $r$.
\end{theorem}
\begin{proof}
\leanok
\end{proof}
\end{supporting}


\begin{theorem}[One fixed deformation with arbitrarily large inner gaps]
\label{thm:fixed-quadratic-obstruction}\uses{thm:cofinite-quadratic-factors}
\lean{Real.exists_Ioo_forall_of_eventually_cofinite, MvPolynomial.exists_fixed_deformation_integralGramFactor_all_distant_entries}\leanok
Fix $\beta>-1$ and $A\ge0$, $A+B>0$, $A+C\ge0$, $C\ne0$. There exists one admissible $0<\delta\le1$ such that, for every $r\ge2$, the canonical factor of the complete degree-$2r$ section has a nonzero entry from $(2r,0)$ to $(2r,2r)$. The same deformed measure is used for every degree.
\end{theorem}
\begin{proof}
\leanok
\end{proof}

\rvnote{thm:fixed-quadratic-obstruction}{
The parameter is not specified numerically, and this does not assert the result for the original uniform disk at $\delta=1$.
}

\chapter{Cartesian disk connections}\label{chap:disk-connections}
\rvsummary{Cartesian disk systems, normalized connections and exact finite sections.}
We identify the separated construction with the normalized Dunkl--Xu basis.
The connection support and degree-padding results determine which finite
matrix sections give the exact coefficients.

\begin{definition}[General Cartesian disk systems]
\label{def:cartesian-disk-system}\uses{def:disk-polynomial-lift,def:disk-outer-measure,def:orthogonal-sequence}
\lean{MvPolynomial.cartesianDiskPolynomial, MvPolynomial.CartesianDiskSystem}\leanok
Let $\mu$ be a measure on $[-1,1]$ and $\nu$ a real measure.
Write $\bar\mu=\iota_*\mu$ for the real-coordinate measure. The $k$th outer
measure is the shifted disk measure of Definition~\ref{def:disk-outer-measure},
using the transport of $\bar\mu$ to $D_k$. Its geometric Cartesian measure
is the pushforward in Definition~\ref{def:disk-measure}.
A real source supported on $[-1,1]$ is transported to the interval without loss. A Cartesian disk system for
$(\mu,\nu)$ consists of an inner polynomial sequence $(q_k)$ and, for each $k$,
an outer polynomial sequence $(p_r^{[k]})$ such that $(q_k)$ is orthogonal for
$\nu$ and $(p_r^{[k]})$ for the $k$th outer measure, each $q_k$ has the parity
of $k$, all leading coefficients are positive, and all squared norms are one.
Its polynomials are
\[
 P_{n,k}(x,y)=p_{n-k}^{[k]}(x)\,L_kq_k(x,y),\qquad 0\le k\le n,
\]
with $L_k$ the parity lift of Definition~\ref{def:disk-polynomial-lift}; on
the disk with $|x|<1$ this is
$p_{n-k}^{[k]}(x)(1-x^2)^{k/2}q_k(y/\sqrt{1-x^2})$.
At the tips $(\pm1,0)$, the values are given by the polynomial lift $L_kq_k$.
\end{definition}

\begin{theorem}[Existence of Cartesian disk systems]
\label{thm:cartesian-disk-system-existence}\uses{def:cartesian-disk-system,thm:univariate-modification-support}
\lean{MvPolynomial.cartesianDiskSystem}\leanok
Let $\mu$ and $\nu$ be measures on $\mathbb R$ with all polynomial moments
finite and positive definite polynomial pairings. Suppose $\mu$ assigns zero
mass to $\mathbb R\setminus[-1,1]$ and $\nu$ is invariant under $x\mapsto-x$.
Then there exists a Cartesian disk system for the restriction of $\mu$ to
$[-1,1]$ and the inner measure $\nu$.
\end{theorem}
\begin{proof}
\leanok
\end{proof}

\begin{supporting}[Reference details: Moments supplied by a Cartesian disk system]
\begin{theorem}[Moments supplied by a Cartesian disk system]
\label{thm:cartesian-disk-system-moments}\uses{def:cartesian-disk-system,thm:koornwinder-moments,thm:orthogonal-moments,thm:disk-coordinate-characterisation,thm:disk-outer-measure}
\lean{MvPolynomial.CartesianDiskSystem.hasFiniteMoments_inner, MvPolynomial.CartesianDiskSystem.hasFiniteMoments_outer, MvPolynomial.CartesianDiskSystem.hasFiniteMoments, MvPolynomial.hasFiniteMoments_cartesianDiskMeasure_of_moments}\leanok
A Cartesian disk system supplies all moments of both one-dimensional factors.
If the outer measure is supported almost everywhere on $[-1,1]$, its Cartesian
measure has all polynomial moments. More generally, finite moments of the two
factors and this outer support condition suffice; no support condition on the
inner factor is needed for this conclusion.
\end{theorem}
\begin{proof}
\leanok
\end{proof}
\end{supporting}


\begin{theorem}[Cartesian orthonormal degree bases]
\label{thm:cartesian-disk-bases}\uses{thm:measure-conversion-cholesky,thm:disk-polynomial-lift,thm:cartesian-disk-system-moments,thm:mathlib-total-degree-finite}
\lean{MvPolynomial.integralInner_cartesianDiskPolynomial, MvPolynomial.CartesianDiskSystem.integralInner_polynomial_eq_zero, MvPolynomial.CartesianDiskSystem.integralInner_polynomial_self, MvPolynomial.CartesianDiskSystem.degreeBasis, MvPolynomial.CartesianDiskSystem.integralGram_degreeBasis, MvPolynomial.CartesianDiskSystem.hasPositiveDefiniteIntegral}\leanok
Assume also that the inner measure $\nu$ is supported in $[-1,1]$.
The Cartesian polynomials are orthonormal and form a basis of every complete
finite total-degree space. Their measure is positive definite on the polynomial
algebra. No Jacobi formula is assumed for the one-dimensional factors.
\end{theorem}
\begin{proof}
\leanok
\end{proof}

\begin{theorem}[Separated polynomial disk modifications]
\label{thm:cartesian-separated-modification}\uses{thm:disk-polynomial-lift,thm:cartesian-disk-system-moments,def:disk-outer-modification,def:subtype-polynomial-density-measure,def:multivariate-integrable-pair,thm:multivariate-moments-basic,thm:cartesian-disk-system-existence}
\lean{MvPolynomial.cartesianDiskModification, MvPolynomial.eval_cartesianDiskModification, MvPolynomial.integral_cartesianDiskModification, MvPolynomial.integralInner_cartesianDiskModification, MvPolynomial.modifiedCartesianDiskSystem, Polynomial.integralInner_diskOuterModification_transfer}\leanok
Write $A(x)=1-x^2$. Let $h,g$ be nonzero nonnegative polynomial densities for
$\mu,\nu$, with $g$ even and $\deg g\le2\ell$. The actual bivariate polynomial
\[
 u(x,\sqrt{A(x)}z)=A(x)^\ell h(x)g(z)
\]
changes the Cartesian measure to the pushforward of
$A^\ell h\,d\mu\otimes g\,d\nu$, provided the inner measure $\nu$ is supported
in $[-1,1]$. Under this support assumption the identity holds against all
continuous real-valued test functions on $\mathbb R^2$, in particular for polynomial pairings. The modified
positive-leading orthonormal Cartesian system exists even when the density
vanishes in the interior.
\end{theorem}
\begin{proof}
\leanok
\end{proof}

\begin{theorem}[Degree-independent Cartesian connection support]
\label{thm:cartesian-connection-support}\uses{thm:cartesian-disk-bases,thm:cartesian-separated-modification}
\lean{MvPolynomial.CartesianDiskSystem.connectionEntry_eq_zero_of_outside_band, MvPolynomial.CartesianDiskSystem.connectionEntry_eq_zero_of_outside_rectangle, MvPolynomial.CartesianDiskSystem.connectionEntry_diag_pos}\leanok
For the separated polynomial disk modification of Theorem~\ref{thm:cartesian-separated-modification}, the coefficient from source $(n,k)$ to target
$(m,j)$ can be nonzero only if
\[
 n-m=s,\quad k-j=2r,\qquad
 0\le s\le\deg h+2\ell,\quad 0\le r\le\ell.
\]
If $\mu$ and $h$ are also even and $\deg h\le2a$, this sharpens to
\[
 (n-m,k-j)=(2s,2r),\qquad 0\le s\le a+\ell,\quad 0\le r\le\ell.
\]
The unchanged-index coefficient is positive. These are support containments,
not assertions that every permitted entry is nonzero.
\end{theorem}
\begin{proof}
\leanok
\end{proof}

\begin{theorem}[Finite Cartesian Gram factors]
\label{thm:cartesian-gram-factors}\uses{thm:cartesian-connection-support}
\lean{MvPolynomial.CartesianDiskSystem.connectionMatrix, MvPolynomial.CartesianDiskSystem.connectionMatrix_apply, MvPolynomial.CartesianDiskSystem.integralGram_eq_connectionMatrix_transpose_mul}\leanok
Assume the source and target factor measures are supported in $[-1,1]$.
On every complete degree space, the connection matrix consists of the mixed integral coefficients of Theorem~\ref{thm:cartesian-connection-support} and factors the modified Gram matrix as $R^\top R$.
The support bounds of Theorem~\ref{thm:cartesian-connection-support} therefore apply to these finite
matrices uniformly in the truncation degree.
\end{theorem}
\begin{proof}
\leanok
\end{proof}

\begin{theorem}[Agreement with the normalized Dunkl--Xu basis]
\label{thm:dunkl-xu-cartesian-normalization}\uses{thm:disk-normalization,def:disk-outer-modification,thm:jacobi-disk-outer-shift,thm:cartesian-disk-system-existence}
\lean{MvPolynomial.dunklXuCartesianSystem, MvPolynomial.dunklXuDisk_eq_pos_smul_cartesian, MvPolynomial.normalizedDunklXuBasis_eq_cartesian, MvPolynomial.dunklXuModifiedConnection_apply}\leanok
For the symmetric Jacobi factors of the standard disk weight, every
positive-leading orthonormal Cartesian realization equals the existing normalized
Dunkl--Xu basis. Raw Jacobi-normalized disk polynomials differ by positive scalars.
Thus the support statements apply in the original total-degree and increasing
inner-index convention, without changing bases within degree blocks.
\end{theorem}
\begin{proof}
\leanok
\end{proof}

\begin{theorem}[Six- and nine-position disk connections]
\label{thm:disk-six-nine-support}\uses{thm:dunkl-xu-cartesian-normalization,thm:cartesian-connection-support}
\lean{MvPolynomial.cartesianDiskModification_X_sq, MvPolynomial.cartesianDiskModification_boundaryProduct, MvPolynomial.dunklXuXYSystem, MvPolynomial.dunklXuBoundaryProductSystem, MvPolynomial.dunklXuXYConnection_eq_zero, MvPolynomial.dunklXuBoundaryProductConnection_eq_zero, MvPolynomial.dunklXuXYConnection_diag_pos, MvPolynomial.dunklXuBoundaryProductConnection_diag_pos}\leanok
For every $\beta>-1$ and every complete finite degree section, multiplying the
disk weight by $x^2y^2$ gives normalized connection support contained in
\[
 (n-m,k-j)\in\{0,2,4\}\times\{0,2\}.
\]
Multiplication by $(1-x^2-y^2)(2-2x^2-y^2)$ gives support contained in
\[
 (n-m,k-j)\in\{0,2,4\}\times\{0,2,4\}.
\]
Both factors have positive diagonal, with at most six or nine permitted positions
per column, respectively; invalid degree/inner-index pairs are omitted.
\end{theorem}
\begin{proof}
\leanok
\end{proof}

\begin{theorem}[Raw disk multiplication and unpivoted LU support]
\label{thm:disk-modification-lu}\uses{thm:cholesky-lu,thm:dunkl-xu-cartesian-normalization,thm:cartesian-connection-support}
\lean{MvPolynomial.dunklXuMultiplicationSection, MvPolynomial.integralInner_mul_normalizedDunklXuDegreeBasis, MvPolynomial.dunklXuModifiedConnection_transpose_mul, MvPolynomial.dunklXuModifiedConnection_eq_zero_of_order, MvPolynomial.dunklXuModifiedConnection_lu_support}\leanok
Let $W$ be the exact leading section of polynomial multiplication in the raw
Dunkl--Xu basis, and let $D$ contain the positive raw basis norms. For a separated
modification as in Theorem~\ref{thm:cartesian-separated-modification} ($h,g$
nonzero nonnegative densities, $g$ even, $\deg g\le2\ell$), its normalized
connection satisfies $R^\top R=DWD^{-1}$.
If $h$ is moreover even, then under a degree/inner-index ordered enumeration
any exact unpivoted $W=LU$ with unit lower triangular $L$ and upper triangular
$U$ has $U_{ij}=0$ if and only if $R_{ij}=0$.
In particular the six- and nine-position bounds transfer to the raw upper factor.
\end{theorem}
\begin{proof}
\leanok
\end{proof}

\begin{supporting}[Reference details: Symmetric polynomial degree projections]
\begin{theorem}[Symmetric polynomial degree projections]
\label{thm:polynomial-integral-projection}\uses{thm:disk-normalization}
\lean{MvPolynomial.integralProjection, MvPolynomial.integralProjection_symmetric, MvPolynomial.integralProjection_eq_self_of_totalDegree_le, MvPolynomial.dunklXuDegreeProjection, MvPolynomial.dunklXuDegreeProjection_eq_self, MvPolynomial.dunklXuDegreeProjection_symmetric}\leanok
The finite orthonormal degree basis defines a polynomial projection by integral
coordinates. It is symmetric for the polynomial pairing and fixes every
polynomial through its degree cutoff.
\end{theorem}
\begin{proof}
\leanok
\end{proof}
\end{supporting}


\begin{theorem}[Two extra blocks for quadratic-product leading sections]
\label{thm:quadratic-product-padding}\uses{thm:polynomial-integral-projection,thm:disk-modification-lu}
\lean{MvPolynomial.compressedMul, MvPolynomial.compressedMul_sq_apply_eq, MvPolynomial.compressedCoordinateQuadratic, MvPolynomial.integralInner_compressed_quadratic_product, MvPolynomial.dunklXu_quadratic_product_padding}\leanok
Compress coordinate multiplication at degree $N+2$. For any two quadratics
$a+bx^2+cy^2$ and $d+ex^2+fy^2$, the leading raw Dunkl--Xu coefficients through
degree $N$ of their product, formed from those compressed coordinate operators,
agree exactly with full polynomial multiplication followed by leading extraction.
Equivalently, $B$ output degree blocks require $B+2$ coordinate blocks.
This applies to both preceding quartic modifications.
\end{theorem}
\begin{proof}
\leanok
\end{proof}

\rvnote{thm:quadratic-product-padding}{
It does not assert exactness of the entire quartic output in the padded space.
}



\part{Integral operators and spectral equations}\label{part:integral-operators}
\rvsummary{Volterra and one-dimensional power-law operators, from formulas to convergent finite solves.}
\chapter{Volterra operators and second-kind equations}\label{chap:volterra}
\rvsummary{Proriol assembly, physical operators and convergent second-kind solves.}
Proriol polynomials on the triangle represent Volterra kernels. We first
compute their action on coefficients, then construct operators on function
spaces and prove convergence of finite solves. The coordinate spaces and
assumptions are stated in each solve theorem.

\begin{theorem}[Proriol marginal integration]
\label{thm:proriol-volterra-marginal}\uses{thm:jacobi-values,thm:unit-jacobi-measure-characterisation,thm:jacobi-pairwise,thm:unit-jacobi-weight,thm:integrable-input,thm:proriol-triangle-basic}
\lean{Polynomial.eval_shiftedJacobi_one_sub,Polynomial.integral_shiftedJacobi_zero_zero,MvPolynomial.eval_proriol_reflected_normalized_scaled,MvPolynomial.integral_proriol_inner,MvPolynomial.integral_proriol_reflected,MvPolynomial.integral_proriol_reflected_normalized,MvPolynomial.triangleMarginal,MvPolynomial.triangleReflect,MvPolynomial.triangleMarginal_proriol,MvPolynomial.triangleMarginal_sum_proriol}\leanok
Let $D_{n,k}$ denote our DLMF-normalized unweighted Proriol polynomial on
$0<y<x<1$, and put
$Q_{n,k}(x,y)=(-1)^{n-k}D_{n,k}(1-x,y)$.
This is precisely the normalization of \cite{gutleb-olver} on $0<y<1-x$:
\[
 Q_{n,k}(x,(1-x)z)=\widetilde P_{n-k}^{(2k+1,0)}(x)
 (1-x)^k\widetilde P_k^{(0,0)}(z).
\]
The linear marginal $I_yf(x)=\int_0^{1-x}f(x,y)\,dy$ satisfies
\[
 I_yQ_{n,k}(x)=
 \begin{cases}(1-x)\widetilde P_n^{(1,0)}(x),&k=0,\\0,&k>0.\end{cases}
\]
Consequently, for every finite expansion,
\[
 I_y\!\left(\sum_{n<N}\sum_{k=0}^n c_{n,k}Q_{n,k}\right)(x)
 =(1-x)\sum_{n<N}c_{n,0}\widetilde P_n^{(1,0)}(x).
\]
The endpoint factor remains explicit; these identities include $x=1$ and
hold for all real $x$ with oriented interval integrals.

References: \cite[Section 3.1]{gutleb-olver};
\cite[Section 2.3]{gutleb-thesis}.
\end{theorem}
\begin{proof}
\leanok
\end{proof}

\begin{theorem}[Exact polynomial Volterra assembly]
\label{thm:volterra-polynomial-assembly}\uses{thm:proriol-volterra-marginal}
\lean{MvPolynomial.triangleExtendY,MvPolynomial.eval_triangleExtendY,MvPolynomial.volterraPolynomial,MvPolynomial.volterraPolynomial_apply,MvPolynomial.volterraPolynomial_add,MvPolynomial.volterraPolynomial_smul}\leanok
For a polynomial kernel $K(x,y)$ and polynomial input $u$, let
$E_yu(x,y)=u(y)$ and $\mathcal RK(x,y)=K(1-x,y)$.
Then extension, multiplication by the reflected kernel, marginal integration,
and output reflection give exactly
\[
 [I_y((\mathcal RK)E_yu)](1-x)
 =\int_0^x K(x,y)u(y)\,dy.
\]
This assembly is linear in the input and in the kernel. The formula holds
for all real $x$, including the collapsed interval at zero.


Reference: \cite[Sections 3.1--3.2]{gutleb-olver}.
\end{theorem}
\begin{proof}
\leanok
\end{proof}

\rvnote{thm:volterra-polynomial-assembly}{
The bounded integral operator and its polynomial coefficient realization are given in Theorems~\ref{thm:volterra-l2} and~\ref{thm:volterra-l2-proriol}.
}

\begin{supporting}[Reference details: Finite Proriol marginal coefficient operator and related results]
\begin{theorem}[Finite Proriol marginal coefficient operator]
\label{thm:proriol-marginal-coefficients}\uses{thm:proriol-volterra-marginal}
\lean{MvPolynomial.proriolVolterraSynthesis,MvPolynomial.proriolMarginalCoefficients,MvPolynomial.proriolMarginalCoefficients_single,MvPolynomial.proriolMarginalCoefficients_entry,MvPolynomial.proriolMarginalPolynomial,MvPolynomial.eval_proriolMarginalPolynomial}\leanok
On finitely supported coefficients indexed by $(n,k)$ with $0\le k\le n$, the
normalized Proriol marginal coefficient operator has entries
\[
 Q_{j,(n,k)}=\begin{cases}1,&k=0\text{ and }j=n,\\0,&\text{otherwise}.\end{cases}
\]
Synthesize its output in $\widetilde P_n^{(1,0)}$ and multiply by $1-X$.
The resulting univariate polynomial agrees pointwise with the analytic marginal
of the synthesized triangle polynomial. Each matrix column has at most one
nonzero; no infinite series or infinite matrix multiplication is used.

References: \cite[Section 3.1]{gutleb-olver}; \cite[equation (2.7)]{gutleb-thesis}.
\end{theorem}
\begin{proof}
\leanok
\end{proof}

\begin{theorem}[Jacobi lowering of the second parameter]
\label{thm:jacobi-beta-lowering}\uses{thm:jacobi-values}
\lean{Polynomial.jacobi_lower_beta}\leanok
Over every commutative $\Q$-algebra, for $n\ge0$,
\[
 (X+1)(P_n^{(\alpha,\beta)})'+\beta P_n^{(\alpha,\beta)}
 =(n+\beta)P_n^{(\alpha+1,\beta-1)}.
\]
No parameter factors are divided out, so exceptional parameters are included.

\end{theorem}
\begin{proof}
\leanok
\end{proof}

\rvnote{thm:jacobi-beta-lowering}{
Derived from the Jacobi derivative formula \cite[18.9.15]{dlmf} and endpoint symmetry.
}

\begin{theorem}[Volterra Jacobi extension/marginal diagonal]
\label{thm:volterra-jacobi-diagonal}\uses{thm:jacobi-beta-lowering,thm:volterra-polynomial-assembly}
\lean{Polynomial.hasDerivAt_mul_shiftedJacobi_zero_one,Polynomial.integral_shiftedJacobi_one_zero,MvPolynomial.triangleMarginal_extendY_shiftedJacobi,MvPolynomial.volterraPolynomial_one_shiftedJacobi,MvPolynomial.volterraJacobiDiagonal,MvPolynomial.volterraJacobiDiagonal_single,MvPolynomial.triangleMarginal_extendY_linearCombination}\leanok
For every $n\ge0$ and real $x$,
\[
 \int_0^x\widetilde P_n^{(1,0)}(y)\,dy
 =\frac{x}{n+1}\widetilde P_n^{(0,1)}(x).
\]
Consequently extension in $y$ followed by the simplex marginal satisfies
\[
 I_yE_y\widetilde P_n^{(1,0)}(x)
 =(1-x)\frac{(-1)^n}{n+1}\widetilde P_n^{(1,0)}(x).
\]
With the endpoint factor kept outside the coefficient operator, the diagonal is
$D_{nn}=(-1)^n/(n+1)$. This coefficient formula is proved for every finitely
supported input, with exact agreement with the analytic marginal. It includes
$n=0$ and the collapsed integration intervals.

References: \cite[Section 3.1]{gutleb-olver}; \cite[Section 2.3]{gutleb-thesis},
with the displayed one-based diagonal converted to zero-based indexing.
\end{theorem}
\begin{proof}
\leanok
\end{proof}

\begin{theorem}[Exact finite Proriol extension coordinates]
\label{thm:proriol-extension-coordinates}\uses{thm:proriol-marginal-coefficients,thm:volterra-polynomial-assembly,thm:proriol-triangle-bases}
\lean{MvPolynomial.proriolVolterraCoordinateTransport,MvPolynomial.proriolVolterraSynthesis_transport,MvPolynomial.proriolVolterraCoordinateTransport_eq_zero,MvPolynomial.proriolExtensionCoefficients,MvPolynomial.triangleReflect_extendY,MvPolynomial.triangleReflect_involutive,MvPolynomial.proriolVolterraSynthesis_injective,MvPolynomial.proriolVolterraSynthesis_extension,MvPolynomial.eval_proriolVolterraSynthesis_extension,MvPolynomial.proriolExtensionCoefficients_eq_zero,MvPolynomial.eval_proriolMarginalPolynomial_extension,MvPolynomial.proriolExtensionCoefficients_cutoff,MvPolynomial.proriolExtensionCoefficients_add,MvPolynomial.proriolExtensionCoefficients_smul}\leanok
For every polynomial $p$ of degree at most $N$, construct finitely supported
coordinates $E_Np$ in the signed, reflected Volterra Proriol basis $Q_{n,k}$.
They satisfy
\[
 \sum_{n,k}(E_Np)_{n,k}Q_{n,k}(x,y)=p(y),\qquad
 (E_Np)_{n,k}=0\quad(n>N).
\]
The synthesis identity is a polynomial identity, hence holds for all real
$x,y$. The coordinates are unique, independent of any admissible cutoff $N$,
and additive and homogeneous in $p$. Applying the existing marginal coefficient
operator, including its endpoint factor, gives exactly
$\int_0^{1-x}p(y)\,dy$.



Reference: \cite[Section 3.1]{gutleb-olver}.
\end{theorem}
\begin{proof}
\leanok
\end{proof}

\rvnote{thm:proriol-extension-coordinates}{
Explicit rational entries for shifted Jacobi inputs and their sharper total-degree-block support are given in Theorems~\ref{thm:proriol-extension-closed-form} and~\ref{thm:proriol-extension-degree-block}.

The construction uses the finite orthogonal-basis coordinate map.
}

\begin{theorem}[Explicit marginal entries and initial extension columns]
\label{thm:proriol-extension-initial-entries}\uses{thm:proriol-extension-coordinates,thm:volterra-jacobi-diagonal}
\lean{MvPolynomial.proriolMarginalCoefficients_apply,MvPolynomial.proriolMarginalCoefficients_extension_shiftedJacobi,MvPolynomial.proriolExtensionCoefficients_shiftedJacobi_inner_zero,MvPolynomial.proriolExtensionCoefficients_shiftedJacobi_zero,MvPolynomial.proriolExtensionCoefficients_shiftedJacobi_one}\leanok
Let $e^{(n)}$ be the actual Proriol extension coordinates of
$\widetilde P_n^{(1,0)}$. For every $n,m\ge0$,
\[
 e^{(n)}_{m,0}=
 \begin{cases}(-1)^n/(n+1),&m=n,\\0,&m\ne n.\end{cases}
\]
The complete first two columns are
\[
 e^{(0)}=\delta_{(0,0)},\qquad
 e^{(1)}=-\tfrac12\delta_{(1,0)}+\tfrac32\delta_{(1,1)}.
\]
Here $\delta_{(n,k)}$ is the coordinate unit vector in the signed, reflected
Volterra Proriol basis. These are equalities of the full finitely supported
vectors; all unlisted entries in the two initial columns vanish.

Reference: \cite[Section 3.1]{gutleb-olver}; \cite[Section 2.3]{gutleb-thesis}.
\end{theorem}
\begin{proof}
\leanok
\end{proof}

\rvnote{thm:proriol-extension-initial-entries}{
The formula for arbitrary positive inner degree is not asserted here.
}

\begin{theorem}[Finite polynomial moments on the Proriol triangle]
\label{thm:proriol-polynomial-moments}\uses{thm:proriol-triangle-bases}
\lean{MvPolynomial.integral_monomial_proriolMeasure_zero,MvPolynomial.proriolTriangleMoment,MvPolynomial.integral_eval_proriolMeasure_zero,MvPolynomial.integralInner_proriolMeasure_zero,MvPolynomial.integral_sq_proriolMeasure_zero,MvPolynomial.proriolDegreeBasis_repr_zero,MvPolynomial.proriolTriangleMoment_sq_proriol_pos}\leanok
For a polynomial $q(x,y)=\sum_{a,b}q_{ab}x^ay^b$, define the finite sum
\[
 \mathcal M(q)=\sum_{a,b}\frac{q_{ab}}{(b+1)(a+b+2)}.
\]
Then $\mathcal M(q)=\int_0^1\int_0^x q(x,y)\,dy\,dx$.
Consequently polynomial pairings and squared norms in the unweighted DLMF
triangle are $\mathcal M(pq)$ and $\mathcal M(p^2)$. For every degree cutoff
$N$, coordinates of a polynomial of total degree at most $N$ in the canonical
unweighted Proriol basis with parameters $(0,0,0)$ are the corresponding moment
quotient, and the squared-norm denominators of the basis members are strictly positive.
Reference: \cite[18.37.8]{dlmf} at zero weight parameters.
\end{theorem}
\begin{proof}
\leanok
\end{proof}

\begin{theorem}[Elementary unweighted Proriol norms]
\label{thm:proriol-unweighted-norm}\uses{thm:proriol-polynomial-moments,thm:jacobi-norm}
\lean{Polynomial.integral_sq_jacobi_zero_left,MvPolynomial.proriolTriangleMoment_sq_proriol}\leanok
For $\beta>-1$, the Jacobi norm with first parameter zero is
\[
 \int_{-1}^1(1+x)^\beta\bigl(P_n^{(0,\beta)}(x)\bigr)^2\,dx
 =\frac{2^{\beta+1}}{2n+\beta+1}.
\]
Writing $P_{m+k,k}$ for the unweighted DLMF Proriol polynomial, for all $m,k\ge0$,
\[
 \mathcal M(P_{m+k,k}^2)=\frac1{2(m+k+1)(2k+1)}.
\]
Reference: \cite[Table 18.3.1 and 18.37.8]{dlmf}.
\end{theorem}
\begin{proof}
\leanok
\end{proof}

\begin{theorem}[Finite-sum formulas for all Proriol extension entries]
\label{thm:proriol-extension-moments}\uses{thm:proriol-extension-coordinates,thm:proriol-unweighted-norm}
\lean{MvPolynomial.proriolVolterraCoordinateTransport_apply,MvPolynomial.proriolExtensionCoefficients_eq_moment,MvPolynomial.proriolExtensionCoefficients_eq_scaled_moment}\leanok
Let $p$ have degree at most $N$, and let $e=E_Np$ be its actual extension
coordinates in the signed, reflected Volterra basis. For $m+k\le N$,
\[
 e_{m+k,k}=(-1)^m\frac{\mathcal M(p(y)P_{m+k,k}(x,y))}
                           {\mathcal M(P_{m+k,k}^2)}
 =(-1)^m\,2(m+k+1)(2k+1)\,\mathcal M(p(y)P_{m+k,k}(x,y)).
\]
Here $P$ is the unreflected DLMF basis and $\mathcal M$ is the explicit finite
coefficient sum above. Entries above the cutoff vanish by the existing extension
theorem. Thus every entry is expressed using finite polynomial arithmetic,
with no integral or abstract basis-coordinate operation on the right-hand side.

Reference: \cite[Section 3.1]{gutleb-olver}, with orthogonal projection in the
normalization of \cite[18.37.8]{dlmf}.
\end{theorem}
\begin{proof}
\leanok
\end{proof}

\rvnote{thm:proriol-extension-moments}{
The explicit shifted-Jacobi formula and its support are Theorems~\ref{thm:proriol-extension-closed-form} and~\ref{thm:proriol-extension-degree-block}.
}

\begin{theorem}[Jacobi extension stays in one triangle degree]
\label{thm:proriol-extension-degree-block}\uses{thm:proriol-extension-moments,thm:orthogonal-lower-degree}
\lean{Polynomial.integral_mul_shiftedJacobi_one_zero_eq_zero,MvPolynomial.integralInner_extendY_monomial,MvPolynomial.integralInner_extendY_shiftedJacobi_eq_zero,MvPolynomial.proriolExtensionCoefficients_shiftedJacobi_eq_zero_of_degree_ne}\leanok
The inner-coordinate extension of $P_n^{(1,0)}(2y-1)$ is orthogonal, for
unnormalized area measure on $0<y<x<1$, to every polynomial of total degree
less than $n$. Consequently its Proriol extension coordinates vanish in every
total-degree block except $n$.
Reference: \cite[Section 3.1]{gutleb-olver}.
\end{theorem}
\begin{proof}
\leanok
\end{proof}
\end{supporting}


\begin{theorem}[Closed-form Jacobi-to-Proriol extension]
\label{thm:proriol-extension-closed-form}\uses{thm:proriol-extension-degree-block,thm:volterra-jacobi-diagonal}
\lean{Polynomial.integral_shiftedJacobi_one_zero_mul,Polynomial.integral_shiftedJacobi_one_zero_mul_zero_zero,Polynomial.integral_mul_shiftedJacobi_zero_zero,MvPolynomial.eval_proriol_outer_one,MvPolynomial.proriolExtensionCoefficients_shiftedJacobi_same_degree,MvPolynomial.proriolExtensionCoefficients_shiftedJacobi_apply}\leanok
In the signed reflected Volterra convention, every extension column is
\[
 E_{(m,k),n}=\begin{cases}
 (-1)^{n-k}\dfrac{2k+1}{n+1},&m=n,\\
 0,&m\ne n.
 \end{cases}
\]
These are the actual previously defined coordinates, including the constant
and linear columns. More generally, for $\deg q\le n$,
$\int_0^1P_n^{(1,0)}(2y-1)q(y)\,dy=q(1)/(n+1)$.
Reference: \cite[Section 3.1]{gutleb-olver}.
\end{theorem}
\begin{proof}
\leanok
\end{proof}

\begin{definition}[Finite Volterra coefficient assembly]
\label{def:volterra-coefficient-assembly}\uses{thm:proriol-extension-coordinates}
\lean{MvPolynomial.proriolVolterraSynthesis_surjective,MvPolynomial.proriolVolterraEquiv,MvPolynomial.proriolExtensionMap,MvPolynomial.proriolExtensionMap_single,MvPolynomial.proriolKernelMap,MvPolynomial.volterraProriolCoefficients}\leanok
Let $S$ be synthesis in the signed reflected Proriol basis $Q_{n,k}$.
 For finite shifted Jacobi $(1,0)$ coefficients $c$, set
\[
 Ec=S^{-1}\!\left(\sum_n c_nP_n^{(1,0)}(2y-1)\right),\qquad
 M_Kb=S^{-1}\bigl(K(1-x,y)Sb\bigr),\qquad B_K=DM_KE,
\]
where $b$ is a finite triangular coefficient vector and $D$ selects inner index zero. The $n$th column of $E$ is the coordinate vector of $P_n^{(1,0)}(2y-1)$.  The endpoint factor and output
reflection are kept outside $B_K$.
Reference: \cite[Sections 3.1--3.2]{gutleb-olver}.
\end{definition}

\begin{theorem}[Algebraic properties of Volterra coefficient assembly]
\label{thm:volterra-coefficient-assembly}\uses{def:volterra-coefficient-assembly}
\lean{MvPolynomial.proriolVolterraSynthesis_surjective,MvPolynomial.proriolVolterraEquiv,MvPolynomial.proriolKernelMap}\leanok
Synthesis $S$ is a linear equivalence between finitely supported triangular coefficients and bivariate polynomials. The kernel map $K\mapsto M_K$ is an algebra homomorphism on finitely supported vectors.
\end{theorem}
\begin{proof}
\leanok
\end{proof}

\begin{supporting}[Reference details: Exact finite Volterra coefficient action]
\begin{theorem}[Exact finite Volterra coefficient action]
\label{thm:volterra-coefficient-exact}\uses{thm:volterra-coefficient-assembly,thm:proriol-extension-initial-entries}
\lean{MvPolynomial.proriolVolterraSynthesis_kernelMap,MvPolynomial.volterraProriolCoefficients_exact,MvPolynomial.volterraProriolCoefficients_one}\leanok
For every real bivariate polynomial kernel $K$, finite coefficient vector $c$,
and real $x$, writing $u(y)=\sum_n c_nP_n^{(1,0)}(2y-1)$,
\[
 x\sum_m(B_Kc)_m P_m^{(1,0)}(1-2x)=\int_0^x K(x,y)u(y)\,dy.
\]
The integral is oriented when $x<0$. Kernel multiplication reconstructs exactly
$K(1-x,y)Sc$. For $K=1$, $B_K$ is the existing diagonal with entries
$(-1)^n/(n+1)$, including degree zero.
Reference: \cite[Sections 3.1--3.2]{gutleb-olver}.
\end{theorem}
\begin{proof}
\leanok
\end{proof}
\end{supporting}


\begin{theorem}[Volterra bandwidth and exact rectangular sections]
\label{thm:volterra-coefficient-bandwidth}\uses{thm:volterra-coefficient-exact,thm:proriol-extension-degree-block}
\lean{MvPolynomial.triangleReflect_proriolVolterraSynthesis_single,MvPolynomial.integralInner_reflect_proriolVolterraSynthesis,MvPolynomial.proriolKernelMap_entry_eq_zero,MvPolynomial.volterraProriolCoefficients_entry_eq_zero,MvPolynomial.volterraProriolSection,MvPolynomial.volterraProriolSection_eq_zero,MvPolynomial.volterraProriolSection_mulVec,MvPolynomial.volterraProriolCoefficients_cutoff}\leanok
Let $d=\operatorname{totalDegree}K$, with $d=0$ for the zero kernel.
The degree blocks of $M_K$ vanish when $|m-n|>d$, and the core Volterra matrix
satisfies $(B_K)_{mn}=0$ under the same condition. These bounds precede the
explicit endpoint factor and reflection. Inputs supported in $0\le n<N$
produce coefficients supported in $0\le m<N+d$. The corresponding
$(N+d)\times N$ section has exactly the full operator's action on the retained
rows and loses no output coefficients.
Reference: \cite[Sections 3.1--3.2]{gutleb-olver}.
\end{theorem}
\begin{proof}
\leanok
\end{proof}

\rvnote{thm:volterra-coefficient-bandwidth}{No same-size truncation of intermediate multiplication operators is asserted.}


\begin{theorem}[Operator Clenshaw for polynomial Volterra kernels]
\label{thm:volterra-operator-clenshaw}\uses{thm:operator-clenshaw,thm:volterra-coefficient-assembly}
\lean{MvPolynomial.proriolKernelMap_aeval,MvPolynomial.proriolKernelMap_X_commute,MvPolynomial.proriolKernelMap_clenshaw,MvPolynomial.volterraProriolCoefficients_clenshaw}\leanok
For $p\in\R[t]$, polynomial functional calculus obeys $M_{p(K)}=p(M_K)$, and the two coordinate
operators commute. Suppose bivariate polynomials satisfy
$p_{k+2}=a_kp_{k+1}+b_kp_k$. Applying the backward Clenshaw recurrence to
$M_{a_k},M_{b_k},M_{c_k}$ with initial operators $M_{p_0},M_{p_1}$ computes
$M_{\sum_{k<N}c_kp_k}$ exactly. Placing it between $E$ and $D$ gives
$B_{\sum_{k<N}c_kp_k}$. This permits polynomial-valued coefficients and nested
coordinate expansions. The identity is on finitely supported vectors;
no floating-point stability or finite-section recurrence is claimed.
References: \cite{clenshaw}; \cite[Sections 3.1--3.2]{gutleb-olver}.
\end{theorem}
\begin{proof}
\leanok
\end{proof}


\begin{definition}[The physical Volterra triangle]
\label{def:volterra-triangle}
\lean{Real.volterraTriangle}\leanok
The physical Volterra triangle is
$\Delta=\{(x,y)\in\mathbb R^2:0\le y\le x\le1\}$.
\end{definition}

\begin{supporting}[Formal details: Integrable Volterra row pairs and related results]
\begin{definition}[Integrable Volterra row pairs]
\label{def:volterra-integral-input}
\lean{Real.VolterraIntegralInput}\leanok
For $x\in\mathbb R$, put $J_x=(\min(0,x),\max(0,x))$.
A Volterra input at $x$ consists of two real-valued functions $k,f$ on $J_x$
whose product is Lebesgue integrable there. Neither factor is required to be
separately integrable, and no endpoint values are part of the input.
\end{definition}

\begin{theorem}[Constructing integrable Volterra rows]
\label{thm:volterra-integral-input}\uses{def:volterra-integral-input,def:volterra-triangle}
\lean{Real.VolterraIntegralInput.ofKernel,Real.intervalIntegrable_volterra,Real.VolterraIntegralInput.ofContinuous,MvPolynomial.intervalIntegrable_polynomialKernel_mul,Real.VolterraIntegralInput.ofPolynomial,Real.intervalIntegrable_continuousVolterra_mul_row,Real.VolterraIntegralInput.ofContinuousKernel,MvPolynomial.intervalIntegrable_polynomialKernel_mul_on,Real.VolterraIntegralInput.ofIntegrableOnPolynomial,Real.intervalIntegrable_locallyIntegrable,Real.intervalIntegrable_of_integrableOn_Icc,Real.intervalIntegrable_absolutelyContinuousKernel_mul,Real.VolterraIntegralInput.ofAbsolutelyContinuous}\leanok
Let $J_x=(\min(0,x),\max(0,x))$.
If $y\mapsto K(x,y)f(y)$ is integrable between $0$ and $x$, its restrictions
to $J_x$ define an integrable Volterra row.
The product is integrable in each of the following cases:
\begin{itemize}
\item $K$ is continuous on $\Delta$, $f$ is continuous on $[0,1]$, and $x\in[0,1]$;
\item $K$ is a polynomial kernel, $f$ is locally integrable on $\R$, and $x\in\R$;
\item $K$ is continuous on $\Delta$, $x\in[0,1]$, and $f$ is integrable on $[0,x]$;
\item $K$ is a polynomial kernel, $f$ is integrable on $[0,b]$, and $x\in[0,b]$;
\item the slice $K(x,\cdot)$ is absolutely continuous between $0$ and $x$,
and $f$ is integrable on that interval.
\end{itemize}
\end{theorem}
\begin{proof}
\leanok
\end{proof}
\end{supporting}


\begin{definition}[Classical Volterra integration]
\label{def:volterra-continuous-kernel}\uses{def:volterra-integral-input}
\lean{Real.volterraIntegral}\leanok
For an integrable row pair $(k,f)$ at $x$, define its oriented value by
\[
 V(k,f;x)=
 \begin{cases}
   \displaystyle\int_{J_x} k(y)f(y)\,dy,&x\ge0,\\
   \displaystyle-\int_{J_x} k(y)f(y)\,dy,&x<0.
 \end{cases}
\]
For the row $k(y)=K(x,y)$ this is denoted $V_Kf(x)$.
The integral exists by the input's product-integrability condition.
\end{definition}

\begin{supporting}[Reference details: Orientation and the zero-length row and related results]
\begin{theorem}[Orientation and the zero-length row]
\label{thm:volterra-integral-value}\uses{def:volterra-continuous-kernel}
\lean{Real.volterraIntegral_of_nonneg,Real.volterraIntegral_of_neg,Real.volterraIntegral_zero}\leanok
For a row with $K(x,\cdot)f$ integrable on
$J_x=(\min(0,x),\max(0,x))$,
\[
 V_Kf(x)=
 \begin{cases}
 \displaystyle\int_{J_x}K(x,y)f(y)\,dy,&x\ge0,\\
 \displaystyle-\int_{J_x}K(x,y)f(y)\,dy,&x<0.
 \end{cases}
\]
In particular $V_Kf(0)=0$; no endpoint values of the factors are used.
\end{theorem}
\begin{proof}
\leanok
\end{proof}

\begin{theorem}[Volterra kernel and input stability]
\label{thm:volterra-approximation-stability}\uses{thm:volterra-integral-input,def:volterra-continuous-kernel}
\lean{Real.abs_volterraIntegral_sub_le,Real.tendstoUniformly_volterraIntegral_of_error_bounds}\leanok
Let $K,L$ be continuous on $\Delta$ and $f,g$ continuous on $[0,1]$.
Their Volterra integrands are integrable for each $x\in[0,1]$.
Suppose $M,F,\delta,\eta\ge0$, with $|K|\le M$, $|f|\le F$,
$|L-K|\le\delta$ and $|g-f|\le\eta$ on the respective domains. Then
\[
 |V_Lg(x)-V_Kf(x)|\le x\bigl(\delta F+(M+\delta)\eta\bigr).
\]
For continuous approximants $L_n,g_n$ with such nonnegative uniform error bounds
$\delta_n,\eta_n\to0$, the outputs converge uniformly to $V_Kf$ on $[0,1]$.
\end{theorem}
\begin{proof}
\leanok
\end{proof}

\begin{theorem}[Approximation by finite Proriol Volterra assembly]
\label{thm:volterra-finite-approximation}\uses{thm:volterra-approximation-stability,thm:multivariate-uniform-approximation,thm:volterra-coefficient-exact}
\lean{MvPolynomial.isCompact_volterraTriangle,MvPolynomial.exists_polynomial_near_volterraKernel,MvPolynomial.volterraProriolCoefficients_error,MvPolynomial.tendstoUniformly_volterraProriolCoefficients,MvPolynomial.exists_volterraProriolCoefficients_approximation}\leanok
Every continuous real kernel on $\Delta$ admits uniform bivariate polynomial
approximations. More generally, polynomial functions uniformly approximate
continuous real functions on any compact subset of $\R^2$.
Let $P$ be a polynomial kernel and $p=\sum_n c_nJ_n$ a finite shifted-Jacobi
input, with $J_n=P_n^{(1,0)}(2x-1)$. If their kernel and input errors relative
to $K,f$ are bounded by $\delta,\eta$, then
\[
 \left|x\sum_j(B_Pc)_jJ_j(1-x)-V_Kf(x)\right|
 \le x\bigl(\delta F+(M+\delta)\eta\bigr),\qquad B_P=D M_P E.
\]
Vanishing uniform approximation errors therefore give uniform convergence of
these assembled physical outputs on all of $[0,1]$, including the endpoints.
For each continuous $K,f$ and each $\varepsilon>0$, there exist a polynomial
$P$ and finite coefficient vector $c$ for which this output error is less than
$\varepsilon$ everywhere on $[0,1]$.
Reference: \cite[Sections 3.1--3.2]{gutleb-olver}.
\end{theorem}
\begin{proof}
\leanok
\end{proof}

\rvnote{thm:volterra-finite-approximation}{These are operator-action approximation results; no coefficient-norm or inverse-solver convergence is asserted.}


\begin{theorem}[Shifted Jacobi Hilbert bases]
\label{thm:unit-jacobi-completeness}\uses{thm:unit-jacobi-measure,thm:general-hilbert-basis,thm:jacobi-positive,thm:compact-polynomial-dense,thm:unit-jacobi-measure-characterisation,thm:unit-jacobi-weight}
\lean{Polynomial.hasPositiveDefiniteIntegral_unitJacobiMeasure,Polynomial.hilbertBasisShiftedJacobi,Polynomial.hilbertBasisShiftedJacobi_apply,Polynomial.unitJacobiMeasure_zero_zero}\leanok
For $a,b>-1$, the affine Jacobi measure has strictly positive polynomial
square integrals, and its normalized shifted Jacobi polynomials form a Hilbert
basis of real $L^2$. The normalization is the reciprocal positive square root
of the squared polynomial norm. In particular,
$\mu_{0,0}^{\mathrm{unit}}=2\,dx|_{[0,1]}$; the affine measure preserves the
original mass rather than being normalized to probability.
\end{theorem}
\begin{proof}
\leanok
\end{proof}

\begin{theorem}[Bounded signed kernels between finite measure spaces]
\label{thm:bounded-kernel-l2}\uses{thm:schur-l2-marginals}
\lean{MeasureTheory.memLp_bounded_mul,MeasureTheory.boundedMulL2,MeasureTheory.boundedMulL2_ae,MeasureTheory.norm_boundedMulL2_le,MeasureTheory.boundedKernelL2Operator,MeasureTheory.norm_boundedKernelL2Operator_le,MeasureTheory.inner_boundedKernelL2Operator,MeasureTheory.integrable_boundedKernelL2_pair}\leanok
For any measure, an almost-everywhere bound $|k|\le M$, $M\ge0$, gives
bounded multiplication on real $L^2$ with norm at most $M$.
Let $\mu,\nu$ be finite measures and let $k$ be an almost-everywhere strongly
measurable real kernel with $|k|\le M$ on their product. There is a bounded
operator $T:L^2(\nu)\to L^2(\mu)$ such that
\[
 \langle g,Tf\rangle=\iint k(x,y)g(x)f(y)\,d\mu(x)d\nu(y),\qquad
 \|T\|\le M\sqrt{\mu(X)\nu(Y)}.
\]
Every displayed pairing is absolutely integrable. The kernel may have either sign.
\end{theorem}
\begin{proof}
\leanok
\end{proof}

\begin{theorem}[Integrable Volterra rows from square-integrable inputs]
\label{thm:volterra-l2-input}\uses{thm:unit-jacobi-completeness,thm:volterra-integral-input}
\lean{Real.integrableOn_coe_unitLegendreLp,Real.VolterraIntegralInput.ofUnitLegendreLp}\leanok
For $\mu=2\,dx|_{[0,1]}$, every $f\in L^2(\mu)$ is integrable on $[0,1]$.
For every kernel $K$ continuous on $\Delta$ and $x\in[0,1]$, the product $y\mapsto K(x,y)f(y)$ is integrable between $0$ and $x$.
\end{theorem}
\begin{proof}
\leanok
\end{proof}
\end{supporting}


\begin{theorem}[Volterra realization on physical L² functions]
\label{thm:volterra-l2}\uses{thm:bounded-kernel-l2,thm:volterra-l2-input,def:volterra-continuous-kernel}
\lean{Real.volterraL2Kernel,Real.measurable_volterraL2Kernel,Real.norm_volterraL2Kernel_le,Real.volterraL2Operator,Real.norm_volterraL2Operator_le,Real.inner_volterraL2Operator,Real.integrable_volterraL2_pair,Real.integral_volterraL2Kernel,Real.inner_volterraL2Operator_eq_integral}\leanok
Let $K$ be continuous on the physical triangle $\Delta$, with $|K|\le M$,
$M\ge0$, and put $\mu=2\,dx|_{[0,1]}$. Extend the kernel by zero outside
$0\le y\le x\le1$, obtaining $\widetilde K$. The bounded operator
$T_K:L^2(\mu)\to L^2(\mu)$ has norm at most $M$ and satisfies
\[
 \langle g,T_Kf\rangle
 =\tfrac12\iint\widetilde K(x,y)g(x)f(y)\,d\mu(x)d\mu(y)
 =\int_{[0,1]} g(x)V_Kf(x)\,d\mu(x).
\]
The double pairing is absolutely integrable for every $f,g\in L^2(\mu)$.
For $x\in[0,1]$, the row integral against $\mu$ is $2V_Kf(x)$.
\end{theorem}
\begin{proof}
\leanok
\end{proof}

\rvnote{thm:volterra-l2}{
For $x\in[0,1]$, the row integral against $\mu$ is $2V_Kf(x)$, explaining the factor $1/2$ in the operator construction.
}

\begin{supporting}[Reference details: Volterra coefficient-section convergence and related results]
\begin{theorem}[Volterra coefficient-section convergence]
\label{thm:volterra-l2-coefficients}\uses{thm:volterra-l2,thm:hilbert-coordinate-operator,thm:l2-finite-sections,thm:volterra-finite-approximation}
\lean{Real.volterraL2CoefficientOperator,Real.norm_volterraL2CoefficientOperator_le,Real.volterraL2CoefficientOperator_repr,Real.tendsto_volterraL2CoefficientOperator_finiteSections,Real.exists_volterraL2_bound}\leanok
Use the normalized shifted Legendre Hilbert basis of $L^2(\mu)$, with
coordinate isometry $U$. Then $A_K=UT_KU^{-1}$ satisfies
$A_K(Uf)=U(T_Kf)$, $\|A_K\|\le M$, and
\[
 P_NA_KP_Nc\longrightarrow A_Kc\quad\hbox{in }\ell^2
 \qquad(c\in\ell^2).
\]
Every continuous triangle kernel admits a finite bound $M$ and hence these
realizations.
\end{theorem}
\begin{proof}
\leanok
\end{proof}

\rvnote{thm:volterra-l2-coefficients}{
Operator-norm convergence and second-kind solves are given in Theorems~\ref{thm:volterra-compact-sections} and~\ref{thm:volterra-finite-section-solves}.
}

\begin{theorem}[Exact L² bridge to finite Proriol assembly]
\label{thm:volterra-l2-proriol}\uses{thm:volterra-l2-coefficients}
\lean{Real.volterraL2Operator_polynomial,MvPolynomial.volterraProriolOutput,MvPolynomial.eval_volterraProriolOutput,MvPolynomial.volterraL2Operator_proriol,MvPolynomial.volterraL2CoefficientOperator_proriol}\leanok
For a polynomial kernel $K$ and finite shifted-Jacobi input
$p=\sum_n c_nP_n^{(1,0)}(2x-1)$, form
\[
 q(x)=x\sum_j(B_Kc)_jP_j^{(1,0)}(1-2x),\qquad B_K=D M_K E.
\]
The polynomial $q$ evaluates to $V_Kp$ and its $L^2(\mu)$ class is exactly
$T_K[p]$. Thus $A_K(U[p])=U[q]$. More generally, any polynomial identity
$V_Kp=q$ on $[0,1]$ identifies the bounded operator's polynomial output.
The normalized shifted Legendre coordinates $U$ describe the physical output;
they are distinct from the unscaled shifted-Jacobi coordinates inside $B_K$.
\end{theorem}
\begin{proof}
\leanok
\end{proof}

\begin{theorem}[Uniform kernel approximation in operator norm]
\label{thm:volterra-kernel-norm}\uses{thm:volterra-l2,thm:volterra-finite-approximation}
\lean{Real.volterraL2Operator_sub,Real.norm_volterraL2Operator_sub_le,Real.exists_polynomial_near_volterraL2Operator,Real.tendsto_volterraL2Operator_of_error_bounds}\leanok
For continuous triangle kernels $K,L$, with finite nonnegative bounds as above,
$T_K-T_L=T_{K-L}$. If $|K-L|\le\delta$ on $\Delta$, $\delta\ge0$, then
$\|T_K-T_L\|\le\delta$. Every continuous $K$ therefore admits polynomial
kernels $p$ with $\|T_p-T_K\|<\varepsilon$, for any $\varepsilon>0$.
More generally, uniform kernel error bounds tending to zero imply convergence
of the corresponding physical operators in operator norm.
\end{theorem}
\begin{proof}
\leanok
\end{proof}

\begin{theorem}[Dominated convergence of real L² classes]
\label{thm:l2-dominated-convergence}
\lean{MeasureTheory.L2.norm_sq_eq_integral_sq,MeasureTheory.tendsto_L2_of_dominated}\leanok
For real $L^2$ representatives, $\|f\|_2^2=\int f^2\,d\mu$.
Let $f_i,g\in L^2(\mu)$, with $f_i\to g$ almost everywhere along a countably
generated filter. If $|f_i-g|\le B$ almost everywhere for every $i$, and
$B^2$ is integrable, then $[f_i]\to[g]$ in $L^2$ norm.
\end{theorem}
\begin{proof}
\leanok
\end{proof}

\begin{theorem}[Compact operators and two-sided coordinate sections]
\label{thm:compact-l2-sections}\uses{thm:l2-finite-sections,thm:compact-projection-approximation}
\lean{lp.isCompactOperator_initialL2Projection,lp.adjoint_initialL2Projection,lp.initialL2Projection_idempotent,lp.tendsto_initialL2Projection_comp,lp.isCompactOperator_adjoint_of_initialL2,lp.tendsto_comp_initialL2Projection,lp.tendsto_initialL2Sections_norm}\leanok
On real or complex $\ell^2$, the initial coordinate maps $P_N$ are compact,
self-adjoint projections. If $A$ is compact, its adjoint is compact and
\[
 P_NA\to A,\qquad AP_N\to A,\qquad P_NAP_N\to A
 \quad\hbox{in operator norm}.
\]
\end{theorem}
\begin{proof}
\leanok
\end{proof}

\begin{theorem}[Compactness from continuous Hilbert-space rows]
\label{thm:continuous-l2-rows-compact}\uses{thm:compact-l2-sections}
\lean{MeasureTheory.isCompactOperator_of_continuous_L2_rows}\leanok
Let $\mu$ be a finite measure supported almost everywhere on a compact set $S$,
and let $T:\ell^2(\mathbb R)\to L^2(\mu)$ be bounded. Suppose
$r:S\to\ell^2(\mathbb R)$ is continuous and
$(Tf)(x)=\langle r(x),f\rangle$ almost everywhere, for every $f$.
Then $T$ is compact.
\end{theorem}
\begin{proof}
\leanok
\end{proof}

\begin{theorem}[Continuous L² rows of the physical Volterra kernel]
\label{thm:volterra-l2-rows}\uses{thm:volterra-l2,thm:l2-dominated-convergence}
\lean{Real.memLp_volterraL2Kernel_row,Real.volterraL2Row,Real.volterraL2Row_ae,Real.inner_volterraL2Row,Real.norm_volterraL2Row_le,Real.memLp_inner_volterraL2Row,Real.volterraL2Operator_ae_inner_row,Real.continuousOn_volterraL2Row}\leanok
Under the continuous-kernel hypotheses above, put
$r_x=[y\mapsto\widetilde K(x,y)]\in L^2(\mu)$, $\mu=2\,dx|_{[0,1]}$.
Every row is square integrable, $\|r_x\|\le\sqrt2 M$, and
\[
 \langle r_x,f\rangle=\int\widetilde K(x,y)f(y)\,d\mu(y),\qquad
 T_Kf(x)=\tfrac12\langle r_x,f\rangle\quad\hbox{almost everywhere}.
\]
The row pairing belongs to $L^2(\mu)$ for every $f\in L^2(\mu)$, and
$x\mapsto r_x$ is continuous on the closed unit interval in $L^2$ norm.
\end{theorem}
\begin{proof}
\leanok
\end{proof}
\end{supporting}


\begin{theorem}[Compact Volterra operators and norm-convergent sections]
\label{thm:volterra-compact-sections}\uses{thm:volterra-l2-rows,thm:continuous-l2-rows-compact,thm:volterra-l2-coefficients}
\lean{Real.isCompactOperator_volterraL2Operator,Real.isCompactOperator_volterraL2CoefficientOperator,Real.tendsto_volterraL2CoefficientOperator_sections_norm}\leanok
For every continuous triangle kernel, $T_K$ and its normalized shifted Legendre
coefficient operator $A_K$ are compact. Consequently
$\|P_NA_KP_N-A_K\|\to0$.
\end{theorem}
\begin{proof}
\leanok
\end{proof}

\begin{supporting}[Reference details: Resolvents from summable or factorially bounded iterates]
\begin{theorem}[Resolvents from summable or factorially bounded iterates]
\label{thm:factorial-operator-resolvent}
\lean{ContinuousLinearMap.summable_pow_of_factorial_bound,ContinuousLinearMap.isUnit_one_sub_of_summable_pow,ContinuousLinearMap.inverse_one_sub_eq_tsum_pow,ContinuousLinearMap.isUnit_one_sub_of_factorial_bound,ContinuousLinearMap.existsUnique_sub_apply_eq_of_summable_pow}\leanok
For a bounded operator $T$ on a real or complex Banach space, summability of
$\sum_{n\ge0}T^n$ implies invertibility of $I-T$, with inverse that series.
Every right-hand side has a unique solution of $u-Tu=f$.
A bound $\|T^n\|\le C M^n/n!$ for all $n$ suffices for summability.
No assumption $\|T\|<1$ is required.
\end{theorem}
\begin{proof}
\leanok
\end{proof}
\end{supporting}


\begin{theorem}[Factorial Volterra iterates and second-kind invertibility]
\label{thm:volterra-factorial-inverse}\uses{thm:volterra-l2-rows,thm:factorial-operator-resolvent}
\lean{Real.volterraL2Operator_ae_bound_of_power,Real.volterraL2Operator_pow_succ_ae_le,Real.norm_volterraL2Operator_pow_succ_le,Real.summable_pow_volterraL2Operator,Real.isUnit_one_sub_volterraL2Operator,Real.existsUnique_volterraL2_solution}\leanok
For $|K|\le M$ on $\Delta$, $M\ge0$, and $n\ge0$,
\[
 |T_K^{n+1}f(x)|\le \frac{\|f\|_{L^2(\mu)}}{\sqrt2}
                  \frac{M^{n+1}x^n}{n!}\quad\hbox{almost everywhere},
 \qquad \|T_K^{n+1}\|\le\frac{M^{n+1}}{n!}.
\]
More generally, an almost-everywhere majorant $|f(x)|\le Cx^n$, $C\ge0$,
gives $|T_Kf(x)|\le MCx^{n+1}/(n+1)$ almost everywhere.
The series of powers is norm summable; $I-T_K$ is invertible and
$u-T_Ku=f$ has exactly one $L^2(\mu)$ solution for each $f$.
\end{theorem}
\begin{proof}
\leanok
\end{proof}

\begin{supporting}[Reference details: Invertibility in normalized Volterra coordinates and related results]
\begin{theorem}[Invertibility in normalized Volterra coordinates]
\label{thm:volterra-coefficient-invertible}\uses{thm:volterra-factorial-inverse,thm:volterra-l2-coefficients}
\lean{Real.isUnit_one_sub_volterraL2CoefficientOperator}\leanok
For a continuous kernel on the closed Volterra triangle, the normalized
coefficient operator $I-A_K$ is invertible.
\end{theorem}
\begin{proof}
\leanok
\end{proof}

\begin{definition}[Second-kind Volterra coefficient inverse]
\label{def:volterra-coefficient-inverse}\uses{thm:volterra-coefficient-invertible}
\lean{Real.volterraCoefficientInverse}\leanok
For a continuous kernel $K$ on the closed Volterra triangle, define
$R_K=(I-A_K)^{-1}$ on $\ell^2$.
\end{definition}

\begin{theorem}[Characterization of the Volterra coefficient inverse]
\label{thm:volterra-coefficient-inverse}\uses{def:volterra-coefficient-inverse}
\lean{Real.volterraCoefficientInverse_mul}\leanok
The coefficient inverse satisfies $R_K(I-A_K)=I$ and $(I-A_K)R_K=I$.
\end{theorem}
\begin{proof}
\leanok
\end{proof}

\begin{definition}[Lifted second-kind Volterra sections]
\label{def:volterra-coefficient-section}\uses{thm:compact-finite-section-convergence,thm:volterra-l2-coefficients}
\lean{Real.volterraCoefficientSection}\leanok
For a continuous kernel $K$ and a natural cutoff $N$, put
$B_{K,N}=I-P_NA_K$ on normalized coefficient space.
\end{definition}

\begin{theorem}[Action of lifted Volterra sections]
\label{thm:volterra-coefficient-section}\uses{def:volterra-coefficient-section}
\lean{Real.volterraCoefficientSection_apply}\leanok
For every $u\in\ell^2$, $B_{K,N}u=u-P_NA_Ku$.
\end{theorem}
\begin{proof}
\leanok
\end{proof}
\end{supporting}


\begin{theorem}[Stable normalized finite-section Volterra solves]
\label{thm:volterra-finite-section-solves}\uses{def:volterra-coefficient-section,def:volterra-coefficient-inverse,thm:volterra-compact-sections}
\lean{Real.volterraCoefficientSection_filter_neBot,Real.eventually_stable_volterraCoefficientSection,Real.eventually_existsUnique_volterra_projected_solution,Real.eventually_norm_volterraCoefficientSection_error_le,Real.tendsto_volterraCoefficientSection_solution}\leanok
Let $A=A_K$ and $B_N=I-P_NA$. The operator $I-A$ is invertible.
For sufficiently large $N$, $B_N$ is invertible with
$\|B_N^{-1}\|\le\|(I-A)^{-1}\|+1$.
For every $f\in\ell^2$, the finite system has exactly one solution
\[
 u_N\in\operatorname{ran}P_N,\qquad P_N(I-A)u_N=P_Nf.
\]
For every cutoff $N$ at which $B_N$ is invertible, the solution is $u_N=B_N^{-1}P_Nf$. Such cutoffs include every sufficiently large $N$. As $N\to\infty$ through these cutoffs, $u_N$ converges in $\ell^2$ to
$u=(I-A)^{-1}f$. Eventually, for every exact solution $u$,
\[
 \|B_N^{-1}P_N(I-A)u-u\|
 \le (\|(I-A)^{-1}\|+1)\|P_Nu-u\|.
\]
On $\operatorname{ran}P_N$, the projected equation is precisely the finite
matrix equation with $I-P_NAP_N$. All coordinates here are normalized shifted
Legendre coordinates of the physical output.
\end{theorem}
\begin{proof}
\leanok
\end{proof}

\begin{supporting}[Reference details: Algebraic and normalized shifted Jacobi coordinates and related results]
\begin{theorem}[Algebraic and normalized shifted Jacobi coordinates]
\label{thm:unit-jacobi-coordinates}\uses{thm:unit-jacobi-completeness}
\lean{Polynomial.shiftedJacobiBasis,Polynomial.shiftedJacobiBasis_apply,Polynomial.linearCombination_shiftedJacobi_repr,Polynomial.normalizedShiftedJacobi,Polynomial.polynomialToL2_normalizedShiftedJacobi}\leanok
For $\alpha,\beta>-1$, the shifted Jacobi polynomials form an algebraic
basis of $\mathbb R[x]$; synthesis of a polynomial's finite coordinates
recovers that polynomial. Write $h_j$ for its squared norm in the affine
Jacobi measure. The polynomial $\widehat P_j=h_j^{-1/2}P_j$ represents
exactly the $j$th vector of the normalized shifted Jacobi Hilbert basis.
\end{theorem}
\begin{proof}
\leanok
\end{proof}

\begin{theorem}[Finite matrices and projected coefficient equations]
\label{thm:l2-finite-matrices}\uses{thm:l2-finite-sections}
\lean{lp.initialL2Vector,lp.initialL2Vector_apply,lp.initialL2Vector_restrict,lp.initialL2Projection_vector,lp.initialL2Matrix,lp.initialL2Matrix_mulVec,lp.initialL2Matrix_equation_iff,lp.existsUnique_initialL2Matrix_solution}\leanok
Let $J_N:\mathbb K^N\to\ell^2$ be zero extension, for $\mathbb K=\mathbb R$
or $\mathbb C$. Restriction followed by $J_N$ is $P_N$, and $P_NJ_N=J_N$.
For bounded $T$, define $G_N(i,j)=(Te_j)_i$, $0\le i,j<N$. Then
\[
 (G_Nv)_i=(TJ_Nv)_i,\qquad
 (I-G_N)v=(f_i)_{i<N}
 \iff P_N(J_Nv-TJ_Nv)=P_Nf.
\]
A unique solution of the projected equation in $\operatorname{ran}P_N$
gives a unique finite matrix solution.
\end{theorem}
\begin{proof}
\leanok
\end{proof}

\begin{theorem}[Finite sections with norm-consistent operator approximations]
\label{thm:approximate-finite-sections}\uses{thm:compact-finite-section-convergence}
\lean{ContinuousLinearMap.approximateLiftedFiniteSection_filter_neBot,ContinuousLinearMap.tendsto_approximateLiftedFiniteSection,ContinuousLinearMap.approximateLiftedFiniteSection_error_eq,ContinuousLinearMap.norm_approximateLiftedFiniteSection_error_le,ContinuousLinearMap.eventually_stable_approximateLiftedFiniteSection}\leanok
Let $P_n$ be contractions on a real or complex Banach space, converging
strongly to $I$, let $K$ be compact, and suppose $L_n\to K$ in operator norm.
Then $B_n=I+P_nL_n\to I+K$ in operator norm. If $I+K$ is invertible,
then eventually $B_n$ is invertible and
$\|B_n^{-1}\|\le\|(I+K)^{-1}\|+1$.
For any contraction $P$ and bounded $K,L$, whenever $B=I+PL$ is invertible,
\[
 B^{-1}P(I+K)u-u=B^{-1}\bigl(Pu-u+P(K-L)u\bigr),
\]
\[
 \|B^{-1}P(I+K)u-u\|
 \le\|B^{-1}\|\bigl(\|Pu-u\|+\|K-L\|\|u\|\bigr).
\]
The fixed-section inverse $(I+PL)^{-1}$ is defined for operators $L$ such that $I+PL$ is invertible. The indices $n$ for which $B_n$ is invertible include every sufficiently large integer. The error identity and estimate do not require compactness.
\end{theorem}
\begin{proof}
\leanok
\end{proof}

\begin{definition}[Volterra sections with varying kernels]
\label{def:approximate-volterra-coefficient-section}\uses{def:volterra-coefficient-section}
\lean{Real.approximateVolterraCoefficientSection}\leanok
For continuous kernels $K_n$ on the closed Volterra triangle, define
$B_n=I-P_nA_{K_n}$ on normalized coefficient space.
\end{definition}

\begin{theorem}[Action of varying-kernel Volterra sections]
\label{thm:approximate-volterra-coefficient-section}\uses{def:approximate-volterra-coefficient-section}
\lean{Real.approximateVolterraCoefficientSection_apply}\leanok
For every $u\in\ell^2$, $B_nu=u-P_nA_{K_n}u$.
\end{theorem}
\begin{proof}
\leanok
\end{proof}

\begin{theorem}[Volterra solves with simultaneous kernel approximation]
\label{thm:volterra-approximate-solves}\uses{thm:approximate-finite-sections,thm:volterra-kernel-norm,def:approximate-volterra-coefficient-section,thm:volterra-compact-sections,def:volterra-coefficient-inverse}
\lean{Real.approximateVolterraSection_filter_neBot,Real.norm_volterraL2CoefficientOperator_sub_le,Real.tendsto_volterraL2CoefficientOperator_of_error_bounds,Real.eventually_stable_approximateVolterraSection,Real.eventually_existsUnique_approximateVolterra_solution,Real.eventually_norm_approximateVolterra_error_le,Real.tendsto_approximateVolterra_solution}\leanok
Let $K,K_n$ be continuous on the closed Volterra triangle, with finite
nonnegative uniform bounds, and suppose
$|K_n-K|\le\varepsilon_n$ there, where $\varepsilon_n\ge0$ and
$\varepsilon_n\to0$. In normalized physical coordinates,
$\|A_{K_n}-A_K\|\le\varepsilon_n$. Put $B_n=I-P_nA_{K_n}$ and
$C=\|(I-A_K)^{-1}\|+1$. Eventually $B_n$ is invertible with
$\|B_n^{-1}\|\le C$, and for every $f\in\ell^2$ the projected system
has a unique solution. The invertible-cutoff domain contains a tail, and along its induced nontrivial filter,
\[
 B_n^{-1}P_nf\longrightarrow (I-A_K)^{-1}f,
\]
and eventually, for every $u\in\ell^2$,
\[
 \|B_n^{-1}P_n(I-A_K)u-u\|
 \le C\bigl(\|P_nu-u\|+\varepsilon_n\|u\|\bigr).
\]
No smallness restriction on $K$ or common bound on all $K_n$ is required.
\end{theorem}
\begin{proof}
\leanok
\end{proof}

\begin{theorem}[Normalized Galerkin matrices from finite Proriol assembly]
\label{thm:volterra-proriol-matrix}\uses{thm:unit-jacobi-coordinates,thm:l2-finite-matrices,thm:volterra-l2-proriol,thm:volterra-finite-section-solves}
\lean{MvPolynomial.volterraNormalizedInput,MvPolynomial.volterraProriolMatrix,MvPolynomial.volterraProriolMatrix_eq,MvPolynomial.volterraProriolMatrix_mulVec,MvPolynomial.volterraProriolMatrix_equation_iff,MvPolynomial.eventually_existsUnique_volterraProriolMatrix_solution}\leanok
For a polynomial kernel $p$, let $c_j$ be the finite shifted-Jacobi $(1,0)$
coordinates of the normalized shifted Legendre polynomial $\widehat P_j$.
Let $q_j$ be the physical polynomial reconstructed from $D M_p E c_j$,
including reflection and multiplication by $x$. With $\mu=2\,dx$ on $[0,1]$,
define the actual $N\times N$ matrix
\[
 G_{p,N}(i,j)=\int\widehat P_i(x)q_j(x)\,d\mu(x).
\]
Then $G_{p,N}(i,j)=(A_pe_j)_i$; its multiplication is exactly the finite
restriction of $A_pJ_Nv$, and
\[
 (I-G_{p,N})v=(f_i)_{i<N}
 \iff P_N(J_Nv-A_pJ_Nv)=P_Nf.
\]
For all sufficiently large $N$, this finite matrix equation has exactly
one solution for every coefficient right-hand side.
\end{theorem}
\begin{proof}
\leanok
\end{proof}

\begin{definition}[Classical assembled Proriol solutions]
\label{def:volterra-proriol-matrix-solution}\uses{def:approximate-volterra-coefficient-section}
\lean{MvPolynomial.volterraProriolMatrixSolution}\leanok
For polynomial kernels $p_n$, let $B_n=I-P_nA_{p_n}$. For an invertible
cutoff $n$ and $f\in\ell^2$, define $v_n$ by restricting $B_n^{-1}P_nf$
to its first $n$ coordinates.
\end{definition}

\begin{theorem}[Characterization of assembled Proriol solutions]
\label{thm:volterra-proriol-matrix-solution}\uses{def:volterra-proriol-matrix-solution,thm:volterra-proriol-matrix}
\lean{MvPolynomial.volterraProriolMatrixSolution_spec}\leanok
At every invertible cutoff, $J_nv_n=B_n^{-1}P_nf$ and
$(I-G_{p_n,n})v_n=(f_i)_{i<n}$.
\end{theorem}
\begin{proof}
\leanok
\end{proof}
\end{supporting}


\begin{theorem}[Convergent assembled Proriol solves for continuous kernels]
\label{thm:volterra-proriol-matrix-solves}\uses{thm:volterra-approximate-solves,thm:volterra-proriol-matrix-solution}
\lean{MvPolynomial.eventually_existsUnique_approximateProriolMatrix_solution,MvPolynomial.eventually_volterraProriolMatrixSolution_spec,MvPolynomial.tendsto_volterraProriolMatrixSolution,MvPolynomial.eventually_norm_volterraProriolMatrixSolution_error_le,MvPolynomial.exists_volterraPolynomialKernel_sequence}\leanok
Let $K$ be continuous on the closed Volterra triangle and let $p_n$ be
polynomial kernels with finite nonnegative bounds and uniform errors
$|p_n-K|\le\varepsilon_n\to0$, $\varepsilon_n\ge0$.
For all sufficiently large $n$, the assembled system
\[
 (I-G_{p_n,n})v_n=(f_i)_{i<n}
\]
has a unique solution. For every cutoff $n$ at which $B_n=I-P_nA_{p_n}$ is invertible, define $v_n$ as in Definition~\ref{def:volterra-proriol-matrix-solution}; its zero extension equals the lifted solution and solves the matrix
equation at every such cutoff. These cutoffs include every sufficiently large $n$. For every $f\in\ell^2$, with $u=(I-A_K)^{-1}f$, the zero extensions converge to $u$ as $n\to\infty$ through these cutoffs, and eventually
\[
 \|J_nv_n-u\|\le (\|(I-A_K)^{-1}\|+1)
       \bigl(\|P_nu-u\|+\varepsilon_n\|u\|\bigr).
\]
Every continuous $K$ admits such a sequence with error at most $1/(n+1)$
and convergent solves for every $f$. No relation between the polynomial
kernel degrees and section sizes is assumed; the columns use the full
finite polynomial assembly before the final row restriction.
\end{theorem}
\begin{proof}
\leanok
\end{proof}

\chapter{One-dimensional power-law operators}\label{chap:power-law}
\rvsummary{Power-law formulas, weighted boundedness, coefficient operators and solves.}
This chapter studies power-law integral operators and their spectral
representations. It covers exact polynomial formulas and approximation
results, with separate treatment of the compactness range and its critical
boundary. Equilibrium minimization is outside its scope.

\begin{supporting}[Supporting results: integrals of real powers]
\begin{theorem}[Integrability and integral of a real power]
\label{thm:mathlib-power-integral}
\lean{Real.integrable_rpow_openInterval,Real.rpow_openInterval_base_pos,Real.integral_rpow_openInterval,Real.integral_natPow_classical}\leanok
For $\alpha>-1$ and $a,b\ge0$, $t^\alpha$ is integrable on the open interval
between $a$ and $b$, all its bases there are positive, and its oriented integral is
$\int_a^b t^\alpha\,dt=(b^{\alpha+1}-a^{\alpha+1})/(\alpha+1)$.
Here orientation means multiplying the integral on $(\min(a,b),\max(a,b))$
by $1$ when $a\le b$ and by $-1$ otherwise.
For every integer $n\ge0$ the polynomial $t^n$ is integrable on arbitrary real
intervals, with $\int_a^b t^n\,dt=(b^{n+1}-a^{n+1})/(n+1)$.
\end{theorem}
\begin{proof}
\leanok
\end{proof}
\end{supporting}

\begin{definition}[Classical power-law integral]
\label{def:power-law-integral}
\lean{Real.PowerLawDomain,Real.PowerLawInput,Real.powerLawIntegral}\leanok
For $\alpha,x\in\R$, let $D_x=\{y\in(-1,1):y\ne x\}$.
An input at $(\alpha,x)$ is a function $f:D_x\to\R$ for which
$y\mapsto |x-y|^\alpha f(y)$ is integrable. Define
\[Q_\alpha f(x)=\int_{D_x}|x-y|^\alpha f(y)\,dy.\]
The parameter and input domain allow every existing weighted integral,
without assigning a kernel value on the diagonal.
Reference: \cite[Section 2]{gutleb-carrillo-olver}.
\end{definition}
\begin{supporting}[Reference details: Characterisation of the power-law integral and related results]
\begin{theorem}[Characterisation of the power-law integral]
\label{thm:power-law-integral}\uses{def:power-law-integral}
\lean{Real.powerLawDomain_base_pos,Real.powerLawIntegral_spec}\leanok
For $f$ in the power-law domain at exponent $\alpha$ and point $x$,
$|x-y|>0$ for $y\in D_x=(-1,1)\setminus\{x\}$, the function
$y\mapsto |x-y|^\alpha f(y)$ is integrable on $D_x$, and
\[
 Q_\alpha f(x)=\int_{D_x}|x-y|^\alpha f(y)\,dy.
\]
\end{theorem}
\begin{proof}
\leanok
\end{proof}

\begin{definition}[Classical left power-law integral]
\label{def:power-law-left}
\lean{Real.PowerLawLeftInput,Real.powerLawLeft}\leanok
For $\alpha\in\R$ and $x\ge-1$, an input is a real-valued function on $(-1,x)$
with integrable weighted product $(x-y)^\alpha f(y)$. Define
\[L_\alpha f(x)=\int_{(-1,x)}(x-y)^\alpha f(y)\,dy.\]
The collapsed interval at $x=-1$ is empty.
\end{definition}
\begin{theorem}[Characterisation of the left power-law integral]
\label{thm:power-law-left}\uses{def:power-law-left}
\lean{Real.powerLawLeft_spec}\leanok
Every kernel base is positive, every input has an integrable weighted product,
and $L_\alpha f(x)$ equals the displayed integral.
\end{theorem}
\begin{proof}
\leanok
\end{proof}

\begin{definition}[Classical right power-law integral]
\label{def:power-law-right}
\lean{Real.PowerLawRightInput,Real.powerLawRight}\leanok
For $\alpha\in\R$ and $x\le1$, an input is a real-valued function on $(x,1)$
with integrable weighted product $(y-x)^\alpha f(y)$. Define
\[R_\alpha f(x)=\int_{(x,1)}(y-x)^\alpha f(y)\,dy.\]
The collapsed interval at $x=1$ is empty.
\end{definition}
\begin{theorem}[Characterisation of the right power-law integral]
\label{thm:power-law-right}\uses{def:power-law-right}
\lean{Real.powerLawRight_spec}\leanok
Every kernel base is positive, every input has an integrable weighted product,
and $R_\alpha f(x)$ equals the displayed integral.
\end{theorem}
\begin{proof}
\leanok
\end{proof}
\end{supporting}


\begin{theorem}[Power-law decomposition and constant input]
\label{thm:power-law-basic}\uses{thm:mathlib-power-integral,thm:power-law-integral,thm:power-law-left,thm:power-law-right}
\lean{Real.PowerLawInput.one,Real.PowerLawLeftInput.one,Real.PowerLawRightInput.one,Real.PowerLawInput.ofContinuous,Real.PowerLawInput.ofContinuous_apply,Real.PowerLawInput.left,Real.PowerLawInput.right,Real.PowerLawInput.reflect,Real.powerLawIntegral_eq_left_add_right,Real.powerLawIntegral_one,Real.powerLawLeft_one,Real.powerLawRight_one,Real.powerLawIntegral_neg}\leanok
For $x\in[-1,1]$, every input in Definition~\ref{def:power-law-integral}
has integrable left and right restrictions, and its integral equals
\[
 \int_{-1}^x(x-y)^\alpha f(y)\,dy+
 \int_x^1(y-x)^\alpha f(y)\,dy.
\]
For $\alpha>-1$ the kernel is integrable, including at $x=\pm1$; multiplying by
an input continuous on $[-1,1]$ preserves integrability. In particular,
\[
 Q_\alpha1(x)=\frac{(1+x)^{\alpha+1}+(1-x)^{\alpha+1}}{\alpha+1}.
\]
Reflection preserves the full weighted-integrability domain and obeys
$Q_\alpha(f\circ(-\mathrm{id}))(-x)=Q_\alpha f(x)$.
This is the unweighted degree-zero base case, not the general weighted
hypergeometric formula.

Reference: \cite[Section 2]{gutleb-carrillo-olver}.
\end{theorem}
\begin{proof}
\leanok
\end{proof}

\begin{supporting}[Reference details: Gegenbauer-weighted functions for power-law integration and related results]
\begin{definition}[Gegenbauer-weighted functions for power-law integration]
\label{def:power-law-gegenbauer-inputs}\uses{thm:power-law-basic,def:gegenbauer,thm:jacobi-integrable,def:gegenbauer-measure,thm:jacobi-weight-characterisation}
\lean{Real.PowerLawEndpointParameters,Real.PowerLawInput.gegenbauerContinuous,Real.PowerLawInput.gegenbauerWeight,Real.PowerLawInput.gegenbauerMode,Real.PowerLawInput.gegenbauerEndpoint}\leanok
Write $w_\lambda(y)=(1-y^2)^{\lambda-1/2}$ on $(-1,1)$.
For $\alpha>-1$, $\lambda>-1/2$, $x\in(-1,1)$ and
$f\in C([-1,1],\R)$, define $g(y)=w_\lambda(y)f(y)$ on $D_x$
as a function in the domain of $u\mapsto Q_\alpha u(x)$.
Taking $f=1$ gives the weighted constant amplitude; taking
$f=C_n^{(\lambda)}$ for $n\in\mathbb N$ gives the weighted polynomial modes.
At $x=1$, define $g(y)=w_\lambda(y)$ on $D_1$ under the endpoint conditions
$\lambda>-1/2$ and $\alpha+\lambda>-1/2$; no separate lower bound on $\alpha$
is imposed. Membership in these domains requires the product
$|x-y|^\alpha g(y)$ to be integrable on $D_x$.
\end{definition}
\begin{theorem}[Characterisation of the Gegenbauer-weighted functions]
\label{thm:power-law-gegenbauer-inputs}\uses{def:power-law-gegenbauer-inputs,thm:gegenbauer-low-degree,thm:gegenbauer-measure-characterisation}
\lean{Real.PowerLawInput.gegenbauerContinuous_apply,Real.PowerLawInput.gegenbauerWeight_apply,Real.PowerLawInput.gegenbauerMode_apply,Real.PowerLawInput.gegenbauerEndpoint_apply,Real.PowerLawInput.gegenbauerEndpoint_reflect_apply,Real.PowerLawInput.gegenbauerMode_zero}\leanok
Let $w_\lambda(y)=(1-y^2)^{\lambda-1/2}$ on $(-1,1)$.
For $\alpha>-1$, $\lambda>-1/2$, $x\in(-1,1)$ and
$f\in C([-1,1],\R)$, the weighted function in
Definition~\ref{def:power-law-gegenbauer-inputs} has value $w_\lambda(y)f(y)$
for every $y\in D_x$. The constant-amplitude case is $w_\lambda(y)$,
and the degree-$n$ case is $w_\lambda(y)C_n^{(\lambda)}(y)$.
At $x=1$, the constant-amplitude case has the same values on $D_1$ under
$\lambda>-1/2$ and $\alpha+\lambda>-1/2$.
Reflection gives $w_\lambda(y)$ on $D_{-1}$ with its weighted integrability
preserved. At degree zero, $w_\lambda C_0^{(\lambda)}=w_\lambda$.
\end{theorem}
\begin{proof}
\leanok
\end{proof}

\begin{theorem}[Gegenbauer-weight power-law endpoint values]
\label{thm:power-law-gegenbauer-endpoints}\uses{thm:jacobi-mass,def:power-law-gegenbauer-inputs}
\lean{Real.powerLawIntegral_gegenbauerEndpoint,Real.powerLawIntegral_gegenbauerEndpoint_neg}\leanok
Let $w_\lambda(y)=(1-y^2)^{\lambda-1/2}$. If $\lambda>-1/2$ and
$\alpha+\lambda>-1/2$, the right-endpoint integrand is integrable and
\[
 Q_\alpha w_\lambda(\pm1)=
 \frac{2^{\alpha+2\lambda}\Gamma(\alpha+\lambda+1/2)
       \Gamma(\lambda+1/2)}{\Gamma(\alpha+2\lambda+1)}.
\]


References: the Jacobi Beta integral, \cite[5.12.1]{dlmf}, and the weighted operator
convention of \cite[Section 2]{gutleb-carrillo-olver}.
\end{theorem}
\begin{proof}
\leanok
\end{proof}

\rvnote{thm:power-law-gegenbauer-endpoints}{The endpoint condition concerns the combined weight and kernel exponent. Interior integrability has different hypotheses; see Theorem~\ref{thm:power-law-weighted-interior}.}


\begin{theorem}[Weighted interior power-law integrability]
\label{thm:power-law-weighted-interior}\uses{def:power-law-gegenbauer-inputs}
\lean{Real.continuous_gegenbauerWeightOn,Real.integrable_powerLaw_gegenbauerContinuous,Real.powerLawIntegral_gegenbauerContinuous_split}\leanok
If $\alpha>-1$, $\lambda>-1/2$, and $-1<x<1$, then
\[
 y\longmapsto |x-y|^\alpha (1-y^2)^{\lambda-1/2}f(y)
\]
is integrable on $D_x$ whenever $f$ is continuous on $[-1,1]$.
In particular every Gegenbauer mode $f=C_n^{(\lambda)}$ is admissible, including
the degenerate parameter $\lambda=0$. These are interior hypotheses; the
endpoint values have the separate combined-exponent condition above.
The left/right fractional-integral decomposition consequently holds for these
weighted inputs without additional integrability assumptions.

Reference: \cite[Section 2]{gutleb-carrillo-olver}.
\end{theorem}
\begin{proof}
\leanok
\end{proof}

\begin{theorem}[First weighted Gegenbauer power-law mode]
\label{thm:power-law-first-mode}\uses{thm:gegenbauer-low-degree,def:power-law-gegenbauer-inputs}
\lean{Real.powerLawLeft_gegenbauer_oneOn,Real.powerLawRight_gegenbauer_oneOn,Real.powerLawIntegral_gegenbauer_oneOn}\leanok
Let $\alpha>-1$, $\lambda>-1/2$, $-1<x<1$, and
$w_\lambda(y)=(1-y^2)^{\lambda-1/2}$. Write $L_\alpha$ and $R_\alpha$
for the unnormalized left and right power-law integrals. Then
\begin{align*}
 L_\alpha[w_\lambda C_1^{(\lambda)}](x)
 &=2\lambda\bigl(xL_\alpha[w_\lambda](x)-L_{\alpha+1}[w_\lambda](x)\bigr),\\
 R_\alpha[w_\lambda C_1^{(\lambda)}](x)
 &=2\lambda\bigl(xR_\alpha[w_\lambda](x)+R_{\alpha+1}[w_\lambda](x)\bigr),\\
 Q_\alpha[w_\lambda C_1^{(\lambda)}](x)
 &=2\lambda\bigl(xQ_\alpha[w_\lambda](x)+R_{\alpha+1}[w_\lambda](x)
       -L_{\alpha+1}[w_\lambda](x)\bigr).
\end{align*}
These reductions include $\lambda=0$.

Reference: \cite[Lemma 2.3]{gutleb-carrillo-olver}.
\end{theorem}
\begin{proof}
\leanok
\end{proof}

\rvnote{thm:power-law-first-mode}{
The degree-zero hypergeometric formula is Theorem~\ref{thm:power-law-hypergeometric-symmetric}.
}

\begin{definition}[Classical power-law Möbius map]
\label{def:power-law-mobius}
\lean{Real.PowerLawMobiusDomain,Real.powerLawMobius,Real.PowerLawMobiusDomain.ofInterval}\leanok
For $x\in\R$, the domain of $\phi_x$ is $\{t\in\R:1-xt\ne0\}$, and
$\phi_x(t)=(x-t)/(1-xt)$.
The rational map also retains all admissible real arguments outside that interval.
\end{definition}
\begin{theorem}[Characterisation of the Möbius map]
\label{thm:power-law-mobius-domain}\uses{def:power-law-mobius,def:open-domain-derivative}
\lean{Real.powerLawMobius_apply,Real.powerLawMobius_den_posOn,Real.powerLawMobius_memOn,Real.powerLawMobius_involutiveOn,Real.powerLawMobius_rangeOn,Real.hasDerivAtOnOpen_powerLawMobius,Real.one_sub_powerLawMobius_sqOn,Real.sub_powerLawMobiusOn}\leanok
The map has the stated rational value throughout its domain. For $x,t\in(-1,1)$,
its denominator is positive, it maps $(-1,1)$ onto itself, and
$\phi_x(\phi_x(t))=t$. Its derivative there is
$(x^2-1)/(1-xt)^2$. Moreover,
\[
 x-\phi_x(t)=\frac{t(1-x^2)}{1-xt},\qquad
 1-\phi_x(t)^2=\frac{(1-x^2)(1-t^2)}{(1-xt)^2}.
\]
\end{theorem}
\begin{proof}
\leanok
\end{proof}

\begin{definition}[Integrable Möbius pullback]
\label{def:power-law-mobius-pullback}\uses{thm:power-law-mobius-domain,def:integrable-input}
\lean{Real.powerLawMobiusPullback}\leanok
For $x\in(-1,1)$ and an integrable real function $f$ on $(-1,1)$, define
its Jacobian-weighted pullback as the integrable input
\[
 t\longmapsto\frac{1-x^2}{(1-xt)^2}f(\phi_x(t)).
\]
\end{definition}
\begin{theorem}[Characterisation of the integrable Möbius pullback]
\label{thm:power-law-mobius-pullback}\uses{def:power-law-mobius-pullback}
\lean{Real.powerLawMobiusPullback_apply,Real.integral_comp_powerLawMobiusOn}\leanok
For $x\in(-1,1)$ and integrable $f$ on $(-1,1)$, the pullback is $t\mapsto(1-x^2)(1-xt)^{-2}f(\phi_x(t))$, and
$\int_{(-1,1)}f=\int_{(-1,1)}(1-x^2)(1-xt)^{-2}f\circ\phi_x$.
This applies to every integrable function $f$, including functions with an integrable
interior singularity; no continuity is required.
\end{theorem}
\begin{proof}
\leanok
\end{proof}

\begin{definition}[Weight after the fixed-singularity substitution]
\label{def:power-law-mobius-weight}\uses{def:power-law-gegenbauer-inputs}
\lean{Real.PowerLawInput.mobiusWeight}\leanok
For $\alpha>-1$, $\lambda>-1/2$ and $x\in(-1,1)$, define on
$D_0=(-1,1)\setminus\{0\}$ the function
\[
 t\longmapsto \frac{w_\lambda(t)}{(1-xt)^{\alpha+2\lambda+1}},
\]
as an element of the domain of $u\mapsto Q_\alpha u(0)$; its product with $|t|^\alpha$ is integrable.
All bases of real powers in this expression are positive on $D_0$.
\end{definition}
\begin{theorem}[Characterisation of the weight after substitution]
\label{thm:power-law-mobius-weight}\uses{def:power-law-mobius-weight}
\lean{Real.PowerLawInput.mobiusWeight_apply}\leanok
For $\alpha>-1$, $\lambda>-1/2$ and $x\in(-1,1)$, the function
\[
 g(t)=\frac{w_\lambda(t)}{(1-xt)^{\alpha+2\lambda+1}},
 \qquad t\in D_0=(-1,1)\setminus\{0\},
\]
is the weighted function of Definition~\ref{def:power-law-mobius-weight}.
Its product $|t|^\alpha g(t)$ is integrable on $D_0$.
\end{theorem}
\begin{proof}
\leanok
\end{proof}

\begin{theorem}[Fixed-singularity power-law substitution]
\label{thm:power-law-mobius}\uses{thm:power-law-gegenbauer-inputs,thm:power-law-mobius-pullback,thm:power-law-mobius-weight}
\lean{Real.powerLawMobius_gegenbauer_integrandOn,Real.powerLawIntegral_gegenbauerWeight_mobiusOn}\leanok
For $\alpha>-1$, $\lambda>-1/2$ and $-1<x<1$, the map $\phi_x(t)=(x-t)/(1-xt)$ is an involution of
$(-1,1)$ with absolute Jacobian $(1-x^2)/(1-xt)^2$. Consequently,
\[
 Q_\alpha w_\lambda(x)=(1-x^2)^{\alpha+\lambda+1/2}
 \int_{-1}^1\frac{|t|^\alpha(1-t^2)^{\lambda-1/2}}
 {(1-xt)^{\alpha+2\lambda+1}}\,dt.
\]
Both integrals are classical and integrable. The substitution holds for every
integrable input by Theorem~\ref{thm:power-law-mobius-pullback}.
The pointwise transformed-kernel identity holds for all real $\alpha,\lambda$
on $D_0$; the integral identity uses the stated integrability range. No value
of a singular kernel at $t=0$ is asserted.
Reference: the degree-zero operator in \cite[Section 2]{gutleb-carrillo-olver}.
\end{theorem}
\begin{proof}
\leanok
\end{proof}

\begin{theorem}[Degree-zero power-law hypergeometric evaluation]
\label{thm:power-law-hypergeometric-split}\uses{thm:hypergeometric-euler-integral,def:power-law-gegenbauer-inputs,thm:gegenbauer-measure-characterisation}
\lean{Real.powerLawLeft_gegenbauerWeight_hypergeometricOn,Real.powerLawRight_gegenbauerWeight_hypergeometricOn,Real.powerLawIntegral_gegenbauerWeight_hypergeometricOn}\leanok
Let $\alpha>-1$, $\lambda>-1/2$, $-1<x<1$, and put
\[
 p=\alpha+\lambda+\tfrac12,\quad
 A=\frac{2^{\lambda-1/2}\Gamma(\lambda+1/2)\Gamma(\alpha+1)}
 {\Gamma(\alpha+\lambda+3/2)},\quad
 H(z)={}_2F_1(\tfrac12-\lambda,\lambda+\tfrac12;
                  \alpha+\lambda+\tfrac32;z).
\]
The left and right weighted fractional integrals are, respectively,
$A(1+x)^pH((1+x)/2)$ and $A(1-x)^pH((1-x)/2)$. In particular,
\[
 Q_\alpha w_\lambda(x)=A\left[(1+x)^pH((1+x)/2)
                                  +(1-x)^pH((1-x)/2)\right].
\]
Both Gauss arguments lie in $(0,1)$. This evaluates the full degree-zero
operator throughout its natural interior range, using the classical real
Gauss function. Theorem~\ref{thm:power-law-hypergeometric-symmetric}
gives the equivalent single Gauss value at $x^2$.
References: \cite[Section 2]{gutleb-carrillo-olver}; \cite[15.6.1]{dlmf}.
\end{theorem}
\begin{proof}
\leanok
\end{proof}



\begin{theorem}[Even quadratic Gauss transformation]
\label{thm:hypergeometric-even-quadratic}\uses{thm:hypergeometric-calculus, thm:gamma-function, thm:open-domain-analytic}
\lean{eqOn_of_secondOrder_linear_ode,Real.ordinaryHypergeometricOn_even_quadratic}\leanok
For real $a$, $p>-1$ and $|x|<1$,
\begin{align*}
 &(1+x)^p F(a,1-a;p+1;(1+x)/2)\\
 &\quad +(1-x)^p F(a,1-a;p+1;(1-x)/2)\\
 &=2F(a,1-a;p+1;1/2)
 F((1-p-a)/2,(a-p)/2;1/2;x^2).
\end{align*}
This includes all terminating cases. A reusable second-order linear ODE
uniqueness theorem compares the transformed functions and their initial data;
all Gauss arguments remain inside the series disk.
Reference: quadratic transformations in \cite[15.8]{dlmf} and the operator
formula in \cite[Section 2]{gutleb-carrillo-olver}.
\end{theorem}
\begin{proof}
\leanok
\end{proof}

\begin{definition}[Classical terminating Gauss polynomial]
\label{def:hypergeometric-polynomial}\uses{def:gauss-coefficients}
\lean{TerminatingGaussParameters,ordinaryHypergeometricPolynomial}\leanok
For $k\in\mathbb N$ and admissible $(a,-k,c)$, define
\[H_k(a,c;X)=\sum_{n=0}^{k}u_n(a,-k;c)X^n.\]
Admissibility allows earlier termination by the other numerator and termination before a later denominator pole.
Reference: \cite[15.2.4--15.2.5]{dlmf}.
\end{definition}

\begin{theorem}[Terminating Gauss polynomials]
\label{thm:hypergeometric-polynomial}\uses{def:hypergeometric-polynomial,thm:ordinary-hypergeometric}
\lean{TerminatingGaussParameters.ofNoPole,ordinaryHypergeometric_eq_eval_polynomial,natDegree_ordinaryHypergeometricPolynomial_le}\leanok
Every non-pole denominator parameter gives admissible terminating data.
For all admissible $(a,-k,c)$, the classical Gauss function equals $H_k(a,c;z)$
at every argument $z$, and $\deg H_k\le k$. The bound permits earlier termination.
Reference: \cite[15.2.4--15.2.5]{dlmf}.
\end{theorem}
\begin{proof}
\leanok
\end{proof}

\begin{theorem}[Symmetric degree-zero hypergeometric formula]
\label{thm:power-law-hypergeometric-symmetric}\uses{thm:power-law-hypergeometric-split,thm:hypergeometric-even-quadratic,thm:integrable-input}
\lean{Real.powerLawIntegral_gegenbauerWeight_originOn,Real.powerLawIntegral_gegenbauerWeight_hypergeometricSqOn}\leanok
For $\alpha>-1$, $\lambda>-1/2$, and $-1<x<1$,
\[
 Q_\alpha w_\lambda(x)=
 \frac{\Gamma((\alpha+1)/2)\Gamma(\lambda+1/2)}
      {\Gamma(\lambda+\alpha/2+1)}
 {}_2F_1(-\alpha/2,-\lambda-\alpha/2;1/2;x^2).
\]
  The full natural parameter range is
retained, including terminating and nonterminating series.
Reference: \cite[Section 2]{gutleb-carrillo-olver}.
\end{theorem}
\begin{proof}
\leanok
\end{proof}

\rvnote{thm:power-law-hypergeometric-symmetric}{
Apply the even quadratic transformation to the two-value formula.

The Beta integral at $x=0$ fixes the normalization.
}

\begin{theorem}[Weighted power-law stability and interior mass bounds]
\label{thm:power-law-stability}\uses{thm:power-law-hypergeometric-symmetric,thm:power-law-gegenbauer-inputs}
\lean{Real.abs_powerLawIntegral_le_weightedMass,Real.abs_powerLawIntegral_sub_le_weightedMass,Real.abs_powerLawIntegral_gegenbauerContinuous_sub_le,Real.continuous_powerLawIntegral_weightedMass,Real.exists_powerLawIntegral_weightedMass_bound}\leanok
Let $\alpha>-1$, $\lambda>-1/2$, and $w_\lambda(y)=(1-y^2)^{\lambda-1/2}$.
For continuous amplitudes $f,g$ on $[-1,1]$ with $|f-g|\le\varepsilon$ there,
\[
 |Q_\alpha(w_\lambda f)(x)-Q_\alpha(w_\lambda g)(x)|
 \le \varepsilon Q_\alpha w_\lambda(x),\qquad -1<x<1.
\]
The corresponding absolute bound holds for every function $u$ in the
domain of $v\mapsto Q_\alpha v(x)$ satisfying $|u(y)|\le\varepsilon w_\lambda(y)$
on $D_x$, with $\varepsilon\ge0$. More generally, two functions $u,v$ in this domain
satisfying $|u-v|\le\varepsilon w_\lambda$ obey the same difference bound.
Thus bounded amplitudes are included whenever their weighted kernel is
integrable; continuous amplitudes satisfy this automatically. The mass $Q_\alpha w_\lambda$ is continuous on $(-1,1)$
and bounded above by a positive constant on every compact subset.
No polynomial compatibility condition is required.
\end{theorem}
\begin{proof}
\leanok
\end{proof}

\begin{theorem}[Non-polynomial power-law approximation]
\label{thm:power-law-approximation}\uses{thm:power-law-stability,def:power-law-polynomial-functional}
\lean{Real.tendstoUniformlyOn_powerLaw_gegenbauerContinuous,Real.exists_polynomial_powerLaw_gegenbauerContinuous_approximation,Real.abs_powerLawIntegral_gegenbauerContinuous_pair_sub_le}\leanok
Under the preceding parameter hypotheses, if continuous amplitudes $f_n$ converge
uniformly to a continuous $f$ on $[-1,1]$, then $Q_\alpha(w_\lambda f_n)$ converges
uniformly to $Q_\alpha(w_\lambda f)$ on every compact subset of $(-1,1)$.
For each such compact set and each positive output tolerance, a polynomial
amplitude achieves that tolerance. For a second power $\beta>-1$ and real $a,b$,
the same input error $\varepsilon$ controls the output error for
$aQ_\alpha+bQ_\beta$ by
\[
 \varepsilon\bigl(|a|Q_\alpha w_\lambda(x)+|b|Q_\beta w_\lambda(x)\bigr).
\]
Both powers use the same weight and amplitudes; neither needs compatibility.
These are interior approximation statements, not endpoint convergence or
finite bandwidth claims.
\end{theorem}
\begin{proof}
\leanok
\end{proof}

\begin{definition}[Weighted functions for power-law integration on the closed output interval]
\label{def:power-law-closed-inputs}\uses{thm:power-law-gegenbauer-inputs}
\lean{Real.PowerLawClosedParameters,Real.PowerLawInput.gegenbauerContinuousClosed,Real.PowerLawInput.gegenbauerWeightClosed}\leanok
Take $\alpha>-1$, $\lambda>-1/2$ and $\alpha+\lambda>-1/2$.
For $x\in[-1,1]$ and $f\in C([-1,1],\R)$, define the functions
$y\mapsto w_\lambda(y)f(y)$ and $y\mapsto w_\lambda(y)$ on $D_x$
as elements of the domain of $u\mapsto Q_\alpha u(x)$, which requires
their products with $|x-y|^\alpha$ to be integrable on $D_x$.
No value of a singular density is assigned on the excluded endpoints or diagonal.
\end{definition}
\begin{theorem}[Characterisation of the closed-interval weighted functions]
\label{thm:power-law-closed-inputs}\uses{def:power-law-closed-inputs}
\lean{Real.PowerLawInput.gegenbauerContinuousClosed_apply,Real.PowerLawInput.gegenbauerWeightClosed_apply,Real.PowerLawInput.gegenbauerContinuousClosed_integrable}\leanok
Suppose $\alpha>-1$, $\lambda>-1/2$ and $\alpha+\lambda>-1/2$.
For every $x\in[-1,1]$ and $f\in C([-1,1],\R)$, the weighted functions of
Definition~\ref{def:power-law-closed-inputs} are $w_\lambda(y)f(y)$ and
$w_\lambda(y)$ on $D_x$; their products with $|x-y|^\alpha$ are integrable.
This includes both endpoints. At a single endpoint, the more general
combined-exponent domain of Definition~\ref{def:power-law-gegenbauer-inputs}
remains available; the three inequalities here give one domain valid for all
$x\in[-1,1]$.
\end{theorem}
\begin{proof}
\leanok
\end{proof}

\begin{theorem}[Closed-interval weighted power-law stability]
\label{thm:power-law-closed-stability}\uses{thm:power-law-closed-inputs}
\lean{Real.abs_powerLawIntegral_le_weightedMassClosed,Real.abs_powerLawIntegral_sub_le_weightedMassClosed,Real.abs_powerLawIntegral_gegenbauerContinuousClosed_sub_le}\leanok
Assume $\alpha>-1$, $\lambda>-1/2$ and $\alpha+\lambda>-1/2$.
Then $|x-y|^\alpha w_\lambda(y)$ is integrable on $D_x=(-1,1)\setminus\{x\}$ for every
$x\in[-1,1]$, including both endpoint choices of $x$. Multiplication by a continuous amplitude
preserves this integrability. If $f,g$ are continuous on $[-1,1]$ and
$|f-g|\le\varepsilon$, then
\[
 |Q_\alpha(w_\lambda f)(x)-Q_\alpha(w_\lambda g)(x)|
 \le\varepsilon Q_\alpha w_\lambda(x),\qquad x\in[-1,1].
\]
The absolute bound also holds for every integrable input $u$ on $D_x$
with $|u|\le\varepsilon w_\lambda$, where $\varepsilon\ge0$. Two integrable
inputs with $|u-v|\le\varepsilon w_\lambda$ satisfy the same difference bound.

\end{theorem}
\begin{proof}
\leanok
\end{proof}

\rvnote{thm:power-law-closed-stability}{
These are sufficient endpoint hypotheses; no necessity or singular asymptotic statement is included here.
}

\begin{theorem}[Uniform power-law mass bounds and closed-interval convergence]
\label{thm:power-law-closed-mass}\uses{thm:power-law-closed-stability}
\lean{Real.powerLawIntegral_weightedMassClosed_le_of_nonneg,Real.powerLawIntegral_weightedMassClosed_le_unweighted,Real.powerLawIntegral_weightedMassClosed_le_of_half_le,Real.exists_powerLawIntegral_weightedMassClosed_bound,Real.tendstoUniformly_powerLaw_gegenbauerContinuousClosed}\leanok
If $\alpha\ge0$ and $\lambda>-1/2$, then, for every $x\in[-1,1]$,
\[
 Q_\alpha w_\lambda(x)\le2^\alpha\int_{-1}^1w_\lambda(y)\,dy.
\]
If $\alpha>-1$ and $\lambda\ge1/2$, then
\[
 Q_\alpha w_\lambda(x)\le
 \frac{(1+x)^{\alpha+1}+(1-x)^{\alpha+1}}{\alpha+1}
 \le\frac{2\,2^{\alpha+1}}{\alpha+1}.
\]
Consequently a positive uniform mass bound exists in either regime.
More generally, under the preceding closed-interval hypotheses, any positive
uniform mass bound upgrades uniform convergence of continuous input amplitudes
to uniform convergence of their weighted integral outputs on all of $[-1,1]$.

\end{theorem}
\begin{proof}
\leanok
\end{proof}

\rvnote{thm:power-law-closed-mass}{The singular-kernel, singular-weight regime is treated in Theorem~\ref{thm:power-law-singular-mass}.}


\begin{theorem}[Schur operators from dominated marginals]
\label{thm:schur-l2-marginals}
\lean{MeasureTheory.dominatedL2Pullback,MeasureTheory.dominatedL2Pullback_ae,MeasureTheory.norm_dominatedL2Pullback_le,MeasureTheory.schurL2Operator,MeasureTheory.norm_schurL2Operator_le,MeasureTheory.inner_schurL2Operator,MeasureTheory.integrable_schurL2_pair,MeasureTheory.map_fst_withDensity_prod,MeasureTheory.map_snd_withDensity_prod}\leanok
Let $\kappa$ be a measure on $X\times Y$, with marginals bounded by
$C\mu$ and $D\nu$, where $C,D$ are finite nonnegative constants.
Pullback along the coordinate projections defines bounded maps
$J_X:L^2(\mu)\to L^2(\kappa)$ and $J_Y:L^2(\nu)\to L^2(\kappa)$
with norms at most $\sqrt C$ and $\sqrt D$.
Then $T=J_X^*J_Y$ satisfies
\[
 \|T\|\le\sqrt C\sqrt D,\qquad
 \langle v,Tu\rangle=\int_{X\times Y}v(x)u(y)\,d\kappa(x,y).
\]
The pairing is absolutely integrable for all $u\in L^2(\nu)$ and
$v\in L^2(\mu)$. For measurable nonnegative densities
$d\kappa=k\,d(\mu\times\nu)$ and sigma-finite $\mu,\nu$, the marginal
densities are $\int k(x,y)\,d\nu(y)$ and $\int k(x,y)\,d\mu(x)$.
\end{theorem}
\begin{proof}
\leanok
\end{proof}

\begin{definition}[Classical joint power-law kernel measures]
\label{def:power-law-kernel-measure}\uses{def:gegenbauer-measure}
\lean{Real.PowerLawKernelDomain,Real.twoWeightPowerLawKernelMeasure,Real.powerLawKernelMeasure}\leanok
Let $D=\{(x,y)\in(-1,1)^2:x\ne y\}$. For arbitrary real
$\alpha,\tau,\lambda$, restrict $\mu_\tau\otimes\mu_\lambda$ to $D$, multiply by its positive density $|x-y|^\alpha$, and push forward
to $\R^2$. This defines the joint two-weight kernel measure. The common-weight
kernel measure is its specialization $\tau=\lambda$.
These measures need not be finite. No diagonal or exterior kernel value is assigned.
\end{definition}
\begin{theorem}[Characterisation of the joint kernel measures]
\label{thm:power-law-kernel-measure}\uses{def:power-law-kernel-measure}
\lean{Real.powerLawKernelDomain_base_pos,Real.twoWeightPowerLawKernelMeasure_lintegral,Real.twoWeightPowerLawKernelMeasure_self}\leanok
The distance base is strictly positive on $D$. For every nonnegative measurable
$\varphi$ on $\R^2$, integration against the joint measure is exactly
\[
 \int_D |x-y|^\alpha\varphi(x,y)\,d(\mu_\tau\otimes\mu_\lambda).
\]
The equality is in $[0,\infty]$, with no finiteness requirement.
Equal input and output weights give the common-weight kernel measure.
\end{theorem}
\begin{proof}
\leanok
\end{proof}

\begin{definition}[Uniform power-law mass bounds]
\label{def:power-law-l2-data}\uses{def:power-law-gegenbauer-inputs}
\lean{Real.PowerLawL2Data}\leanok
Mass-bound data consist of $\alpha>-1$, $\lambda>-1/2$, a finite nonnegative
real bound $M$, and $Q_\alpha w_\lambda(x)\le M$ for every $x\in(-1,1)$.
Here $Q_\alpha w_\lambda(x)=\int_{D_x}|x-y|^\alpha w_\lambda(y)\,dy$, with an integrable integrand on $D_x$.
\end{definition}
\begin{theorem}[Characterisation of the uniform mass bound]
\label{thm:power-law-l2-data}\uses{def:power-law-l2-data}
\lean{Real.PowerLawL2Data.mass_le}\leanok
For $(\alpha,\lambda,M)$ in the mass-bound domain of
Definition~\ref{def:power-law-l2-data}, every $x\in(-1,1)$ satisfies
\[
 Q_\alpha w_\lambda(x)
 =\int_{D_x}|x-y|^\alpha w_\lambda(y)\,dy\le M.
\]
\end{theorem}
\begin{proof}
\leanok
\end{proof}

\begin{definition}[Classical bounded power-law operator]
\label{def:power-law-l2-operator}\uses{thm:gegenbauer-measure-bridge,thm:power-law-kernel-measure,thm:schur-l2-marginals,thm:power-law-gegenbauer-inputs,thm:power-law-l2-data}
\lean{Real.powerLawL2Operator}\leanok
For $(\alpha,\lambda,M)$ satisfying Definition~\ref{def:power-law-l2-data}, define the bounded operator on
$L^2(\mu_\lambda;\R)$ by the Schur construction for the classical joint
power-law kernel measure $K_{\alpha,\lambda}$.
\end{definition}
\begin{theorem}[Characterisation of the bounded power-law operator]
\label{thm:power-law-l2-characterisation}\uses{def:power-law-l2-operator}
\lean{Real.powerLawL2Operator_spec,Real.powerLawL2Operator_eqOn}\leanok
For every $f,g\in L^2(\mu_\lambda)$ the pairing $g(x)f(y)$ is integrable
against $K_{\alpha,\lambda}$ and
$\langle g,Tf\rangle=\int g(x)f(y)\,dK_{\alpha,\lambda}$.
The operator is independent of the chosen uniform mass bound.
\end{theorem}
\begin{proof}
\leanok
\end{proof}

\begin{theorem}[Power-law operators on a common Gegenbauer L² space]
\label{thm:power-law-l2-operator}\uses{thm:power-law-closed-mass,thm:power-law-l2-characterisation}
\lean{Real.lintegral_powerLaw_gegenbauerMeasureOn,Real.powerLawKernelMeasure_fst_leOn,Real.powerLawKernelMeasure_snd_leOn,Real.norm_powerLawL2Operator_leOn,Real.powerLawL2Operator_hasPowerLawIntegral,Real.exists_powerLawL2Data}\leanok
Let $\alpha>-1$, $\lambda>-1/2$, $d\mu=w_\lambda(x)\,dx$ on $(-1,1)$,
and suppose $Q_\alpha w_\lambda(x)\le M$ throughout $(-1,1)$, with $M\ge0$.
There is a bounded operator $T_\alpha:L^2(\mu)\to L^2(\mu)$ with
$\|T_\alpha\|\le M$, characterized by
\[
 \langle g,T_\alpha f\rangle
 =\iint_{\substack{(x,y)\in(-1,1)^2\\x\ne y}} |x-y|^\alpha g(x)f(y)\,d\mu(x)d\mu(y)
 =\int F(x)\,d\mu(x).
\]
The double-integral pairing is absolutely integrable for all $f,g\in L^2(\mu)$.
Here $F$ is integrable and, for almost every $x$, is the actual power-law
integral of $y\mapsto g(x)w_\lambda(y)f(y)$ on $D_x$, whose product with $|x-y|^\alpha$ is integrable.
This is the classical weak Fubini identity; no value is asserted for a
divergent slice on an exceptional set.
The construction is independent of the chosen uniform mass bound.
The preceding explicit estimates supply such a bound whenever $\alpha\ge0$
or $\lambda\ge1/2$, retaining $\alpha>-1$ and $\lambda>-1/2$.
The operator itself allows $\lambda=0$ and does not require polynomial compatibility.
Its identification with the integral is stated in weak form.
\end{theorem}
\begin{proof}
\leanok
\end{proof}

\begin{theorem}[Extension of the compatible polynomial operator]
\label{thm:power-law-l2-compatible}\uses{thm:power-law-compatible-operator,thm:general-polynomial-L2-basic,thm:power-law-l2-characterisation}
\lean{Real.powerLawL2Operator_polynomial_compatibleOn}\leanok
For $(\alpha,\lambda,M)$ in the mass-bound domain of Definition~\ref{def:power-law-l2-data}, if $\lambda\ne0$ and
$\lambda+\alpha/2=k\in\mathbb N$, then $T_\alpha$ applied to the $L^2$ class
of a polynomial $p$ equals the $L^2$ class of the previously constructed exact
polynomial operator applied to $p$. Thus the analytic and algebraic operators agree.
\end{theorem}
\begin{proof}
\leanok
\end{proof}


\begin{theorem}[Strong convergence of finite sections on $\ell^2$]
\label{thm:l2-finite-sections}\uses{thm:mathlib-hilbert-sums}
\lean{lp.initialL2Projection,lp.initialL2Projection_apply,lp.norm_initialL2Projection_apply_le,lp.norm_initialL2Projection_le,lp.tendsto_initialL2Projection,lp.norm_initialL2Sections_sub_le,lp.tendsto_initialL2Sections}\leanok
Over $\mathbb R$ or $\mathbb C$, let $P_N$ retain the first $N$ coordinates
in $\ell^2(\mathbb N)$. Then $\|P_N\|\le1$ and $P_Nc\to c$ for every
$c\in\ell^2$. For any bounded linear operator $A$ on this space,
\[
 \|P_NAP_Nc-Ac\|
 \le \|A\|\,\|P_Nc-c\|+\|P_NAc-Ac\|,
 \qquad P_NAP_Nc\longrightarrow Ac.
\]
The convergence is strong operator convergence.
\end{theorem}
\begin{proof}
\leanok
\end{proof}

\rvnote{thm:l2-finite-sections}{
This is strong convergence for each input, without an assertion of operator-norm convergence or convergence of inverse solves.
}

\begin{theorem}[Bounded operators in Hilbert coordinates]
\label{thm:hilbert-coordinate-operator}
\lean{HilbertBasis.coefficientOperator,HilbertBasis.coefficientOperator_repr,HilbertBasis.norm_coefficientOperator,HilbertBasis.coefficientOperator_single,HilbertBasis.coefficientOperator_sub,HilbertBasis.norm_coefficientOperator_sub}\leanok
Let $(b_n)_{n\ge0}$ be a real or complex Hilbert basis, with coordinate
isometry $U$, and let $T$ be a bounded linear operator. The coordinate
operator $A=UTU^{-1}$ satisfies
\[
 A(Uf)=U(Tf),\qquad \|A\|=\|T\|,\qquad
 (Ae_j)_i=\langle b_i,Tb_j\rangle.
\]
Transport also commutes with subtraction, so $\|A_S-A_T\|=\|S-T\|$.
\end{theorem}
\begin{proof}
\leanok
\end{proof}

\begin{definition}[Power-law coefficient operator]
\label{def:power-law-l2-coefficient-operator}\uses{def:power-law-l2-operator,thm:hilbert-coordinate-operator,thm:gegenbauer-basis}
\lean{Real.powerLawCoefficientOperator}\leanok
For $(\alpha,\lambda,M)$ satisfying Definition~\ref{def:power-law-l2-data}, with $\lambda\ne0$, define $A=UTU^{-1}$,
where $U$ is the isometry furnished by the normalized Gegenbauer Hilbert basis.
\end{definition}
\begin{theorem}[Characterisation of the coefficient operator]
\label{thm:power-law-l2-coefficient-operator}\uses{def:power-law-l2-coefficient-operator}
\lean{Real.powerLawCoefficientOperator_reprOn}\leanok
Let $T$ be the mass-bounded power-law operator on $L^2(\mu_\lambda;\R)$
and let $U$ be its normalized Gegenbauer coordinate isometry.
The coefficient operator $A$ satisfies $A(Uf)=U(Tf)$ for every
$f\in L^2(\mu_\lambda;\R)$, on the parameter domain of
Definition~\ref{def:power-law-l2-coefficient-operator}.
\end{theorem}
\begin{proof}
\leanok
\end{proof}

\begin{theorem}[Power-law coefficient-space and finite-section convergence]
\label{thm:power-law-l2-coefficients}\uses{thm:l2-finite-sections,thm:power-law-l2-operator,def:power-law-l2-coefficient-operator,thm:gegenbauer-L2-expansion}
\lean{Real.tendsto_powerLawL2Operator_gegenbauer_sumsOn,Real.norm_powerLawCoefficientOperator_leOn,Real.tendsto_powerLawCoefficientOperator_finiteSectionsOn}\leanok
Assume the bounded-operator hypotheses and $\lambda\ne0$. Let $U$ be the
isometry from $L^2(\mu)$ to real $\ell^2$ furnished by the normalized Gegenbauer
Hilbert basis. The coefficient operator $A=UT_\alpha U^{-1}$ has $\|A\|\le M$
and represents $T_\alpha$ exactly. These are orthonormal Hilbert coordinates,
not the unscaled classical Gegenbauer coefficients.
Applying $T_\alpha$ to the classical Gegenbauer partial sums of any $L^2$ input
converges in $L^2$ to its image. Moreover, if $P_N$ retains the first $N$
normalized coefficients, then
\[
 P_N A P_N c\longrightarrow Ac\quad\hbox{in }\ell^2
 \qquad(c\in\ell^2).
\]
The convergence is strong operator convergence.
\end{theorem}
\begin{proof}
\leanok
\end{proof}

\rvnote{thm:power-law-l2-coefficients}{
This is strong convergence of operator actions; no operator-norm convergence, finite bandwidth for incompatible parameters or inverse-solver convergence is asserted.
}

\begin{theorem}[Mass bounds for simultaneous kernel and endpoint singularities]
\label{thm:power-law-singular-mass}\uses{thm:power-law-closed-mass}
\lean{Real.rpow_mul_rpow_le_add_rpowOn,Real.powerLaw_gegenbauerWeight_le_three_powersOn,Real.powerLawIntegral_one_leOn,Real.powerLawIntegral_weightedMassClosed_le_singular,Real.exists_powerLawIntegral_weightedMassClosed_bound_of_endpoint}\leanok
For positive $a,b$ and nonpositive $p,q$,
$a^pb^q\le a^{p+q}+b^{p+q}$.
For $\alpha<0$, $\lambda\le1/2$, $x\in\R$, $y\in D_x=(-1,1)\setminus\{x\}$ and $s=\alpha+\lambda-1/2$,
\[
 |x-y|^\alpha w_\lambda(y)
 \le |x-y|^s+|1-y|^s+|-1-y|^s.
\]
If additionally $\alpha>-1$, $\lambda>-1/2$ and $\alpha+\lambda>-1/2$, this gives
\[
 Q_\alpha w_\lambda(x)\le
 \frac{6\,2^{\alpha+\lambda+1/2}}{\alpha+\lambda+1/2},
 \qquad x\in[-1,1].
\]
 Together with the earlier nonsingular regimes,
a positive uniform mass bound therefore exists for every
$\alpha>-1$, $\lambda>-1/2$, $\alpha+\lambda>-1/2$.

\end{theorem}
\begin{proof}
\leanok
\end{proof}

\rvnote{thm:power-law-singular-mass}{No claim about the critical or more singular ranges is made.}


\rvnote{thm:power-law-singular-mass}{
The constant is not claimed sharp.
}

\begin{theorem}[Power-law convergence throughout the strict endpoint-integrable range]
\label{thm:power-law-endpoint-l2}\uses{thm:power-law-singular-mass,thm:power-law-l2-coefficients}
\lean{Real.tendstoUniformly_powerLaw_gegenbauerContinuousClosed_of_endpoint,Real.exists_powerLawL2Data_of_endpoint,Real.exists_powerLawCoefficientOperator_of_endpointOn}\leanok
Assume $\alpha>-1$, $\lambda>-1/2$ and $\alpha+\lambda>-1/2$.
Uniform convergence of continuous amplitudes on $[-1,1]$ gives uniform convergence
of their weighted power-law integral outputs on the entire closed interval.
There is a bounded power-law operator on the common space $L^2(\mu_\lambda)$,
with the weak integral characterization and a finite norm bound supplied above.
If $\lambda\ne0$, the normalized Gegenbauer coefficient operator has the same norm bound and its leading finite sections satisfy $P_NAP_Nc\to Ac$ in
$\ell^2$ for every $c\in\ell^2$. No polynomial compatibility is needed.
These remain statements about operator actions, not inverse solves or optimal
boundedness thresholds.
\end{theorem}
\begin{proof}
\leanok
\end{proof}


\begin{theorem}[Binomial specialization of the Gauss function]
\label{thm:hypergeometric-cancellation}\uses{thm:hypergeometric-calculus,thm:gamma-function}
\lean{Real.ordinaryHypergeometricOn_cancel,Real.ordinaryHypergeometricOn_cancel_of_pos}\leanok
For real $a,c,z$, if $(c)_n\ne0$ for every $n\ge0$ and $|z|<1$, then
${}_2F_1(a,c;c;z)=(1-z)^{-a}$. In particular, the denominator condition
holds for $c>0$. Reference: \cite[15.4.6]{dlmf}.
\end{theorem}
\begin{proof}
\leanok
\end{proof}

\begin{theorem}[Changing a Gegenbauer density]
\label{thm:gegenbauer-density-shift}\uses{thm:gegenbauer-measure-characterisation}
\lean{Polynomial.gegenbauerMeasure_shift}\leanok
For any real $\lambda,p$, multiplication of $\mu_\lambda$ by
$(1-x^2)^p$ gives $\mu_{\lambda+p}$. These are identities of measures
on $(-1,1)$, without a finiteness assumption. The multiplying factor is
the Gegenbauer weight $w_{p+1/2}(x)=(1-x^2)^p$.
\end{theorem}
\begin{proof}
\leanok
\end{proof}

\begin{definition}[Measurable positive weights]
\label{def:positive-measurable-weight}
\lean{MeasureTheory.PositiveWeight}\leanok
For a measure $\mu$, a positive measurable weight is a measurable real function $h$ such that $h>0$ $\mu$-almost everywhere.
No pointwise positivity on a null set is imposed.
\end{definition}
\begin{theorem}[Characterisation of positive measurable weights]
\label{thm:positive-measurable-weight}\uses{def:positive-measurable-weight}
\lean{MeasureTheory.PositiveWeight.spec}\leanok
A weight $h$ in the positive-measurable-weight domain is measurable and
satisfies $h(x)>0$ for $\mu$-almost every $x$.
\end{theorem}
\begin{proof}
\leanok
\end{proof}

\begin{definition}[Classical density-change isometry]
\label{def:division-weight-l2}\uses{def:positive-measurable-weight}
\lean{MeasureTheory.divWeightL2}\leanok
For a positive measurable weight $h$, define $D_h:L^2(\mu;\R)\to L^2(h^2\mu;\R)$
as the linear isometry characterized by $(D_hu)h=u$ $h^2\mu$-almost everywhere
for every $u\in L^2(\mu;\R)$. This defines an $L^2$ equivalence class and
assigns no pointwise quotient at a zero of $h$.
\end{definition}
\begin{theorem}[Characterisation of the density-change isometry]
\label{thm:division-weight-l2}\uses{def:division-weight-l2}
\lean{MeasureTheory.divWeightL2_spec,MeasureTheory.divWeightL2_unique}\leanok
The isometry satisfies $(D_hu)h=u$ $h^2\mu$-almost everywhere and is uniquely
determined by this identity. Equivalently, it is division by $h$ on the conull
positive-weight domain. Neither $h\in L^2(\mu)$ nor finiteness of $h^2\mu$ is required.
\end{theorem}
\begin{proof}
\leanok
\end{proof}

\begin{definition}[Classical weighted Schur operator]
\label{def:weighted-schur-l2}\uses{def:division-weight-l2,thm:schur-l2-marginals}
\lean{MeasureTheory.weightedSchurL2Operator}\leanok
For a positive measurable weight $h$ and a joint measure whose marginals are
bounded by $C h^2\mu$, with $0\le C<\infty$, let $S$ be the ordinary Schur
operator on $L^2(h^2\mu)$. Define $T=D_h^* S D_h$ on $L^2(\mu)$.
\end{definition}

\begin{theorem}[Weighted Schur construction]
\label{thm:weighted-schur-l2}\uses{def:weighted-schur-l2}
\lean{MeasureTheory.norm_weightedSchurL2Operator_le,MeasureTheory.inner_weightedSchurL2Operator,MeasureTheory.integrable_weightedSchurL2_pair}\leanok
Let $h$ be real measurable and positive $\mu$-almost everywhere.
Division by $h$ is an isometry $D:L^2(\mu;\mathbb R)\to L^2(h^2\mu;\mathbb R)$.
If a measure $\kappa$ on $X\times X$ has both marginals at most
$C h^2\mu$, with $0\le C<\infty$, then there is a bounded operator $T$ on
$L^2(\mu;\mathbb R)$ with
\[
 \|T\|\le C,\qquad
 \langle v,Tu\rangle=\int (D_hv)(x)(D_hu)(y)\,d\kappa(x,y).
\]
This pairing is absolutely integrable for every pair of $L^2$ functions.
Neither $h\in L^2(\mu)$ nor finiteness of $h^2\mu$ is required.
\end{theorem}
\begin{proof}
\leanok
\end{proof}

\begin{definition}[Critical Schur auxiliary weight]
\label{def:power-law-schur-weight}\uses{def:gegenbauer-measure}
\lean{Real.powerLawSchurWeight}\leanok
For every real $\alpha$, define the continuous positive weight
$h_\alpha(x)=(1-x^2)^{\alpha/2}$ on $x\in(-1,1)$.
\end{definition}
\begin{theorem}[Characterisation of the auxiliary weight]
\label{thm:power-law-schur-weight}\uses{def:power-law-schur-weight,thm:gegenbauer-measure-characterisation}
\lean{Real.powerLawSchurWeight_apply,Real.powerLawSchurWeight_pos}\leanok
For every real $\alpha$ and $x\in(-1,1)$,
\[
 h_\alpha(x)=(1-x^2)^{\alpha/2}>0.
\]
The domain of $h_\alpha$ excludes the endpoints and exterior of the interval.
\end{theorem}
\begin{proof}
\leanok
\end{proof}

\begin{definition}[Classical critical Schur constant]
\label{def:power-law-schur-constant}\uses{thm:gamma-function}
\lean{Real.PowerLawSchurParameter,Real.powerLawSchurConstant}\leanok
For real $\alpha$ such that neither $(\alpha+1)/2$ nor $-\alpha/2$ is a
nonpositive integer, define
$C_\alpha=\Gamma((\alpha+1)/2)\Gamma(-\alpha/2)/\Gamma(1/2)$.
\end{definition}
\begin{theorem}[Characterisation of the Schur constant]
\label{thm:power-law-schur-constant}\uses{def:power-law-schur-constant}
\lean{Real.PowerLawSchurParameter.ofCritical,Real.powerLawSchurConstant_apply,Real.powerLawSchurConstant_pos}\leanok
Whenever neither $(\alpha+1)/2$ nor $-\alpha/2$ is a nonpositive integer,
\[
 C_\alpha=\frac{\Gamma((\alpha+1)/2)\Gamma(-\alpha/2)}{\Gamma(1/2)}.
\]
Every $-1<\alpha<0$ satisfies these conditions and gives $C_\alpha>0$.
\end{theorem}
\begin{proof}
\leanok
\end{proof}

\begin{theorem}[Exact critical power-law Schur weight]
\label{thm:power-law-critical-weight}\uses{thm:hypergeometric-cancellation,thm:gegenbauer-density-shift,thm:power-law-schur-weight,thm:power-law-schur-constant,thm:power-law-hypergeometric-symmetric,thm:power-law-gegenbauer-inputs,thm:power-law-kernel-measure,thm:gegenbauer-measure-bridge}
\lean{Real.powerLawIntegral_gegenbauerWeight_criticalOn,Real.lintegral_powerLaw_critical_weightOn}\leanok
For $-1<\alpha<0$ and $x\in(-1,1)$, set
\[
 h_\alpha(x)=(1-x^2)^{\alpha/2},\qquad
 C_\alpha=\frac{\Gamma((\alpha+1)/2)\Gamma(-\alpha/2)}{\Gamma(1/2)}>0.
\]
Then $Q_\alpha w_{-\alpha/2-1/2}(x)=C_\alpha h_\alpha(x)$.
Consequently, if $\lambda>-1/2$ and $\alpha+\lambda=-1/2$,
\[
 \int_{-1}^1 |x-y|^\alpha h_\alpha(y)\,d\mu_\lambda(y)
 = C_\alpha h_\alpha(x).
\]
The nonnegative integral is taken on $D_x=(-1,1)\setminus\{x\}$, where all factors have their classical values.
The auxiliary weight is positive throughout the support domain of every Gegenbauer measure.
\end{theorem}
\begin{proof}
\leanok
\end{proof}

\begin{definition}[Tilted critical power-law kernel measure]
\label{def:power-law-schur-measure}\uses{def:power-law-kernel-measure}
\lean{Real.powerLawSchurMeasure}\leanok
For real $\alpha,\lambda$, define the joint kernel measure with exponent
$\alpha$ and Gegenbauer parameter $\lambda+\alpha/2$ in both variables.
It is the classical measure obtained by tilting the original kernel by
$h_\alpha(x)h_\alpha(y)$ on its off-diagonal open-square domain.
\end{definition}
\begin{theorem}[Characterisation of the tilted kernel measure]
\label{thm:power-law-schur-measure}\uses{def:power-law-schur-measure,thm:power-law-critical-weight,thm:schur-l2-marginals}
\lean{Real.powerLawSchurMeasure_apply,Real.powerLawSchurMeasure_fstOn,Real.powerLawSchurMeasure_sndOn}\leanok
The measure is exactly the kernel measure with parameter $\lambda+\alpha/2$.
For $-1<\alpha<0$ and $\lambda=-\alpha-1/2$, both marginals are
$C_\alpha\mu_{-1/2}=C_\alpha h_\alpha^2\mu_\lambda$.
\end{theorem}
\begin{proof}
\leanok
\end{proof}

\begin{definition}[Critical weak power-law operator]
\label{def:power-law-critical-operator}\uses{thm:power-law-schur-measure,thm:weighted-schur-l2,thm:division-weight-l2}
\lean{Real.criticalPowerLawL2Operator}\leanok
For $\alpha>-1$, $\lambda>-1/2$, and $\alpha+\lambda=-1/2$, define the
bounded operator on $L^2(\mu_\lambda;\R)$ by its integrable weak pairing
$\langle g,Tf\rangle=\int g(x)f(y)\,dK_{\alpha,\lambda}(x,y)$,
where $K_{\alpha,\lambda}$ is the classical joint power-law kernel measure.
\end{definition}
\begin{theorem}[Characterisation of the critical weak operator]
\label{thm:power-law-critical-operator}\uses{def:power-law-critical-operator}
\lean{Real.criticalPowerLawL2Operator_spec,Real.criticalPowerLawL2Operator_unique}\leanok
Assume $\alpha>-1$, $\lambda>-1/2$ and $\alpha+\lambda=-1/2$. For every $f,g\in L^2(\mu_\lambda;\R)$,
the function $(x,y)\mapsto g(x)f(y)$ is integrable against the classical
joint kernel measure $K_{\alpha,\lambda}$, and
\[
 \langle g,Tf\rangle=\int g(x)f(y)\,dK_{\alpha,\lambda}(x,y).
\]
This identity for all $f,g$ uniquely determines the bounded operator $T$.
\end{theorem}
\begin{proof}
\leanok
\end{proof}

\begin{theorem}[Critical same-space power-law boundedness]
\label{thm:power-law-critical-l2}\uses{thm:power-law-endpoint-l2,thm:power-law-critical-operator}
\lean{Real.norm_criticalPowerLawL2Operator_leOn,Real.inner_criticalPowerLawL2Operator_hasPowerLawIntegral,Real.exists_powerLawL2Operator_of_le_endpointOn}\leanok
Assume $\alpha>-1$, $\lambda>-1/2$ and $\alpha+\lambda=-1/2$.
The measure with density $|x-y|^\alpha h_\alpha(x)h_\alpha(y)$
on $\{(x,y)\in(-1,1)^2:x\ne y\}$ against $\mu_\lambda\otimes\mu_\lambda$
has both marginals exactly $C_\alpha h_\alpha^2\mu_\lambda$.
It defines a bounded operator $T$ on $L^2(\mu_\lambda;\mathbb R)$ with
$\|T\|\le C_\alpha$ and
\[
 \langle g,Tf\rangle=
 \iint_{\substack{(x,y)\in(-1,1)^2\\x\ne y}} |x-y|^\alpha g(x)f(y)\,d\mu_\lambda(x)d\mu_\lambda(y).
\]
The double pairing is absolutely integrable for every $f,g\in L^2$;
there is a $\mu_\lambda$-integrable function $F$ such that, for almost every $x$,
the function $y\mapsto |x-y|^\alpha g(x)w_\lambda(y)f(y)$ is integrable on $D_x$, $F(x)$ is its integral, and
$\langle g,Tf\rangle=\int F\,d\mu_\lambda$.
Together with the strict case, a bounded weak kernel realization exists whenever
$\alpha+\lambda\ge-1/2$.
\end{theorem}
\begin{proof}
\leanok
\end{proof}

\rvnote{thm:power-law-critical-l2}{
No sharp operator norm or uniform endpoint convergence is asserted here.
}


\begin{theorem}[Critical power-law coefficient and finite-section convergence]
\label{thm:power-law-critical-coefficients}\uses{thm:power-law-critical-l2}
\lean{Real.criticalPowerLawCoefficientOperator,Real.norm_criticalPowerLawCoefficientOperator,Real.norm_criticalPowerLawCoefficientOperator_le,Real.criticalPowerLawCoefficientOperator_repr,Real.criticalPowerLawCoefficientOperator_singleOn,Real.tendsto_criticalPowerLawCoefficientOperator_finiteSections,Real.tendsto_criticalPowerLawL2Operator_gegenbauer_sums}\leanok
Assume $\alpha>-1$, $\lambda>-1/2$, $\lambda\ne0$ and
$\alpha+\lambda=-1/2$. Let $T$ be the critical weighted Schur operator,
and let $U$ be the normalized Gegenbauer Hilbert-coordinate isometry.
Then $A=UTU^{-1}$ satisfies
\[
 \|A\|=\|T\|\le C_\alpha,\qquad A(Uf)=U(Tf),\qquad
 P_NAP_Nc\longrightarrow Ac\quad(c\in\ell^2).
\]
Writing $b_i$ for the normalized $L^2$ basis vectors, its entries are
\[
 (Ae_j)_i=\iint_{\substack{(x,y)\in(-1,1)^2\\x\ne y}} |x-y|^\alpha b_i(x)b_j(y)\,
 d\mu_\lambda(x)d\mu_\lambda(y).
\]
Applying $T$ to the classical Gegenbauer partial sums of any $L^2$ input
also converges in $L^2$ to its image. Together with the strict-range result,
normalized coefficient actions and finite-section convergence therefore cover
$\alpha+\lambda\ge-1/2$ whenever $\lambda\ne0$.
The coefficient restriction comes from the classical Gegenbauer basis;
the underlying $L^2$ boundedness theorem includes $\lambda=0$.
\end{theorem}
\begin{proof}
\leanok
\end{proof}

\begin{definition}[Classical weak L² kernel data]
\label{def:weak-l2-kernel}\uses{def:product-measure-input}
\lean{MeasureTheory.L2WeakKernel,MeasureTheory.L2WeakKernel.productInput}\leanok
For a measure $\mu$, weak kernel data include almost-everywhere measurability
of $x\mapsto (y\mapsto(x,y))_*\mu$ as a measure-valued map, certifying the
product $\mu\times\mu$. They consist further of a real kernel $k(x,y)$,
a bounded operator $T$ on $L^2(\mu;\R)$, and, for every $L^2$ pair $f,g$,
absolute integrability of $k(x,y)g(x)f(y)$ and the identity
$\langle g,Tf\rangle=\int k(x,y)g(x)f(y)\,d(\mu\times\mu)$.
\end{definition}
\begin{theorem}[Characterisation of weak L² kernel data]
\label{thm:weak-l2-kernel}\uses{def:weak-l2-kernel,thm:product-measure-input}
\lean{MeasureTheory.L2WeakKernel.spec,MeasureTheory.L2WeakKernel.ofSFinite,MeasureTheory.L2WeakKernel.pairing_representatives}\leanok
Let $(\mu,k,T)$ satisfy the weak-kernel conditions of
Definition~\ref{def:weak-l2-kernel}. For every $f,g\in L^2(\mu;\R)$,
$k(x,y)g(x)f(y)$ is integrable and
\[
 \langle g,Tf\rangle=\int k(x,y)g(x)f(y)\,d(\mu\otimes\mu)(x,y).
\]
The integral does not depend on the pointwise representatives of $f,g$.
For s-finite $\mu$, the measure-valued product kernel is automatically
almost-everywhere measurable; in particular the product-measurability
condition adds no restriction for finite measures.
\end{theorem}
\begin{proof}
\leanok
\end{proof}

\begin{theorem}[Rectangle tests for an L² kernel operator]
\label{thm:kernel-rectangle-test}\uses{def:weak-l2-kernel}
\lean{MeasureTheory.L2WeakKernel.integrableOn_rectangle,MeasureTheory.L2WeakKernel.rectangle_bound}\leanok
Let $\mu$ be finite and let a bounded real $L^2(\mu)$ operator $T$ have
weak kernel $k(x,y)$ as in Definition~\ref{def:weak-l2-kernel}.
For measurable $A,B$, the rectangle integral exists and
\[
 \left|\int_{A\times B}k(x,y)\,d(\mu\times\mu)\right|
 \le \|T\|\sqrt{\mu(A)}\sqrt{\mu(B)}.
\]
\end{theorem}
\begin{proof}
\leanok
\end{proof}


\begin{definition}[Domain of classical weak power-law realizations]
\label{def:power-law-weak-realization}\uses{def:power-law-kernel-measure}
\lean{Real.PowerLawWeakRealization}\leanok
For real $\alpha,\lambda$, a weak realization is a bounded operator $T$ on
$L^2(\mu_\lambda;\R)$ such that the function $g(x)f(y)$ is integrable against the classical joint kernel measure for every $L^2$ pair $f,g$, and
$\langle g,Tf\rangle=\int g(x)f(y)\,dK_{\alpha,\lambda}$.
This domain may be empty.
\end{definition}
\begin{theorem}[Characterisation of weak power-law realizations]
\label{thm:power-law-weak-realization}\uses{def:power-law-weak-realization}
\lean{Real.PowerLawWeakRealization.spec}\leanok
For real $\alpha,\lambda$ and a weak realization $T$ in the sense of
Definition~\ref{def:power-law-weak-realization}, every
$f,g\in L^2(\mu_\lambda;\R)$ satisfies
\[
 \langle g,Tf\rangle=\int g(x)f(y)\,dK_{\alpha,\lambda}(x,y),
\]
with an integrable integrand. This assertion is conditional on the existence
of $T$ in that domain.
\end{theorem}
\begin{proof}
\leanok
\end{proof}

\begin{theorem}[Endpoint obstruction to power-law boundedness]
\label{thm:power-law-endpoint-obstruction}\uses{thm:kernel-rectangle-test,thm:gegenbauer-moments,thm:power-law-weak-realization,thm:weak-l2-kernel,thm:power-law-kernel-measure,thm:gegenbauer-measure-characterisation}
\lean{Real.gegenbauerMeasure_endpoint_interval_lowerOn,Real.powerLaw_rectangle_lower_le_normOn,Real.powerLaw_endpoint_scaling_le_normOn,Real.powerLaw_rectangle_lower_le_norm_of_agreesOn,Real.powerLaw_endpoint_scaling_le_norm_of_agreesOn,Real.not_exists_powerLawL2Operator_below_endpointOn}\leanok
If $-1/2<\lambda\le1/2$, $0<\delta\le1$, and
$1-\delta\le a<b\le1$, then
\[
 \mu_\lambda((a,b))\ge(2\delta)^{\lambda-1/2}(b-a).
\]
For $\alpha\le0$, any bounded weak kernel realization $T$ on
$L^2(\mu_\lambda;\mathbb R)$ satisfies
\[
 \|T\|\ge \frac{2^{\lambda-1/2}}4\,
 \delta^{\alpha+\lambda+1/2}\qquad(0<\delta\le1).
\]
More generally, separated measurable sets with
$\delta/4\le|x-y|\le\delta$ give the lower bound
$\delta^\alpha\sqrt{\mu_\lambda(A)\mu_\lambda(B)}\le\|T\|$
when both masses are positive. Both bounds also hold for a weak kernel
that agrees with the power kernel on the tested rectangle, or on the
endpoint strip squared away from the diagonal, respectively. No diagonal
value of a singular kernel is prescribed.
Thus, for $\alpha>-1$, $\lambda>-1/2$ and $\alpha+\lambda<-1/2$,
no bounded same-space $L^2$ operator has the stated weak power-law kernel.
\end{theorem}
\begin{proof}
\leanok
\end{proof}
\end{supporting}


\begin{theorem}[Exact same-space power-law boundedness criterion]
\label{thm:power-law-l2-classification}\uses{thm:power-law-critical-l2,thm:power-law-endpoint-obstruction}
\lean{Real.exists_powerLawL2Operator_iffOn}\leanok
Assume $\alpha>-1$ and $\lambda>-1/2$.
There exists a bounded operator $T$ on $L^2(\mu_\lambda;\mathbb R)$ satisfying
\[
 \langle g,Tf\rangle=
 \iint_{\substack{(x,y)\in(-1,1)^2\\x\ne y}} |x-y|^\alpha g(x)f(y)\,d\mu_\lambda(x)d\mu_\lambda(y)
\]
with absolute integrability for every $f,g\in L^2(\mu_\lambda;\mathbb R)$
if and only if $\alpha+\lambda\ge-1/2$.
This includes $\lambda=0$. It classifies this common-space realization;
it is not a two-weight mapping criterion or an inverse-solver theorem.
\end{theorem}
\begin{proof}
\leanok
\end{proof}


\begin{supporting}[Reference details: Normalized coordinates including the Chebyshev limit and related results]
\begin{definition}[Normalized coordinates including the Chebyshev limit]
\label{def:power-law-full-basis}\uses{thm:T-basis,thm:classical-jacobi-measures,thm:gegenbauer-basis}
\lean{Polynomial.gegenbauerMeasure_zero,Polynomial.hilbertBasisChebyshevGegenbauerZero,Polynomial.powerLawHilbertBasis,Polynomial.powerLawHilbertBasis_zero,Polynomial.powerLawHilbertBasis_of_ne_zero}\leanok
The measures satisfy $\mu_0=\mu_T$. For $\lambda>-1/2$, choose the
normalized Gegenbauer Hilbert basis when $\lambda\ne0$, and the
normalized Chebyshev $T$ Hilbert basis when $\lambda=0$, transported along
this equality of measures. Denote this basis by $(b_n)$ and its coordinate
isometry by $U$.
\end{definition}

\rvnote{def:power-law-full-basis}{
This is a case distinction in the choice of coordinates, not an assertion of continuity in $\lambda$.
}

\begin{definition}[Classical full-range power-law realization]
\label{def:power-law-full-realization}\uses{thm:power-law-l2-classification}
\lean{Real.boundedPowerLawL2Operator}\leanok
For $\alpha>-1$, $\lambda>-1/2$, and $\alpha+\lambda\ge-1/2$, select the
bounded operator on $L^2(\mu_\lambda;\R)$ characterized by integrable pairings
$\langle g,Tf\rangle=\int g(x)f(y)\,dK_{\alpha,\lambda}$ against the classical
joint kernel measure.
\end{definition}
\begin{theorem}[Characterisation of the full-range realization]
\label{thm:power-law-full-characterisation}\uses{def:power-law-full-realization}
\lean{Real.boundedPowerLawL2Operator_spec,Real.boundedPowerLawL2Operator_eqOn}\leanok
Assume $\alpha>-1$, $\lambda>-1/2$ and $\alpha+\lambda\ge-1/2$. For every $f,g\in L^2(\mu_\lambda;\R)$,
the function $(x,y)\mapsto g(x)f(y)$ is integrable against the classical
joint kernel measure $K_{\alpha,\lambda}$, and
\[
 \langle g,Tf\rangle=\int g(x)f(y)\,dK_{\alpha,\lambda}(x,y).
\]
This identity for all $f,g$ uniquely determines the bounded operator $T$.
\end{theorem}
\begin{proof}
\leanok
\end{proof}

\begin{theorem}[Unique full-range weak power-law realization]
\label{thm:power-law-full-realization}\uses{thm:power-law-full-characterisation}
\lean{Real.boundedPowerLawL2Operator_eq_massOn,Real.boundedPowerLawL2Operator_eq_critical}\leanok
Assume $\alpha>-1$, $\lambda>-1/2$ and $\alpha+\lambda\ge-1/2$.
There is a unique bounded operator $T$ on $L^2(\mu_\lambda;\mathbb R)$ with
\[
 \langle g,Tf\rangle=\iint_{\substack{(x,y)\in(-1,1)^2\\x\ne y}} |x-y|^\alpha g(x)f(y)\,
 d\mu_\lambda(x)d\mu_\lambda(y).
\]
Every such pairing is absolutely integrable. This realization agrees with
any earlier mass-bound construction and, on the critical line, with the
weighted Schur construction.
\end{theorem}
\begin{proof}
\leanok
\end{proof}
\end{supporting}


\begin{theorem}[Full-range coefficient convergence, including Chebyshev T]
\label{thm:power-law-full-coefficients}\uses{def:power-law-full-basis,thm:power-law-full-realization,thm:power-law-critical-coefficients}
\lean{Real.boundedPowerLawCoefficientOperator,Real.norm_boundedPowerLawCoefficientOperator,Real.norm_boundedPowerLawCoefficientOperator_le_critical,Real.boundedPowerLawCoefficientOperator_repr,Real.boundedPowerLawCoefficientOperator_singleOn,Real.tendsto_boundedPowerLawCoefficientOperator_finiteSections,Real.boundedPowerLawCoefficientOperator_eq_massOn,Real.boundedPowerLawCoefficientOperator_eq_critical,Real.norm_boundedPowerLawCoefficientOperator_le_massOn}\leanok
For every $\alpha>-1$, $\lambda>-1/2$ and $\alpha+\lambda\ge-1/2$,
use the preceding basis and weak realization. Then $A=UTU^{-1}$ satisfies
\[
 \|A\|=\|T\|,\qquad A(Uf)=U(Tf),\qquad
 P_NAP_Nc\longrightarrow Ac\quad(c\in\ell^2),
\]
and its entries are the absolutely integrable kernel pairings
\[
 (Ae_j)_i=\iint_{\substack{(x,y)\in(-1,1)^2\\x\ne y}} |x-y|^\alpha b_i(x)b_j(y)\,
 d\mu_\lambda(x)d\mu_\lambda(y).
\]
A uniform nonnegative mass bound $M$ gives $\|A\|\le M$; on the critical
line $\|A\|\le C_\alpha$. For $\lambda\ne0$, this operator equals the
previous Gegenbauer coefficient constructions. At $\lambda=0$ it uses
normalized Chebyshev $T$ coordinates and covers $\alpha\ge-1/2$, including
$\alpha=-1/2$. There is no exclusion of zero from this full-range theorem.
Convergence is strong for each input; no operator-norm convergence, new
bandwidth assertion or inverse-solver convergence is claimed.
\end{theorem}
\begin{proof}
\leanok
\end{proof}

\begin{supporting}[Reference details: Classical two-weight Schur operator and related results]
\begin{definition}[Classical two-weight Schur operator]
\label{def:two-weight-schur-l2}\uses{def:division-weight-l2,thm:schur-l2-marginals}
\lean{MeasureTheory.twoWeightSchurL2Operator}\leanok
For positive measurable weights $h$ for $\mu$ and $g$ for $\nu$, and a joint
measure with marginals bounded by $C h^2\mu$ and $D g^2\nu$, where $C,D$ are
finite and nonnegative, define $T=D_h^* S D_g$. Here $S$ is the ordinary Schur
operator between the two tilted $L^2$ spaces.
\end{definition}

\begin{theorem}[Asymmetric weighted Schur construction]
\label{thm:two-weight-schur-l2}\uses{def:two-weight-schur-l2}
\lean{MeasureTheory.norm_twoWeightSchurL2Operator_le,MeasureTheory.inner_twoWeightSchurL2Operator,MeasureTheory.integrable_twoWeightSchurL2_pair}\leanok
Let $\mu$ on $X$ and $\nu$ on $Y$ have measurable positive almost-everywhere
auxiliary weights $h,g$. If a joint measure $\kappa$ satisfies
\[
 (\pi_X)_*\kappa\le C h^2\mu,\qquad
 (\pi_Y)_*\kappa\le D g^2\nu
\]
for finite nonnegative $C,D$, then there is a bounded
$T:L^2(\nu;\mathbb R)\to L^2(\mu;\mathbb R)$ with
\[
 \|T\|\le\sqrt{CD},\qquad
 \langle v,Tu\rangle=\int (D_hv)(x)(D_gu)(y)\,d\kappa.
\]
The pairing is absolutely integrable for all $u,v$. No finiteness of the
original or tilted measures, or square-integrability of $h,g$, is required.
\end{theorem}
\begin{proof}
\leanok
\end{proof}

\begin{theorem}[Coordinate transport between different Hilbert spaces]
\label{thm:hilbert-coordinate-map}
\lean{HilbertBasis.coefficientMap,HilbertBasis.coefficientMap_repr,HilbertBasis.norm_coefficientMap,HilbertBasis.coefficientMap_single}\leanok
For real or complex Hilbert bases $(a_j)$ of $E$ and $(b_i)$ of $F$ and
bounded $T:E\to F$, the coefficient map $A=U_F T U_E^{-1}$ satisfies
\[
 A(U_E f)=U_F(Tf),\qquad \|A\|=\|T\|,\qquad
 (Ae_j)_i=\langle b_i,Ta_j\rangle.
\]
The source and target spaces and bases may differ.
\end{theorem}
\begin{proof}
\leanok
\end{proof}

\begin{definition}[Two-weight power-law operator]
\label{def:power-law-two-weight-l2}\uses{thm:power-law-kernel-measure,thm:schur-l2-marginals,thm:gegenbauer-measure-bridge,thm:power-law-gegenbauer-inputs,thm:power-law-l2-data}
\lean{Real.twoWeightPowerLawL2Operator}\leanok
For a common exponent $\alpha>-1$, assume the uniform mass bounds of Definition~\ref{def:power-law-l2-data} for input
weight $\lambda>-1/2$ and output weight $\tau>-1/2$, with respective bounds
$M,N\ge0$. Define $T:L^2(\mu_\lambda)\to L^2(\mu_\tau)$ by the Schur
construction for the classical two-weight joint kernel measure.
\end{definition}
\begin{theorem}[Characterisation of the two-weight operator]
\label{thm:power-law-two-weight-characterisation}\uses{def:power-law-two-weight-l2}
\lean{Real.twoWeightPowerLawL2Operator_spec}\leanok
For $T:L^2(\mu_\lambda;\R)\to L^2(\mu_\tau;\R)$ in the two-weight realization domain, every $f\in L^2(\mu_\lambda;\R)$ and $g\in L^2(\mu_\tau;\R)$ satisfy:
$g(x)f(y)$ is integrable against the joint kernel measure and its integral
equals $\langle g,Tf\rangle$.
\end{theorem}
\begin{proof}
\leanok
\end{proof}
\end{supporting}


\begin{theorem}[Sufficient two-weight power-law bounds]
\label{thm:power-law-two-weight-l2}\uses{thm:power-law-endpoint-l2,thm:power-law-two-weight-characterisation}
\lean{Real.twoWeightPowerLawKernelMeasure_fst_leOn,Real.twoWeightPowerLawKernelMeasure_snd_leOn,Real.norm_twoWeightPowerLawL2Operator_leOn,Real.twoWeightPowerLawL2Operator_hasPowerLawIntegral,Real.exists_twoWeightPowerLawL2Operator_of_endpointOn}\leanok
Let $\alpha>-1$, $\lambda,\tau>-1/2$. Suppose the input-weight row masses
are bounded by $M\ge0$ and the output-weight row masses by $N\ge0$ throughout
$(-1,1)$. Then the kernel measure
with density $|x-y|^\alpha$ on $\{(x,y)\in(-1,1)^2:x\ne y\}$
against $\mu_\tau\otimes\mu_\lambda$ defines
$T:L^2(\mu_\lambda)\to L^2(\mu_\tau)$, with
\[
 \|T\|\le\sqrt{MN},\qquad
 \langle g,Tf\rangle=\iint_{\substack{(x,y)\in(-1,1)^2\\x\ne y}} |x-y|^\alpha g(x)f(y)\,
 d\mu_\tau(x)d\mu_\lambda(y).
\]
Every weak pairing is absolutely integrable. Fubini provides an integrable
function $F$ on the output space and, for $\mu_\tau$-almost every $x$, an integrable function $y\mapsto |x-y|^\alpha g(x)w_\lambda(y)f(y)$ on $D_x$ whose integral is $F(x)$, with
$\langle g,Tf\rangle=\int F\,d\mu_\tau$.
Such positive mass bounds exist whenever
$\alpha+\lambda>-1/2$ and $\alpha+\tau>-1/2$.

\end{theorem}
\begin{proof}
\leanok
\end{proof}

\rvnote{thm:power-law-two-weight-l2}{
These are sufficient conditions, not a sharp two-weight classification.
}

\begin{supporting}[Reference details: Two-weight coefficient map and related results]
\begin{definition}[Two-weight coefficient map]
\label{def:power-law-two-weight-coefficients}\uses{def:power-law-two-weight-l2,thm:hilbert-coordinate-map,def:power-law-full-basis}
\lean{Real.twoWeightPowerLawCoefficientOperator}\leanok
Define $A=U_\tau T U_\lambda^{-1}$ using the normalized source and target
Hilbert bases, including the Chebyshev $T$ choice at a zero parameter.
\end{definition}
\begin{theorem}[Characterisation of the two-weight coefficient map]
\label{thm:power-law-two-weight-coefficient-characterisation}\uses{def:power-law-two-weight-coefficients}
\lean{Real.twoWeightPowerLawCoefficientOperator_reprOn}\leanok
For the two-weight operator $T:L^2(\mu_\lambda;\R)\to L^2(\mu_\tau;\R)$
and coordinate isometries $U_\lambda,U_\tau$, its coefficient map $A$ satisfies
\[
 A(U_\lambda f)=U_\tau(Tf)\qquad(f\in L^2(\mu_\lambda;\R)).
\]
\end{theorem}
\begin{proof}
\leanok
\end{proof}

\begin{theorem}[Two-weight power-law coefficient sections]
\label{thm:power-law-two-weight-coefficients}\uses{def:power-law-two-weight-coefficients,thm:power-law-two-weight-characterisation,thm:l2-finite-sections}
\lean{Real.norm_twoWeightPowerLawCoefficientOperatorOn,Real.twoWeightPowerLawCoefficientOperator_singleOn,Real.tendsto_twoWeightPowerLawCoefficientOperator_finiteSectionsOn}\leanok
Under the preceding mass-bound hypotheses, use normalized Gegenbauer coordinates
at each nonzero parameter and normalized Chebyshev $T$ coordinates at zero.
The coefficient map $A=U_\tau T U_\lambda^{-1}$ preserves the operator norm,
intertwines the action, and has entries
\[
 (Ae_j)_i=\iint_{\substack{(x,y)\in(-1,1)^2\\x\ne y}} |x-y|^\alpha b_i^\tau(x)b_j^\lambda(y)\,
 d\mu_\tau(x)d\mu_\lambda(y).
\]
For every $c\in\ell^2$, $P_NAP_Nc\to Ac$ strongly. Either or both weight
parameters may be zero.
\end{theorem}
\begin{proof}
\leanok
\end{proof}

\rvnote{thm:power-law-two-weight-coefficients}{
No additional bandwidth or inverse-solver convergence is asserted.
}
\end{supporting}


\begin{theorem}[Self-adjoint weighted power-law operators]
\label{thm:power-law-self-adjoint}\uses{thm:power-law-full-coefficients}
\lean{Real.inner_boundedPowerLawL2Operator_symm,Real.isSelfAdjoint_boundedPowerLawL2Operator,Real.isSelfAdjoint_boundedPowerLawCoefficientOperator}\leanok
Assume $\alpha>-1$, $\lambda>-1/2$ and $\alpha+\lambda\ge-1/2$.
Let $A$ be the weak realization of
\[
 (Af)(x)=\int_{-1}^1 |x-y|^\alpha f(y)\,d\mu_\lambda(y).
\]
Then $\langle g,Af\rangle=\langle f,Ag\rangle$ for all real
$L^2(\mu_\lambda)$ vectors. Both $A$ and its normalized coefficient
operator are self-adjoint, including on the critical boundary and at
$\lambda=0$, where the normalized basis is Chebyshev $T$.
\end{theorem}
\begin{proof}
\leanok
\end{proof}

\begin{supporting}[Reference details: Domination of weak kernels controls operator norms and related results]
\begin{theorem}[Domination of weak kernels controls operator norms]
\label{thm:dominated-weak-kernel}
\lean{MeasureTheory.norm_le_of_dominated_weak_kernel}\leanok
Let $\mu$ be an s-finite measure. Let bounded real $L^2(\mu)$ operators
$T,S$ have weak kernels $k,q$, respectively, with all their weak pairings
absolutely integrable. If $c\ge0$ and $|k(x,y)|\le c q(x,y)$ almost
everywhere for $\mu\otimes\mu$, then $\|T\|\le c\|S\|$.
\end{theorem}
\begin{proof}
\leanok
\end{proof}

\begin{theorem}[Compact continuous bounded kernels]
\label{thm:continuous-bounded-kernel-compact}\uses{thm:bounded-kernel-l2,thm:continuous-l2-rows-compact,thm:l2-dominated-convergence}
\lean{MeasureTheory.memLp_continuousKernel_row,MeasureTheory.continuousKernelRow,MeasureTheory.continuousKernelRow_ae,MeasureTheory.inner_continuousKernelRow,MeasureTheory.norm_continuousKernelRow_le,MeasureTheory.memLp_inner_continuousKernelRow,MeasureTheory.boundedKernelL2Operator_ae_inner_row,MeasureTheory.continuous_continuousKernelRow,MeasureTheory.isCompactOperator_boundedKernelL2Operator_of_continuous}\leanok
Let $\mu$ be a finite Borel measure on $\mathbb R$ and $k$ a continuous
real kernel on $\mathbb R^2$ with $|k|\le M$, $M\ge0$.
Each row $r_x=k(x,\cdot)$ defines an $L^2(\mu)$ vector, continuously
in $x$, and $\|r_x\|\le\sqrt{\mu(\mathbb R)}M$.
For $f\in L^2(\mu)$, its row pairing belongs to $L^2(\mu)$ and represents
the bounded-kernel operator almost everywhere:
\[
 T_k f(x)=\langle r_x,f\rangle=\int k(x,y)f(y)\,d\mu(y).
\]
If $\mu$ is concentrated on a compact set and a countable Hilbert basis
of $L^2(\mu)$ is given, then $T_k$ is compact.
\end{theorem}
\begin{proof}
\leanok
\end{proof}

\begin{definition}[Classical regularized power-law kernel]
\label{def:power-law-regularized-kernel}
\lean{Real.PowerLawCutoff,Real.PowerLawOffDiagonalSquare,Real.regularizedPowerLawKernel}\leanok
For a real exponent $\alpha$, admissible cutoffs satisfy $\delta>0$ or
$\alpha\ge0$. On $\R^2$ define
\[
 k_{\alpha,\delta}(x,y)=\max\{\delta,\min\{|x-y|,2\}\}^{\alpha}.
\]
A zero base is allowed only for a nonnegative exponent, with $0^0=1$.
For comparison with the unregularized kernel, the physical off-diagonal domain
consists of pairs $x,y\in[-1,1]$ with $x\ne y$.
\end{definition}
\begin{theorem}[Characterisation of regularized kernels]
\label{thm:power-law-regularized-domain}\uses{def:power-law-regularized-kernel}
\lean{Real.regularizedPowerLawKernel_apply,Real.regularizedPowerLawKernel_base_domain,Real.continuous_regularizedPowerLawKernel_on}\leanok
For real $\alpha$ and an admissible cutoff $\delta$ (that is,
$\delta>0$ or $\alpha\ge0$), the regularized kernel is
\[
 k_{\alpha,\delta}(x,y)=\bigl(\max\{\delta,\min\{|x-y|,2\}\}\bigr)^\alpha.
\]
Its base is positive if $\delta>0$; otherwise it is nonnegative and
$\alpha\ge0$. For every admissible cutoff, $k_{\alpha,\delta}$ is continuous
on $\R^2$.
\end{theorem}
\begin{proof}
\leanok
\end{proof}

\begin{theorem}[Continuous regularization of power-law kernels]
\label{thm:power-law-regularized-kernel}\uses{def:power-law-regularized-kernel}
\lean{Real.norm_regularizedPowerLawKernel_le_on,Real.abs_sub_regularizedPowerLawKernel_le_on,Real.exists_lower_powerLaw_exponent}\leanok
Define $k_{\alpha,\delta}(x,y)=(\max\{\delta,\min\{|x-y|,2\}\})^\alpha$.
For $\delta>0$ this is continuous; if $\delta\le2$, then
$|k_{\alpha,\delta}|\le\delta^\alpha+2^\alpha$ globally.
For $\beta<\alpha$, $\beta\le0$, and distinct $x,y\in[-1,1]$,
\[
 \bigl||x-y|^\alpha-k_{\alpha,\delta}(x,y)\bigr|
 \le2\delta^{\alpha-\beta}|x-y|^\beta.
\]
Whenever $\alpha>-1$, $\lambda>-1/2$, and $\alpha+\lambda>-1/2$,
there exists $\beta<\min\{\alpha,0\}$ with $\beta>-1$ and
$\beta+\lambda>-1/2$.
\end{theorem}
\begin{proof}
\leanok
\end{proof}

\begin{theorem}[Strict-range compactness and norm-convergent power-law sections]
\label{thm:power-law-compact-sections}\uses{thm:power-law-regularized-kernel,thm:dominated-weak-kernel,thm:continuous-bounded-kernel-compact,thm:power-law-full-coefficients}
\lean{Real.regularizedPowerLawL2Operator,Real.isCompactOperator_regularizedPowerLawL2Operator,Real.norm_sub_regularizedPowerLawL2Operator_le,Real.tendsto_regularizedPowerLawL2Operator,Real.isCompactOperator_boundedPowerLawL2Operator,Real.isCompactOperator_boundedPowerLawCoefficientOperator,Real.tendsto_boundedPowerLawCoefficientOperator_sections_norm}\leanok
Let $R_{\alpha,\delta}$ be the bounded realization of $k_{\alpha,\delta}$
on $L^2(\mu_\lambda)$, $\lambda>-1/2$, $0<\delta\le2$. It is compact.
If both exponents $\alpha,\beta$ belong to the boundedness range,
$\beta<\alpha$, and $\beta\le0$, then
\[
 \|A_\alpha-R_{\alpha,\delta}\|
 \le2\delta^{\alpha-\beta}\|A_\beta\|.
\]
Consequently, for $\alpha>-1$ and $\alpha+\lambda>-1/2$,
$R_{\alpha,1/(n+1)}\to A_\alpha$ in operator norm and $A_\alpha$ is compact.
Its normalized coefficient realization $\widehat A_\alpha$ is compact too, and
\[
 \|P_n\widehat A_\alpha P_n-\widehat A_\alpha\|\longrightarrow0.
\]

\end{theorem}
\begin{proof}
\leanok
\end{proof}

\rvnote{thm:power-law-compact-sections}{
No assertion of compactness on the critical boundary is made here.
}
\end{supporting}


\begin{theorem}[Conditional second-kind power-law solves]
\label{thm:power-law-second-kind-solves}\uses{thm:power-law-compact-sections,thm:compact-finite-section-convergence,thm:l2-finite-matrices}
\lean{Real.isUnit_one_sub_smul_powerLawCoefficientOperator_iff,Real.isUnit_one_sub_smul_powerLawCoefficientOperator_of_norm_lt_one,Real.existsUnique_powerLawL2_solution,Real.powerLawCoefficientSection,Real.invertible_powerLawCoefficientSection_filter_neBot,Real.eventually_stable_powerLawCoefficientSectionOn,Real.eventually_existsUnique_powerLaw_projected_solution,Real.eventually_existsUnique_powerLaw_matrix_solution,Real.eventually_norm_powerLawCoefficientSection_error_leOn,Real.tendsto_powerLawCoefficientSection_solutionOn}\leanok
Throughout the boundedness range, $I-cA$ is invertible precisely when
$I-c\widehat A$ is invertible. The condition $|c|\|A\|<1$ is sufficient.
Assuming invertibility, $u-cAu=f$ has a unique physical $L^2$ solution
for every right-hand side.
In the strict range $\alpha+\lambda>-1/2$, let
$B_N=I-cP_N\widehat A$ and assume $I-c\widehat A$ invertible.
For all sufficiently large $N$, $B_N$ is invertible with
\[
 \|B_N^{-1}\|\le C:=\|(I-c\widehat A)^{-1}\|+1.
\]
For every coefficient right-hand side $f$, the projected equations
\[
 u_N\in\operatorname{ran}P_N,\qquad P_N(u_N-c\widehat A u_N)=P_Nf
\]
have exactly one solution, as does the corresponding actual $N\times N$
matrix equation. On the cofinal set of invertible cutoffs, the solutions $u_N=B_N^{-1}P_Nf$ converge to
$u=(I-c\widehat A)^{-1}f$ and eventually satisfy
\[
 \|B_N^{-1}P_N(I-c\widehat A)u-u\|\le C\|P_Nu-u\|.
\]
 These results concern linear integral equations, not equilibrium or
optimisation problems.
\end{theorem}
\begin{proof}
\leanok
\end{proof}

\rvnote{thm:power-law-second-kind-solves}{Invertibility is an explicit hypothesis: no claim is made for every coupling $c$.}


\begin{supporting}[Reference details: Vanishing multipliers and compact operators]
\begin{theorem}[Vanishing multipliers and compact operators]
\label{thm:vanishing-l2-multipliers}\uses{thm:bounded-kernel-l2,thm:l2-dominated-convergence,thm:compact-projection-approximation,def:weak-l2-kernel}
\lean{MeasureTheory.inner_boundedMulL2,MeasureTheory.tendsto_boundedMulL2_zero,MeasureTheory.tendsto_boundedMulL2_comp_compact_zero,MeasureTheory.inner_boundedMulL2_comp_kernel}\leanok
On real $L^2(\mu)$, let $Q_n$ multiply by measurable $q_n$ with
$|q_n|\le M$ almost everywhere, uniformly in $n$. If $q_n\to0$ almost
everywhere, then $Q_nf\to0$ for every $f\in L^2(\mu)$.
If $T$ is compact, $\|Q_nT\|\to0$.
A real bounded multiplier is self-adjoint. For finite $\mu$, if $T$ has weak
kernel $k(x,y)$ whose pairings $k(x,y)g(x)f(y)$ are integrable for every $f,g\in L^2(\mu;\mathbb R)$, multiplication of its
output by $q(x)$ gives integrable weak pairings with kernel $q(x)k(x,y)$.
\end{theorem}
\begin{proof}
\leanok
\end{proof}
\end{supporting}


\begin{theorem}[Sharp common-space power-law compactness threshold]
\label{thm:power-law-compactness-classification}\uses{thm:vanishing-l2-multipliers,thm:power-law-compact-sections}
\lean{Real.not_isCompactOperator_powerLaw_criticalOn,Real.isCompactOperator_boundedPowerLawL2Operator_iff,Real.isCompactOperator_boundedPowerLawCoefficientOperator_iff}\leanok
Let $\alpha>-1$, $\lambda>-1/2$ and $\alpha+\lambda\ge-1/2$.
The bounded power-law operator $A$ on $L^2(\mu_\lambda)$ and its normalized
coefficient operator $\widehat A$ satisfy
\[
 A\text{ compact}\quad\Longleftrightarrow\quad
 \widehat A\text{ compact}\quad\Longleftrightarrow\quad
 \alpha+\lambda>-1/2.
\]
More generally, every bounded weak realization of the kernel on the critical
line $\alpha+\lambda=-1/2$ is noncompact.
\end{theorem}
\begin{proof}
\leanok
\end{proof}

\begin{supporting}[Reference details: Small-norm finite sections without compactness and related results]
\begin{theorem}[Small-norm finite sections without compactness]
\label{thm:small-norm-finite-sections}\uses{thm:compact-finite-section-convergence}
\lean{ContinuousLinearMap.norm_inverse_one_sub_le,ContinuousLinearMap.isUnit_liftedFiniteSection_of_norm_lt_one,ContinuousLinearMap.norm_inverse_liftedFiniteSection_le,ContinuousLinearMap.norm_liftedFiniteSection_error_le_of_norm_lt_one,ContinuousLinearMap.tendsto_liftedFiniteSection_solution_of_norm_lt_one}\leanok
Let $E$ be a real or complex Banach space and $T:E\to E$ bounded with
$\|T\|\le q<1$. Then $I-T$ is invertible and
$\|(I-T)^{-1}\|\le(1-q)^{-1}$.
For any bounded $P$ with $\|P\|\le1$, the lifted section $B_P=I-PT$ is
invertible with the same inverse bound, and
\[
 \|B_P^{-1}P(I-T)u-u\|\le(1-q)^{-1}\|Pu-u\|.
\]
If $P_n$ are contractions converging strongly to $I$, then for every $f$,
$B_{P_n}^{-1}P_nf\to(I-T)^{-1}f$.
These estimates need neither compactness nor idempotence; for projections,
the earlier range and projected-equation lemmas identify the actual finite solve.
\end{theorem}
\begin{proof}
\leanok
\end{proof}

\begin{theorem}[Full residual error certificates]
\label{thm:inverse-residual-error}\uses{thm:small-norm-finite-sections}
\lean{ContinuousLinearMap.norm_sub_inverse_apply_le_residual,ContinuousLinearMap.norm_sub_inverse_one_sub_apply_le_residual}\leanok
For an invertible bounded operator $B$ on a real or complex normed space,
any right-hand side $f$ and any candidate $v$ satisfy
\[
 \|v-B^{-1}f\|\le\|B^{-1}\|\,\|Bv-f\|.
\]
In a Banach space, if $B=I-T$ and $\|T\|\le q<1$, the multiplier
$\|B^{-1}\|$ can be replaced by $(1-q)^{-1}$.
The residual here is the full-space residual, not merely its finite projection.
\end{theorem}
\begin{proof}
\leanok
\end{proof}
\end{supporting}


\begin{theorem}[All-resolution power-law solves at small coupling]
\label{thm:power-law-small-coupling}\uses{thm:small-norm-finite-sections,thm:power-law-second-kind-solves}
\lean{Real.isUnit_powerLaw_section_of_small_coupling,Real.norm_inverse_powerLaw_section_le_of_small_couplingOn,Real.existsUnique_powerLaw_projected_solution_of_small_coupling,Real.existsUnique_powerLaw_matrix_solution_of_small_coupling,Real.norm_powerLaw_section_error_le_of_small_couplingOn,Real.tendsto_powerLaw_section_solution_of_small_couplingOn}\leanok
Let $\alpha>-1$, $\lambda>-1/2$, $\alpha+\lambda\ge-1/2$ and
$q=|c|\|A\|<1$. Write $B_N=I-cP_N\widehat A$.
For every $N$, $B_N$ is invertible with $\|B_N^{-1}\|\le(1-q)^{-1}$.
For every $f\in\ell^2$, there is a unique solution
\[
 P_Nu_N=u_N,\qquad P_N(u_N-c\widehat A u_N)=P_Nf,
\]
and a unique solution of the corresponding actual $N\times N$ matrix equation.
The solutions $u_N=B_N^{-1}P_Nf$ converge to
$u=(I-c\widehat A)^{-1}f$, with
\[
 \|u_N-u\|\le(1-q)^{-1}\|P_Nu-u\|
\]
at every resolution. This includes the noncompact critical line.
\end{theorem}
\begin{proof}
\leanok
\end{proof}

\begin{theorem}[Power-law residual error certificates]
\label{thm:power-law-residual-certificates}\uses{thm:inverse-residual-error,thm:power-law-small-coupling}
\lean{Real.norm_powerLawL2_error_le_residualOn,Real.norm_powerLawCoefficient_error_le_residualOn,Real.norm_powerLawL2_error_le_residual_of_small_couplingOn,Real.norm_powerLawCoefficient_error_le_residual_of_small_couplingOn}\leanok
Throughout $\alpha>-1$, $\lambda>-1/2$, $\alpha+\lambda\ge-1/2$, if
$I-cA$ is invertible, then every physical candidate $v$ satisfies
\[
 \|v-(I-cA)^{-1}f\|
 \le\|(I-cA)^{-1}\|\,\|v-cAv-f\|.
\]
The same statement holds in normalized coefficient coordinates with
$\widehat A$. If $q=|c|\|A\|<1$, the inverse-norm factor in either version
can be replaced by $(1-q)^{-1}$. The residual is the full residual, including unresolved coordinates.
\end{theorem}
\begin{proof}
\leanok
\end{proof}

\rvnote{thm:power-law-residual-certificates}{
Unresolved residual coordinates are included; these theorems do not assert a computable numerical bound without a certified bound for the full residual and the inverse norm.
}

\begin{supporting}[Reference details: Gegenbauer lowering and weighted derivative and related results]
\begin{theorem}[Gegenbauer lowering and weighted derivative]
\label{thm:gegenbauer-weighted-derivative}\uses{thm:gegenbauer-ode,thm:gegenbauer-value-one,thm:open-domain-derivative,thm:jacobi-weight-calculus,def:gegenbauer-measure}
\lean{Polynomial.gegenbauer_euler_derivative,Polynomial.gegenbauer_lowering,Polynomial.gegenbauerWeight_add_one,Polynomial.hasDerivAt_gegenbauerWeight_add_one,Polynomial.hasDerivAt_gegenbauerWeight_mul_derivative}\leanok
For $\lambda>-1/2$, the cleared lowering identity is
$(1-x^2)(C_{n+1}^{(\lambda)})'+(n+1)xC_{n+1}^{(\lambda)}
=(n+2\lambda)C_n^{(\lambda)}$.
On $(-1,1)$,
\[
 \frac{d}{dx}\left[w_{\lambda+1}(x)(C_n^{(\lambda)})'(x)\right]
 =-n(n+2\lambda)w_\lambda(x)C_n^{(\lambda)}(x).
\]
The latter interior derivative identity itself needs no parameter restriction.
Reference: \cite[18.9.19--18.9.20]{dlmf}.
\end{theorem}
\begin{proof}
\leanok
\end{proof}

\begin{definition}[Classical left polynomial power-law functional]
\label{def:power-law-left-polynomial}\uses{def:power-law-gegenbauer-inputs}
\lean{Real.PowerLawLeftInput.gegenbauerPolynomial,Real.powerLawLeftPolynomial}\leanok
For $\alpha>-1$, $\lambda>-1/2$, $x\in(-1,1)$ and $P\in\R[X]$, let
$f_P(y)=w_\lambda(y)P(y)$ on $(-1,x)$. Define the linear functional
\[
 L_{\alpha,\lambda,x}(P)=\int_{y\in(-1,x)}(x-y)^\alpha w_\lambda(y)P(y)\,dy.
\]
\end{definition}
\begin{theorem}[Characterisation of the left polynomial functional]
\label{thm:power-law-left-polynomial}\uses{def:power-law-left-polynomial}
\lean{Real.PowerLawLeftInput.gegenbauerPolynomial_apply,Real.powerLawLeftPolynomial_apply}\leanok
Let $\alpha>-1$, $\lambda>-1/2$, $x\in(-1,1)$ and $P\in\R[X]$.
The weighted polynomial restricted to $(-1,x)$ has values
$f_P(y)=w_\lambda(y)P(y)$, and
\[
 L_{\alpha,\lambda,x}(P)=\int_{y\in(-1,x)}(x-y)^\alpha w_\lambda(y)P(y)\,dy.
\]
\end{theorem}
\begin{proof}
\leanok
\end{proof}

\begin{theorem}[Fixed-exponent weighted power-law recurrence]
\label{thm:power-law-mode-recurrence}\uses{thm:gegenbauer-weighted-derivative,thm:gegenbauer-parity,thm:power-law-left-polynomial,thm:jacobi-weight-closed,thm:gegenbauer-measure-characterisation}
\lean{Real.powerLawLeftPolynomial_gegenbauer_recurrenceOn,Real.powerLawRight_gegenbauer_recurrenceOn,Real.powerLawIntegral_gegenbauer_recurrenceOn}\leanok
Let $U_n(x)=Q_\alpha(w_\lambda C_n^{(\lambda)})(x)$, with
$\alpha>-1$, $\lambda>-1/2$ and $|x|<1$. For $n\ge1$,
\begin{align*}
 2(n+\lambda)n(n+2\lambda)xU_n
 &=n(n+1)(n+2\lambda+\alpha+1)U_{n+1}\\
 &\quad +(n+2\lambda)(n-\alpha-1)(n+2\lambda-1)U_{n-1}.
\end{align*}
The same cleared recurrence holds separately for the left and right integrals.
It includes $\lambda=0$ and does not divide by a polynomial norm.
Reference: \cite[Lemma 2.2 and Theorem 2.4]{gutleb-carrillo-olver}.
\end{theorem}
\begin{proof}
\leanok
\end{proof}

\begin{definition}[Continuously differentiable convolution kernels]
\label{def:c1-convolution-kernel}
\lean{Real.C1ConvolutionKernel}\leanok
A continuously differentiable convolution kernel consists of continuous real
functions $g$ and $s$ on $\R$, with $g'(t)=s(t)$ at every $t\in\R$.
\end{definition}
\begin{theorem}[Characterisation of differentiable kernels]
\label{thm:c1-convolution-kernel}\uses{def:c1-convolution-kernel}
\lean{Real.C1ConvolutionKernel.hasDeriv}\leanok
For a continuously differentiable convolution kernel $(g,s)$ as in
Definition~\ref{def:c1-convolution-kernel}, $g$ is differentiable at every
$t\in\R$ and $g'(t)=s(t)$.
\end{theorem}
\begin{proof}
\leanok
\end{proof}

\begin{definition}[Classical interval convolution]
\label{def:interval-convolution}\uses{def:integrable-input}
\lean{Real.convolutionInput,Real.intervalConvolution}\leanok
For $g\in C(\R;\R)$, an integrable real function $f$ on $(-1,1)$, and $x\in\R$,
consider the integrable function $y\mapsto g(x-y)f(y)$.
Define $V_gf(x)=\int_{-1}^1 g(x-y)f(y)\,dy$ with this integrable product as integrand.
\end{definition}
\begin{theorem}[Characterisation of interval convolution]
\label{thm:interval-convolution}\uses{def:interval-convolution,thm:c1-convolution-kernel}
\lean{Real.convolutionInput_apply,Real.intervalConvolution_apply,Real.hasDerivAt_intervalConvolution}\leanok
For $g\in C(\R;\R)$, $f\in L^1((-1,1);\R)$ with a pointwise
representative, and $x\in\R$, the convolution integrand is $g(x-y)f(y)$ and
\[
 V_gf(x)=\int_{y\in(-1,1)}g(x-y)f(y)\,dy.
\]
If $(g,s)$ is a continuously differentiable convolution kernel, then
$(V_gf)'(x)=V_sf(x)$ for $-1<x<1$.
Integrable endpoint singularities of $f$ are allowed.
\end{theorem}
\begin{proof}
\leanok
\end{proof}

\begin{definition}[Nonnegative-exponent power kernels]
\label{def:abs-power-kernel}
\lean{Real.absPowerKernel}\leanok
For $\beta\ge0$, let $k_\beta(t)=|t|^\beta$ on $\R$, a continuous real function,
with $0^0=1$.
\end{definition}
\begin{theorem}[Characterisation of power kernels]
\label{thm:abs-power-kernel}\uses{def:abs-power-kernel}
\lean{Real.absPowerKernel_apply,Real.hasDerivAt_absPowerKernel}\leanok
For $\beta\ge0$, $k_\beta(t)=|t|^\beta$ on $\R$, with $0^0=1$. For $\beta>1$ its derivative exists at
every real $t$ and is $\beta|t|^{\beta-1}\operatorname{sign}(t)$, including value
zero at $t=0$.
\end{theorem}
\begin{proof}
\leanok
\end{proof}

\begin{definition}[Differentiable power kernel data]
\label{def:c1-power-kernel}\uses{thm:abs-power-kernel,def:c1-convolution-kernel}
\lean{Real.C1ConvolutionKernel.power}\leanok
For $\beta>1$, equip $k_\beta$ with its continuous slope
$s_\beta(t)=\beta|t|^{\beta-1}\operatorname{sign}(t)$.
\end{definition}
\begin{theorem}[Characterisation of differentiable power data]
\label{thm:c1-power-kernel}\uses{def:c1-power-kernel}
\lean{Real.C1ConvolutionKernel.power_value,Real.C1ConvolutionKernel.power_slope}\leanok
For $\beta>1$, the continuously differentiable power kernel has value
$k_\beta(t)=|t|^\beta$ and slope
$s_\beta(t)=\beta|t|^{\beta-1}\operatorname{sign}(t)$ for every $t\in\R$.
\end{theorem}
\begin{proof}
\leanok
\end{proof}

\begin{definition}[Integrable functions for nonnegative power kernels]
\label{def:power-law-l1-input}\uses{def:abs-power-kernel,def:interval-convolution,def:power-law-integral}
\lean{Real.PowerLawInput.ofL1}\leanok
For $\beta\ge0$, $x\in\R$, and an integrable real function $f$ on $(-1,1)$, regard $f|_{D_x}$ as an element of the domain of $u\mapsto Q_\beta u(x)$: its product with $|x-y|^\beta$ is integrable on $D_x$.
\end{definition}
\begin{theorem}[Characterisation of the restricted integrable functions]
\label{thm:power-law-l1-input}\uses{def:power-law-l1-input,thm:interval-convolution}
\lean{Real.PowerLawInput.ofL1_apply,Real.powerLawIntegral_ofL1_eq_convolution}\leanok
For $\beta\ge0$, $x\in\R$ and integrable $f$ on $(-1,1)$, the restriction
$f|_{D_x}$ retains the values $f(y)$ at every $y\in D_x$, and
\[
 Q_\beta(f|_{D_x})(x)=V_{k_\beta}f(x),\qquad k_\beta(t)=|t|^\beta.
\]
\end{theorem}
\begin{proof}
\leanok
\end{proof}
\end{supporting}


\begin{theorem}[Differentiation of power-law integrals]
\label{thm:power-law-differentiation}\uses{def:c1-power-kernel,thm:power-law-l1-input,def:open-domain-derivative,thm:power-law-gegenbauer-inputs,thm:gegenbauer-moments}
\lean{Real.hasDerivAt_powerLawIntegralOn,Real.powerLawIntegral_gegenbauer_one_hasDerivOn}\leanok
For $\beta>1$ and any integrable input $f$ on $[-1,1]$,
\[
 (Q_\beta f)'(x)=\int_{-1}^1\beta|x-y|^{\beta-1}\operatorname{sign}(x-y)f(y)\,dy,
 \qquad |x|<1.
\]
The derivative exists. In the weighted Gegenbauer parameter range the shifted
zeroth mode also has a derivative and this gives
$U_1=2\lambda(xU_0-(Q_{\alpha+2}w_\lambda)'/(\alpha+2))$.

Reference: \cite[Section 2]{gutleb-carrillo-olver}.
\end{theorem}
\begin{proof}
\leanok
\end{proof}

\rvnote{thm:power-law-differentiation}{
The proof uses dominated differentiation for a smooth convolution kernel; it allows integrable endpoint-singular weights.
}

\begin{supporting}[Reference details: Classical initial power-law polynomial and related results]
\begin{definition}[Classical initial power-law polynomial]
\label{def:power-law-initial-polynomial}\uses{def:gamma-function,thm:hypergeometric-polynomial}
\lean{Real.PowerLawInitialData,Real.powerLawGegenbauerInitialPolynomial}\leanok
Initial data $(\alpha,\lambda)$ require all three arguments
$(\alpha+1)/2$, $\lambda+1/2$, and $\lambda+\alpha/2+1$ to avoid the
nonpositive integers. For every $k\in\N$, define
\[
 q_0^{\alpha,\lambda,k}(X)=
 \frac{\Gamma((\alpha+1)/2)\Gamma(\lambda+1/2)}{\Gamma(\lambda+\alpha/2+1)}
 H_k(-\alpha/2,1/2;X^2).
\]
Here $H_k$ is the terminating polynomial of Definition~\ref{def:hypergeometric-polynomial}.
No compatibility or positivity is required for this polynomial definition.
\end{definition}
\begin{theorem}[Characterisation of initial power-law data]
\label{thm:power-law-initial-polynomial}\uses{def:power-law-initial-polynomial,thm:gamma-function}
\lean{Real.PowerLawInitialData.shift,Real.PowerLawInitialData.ofIntegrable,Real.natDegree_powerLawGegenbauerInitialPolynomial_le}\leanok
The initial polynomial has degree at most $2k$ throughout this non-pole domain.
Shifting $\alpha$ to $\alpha+2$ preserves admissibility. Every
$\alpha>-1$, $\lambda>-1/2$ gives admissible initial data.
\end{theorem}
\begin{proof}
\leanok
\end{proof}

\begin{definition}[Classical finite power-law recurrence]
\label{def:power-law-recurrence}\uses{def:polynomial-linear-maps,thm:power-law-initial-polynomial}
\lean{Real.PowerLawRecurrenceData,Real.PowerLawRecurrenceData.truncate,Real.PowerLawFamilyData,Real.powerLawGegenbauerPolynomial}\leanok
For a mode $n$, take initial data and require $\alpha+2\ne0$ if $n\ge1$,
and $m+2\lambda+\alpha+2\ne0$ for every $m+2\le n$.
Define the finite recurrence by the preceding $q_0$,
\[
 q_1=2\lambda\left(Xq_0-\frac{(q_0^{\alpha+2,\lambda,k+1})'}{\alpha+2}\right),
\]
and, through the requested mode,
\[
 q_{m+2}=\frac{2(m+\lambda+1)(m+1)(m+1+2\lambda)Xq_{m+1}
 -(m+1+2\lambda)(m-\alpha)(m+2\lambda)q_m}
 {(m+1)(m+2)(m+2\lambda+\alpha+2)}.
\]
Whole-family data require these nonzero-denominator conditions at every step.
A finite mode requires no conditions on later unused steps.
\end{definition}
\begin{theorem}[Characterisation of the classical recurrence domain]
\label{thm:power-law-recurrence}\uses{def:power-law-recurrence}
\lean{Real.PowerLawRecurrenceData.ofIntegrable,Real.PowerLawRecurrenceData.denominator_ne_zero,Real.PowerLawFamilyData.atDegree,Real.PowerLawFamilyData.ofIntegrable,Real.natDegree_powerLawGegenbauerPolynomial_le,Real.powerLawGegenbauer_recurrence_denominator_pos}\leanok
Admissibility restricts to every earlier mode, and whole-family data give every
finite prefix. Each used denominator is nonzero and $\deg q_n\le n+2k$.
The natural range $\alpha>-1$, $\lambda>-1/2$ gives whole-family data;
there $(m+1)(m+2)(m+2\lambda+\alpha+2)>0$ for every $m\ge0$.
\end{theorem}
\begin{proof}
\leanok
\end{proof}
\end{supporting}


\begin{theorem}[Polynomial outputs for compatible power laws]
\label{thm:power-law-compatible-polynomials}\uses{thm:power-law-hypergeometric-symmetric,thm:power-law-differentiation,thm:power-law-mode-recurrence,thm:power-law-recurrence}
\lean{Real.eval_powerLawGegenbauerInitialPolynomial,Real.eval_powerLawGegenbauerPolynomial}\leanok
Suppose $\alpha>-1$, $\lambda>-1/2$ and $k=\lambda+\alpha/2\in\mathbb N$.
Each $U_n$ equals an explicitly defined polynomial $q_n$ on $(-1,1)$, with
$\deg q_n\le n+2k$. The zeroth polynomial is the terminating Gauss expression,
the first uses the derivative of the exponent-shifted zeroth polynomial,
and all later polynomials use the fixed-exponent recurrence. Its leading
recurrence denominator is strictly positive in this parameter range.
Reference: \cite[Lemma 2.7 and Theorem 2.4]{gutleb-carrillo-olver}.
\end{theorem}
\begin{proof}
\leanok
\end{proof}

\begin{supporting}[Reference details: Classical polynomial power-law functional and related results]
\begin{definition}[Classical polynomial power-law functional]
\label{def:power-law-polynomial-functional}\uses{def:power-law-gegenbauer-inputs}
\lean{Real.PowerLawInput.gegenbauerPolynomial,Real.powerLawPolynomial}\leanok
For $\alpha>-1$, $\lambda>-1/2$ and $x\in(-1,1)$, define the linear functional
on real polynomials by $P\mapsto Q_\alpha(w_\lambda P)(x)$, where $w_\lambda P$ is restricted to $D_x=(-1,1)\setminus\{x\}$.
\end{definition}
\begin{theorem}[Characterisation of the polynomial functional]
\label{thm:power-law-polynomial-functional}\uses{def:power-law-polynomial-functional}
\lean{Real.PowerLawInput.gegenbauerPolynomial_apply,Real.powerLawPolynomial_apply}\leanok
For $\alpha>-1$, $\lambda>-1/2$, $x\in(-1,1)$ and $P\in\R[X]$,
the weighted polynomial on $D_x=(-1,1)\setminus\{x\}$ has values
$w_\lambda(y)P(y)$. The polynomial functional of
Definition~\ref{def:power-law-polynomial-functional} therefore satisfies
\[
 Q_\alpha(w_\lambda P)(x)=\int_{D_x}|x-y|^\alpha w_\lambda(y)P(y)\,dy.
\]
\end{theorem}
\begin{proof}
\leanok
\end{proof}

\begin{definition}[Classical polynomial power-law operator and columns]
\label{def:power-law-polynomial-operator}\uses{def:gegenbauer-polynomial-basis,thm:power-law-recurrence}
\lean{Real.PowerLawOperatorData,Real.PowerLawColumnData,Real.powerLawGegenbauerOperator,Real.powerLawGegenbauerCoefficients}\leanok
For whole-family data with $\lambda>-1/2$ and $\lambda\ne0$, define the
polynomial linear operator $C_n^{(\lambda)}\mapsto q_n$.
An individual coefficient column is the Gegenbauer-basis representation of
$q_n$ and requires only admissible data through degree $n$, together with the
same nondegenerate basis conditions. It does not require later recurrence steps.
\end{definition}
\begin{theorem}[Characterisation of polynomial power-law columns]
\label{thm:power-law-polynomial-operator}\uses{def:power-law-polynomial-operator,thm:coordinate-basis-vectors,thm:sequence-finite-coordinates}
\lean{Real.PowerLawOperatorData.ofIntegrable,Real.PowerLawOperatorData.atDegree,Real.powerLawGegenbauerOperator_gegenbauer,Real.powerLawGegenbauerCoefficients_eq_operator,Real.powerLawGegenbauerCoefficients_eq_zero_of_degree}\leanok
The parameter range $\alpha>-1$, $\lambda>-1/2$, $\lambda\ne0$ satisfies
the whole-family admissibility conditions of
Definition~\ref{def:power-law-polynomial-operator}.
On that definition's full domain, restriction to degree $n$ satisfies the
finite-column conditions. For each $k\in\N$, the polynomial operator sends
$C_n^{(\lambda)}$ to the recurrence polynomial $q_n$, and the finite-column
coefficients are exactly its Gegenbauer coordinate coefficients.
For every admissible finite column, the coefficient at degree $j>n+2k$ is zero.
\end{theorem}
\begin{proof}
\leanok
\end{proof}
\end{supporting}


\begin{theorem}[Exact compatible coefficient operator and bandwidth]
\label{thm:power-law-compatible-operator}\uses{thm:power-law-compatible-polynomials,thm:gegenbauer-orthogonal,def:operators,thm:general-coordinates,thm:orthogonal-lower-degree,thm:power-law-polynomial-operator,thm:power-law-polynomial-functional,thm:power-law-kernel-measure}
\lean{Real.eval_powerLawGegenbauerOperator,Real.integralInner_powerLawGegenbauerOperator_comm,Real.powerLawGegenbauerCoefficients_eq_zero_of_gap,Real.powerLawGegenbauerCoefficients_isBanded}\leanok
Under compatibility as above and $\lambda\ne0$, the polynomial operator
$C_n^{(\lambda)}\mapsto q_n$ agrees with the actual weighted integral on every
polynomial input. Its coefficient matrix in the original Gegenbauer basis
has lower and upper bandwidth $2k$: $A_{j,n}=0$ if $j>n+2k$ or $n>j+2k$.
The operator is symmetric for the Gegenbauer inner product; in this
unnormalized basis this means weighted symmetry, not ordinary matrix symmetry.
Reference: \cite[Section 2]{gutleb-carrillo-olver}.
\end{theorem}
\begin{proof}
\leanok
\end{proof}













\chapter{Toeplitz operators and finite sections}\label{chap:toeplitz}
\rvsummary{Symbols, compact defects and stability of finite sections.}
The first-kind Volterra analysis uses the Toeplitz and finite-section results
collected here. Convergence requires stability in addition to a formula for
the matrix; the theorems state the needed assumptions.

\begin{supporting}[Reference details: Unilateral shifts on real and complex $\ell^2$]
\begin{theorem}[Unilateral shifts on real and complex $\ell^2$]
\label{thm:toeplitz-shifts}\uses{thm:l2-finite-sections}
\lean{lp.norm_sq_eq_tsum_sq,lp.memℓp_cons_zero,lp.forwardShiftL2,lp.forwardShiftL2_zero,lp.forwardShiftL2_succ,lp.backwardShiftL2,lp.forwardShiftL2_single,lp.backwardShiftL2_apply,lp.backwardShiftL2_comp_forwardShiftL2,lp.forwardShiftL2_comp_backwardShiftL2}\leanok
For $\mathbb K=\mathbb R$ or $\mathbb C$, the forward shift $S$ on
$\ell^2(\mathbb N;\mathbb K)$ is an isometry with $(Su)_0=0$ and
$(Su)_{n+1}=u_n$. Its adjoint $B=S^*$ satisfies $(Bu)_n=u_{n+1}$,
$BS=I$ and $SB=I-P_1$, where $P_1$ projects onto coordinate zero.
In particular $Se_n=e_{n+1}$.
Reference: Böttcher--Silbermann, \emph{Introduction to Large Truncated
Toeplitz Matrices}, Chapter 1.
\end{theorem}
\begin{proof}
\leanok
\end{proof}

\rvnote{thm:toeplitz-shifts}{
The construction uses square summability and $\|u\|^2=\sum_n\|u_n\|^2$.
}
\end{supporting}


\begin{theorem}[Bounded Toeplitz operators]
\label{thm:toeplitz-bounded}\uses{thm:toeplitz-shifts}
\lean{lp.continuousLinearMap_ext_single,lp.IsToeplitz,lp.isToeplitz_iff_shift,lp.toeplitzSubmodule,lp.isClosed_isToeplitz,lp.isToeplitz_one,lp.isToeplitz_forwardShiftL2,lp.isToeplitz_backwardShiftL2,lp.IsToeplitz.adjoint,lp.isToeplitz_forwardShiftL2_pow,lp.isToeplitz_backwardShiftL2_pow}\leanok
A bounded operator $T$ on real or complex $\ell^2$ is called Toeplitz when
$(Te_{j+1})_{i+1}=(Te_j)_i$ for all $i,j\ge0$. This is equivalent to
$S^*TS=T$. These operators form a closed linear subspace, preserved by
Hilbert adjoints, containing $I$, $S$, $S^*$ and every nonnegative power
of either shift. Bounded operators are determined by their coordinate columns.
No finite-band or polynomial-symbol assumption is imposed in this definition.
\end{theorem}
\begin{proof}
\leanok
\end{proof}

\begin{theorem}[Laurent-polynomial Toeplitz quantization]
\label{thm:toeplitz-laurent}\uses{thm:toeplitz-bounded,thm:compact-l2-sections}
\lean{lp.laurentShiftL2,lp.isToeplitz_laurentShiftL2,lp.norm_laurentShiftL2_le,lp.laurentToeplitz,lp.laurentToeplitz_apply,lp.norm_laurentToeplitz_le,lp.laurentToeplitz_single,lp.isToeplitz_laurentToeplitz,lp.laurentShiftL2_mul_neg_one,lp.isCompactOperator_laurentShiftL2_defect}\leanok
Put $V_k=S^k$ for $k\ge0$, and $V_k=(S^*)^{-k}$ for $k<0$.
For a Laurent polynomial $p(z)=\sum_k a_kz^k$ over $\mathbb R$ or
$\mathbb C$, define $T_p=\sum_k a_kV_k$. This map is linear,
$T_{az^k}=aV_k$, and $T_p$ is bounded and Toeplitz, with
\[
 \|V_k\|\le1,\qquad \|T_p\|\le\sum_k|a_k|.
\]
The basic product defect is $V_1V_{-1}=I-P_1$; hence
$I-V_1V_{-1}$ is compact.
\end{theorem}
\begin{proof}
\leanok
\end{proof}

\rvnote{thm:toeplitz-laurent}{
This does not assert multiplicativity of the symbol map or a general finite-section stability theorem.
}



\begin{theorem}[Finite-rank Laurent product defects]
\label{thm:toeplitz-product}\uses{thm:toeplitz-laurent}
\lean{lp.forwardShiftL2_pow_apply,lp.forwardShiftL2_pow_apply_of_lt,lp.backwardShiftL2_pow_apply,lp.laurentShiftL2_apply_of_le,lp.laurentShiftL2_single_entry,lp.laurentShiftL2_mul_sub_apply_eq_zero,lp.finiteCoordinateRange,lp.initialL2Projection_comp_eq_of_mem_finiteCoordinateRange,lp.isCompactOperator_of_mem_finiteCoordinateRange,lp.finiteDimensional_range_of_mem_finiteCoordinateRange,lp.laurentToeplitz_mul_sub_mem_finiteCoordinateRange,lp.finiteDimensional_range_laurentToeplitz_mul_sub,lp.isCompactOperator_laurentToeplitz_mul_sub}\leanok
For real or complex Laurent polynomials $p,q$, $T_pT_q-T_{pq}$ has
finite-dimensional range and is compact. More precisely, its output is supported
on finitely many initial coordinates. For individual shifts $V_k,V_l$, all
coordinates $i\ge |k|+|l|$ vanish. The finite-coordinate-range operators form a
linear subspace; an initial coordinate projection fixes every such operator.
\end{theorem}
\begin{proof}
\leanok
\end{proof}


\begin{theorem}[Continuous-symbol Toeplitz operators]
\label{thm:toeplitz-continuous}\uses{thm:toeplitz-bounded,thm:bounded-kernel-l2,thm:mathlib-orthogonal-synthesis}
\lean{AddCircle.orthonormal_nonnegativeFourier,AddCircle.nonnegativeFourierIsometry,AddCircle.nonnegativeFourierIsometry_single,AddCircle.continuousSymbolMultiplier,AddCircle.norm_continuousSymbolMultiplier_le,AddCircle.continuousSymbolMultiplier_ae,AddCircle.continuousToeplitz,AddCircle.norm_continuousToeplitz_le,AddCircle.continuousSymbolMultiplier_toLp,AddCircle.continuousToeplitz_entry,AddCircle.isToeplitz_continuousToeplitz,AddCircle.continuousToeplitz_add,AddCircle.continuousToeplitz_smul,AddCircle.continuousToeplitzCLM,AddCircle.tendsto_continuousToeplitz}\leanok
On a circle of period $T>0$ with normalized Haar measure, let $E$ synthesize
the nonnegative orthonormal Fourier modes into complex $L^2$. For a continuous
complex symbol $f$, multiplication $M_f$ is bounded and
$T_f=E^*M_fE$ satisfies $\|T_f\|\le\|f\|_\infty$ and
$(T_f)_{ij}=\widehat f(i-j)$. This defines a bounded linear map from continuous
symbols to Toeplitz operators; uniform convergence of symbols implies operator-norm
convergence.
\end{theorem}
\begin{proof}
\leanok
\end{proof}



\begin{supporting}[Reference details: Laurent symbols on a circle and related results]
\begin{definition}[Laurent symbols on a circle]
\label{def:laurent-symbol}
\lean{AddCircle.laurentSymbol}\leanok
Let $T\in\mathbb R\setminus\{0\}$ and let $\mathbb R/T\mathbb Z$ be the circle
with signed period $T$. For a Laurent polynomial $p=\sum_{k\in\mathbb Z}p_kz^k$
with finite support, its continuous circle symbol is
\[
 \mathcal L_T(p)(x)=\sum_{k\in\mathbb Z}p_k\exp(2\pi i kx/T).
\]
This defines a complex-linear map from Laurent polynomials over $\mathbb C$ to continuous
complex-valued functions on the circle.
\end{definition}

\begin{theorem}[Algebra of Laurent symbols]
\label{thm:laurent-symbol-algebra}\uses{def:laurent-symbol}
\lean{AddCircle.laurentSymbol_single,AddCircle.laurentSymbol_mul,AddCircle.laurentSymbol_one,AddCircle.laurentSymbol_pow}\leanok
For every nonzero signed period $T$,
$\mathcal L_T(az^k)(x)=a\exp(2\pi i kx/T)$,
$\mathcal L_T(pq)=\mathcal L_T(p)\mathcal L_T(q)$,
$\mathcal L_T(1)=1$, and $\mathcal L_T(p^n)=\mathcal L_T(p)^n$
for every $n\ge0$.
\end{theorem}
\begin{proof}
\leanok
\end{proof}
\end{supporting}



\begin{theorem}[Continuous symbols modulo compact operators]
\label{thm:toeplitz-symbol-product}\uses{thm:toeplitz-product,thm:toeplitz-continuous,thm:laurent-symbol-algebra}
\lean{AddCircle.continuousToeplitz_fourier,AddCircle.continuousToeplitz_laurentSymbol,AddCircle.norm_laurentToeplitz_le_symbol,AddCircle.denseRange_laurentSymbol,AddCircle.isCompactOperator_continuousToeplitz_mul_sub,AddCircle.continuousToeplitz_one,AddCircle.isCompactOperator_add_sub_continuousToeplitz,AddCircle.isCompactOperator_mul_sub_continuousToeplitz,AddCircle.isCompactOperator_pow_sub_continuousToeplitz}\leanok
For $T>0$, Laurent symbols are dense in the continuous complex-valued functions on the circle, and their
Fourier compression agrees with Laurent quantization. Consequently
$T_fT_g-T_{fg}$ is compact for every pair of continuous symbols. If $A-T_f$
and $B-T_g$ are compact, so are $A+B-T_{f+g}$, $AB-T_{fg}$, and
$A^n-T_{f^n}$. The constant-one symbol gives the identity.
\end{theorem}
\begin{proof}
\leanok
\end{proof}


\begin{supporting}[Reference details: Stability under compact perturbations and related results]
\begin{theorem}[Stability under compact perturbations]
\label{thm:stable-compact-perturbation}\uses{thm:compact-projection-approximation}
\lean{ContinuousLinearMap.tendsto_apply_of_strongly_tendsto_of_tendsto,ContinuousLinearMap.tendsto_inverse_apply_of_strongly_tendsto,ContinuousLinearMap.tendsto_compact_comp_of_strongly_tendsto,ContinuousLinearMap.tendsto_inverse_comp_compact_of_strongly_tendsto,ContinuousLinearMap.tendsto_comp_of_strongly_tendsto_of_tendsto_compact,ContinuousLinearMap.eventually_isUnit_add_compact_of_stable,ContinuousLinearMap.tendsto_inverse_add_compact_apply_of_stable,ContinuousLinearMap.eventually_norm_inverse_add_compact_le_of_stable}\leanok
Let $A_n$ and $B$ be invertible bounded operators on a real or complex Banach
space, with $A_n\to B$ strongly and $\sup_n\|A_n^{-1}\|\le C$.
Then $A_n^{-1}\to B^{-1}$ strongly. Uniformly bounded strong convergence
becomes norm convergence after a compact right factor, including norm-convergent
right factors with compact limit. If $K_n\to K$ in norm, $K$ is compact,
and $B+K$ is invertible, then $A_n+K_n$ are eventually invertible with
uniformly bounded inverses. For $f_n\to f$,
$(A_n+K_n)^{-1}f_n\to(B+K)^{-1}f$.
An eventual inverse bound is $(\|(I+B^{-1}K)^{-1}\|+1)C$.
\end{theorem}
\begin{proof}
\leanok
\end{proof}


\begin{theorem}[Coercive orthogonal sections]
\label{thm:coercive-sections}
\lean{ContinuousLinearMap.coerciveFiniteSection,ContinuousLinearMap.inner_coerciveFiniteSection,ContinuousLinearMap.re_inner_coerciveFiniteSection_le,ContinuousLinearMap.isUnit_of_re_inner_le,ContinuousLinearMap.isUnit_coerciveFiniteSection,ContinuousLinearMap.norm_inverse_le_of_re_inner_le,ContinuousLinearMap.norm_inverse_coerciveFiniteSection_le}\leanok
On a real or complex Hilbert space, define $A_P=I-P+PAP$ for a self-adjoint
projection $P$. Its quadratic form splits into
$\|u-Pu\|^2+\langle Pu,APu\rangle$.
If $0<a\le1$ and $\operatorname{Re}\langle u,Au\rangle\ge a\|u\|^2$,
then every $A_P$ is invertible and $\|A_P^{-1}\|\le a^{-1}$.
A positive coercivity bound also makes $A$ itself invertible, with the same
inverse bound.
\end{theorem}
\begin{proof}
\leanok
\end{proof}


\begin{theorem}[Convergence of coercive-plus-compact sections]
\label{thm:coercive-compact-sections}\uses{thm:coercive-sections,thm:stable-compact-perturbation,thm:compact-l2-sections}
\lean{lp.tendsto_coerciveFiniteSection,lp.coerciveFiniteSection_add,lp.isUnit_initialCoerciveFiniteSection,lp.norm_inverse_initialCoerciveFiniteSection_le,lp.eventually_isUnit_coerciveFiniteSection_add_compact,lp.tendsto_coerciveFiniteSection_add_compact_solution,lp.eventually_norm_inverse_coerciveFiniteSection_add_compact_le,lp.existsUnique_coerciveFiniteSection_solution}\leanok
For initial-coordinate projections $P_N$ on $\ell^2$, the sections
$I-P_N+P_NAP_N$ converge strongly to $A$. Suppose $A$ has coercivity
constant $0<a\le1$, $K$ is compact, and $A+K$ is invertible. Then the
sections of $A+K$ are eventually invertible, their inverses are uniformly
bounded, and their solutions for right-hand sides $P_Nf$ converge to
$(A+K)^{-1}f$. Invertibility of a section gives a unique solution of
$P_Nu=u$, $P_N(A+K)u=P_Nf$.
\end{theorem}
\begin{proof}
\leanok
\end{proof}


\begin{theorem}[Positive-real-part symbols]
\label{thm:toeplitz-positive}\uses{thm:toeplitz-continuous}
\lean{AddCircle.re_inner_continuousSymbolMultiplier_le,AddCircle.re_inner_continuousToeplitz_le,AddCircle.isUnit_continuousToeplitz_of_re_pos}\leanok
If $\operatorname{Re}f\ge a$ pointwise, the multiplier and its Toeplitz
compression have quadratic forms bounded below by $a\|u\|^2$.
For $a>0$, $T_f$ is invertible.
\end{theorem}
\begin{proof}
\leanok
\end{proof}
\end{supporting}



\begin{theorem}[Stable continuous Toeplitz-plus-compact sections]
\label{thm:toeplitz-stable-sections}\uses{thm:toeplitz-positive,thm:coercive-compact-sections,thm:compact-injective}
\lean{AddCircle.isUnit_continuousToeplitz_add_of_injective,AddCircle.re_inner_continuousToeplitz_min_le,AddCircle.eventually_isUnit_continuousToeplitz_section,AddCircle.eventually_norm_inverse_continuousToeplitz_section_le,AddCircle.tendsto_continuousToeplitz_section_solution,AddCircle.eventually_existsUnique_continuousToeplitz_section_solution}\leanok
Suppose $f$ is continuous with $\operatorname{Re}f\ge a>0$, $K$ is compact,
and $T_f+K$ is invertible. Its initial-coordinate sections are eventually
invertible and their inverse norms are eventually bounded by
$(\|(I+T_f^{-1}K)^{-1}\|+1)/\min(a,1)$.
The uniquely solvable projected equations converge in $\ell^2$ to the exact
solution for every right-hand side. Invertibility of the full perturbed
operator is an explicit hypothesis here; equivalently, it suffices to prove
injectivity using the compact-perturbation criterion.
\end{theorem}
\begin{proof}
\leanok
\end{proof}


\begin{supporting}[Reference details: Real Laurent operators and complex symbols and related results]
\begin{theorem}[Real Laurent operators and complex symbols]
\label{thm:toeplitz-real-symbol}\uses{thm:toeplitz-symbol-product,thm:toeplitz-positive}
\lean{lp.complexifyL2,lp.complexifyL2_apply,lp.complexifyL2_forwardShift,lp.complexifyL2_backwardShift,lp.complexifyL2_laurentShift,lp.complexifyL2_laurentToeplitz,lp.re_inner_complexifyL2,lp.inner_laurentToeplitz_le}\leanok
Coefficientwise complexification is a real-linear isometry, intertwines both
unilateral shifts and every real Laurent operator, and preserves the real
part of the inner product. Thus a lower bound $\operatorname{Re}p\ge a$
on the complex circle symbol gives
$\langle u,T_pu\rangle_{\mathbb R}\ge a\|u\|^2$ on real $\ell^2$.
\end{theorem}
\begin{proof}
\leanok
\end{proof}


\begin{theorem}[Injective compact perturbations are invertible]
\label{thm:compact-injective}
\lean{ContinuousLinearMap.isUnit_one_add_of_injective_of_compact,ContinuousLinearMap.isUnit_add_of_injective_of_compact}\leanok
On a Banach space over a nontrivially normed field, an injective operator
$I+K$ with $K$ compact is invertible. More generally, if $A$ is invertible
and $K$ is compact, injectivity of $A+K$ implies invertibility.

\end{theorem}
\begin{proof}
\leanok
\end{proof}

\rvnote{thm:compact-injective}{
This is a direct application of the Fredholm alternative; it does not assert injectivity of a particular integral operator.
}


\begin{theorem}[Compact reciprocal affine diagonals]
\label{thm:reciprocal-diagonals}\uses{thm:lp-compact-diagonal}
\lean{lp.norm_reciprocalAffine_le_one,lp.reciprocalAffineDiagonal,lp.reciprocalAffineDiagonal_apply,lp.isCompactOperator_reciprocalAffineDiagonal}\leanok
For integers $a\ge0$, $b>0$, the real $\ell^2$ diagonal with entries
$(an+b)^{-1}$ has entries bounded by one. If $a>0$, this diagonal operator
is compact.
\end{theorem}
\begin{proof}
\leanok
\end{proof}
\end{supporting}



\chapter{First-kind Volterra equations}\label{chap:volterra-first-kind}
\rvsummary{Preconditioning, uniqueness and convergence for polynomial and continuous kernels.}
The integration diagonal provides a preconditioner for the Jacobi operator.
We use it to prove uniqueness and finite-section convergence in the weighted
solution spaces below, under the stated nonvanishing-diagonal conditions.

\begin{definition}[First-order weighted coefficient space]
\label{def:first-order-l2}
\lean{lp.firstOrderL2Submodule,lp.FirstOrderL2,lp.firstOrderL2Equiv,lp.firstOrderL2Isometry}\leanok
Let $\ell^2_1$ be the real sequences $q$ for which $((n+1)q_n)_{n\ge0}$
belongs to $\ell^2$, equipped with norm
\[
 \|q\|_{\ell^2_1}=\|((n+1)q_n)_{n\ge0}\|_{\ell^2}.
\]

Reference: \cite[Definition 5.3]{gutleb-olver}; \cite[Section 2.7]{gutleb-thesis}.
\end{definition}

\begin{theorem}[Degree weighting on the first-order coefficient space]
\label{thm:first-order-l2}\uses{def:first-order-l2}
\lean{lp.firstOrderL2Equiv,lp.firstOrderL2Isometry,lp.norm_firstOrderL2}\leanok
Degree weighting is a linear isometric equivalence $\ell^2_1\to\ell^2$; in particular this is a Banach space.
\end{theorem}
\begin{proof}
\leanok
\end{proof}

\rvnote{def:first-order-l2}{
The norm is on raw coefficients; it is not the norm induced by normalized physical $L^2$ coordinates.
}

\begin{theorem}[Boundedly invertible first-kind integration diagonal]
\label{thm:volterra-first-kind-diagonal}\uses{thm:first-order-l2,thm:volterra-jacobi-diagonal}
\lean{lp.norm_alternating_mul,lp.alternatingL2Isometry,MvPolynomial.volterraJacobiL2Equiv,MvPolynomial.volterraJacobiL2Equiv_apply,MvPolynomial.volterraJacobiL2Equiv_symm_apply,MvPolynomial.volterraJacobiL2Equiv_single,MvPolynomial.existsUnique_volterraJacobiL2_solution}\leanok
The endpoint-factor-free constant-kernel diagonal extends to a linear
isometric equivalence $D:\ell^2\to\ell^2_1$, with
\[
 (Du)_n=\frac{(-1)^n}{n+1}u_n,\qquad
 (D^{-1}q)_n=(-1)^n(n+1)q_n.
\]
It agrees with the existing finite coefficient map on every coordinate
vector. Every $q\in\ell^2_1$ therefore has exactly one solution $u\in\ell^2$
to $Du=q$. The endpoint factor and reflection remain outside this map.
Reference: \cite[Lemma 5.6]{gutleb-olver}; \cite[Lemma 2.3]{gutleb-thesis}.
\end{theorem}
\begin{proof}
\leanok
\end{proof}

\begin{theorem}[Constant-kernel first-kind finite solves]
\label{thm:volterra-first-kind-diagonal-solves}\uses{thm:volterra-first-kind-diagonal,thm:l2-finite-sections}
\lean{MvPolynomial.volterraJacobiFiniteSolution,MvPolynomial.volterraJacobiFiniteSolution_apply,MvPolynomial.volterraJacobiFiniteSolution_equations,MvPolynomial.tendsto_volterraJacobiL2_solution,MvPolynomial.norm_volterraJacobiL2_error_eq_residual}\leanok
For $q\in\ell^2_1$, the truncated solution $u_N=P_ND^{-1}q$ has coefficients
\[
 (u_N)_n=\begin{cases}(-1)^n(n+1)q_n&n<N,\\0&n\ge N.\end{cases}
\]
It satisfies $(Du_N)_n=q_n$ for every $n<N$ and converges in raw $\ell^2$
to the exact solution. Every candidate $v\in\ell^2$ has the exact error identity
\[
 \|v-D^{-1}q\|_{\ell^2}=\|Dv-q\|_{\ell^2_1}.
\]

\end{theorem}
\begin{proof}
\leanok
\end{proof}

\rvnote{thm:volterra-first-kind-diagonal-solves}{
This is the constant-kernel diagonal problem, not yet the general polynomial-kernel solver.
}

\begin{supporting}[Reference details: Raw shifted Jacobi multiplication entries and related results]
\begin{theorem}[Raw shifted Jacobi multiplication entries]
\label{thm:shifted-jacobi-multiplication}\uses{thm:jacobi-recurrence,def:unit-jacobi-measure}
\lean{Polynomial.shiftedJacobiOneZeroUpper,Polynomial.shiftedJacobiOneZeroDiag,Polynomial.shiftedJacobiOneZeroLower,Polynomial.X_mul_shiftedJacobi_one_zero_succ,Polynomial.X_mul_shiftedJacobi_one_zero_zero}\leanok
Writing $Q_n(x)=P_n^{(1,0)}(2x-1)$, multiplication by $x$ has upper,
diagonal and lower column coefficients
\[
 a_n=\frac{n+2}{2(2n+3)},\qquad
 b_n=\frac12-\frac1{2(2n+1)(2n+3)},\qquad
 c_n=\frac n{2(2n+1)}.
\]
For $n\ge1$, $xQ_n=a_nQ_{n+1}+b_nQ_n+c_nQ_{n-1}$; the constant
column is $xQ_0=a_0Q_1+b_0Q_0$.
\end{theorem}
\begin{proof}
\leanok
\end{proof}


\begin{theorem}[Laurent symbols modulo compact operators]
\label{thm:toeplitz-laurent-algebra}\uses{thm:toeplitz-product}
\lean{lp.laurentToeplitz_one,lp.isCompactOperator_add_sub_laurentToeplitz,lp.isCompactOperator_mul_sub_laurentToeplitz,lp.isCompactOperator_pow_sub_laurentToeplitz}\leanok
Over real or complex scalars, $T_1=I$. If $A-T_p$ and $B-T_q$ are compact,
then $A+B-T_{p+q}$, $AB-T_{pq}$, and $A^n-T_{p^n}$ are compact.
This keeps the real coefficient convention available without complexifying
the underlying integral equation.
\end{theorem}
\begin{proof}
\leanok
\end{proof}


\begin{theorem}[Finite coefficients in square-summable space]
\label{thm:finite-l2-embedding}
\lean{lp.finsuppL2,lp.finsuppL2_apply,lp.finsuppL2_single,lp.finsuppL2_injective}\leanok
The canonical linear embedding of finitely supported real or complex
sequences into $\ell^2$ preserves every coordinate and every coordinate vector,
and is injective.
\end{theorem}
\begin{proof}
\leanok
\end{proof}


\begin{theorem}[Bounded raw Jacobi multiplication]
\label{thm:raw-jacobi-l2}\uses{thm:reciprocal-diagonals,thm:shifted-jacobi-multiplication,thm:finite-l2-embedding,thm:toeplitz-shifts,thm:unit-jacobi-coordinates}
\lean{Polynomial.shiftedJacobiToeplitzPart,Polynomial.shiftedJacobiOneZeroCorrection,Polynomial.shiftedJacobiOneZeroL2,Polynomial.isCompactOperator_shiftedJacobiOneZeroCorrection,Polynomial.reflectedJacobiOneZeroCorrection,Polynomial.reflectedJacobiOneZeroL2,Polynomial.isCompactOperator_reflectedJacobiOneZeroCorrection,Polynomial.shiftedJacobiOneZeroL2_apply_zero,Polynomial.shiftedJacobiOneZeroL2_apply_succ,Polynomial.shiftedJacobiOneZeroL2_single_zero,Polynomial.shiftedJacobiOneZeroL2_single_succ,Polynomial.shiftedJacobiOneZeroCoefficients,Polynomial.shiftedJacobiOneZeroCoefficients_basis,Polynomial.shiftedJacobiOneZeroL2_coefficients_basis,Polynomial.shiftedJacobiOneZeroL2_coefficients}\leanok
Let $J_0=\tfrac12 I+\tfrac14(S+S^*)$. The raw shifted Jacobi $(1,0)$
multiplication matrix extends to a bounded real $\ell^2$ operator $J=J_0+C_J$
with $C_J$ compact. Its columns are exactly the coefficients $a_n,b_n,c_n$
of the shifted Jacobi recurrence, and the embedding of every polynomial
satisfies $J\,\operatorname{coeff}(p)=\operatorname{coeff}(xp)$.
There is a bounded tridiagonal operator $H=J_0+C_H$ with $C_H$ compact.
\end{theorem}
\begin{proof}
\leanok
\end{proof}

\rvnote{thm:raw-jacobi-l2}{
The reflected preconditioner identity for $H$ is Theorem~\ref{thm:raw-jacobi-preconditioner}.
}


\begin{theorem}[Exact reflected Jacobi preconditioning]
\label{thm:raw-jacobi-preconditioner}\uses{thm:raw-jacobi-l2,thm:volterra-first-kind-diagonal}
\lean{Polynomial.reflectedJacobiOneZeroL2_apply_zero,Polynomial.reflectedJacobiOneZeroL2_apply_succ,Polynomial.volterraRawDiagonal,Polynomial.volterraRawDiagonal_apply,Polynomial.volterraRawDiagonal_eq_weighted,Polynomial.volterraRawDiagonal_mul_reflectedJacobiOneZeroL2}\leanok
The bounded raw diagonal $D_n=(-1)^n/(n+1)$ agrees coefficientwise with
the existing equivalence into weighted $\ell^2_1$. The second Jacobi operator
satisfies $DH=(I-J)D$. Its upper, diagonal and lower column entries are
\[
 \frac{(n+2)^2}{2(n+1)(2n+3)},\qquad
 \frac12+\frac1{2(2n+1)(2n+3)},\qquad
 \frac{n^2}{2(n+1)(2n+1)}.
\]
Thus the reflected expression is $D^{-1}(I-J)D$; the statement does not treat
$D^{-1}$ as a bounded endomorphism of raw $\ell^2$.
\end{theorem}
\begin{proof}
\leanok
\end{proof}


\begin{theorem}[The common Jacobi Laurent symbol]
\label{thm:raw-jacobi-symbol}\uses{thm:toeplitz-laurent-algebra,thm:raw-jacobi-l2,thm:laurent-symbol-algebra}
\lean{Polynomial.shiftedJacobiLaurentSymbol,Polynomial.laurentToeplitz_shiftedJacobiLaurentSymbol,Polynomial.isCompactOperator_shiftedJacobiOneZeroL2_sub_symbol,Polynomial.isCompactOperator_reflectedJacobiOneZeroL2_sub_symbol,Polynomial.isCompactOperator_reflectedJacobi_pow_mul_sub_symbol,Polynomial.shiftedJacobiLaurentSymbol_circle,Polynomial.shiftedJacobiLaurentSymbol_circle_mem_Icc}\leanok
Both $J$ and $H$ differ compactly from $T_t$, where
$t(z)=(2+z+z^{-1})/4$. Consequently $H^rJ^s-T_{t^{r+s}}$ is compact.
On a circle with nonzero signed period the symbol is real, $t(z)=(1+\operatorname{Re}z)/2\in[0,1]$.
\end{theorem}
\begin{proof}
\leanok
\end{proof}


\begin{theorem}[Polynomial Volterra factors and diagonal symbol]
\label{thm:volterra-diagonal-toeplitz}\uses{thm:raw-jacobi-symbol}
\lean{MvPolynomial.volterraPreconditionedL2,MvPolynomial.volterraDiagonalLaurentSymbol,MvPolynomial.isCompactOperator_volterraPreconditionedL2_sub_symbol,MvPolynomial.volterraDiagonalLaurentSymbol_eq_eval₂,MvPolynomial.volterraDiagonalLaurentSymbol_circle}\leanok
For a polynomial kernel $K(x,y)=\sum_{r,s}k_{rs}x^ry^s$, define the bounded
ordered coefficient operator $A_K=\sum_{r,s}k_{rs}H^rJ^s$ and the Laurent
symbol $p_K=K(t,t)$. Then $A_K-T_{p_K}$ is compact. On a circle with nonzero signed period,
$p_K(z)=K(\tau,\tau)$ where $\tau=(1+\operatorname{Re}z)/2\in[0,1]$.

\end{theorem}
\begin{proof}
\leanok
\end{proof}

\rvnote{thm:volterra-diagonal-toeplitz}{
The integral realization is given in Theorem~\ref{thm:volterra-raw-l2-agreement}.
}


\begin{theorem}[Conditional first-kind coefficient stability]
\label{thm:volterra-conditional-toeplitz-sections}\uses{thm:volterra-diagonal-toeplitz,thm:compact-injective,thm:coercive-compact-sections,thm:toeplitz-real-symbol}
\lean{MvPolynomial.inner_volterraDiagonalToeplitz_le,MvPolynomial.isUnit_volterraPreconditionedL2_of_injective,MvPolynomial.eventually_isUnit_volterraPreconditionedL2_section,MvPolynomial.tendsto_volterraPreconditionedL2_section_solution}\leanok
Assume $K(x,x)\ge a>0$ on $[0,1]$ and that the ordered coefficient operator
$A_K$ is injective. The Toeplitz part has quadratic form bounded below by
$a\|u\|^2$, and $A_K$ is invertible. Its leading sections
$I-P_N+P_NA_KP_N$ are eventually invertible, and their solutions for
right-hand sides $P_Ng$ converge in raw $\ell^2$ to $A_K^{-1}g$.

\end{theorem}
\begin{proof}
\leanok
\end{proof}

\rvnote{thm:volterra-conditional-toeplitz-sections}{
Theorem~\ref{thm:volterra-nonvanishing-sections} removes the injectivity hypothesis for polynomial kernels with nonvanishing diagonal.
}


\begin{theorem}[Differentiating polynomial Volterra kernels]
\label{thm:volterra-polynomial-differentiation}\uses{def:volterra-continuous-kernel,thm:volterra-integral-input}
\lean{MvPolynomial.polynomialVolterraIntegral_add,MvPolynomial.polynomialVolterraIntegral_mul_X_zero,MvPolynomial.polynomialVolterraIntegral_mul_X_one,MvPolynomial.continuous_polynomialVolterraIntegral,MvPolynomial.ae_hasDerivAt_polynomialVolterraIntegral}\leanok
For a real polynomial kernel $K(x,y)$, write the classical oriented integral
$V_K f(x)=\int_0^xK(x,y)f(y)\,dy$. If $f$ is locally integrable on $\mathbb R$,
then $V_K f$ is continuous and, almost everywhere,
\[
 (V_K f)'(x)=K(x,x)f(x)+V_{\partial_x K}f(x).
\]
Polynomial kernel slices times $f$ are interval-integrable. Kernel addition
commutes with integration; multiplying the kernel by $x$ multiplies the output
by $x$, and multiplying it by $y$ multiplies the input by its coordinate.
The algebraic identities also hold at any fixed endpoint whenever their
row products are integrable; the factors need not be separately integrable.
Reference: \cite{linz-volterra}, first-kind equations with differentiable kernels.
\end{theorem}
\begin{proof}
\leanok
\end{proof}
\end{supporting}


\begin{theorem}[First-kind uniqueness for integrable inputs]
\label{thm:volterra-first-kind-uniqueness}\uses{thm:volterra-polynomial-differentiation}
\lean{Real.eq_zero_of_norm_le_volterra_integral,MvPolynomial.ae_diagonal_mul_eq_neg_polynomialVolterraIntegral,MvPolynomial.eq_zero_of_polynomialVolterraIntegral_eq_zero,MvPolynomial.polynomialVolterraIntegral_congr_ae_on,MvPolynomial.ae_eq_zero_of_polynomialVolterraIntegral_eq_zero,MvPolynomial.ae_eq_zero_of_integrableOn_polynomialVolterraIntegral_eq_zero}\leanok
Let $b>0$, let $K$ be a real polynomial kernel, and suppose
$K(x,x)\ne0$ on $[0,b]$. For every $f\in L^1([0,b])$,
\[
 (\forall x\in[0,b],\ V_K f(x)=0)\quad\Longrightarrow\quad f=0
 \text{ almost everywhere on }[0,b].
\]
The diagonal may have either sign. The same conclusion is proved for locally
integrable inputs; globally continuous inputs vanish pointwise on the interval.
A zero output satisfies
$K(x,x)f(x)=-V_{\partial_xK}f(x)$ almost everywhere. Volterra integration on any
$[0,x]\subseteq[0,b]$ respects almost-everywhere equality of inputs
whenever the kernel--input products are integrable on that row.
The supporting integral inequality states that a globally continuous real $g$ with
$\|g(x)\|\le M\int_0^x\|g(y)\|\,dy$ on $[0,b]$ vanishes there.
\end{theorem}
\begin{proof}
\leanok
\end{proof}

\begin{supporting}[Reference details: Injectivity of exact finite Proriol assembly and related results]
\begin{theorem}[Injectivity of exact finite Proriol assembly]
\label{thm:volterra-finite-injectivity}\uses{thm:volterra-first-kind-uniqueness,thm:volterra-coefficient-exact,thm:unit-jacobi-coordinates}
\lean{MvPolynomial.eq_zero_of_volterraProriolCoefficients_eq_zero,MvPolynomial.injective_volterraProriolCoefficients}\leanok
If $K(x,x)\ne0$ on some interval $[0,b]$ with $b>0$, then the exact
Proriol Volterra coefficient map is injective on finitely supported shifted
Jacobi $(1,0)$ inputs. In particular, a finite input with zero assembled output
is zero. This concerns the full finite-support map, not injectivity of every
square truncation or of its extension to infinite raw coefficient sequences.
\end{theorem}
\begin{proof}
\leanok
\end{proof}

\begin{theorem}[Raw shifted Jacobi synthesis]
\label{thm:raw-jacobi-synthesis}\uses{thm:unit-jacobi-completeness,thm:finite-l2-embedding,thm:lp-compact-diagonal,thm:jacobi-norm}
\lean{Polynomial.integral_sq_shiftedJacobi_one_zero,Polynomial.norm_shiftedJacobiOneZeroScale_le,Polynomial.shiftedJacobiOneZeroSynthesis,Polynomial.shiftedJacobiOneZeroSynthesis_single,Polynomial.shiftedJacobiOneZeroSynthesis_injective,Polynomial.shiftedJacobiOneZeroSynthesis_finsupp,Polynomial.tendsto_shiftedJacobiOneZeroSynthesis_projection}\leanok
For $\mu=\mathrm{unitJacobiMeasure}(1,0)$ the raw shifted Jacobi polynomial
$\widetilde P_n^{(1,0)}$ has squared norm $2/(n+1)$. Scaling normalized
Hilbert coordinates by $\sqrt{2/(n+1)}\le\sqrt2$ defines a bounded injective
synthesis map $S:\ell^2\to L^2(\mu)$. It sends each coordinate vector to
its raw polynomial, agrees with finite polynomial synthesis, and sends
coefficient truncations to a sequence converging to $Sc$ in $L^2(\mu)$.

\end{theorem}
\begin{proof}
\leanok
\end{proof}

\rvnote{thm:raw-jacobi-synthesis}{No bounded inverse on its range or surjectivity is asserted.}


\begin{theorem}[Ordered Volterra kernel algebra]
\label{thm:volterra-ordered-algebra}\uses{thm:volterra-diagonal-toeplitz}
\lean{MvPolynomial.volterraPreconditionedLinear,MvPolynomial.volterraPreconditionedLinear_apply,MvPolynomial.volterraPreconditionedL2_add,MvPolynomial.volterraPreconditionedL2_monomial,MvPolynomial.volterraPreconditionedL2_C,MvPolynomial.volterraPreconditionedL2_mul_X_zero,MvPolynomial.volterraPreconditionedL2_mul_X_one}\leanok
Ordered polynomial quantization $A_K=\sum_{i,j}K_{ij}H^iJ^j$ is linear in
$K$, has $A_a=aI$, and obeys $A_{xK}=HA_K$ and $A_{yK}=A_KJ$.
A monomial $ax^iy^j$ gives exactly $aH^iJ^j$; the order of these factors
is retained.
\end{theorem}
\begin{proof}
\leanok
\end{proof}

\begin{theorem}[Raw operators recover finite Proriol assembly]
\label{thm:volterra-raw-finite-agreement}\uses{thm:volterra-ordered-algebra,thm:volterra-coefficient-exact,thm:volterra-polynomial-differentiation,thm:raw-jacobi-preconditioner}
\lean{Polynomial.shiftedJacobiOneZeroCoefficients_linearCombination,Polynomial.volterraRawDiagonal_finsupp,MvPolynomial.volterraProriolCore_mul_X_zero,MvPolynomial.volterraProriolCoefficients_raw_mul_X_zero,MvPolynomial.volterraProriolCore_mul_X_one,MvPolynomial.volterraProriolCoefficients_add_kernel,MvPolynomial.volterraProriolCoefficients_smul_kernel,MvPolynomial.volterraProriolCoefficients_C,MvPolynomial.volterraRawDiagonal_preconditioned_finsupp}\leanok
Let $V_K$ denote the exact finite Proriol coefficient assembly and $\iota$
the embedding of finite coefficients into raw $\ell^2$. Then
\[
 D A_K\iota c=\iota(V_Kc).
\]
The raw coefficient map recovers $c$ from its polynomial synthesis, and
$D\iota c$ agrees with the finite integration diagonal. Assembly is linear
in the kernel and sends a constant kernel $a$ to $aD$. Multiplication of
the kernel by $x$ multiplies the reflected output polynomial by $1-X$
(and its raw coefficient vector by $I-J$); multiplication by $y$ replaces
the input polynomial by $X$ times that polynomial.
\end{theorem}
\begin{proof}
\leanok
\end{proof}

\begin{theorem}[Bounded kernels with prescribed marginals]
\label{thm:marginal-kernel-l2}\uses{thm:bounded-kernel-l2}
\lean{MeasureTheory.marginalKernelL2Operator,MeasureTheory.inner_marginalKernelL2Operator,MeasureTheory.norm_marginalKernelL2Operator_le,MeasureTheory.integrable_marginalKernelL2_pair}\leanok
Let a joint measure $\kappa$ have coordinate marginals $\mu$ and $\nu$.
If a real kernel $k$ is almost-everywhere strongly measurable and
$|k|\le M$ almost everywhere, with $M\ge0$, then it defines a bounded
operator $B:L^2(\nu)\to L^2(\mu)$ with $\|B\|\le M$ and
\[
 \langle g,Bf\rangle=\int k(x,y)g(x)f(y)\,d\kappa(x,y).
\]
The displayed integrand is absolutely integrable for every pair of inputs.
\end{theorem}
\begin{proof}
\leanok
\end{proof}

\begin{theorem}[The reflected Volterra joint measure]
\label{thm:jacobi-volterra-joint}\uses{thm:unit-jacobi-measure-characterisation,thm:jacobi-moments,thm:unit-jacobi-weight,thm:jacobi-weight-basic}
\lean{Real.reflectedVolterraTriangleMeasure,Real.map_fst_reflectedVolterraTriangleMeasure,Real.map_snd_reflectedVolterraTriangleMeasure,Polynomial.unitJacobiMeasure_one_zero_triangle,Real.jacobiVolterraJoint,Real.map_fst_jacobiVolterraJoint,Real.map_snd_jacobiVolterraJoint,Real.ae_mem_jacobiVolterraJoint,Real.integral_jacobiVolterraJoint,Polynomial.integrable_unitJacobiMeasure_one_zero,Polynomial.integral_unitJacobiMeasure_one_zero}\leanok
Area measure on $\{(s,y)\in[0,1]^2:s+y\le1\}$ has both coordinate
marginals $(1-t)\,dt$. Thus the joint measure $\kappa$ obtained by multiplying
this area measure by four has both marginals
$\mu=\mathrm{unitJacobiMeasure}(1,0)=4(1-t)\,dt$. It is supported on that
simplex and, for every $\kappa$-integrable real function $F$,
\[
 \int F\,d\kappa=4\int_0^1\int_0^{1-s}F(s,y)\,dy\,ds.
\]
For a real integrable input $f$ against $\mu$, the function
$t\mapsto(1-t)f(t)$ is integrable on $[0,1]$, and
$\int f\,d\mu=4\int_0^1(1-t)f(t)\,dt$.
\end{theorem}
\begin{proof}
\leanok
\end{proof}

\begin{theorem}[Local integrability in raw Jacobi function space]
\label{thm:jacobi-volterra-local-integrability}\uses{thm:jacobi-volterra-joint}
\lean{Polynomial.ae_unitJacobiMeasure_one_zero_iff,Polynomial.integrableOn_coe_jacobiOneZeroLp}\leanok
The measure $\mu=\mathrm{unitJacobiMeasure}(1,0)$ and Lebesgue measure
restricted to $[0,1]$ have the same null sets. Every $f\in L^2(\mu)$ is
Lebesgue integrable on $[0,b]$ for every $b<1$.
\end{theorem}
\begin{proof}
\leanok
\end{proof}

\begin{theorem}[Classical rows for weighted Jacobi inputs]
\label{thm:jacobi-volterra-input}\uses{thm:jacobi-volterra-local-integrability,thm:volterra-integral-input}
\lean{Real.VolterraIntegralInput.ofJacobiOneZeroLp}\leanok
For a continuous kernel on $\Delta$ and $f\in L^2(\mathrm{unitJacobiMeasure}(1,0))$,
the product $K(x,\cdot)f$ is integrable on every row with $0\le x<1$.
Thus $V_Kf(x)$ is the integral of this product over $(0,x)$, with the zero row equal to zero.
\end{theorem}
\begin{proof}
\leanok
\end{proof}

\begin{theorem}[Weighted reflected Volterra realization]
\label{thm:jacobi-volterra-l2}\uses{thm:jacobi-volterra-input,thm:marginal-kernel-l2,thm:general-polynomial-L2-basic,def:volterra-continuous-kernel}
\lean{MvPolynomial.ae_norm_reflectedVolterraKernel_le,MvPolynomial.jacobiVolterraL2Operator,MvPolynomial.norm_jacobiVolterraL2Operator_le,MvPolynomial.inner_jacobiVolterraL2Operator,MvPolynomial.jacobiVolterraL2Operator_polynomial}\leanok
For a real polynomial kernel $K$ with $|K|\le M$ on the physical unit
triangle, $M\ge0$, the reflected joint kernel $K(1-s,y)$ defines a bounded
operator $B_K:L^2(\mu)\to L^2(\mu)$, $\mu=\mathrm{unitJacobiMeasure}(1,0)$,
with $\|B_K\|\le M$. For all inputs $f,g$,
\[
 \langle g,B_Kf\rangle
 =4\int_{(0,1)} g(s)V_Kf(1-s)\,ds.
\]
If polynomials $p,q$ satisfy $V_Kp(x)=xq(1-x)$ on $[0,1]$, then
$B_K[p]=[q]$ in weighted $L^2$. Thus the polynomial output retains both
the reflection and the division by the endpoint factor.
\end{theorem}
\begin{proof}
\leanok
\end{proof}

\begin{theorem}[Infinite raw coefficient agreement]
\label{thm:volterra-raw-l2-agreement}\uses{thm:jacobi-volterra-l2,thm:raw-jacobi-synthesis,thm:volterra-raw-finite-agreement}
\lean{MvPolynomial.jacobiVolterraL2Operator_raw_synthesis}\leanok
For every polynomial kernel $K$ and triangle bound $M\ge0$, raw Jacobi
synthesis $S$ intertwines the full ordered coefficient operator and the
weighted integral realization:
\[
 S D A_K=B_K S.
\]
This is an equality of bounded operators on all of raw $\ell^2$, not only
on finite coefficient sequences.
\end{theorem}
\begin{proof}
\leanok
\end{proof}



\begin{theorem}[Continuous kernels in the weighted reflected realization]
\label{thm:continuous-jacobi-volterra}\uses{thm:volterra-l2,thm:jacobi-volterra-l2}
\lean{Real.ae_reflected_mem_volterraTriangle,Real.jacobiVolterraL2Operator,Real.norm_jacobiVolterraL2Operator_le,Real.integrable_jacobiVolterra_pair,Real.inner_jacobiVolterraL2Operator_joint,Real.inner_jacobiVolterraL2Operator,Real.jacobiVolterraL2Operator_polynomial,Real.jacobiVolterraL2Operator_sub,Real.norm_jacobiVolterraL2Operator_sub_le}\leanok
Let $\Delta=\{(x,y):0\le y\le x\le1\}$ and $\mu(ds)=4(1-s)\,ds$ on $[0,1]$.
For a real kernel $K$ continuous on $\Delta$ and $|K|\le M$, $M\ge0$,
the averaged reflected operator $B_K$ is bounded on $L^2(\mu)$, with
$\|B_K\|\le M$. Its weak formula is
\[
 \langle g,B_Kf\rangle=4\int_{(0,1)}g(s)V_Kf(1-s)\,ds.
\]
Each inner Volterra integral exists for $0<s<1$; no endpoint value is needed.
The joint-measure pairing is integrable, and reflection lands in $\Delta$
almost everywhere for that measure. This construction agrees with the polynomial
one. Subtraction of kernels subtracts operators; if $|K-L|\le\delta$ on
$\Delta$, $\delta\ge0$, then $\|B_K-B_L\|\le\delta$.
\end{theorem}
\begin{proof}
\leanok
\end{proof}

\begin{theorem}[Continuous Volterra primitives of integrable inputs]
\label{thm:continuous-volterra-primitives}\uses{thm:volterra-polynomial-differentiation,thm:volterra-finite-approximation}
\lean{Real.intervalIntegrable_continuousVolterra_mul,Real.norm_volterraIntegral_sub_le_of_integrable,Real.tendstoUniformly_volterraIntegral_of_integrable,Real.continuous_volterraIntegral_of_integrable}\leanok
For $K$ continuous on $\Delta$ and $f\in L^1(\mathbb R)$, every slice
$K(x,\cdot)f$ is integrable on $[0,x]$, $0\le x\le1$, and $V_Kf$ is continuous
on $[0,1]$. If $K,L$ are continuous and $|K-L|\le e$, $e\ge0$, on $\Delta$, then
\[
 |V_Kf(x)-V_Lf(x)|\le e\int_0^1|f(y)|\,dy.
\]
Consequently, uniformly convergent polynomial kernel approximants give uniformly
convergent primitives for each such input.
\end{theorem}
\begin{proof}
\leanok
\end{proof}
\end{supporting}


\begin{theorem}[First-kind uniqueness from absolutely continuous slices]
\label{thm:volterra-ac-uniqueness}\uses{thm:volterra-first-kind-uniqueness,thm:integrable-derivative}
\lean{Real.volterraIntegral_congr_ae_on,Real.primitive_eq_zero_of_volterraIntegral_eq_zero,Real.ae_eq_zero_of_volterraIntegral_eq_zero,Real.ae_eq_zero_of_integrableOn_volterraIntegral_eq_zero}\leanok
Let $f$ be integrable on $[0,b]$. Suppose each slice $y\mapsto K(x,y)$ is
absolutely continuous on $[0,x]$ and $K(x,x)\ne0$. Suppose that for one real
constant $C$, each row has an integrable derivative representative $h_x$ with
\[
 |h_x(y)|\le C|K(x,x)|\quad\text{for almost every }y\in[0,x],\qquad 0\le x\le b.
\]
The condition is independent of the derivative representative.
If $V_Kf$ vanishes throughout $[0,b]$, then $f=0$ almost everywhere there.
For locally integrable $f$, its primitive $\int_0^x f$ also vanishes throughout
that interval. Changing the input on a null subset of $[0,b]$ does not change
its Volterra integrals on the interval.
\end{theorem}
\begin{proof}
\leanok
\end{proof}

\begin{supporting}[Reference details: Injectivity of the continuous weighted Volterra realization and related results]
\begin{theorem}[Injectivity of the continuous weighted Volterra realization]
\label{thm:continuous-volterra-injective}\uses{thm:continuous-jacobi-volterra,thm:continuous-volterra-primitives,thm:volterra-ac-uniqueness}
\lean{Real.volterraIntegral_eq_zero_of_jacobiVolterraL2Operator_eq_zero,Real.eq_zero_of_jacobiVolterraL2Operator_eq_zero,Real.injective_jacobiVolterraL2Operator}\leanok
Let $K$ be continuous and bounded on $\Delta$. If $B_Kf=0$, then $V_Kf=0$
on every $[0,b]$ with $0<b<1$.
If additionally the absolutely continuous slice, derivative-bound and
nonvanishing-diagonal hypotheses of Theorem~\ref{thm:volterra-ac-uniqueness}
hold on $[0,1]$, then $B_K$ is injective on $L^2(\mu)$.
\end{theorem}
\begin{proof}
\leanok
\end{proof}

\begin{theorem}[Injectivity of the weighted reflected Volterra operator]
\label{thm:jacobi-volterra-injectivity}\uses{thm:continuous-volterra-injective}
\lean{MvPolynomial.polynomialVolterraIntegral_eq_zero_of_jacobiVolterraL2Operator_eq_zero,MvPolynomial.eq_zero_of_jacobiVolterraL2Operator_eq_zero,MvPolynomial.injective_jacobiVolterraL2Operator}\leanok
For a polynomial kernel with $K(x,x)\ne0$ on $[0,1]$, the weighted
operator $B_K$ is injective on $L^2(\mathrm{unitJacobiMeasure}(1,0))$.
In particular a zero weighted output forces a zero input. Independently
of the diagonal condition, $B_Kf=0$ forces $V_Kf(x)=0$ throughout every
$[0,b]$, $0<b<1$.
\end{theorem}
\begin{proof}
\leanok
\end{proof}

\begin{theorem}[Genuine raw injectivity and positive-diagonal convergence]
\label{thm:volterra-raw-injective-sections}\uses{thm:volterra-raw-l2-agreement,thm:jacobi-volterra-injectivity,thm:volterra-conditional-toeplitz-sections}
\lean{MvPolynomial.injective_volterraPreconditionedL2,MvPolynomial.isUnit_volterraPreconditionedL2_of_pos,MvPolynomial.eventually_isUnit_volterraPreconditionedL2_section_of_pos,MvPolynomial.tendsto_volterraPreconditionedL2_section_solution_of_pos}\leanok
If $K(x,x)\ne0$ on $[0,1]$, the ordered operator $A_K$ is injective
on raw $\ell^2$. If moreover $K(x,x)\ge a>0$ there, then $A_K$ is
invertible and its leading finite sections
$A_{K,N}=I-P_N+P_NA_KP_N$ are eventually invertible, with
\[
 A_{K,N}^{-1}P_Ng\longrightarrow A_K^{-1}g
 \quad\hbox{in raw }\ell^2.
\]
No separate injectivity assumption is required.
\end{theorem}
\begin{proof}
\leanok
\end{proof}

\rvnote{thm:volterra-raw-injective-sections}{
Theorem~\ref{thm:volterra-nonvanishing-sections} also covers a negative diagonal.
}

\begin{theorem}[Sign changes in finite sections]
\label{thm:signed-finite-sections}\uses{thm:coercive-sections}
\lean{ContinuousLinearMap.coerciveFiniteSection_neg,ContinuousLinearMap.isUnit_coerciveFiniteSection_neg_iff,ContinuousLinearMap.inverse_coerciveFiniteSection_neg_apply,ContinuousLinearMap.inverse_neg_apply_neg,ContinuousLinearMap.neg_apply_eq_neg_iff,ContinuousLinearMap.coerciveFiniteSection_neg_apply_eq_iff}\leanok
For a bounded projection $P$ on a real inner-product space and a bounded
operator $A$, put $A_P=I-P+PAP$. Then
\[
 (-A)_P=(I-2P)A_P.
\]
Consequently the two sections are invertible simultaneously. On this
invertible domain, $(-A)_P^{-1}(-Pg)=A_P^{-1}Pg$. For invertible $A$,
likewise $(-A)^{-1}(-g)=A^{-1}g$. For every $A$, including singular
operators and sections,
\[
 (-A)_P u=-Pg\iff A_Pu=Pg,\qquad (-A)u=-g\iff Au=g.
\]
\end{theorem}
\begin{proof}
\leanok
\end{proof}
\end{supporting}


\begin{theorem}[Polynomial first-kind convergence for a nonvanishing diagonal]
\label{thm:volterra-nonvanishing-sections}\uses{thm:volterra-raw-injective-sections,thm:signed-finite-sections}
\lean{MvPolynomial.exists_pos_volterra_diagonal_bound_or_neg,MvPolynomial.volterraPreconditionedL2_neg,MvPolynomial.isUnit_volterraPreconditionedL2,MvPolynomial.eventually_isUnit_volterraPreconditionedL2_section_of_nonvanishing,MvPolynomial.tendsto_volterraPreconditionedL2_section_solution_of_nonvanishing}\leanok
Let $K$ be a real polynomial kernel with $K(x,x)\ne0$ for all $x\in[0,1]$.
Then its raw ordered operator $A_K$ is invertible, its leading finite
sections $A_{K,N}=I-P_N+P_NA_KP_N$ are eventually invertible, and for
every raw $\ell^2$ right-hand side $g$,
\[
 A_{K,N}^{-1}P_Ng\longrightarrow A_K^{-1}g\quad\hbox{in }\ell^2.
\]
The real diagonal has a uniform positive bound after possibly negating
$K$, and $A_{-K}=-A_K$. Either sign is covered, without an additional
injectivity assumption.
\end{theorem}
\begin{proof}
\leanok
\end{proof}


\begin{supporting}[Reference details: Weighted coefficient operator and classical solution and related results]
\begin{definition}[Weighted coefficient operator and classical solution]
\label{def:volterra-weighted-solution}\uses{def:invertible-index,thm:volterra-diagonal-toeplitz,thm:volterra-first-kind-diagonal}
\lean{MvPolynomial.volterraWeightedL2,MvPolynomial.volterraWeightedSolution}\leanok
For a real polynomial kernel $K$, let $A_K$ be its preconditioned coefficient operator and $D:\ell^2\to\ell^2_1$ the isometric integration diagonal. Define $W_K=DA_K$. For a kernel in the domain where $A_K$ is invertible and weighted data $q\in\ell^2_1$, define $u_K(q)=A_K^{-1}D^{-1}q$.
\end{definition}

\begin{theorem}[Equation for the classical weighted solution]
\label{thm:volterra-weighted-solution}\uses{def:volterra-weighted-solution,thm:invertible-index}
\lean{MvPolynomial.volterraWeightedSolution_property}\leanok
For every kernel in the invertible domain and every $q\in\ell^2_1$, $W_Ku_K(q)=q$.
\end{theorem}
\begin{proof}
\leanok
\end{proof}

\begin{definition}[Classical weighted finite solution]
\label{def:volterra-weighted-finite-solution}\uses{def:invertible-index,thm:volterra-diagonal-toeplitz,thm:coercive-sections,thm:volterra-first-kind-diagonal}
\lean{MvPolynomial.volterraWeightedSection,MvPolynomial.volterraWeightedFiniteSolution}\leanok
For a real polynomial kernel $K$, set $S_N(A_K)=I-P_N+P_NA_KP_N$. At a cutoff in the domain where $S_N(A_K)$ is invertible, define $u_{K,N}(q)=S_N(A_K)^{-1}P_ND^{-1}q$ for $q\in\ell^2_1$.
\end{definition}

\begin{theorem}[Equations for the classical weighted finite solution]
\label{thm:volterra-weighted-finite-solution}\uses{def:volterra-weighted-finite-solution,def:volterra-weighted-solution}
\lean{MvPolynomial.volterraWeightedFiniteSolution_equations}\leanok
At every invertible cutoff, $u_{K,N}(q)$ has zero coefficients at indices $i\ge N$, and $(W_Ku_{K,N}(q))_i=q_i$ for $i<N$.
\end{theorem}
\begin{proof}
\leanok
\end{proof}
\end{supporting}


\begin{theorem}[Weighted polynomial first-kind Volterra solves]
\label{thm:volterra-weighted-solves}\uses{thm:inverse-residual-error,thm:volterra-nonvanishing-sections,thm:volterra-weighted-finite-solution}
\lean{MvPolynomial.volterraWeightedSolution_spec,MvPolynomial.existsUnique_volterraWeightedL2_solution,MvPolynomial.norm_volterraWeightedSolution_error_le,MvPolynomial.volterraWeightedFiniteSolution_filter_neBot,MvPolynomial.tendsto_volterraWeightedFiniteSolution,MvPolynomial.tendsto_volterraWeightedFiniteSolution_synthesis,MvPolynomial.volterraWeightedL2_finsupp,MvPolynomial.eventually_volterraWeightedFiniteSolution_equations,MvPolynomial.volterraProriolCoefficients_equations_iff,MvPolynomial.norm_volterraWeightedSolution_synthesis_error_le}\leanok
Write $D:\ell^2\to\ell^2_1$ for the isometric integration diagonal and
$W_K=DA_K$ for the endpoint-factor-free weighted coefficient operator.
For a real polynomial kernel with $K(x,x)\ne0$ on $[0,1]$ and any
$q\in\ell^2_1$, the equation $W_Ku=q$ has the unique solution
$u=A_K^{-1}D^{-1}q$. Its finite solves
\[
 u_N=(I-P_N+P_NA_KP_N)^{-1}P_ND^{-1}q
\]
are defined for every cutoff $N$ at which $I-P_N+P_NA_KP_N$ is invertible. These cutoffs include every sufficiently large $N$. As $N\to\infty$ through them, the finite solves converge to $u$ in raw $\ell^2$, and their synthesized input functions converge
in $L^2(4(1-x)\,dx)$ on $[0,1]$.
Whenever the section is invertible, $u_N$ has zero tail beyond $N$ and
$(W_Ku_N)_i=q_i$ for $i<N$; this holds eventually.
For every finite coefficient vector $c$, $W_Kc$ agrees coefficientwise
with the exact reflected Proriol assembly, so the retained finite assembly
equations are precisely these weighted equations.
For arbitrary raw coefficients $v$ the full weighted residual gives
\[
 \|v-u\|_{\ell^2}\le
 \|A_K^{-1}\|\,\|W_Kv-q\|_{\ell^2_1}.
\]
Multiplying this bound by the norm of raw Jacobi synthesis bounds the
weighted physical input error. The residual includes unresolved coordinates.
The output convention retains reflection and removes the endpoint factor,
as in the exact assembly theorem.
\end{theorem}
\begin{proof}
\leanok
\end{proof}

\rvnote{thm:volterra-weighted-solves}{
The residual includes unresolved coordinates; no computable bound on the inverse norm is asserted here.
}


\begin{supporting}[Reference details: Quantitative perturbation and stable operator families and related results]
\begin{theorem}[Quantitative perturbation and stable operator families]
\label{thm:norm-operator-perturbation}\uses{thm:inverse-residual-error}
\lean{ContinuousLinearMap.isUnit_of_inverse_norm_mul_sub_lt_one,ContinuousLinearMap.norm_inverse_le_of_inverse_norm_mul_sub,ContinuousLinearMap.norm_inverse_solution_sub_le,ContinuousLinearMap.eventually_stable_of_norm_sub_tendsto_zero,ContinuousLinearMap.tendsto_inverse_apply_of_eventually_stable}\leanok
Let $A,B$ be bounded endomorphisms of a real or complex Banach space.
If $A$ is invertible and $\|A^{-1}\|\|A-B\|\le r<1$, then $B$ is invertible and
$\|B^{-1}\|\le\|A^{-1}\|/(1-r)$.
For invertible $B$,
\[
 \|B^{-1}g-u\|\le\|B^{-1}\|
 (\|g-Au\|+\|A-B\|\|u\|).
\]
If $A_n$ are eventually invertible with inverse norms bounded by $C$ and
$\|A_n-B_n\|\to0$, then $B_n$ are eventually invertible with inverse norms
bounded by $2C$. For any eventually stable family $B_n$, consistency
$g_n-B_nu\to0$ implies $B_n^{-1}g_n\to u$.
\end{theorem}
\begin{proof}
\leanok
\end{proof}

\begin{theorem}[Approximate double sections with varying data]
\label{thm:approximate-coercive-sections}\uses{thm:norm-operator-perturbation,thm:coercive-sections,thm:stable-compact-perturbation,thm:l2-finite-sections}
\lean{lp.norm_coerciveFiniteSection_sub_le,lp.tendsto_approximateCoerciveFiniteSection,lp.eventually_stable_approximateCoerciveFiniteSection,lp.tendsto_approximateCoerciveFiniteSection_solution,lp.norm_approximateCoerciveFiniteSection_error_le}\leanok
For the initial coordinate projections on real or complex $\ell^2$, write
$S_N(A)=I-P_N+P_NAP_N$. Then
$\|S_N(A)-S_N(B)\|\le\|A-B\|$.
If $B_N\to A$ in operator norm, then $S_N(B_N)u\to Au$ for every $u$.
An eventual inverse bound $C$ for $S_N(A)$ transfers to an eventual bound
$2C$ for $S_N(B_N)$. Under that stability hypothesis and $g_N\to Au$,
$S_N(B_N)^{-1}P_Ng_N\to u$.
At each invertible approximate section,
\[
 \|S_N(B)^{-1}P_Ng-u\|
 \le\|S_N(B)^{-1}\|
 \big((1+\|A\|)\|P_Nu-u\|+\|g-Au\|+\|A-B\|\|u\|\big).
\]
\end{theorem}
\begin{proof}
\leanok
\end{proof}

\begin{theorem}[Coefficient majorants for first-kind operator approximation]
\label{thm:volterra-kernel-majorant}\uses{thm:volterra-ordered-algebra}
\lean{MvPolynomial.volterraKernelMajorant,MvPolynomial.norm_volterraPreconditionedL2_le_majorant,MvPolynomial.tendsto_volterraPreconditionedL2_of_majorant,MvPolynomial.summable_volterraPreconditionedL2_of_majorant,MvPolynomial.tendsto_volterraPreconditionedL2_sum_of_majorant}\leanok
Write $H,J$ for the reflected and raw Jacobi factors and
$K(x,y)=\sum k_{ij}x^iy^j$. Define
\[
 m(K)=\sum_{i,j}|k_{ij}|\|H\|^i\|J\|^j.
\]
Then $\|A_K\|\le m(K)$. In particular, $m(L_N-K)\to0$ implies
$A_{L_N}\to A_K$ in operator norm.
For polynomial increments $p_n$, if $\sum_n m(p_n)<\infty$,
the operator series $\sum_n A_{p_n}$ converges, and
$A_{\sum_{n<N}p_n}$ converges in norm to its sum.

\end{theorem}
\begin{proof}
\leanok
\end{proof}

\rvnote{thm:volterra-kernel-majorant}{
Theorem~\ref{thm:continuous-volterra-raw-bridge} identifies the coefficient limit with the integral operator under its physical and operator-norm convergence hypotheses.
}
\end{supporting}


\begin{theorem}[First-kind solves with approximate kernels and weighted data]
\label{thm:volterra-approximate-weighted-solves}\uses{thm:volterra-weighted-solves,thm:approximate-coercive-sections,thm:invertible-index-continuity}
\lean{MvPolynomial.exists_eventually_stable_volterraPreconditionedL2,MvPolynomial.isUnit_volterraPreconditionedL2_of_near,MvPolynomial.eventually_stable_approximateVolterraWeightedSection,MvPolynomial.approximateVolterraWeightedSolution_filter_neBot,MvPolynomial.tendsto_approximateVolterraWeightedSolution,MvPolynomial.volterraWeightedOperatorLimit_filter_neBot,MvPolynomial.tendsto_volterraWeightedSolution_of_operator_limit,MvPolynomial.norm_approximateVolterraWeightedSolution_error_le,MvPolynomial.tendsto_approximateVolterraWeightedSolution_synthesis,MvPolynomial.norm_volterraWeightedSolution_sub_le}\leanok
Let $K$ be a real polynomial with $K(x,x)\ne0$ on $[0,1]$.
The sections $S_N(A_K)$ have a uniform eventual inverse bound, for either diagonal sign.
If polynomial kernels $L_N$ satisfy $A_{L_N}\to A_K$ in operator norm and
$q_N\to r$ in $\ell^2_1$, their finite weighted solves are eventually stable and
\[
 u_N=S_N(A_{L_N})^{-1}P_ND^{-1}q_N
 \longrightarrow u=A_K^{-1}D^{-1}r .
\]
Here the finite solves are indexed by the invertible sections, a domain containing every sufficiently large cutoff, and convergence uses its induced nontrivial filter of tails. The synthesized inputs also converge along this filter in the weighted physical Jacobi $L^2$ space.
For any polynomial $L$, weighted data $q,r$, and invertible section $S_N(A_L)$,
\[
 \|u_N-u\|\le\|S_N(A_L)^{-1}\|
 \big((1+\|A_K\|)\|P_Nu-u\|+\|q-r\|_{\ell^2_1}
       +\|A_K-A_L\|\|u\|\big),
 \qquad u_N=S_N(A_L)^{-1}P_ND^{-1}q .
\]
The analogous full-solve bound omits the truncation term and uses $\|A_L^{-1}\|$,
assuming $A_L$ invertible. Small relative kernel error
$\|A_K^{-1}\|\|A_K-A_L\|<1$ suffices for this invertibility.
More generally, if $A_{L_N}\to A$ in norm for an arbitrary invertible bounded
operator $A$, full weighted solves converge to $A^{-1}D^{-1}r$ through the indices where $A_{L_N}$ is invertible. This domain also contains a tail and has a nontrivial induced filter.
This last assertion neither identifies a nonpolynomial integral kernel nor
asserts finite-section stability for an arbitrary invertible $A$.
Reference: \cite[Section 2.7, Theorem 2.5 and Lemma 2.5]{gutleb-thesis}.
\end{theorem}
\begin{proof}
\leanok
\end{proof}










\begin{supporting}[Reference details: Continuous-kernel identification of raw operator limits and related results]
\begin{theorem}[Continuous-kernel identification of raw operator limits]
\label{thm:continuous-volterra-raw-bridge}\uses{thm:continuous-jacobi-volterra,thm:volterra-raw-l2-agreement,thm:volterra-kernel-majorant}
\lean{Real.tendsto_polynomial_jacobiVolterraL2Operator,Real.jacobiVolterraL2Operator_raw_limit,Real.jacobiVolterraL2Operator_raw_series}\leanok
Let $K$ be continuous on $\Delta$, let polynomial kernels $p_n$ converge
uniformly to $K$, and suppose their raw preconditioned operators $A_{p_n}$
converge in operator norm to $A$. Then $B_{p_n}\to B_K$ in operator norm and
\[
 S D A=B_K S,
\]
where $S$ is raw shifted-Jacobi synthesis and $D$ the integration diagonal.
In particular, for polynomial increments $q_n$ with $\sum_n m(q_n)<\infty$
and partial sums converging uniformly to $K$, this identity holds with
$A=\sum_n A_{q_n}$.

\end{theorem}
\begin{proof}
\leanok
\end{proof}

\rvnote{thm:continuous-volterra-raw-bridge}{
Uniform physical convergence and preconditioned norm convergence are separate hypotheses; neither is silently substituted for the other.
}

\begin{theorem}[Real Laurent Toeplitz norm bound]
\label{thm:real-toeplitz-norm}\uses{thm:toeplitz-real-symbol}
\lean{lp.norm_laurentToeplitz_le_circle}\leanok
For a real Laurent polynomial $p$ and $M\ge0$, a bound $|p(e^{it})|\le M$
on the circle gives $\|T_p\|\le M$ on real $\ell^2$.
\end{theorem}
\begin{proof}
\leanok
\end{proof}

\begin{theorem}[Toeplitz decomposition of Volterra operator limits]
\label{thm:volterra-toeplitz-limit}\uses{thm:volterra-conditional-toeplitz-sections,thm:real-toeplitz-norm}
\lean{MvPolynomial.norm_volterraDiagonalToeplitz_le,MvPolynomial.norm_volterraDiagonalToeplitz_sub_le,MvPolynomial.exists_volterraDiagonalToeplitz_limit,MvPolynomial.exists_coercive_volterraToeplitz_limit}\leanok
The norm of the Toeplitz part of a polynomial Volterra operator is bounded
by a uniform bound for its diagonal $p(x,x)$ on $[0,1]$; the same holds for
differences. Hence uniform convergence $p_n(x,x)\to K(x,x)$ gives a norm limit
$T$ of the Toeplitz parts. If also $A_{p_n}\to A$ in norm and
$K(x,x)\ge a>0$, then $A-T$ is compact and
\[
 (a/2)\|u\|^2\le\langle u,Tu\rangle.
\]
\end{theorem}
\begin{proof}
\leanok
\end{proof}

\begin{theorem}[Stable Volterra limits with either diagonal sign]
\label{thm:volterra-limit-stability}\uses{thm:volterra-nonvanishing-sections,thm:volterra-toeplitz-limit}
\lean{MvPolynomial.isUnit_and_stable_volterra_limit}\leanok
Suppose $K(x,x)$ is continuous and nonzero on $[0,1]$, polynomial diagonals
converge uniformly to it, $A_{p_n}\to A$ in operator norm, and $A$ is injective.
Then $A$ is invertible and there is $B\ge0$ such that, eventually,
$S_N(A)=I-P_N+P_NAP_N$ is invertible and $\|S_N(A)^{-1}\|\le B$.
\end{theorem}
\begin{proof}
\leanok
\end{proof}
\end{supporting}


\begin{theorem}[Nonpolynomial first-kind Volterra convergence]
\label{thm:continuous-volterra-solves}\uses{thm:continuous-volterra-raw-bridge,thm:volterra-limit-stability,thm:approximate-coercive-sections}
\lean{Real.isUnit_and_stable_continuousVolterra,Real.tendsto_continuousVolterra_section_solution,Real.isUnit_and_stable_continuousVolterra_series}\leanok
Let $K$ be continuous on $\Delta$, have absolutely continuous $y$-slices,
nonzero diagonal, and satisfy $|\partial_yK(x,y)|\le C|K(x,x)|$ on $\Delta$.
Suppose polynomial kernels $p_n$ converge uniformly to $K$ and $A_{p_n}\to A$
in operator norm. Then $A$ is invertible and its leading sections have a
uniform eventual inverse bound. For $g_N\to g$ in raw $\ell^2$,
\[
 S_N(A_{p_N})^{-1}P_Ng_N\longrightarrow A^{-1}g.
\]
For polynomial increments $q_n$ whose partial sums converge uniformly to $K$,
summability of $m(q_n)$ suffices for invertibility and uniform eventual section
stability of $A=\sum_n A_{q_n}$.

\end{theorem}
\begin{proof}
\leanok
\end{proof}

\rvnote{thm:continuous-volterra-solves}{
These hypotheses give a sufficient nonpolynomial class; they do not assert preconditioned norm convergence for arbitrary continuous kernels.
}


\chapter{Series kernels and finite-data certificates}\label{chap:volterra-series}
\rvsummary{Joint kernel/operator approximation and the exponential-kernel example.}
Coefficient majorants bound both the kernel approximation error and the
operator error. For the exponential kernel, these bounds give a residual
certificate from finite data.

\begin{supporting}[Reference details: Joint weights for kernel monomials and related results]
\begin{theorem}[Joint weights for kernel monomials]
\label{thm:volterra-series-terms}\uses{thm:volterra-ordered-algebra}
\lean{MvPolynomial.volterraSeriesTerm,MvPolynomial.eval_volterraSeriesTerm,MvPolynomial.volterraSeriesWeight,MvPolynomial.volterraSeriesWeight_nonneg,MvPolynomial.norm_eval_volterraSeriesTerm_le,MvPolynomial.volterraPreconditionedL2_seriesTerm,MvPolynomial.norm_volterraPreconditionedL2_seriesTerm_le,MvPolynomial.norm_volterraSeriesTerm_derivative_le}\leanok
For a monomial $p(x,y)=cx^iy^j$, define
\[
 w_{R,S}(i,j,c)=|c|R^iS^j(j+1).
\]
This weight is nonnegative for $R,S\ge0$ and bounds $|p(x,y)|$ when
$|x|\le R$, $|y|\le S$. If $S\ge1$, it also bounds
$|cx^i j y^{j-1}|$, with the derivative zero when $j=0$.
Its ordered operator is $cH^iJ^j$. For
$R\ge\|H\|$, $S\ge\|J\|$, its operator norm is at most $w_{R,S}(i,j,c)$.
\end{theorem}
\begin{proof}
\leanok
\end{proof}

\begin{definition}[Closed rectangles for kernel series]
\label{def:volterra-series-box}
\lean{MvPolynomial.volterraSeriesBox}\leanok
For real $R,S$, put
$B_{R,S}=\{(x,y)\in\mathbb R^2:|x|\le R,\ |y|\le S\}$.
This set is empty if either radius is negative.
\end{definition}

\begin{theorem}[Closedness of the convergence rectangle]
\label{thm:volterra-series-box}\uses{def:volterra-series-box}
\lean{MvPolynomial.isClosed_volterraSeriesBox}\leanok
Every $B_{R,S}$ is closed.
\end{theorem}
\begin{proof}
\leanok
\end{proof}

\begin{definition}[Convergent physical kernel-series input]
\label{def:volterra-series-kernel-input}
\lean{MvPolynomial.VolterraSeriesKernelInput}\leanok
Fix exponent sequences $i_n,j_n\in\mathbb N$ and a point $z=(x,y)\in\mathbb R^2$.
A kernel-series input at $z$ is a real coefficient sequence $c_n$ for which
$\sum_n |c_nx^{i_n}y^{j_n}|<\infty$.
Convergence is required at the selected point; no common rectangle is required
to define the series at $z$.
\end{definition}

\begin{theorem}[Kernel-series inputs from a coefficient majorant]
\label{thm:volterra-series-kernel-input}\uses{def:volterra-series-kernel-input,def:volterra-series-box,thm:volterra-series-terms,def:volterra-triangle}
\lean{MvPolynomial.VolterraSeriesKernelInput.ofWeight,MvPolynomial.VolterraSeriesKernelInput.ofTriangle}\leanok
If $R,S\ge0$ and $\sum_n w_{R,S}(i_n,j_n,c_n)<\infty$, the scalar series $\sum_n c_n x^{i_n}y^{j_n}$ converges absolutely at every $(x,y)\in B_{R,S}$.
If $R,S\ge1$, this convergence holds at every point of the physical triangle $\Delta$.
\end{theorem}
\begin{proof}
\leanok
\end{proof}

\begin{definition}[Classical physical kernel series]
\label{def:volterra-series-kernel}\uses{def:volterra-series-kernel-input}
\lean{MvPolynomial.volterraSeriesKernel}\leanok
For a kernel-series input at $z=(x,y)$, define
$K(z)=\sum_{n\ge0}c_nx^{i_n}y^{j_n}$, the sum of its absolutely convergent series.
\end{definition}

\begin{theorem}[Value of the physical kernel series]
\label{thm:volterra-series-kernel-value}\uses{def:volterra-series-kernel}
\lean{MvPolynomial.hasSum_volterraSeriesKernel}\leanok
At every point $(x,y)$ where the monomial series is absolutely convergent,
its kernel value is its sum:
\[
 \sum_{n\ge0}c_n x^{i_n}y^{j_n}=K(x,y).
\]
\end{theorem}
\begin{proof}
\leanok
\end{proof}

\begin{theorem}[Continuity on the closed convergence rectangle]
\label{thm:volterra-series-kernel-continuity}\uses{thm:volterra-series-kernel-input,def:volterra-series-kernel}
\lean{MvPolynomial.continuous_volterraSeriesKernel_box}\leanok
If $R,S\ge0$ and the joint weights are summable, the function
$K:B_{R,S}\to\mathbb R$ is continuous.
\end{theorem}
\begin{proof}
\leanok
\end{proof}

\begin{definition}[Continuous extensions of a kernel series]
\label{def:volterra-series-kernel-extension}\uses{thm:volterra-series-kernel-input,def:volterra-series-kernel}
\lean{MvPolynomial.VolterraSeriesKernelExtension}\leanok
For nonnegative radii and a summable joint majorant, a kernel-series extension
is a continuous function $F:\mathbb R^2\to\mathbb R$ such that
$F(z)=K(z)$ for every $z\in B_{R,S}$.
Outside the rectangle, $F$ is an extension and is not asserted to be a series value.
\end{definition}

\begin{theorem}[Existence of continuous kernel-series extensions]
\label{thm:volterra-series-kernel-extension}\uses{def:volterra-series-kernel-extension,thm:volterra-series-kernel-continuity,thm:volterra-series-box}
\lean{MvPolynomial.nonempty_volterraSeriesKernelExtension}\leanok
Every kernel series with nonnegative radii and a summable joint majorant
admits a continuous extension to $\mathbb R^2$ agreeing on $B_{R,S}$.
\end{theorem}
\begin{proof}
\leanok
\end{proof}

\begin{definition}[Convergent Volterra operator-series data]
\label{def:volterra-operator-series-input}\uses{thm:volterra-series-terms}
\lean{MvPolynomial.VolterraSeriesOperatorInput}\leanok
For exponent sequences $i_n,j_n\in\mathbb N$, an operator-series input is
a real coefficient sequence $c_n$ for which the ordered operator family
$A_{c_nx^{i_n}y^{j_n}}$ is unconditionally summable in the Banach space of bounded operators
on real $\ell^2$.
\end{definition}

\begin{theorem}[Operator-series input from a coefficient majorant]
\label{thm:volterra-operator-series-input}\uses{def:volterra-operator-series-input}
\lean{MvPolynomial.VolterraSeriesOperatorInput.ofWeight}\leanok
If $R\ge\|H\|$, $S\ge\|J\|$, and
$(w_{R,S}(i_n,j_n,c_n))_{n\ge0}$ is summable, then the ordered operator series $\sum_n c_nH^{i_n}J^{j_n}$ is summable in the Banach space of bounded operators.
\end{theorem}
\begin{proof}
\leanok
\end{proof}

\begin{definition}[Sum of an ordered Volterra operator series]
\label{def:volterra-series-operator}\uses{def:volterra-operator-series-input}
\lean{MvPolynomial.volterraSeriesOperator}\leanok
For a convergent operator-series input, define the unconditional sum
$A=\sum_{n\ge0}A_{c_nx^{i_n}y^{j_n}}$ in the space of bounded operators
on real $\ell^2$.
\end{definition}

\begin{theorem}[Value of a convergent Volterra operator series]
\label{thm:volterra-series-operator}\uses{def:volterra-series-operator}
\lean{MvPolynomial.hasSum_volterraSeriesOperator}\leanok
On the domain where $\sum_n c_nH^{i_n}J^{j_n}$ is unconditionally summable in the Banach space of bounded operators, the ordered Volterra series operator $A$ is its sum:
\[
 \sum_{n\ge0}c_nH^{i_n}J^{j_n}=A.
\]
\end{theorem}
\begin{proof}
\leanok
\end{proof}


\begin{definition}[Summable Volterra tail data]
\label{def:volterra-tail-input}\uses{thm:volterra-series-terms}
\lean{MvPolynomial.VolterraSeriesTailInput}\leanok
For exponent sequences $i_n,j_n\in\mathbb N$ and real radii $R,S$,
a Volterra tail input is a real coefficient sequence $c_n$ for which
the real family $w_n=w_{R,S}(i_n,j_n,c_n)$ is absolutely summable.
No positivity of the radii is needed to specify this convergence domain.
\end{definition}

\begin{theorem}[Convergence after a finite cutoff]
\label{thm:volterra-tail-input}\uses{def:volterra-tail-input}
\lean{MvPolynomial.VolterraSeriesTailInput.summable_tail}\leanok
Let $w_n=w_{R,S}(i_n,j_n,c_n)$ be the joint coefficient weights.
If $(w_n)_{n\ge0}$ is summable, then for every $N\ge0$ the shifted family
$(w_{n+N})_{n\ge0}$ is summable.
\end{theorem}
\begin{proof}
\leanok
\end{proof}

\begin{definition}[Volterra coefficient tail]
\label{def:volterra-series-tail}\uses{def:volterra-tail-input}
\lean{MvPolynomial.volterraSeriesTail}\leanok
For a Volterra tail input, define the convergent coefficient tail by
$E_N=\sum_{n\ge0}w_{n+N}$.
\end{definition}

\begin{theorem}[Value and convergence of Volterra tails]
\label{thm:volterra-series-tail}\uses{def:volterra-series-tail,thm:volterra-tail-input}
\lean{MvPolynomial.hasSum_volterraSeriesTail,MvPolynomial.volterraSeriesTail_nonneg,MvPolynomial.tendsto_volterraSeriesTail}\leanok
The shifted family $(w_{n+N})_{n\ge0}$ has sum $E_N$, and $E_N\to0$.
If $R,S\ge1$, then $E_N\ge0$ for every $N$.
\end{theorem}
\begin{proof}
\leanok
\end{proof}
\end{supporting}



\begin{theorem}[Joint physical and operator approximation from coefficients]
\label{thm:volterra-series-approximation}\uses{thm:volterra-series-kernel-input,def:volterra-series-kernel,thm:volterra-series-tail,thm:volterra-operator-series-input,def:volterra-series-operator}
\lean{MvPolynomial.volterraSeriesPartial,MvPolynomial.norm_volterraSeriesPartial_sub_kernel_le,MvPolynomial.continuous_volterraSeriesKernel,MvPolynomial.norm_volterraSeriesPartial_sub_operator_le,MvPolynomial.tendsto_volterraSeriesPartial_operator}\leanok
Let $i_n,j_n\in\mathbb N$, $c_n\in\mathbb R$, and assume $R,S\ge1$ and
$\sum_nw_n<\infty$, where $w_n=w_{R,S}(i_n,j_n,c_n)$.
Put $p_n(x,y)=c_nx^{i_n}y^{j_n}$ and $P_N=\sum_{n<N}p_n$.
On $B_{R,S}$ use the classical kernel sum $K$ from
Definition~\ref{def:volterra-series-kernel}.
For the coefficient tail $E_N$ from Definition~\ref{def:volterra-series-tail},
we have $E_N\ge0$, $E_N\to0$, and $K$ is continuous on $\Delta$, with
$|P_N-K|\le E_N$ uniformly there.
If additionally $R\ge\|H\|$, $S\ge\|J\|$, the operator-series input is
convergent. Writing $A=\sum_n A_{p_n}$, we have
$\|A_{P_N}-A\|\le E_N$ and $A_{P_N}\to A$ in operator norm.
The enumeration is arbitrary; a bivariate power series is obtained by enumerating
its exponent pairs. No injectivity of the enumeration is needed.
\end{theorem}
\begin{proof}
\leanok
\end{proof}

\begin{theorem}[Slice regularity from weighted coefficients]
\label{thm:volterra-series-regularity}\uses{thm:volterra-series-kernel-extension}
\lean{MvPolynomial.hasDerivAt_volterraSeriesKernel,MvPolynomial.norm_deriv_volterraSeriesKernel_le,MvPolynomial.absolutelyContinuousOnInterval_volterraSeriesKernel,MvPolynomial.exists_volterraSeriesKernel_relative_derivative_bound}\leanok
Under the preceding summability condition, assume $R\ge1$ and $S>1$.
For $|x|\le R$, $|y|<S$, the kernel is differentiable in its inner variable and
\[
 \partial_yK(x,y)=\sum_n c_nx^{i_n}j_n y^{j_n-1}.
\]
The summand is zero when $j_n=0$; when $j_n>0$, its exponent is nonnegative.
On $\Delta$, $|\partial_yK|\le W:=\sum_n w_n$ and every slice on $[0,x]$
is absolutely continuous. These local derivative and row statements hold for
 every continuous extension agreeing on $B_{R,S}$; the derivative of the
 domain-valued kernel is independent of the chosen extension.
 If $K(x,x)\ne0$ throughout $[0,1]$, there exists
$C\in\mathbb R$ such that $|\partial_yK(x,y)|\le C|K(x,x)|$ on $\Delta$.
\end{theorem}
\begin{proof}
\leanok
\end{proof}

\begin{theorem}[First-kind solves from the joint coefficient criterion]
\label{thm:volterra-series-solves}\uses{thm:continuous-volterra-solves,thm:volterra-series-regularity,thm:volterra-series-approximation}
\lean{MvPolynomial.isUnit_and_stable_volterraSeriesOperator,MvPolynomial.tendsto_volterraSeriesSection_solution,MvPolynomial.norm_volterraSeriesSection_error_le}\leanok
Suppose $R\ge\max(1,\|H\|)$, $S>1$, $S\ge\|J\|$, $\sum_nw_n<\infty$,
and $K(x,x)\ne0$ on $[0,1]$. Then $A$ is invertible and its padded leading
sections $S_N(A)=I-P_N^{\ell^2}+P_N^{\ell^2}AP_N^{\ell^2}$ have a uniform
eventual inverse bound. For raw data $g_N\to g$,
\[
 S_N(A_{P_N})^{-1}P_N^{\ell^2}g_N\to A^{-1}g.
\]
Here the kernel truncation $P_N$ and coordinate projection $P_N^{\ell^2}$
are different objects. More generally, at any invertible section
$B=S_N(A_{P_L})$, for arbitrary $u,g\in\ell^2$,
\[
 \|B^{-1}P_N^{\ell^2}g-u\|\le\|B^{-1}\|
 \big((1+\|A\|)\|P_N^{\ell^2}u-u\|+\|g-Au\|+E_L\|u\|\big).
\]
 The error estimate itself only requires the summability,
operator-radius bounds and invertibility of that section.
\end{theorem}
\begin{proof}
\leanok
\end{proof}

\rvnote{thm:volterra-series-solves}{The inverse norm and solution tail remain explicit; no numerical certificate for them is asserted.}


\begin{supporting}[Reference details: The exponential kernel satisfies the coefficient criterion]
\begin{theorem}[The exponential kernel satisfies the coefficient criterion]
\label{thm:volterra-exponential-series}\uses{thm:volterra-series-kernel-input,def:volterra-series-kernel}
\lean{MvPolynomial.volterraExponentialCoefficient,MvPolynomial.summable_volterraExponentialWeight,MvPolynomial.VolterraSeriesKernelInput.ofExponential,MvPolynomial.volterraSeriesKernel_exponential}\leanok
For any real $\lambda$, enumerate $(i_n,j_n)=\operatorname{unpair}(n)$ and set
\[
 c_n=\frac{\lambda^{i_n}(-\lambda)^{j_n}}{i_n!j_n!}.
\]
The series $\sum_n c_n x^{i_n}y^{j_n}$ converges absolutely for every $(x,y)\in\R^2$.
The classical kernel value equals $e^{\lambda(x-y)}$ for all real $x,y$.
For every $R,S\ge0$, its joint weights are summable.
The pairing enumerates all exponent pairs.
\end{theorem}
\begin{proof}
\leanok
\end{proof}
\end{supporting}


\begin{theorem}[End-to-end exponential first-kind solves]
\label{thm:volterra-exponential-solves}\uses{thm:volterra-series-solves,thm:volterra-exponential-series}
\lean{MvPolynomial.volterraExponentialPartial,MvPolynomial.volterraExponentialOperator,MvPolynomial.norm_volterraExponentialPartial_sub_le,MvPolynomial.isUnit_and_stable_volterraExponentialOperator,MvPolynomial.tendsto_volterraExponentialSection_solution,MvPolynomial.norm_volterraExponentialKernel_le,MvPolynomial.volterraExponentialOperator_raw_synthesis,MvPolynomial.jacobiVolterraL2Operator_exponential_solution}\leanok
Let $K_\lambda(x,y)=e^{\lambda(x-y)}$, for any $\lambda\in\mathbb R$, and let
$P_N$ be its paired polynomial series truncations and $A_\lambda$ the ordered
operator series. The physical bound is $|K_\lambda|\le e^{|\lambda|}$ on
$\Delta$; its truncation error is bounded by the coefficient tail for every
$R,S\ge1$. The diagonal is identically one.
The operator $A_\lambda$ is invertible, its leading sections have a uniform
eventual inverse bound, and for $g_N\to g$ in raw $\ell^2$,
\[
 S_N(A_{P_N})^{-1}P_N^{\ell^2}g_N\to A_\lambda^{-1}g.
\]
The exact integral realization satisfies $SD A_\lambda=B_{K_\lambda}S$.
Consequently, the limiting coefficient solution satisfies the genuine weighted
integral equation
\[
 B_{K_\lambda}S(A_\lambda^{-1}g)=SDg.
\]
Reflection and endpoint averaging are those of $B_K$, and $S$ is synthesis from raw coefficients.
\end{theorem}
\begin{proof}
\leanok
\end{proof}


\begin{theorem}[Ordered exponential factorisation and full inverse bound]
\label{thm:volterra-exponential-factor}\uses{thm:volterra-exponential-solves}
\lean{MvPolynomial.volterraExponentialOperator_eq,MvPolynomial.inverse_volterraExponentialOperator,MvPolynomial.norm_inverse_volterraExponentialOperator_le,MvPolynomial.norm_sub_volterraExponential_solution_le}\leanok
For every real $\lambda$, the raw coefficient operator factors in the stated order:
\[
 A_\lambda=e^{\lambda H}e^{-\lambda J},\qquad
 A_\lambda^{-1}=e^{\lambda J}e^{-\lambda H}.
\]
No commutation of $H$ and $J$ is assumed. In particular, with
$C_\lambda=\exp(|\lambda|(\|H\|+\|J\|))$,
\[
 \|A_\lambda^{-1}\|\le C_\lambda,\qquad
 \|v-A_\lambda^{-1}g\|\le C_\lambda\|g-A_\lambda v\|
 \quad(g,v\in\ell^2).
\]

\end{theorem}
\begin{proof}
\leanok
\end{proof}

\rvnote{thm:volterra-exponential-factor}{
This is a bound for the full inverse; it does not assert the same constant for finite-section inverses.
}

\begin{theorem}[Explicit exponential truncation bounds]
\label{thm:volterra-exponential-tail}\uses{thm:volterra-exponential-solves}
\lean{MvPolynomial.tsum_volterraExponentialWeight_le,MvPolynomial.volterraExponentialTail_le_geometric,MvPolynomial.norm_volterraExponentialPartial_sub_le_geometric,MvPolynomial.norm_volterraExponentialPartial_operator_sub_le_geometric}\leanok
For $R,S\ge0$, the total joint coefficient weight of the exponential series satisfies
\[
 \sum_n w_{R,S}(i_n,j_n,c_n)\le e^{|\lambda|R+2|\lambda|S}.
\]
For $0<q\le1$ and natural numbers $m,N$ with $m^2\le N$, its tail satisfies
\[
 E_N(R,S)\le q^m\exp\!\left(\frac{|\lambda|R+2|\lambda|S}{q}\right).
\]
Consequently, on the physical triangle,
\[
 |P_N(x,y)-e^{\lambda(x-y)}|\le q^m e^{3|\lambda|/q}.
\]
For $R\ge\|H\|$, $S\ge\|J\|$, the raw operator error satisfies
\[
 \|A_{P_N}-A_\lambda\|\le
 q^m\exp\!\left(\frac{|\lambda|R+2|\lambda|S}{q}\right).
\]
For fixed $q<1$ these are geometric bounds in $m$; $N$ is the number of
paired terms, not a total-degree cutoff. Choosing $q=1/2$ gives entirely
explicit bounds without an infinite tail sum.
\end{theorem}
\begin{proof}
\leanok
\end{proof}


\begin{supporting}[Reference details: Numerical bounds for the raw Jacobi factors and related results]
\begin{theorem}[Numerical bounds for the raw Jacobi factors]
\label{thm:raw-jacobi-norm-bounds}\uses{thm:raw-jacobi-l2}
\lean{Polynomial.norm_shiftedJacobiToeplitzPart_le,Polynomial.norm_shiftedJacobiOneZeroL2_le,Polynomial.norm_reflectedJacobiOneZeroL2_le}\leanok
For the constant tridiagonal part $T$, raw shifted Jacobi multiplication $J$,
and reflected preconditioned factor $H$,
\[
 \|T\|\le1,\qquad \|J\|\le2,\qquad \|H\|\le3.
\]

\end{theorem}
\begin{proof}
\leanok
\end{proof}

\rvnote{thm:raw-jacobi-norm-bounds}{
These are conservative bounds in the raw coefficient $\ell^2$ norm.
}

\begin{theorem}[Full residual from rectangular Proriol assembly]
\label{thm:volterra-finite-residual}\uses{thm:volterra-coefficient-bandwidth,thm:volterra-raw-finite-agreement}
\lean{lp.norm_finsuppL2_eq_sqrt_sum,MvPolynomial.volterraPreconditionedL2_finsupp_apply,MvPolynomial.volterraPreconditionedL2_finsupp_cutoff,MvPolynomial.volterraFiniteResidual_cutoff,MvPolynomial.norm_volterraFiniteResidual_eq,MvPolynomial.volterraPreconditionedL2_rectangular_mulVec}\leanok
Let $K$ be a polynomial kernel of total degree $d_K$ (zero for the zero kernel),
and let $c,d$ be finite
raw coefficient vectors supported below $N,M$, respectively. Write
$b_i=(\mathcal V_Kc)_i$ for the Proriol assembled coefficients before the
endpoint factor. The raw preconditioned action is
\[
 (A_Kc)_i=(-1)^i(i+1)b_i.
\]
It vanishes for $i\ge N+d_K$. Thus the residual $r=d-A_Kc$ vanishes for
$i\ge Q:=\max(M,N+d_K)$, and its full norm is exactly
\[
 \|r\|_2=\left(\sum_{i<Q}[d_i-(-1)^i(i+1)b_i]^2\right)^{1/2},\qquad
 \|c\|_2=\left(\sum_{i<N}c_i^2\right)^{1/2}.
\]
For input $v\in\R^N$, each of the first $N+d_K$ raw output rows is the
corresponding row of the rectangular Proriol section multiplied by
$(-1)^i(i+1)$. The residual sum includes rows beyond the retained square
system; those rows cannot generally be discarded.
\end{theorem}
\begin{proof}
\leanok
\end{proof}
\end{supporting}


\begin{theorem}[Finite-data certificate for the exponential first-kind solve]
\label{thm:volterra-exponential-certificate}\uses{thm:volterra-exponential-factor,thm:volterra-exponential-tail,thm:raw-jacobi-norm-bounds,thm:volterra-finite-residual}
\lean{MvPolynomial.norm_inverse_volterraExponentialOperator_le_numeric,MvPolynomial.norm_volterraExponentialPartial_operator_sub_le_numeric,MvPolynomial.norm_volterraExponential_solution_error_le,MvPolynomial.norm_volterraExponential_polynomial_certificate,MvPolynomial.norm_volterraExponential_finite_certificate}\leanok
For every real $\lambda$,
\[
 \|A_\lambda^{-1}\|\le e^{5|\lambda|}.
\]
For any raw vectors $g,d,v$, bounded operator $B$, and error bounds
$\|g-d\|\le\delta$, $\|B-A_\lambda\|\le\varepsilon$,
\[
 \|v-A_\lambda^{-1}g\|\le
 e^{5|\lambda|}\bigl(\delta+\|d-Bv\|+\varepsilon\|v\|\bigr).
\]
In particular, for the paired kernel truncation $P_L$, $0<q\le1$, and
$m^2\le L$, one may take
\[
 B=A_{P_L},\qquad \varepsilon=q^m e^{7|\lambda|/q}.
\]
For finite input coefficients $c$ supported below $N$ and finite data $d$
supported below $M$, put $Q=\max(M,N+\operatorname{totalDegree}P_L)$ and
$b_i=(\mathcal V_{P_L}c)_i$. Given $\|g-d\|\le\delta$, the certificate is
\[
 \|c-A_\lambda^{-1}g\|_2\le e^{5|\lambda|}
 \left(\delta+
 \sqrt{\sum_{i<Q}[d_i-(-1)^i(i+1)b_i]^2}
 +q^m e^{7|\lambda|/q}\sqrt{\sum_{i<N}c_i^2}\right).
\]
Finite vectors are embedded by zero extension. The candidate need not satisfy
its retained equations exactly.
\end{theorem}
\begin{proof}
\leanok
\end{proof}

\rvnote{thm:volterra-exponential-certificate}{The data-tail bound is supplied explicitly; no finite-section inverse bound or floating-point rounding model is assumed.}


Volterra operators through the triangle: see Theorems~\ref{thm:volterra-coefficient-exact}, \ref{thm:volterra-finite-approximation}, \ref{thm:volterra-l2-proriol}, \ref{thm:volterra-compact-sections}, \ref{thm:volterra-factorial-inverse}, \ref{thm:volterra-finite-section-solves}, \ref{thm:volterra-proriol-matrix-solves}, \ref{thm:volterra-first-kind-diagonal}, \ref{thm:volterra-first-kind-diagonal-solves}, \ref{thm:volterra-raw-l2-agreement}, \ref{thm:volterra-raw-injective-sections}, \ref{thm:volterra-nonvanishing-sections}, \ref{thm:volterra-series-solves}, \ref{thm:volterra-exponential-solves}, \ref{thm:volterra-exponential-factor}, \ref{thm:volterra-exponential-tail} and~\ref{thm:volterra-exponential-certificate}.


\chapter{Roadmap and open targets}\label{chap:not-yet-covered}
\rvsummary{Targets that are not yet fully formalised.}
The results in this chapter are targets for future work. Formalised subcases
appear in the relevant subject chapters.

\section{Rational modifications and infinite factorizations}\label{thm:rational-infinite-modification}
Extend polynomial measure conversion to rational modifications and
infinite-dimensional factorizations, with explicit operator domains,
invertibility assumptions and finite-section convergence statements.

Reference: \cite[Section 2]{gutleb-olver-slevinsky}.



\section{Proriol families on triangles and simplices}\label{thm:proriol-polynomials}
Develop Proriol families on triangles and their simplex generalizations, including
orthogonality, normalizations and structured polynomial operators.

\section{Dunkl--Xu disk polynomials}\label{thm:dunkl-xu-disk-polynomials}
Develop the Cartesian Koornwinder construction of Dunkl--Xu disk polynomials,
including orthogonality, normalization, coordinate-multiplication matrices and
weight raising/lowering. Record the basis ordering explicitly and distinguish
this construction from the Zernike disk bases when describing operator structure.

Reference: \cite{dunkl-xu}.

\section{Higher-dimensional measure modification}\label{thm:multivariate-measure-modification}
Build on the multivariate Cholesky conversion foundation to study the resulting
block structure and coordinate-multiplication operators. Distinguish sparsity of
modified Gram matrices from sparsity of their connection factors. Specific
results will receive separate statements and hypotheses as this theory develops.

\section{One-dimensional power-law integral operators}\label{thm:power-law-1d-operator}
Develop the weighted operator
\[
 (Q_\alpha f)(x)=\int_{-1}^1 |x-y|^\alpha f(y)\,dy
\]
on $f(y)=(1-y^2)^{\lambda-1/2}C_n^{(\lambda)}(y)$.
Integrability, the left/right decomposition, initial modes, recurrence and the
exact compatible coefficient operator are proved above. Continuous amplitudes now
have polynomial approximation with locally uniform output convergence, including
a common weight for two incompatible powers. Closed-interval stability now has
explicit combined-exponent hypotheses, with uniform mass bounds throughout
$\alpha>-1$, $\lambda>-1/2$, $\alpha+\lambda>-1/2$. These bounds give a common
Gegenbauer $L^2$ operator and strong finite coefficient-section convergence.
The common-space $L^2$ boundedness threshold is now classified exactly:
$\alpha+\lambda\ge-1/2$, including the critical line by a weighted Schur test.
Normalized coefficient-section transport now covers the full boundedness range,
using Chebyshev $T$ at $\lambda=0$ and Gegenbauer elsewhere. Sufficient two-weight bounds and their coefficient sections are now proved;
a sharp two-weight classification remains a later target. In the common-weight
space, self-adjointness now holds throughout the boundedness range. In the
strict range, compactness and operator-norm section convergence give eventually
stable, convergent second-kind solves under an explicit invertibility hypothesis.
Compactness is now classified sharply: the critical line is noncompact.
Small coupling gives stable, convergent finite solves at every resolution even
on that line, and full residuals give error certificates. Critical-line solves
under invertibility alone, sharp two-weight results and first-kind inversion
remain separate. Keep exact compatible
sparsity separate from these approximation statements.
This target concerns the integral operator, not equilibrium minimizers.

References: \cite[Section 2]{gutleb-carrillo-olver};
\cite[Sections 3.4--3.5]{gutleb-thesis} (operator construction only).
